% Canonical single-file TeX anthology for: Bernhard Riemann, selected papers
% Author/editor lane: Bernhard Riemann
% Generated by Build-ZenodoV3NormalizedCorpus.py.
% Multiple upstream TeX roots were flattened below in natural order.

\documentclass[12pt,a4paper]{book}
\usepackage[T1]{fontenc}
\usepackage[utf8]{inputenc}
\usepackage{lmodern}
\usepackage{textcomp}
\usepackage{amsmath,amssymb,amsthm}
\usepackage{mathtools}
\usepackage{geometry}
\usepackage{hyperref}
\usepackage{microtype}
\usepackage{graphicx}
\usepackage{array}
\usepackage{booktabs}
\usepackage{longtable}
\usepackage{enumitem}
\usepackage{xcolor}
\geometry{margin=1in}
\newtheorem{theorem}{Theorem}[chapter]
\newtheorem{lemma}[theorem]{Lemma}
\newtheorem{corollary}[theorem]{Corollary}
\newtheorem{proposition}[theorem]{Proposition}
\theoremstyle{definition}
\newtheorem{definition}[theorem]{Definition}
\theoremstyle{remark}
\newtheorem{remark}[theorem]{Remark}
\newtheorem*{scholium}{Scholium}
\newenvironment{note}{\par\medskip\noindent\textbf{Note.}\ }{\par\medskip}

\begin{document}
\title{Bernhard Riemann, selected papers}
\author{Bernhard Riemann}
\maketitle
\tableofcontents

\clearpage
\chapter*{Source segment 1: riemann pages 001 050}
\addcontentsline{toc}{chapter}{Source segment 1: riemann pages 001 050}
% ============================================
% Title page (Page 7)
% ============================================
\thispagestyle{empty}
\begin{center}
\vspace*{2cm}
{\Large\textsc{Bernhard Riemann's}}\\[0.5cm]
{\LARGE\textsc{Gesammelte}}\\[0.3cm]
{\LARGE\textsc{Mathematische Werke}}\\[0.5cm]
{\large\textsc{und}}\\[0.5cm]
{\Large\textsc{Wissenschaftlicher Nachlass.}}\\[1.5cm]
{\large\textsc{Herausgegeben}}\\[0.3cm]
{\large\textsc{unter Mitwirkung von Richard Dedekind}}\\[0.3cm]
{\large\textsc{von}}\\[0.3cm]
{\Large\textsc{Heinrich Weber.}}\\[1.5cm]
{\large\textsc{Zweite Auflage}}\\[0.5cm]
{\normalsize\textsc{Bearbeitet von}}\\[0.3cm]
{\large\textsc{Heinrich Weber,}}\\[0.3cm]
{\small\textsc{mit einem Bildniss Riemann's.}}\\[1.5cm]
{\large\textsc{Leipzig,}}\\[0.3cm]
{\large\textsc{Druck und Verlag von B.\ G.\ Teubner.}}\\[0.5cm]
{\large\textsc{1892.}}\\[0.3cm]
{\small Printed in Germany}
\end{center}
\clearpage

% ============================================
% Vorrede zur ersten Auflage (Pages 9-11)
% ============================================
\selectlanguage{ngerman}
\section*{Vorrede zur ersten Auflage.}
\addcontentsline{toc}{section}{Vorrede zur ersten Auflage}

Das Werk, welches hiermit in die Oeffentlichkeit tritt, ist die endliche Ausf\"uhrung eines seit lange geplanten Unternehmens. Bei der Bedeutung, welche die grossen Sch\"opfungen Riemann's f\"ur die Entwicklung der neueren Mathematik haben, geh\"oren die meisten der Riemann'schen Abhandlungen zu den unentbehrlichsten H\"ulfsmitteln des Mathematikers, und eine Sammlung seiner Werke d\"urfte daher einem allgemein gehegten Wunsche um so mehr entgegen kommen, als die meisten derselben im Buchhandel nicht oder nur schwer zu erhalten sind. Es kommt dazu die dringende Pflicht gegen die Wissenschaft, die im handschriftlichen Nachlass noch verborgenen Untersuchungen und Gedanken der Oeffentlichkeit nicht l\"anger vorzuenthalten.

Schon im Fr\"uhjahr 1872 war daher unter mehreren Freunden Riemann's der Plan zu einer solchen Sammlung entstanden und Clebsch hatte mit seiner ganzen Thatkraft die Leitung des Unternehmens in die Hand genommen und sich mit Dedekind vereinigt, in dessen Besitz nach Riemann's Wunsch der handschriftliche Nachlass nach des Verfassers Tod gekommen war, und der bereits mehrere Abhandlungen aus demselben herausgegeben hatte.

Durch den beklagenswerthen und unerwarteten Tod von Clebsch gerieth leider das Vorhaben ins Stocken und blieb l\"angere Zeit g\"anzlich liegen. Als mir im November 1874 Dedekind im Namen der Frau Professorin Riemann den Vorschlag machte, die Leitung der Herausgabe zu \"ubernehmen, bin ich nicht ohne schwere Bedenken darauf eingegangen. Denn obwohl ich von dem Umfang der damit verbundenen Arbeit damals noch keine richtige Vorstellung hatte, war ich mir der zu \"ubernehmenden Verantwortung wohl bewusst. Nur die Erw\"agung, dass im Falle meiner Weigerung die Ausf\"uhrung abermals auf lange Zeit hinausgeschoben zu werden, wenn nicht g\"anzlich zu scheitern drohte, half mir meine Bedenken \"uberwinden, und so entschloss ich mich, was an mir l\"age, zu thun, um das Unternehmen zu einem befriedigenden Abschluss zu bringen, da Dedekind mir die Versicherung gab, mich bei der Arbeit nach Kr\"aften zu unterst\"utzen, ein Versprechen, welches er treulich gehalten hat.

\clearpage

Die von Riemann selbst oder nach seinem Tode bereits ver\"offentlichten Arbeiten wurden revidirt, hin und wieder durch einen im Nachlass aufgefundenen Zusatz bereichert, und in kleinen Ungenauigkeiten verbessert, sonst aber in unver\"anderter Form aufgenommen. Nur die Abhandlung \"uber die Fl\"achen vom kleinsten Inhalt hat in Folge einer von K.\ Hattendorff auf meinen Wunsch ausgef\"uhrten Ueberarbeitung einige wesentlichere Aenderungen erfahren.

Von den im Nachlass enthaltenen Entw\"urfen fanden sich einige in fast druckfertiger Form vor, andere aber in einem so fragmentarischen Zustande, dass die Verkn\"upfung und Darstellung erhebliche Schwierigkeiten machte. Von der grossen Menge nur Formeln ohne Text enthaltender Papiere war wenig f\"ur den Druck zu verwerthen.

Besonders hervorzuheben ist unter den ersteren die Arbeit \"uber den R\"uckstand in der Leidener Flasche, welche Riemann schon im Anschluss an die Mittheilung in der G\"ottinger Naturforscher-Versammlung zur Publication vorbereitet hatte, ferner die in lateinischer Sprache geschriebene Beantwortung einer Preisfrage der Pariser Akademie \"uber isotherme Curven, welche besonders deshalb von hohem Interesse ist, weil darin Riemann's Untersuchungen \"uber die allgemeinen Eigenschaften der mehrfach ausgedehnten Mannigfaltigkeiten in den Grundz\"ugen niedergelegt sind und eine merkw\"urdige Verwendung finden. Die Darstellung in dieser Abhandlung ist eine \"ausserst knappe, und die Wege, auf denen die endlichen Resultate erhalten wurden, finden sich darin nur im Allgemeinen angedeutet. Von der Ausf\"uhrung einer beabsichtigten zweiten eingehenderen Darstellung des Gegenstandes wurde Riemann durch seinen Gesundheitszustand abgehalten. Dass ich im Stande bin, diese sch\"one Untersuchung in der letzten von Riemann herr\"uhrenden Redaction zum Abdruck zu bringen, verdanke ich der G\"ute des best\"andigen Secret\"ars der Pariser Akademie, Herrn Dumas, welcher auf ein Namens der G\"ottinger Gesellschaft der Wissenschaften von Herrn Wohler an ihn gerichtetes Ansuchen mit der dankenswerthesten Bereitwilligkeit mir das Originalmanuscript zur Verf\"ugung stellte.

Von Riemann's Untersuchungen \"uber lineare Differentialgleichungen mit algebraischen Coefficienten liegt der erste Theil in ziemlich druckfertiger Form von Riemann's Hand vor und war vermuthlich zu der Publication bestimmt, die in der Abhandlung \"uber Abel'sche Functionen angek\"undigt ist, aber nicht zur Ausf\"uhrung kam. Ein zweiter Theil, der die wahre Verallgemeinerung der Theorie der hypergeometrischen Reihen enth\"alt, fand sich nur im ersten Entw\"urfe vor, jedoch so, dass der Gedankengang vollst\"andig hergestellt werden konnte.

\clearpage

Ferner ist hier noch der in italienischer Sprache geschriebene Anfang zu einer Untersuchung \"uber die Darstellbarkeit des Quotienten zweier hypergeometrischer Reihen durch einen Kettenbruch zu erw\"ahnen, deren Bearbeitung H.\ A.\ Schwarz in G\"ottingen \"ubernommen hat, dem ich hierf\"ur sowie f\"ur manchen Rath in anderen Stellen hier meinen Dank ausspreche.

Obwohl die Vorlesungen Riemann's dem urspr\"unglichen Plane nach von dieser Sammlung ausgeschlossen sind, so habe ich mich doch zur Aufnahme zweier kleinerer, in sich abgeschlossener Untersuchungen \"uber die Convergenz der $p$-fach unendlichen Theta-Reihe und \"uber die Abel'schen Functionen f\"ur den Fall $p = 3$ entschlossen, bei deren Bearbeitung ein von G.\ Roch gef\"uhrtes Vorlesungsheft zu Grunde gelegt werden konnte, theils wegen des grossen Interesses, welches die Gegenst\"ande haben, theils weil eine zusammenh\"angende Ver\"offentlichung dieser Vorlesungen, wie es scheint, vorl\"aufig nicht in Aussicht steht.

Ich erw\"ahne hier noch die den Anhang bildenden naturphilosophischen Fragmente, welche wenigstens eine ungef\"ahre Vorstellung von dem Inhalt der Speculationen geben k\"onnen, denen Riemann einen grossen Theil seiner Gedankenarbeit widmete und die ihn viele Jahre seines Lebens hindurch begleitet haben. Diese Bruchst\"ucke d\"urften trotz ihrer L\"uckenhaftigkeit und Unvollst\"andigkeit geeignet sein, auch in weiteren Kreisen Aufmerksamkeit zu erregen, wenn sie auch nicht viel mehr als die Anf\"ange und die allgemeinsten Grundz\"uge einer eigenth\"umlichen und tiefsinnigen Weltanschauung enthalten.

Eine willkommene Beigabe f\"ur die Freunde und Verehrer Riemann's wird endlich die biographische Skizze sein, welche Dedekind auf meinen Wunsch auf der Grundlage von Briefen und anderen Mittheilungen der Riemann'schen Familie, unterst\"utzt durch seine eigenen Erinnerungen verfasst hat.

Was die Anordnung des Stoffes betrifft, so ist in den beiden ersten Abtheilungen die chronologische Reihenfolge streng inne gehalten worden; in der dritten Abtheilung, welche den Nachlass enth\"alt, konnte diese Anordnung nicht ganz consequent durchgef\"uhrt werden, theils weil sich die Entstehungszeit hier nicht immer vollst\"andig feststellen liess, theils weil die mehr ausgef\"uhrten Untersuchungen dem Fragmentarischen vorangestellt werden sollten.

\begin{flushright}
K\"onigsberg, im M\"arz 1876.\\
H.\ Weber.
\end{flushright}

\clearpage

% ============================================
% Vorrede zur zweiten Auflage (Pages 12-13)
% ============================================
\section*{Vorrede zur zweiten Auflage.}
\addcontentsline{toc}{section}{Vorrede zur zweiten Auflage}

Sechzehn Jahre sind seit dem Erscheinen der ersten Auflage der Gesammtausgabe von Riemann's Werken verstrichen. Der Entwicklungsgang, den in diesem Zeitraum die mathematische Wissenschaft genommen hat, laesst vielfach und deutlich die Spuren von Riemann's Wirken erkennen; wir brauchen nur an den Ausbau der Theorie der Abel'schen Functionen und der linearen Differentialgleichungen, an die Lehre von den mehrfach ausgedehnten Mannigfaltigkeiten und die nicht-euklidische Geometrie zu erinnern, die mit Allem, was damit zusammenhaengt, jetzt im Vordergrund des wissenschaftlichen Interesses stehen. Nicht nur Riemann's ausgefuehrte Arbeiten, sondern auch manche der im Nachlass vorgefundenen, in den gesammten Werken mitgetheilten Andeutungen und Fragmente haben den Anstoss zu weitergehenden Forschungen gegeben.

Die Form und Art der Ausgabe der Gesammtwerke hat die Zustimmung der Mathematiker gefunden. Bedenken und Einwendungen, die hier und da in der Literatur hervorgetreten sind, und sich meist auf die von den Herausgebern zugefuegten Noten beziehen, sind in der neuen Auflage nach Moeglichkeit beruecksichtigt, und erledigen sich wohl durch eine etwas ausfuehrlichere Darstellung.

Der handschriftliche Nachlass wurde im Laufe der Jahre mehrfach und besonders bei der Vorbereitung der neuen Ausgabe einer Durchsicht unterworfen. Die Ausbeute war zwar nicht sehr gross, lieferte aber doch manchen schaetzbaren Zusatz, der unter die Anmerkungen aufgenommen werden konnte. Neu hinzugefuegt ist das kleine Fragment XXV ueber die Bewegung der Waerme im Ellipsoid, ferner ein Zusatz zu Nr.\ XXX (jetzt XXXI) ueber die quadratischen Relationen, die zwischen den Functionen $\varphi$ der Theorie der Abel'schen Functionen bestehen. Dem XXV.\ (jetzt XXVI.)\ Fragment wurde ein Zusatz im Titel gegeben, wodurch deutlicher auf seine grosse allgemeine Bedeutung hingewiesen werden sollte.

Sorgfaeltig durchgearbeitet und erweitert sind die Anmerkungen, wodurch wir hoffen, ihre Brauchbarkeit zu erhoehen. Die Erlaeuterungen von Dedekind zu dem Fragment ueber die Grenzfaelle der elliptischen Modulfunctionen sind ganz neu redigirt und erleichtern in dieser Form noch mehr den Zugang zu den Formeln Riemann's.

Auch die Erlaeuterungen zu Nr.\ XXII ``Commentatio mathematica etc.'' sind etwas ausfuehrlicher gestaltet worden, da die Darstellung in der ersten Auflage das Verstaendniss noch nicht hinlaenglich zu foerdern schien. Der Herausgeber hat es dagegen nicht unternommen, die Anwendung auf die Preisaufgabe der Pariser Akademie weiter zu verfolgen; auch ist ihm kein Versuch bekannt geworden, diese Frage weiter zu foerdern, so dass das Problem immer noch nicht als vollstaendig geloest betrachtet werden kann. Vielleicht ermuthigen diese Zeilen juengere Fachgenossen, die vor muehsamen Rechnungen nicht zurueckschrecken, das Problem aufs Neue in Angriff zu nehmen, das nicht nur durch sich selbst, sondern besonders auch durch die tiefen und eigenartigen Hulfsmittel, die Riemann zu seiner Loesung geschaffen hat, von hohem wissenschaftlichem Werthe ist.

\begin{flushright}
Marburg, im Juli 1892.\\
H.\ Weber.
\end{flushright}

\clearpage

% ============================================
% Inhalt (Pages 14-16) - Table of Contents
% ============================================
\section*{Inhalt.}
\addcontentsline{toc}{section}{Inhalt}

\subsection*{Erste Abtheilung.}
\addcontentsline{toc}{subsection}{Erste Abtheilung}

\textit{Abhandlungen, die von Riemann selbst veroeffentlicht sind.}

\begin{longtable}{p{12cm}r}
I.\ Grundlagen fuer eine allgemeine Theorie der Functionen einer veraenderlichen complexen Groesse. (Inauguraldissertation, Goettingen 1851.) & 3\\
\quad Anmerkungen zur vorstehenden Abhandlung & 46\\[0.5em]

II.\ Ueber die Gesetze der Vertheilung von Spannungselectricitaet in ponderablen Koerpern, wenn diese nicht als vollkommene Leiter oder Nichtleiter, sondern als dem Enthalten von Spannungselectricitaet mit endlicher Kraft widerstrebend betrachtet werden & 49\\
\quad (Amtlicher Bericht ueber die 31.\ Versammlung deutscher Naturforscher und Aerzte zu Goettingen im September 1854.)\\[0.5em]

III.\ Zur Theorie der Nobili'schen Farbenringe & 55\\
\quad (Aus Poggendorff's Annalen der Physik und Chemie, Bd.\ 95. 28.\ Maerz 1855.)\\
\quad Anmerkungen zur vorstehenden Abhandlung & 63\\[0.5em]

IV.\ Beitraege zur Theorie der durch die Gauss'sche Reihe $F(alpha, beta, gamma, x)$ darstellbaren Functionen & 67\\
\quad (Aus dem siebenten Bande der Abhandlungen der Koeniglichen Gesellschaft der Wissenschaften zu Goettingen. 1857.)\\[0.5em]

V.\ Selbstanzeige der vorstehenden Abhandlung & 84\\
\quad (Goettinger Nachrichten 1857. No.\ 1.)\\
\quad Anmerkungen zur Abhandlung IV & 86\\[0.5em]

VI.\ Theorie der Abel'schen Functionen. & \\
\quad (Aus Borchardt's Journal fuer reine und angewandte Mathematik, Bd.\ 54. 1857.)\\
\quad 1.\ Allgemeine Voraussetzungen und Hulfsmittel fuer die Untersuchung von Functionen unbeschraenkt veraenderlicher Groessen & 88\\
\quad 2.\ Lehrsaetze aus der analysis situs fuer die Theorie der Integrale von zweigliedrigen vollstaendigen Differentialien & 91\\
\quad 3.\ Bestimmung einer Function einer veraenderlichen complexen Groesse durch Grenz- und Unstetigkeitsbedingungen & 96\\
\quad 4.\ Theorie der Abel'schen Functionen & 100\\
\quad Anmerkungen zur vorstehenden Abhandlung & 143\\[0.5em]

VII.\ Ueber die Anzahl der Primzahlen unter einer gegebenen Groesse & 145\\
\quad (Monatsberichte der Berliner Akademie, November 1859.)\\
\quad Anmerkungen zur vorstehenden Abhandlung & 154\\
\end{longtable}

\clearpage

\begin{longtable}{p{12cm}r}
VIII.\ Ueber die Fortpflanzung ebener Luftwellen von endlicher Schwingungsweite & 157\\
\quad (Aus dem achten Bande der Abhandlungen der Koeniglichen Gesellschaft der Wissenschaften zu Goettingen. 1860.)\\[0.5em]

IX.\ Selbstanzeige der vorstehenden Abhandlung & 176\\
\quad (Goettinger Nachrichten, 1859. No.\ 19.)\\
\quad Anmerkungen zur Abhandlung VIII & 179\\[0.5em]

X.\ Ein Beitrag zu den Untersuchungen ueber die Bewegung eines fluessigen gleichartigen Ellipsoides & 182\\
\quad (Aus dem neunten Bande der Abhandlungen der Koeniglichen Gesellschaft der Wissenschaften zu Goettingen, 1861.)\\[0.5em]

XI.\ Ueber das Verschwinden der Theta-Functionen & 212\\
\quad (Aus Borchardt's Journal fuer reine und angewandte Mathematik, Bd.\ 65. 1865.)\\[0.5em]

\multicolumn{2}{l}{\textit{Zweite Abtheilung.}}\\
\multicolumn{2}{l}{\textit{Abhandlungen, die nach Riemann's Tode bereits herausgegeben sind.}}\\[0.5em]

XII.\ Ueber die Darstellbarkeit einer Function durch eine trigonometrische Reihe & 227\\
\quad (Habilitationsschrift, 1854, aus dem dreizehnten Bande der Abhandlungen der Koeniglichen Gesellschaft der Wissenschaften zu Goettingen.)\\
\quad Anmerkungen zur vorstehenden Abhandlung & 266\\[0.5em]

XIII.\ Ueber die Hypothesen, welche der Geometrie zu Grunde liegen & 272\\
\quad (Habilitationsschrift, 1854, aus dem dreizehnten Bande der Abhandlungen der Koeniglichen Gesellschaft der Wissenschaften zu Goettingen.)\\[0.5em]

XIV.\ Ein Beitrag zur Electrodynamik. (1858.) & 288\\
\quad (Aus Poggendorff's Annalen der Physik und Chemie. Bd.\ CXXXI.)\\[0.5em]

XV.\ Beweis des Satzes, dass eine einwerthige mehr als $2n$ fach periodische Function von $n$ Veraenderlichen unmoeglich ist & 294\\
\quad (Auszug aus einem Schreiben Riemann's an Herrn Weierstrass vom 26.\ October 1859. Aus Borchardt's Journal fuer reine und angewandte Mathematik, Bd.\ 71.)\\[0.5em]

XVI.\ Estratto di una lettera scritta in lingua Italiana il di 21.\ Gennaio 1864 al Sig.\ Professore Enrico Betti & 298\\
\quad (Annali di Matematica, Ser.\ 1.\ T.\ VII.)\\[0.5em]

XVII.\ Ueber die Flaeche vom kleinsten Inhalt bei gegebener Begrenzung & 301\\
\quad (Aus dem dreizehnten Bande der Abhandlungen der Koeniglichen Gesellschaft der Wissenschaften zu Goettingen.)\\
\quad Anmerkungen zur vorstehenden Abhandlung & 334\\[0.5em]

XVIII.\ Mechanik des Ohres & 338\\
\quad (Aus Henle und Pfeuffer's Zeitschrift fuer rationelle Medicin, dritte Reihe. Bd.\ 29.)\\
\end{longtable}

\clearpage

\begin{longtable}{p{12cm}r}
\multicolumn{2}{l}{\textit{Dritte Abtheilung.}}\\
\multicolumn{2}{l}{\textit{Nachlass.}}\\[0.5em]

XIX.\ Versuch einer allgemeinen Auffassung der Integration und Differentiation. (1847.) & 353\\[0.5em]

XX.\ Neue Theorie des Rueckstandes in electrischen Bindungsapparaten. (1854.) & 367\\[0.5em]

XXI.\ Zwei allgemeine Lehrsaetze ueber lineare Differentialgleichungen mit algebraischen Coefficienten. (20.\ Febr.\ 1857.) & 379\\[0.5em]

XXII.\ Commentatio mathematica, qua respondere tentatur quaestioni ab $Ill^{ma}$ Academia Parisiensi propositae: & 391\\
\quad ``Trouver quel doit etre l'etat calorifique d'un corps solide homogene indefini pour qu'un Systeme de courbes isothermes, a un instant donne, restent isothermes apres un temps quelconque, de telle sorte que la temperature d'un point puisse s'exprimer en fonction du temps et de deux autres variables independantes.'' (1861.)\\
\quad Anmerkungen zur vorstehenden Abhandlung & 405\\[0.5em]

XXIII.\ Sullo svolgimento del quoziente di due serie ipergeometriche in frazione continua infinita. (1863.) & 424\\[0.5em]

XXIV.\ Ueber das Potential eines Ringes & 431\\[0.5em]

XXV.\ Verbreitung der Waerme im Ellipsoid & 437\\[0.5em]

XXVI.\ Gleichgewicht der Electricitaet auf Cylindern mit kreisfoermigem Querschnitt und parallelen Axen. Conforme Abbildung von durch Kreise begrenzten Figuren & 440\\[0.5em]

XXVII.\ Beispiele von Flaechen kleinsten Inhalts bei gegebener Begrenzung & 445\\[0.5em]

XXVIII.\ Fragmente ueber die Grenzfaelle der elliptischen Modulfunctionen. (1852.) & 455\\
\quad Erlaeuterungen zu den vorstehenden Fragmenten von R.\ Dedekind & 466\\[0.5em]

XXIX.\ Fragment aus der Analysis Situs & 479\\[0.5em]

XXX.\ Convergenz der $p$-fach unendlichen Theta-Reihe & 483\\[0.5em]

XXXI.\ Zur Theorie der Abel'schen Functionen & 487\\
\end{longtable}

\subsection*{Anhang.}
\addcontentsline{toc}{subsection}{Anhang}

\textit{Fragmente philosophischen Inhalts.}

\begin{longtable}{p{12cm}r}
I.\ Zur Psychologie und Metaphysik & 509\\
II.\ Erkenntnisstheoretisches & 521\\
III.\ Naturphilosophie & 526\\[0.5em]
Bernhard Riemann's Lebenslauf & 539\\
\end{longtable}

\clearpage

% ============================================
% Erste Abtheilung (Page 17)
% ============================================
\begin{center}
{\Large\textsc{Erste Abtheilung.}}\\[1em]
{\large\textsc{Riemann's gesammelte mathematische Werke.}}
\end{center}

\clearpage

% ============================================
% Paper I: Grundlagen (Pages 19-50)
% ============================================
\section*{I.}
\addcontentsline{toc}{section}{I. Grundlagen fuer eine allgemeine Theorie der Functionen einer veraenderlichen complexen Groesse.}

\begin{center}
{\Large\textbf{Grundlagen fuer eine allgemeine Theorie der Functionen einer veraenderlichen complexen Groesse.}}\\[0.5em]
(Inauguraldissertation, Goettingen, 1851; zweiter unveraenderter Abdruck, Goettingen 1867.)
\end{center}

\subsection*{1.}

Denkt man sich unter $z$ eine veraenderliche Groesse, welche nach und nach alle moeglichen reellen Werthe annehmen kann, so wird, wenn jedem ihrer Werthe ein einziger Werth der unbestimmten Groesse $w$ entspricht, $w$ eine \emph{Function} von $z$ genannt, und wenn, waehrend $z$ alle zwischen zwei festen Werthen gelegenen Werthe stetig durchlaeuft, $w$ ebenfalls stetig sich aendert, so heisst diese Function innerhalb dieses Intervalls \emph{stetig} oder \emph{continuirlich}.\footnote{Diese Definition setzt offenbar zwischen den einzelnen Werthen der Function durchaus kein Gesetz fest, indem, wenn ueber diese Function fuer ein bestimmtes Intervall verfuegt ist, die Art ihrer Fortsetzung ausserhalb desselben ganz der Willkuer ueberlassen bleibt.}

Die Abhaengigkeit der Groesse $w$ von $z$ kann durch ein mathematisches Gesetz gegeben sein, so dass durch bestimmte Groessenoperationen zu jedem Werthe von $z$ das ihm entsprechende $w$ gefunden wird. Die Faehigkeit, fuer alle innerhalb eines gegebenen Intervalls liegenden Werthe von $z$ durch dasselbe Abhaengigkeitsgesetz bestimmt zu werden, schrieb man frueher nur einer gewissen Gattung von Functionen zu (\textit{functiones continuae} nach Euler's Sprachgebrauch); neuere Untersuchungen haben indess gezeigt, dass es analytische Ausdruecke giebt, durch welche eine jede stetige Function fuer ein gegebenes Intervall dargestellt werden kann. Es ist daher einerlei, ob man die Abhaengigkeit der Groesse $w$ von der Groesse $z$ als eine willkuerlich gegebene oder als eine durch bestimmte Groessenoperationen bedingte definirt. Beide Begriffe sind in Folge der erwaenten Theoreme congruent.

Anders verhaelt es sich aber, wenn die Veraenderlichkeit der Groesse $z$ nicht auf reelle Werthe beschraenkt wird, sondern auch complexe von der Form $x + yi$ (wo $i = \sqrt{-1}$) zugelassen werden.

Es seien $x + yi$ und $x + yi + dx + dyi$ zwei unendlich wenig verschiedene Werthe der Groesse $z$, welchen die Werthe $u + vi$ und $u + vi + du + dvi$ der Groesse $w$ entsprechen. Alsdann wird, wenn die Abhaengigkeit der Groesse $w$ von $z$ eine willkuerlich angenommene ist, das Verhaeltniss $\frac{du + dvi}{dx + dyi}$ sich mit den Werthen von $dx$ und $dy$, allgemein zu reden, aendern, indem, wenn man $dx + dyi = \varepsilon e^{\varphi i}$ setzt,
\begin{multline*}
\frac{du + dvi}{dx + dyi} = \frac{1}{2}\left(\frac{\partial u}{\partial x} + \frac{\partial v}{\partial y}\right) + \frac{1}{2}\left(\frac{\partial u}{\partial x} - \frac{\partial v}{\partial y}\right) \frac{dx - dyi}{dx + dyi} \\
+ \frac{i}{2}\left(\frac{\partial u}{\partial y} + \frac{\partial v}{\partial x}\right) + \frac{i}{2}\left(\frac{\partial v}{\partial x} - \frac{\partial u}{\partial y}\right) \frac{dx - dyi}{dx + dyi}
\end{multline*}
wird. Auf welche Art aber auch $w$ als Function von $z$ durch Verbindung der einfachen Groessenoperationen bestimmt werden moege, immer wird der Werth des Differentialquotienten $\frac{dw}{dz}$ von dem besondern Werthe des Differentials $dz$ unabhaengig sein.\footnote{Diese Behauptung ist offenbar in allen Faellen gerechtfertigt, wo sich aus dem Ausdrucke von $w$ durch $z$ mittelst der Regeln der Differentiation ein Ausdruck von $\frac{dw}{dz}$ durch $z$ finden laesst; ihre streng allgemeine Gueltigkeit bleibt fuer jetzt dahin gestellt.} Offenbar kann also auf diesem Wege nicht jede beliebige Abhaengigkeit der complexen Groesse $w$ von der complexen Groesse $z$ ausgedrueckt werden.

Das eben hervorgehobene Merkmal aller irgendwie durch Groessenoperationen bestimmbaren Functionen werden wir fuer die folgende Untersuchung, wo eine solche Function unabhaengig von ihrem Ausdrucke betrachtet werden soll, zu Grunde legen, indem wir, ohne jetzt dessen Allgemeingueltigkeit und Zulaenglichkeit fuer den Begriff einer durch Groessenoperationen ausdrueckbaren Abhaengigkeit zu beweisen, von folgender Definition ausgehen:

\subsection*{2.}

Eine veraenderliche complexe Groesse $w$ heisst eine \emph{Function} einer andern veraenderlichen complexen Groesse $z$, wenn sie mit ihr sich so aendert, dass der Werth des Differentialquotienten $\frac{dw}{dz}$ unabhaengig von dem Werthe des Differentials $dz$ ist.

Sowohl die Groesse $z$, als die Groesse $w$ werden als veraenderliche Groessen betrachtet, die jeden complexen Werth annehmen koennen. Die Auffassung einer solchen Veraenderlichkeit, welche sich auf ein zusammenhaengendes Gebiet von zwei Dimensionen erstreckt, wird wesentlich erleichtert durch eine Anknuepfung an raeumliche Anschauungen.

Man denke sich jeden Werth $x + yi$ der Groesse $z$ repraesentirt durch einen Punkt $O$ der Ebene $A$, dessen rechtwinklige Coordinaten $x$, $y$, jeden Werth $u + vi$ der Groesse $w$ durch einen Punkt $Q$ der Ebene $B$, dessen rechtwinklige Coordinaten $u$, $v$ sind. Eine jede Abhaengigkeit der Groesse $w$ von $z$ wird sich dann darstellen als eine Abhaengigkeit der Lage des Punktes $Q$ von der des Punktes $O$. Entspricht jedem Werthe von $z$ ein bestimmter mit $z$ stetig sich aendernder Werth von $w$, mit andern Worten, sind $u$ und $v$ stetige Functionen von $x$, $y$, so wird jedem Punkte der Ebene $A$ ein Punkt der Ebene $B$, jeder Linie, allgemein zu reden, eine Linie, jedem zusammenhaengenden Flaechenstuecke ein zusammenhaengendes Flaechenstueck entsprechen. Man wird sich also diese Abhaengigkeit der Groesse $w$ von $z$ vorstellen koennen als eine \emph{Abbildung der Ebene $A$ auf der Ebene $B$}.

\subsection*{3.}

Es soll nun untersucht werden, welche Eigenschaft diese Abbildung erhaelt, wenn $w$ eine Function der complexen Groesse $z$, d.\ h. wenn $\frac{dw}{dz}$ von $dz$ unabhaengig ist.

Wir bezeichnen durch $o$ einen unbestimmten Punkt der Ebene $A$ in der Naehe von $O$, sein Bild in der Ebene $B$ durch $q$, ferner durch $x + yi + dx + dyi$ und $u + vi + du + dvi$ die Werthe der Groessen $z$ und $w$ in diesen Punkten. Es koennen dann $dx$, $dy$ und $du$, $dv$ als rechtwinklige Coordinaten der Punkte $o$ und $q$ in Bezug auf die Punkte $O$ und $Q$ als Anfangspunkte angesehen werden, und wenn man $dx + dyi = \varepsilon e^{\varphi i}$ und $du + dvi = \eta e^{\psi i}$ setzt, so werden die Groessen $\varepsilon$, $\varphi$, $\eta$, \psi$ Polarcoordinaten dieser Punkte fuer dieselben Anfangspunkte sein.

Sind nun $o$ und $o'$ irgend zwei bestimmte Lagen des Punktes $o$ in unendlicher Naehe von $O$, und drueckt man die von ihnen abhaengigen Bedeutungen der uebrigen Zeichen durch entsprechende Indices aus, so giebt die Voraussetzung
\[
\frac{du' + dv'i}{dx' + dy'i} = \frac{du'' + dv''i}{dx'' + dy''i}
\]
und folglich
\[
\frac{du' + dv'i}{du'' + dv''i} = \frac{\eta'}{\eta''} e^{(\psi' - \psi'')i} = \frac{dx' + dy'i}{dx'' + dy''i} = \frac{\varepsilon'}{\varepsilon''} e^{(\varphi' - \varphi'')i},
\]
woraus $\frac{\eta'}{\eta''} = \frac{\varepsilon'}{\varepsilon''}$ und $\psi' - \psi'' = \varphi' - \varphi''$, d.\ h. in den Dreiecken $oOo''$ und $qQq''$ sind die Winkel $oOo''$ und $qQq''$ gleich und die sie einschliessenden Seiten einander proportional.

Es findet also zwischen zwei einander entsprechenden unendlich kleinen Dreiecken und folglich allgemein zwischen den kleinsten Theilen der Ebene $A$ und ihres Bildes auf der Ebene $B$ \emph{Aehnlichkeit Statt}. Eine Ausnahme von diesem Satze tritt nur in den besonderen Faellen ein, wenn die einander entsprechenden Aenderungen der Groessen $z$ und $w$ nicht in einem endlichen Verhaeltnisse zu einander stehen, was bei Herleitung desselben stillschweigend vorausgesetzt ist.\footnote{Ueber diesen Gegenstand sehe man:\\
\null\quad``Allgemeine Aufloesung der Aufgabe: Die Theile einer gegebenen Flaeche so abzubilden, dass die Abbildung dem Abgebildeten in den kleinsten Theilen aehnlich wird, von C.\ F.\ Gauss. (Als Beantwortung der von der koeniglichen Societaet der Wissenschaften in Copenhagen fuer 1822 aufgegebenen Preisfrage, abgedruckt in:\\
\null\quad``Astronomische Abhandlungen, herausgegeben von Schumacher. Drittes Heft. Altona. 1825.'') (Gauss Werke Bd.\ IV, p.\ 189.)}

\subsection*{4.}

Bringt man den Differentialquotienten $\frac{du + dvi}{dx + dyi}$ in die Form
\[
\frac{\left(\frac{\partial u}{\partial x} + \frac{\partial v}{\partial y}i\right) + \left(\frac{\partial u}{\partial y} - \frac{\partial v}{\partial x}\right)i}{dx + dyi},
\]
so erhellt, dass er und zwar nur dann fuer je zwei Werthe von $dx$ und $dy$ denselben Werth haben wird, wenn
\begin{equation}\label{eq:cauchy-riemann}
\frac{\partial u}{\partial x} = \frac{\partial v}{\partial y}, \qquad \frac{\partial v}{\partial x} = -\frac{\partial u}{\partial y}
\end{equation}
ist. Diese Bedingungen sind also hinreichend und nothwendig, damit $w = u + vi$ eine Function von $z = x + yi$ sei. Fuer die einzelnen Glieder dieser Function fliessen aus ihnen die folgenden:
\begin{equation}
\frac{\partial^2 u}{\partial x^2} + \frac{\partial^2 u}{\partial y^2} = 0, \qquad \frac{\partial^2 v}{\partial x^2} + \frac{\partial^2 v}{\partial y^2} = 0,
\end{equation}
welche fuer die Untersuchung der Eigenschaften, die Einem Gliede einer solchen Function einzeln betrachtet zukommen, die Grundlage bilden. Wir werden den Beweis fuer die wichtigsten dieser Eigenschaften einer eingehenderen Betrachtung der vollstaendigen Function voraufgehen lassen, zuvor aber noch einige Punkte, welche allgemeineren Gebieten angehoeren, heroertern und festlegen, um uns den Boden fuer jene Untersuchungen zu ebenen.

Fuer die folgenden Betrachtungen beschraenken wir die Veraenderlichkeit der Groessen $x$, $y$ auf ein endliches Gebiet, indem wir als Ort des Punktes $O$ nicht mehr die Ebene $A$ selbst, sondern eine ueber dieselbe ausgebreitete Flaeche $T$ betrachten. Wir waehlen diese Einkleidung, bei der es unanstoessig sein wird, von auf einander liegenden Flaechen zu reden, um die Moeglichkeit offen zu lassen, dass der Ort des Punktes $O$ ueber denselben Theil der Ebene sich mehrfach erstrecke, setzen jedoch fuer einen solchen Fall voraus, dass die auf einander liegenden Flaechentheile nicht laengs einer Linie zusammenhaengen, so dass eine Umfaltung der Flaeche, oder eine Spaltung in auf einander liegende Theile nicht vorkommt.

Die Anzahl der in jedem Theile der Ebene auf einander liegenden Flaechentheile ist alsdann vollkommen bestimmt, wenn die Begrenzung der Lage und dem Sinne nach (d.\ h. ihre innere und aussere Seite) gegeben ist; ihr Verlauf kann sich jedoch noch verschieden gestalten.

In der That, ziehen wir durch den von der Flaeche bedeckten Theil der Ebene eine beliebige Linie $l$, so aendert sich die Anzahl der ueber einander liegenden Flaechentheile nur beim Ueberschreiten der Begrenzung, und zwar beim Uebertritt von Aussen nach Innen um $+1$, im entgegengesetzten Falle um $-1$, und ist also ueberall bestimmt. Laengs des Ufers dieser Linie setzt sich nun jeder angrenzende Flaechentheil auf ganz bestimmte Art fort, so lange die Linie die Begrenzung nicht trifft, da eine Unbestimmtheit jedenfalls nur in einem einzelnen Punkte und also entweder in einem Punkte der Linie selbst oder in einer endlichen Entfernung von derselben Statt hat; wir koennen daher, wenn wir unsere Betrachtung auf einen im Innern der Flaeche verlaufenden Theil der Linie $l$ und zu beiden Seiten auf einen hinreichend kleinen Flaechenstreifen beschraenken, von bestimmten angrenzenden Flaechentheilen reden, deren Anzahl auf jeder Seite gleich ist, und die wir, indem wir der Linie eine bestimmte Richtung beilegen, auf der Linken mit $a_1$, $a_2$, \dots, $a_n$ und auf der Rechten mit $a_1'$, $a_2'$, \dots, $a_n'$ bezeichnen.

Jeder Flaechentheil $a$ wird sich dann in einen der Flaechentheile $a'$ fortsetzen; dieser wird zwar im Allgemeinen fuer den ganzen Lauf der Linie $l$ derselbe sein, kann sich jedoch fuer besondere Lagen von $l$ in einem ihrer Punkte aendern. Nehmen wir an, dass oberhalb eines solchen Punktes $c$ (d.\ h. laengs des vorhergehenden Theils von $l$) mit den Flaechentheilen $a_1$, $a_2$, \dots, $a_n$ der Reihe nach die Flaechentheile $a_1'$, $a_2'$, \dots, $a_n'$ verbunden seien, unterhalb desselben aber die Flaechentheile $a_{\alpha_1}'$, $a_{\alpha_2}'$, \dots, $a_{\alpha_n}'$, wo $\alpha_1$, $\alpha_2$, \dots, $\alpha_n$ nur in der Anordnung von $1$, $2$, \dots, $n$ verschieden sind, so wird ein oberhalb $c$ von $a_1$ in $a_1'$ eintretender Punkt, wenn er unterhalb $c$ auf die linke Seite zuruecktritt, in den Flaechentheil $a_{\alpha_1}$ gelangen, und wenn er den Punkt $c$ von der Linken zur Rechten umkreiset, wird der Index des Flaechentheils, in welchem er sich befindet, der Reihe nach die Zahlen
\[
1, \alpha_1, \alpha_{\alpha_1}, \dots, \beta, \alpha_\beta, \dots
\]
durchlaufen. In dieser Reihe sind, so lange das Glied $1$ nicht wiederkehrt, nothwendig alle Glieder von einander verschieden, weil einem beliebigen mittlern Gliede $\alpha_\mu$ nothwendig $\mu$ und nach einander alle frueheren Glieder bis $1$ in unmittelbarer Folge vorhergehen; wenn aber nach einer Anzahl von Gliedern, die offenbar kleiner als $n$ sein muss und $= m$ sei, das Glied $1$ wiederkehrt, so muessen die uebrigen Glieder in derselben Ordnung folgen. Der um $c$ sich bewegende Punkt kommt alsdann nach je $m$ Umlaeufen in denselben Flaechentheil zurueck und ist auf $m$ der auf einander liegenden Flaechentheile eingeschraenkt, welche sich ueber $c$ zu einem einzigen Punkte vereinigen. Wir nennen diesen Punkt einen \emph{Windungspunkt $(m-1)$ter Ordnung} der Flaeche $T$.

Durch Anwendung desselben Verfahrens auf die uebrigen $n - m$ Flaechentheile werden diese, wenn sie nicht gesondert verlaufen, in Systeme von $m_2$, $m_3$, \dots\ Flaechentheilen zerfallen, in welchem Falle auch noch Windungspunkte $(m_2 - 1)$ter, $(m_3 - 1)$ter \dots\ Ordnung in dem Punkte $c$ liegen.

Wenn die Lage und der Sinn der Begrenzung von $T$ und die Lage ihrer Windungspunkte gegeben ist, so ist $T$ entweder vollkommen bestimmt oder doch auf eine endliche Anzahl verschiedener Gestalten beschraenkt; Letzteres, in so fern sich diese Bestimmungsstuecke auf verschiedene der auf einander liegenden Flaechentheile beziehen koennen.

\subsection*{5.}

Eine veraenderliche Groesse, die fuer jeden Punkt $O$ der Flaeche $T$ allgemein zu reden, d.\ h. ohne eine Ausnahme in einzelnen Linien und Punkten\footnote{Diese Beschraenkung ist zwar nicht durch den Begriff einer Function an sich geboten, aber um Infinitesimalrechnung auf sie anwenden zu koennen erforderlich: eine Function, die in allen Punkten einer Flaeche unstetig ist, wie z.\ B. eine Function, die fuer ein commensurables $x$ und ein commensurables $y$ den Werth $1$, sonst aber den Werth $2$ hat, kann weder einer Differentiation, noch einer Integration, also (unmittelbar) der Infinitesimalrechnung ueberhaupt nicht unterworfen werden. Die fuer die Flaeche $T$ hier willkuerlich gemachte Beschraenkung wird sich spaeter (Art.\ 15) rechtfertigen.} auszuschliessen, Einen bestimmten mit der Lage desselben stetig sich aendernden Werth annimmt, kann offenbar als eine Function von $x$, $y$ angesehen werden, und ueberall, wo in der Folge von Functionen von $x$, $y$ die Rede sein wird, werden wir den Begriff derselben auf diese Art festlegen.

\subsection*{6.}

Ehe wir uns jedoch zur Betrachtung solcher Functionen wenden, schalten wir noch einige Eroerterungen ueber den Zusammenhang einer Flaeche ein. Wir beschraenken uns dabei auf solche Flaechen, die sich nicht laengs einer Linie spalten.

Wir betrachten zwei Flaechentheile als \emph{zusammenhaengend} oder Einem Stuecke angehoerig, wenn sich von einem Punkte des einen durch das Innere der Flaeche eine Linie nach einem Punkte des andern ziehen laesst, als \emph{getrennt}, wenn diese Moeglichkeit nicht Statt findet.

Die Untersuchung des Zusammenhangs einer Flaeche beruht auf ihrer Zerlegung durch \emph{Querschnitte}, d.\ h. Linien, welche von einem Begrenzungspunkte das Innere einfach --- keinen Punkt mehrfach --- bis zu einem Begrenzungspunkte durchschneiden. Letzterer kann auch in dem zur Begrenzung hinzugekommenen Theile, also in einem fruehern Punkte des Querschnitts, liegen.

Eine zusammenhaengende Flaeche heisst, wenn sie durch jeden Querschnitt in Stuecke zerfaellt, eine \emph{einfach zusammenhaengende}, andernfalls eine \emph{mehrfach zusammenhaengende}.

\textbf{Lehrsatz I.} \emph{Eine einfach zusammenhaengende Flaeche $A$ zerfaellt durch jeden Querschnitt $ab$ in zwei einfach zusammenhaengende Stuecke.}

Gesetzt, eins dieser Stuecke wuerde durch einen Querschnitt $cd$ nicht zerstueckt, so erhielte man offenbar, je nachdem keiner seiner Endpunkte oder der Endpunkt $c$ oder beide Endpunkte in $ab$ fielen, durch Herstellung der Verbindung laengs der ganzen Linie $ab$ oder laengs des Theils $cb$ oder des Theils $cd$ derselben eine zusammenhaengende Flaeche, welche durch einen Querschnitt aus $A$ entstaende, gegen die Voraussetzung.

\textbf{Lehrsatz II.} \emph{Wenn eine Flaeche $T$ durch $n_1$ Querschnitte $q_1$ in ein System $T_1$ von $m_1$ einfach zusammenhaengenden Flaechenstuecken und durch $n_2$ Querschnitte $q_2$ in ein System $T_2$ von $m_2$ Flaechenstuecken zerfaellt, so kann $n_2 - m_2$ nicht $> n_1 - m_1$ sein.}

Jede Linie $q_2$ bildet, wenn sie nicht ganz in das Querschnittsystem $q_1$ faellt, zugleich einen oder mehrere Querschnitte $q_2'$ der Flaeche $T_1$. Als Endpunkte der Querschnitte $q_2'$ sind anzusehen:

1) die $2n_2$ Endpunkte der Querschnitte $q_2$, ausgenommen, wenn ihre Enden mit einem Theil des Liniensystems $q_1$ zusammenfallen,

2) jeder mittlere Punkt eines Querschnitts $q_2$, in welchem er in einen mittlern Punkt einer Linie $q_1$ eintritt, ausgenommen, wenn er sich schon in einer andern Linie $q_2$ befindet, d.\ h. wenn ein Ende eines Querschnitts $q_2$ mit ihm zusammenfaellt.

Bezeichnet nun $\mu$, wie oft Linien beider Systeme waehrend ihres Laufes zusammentreffen oder auseinandergehen (wo also ein einzelner gemeinsamer Punkt doppelt zu rechnen ist), $\nu_2$, wie oft ein Endstueck der $q_2$ mit einem mittlern Stuecke der $q_1$, $\nu_1$, wie oft ein Endstueck der $q_1$ mit einem mittlern Stuecke der $q_2$, endlich $\nu_3$, wie oft ein Endstueck der $q_2$ mit einem Endstuecke der $q_1$ zusammenfaellt, so liefert Nr.\ 1 $2n_2 - \nu_2 - \nu_3$ und Nr.\ 2 $\mu - \nu_1$ Endpunkte der Querschnitte $q_2'$; beide Faelle zusammengenommen aber umfassen saemmtliche Endpunkte und jeden nur einmal, und die Anzahl dieser Querschnitte ist daher
\[
\frac{2n_2 - \nu_2 - \nu_3 + \mu - \nu_1}{2} = n_2 + s.
\]

Durch ganz aehnliche Schluesse ergiebt sich die Anzahl der Querschnitte $q_1'$ der Flaeche $T_2$, welche durch die Linien $q_1$ gebildet werden,
\[
\frac{2n_1 - \nu_1 - \nu_3 + \mu - \nu_2}{2} = n_1 + s.
\]
Die Flaeche $T_1$ wird nun offenbar durch die $n_2 + s$ Querschnitte $q_2'$ in dieselbe Flaeche verwandelt, in welche $T_2$ durch die $n_1 + s$ Querschnitte $q_1'$ zerfaellt wird. Es besteht aber $T_1$ aus $m_1$ einfach zusammenhaengenden Stuecken und zerfaellt daher nach Satz I durch $n_2 + s$ Querschnitte in $m_1 + n_2 + s$ Flaechenstuecke; folglich muesste, waere $m_2 < m_1 + n_2 - n_1$, die Zahl der Flaechenstuecke $T_2$ durch $n_1 + s$ Querschnitte um mehr als $n_1 + s$ vermehrt werden, was ungereimt ist.\footnote{Unter einer Zerlegung durch mehrere Querschnitte ist stets eine successive zu verstehen, d.\ h. eine solche, wo die durch einen Querschnitt entstandene Flaeche durch einen neuen Querschnitt weiter zerlegt wird.}

\subsection*{7.}

Zufolge dieses Lehrsatzes ist, wenn die Anzahl der Querschnitte unbestimmt durch $n$, die Anzahl der Stuecke durch $m$ bezeichnet wird, $n - m$ fuer alle Zerlegungen einer Flaeche in einfach zusammenhaengende Stuecke constant; denn betrachten wir irgend zwei bestimmte Zerlegungen durch $n_1$ Querschnitte in $m_1$ Stuecke und durch $n_2$ Querschnitte in $m_2$ Stuecke, so muss, wenn erstere einfach zusammenhaengend sind, $n_2 - m_2 \leq n_1 - m_1$, und wenn letztere einfach zusammenhaengend sind, $n_1 - m_1 \leq n_2 - m_2$; also wenn Beides zutrifft, $n_2 - m_2 = n_1 - m_1$ sein.

Diese Zahl kann fueglich mit dem Namen \emph{``Ordnung des Zusammenhangs''} einer Flaeche belegt werden; sie wird
\begin{itemize}[nosep]
\item durch jeden Querschnitt um $1$ erniedrigt --- nach der Definition ---,
\item durch eine von einem Innern Punkte das Innere einfach bis zu einem Begrenzungspunkte oder einem fruehern Schnittpunkte durchschneidende Linie nicht geaendert und
\item durch einen Innern allenthalben einfachen in zwei Punkten endenden Schnitt um $1$ erhoeht,
\end{itemize}
weil erstere durch Einen, letztere aber durch zwei Querschnitte in Einen Querschnitt verwandelt werden kann.

Endlich wird die Ordnung des Zusammenhangs einer aus mehreren Stuecken bestehenden Flaeche erhalten, wenn man die Ordnungen des Zusammenhangs dieser Stuecke zu einander addirt.

Wir werden uns indess in der Folge meistens auf eine aus Einem Stuecke bestehende Flaeche beschraenken, und uns fuer ihren Zusammenhang der kunstloseren Bezeichnung eines \emph{einfachen}, \emph{zweifachen} etc.\ bedienen, indem wir unter einer \emph{$n$fach zusammenhaengenden} Flaeche eine solche verstehen, die durch $n - 1$ Querschnitte in eine einfach zusammenhaengende zerlegbar ist.

In Bezug auf die Abhaengigkeit des Zusammenhangs der Begrenzung von dem Zusammenhang einer Flaeche erhellt leicht:

1) Die Begrenzung einer einfach zusammenhaengenden Flaeche besteht nothwendig aus Einer in sich zuruecklaufenden Linie.

Bestaende die Begrenzung aus getrennten Stuecken, so wuerde ein Querschnitt $q$, der einen Punkt eines Stuecks $a$ mit einem Punkte eines andern $b$ verbaende, nur zusammenhaengende Flaechentheile von einander scheiden, da sich im Innern der Flaeche laengs $a$ eine Linie von der einen Seite des Querschnitts $q$ an die entgegengesetzte fuehren liesse; und folglich wuerde $q$ die Flaeche nicht zerstuecken, gegen die Voraussetzung.

2) Durch jeden Querschnitt wird die Anzahl der Begrenzungsstuecke entweder um $1$ vermindert oder um $1$ vermehrt.

Ein Querschnitt $q$ verbindet entweder einen Punkt eines Begrenzungsstuecks $a$ mit einem Punkte eines andern $b$, --- in diesem Falle bilden alle diese Linien zusammengenommen in der Folge $a$, $q$, $b$, $q$ ein einziges in sich zuruecklaufendes Stueck der Begrenzung --- oder er verbindet zwei Punkte eines Stuecks der Begrenzung, --- in diesem Falle zerfaellt dieses durch seine beiden Endpunkte in zwei Stuecke, deren jedes mit dem Querschnitte zusammengenommen ein in sich zuruecklaufendes Begrenzungsstueck bildet --- oder endlich, er endet in einem seiner frueheren Punkte und kann betrachtet werden als zusammengesetzt aus einer in sich zuruecklaufenden Linie $o$ und einer andern $l$, welche einen Punkt von $o$ mit einem Punkte eines Begrenzungsstuecks $a$ verbindet, --- in welchem Falle $o$ eines Theils, und $a$, $l$, $o$, $l$ andern Theils je ein in sich zuruecklaufendes Begrenzungsstueck bilden.

Es treten also entweder --- im erstem Falle --- an die Stelle zweier Eins, oder --- in den beiden letzteren Faellen --- an die Stelle Eines zwei Begrenzungsstuecke, woraus unser Satz folgt.

Die Anzahl der Stuecke, aus welchen die Begrenzung eines $n$fach zusammenhaengenden Flaechenstuecks besteht, ist daher entweder $= n$ oder um eine gerade Zahl kleiner.

Hieraus ziehen wir noch das \textbf{Corollar}:

\emph{Wenn die Anzahl der Begrenzungsstuecke einer $n$fach zusammenhaengenden Flaeche $= n$ ist, so zerfaellt diese durch jeden ueberall einfachen im Innern in sich zuruecklaufenden Schnitt in zwei getrennte Stuecke.}

Denn die Ordnung des Zusammenhangs wird dadurch nicht geaendert, die Anzahl der Begrenzungsstuecke um $2$ vermehrt; die Flaeche wuerde also, wenn sie eine zusammenhaengende waere, einen $n$fach Zusammenhang und $n + 2$ Begrenzungsstuecke haben, was unmoeglich ist.

\subsection*{8.}

Sind $X$ und $Y$ zwei in allen Punkten der ueber $A$ ausgebreiteten Flaeche $T$ stetige Functionen von $x$, $y$, so ist das ueber alle Elemente $dT$ dieser Flaeche ausgedehnte Integral
\begin{equation}
\int \left(\frac{\partial X}{\partial x} + \frac{\partial Y}{\partial y}\right) dT = -\int (X \cos\xi + Y \cos\eta) \, ds,
\end{equation}
wenn in jedem Punkte der Begrenzung die Neigung einer auf sie nach Innen gezogenen Normale gegen die $x$-Axe durch $\xi$, gegen die $y$-Axe durch $\eta$ bezeichnet wird, und sich diese Integration auf saemmtliche Elemente $ds$ der Begrenzungslinie erstreckt.

Um dieses Integral zu transformiren, zerlegen wir den von der Flaeche $T$ bedeckten Theil der Ebene $A$ durch ein System der $x$-Axe paralleler Linien in Elementarstreifen, und zwar so, dass jeder Windungspunkt der Flaeche $T$ in eine dieser Linien faellt. Unter dieser Voraussetzung besteht der auf jeden derselben fallende Theil von $T$ aus einem oder mehreren abgesondert verlaufenden trapezfoermigen Stuecken. Der Beitrag eines unbestimmten dieser Flaechenstreifen, welcher aus der $y$-Axe das Element $dy$ ausscheidet, zu dem Werthe $\int \frac{\partial X}{\partial x} \, dT$ wird dann offenbar $= dy \int \frac{\partial X}{\partial x} \, dx$, wenn diese Integration durch diejenige oder diejenigen der Flaeche $T$ angehoerigen geraden Linien ausgedehnt wird, welche auf eine durch einen Punkt von $dy$ gehende Normale fallen. Sind nun die unteren Endpunkte derselben (d.\ h. welchen die kleinsten Werthe von $x$ entsprechen) $O_1$, $O_2$, $O_3$, \dots, die oberen $O'$, $O''$, $O'''$, \dots\ und bezeichnen wir mit $X_1$, $X_2$, \dots, $X'$, $X''$, \dots\ die Werthe von $X$ in diesen Punkten, mit $ds_1$, $ds_2$, \dots, $ds'$, $ds''$, \dots\ die entsprechenden von dem Flaechenstreifen aus der Begrenzung ausgeschiedenen Elemente, mit $\xi_1$, $\xi_2$, \dots, $\xi'$, $\xi''$, \dots\ die Werthe von $\xi$ an diesen Elementen, so wird
\[
\int \frac{\partial X}{\partial x} \, dx = -X_1 - X_2 - \dots + X' + X'' + X''' + \dots
\]
Die Winkel $\xi$ werden offenbar spitz an den unteren, stumpf an den oberen Endpunkten, und es wird daher
\begin{align*}
dy &= \cos\xi_1 \, ds_1 = \cos\xi_2 \, ds_2 = \dots \\
   &= -\cos\xi' \, ds' = -\cos\xi'' \, ds'' = \dots
\end{align*}
Durch Substitution dieser Werthe ergiebt sich
\[
\int \frac{\partial X}{\partial x} \, dT = -\int X \cos\xi \, ds,
\]
wo sich die Summation auf alle Begrenzungselemente bezieht, welche in der $y$-Axe $dy$ zur Projection haben.

Durch Integration ueber saemmtliche in Betracht kommende $dy$ werden offenbar saemmtliche Elemente der Flaeche $T$ und saemmtliche Elemente der Begrenzung erschoepft, und man erhaelt daher, in diesem Umfange genommen,
\[
\int \frac{\partial X}{\partial x} \, dT = -\int X \cos\xi \, ds.
\]

Durch ganz aehnliche Schluesse findet man
\[
\int \frac{\partial Y}{\partial y} \, dT = -\int Y \cos\eta \, ds,
\]
und folglich
\begin{equation}
\int \left(\frac{\partial X}{\partial x} + \frac{\partial Y}{\partial y}\right) dT = -\int (X \cos\xi + Y \cos\eta) \, ds, \quad \text{w.\ z.\ b.\ w.}
\end{equation}

\clearpage
\chapter*{Source segment 2: riemann pages 051 100}
\addcontentsline{toc}{chapter}{Source segment 2: riemann pages 051 100}
\selectlanguage{ngerman}

\begin{center}
{\LARGE\bfseries Bernhard Riemann --- Gesammelte Mathematische Werke}\\[0.5em]
{\large (pp.~51--100)}\\[1em]
\rule{\textwidth}{0.4pt}
\end{center}

\vspace{1em}

% ============================================================
% I. Grundlagen für eine allgemeine Theorie 
%    der Functionen einer veränderlichen complexen Grösse
%    (Fortsetzung --- Schluss)
% ============================================================

\section*{I. Grundlagen für eine allgemeine Theorie\\
der Functionen einer veränderlichen complexen Grösse.}
\addcontentsline{toc}{section}{I. Grundlagen für eine allgemeine Theorie der Functionen einer veränderlichen complexen Grösse}

\markboth{\small I. Grundlagen für eine allgemeine Theorie}{}

\subsection*{19.}
\addcontentsline{toc}{subsection}{Art. 19}

Die Principien, welche dem Lehrsätze am Schlüsse des vorigen Art.\ zu Grunde liegen, eröffnen den Weg, bestimmte Functionen einer veränderlichen complexen Grösse (unabhängig von einem Ausdrucke für dieselben) zu untersuchen.

Zur Orientirung auf diesem Felde wird ein Ueberschlag über den Umfang der zur Bestimmung einer solchen Function innerhalb eines gegebenen Grössengebiets erforderlichen Bedingungen dienen.

Halten wir uns zunächst an einen bestimmten Fall, so kann, wenn die über $A$ ausgebreitete Fläche, durch welche dies Grössengebiet dargestellt wird, eine einfach zusammenhängende ist, die Function $w = u + vi$ von den folgenden Bedingungen gemäss bestimmt werden:
\begin{enumerate}
\item für $u$ ist in allen Begrenzungspunkten ein Werth gegeben, der sich für eine unendlich kleine Ortsänderung um eine unendlich kleine Grösse von derselben Ordnung, übrigens aber beliebig ändet\footnote{An sich sind die Aenderungen dieses Werthes nur der Beschränkung unterworfen, nicht längs eines Theils der Begrenzung unstetig zu sein; eine weitere Beschränkung ist nur gemacht, um hier unnöthige Weitläufigkeiten zu vermeiden.};
\item der Werth von $v$ ist in irgend einem Punkte beliebig gegeben;
\item die Function soll in allen Punkten endlich und stetig sein.
\end{enumerate}
Durch diese Bedingungen aber ist sie vollkommen bestimmt.

In der That folgt dies aus dem Lehrsatze des vorigen Art., wenn man, was immer möglich sein wird, $\alpha + \beta i$ so bestimmt, dass $\alpha$ am Rande dem gegebenen Werth gleich und in der ganzen Fläche für jede unendlich kleine Ortsänderung die Aenderung von $\alpha + \beta i$ unendlich klein von derselben Ordnung ist.

Es kann also, allgemein zu reden, $u$ am Rande als eine ganz willkürliche Function von $s$ gegeben werden, und dadurch ist $v$ überall mit bestimmt; umgekehrt kann aber auch $v$ in jedem Begrenzungspunkte beliebig angenommen werden, woraus dann der Werth von $u$ folgt. Der Spielraum für die Wahl der Werthe von $w$ am Rande umfasst daher eine Mannigfaltigkeit von Einer Dimension für jeden Begrenzungspunkt, und die vollständige Bestimmung derselben erfordert für jeden Begrenzungspunkt Eine Gleichung, wobei es indess nicht wesentlich sein wird, dass jede dieser Gleichungen sich auf den Werth Eines Gliedes in Einem Begrenzungspunkte allein bezieht. Es wird diese Bestimmung auch so geschehen können, dass für jeden Begrenzungspunkt Eine mit der Lage dieses Punktes ihre Form stetig ändernde, beide Glieder enthaltende Gleichung gegeben ist, oder für mehrere Theile der Begrenzung gleichzeitig so, dass jedem Punkte eines dieser Theile $n - 1$ bestimmte Punkte, aus jedem der übrigen Theile einer, zugesellt und für je $n$ solcher Punkte gemeinschaftlich $n$ mit ihrer Lage stetig veränderliche Gleichungen gegeben sind. Diese Bedingungen, deren Gesammtheit eine stetige Mannigfaltigkeit bildet und welche durch Gleichungen zwischen willkürlichen Functionen ausgedrückt werden, werden aber, um für die Bestimmung einer im Innern des Grössengebiets überall stetigen Function zulässig und hinreichend zu sein, allgemein zu reden, noch einer Beschränkung oder Ergänzung durch einzelne Bedingungsgleichungen --- Gleichungen für willkürliche Constanten --- bedürfen, indem bis auf diese sich die Genauigkeit unserer Schätzung offenbar nicht erstreckt.

Für den Fall, wo das Gebiet der Veränderlichkeit der Grösse $z$ durch eine mehrfach zusammenhängende Fläche dargestellt wird, erleiden diese Betrachtungen keine wesentliche Abänderung, indem die Anwendung des Lehrsatzes in Art.~18 eine bis auf die Aenderungen beim Ueberschreiten der Querschnitte ebenso wie vorhin beschaffene Function liefert --- Aenderungen, welche ${}= 0$ gemacht werden können, wenn die Grenzbedingungen eine der Anzahl der Querschnitte gleiche Anzahl verfügbarer Constanten enthalten.

Der Fall, wo im Innern längs einer Linie auf Stetigkeit verzichtet wird, ordnet sich dem vorigen unter, wenn man diese Linie als einen Schnitt der Fläche betrachtet.

Wenn endlich in einem einzelnen Punkte eine Verletzung der Stetigkeit, also nach Art.~12 ein Unendlichwerden der Function, zugelassen wird, so kann unter Beibehaltung der sonstigen in unserm Anfangsfalle gemachten Voraussetzungen für diesen Punkt eine Function von $z$, nach deren Subtraction die zu bestimmende Function stetig werden soll, beliebig gegeben werden; dadurch aber ist sie völlig bestimmt. Denn nimmt man die Grösse $\alpha + \beta i$ in einem beliebig kleinen um den Unstetigkeitspunkt beschriebenen Kreise gleich dieser gegebenen Function, übrigens aber den früheren Vorschriften gemäss an, so wird das Integral
\[
\int \left[\left(\frac{\partial \alpha}{\partial x} - \frac{\partial \beta}{\partial y}\right)^2 + \left(\frac{\partial \alpha}{\partial y} + \frac{\partial \beta}{\partial x}\right)^2\right] dT
\]
über diesen Kreis erstreckt ${}= 0$, über den übrigen Theil erstreckt einer endlichen Grösse gleich, und man kann also den Lehrsatz des vorigen Art.\ anwenden, wodurch man eine Function mit den verlangten Eigenschaften erhält. Hieraus kann man mit Hülfe des Lehrsatzes im Art.~13 folgern, dass im Allgemeinen, wenn in einem einzelnen Unstetigkeitspunkte die Function unendlich gross von der Ordnung $n$ werden darf, eine Anzahl von $2n$ Constanten verfügbar wird.

Geometrisch dargestellt liefert (nach Art.~15) eine Function $w$ einer innerhalb eines gegebenen Grössengebiets von zwei Dimensionen veränderlichen complexen Grösse $z$, von einer gegebenen $A$ bedeckenden Fläche $T$ ein ihr in den kleinsten Theilen, einzelne Punkte ausgenommen, ähnliches, $B$ bedeckendes Abbild $S$. Die Bedingungen, welche so eben zur Bestimmung der Function hinreichend und nothwendig befunden worden sind, beziehen sich auf ihren Werth entweder in Begrenzungs- oder in Unstetigkeitspunkten; sie erscheinen also (Art.~15) sämmtlich als Bedingungen für die Lage der Begrenzung von $S$ und zwar geben sie für jeden Begrenzungspunkt Eine Bedingungsgleichung. Bezieht sich jede derselben nur auf Einen Begrenzungspunkt, so werden sie durch eine Schaar von Curven repräsentirt, von denen für jeden Begrenzungspunkt Eine den geometrischen Ort bildet. Werden zwei mit einander stetig fortrückende Begrenzungspunkte gemeinschaftlich zwei Bedingungsgleichungen unterworfen, so entsteht dadurch zwischen zwei Begrenzungstheilen eine solche Abhängigkeit, dass, wenn die Lage des einen willkürlich angenommen wird, die Lage des andern daraus folgt. Aehnlicher Weise ergiebt sich für andere Formen der Bedingungsgleichungen eine geometrische Bedeutung, was wir indess nicht weiter verfolgen wollen.

\subsection*{20.}
\addcontentsline{toc}{subsection}{Art. 20}

Die Einführung der complexen Grössen in die Mathematik hat ihren Ursprung und nächsten Zweck in der Theorie einfacher\footnote{Wir betrachten hier als Elementaroperationen Addition und Subtraction, Multiplication und Division, Integration und Differentiation, und ein Abhängigkeitsgesetz als desto einfacher, durch je weniger Elementaroperationen die Abhängigkeit bedingt wird. In der That lassen sich durch eine endliche Anzahl dieser Operationen alle bis jetzt in der Analysis benutzten Functionen definiren.} durch Grössenoperationen ausgedrückter Abhängigkeitsgesetze zwischen veränderlichen Grössen. Wendet man nämlich diese Abhängigkeitsgesetze in einem erweiterten Umfange an, indem man den veränderlichen Grössen, auf welche sie sich beziehen, complexe Werthe giebt, so tritt eine sonst versteckt bleibende Harmonie und Regelmässigkeit hervor. Die Fälle, in denen dies geschehen ist, umfassen zwar bis jetzt erst ein kleines Gebiet --- sie lassen sich fast sämmtlich auf diejenigen Abhängigkeitsgesetze zwischen zwei veränderlichen Grössen zurückführen, wo die eine entweder eine algebraische\footnote{D.~h. wo zwischen beiden eine algebraische Gleichung Statt findet.} Function der andern ist oder eine solche Function, deren Differentialquotient eine algebraische Function ist ---, aber beinahe jeder Schritt, der hier gethan ist, hat nicht bloss den ohne Hülfe der complexen Grössen gewonnenen Resultaten eine einfachere, geschlossenere Gestalt gegeben, sondern auch zu neuen Entdeckungen die Bahn gebrochen, wozu die Geschichte der Untersuchungen über algebraische Functionen, Kreis- oder Exponentialfunctionen, elliptische und Abel'sche Functionen den Beleg liefert.

Es soll kurz angedeutet werden, was durch unsere Untersuchung für die Theorie solcher Functionen gewonnen ist.

Die bisherigen Methoden, diese Functionen zu behandeln, legten stets als Definition einen Ausdruck der Function zu Grunde, wodurch ihr Werth für jeden Werth ihres Arguments gegeben wurde; durch unsere Untersuchung ist gezeigt, dass, in Folge des allgemeinen Charakters einer Function einer veränderlichen complexen Grösse, in einer Definition dieser Art ein Theil der Bestimmungsstücke eine Folge der übrigen ist, und zwar ist der Umfang der Bestimmungsstücke auf die zur Bestimmung nothwendigen zurückgeführt worden. Dies vereinfacht die Behandlung derselben wesentlich. Um z.~B. die Gleichheit zweier Ausdrücke derselben Function zu beweisen, musste man sonst den einen in den andern transformiren, d.~h. zeigen, dass beide für jeden Werth der veränderlichen Grösse übereinstimmten; jetzt genügt der Nachweis ihrer Uebereinstimmung in einem weit geringern Umfange.

Eine Theorie dieser Functionen auf den hier gelieferten Grundlagen würde die Gestaltung der Function (d.~h. ihren Werth für jeden Werth ihres Arguments) unabhängig von einer Bestimmungsweise derselben durch Grössenoperationen festlegen, indem zu dem allgemeinen Begriffe einer Function einer veränderlichen complexen Grösse nur die zur Bestimmung der Function nothwendigen Merkmale hinzugefügt würden, und dann erst zu den verschiedenen Ausdrücken deren die Function fähig ist übergehen. Der gemeinsame Charakter einer Gattung von Functionen, welche auf ähnliche Art durch Grössenoperationen ausgedrückt werden, stellt sich dann dar in der Form der ihnen auferlegten Grenz- und Unstetigkeitsbedingungen. Wird z.~B. das Gebiet der Veränderlichkeit der Grösse $z$ über die ganze unendliche Ebene $A$ einfach oder mehrfach erstreckt, und innerhalb derselben der Function nur in einzelnen Punkten eine Unstetigkeit, und zwar nur ein Unendlichwerden, dessen Ordnung endlich ist, gestattet (wobei für ein unendliches $z$ diese Grösse selbst, für jeden endlichen Werth $z'$ derselben aber $\frac{1}{z-z'}$ als ein unendlich Grosses erster Ordnung gilt), so ist die Function nothwendig algebraisch, und umgekehrt erfüllt diese Bedingung jede algebraische Function.

Die Ausführung dieser Theorie, welche, wie bemerkt, einfache durch Grössenoperationen bedingte Abhängigkeitsgesetze ins Licht zu setzen bestimmt ist, unterlassen wir indess jetzt, da wir die Betrachtung des Ausdruckes einer Function gegenwärtig ausschliessen.

Aus demselben Grunde befassen wir uns hier auch nicht damit, die Brauchbarkeit unserer Sätze als Grundlagen einer allgemeinen Theorie dieser Abhängigkeitsgesetze darzuthun, wozu der Beweis erfordert wird, dass der hier zu Grunde gelegte Begriff einer Function einer veränderlichen complexen Grösse mit dem einer durch Grössenoperationen ausdrückbaren Abhängigkeit\footnote{Es wird darunter jede durch eine endliche oder unendliche Anzahl der vier einfachsten Rechnungsoperationen, Addition und Subtraction, Multiplication und Division, ausdrückbare Abhängigkeit begriffen. Der Ausdruck Grössenoperationen soll (im Gegensatze zu Zahlenoperationen) solche Rechnungsoperationen andeuten, bei denen die Commensurabilität der Grössen nicht in Betracht kommt.} völlig zusammenfällt.

\subsection*{21.}
\addcontentsline{toc}{subsection}{Art. 21}

Es wird jedoch zur Erläuterung unserer allgemeinen Sätze ein ausgeführtes Beispiel ihrer Anwendung von Nutzen sein.

Die im vorigen Artikel bezeichnete Anwendung derselben ist, obwohl die bei ihrer Aufstellung zunächst beabsichtigte, doch nur eine specielle. Denn wenn die Abhängigkeit durch eine endliche Anzahl der dort als Elementaroperationen betrachteten Grössenoperationen bedingt ist, so enthält die Function nur eine endliche Anzahl von Parametern, was für die Form eines Systems von einander unabhängiger Grenz- und Unstetigkeitsbedingungen, die zu ihrer Bestimmung hinreichen, den Erfolg hat, dass unter ihnen längs einer Linie in jedem Punkte willkürlich zu bestimmende Bedingungen gar nicht vorkommen können. Für unsern jetzigen Zweck schien es daher geeigneter, nicht ein dorther entnommenes Beispiel zu wählen, sondern vielmehr ein solches, wo die Function der complexen Veränderlichen von einer willkürlichen Function abhängt.

Zur Veranschaulichung und bequemeren Fassung geben wir demselben die am Schlüsse des Art.~19 gebrauchte geometrische Einkleidung. Es erscheint dann als eine Untersuchung über die Möglichkeit, von einer gegebenen Fläche ein zusammenhängendes in den kleinsten Theilen ähnliches Abbild zu liefern, dessen Gestalt gegeben ist, wo also in obiger Form ausgedrückt, für jeden Begrenzungspunkt des Abbildes eine Ortscurve, und zwar für alle dieselbe, ausserdem aber (Art.~5) der Sinn der Begrenzung und die Windungspunkte desselben gegeben sind. Wir beschränken uns auf die Lösung dieser Aufgabe in dem Falle, wo jedem Punkte der einen Fläche nur Ein Punkt der andern entsprechen soll und die Flächen einfach zusammenhängend sind, für welchen Fall sie in folgendem Lehrsatz enthalten ist.

\medskip
\noindent\textbf{Lehrsatz.} \textit{Zwei gegebene einfach zusammenhängende ebene Flächen können stets so auf einander bezogen werden, dass jedem Punkte der einen Ein mit ihm stetig fortrückender Punkt der andern entspricht und ihre entsprechenden kleinsten Theile ähnlich sind; und zwar kann zu Einem innern Punkte und zu Einem Begrenzungspunkte der entsprechende beliebig gegeben werden; dadurch aber ist für alle Punkte die Beziehung bestimmt.}
\medskip

Wenn zwei Flächen $T$ und $R$ auf eine dritte $S$ so bezogen sind, dass zwischen den entsprechenden kleinsten Theilen Aehnlichkeit Statt findet, so ergiebt sich daraus eine Beziehung zwischen den Flächen $T$ und $R$, von welcher offenbar dasselbe gilt. Die Aufgabe, zwei beliebige Flächen auf einander so zu beziehen, dass Aehnlichkeit in den kleinsten Theilen Statt findet, ist dadurch auf die zurückgeführt, jede beliebige Fläche durch Eine bestimmte in den kleinsten Theilen ähnlich abzubilden. Wir haben hiernach, wenn wir in der Ebene $B$ um den Punkt, wo $w = 0$ ist, mit dem Radius $1$ einen Kreis $K$ beschreiben, um unsern Lehrsatz darzuthun, nur nöthig zu beweisen:

\medskip
\noindent\textit{Eine beliebige einfach zusammenhängende $A$ bedeckende Fläche $T$ kann durch den Kreis $K$ stets zusammenhängend und in den kleinsten Theilen ähnlich abgebildet werden und zwar nur auf Eine Art so, dass dem Mittelpunkte ein beliebig gegebener innerer Punkt $O_0$ und einem beliebig gegebenen Punkte der Peripherie ein beliebig gegebener Begrenzungspunkt $O'$ der Fläche $T$ entspricht.}
\medskip

Wir bezeichnen die bestimmten Bedeutungen von $z$, $\varrho$ für die Punkte $O_0$, $O'$ durch entsprechende Indices und beschreiben in $T$ um $O_0$ als Mittelpunkt einen beliebigen Kreis $\mathfrak{S}$, welcher sich nicht bis zur Begrenzung von $T$ erstreckt und keinen Windungspunkt enthält. Führen wir Polarcoordinaten ein, indem wir $z - z_0 = re^{\varphi i}$ setzen, so wird die Function $\log(z - z_0) = \log r + \varphi i$. Der reelle Werth ändert sich daher im ganzen Kreise mit Ausnahme des Punktes $O_0$, wo er unendlich wird, stetig. Der imaginäre aber erhält, wenn überall unter den möglichen Werthen von $\varphi$ der kleinste positive gewählt wird, längs des Radius, wo $z - z_0$ reelle positive Werthe annimmt, auf der einen Seite den Werth $0$, auf der andern den Werth $2\pi$, ändert sich aber dann in allen übrigen Punkten stetig. Offenbar kann dieser Radius durch eine ganz beliebige vom Mittelpunkte nach der Peripherie gezogene Linie $l$ ersetzt werden, so dass die Function $\log(z - z_0)$ beim Übertritt des Punktes $O$ von der negativen (d.~h. wo nach Art.~8 $p$ negativ wird) auf die positive Seite dieser Linie eine plötzliche Verminderung um $2\pi i$ erleidet, übrigens aber sich mit dessen Lage im ganzen Kreise $\mathfrak{S}$ stetig ändert. Nehmen wir nun die complexe Function $\alpha + \beta i$ von $x$, $y$ im Kreise $\mathfrak{S} = \log(z - z_0)$, ausserhalb desselben aber, indem wir $l$ beliebig bis an den Rand verlängern, so an, dass sie
\begin{enumerate}
\item an der Peripherie von $\mathfrak{S}$ $= \log(z - z_0)$, am Rande von $T$ bloss imaginär wird,
\item beim Übertritt von der negativen auf die positive Seite der Linie $l$ sich um $-2\pi i$, sonst aber bei jeder unendlich kleinen Ortsänderung um eine unendlich kleine Grösse von derselben Ordnung ändert,
\end{enumerate}
was immer möglich sein wird, so erhält das Integral
\[
\int \left[\left(\frac{\partial \alpha}{\partial x} - \frac{\partial \beta}{\partial y}\right)^2 + \left(\frac{\partial \alpha}{\partial y} + \frac{\partial \beta}{\partial x}\right)^2\right] dT
\]
über $\mathfrak{S}$ ausgedehnt den Werth Null, über den ganzen übrigen Theil erstreckt einen endlichen Werth, und es kann daher $\alpha + \beta i$ durch Hinzufügung einer bis auf einen bloss imaginären constanten Rest bestimmten stetigen Function von $x$, $y$, welche am Rande bloss imaginär ist, in eine Function $t = m + ni$ von $z$ verwandelt werden. Der reelle Theil $m$ dieser Function wird am Rande $= 0$, im Punkte $O_0 = -\infty$ und ändert sich im ganzen übrigen $T$ stetig. Für jeden zwischen $0$ und $-\infty$ liegenden Werth $a$ von $m$ zerfällt daher $T$ durch eine Linie, wo $m = a$ ist, in Theile, wo $m < a$ ist und die $O_0$ im Innern enthalten, einerseits und andererseits in Theile, wo $m > a$ ist und deren Begrenzung theils durch den Rand von $T$, theils durch Linien, wo $m = a$ ist, gebildet wird. Die Ordnung des Zusammenhangs der Fläche $T$ wird durch diese Zerfällung entweder nicht geändert oder erniedrigt, die Fläche zerfällt daher, da diese Ordnung $= -1$ ist, entweder in zwei Stücke von der Ordnung des Zusammenhangs $0$ und $-1$, oder in mehr als zwei Stücke. Letzteres aber ist unmöglich, weil dann wenigstens in Einem dieser Stücke $m$ überall endlich und stetig und in allen Theilen der Begrenzung constant sein müsste, folglich entweder in einem Flächentheil einen constanten Werth, oder irgendwo --- in einem Punkte oder längs einer Linie --- einen Maximum- oder Minimumwerth haben müsste, gegen Art.~11, III. Die Punkte, wo $m$ constant ist, bilden also in sich zurücklaufende allenthalben einfache Linien, welche ein den Punkt $O_0$ einschliessendes Stück begrenzen, und zwar nimmt $m$ nach Innen zu nothwendig ab, woraus folgt, dass bei einem positiven Umlaufe (wo nach Art.~8 $s$ wächst) $n$ soweit es stetig ist, stets zunimmt, und also, da es nur beim Uebertritt von der negativen auf die positive Seite der Linie $l$ eine plötzliche Aenderung um $-2\pi$) erleidet, jedem Werth zwischen $0$ und $2\pi$ Einmal (von einem Vielfachen von $2\pi$ abgesehen) gleich wird.

Setzen wir nun $e^t = w$, so werden $e^m$ und $n$ Polarcoordinaten des Punktes $Q$ in Bezug auf den Mittelpunkt des Kreises $K$. Die Gesammtheit der Punkte $Q$ bildet dann offenbar eine über $K$ allenthalben einfach ausgebreitete Fläche $S$; der Punkt $Q_0$ derselben fällt auf den Mittelpunkt des Kreises; der Punkt $Q'$ aber kann vermittelst der in $n$ noch verfügbaren Constante auf einen beliebig gegebenen Punkt der Peripherie gerückt werden, w.~z.~b.~w.

In dem Falle, wo der Punkt $O_0$ ein Windungspunkt $(n - 1)$ter Ordnung ist, gelangt man, wenn nur $\log(z - z_0)$ durch $-\log(z - z_0)$ ersetzt wird, durch ganz ähnliche Schlüsse zum Ziele, deren weitere Ausführung man indess aus Art.~14 leicht ergänzen wird.

\subsection*{22.}
\addcontentsline{toc}{subsection}{Art. 22}

Die vollständige Durchführung der Untersuchung des vorigen Artikels für den allgemeinern Fall, wo Einem Punkte der einen Fläche mehrere Punkte der andern entsprechen sollen, und ein einfacher Zusammenhang für dieselben nicht vorausgesetzt wird, unterlassen wir hier, zumal da, aus geometrischem Gesichtspunkte aufgefasst, unsere ganze Untersuchung sich in einer allgemeinern Gestalt hätte führen lassen. Die Beschränkung auf ebene, einzelne Punkte ausgenommen, schlichte Flächen, ist nämlich für dieselbe nicht wesentlich; vielmehr gestattet die Aufgabe, eine beliebig gegebene Fläche auf einer andern beliebig gegebenen in den kleinsten Theilen ähnlich abzubilden, eine ganz ähnliche Behandlung. Wir begnügen uns, hierüber auf zwei Gauss'sche Abhandlungen, die zu Art.~3 citirte und die \textit{disquis.\ gen.\ circa superf.\ art.~13}, zu verweisen.

\bigskip
\begin{center}
\rule{0.5\textwidth}{0.4pt}
\end{center}

\vspace{1em}
% ============================================================
% Inhalt (Table of Contents)
% ============================================================

\section*{Inhalt.\footnote{Diese Inhaltsübersicht rührt fast vollständig von Riemann her.}}
\addcontentsline{toc}{section}{Inhalt}

\begin{center}
\begin{tabular}{p{0.8\textwidth}r}
1. Eine veränderliche complexe Grösse $w = u + vi$ heisst eine Function einer andern veränderlichen Grösse $z = x + yi$, wenn sie mit ihr sich so ändert, dass $\frac{dw}{dz}$ von $dz$ unabhängig ist. Diese Definition wird begründet durch die Bemerkung, dass dies immer stattfindet, wenn die Abhängigkeit der Grösse $w$ von $z$ durch einen analytischen Ausdruck gegeben ist \dotfill & 3\\[0.5em]
2. Die Werthe der veränderlichen complexen Grössen $z$ und $w$ werden dargestellt durch die Punkte $O$ und $Q$ zweier Ebenen $A$ und $B$, ihre Abhängigkeit von einander als eine Abbildung der einen Ebene auf die andere \dotfill & 5\\[0.5em]
3. Ist die Abhängigkeit eine solche (Art.~1), dass $\frac{dw}{dz}$ von $dz$ unabhängig ist, so findet zwischen dem Original und seinem Bilde Aehnlichkeit in den kleinsten Theilen statt \dotfill & 5\\[0.5em]
4. Die Bedingung, dass $\frac{dw}{dz}$ von $dz$ unabhängig ist, ist identisch mit folgenden: $\frac{\partial u}{\partial x} = \frac{\partial v}{\partial y}$, $\frac{\partial u}{\partial y} = -\frac{\partial v}{\partial x}$. Aus ihnen folgen $\frac{\partial^2 u}{\partial x^2} + \frac{\partial^2 u}{\partial y^2} = 0$, $\frac{\partial^2 v}{\partial x^2} + \frac{\partial^2 v}{\partial y^2} = 0$ \dotfill & 5\\[0.5em]
5. Als Ort des Punktes $O$ wird für die Ebene $A$ eine begrenzte über dieselbe ausgebreitete Fläche $T$ substituirt. Windungspunkte dieser Fläche \dotfill & 7\\[0.5em]
6. Ueber den Zusammenhang einer Fläche \dotfill & 9\\[0.5em]
7. Das Integral $\int \left(\frac{\partial X}{\partial x} + \frac{\partial Y}{\partial y}\right) dT$, durch die ganze Fläche $T$ erstreckt, ist gleich $-\int (X \cos\xi + Y \cos\eta)\, ds$ durch ihre ganze Begrenzung, wenn $X$ und $Y$ beliebige in allen Punkten von $T$ stetige Functionen von $x$ und $y$ sind \dotfill & 12\\[0.5em]
8. Einführung der Coordinaten $s$ und $p$ des Punktes $O$ in Bezug auf eine beliebige Linie. Die gegenseitige Abhängigkeit des Vorzeichens von $ds$ und $dp$ wird so festgesetzt, dass $\frac{ds}{dx} = \frac{dp}{dy}$ ist \dotfill & 14\\[0.5em]
9. Anwendung des Satzes im Art.~7, wenn in allen Flächentheilen $\frac{\partial X}{\partial x} + \frac{\partial Y}{\partial y} = 0$ ist \dotfill & 15\\[0.5em]
10. Bedingungen, unter welchen im Innern einer $A$ einfach bedeckenden Fläche $T$ eine Function $u$, welche, allgemein zu reden, der Gleichung $\frac{\partial^2 u}{\partial x^2} + \frac{\partial^2 u}{\partial y^2} = 0$ genügt, nebst allen ihren Differentialquotienten überall endlich und stetig ist \dotfill & 18\\[0.5em]
\end{tabular}
\end{center}

\newpage
\begin{center}
\begin{tabular}{p{0.8\textwidth}r}
11. Eigenschaften einer solchen Function \dotfill & 21\\[0.5em]
12. Bedingungen, unter welchen im Innern einer $A$ einfach bedeckenden einfach zusammenhängenden Fläche $T$ eine Function $w$ von $z$ überall nebst allen ihren Differentialquotienten endlich und stetig ist \dotfill & 23\\[0.5em]
13. Unstetigkeiten einer solchen Function in einem inneren Punkte \dotfill & 24\\[0.5em]
14. Ausdehnung der Sätze der Art.~12 und 13 auf Punkte im Innern einer beliebigen ebenen Fläche \dotfill & 25\\[0.5em]
15. Allgemeine Eigenschaften der Abbildung einer in der Ebene $A$ ausgebreiteten Fläche $T$ auf eine in der Ebene $B$ ausgebreitete Fläche $S$, durch welche die Werthe einer Function $w$ von $z$ geometrisch dargestellt werden \dotfill & 28\\[0.5em]
16. Das Integral $\int \left[\left(\frac{\partial \mu}{\partial x}\right)^2 + \left(\frac{\partial \mu}{\partial y}\right)^2\right] dT$, durch die ganze Fläche $T$ erstreckt, erhält bei Aenderung von $\mu$ um stetige oder doch nur in einzelnen Punkten unstetige Functionen, die am Rande $= 0$ sind, immer für Eine einen Minimumswerth und wenn man durch Abänderung in einzelnen Punkten hebbare Unstetigkeiten ausschliesst, nur für Eine \dotfill & 30\\[0.5em]
17. Begründung eines im vorigen Art.\ vorausgesetzten Satzes mittelst der Grenzmethode \dotfill & 31\\[0.5em]
18. Ist in einer beliebigen zusammenhängenden, durch Querschnitte in eine einfach zusammenhängende $T^*$ zerlegten ebenen Fläche $T$ eine Function $\alpha + \beta i$ von $x$, $y$ gegeben, für welche
\[
\int \left[\left(\frac{\partial \alpha}{\partial x} - \frac{\partial \beta}{\partial y}\right)^2 + \left(\frac{\partial \alpha}{\partial y} + \frac{\partial \beta}{\partial x}\right)^2\right] dT
\]
durch die ganze Fläche erstreckt, endlich ist, so kann sie immer und nur auf eine Art in eine Function von $z$ verwandelt werden durch Hinzufügung einer Function $\mu + \nu i$ von $x$, $y$, welche so bedingt ist: 1)~$\mu$ ist am Rande $= 0$, $\nu$ in Einem Punkte gegeben. 2)~Die Aenderungen von $\mu$ sind in $T$, die von $\nu$ in $T^*$ nur in einzelnen Punkten und nur so unstetig, dass
\[
\int \left[\left(\frac{\partial \mu}{\partial x}\right)^2 + \left(\frac{\partial \mu}{\partial y}\right)^2\right] dT \quad \text{und} \quad \int \left[\left(\frac{\partial \nu}{\partial x}\right)^2 + \left(\frac{\partial \nu}{\partial y}\right)^2\right] dT
\]
durch die ganze Fläche endlich bleiben und letztere an den Querschnitten beiderseits gleich \dotfill & 33\\[0.5em]
19. Ueberschlag über die hinreichenden und nothwendigen Bedingungen zur Bestimmung einer Function complexen Arguments innerhalb eines gegebenen Grössengebiets \dotfill & 35\\[0.5em]
20. Die frühere Bestimmungsweise einer Function durch Grössenoperationen enthält überflüssige Bestandtheile. Durch die hier durchgeführten Betrachtungen ist der Umfang der Bestimmungsstücke einer Function auf das nothwendige Mass zurückgeführt \dotfill & 37\\[0.5em]
21. Zwei gegebene einfach zusammenhängende Flächen können stets so auf einander bezogen werden, dass jedem Punkte der einen Ein mit ihm stetig fortrückender Punkt der andern entspricht und ihre entsprechenden kleinsten Theile ähnlich sind; und zwar kann zu Einem inneren Punkt und zu Einem Begrenzungspunkt der entsprechende beliebig gegeben werden. Dadurch ist für alle Punkte die Beziehung bestimmt \dotfill & 39\\[0.5em]
22. Schlussbemerkungen \dotfill & 42\\
\end{tabular}
\end{center}

\vspace{2em}
\begin{center}
\rule{0.5\textwidth}{0.4pt}
\end{center}
\vspace{1em}

% ============================================================
% Anmerkungen (zu Paper I)
% ============================================================

\section*{Anmerkungen.}
\addcontentsline{toc}{section}{Anmerkungen (zu Paper I)}

\textbf{(1)} (zu Seite 3.) In Riemann's Papieren findet sich der folgende an diese Stelle gehörige Zusatz:

\smallskip
\textit{``Unter dem Ausdruck: die Grösse $w$ ändert sich stetig mit $z$ zwischen den Grenzen $z = a$ und $z = b$ verstehen wir: in diesem Intervall entspricht jeder unendlich kleinen Aenderung von $z$ eine unendlich kleine Aenderung von $w$ oder, greiflicher ausgedrückt: für eine beliebig gegebene Grösse $\varepsilon$ lässt sich stets die Grösse $\alpha$ so annehmen, dass innerhalb eines Intervalls für $z$, welches kleiner als $\alpha$ ist, der Unterschied zweier Werthe von $w$ nie grösser als $\varepsilon$ ist. Die Stetigkeit einer Function führt hiernach, auch wenn dies nicht besonders hervorgehoben ist, ihre beständige Endlichkeit mit sich.''}
\smallskip

\textbf{(2)} (zu Seite 8.) Wenn hier nicht ein Versehen vorliegt, so ist der Ausdruck ``von der Linken zur Rechten'' in einer der gewöhnlichen entgegengesetzten Bedeutung gebraucht, wonach der Sinn des Umkreisens vom Standpunkt eines im Mittelpunkt aufgestellten den kreisenden Punkt mit den Augen verfolgenden Beobachters beurtheilt wird.

\textbf{(3)} (zu Seite 18.) Zur Erläuterung dieser im Ausdruck etwas dunkeln Stelle kann folgendes Beispiel dienen:

In der beistehenden Figur ist $T$ eine dreifach zusammenhängende Fläche. $(ab)$ sei der erste Querschnitt $q_1$, $(cd)$ der zweite $q_2$. Man hat hier drei verschiedene constante Werthdifferenzen der Function
\[
Z = \int \left(\frac{\partial \alpha}{\partial x} - \frac{\partial \beta}{\partial y}\right) ds
\]
zu unterscheiden. Diese seien: an der Strecke $(ad)$: $A$, an der Strecke $(cb)$: $B$, an der Strecke $(cd)$: $C$. Durchläuft man also zuerst $(cd)$, so kann hier $C$ irgend einen Werth haben. Durchläuft man hierauf $(bc)$, so kann hier $B$ einen andern beliebigen Werth haben. An $(ac)$ ist aber hiernach die constante Werthdifferenz $A$ der Function $Z$ völlig bestimmt, nämlich (wenn die Vorzeichen passend bestimmt werden) $A = B + C$. Auf ähnliche Weise schliesst man allgemein, dass, so oft beim Rückwärtsdurchlaufen des Querschnittsystems ein schon durchlaufener Querschnitt einmündet, die Aenderung, welche die constante Werthdifferenz der Function dadurch erfährt, vollkommen bestimmt ist.

\textbf{(4)} (zu Seite 20.) Die Formel
\[
\int \frac{\partial u}{\partial p}\, ds = 0
\]
wird erhalten, wenn man in dem Integral
\[
\int \left( u \frac{\partial u}{\partial p} - u \frac{\partial u}{\partial p} \right) ds
\]
$u = 1$ annimmt, wodurch es, über die Begrenzung eines Flächenstücks ausgedehnt, in dem $u$ die Voraussetzungen des Art.~10 erfüllt, verschwindet.

\textbf{(5)} (zu Seite 30.) Das Beweisverfahren des Art.~16 wird von Riemann später (Theorie der Abel'schen Functionen, Abh.~VI dieser Ausgabe Nr.~3 und Nr.~4, Art.~1) als \textit{Dirichlet'sches Princip} bezeichnet (auf Grund Dirichlet'scher Vorlesungen). Auch Gauss wendet ähnliche Schlüsse an (Allgemeine Lehrsätze in Beziehung auf die im verkehrten Verhältnisse des Quadrats der Entfernung wirkenden Anziehungs- und Abstossungskräfte, Werke Bd.~V). In späterer Zeit ist die Bündigkeit dieser Schlussweise angefochten worden; besonders wird, und mit Recht, die Evidenz der Existenz eines Minimums für das Integral $\Omega$ bestritten. Die Richtigkeit des Satzes selbst, der durch diesen Schluss bewiesen werden soll, der den functionentheoretischen Arbeiten von Riemann ihren eigenthümlich einfachen und allgemeinen Charakter verleiht, ist durch neuere Forschungen auf anderer Grundlage bewiesen. (Vgl. besonders die einschlägenden Arbeiten von H.~A.~Schwarz, Monatsberichte der Berliner Akademie, October 1870, Journal f.~Mathematik Bd.~74, auch gesammelte Abhandlungen, und C.~Neumann, Untersuchungen über das logarithmische und Newton'sche Potential, Leipzig 1877; Vorlesungen über Riemann's Theorie der Abel'schen Integrale, 2.~Auflage, Leipzig 1884.)

\textbf{(6)} (zu Seite 33.) Die folgenden Bemerkungen sind fast wörtlich den in Riemann's handschriftlichem Nachlass gefundenen Entwürfen zu Art.~17 entnommen und dienen theils zur Erläuterung, theils zur Ergänzung der Untersuchung.

Von den Werthen $P_1$ und $P_2$ kann auch einer überall $= 0$ genommen werden, wenn nur $T'$ eine endliche Breite behält, wodurch unser Beweis auf den Fall anwendbar wird, wo die Unstetigkeit längs eines Theils der Begrenzung einträte, oder durch Abänderung von $y$ längs einer Linie im Innern entstanden wäre. Für $m$ ist deshalb nicht geradezu der kleinste Werth von $(Y_1 - Y_2)^2$ in dem angegebenen Intervall von $p_1$ und $p_2$ gesetzt, damit der Beweis auch auf den Fall anwendbar ist, wo $y$ unendlich viele Maxima und Minima, also z.~B. in der Nähe der Unstetigkeitslinie den Werth $\sin \frac{1}{x}$, hätte.

In ähnlicher Weise lässt sich zeigen, dass $L$ über alle Grenzen wächst, wenn $X$ sich einer Function $y$ unbegrenzt nähert, die in einem Punkt $O'$ so unstetig wird, dass in einem Theil einer mit dem Radius $\varrho$ um $O'$ beschriebenen Kreislinie $\varrho \frac{\partial Y}{\partial x}$, $\varrho \frac{\partial Y}{\partial y}$ für ein unendlich kleines $\varrho$ sich einer endlichen Grenze nähern oder unendlich werden.

Es lässt sich in diesem Fall ein Werth $R$ von $\varrho$ so annehmen, dass unterhalb desselben
\[
\left(\frac{\partial Y}{\partial x}\right)^2 + \left(\frac{\partial Y}{\partial y}\right)^2 - \frac{\partial (\varphi + \psi)}{\partial \varphi}
\]
nicht $0$ wird. Bezeichnen wir den kleinsten Werth dieser Grösse in diesem Intervall durch $a$, so wird der Beitrag eines zwischen $\varrho = R$ und $\varrho = r$ (wo $r < R$) enthaltenen Kreisringes zu $L$
\[
> \int_0^R a \frac{d\varrho}{\varrho} = a \log \frac{R}{r}
\]
und folglich, wenn man $r = R e^{-\frac{C}{a}}$ annimmt, $> C$. Wählt man also zur Begrenzung von $T'$ einen Kreis, wo $\varrho < R e^{-\frac{C}{a}}$, so wird der aus dem übrigen $T$ stammende Theil von $L$ und folglich $L$ selbst, wie auch $l$ im Innern des Kreises angenommen werden möge, $< C$.

(Diese Untersuchung bezieht sich zwar zunächst auf einen Punkt, der kein Windungspunkt und kein Begrenzungspunkt ist, erleidet aber eine wesentliche Aenderung nur für einen Begrenzungspunkt, wo die Fläche eine Spitze, d.~h. ihre Begrenzung einen Rückkehrpunkt hat. Die Bestimmung eines Grades der Unstetigkeit, welchen $l$ nicht erreichen kann, beruht indess auch hier auf denselben Principien und wir begnügen uns daher mit der Andeutung dieses Falles.)

Es liefert also, wenn der Flächentheil, wo $l$ und $y$ verschieden sind, unendlich klein wird, im Fall einer Unstetigkeitslinie $T'$ selbst, im Fall eines Unstetigkeitspunktes der übrige Theil von $T$ einen unendlichen Beitrag zu $L$, und unsere Behauptung ist daher, wenn die Unstetigkeit den hier vorausgesetzten Grad erreicht, gerechtfertigt. Ihre Gültigkeit in diesem Umfang genügt für uns und in der That wird sie für leichtere Unstetigkeiten unrichtig, wie z.~B. wenn $y$ in der Entfernung $\varrho$ des Punktes $O$ vom Unstetigkeitspunkt $= \left(\log \frac{1}{\varrho}\right)^\mu$ und $\mu < 1$ ist. Wir geben daher dem ersten Theil des Satzes im Art.~16 folgende Beschränkung: Das Integral $\Omega$ hat, $\omega = \alpha + X$ gesetzt, entweder für eine der Functionen $X$ ein Minimum, oder $l$ nimmt, während $X$ sich einem kleinsten Grenzwerth nähert, doch nur in einzelnen Punkten eine Unstetigkeit an, bei welcher die Ordnung von $\varrho \frac{\partial Y}{\partial x}$, $\varrho \frac{\partial Y}{\partial y}$, wenn sie unendlich werden, die Einheit nicht erreicht.

Eine Unstetigkeit der Function $\omega$, die durch Abänderung eines Werthes in einem Punkt hebbar ist, muss z.~B. eintreten, wenn in der Fläche irgendwo ein Stich, also ein einzelner Begrenzungspunkt, wo $z = 0$ sein müsste, angenommen würde.

\textbf{(7)} (zu Seite 39.) Spätere Untersuchungen haben dargethan, dass die Kraft analytischer Ausdrücke weiter reicht, als es nach diesem Ausspruch von Riemann den Anschein hat. Merkwürdige Beispiele hiervon hat zuerst Seidel gegeben; (Crelles Journal Bd.~73, S.~279), der unter anderem analytische Ausdrücke, die von $z$ abhängen, aufgestellt hat, die in einem Kreis gleich einer beliebigen Function von $z$, ausserhalb $= 0$ sind, oder die überall mit Ausnahme einer Kreisperipherie $= 0$, auf der Kreisperipherie $= 1$ sind. Lässt man bestimmte Integrale zu, so kann man sogar noch weiter gehen und z.~B. $x$ oder $y$ oder $\sqrt{x^2 + y^2}$ als Function von $z = x + yi$ darstellen.

Weierstrass hat gezeigt (zur Functionentheorie, Monatsberichte der Berliner Akademie, August 1880, auch in der Sammlung von Abhandlungen aus der Functionenlehre, Berlin 1886), wie man unendliche Reihen finden kann, deren Glieder rationale Functionen von $z$ sind, die in einer beliebigen Anzahl verschiedener Gebiete der Variablen $z$ verschiedene beliebig gegebene Functionen von $z$ darstellen.

\vspace{2em}
\begin{center}
\rule{0.5\textwidth}{0.4pt}
\end{center}
\vspace{1em}

% ============================================================
% II. Ueber die Gesetze der Vertheilung von Spannungselectricität
% ============================================================

\section*{II. Ueber die Gesetze der Vertheilung von Spannungselectricität\\
in ponderabeln Körpern, wenn diese nicht als vollkommene\\
Leiter oder Nichtleiter, sondern als dem Enthalten von\\
Spannungselectricität mit endlicher Kraft widerstrebend\\
betrachtet werden.}
\addcontentsline{toc}{section}{II. Ueber die Gesetze der Vertheilung von Spannungselectricität}

\begin{center}
\small
(Amtlicher Bericht über die 31.~Versammlung deutscher Naturforscher\\
und Aerzte zu Göttingen im September 1854.\footnote{Vortrag gehalten am 21.~Sept.~1854.})
\end{center}
\vspace{0.5em}

Mittelst der sinnreichen Werkzeuge für Spannungselectricität, welche Herr Prof.\ Kohlrausch in der gestrigen Sitzung dieser Section erwähnte, hat derselbe auch die Bildung des Rückstandes in der Leydener Flasche und in andern Apparaten zur Bindung von Electricität untersucht. Diese Erscheinung ist im Wesentlichen folgende: Wenn man eine Leydener Flasche, nachdem sie längere Zeit geladen gestanden hat, entladet und sie dann eine Zeit lang isolirt stehen lässt, so tritt nach einiger Zeit eine merkliche Ladung wieder auf. Sie führt zu der Annahme, dass bei der ersten Entladung nur ein Theil der geschiedenen Electricitätsmenge sich wieder vereinigte, ein Theil aber in der Flasche zurückblieb. Den ersten Theil nennt man die \textit{disponible} Ladung, den zweiten den \textit{Rückstand}. Die Genauigkeit der Messungen, welche Herr Prof.~Kohlrausch über das Sinken der disponibeln Ladung und über das Wiederauftreten des Rückstandes angestellt hat, reizte mich, an derselben ein aus andern Gründen wahrscheinliches Gesetz zu prüfen, welches eine in der bisherigen Theorie der Spannungselectricität vorhandene Lücke ausfüllt.

Bekanntlich beziehen sich die mathematischen Untersuchungen über Spannungselectricität auf ihre Vertheilung in vollkommenen und völlig isolirten Leitern; man betrachtet also die ponderabeln Körper entweder als absolute Leiter oder als absolute Nichtleiter. Eine Folge davon ist, dass nach dieser Theorie sich beim Gleichgewicht die gesammte Spannungselectricität nur an den Grenzflächen der Leiter und Isolatoren ansammelt. Zugestandenermassen aber ist dies eine blosse Fiction. In der Natur wird es weder einen Körper geben, in welchen durchaus keine Spannungselectricität eindringen kann, noch einen Körper, in welchem sich die gesammte Spannungselectricität auf eine mathematische Fläche zusammenziehen kann. Man muss vielmehr annehmen, dass die ponderabeln Körper dem Aufnehmen oder dem Enthalten von Spannungselectricität mit endlicher Kraft widerstreben, und zwar ist die Annahme, deren Consequenzen sich der Erfahrung gemäss zeigen, die, dass sie nicht dem electrisch Werden oder dem Aufnehmen von Spannungselectricität, sondern dem electrisch Sein oder dem Enthalten von Spannungselectricität widerstreben. Das Gesetz dieses Widerstrebens ist, je nach der dualistischen oder unitarischen Vorstellungsart, folgendes. Nach der dualistischen Vorstellungsart, nach welcher die Spannungselectricität der Ueberschuss der positiven Electricität über die negative ist, muss man in jedem Punkte des ponderabeln Körpers eine Ursache annehmen, welche mit einer der Dichtigkeit dieses Ueberschusses proportionalen Intensität die Dichtigkeit der Electricität gleichen Zeichens --- derjenigen, welche im Ueberschuss vorhanden ist --- zu vermindern und die der entgegengesetzten zu vermehren strebt. Nach der unitarischen Auffassungsweise, nach welcher die Spannungselectricität der Ueberschuss der in dem Körper enthaltenen Electricität über die ihm natürliche ist, muss man in jedem Punkte desselben eine Ursache annehmen, welche mit einer der Dichtigkeit dieses Ueberschusses proportionalen Intensität die Dichtigkeit der Electricität zu vermindern oder bei negativem Ueberschuss zu vermehren strebt. Ausser dieser Bewegungsursache hat man nun, wenn keine merklichen thermischen oder magnetischen oder voltainductorischen Wirkungen und Einflüsse stattfinden, und die ponderabeln Körper gegen einander ruhen, nur noch die dem Coulomb'schen Gesetz gemässe electromotorische Kraft in Rechnung zu ziehen. Unter denselben Umständen kann man für die Abhängigkeit der erfolgten Bewegung von den Bewegungsursachen Proportionalität zwischen electromotorischer Kraft und Stromintensität annehmen.

Um diese Bewegungsgesetze in Formeln auszudrücken, seien $x$, $y$, $z$ rechtwinklige Coordinaten und im Punkte $(x, y, z)$ zur Zeit $t$ die Dichtigkeit der Spannungselectricität $\varrho$, und $u$ der $4\pi$te Theil des Potentials der gesammten Spannungselectricität nach Gauss'scher Definition, nach welcher das Potential in einem bestimmten Punkte gleich ist dem Integral über sämmtliche Massen Spannungselectricität, jede dividirt durch die Entfernung von diesem Punkte. Die dem Coulomb'schen Gesetz gemässe electromotorische Kraft ist dann, nach den Richtungen der drei Axen zerlegt, proportional
\[
\frac{\partial u}{\partial x}, \quad \frac{\partial u}{\partial y}, \quad \frac{\partial u}{\partial z},
\]
die von der Reaction des ponderabeln Körpers herrührende proportional
\[
\frac{\partial \varrho}{\partial x}, \quad \frac{\partial \varrho}{\partial y}, \quad \frac{\partial \varrho}{\partial z}.
\]
Die Componenten der electromotorischen Kraft können also gleich gesetzt werden
\[
\frac{\partial u}{\partial x} - \beta^2 \frac{\partial \varrho}{\partial x}, \quad
\frac{\partial u}{\partial y} - \beta^2 \frac{\partial \varrho}{\partial y}, \quad
\frac{\partial u}{\partial z} - \beta^2 \frac{\partial \varrho}{\partial z},
\]
wo $\beta^2$ nur von der Natur des ponderabeln Körpers abhängt. Diesen sind nun die Componenten der Stromintensität proportional, sie sind also $= \alpha \xi$, $\alpha \eta$, $\alpha \zeta$, wenn man durch $\xi$, $\eta$, $\zeta$ die Componenten der Stromintensität und durch $\alpha$ eine von der Natur des ponderabeln Körpers abhängige Constante bezeichnet.

Verbindet man hiermit die phoronomische Gleichung
\[
\frac{\partial \varrho}{\partial t} + \frac{\partial \xi}{\partial x} + \frac{\partial \eta}{\partial y} + \frac{\partial \zeta}{\partial z} = 0,
\]
welche man erhält, indem man die in das Raumelement $dx\,dy\,dz$ im Zeitelement $dt$ einströmende Electricitätsmenge auf doppelte Weise ausdrückt, und die Gleichung
\[
\frac{\partial^2 u}{\partial x^2} + \frac{\partial^2 u}{\partial y^2} + \frac{\partial^2 u}{\partial z^2} = -4\pi \varrho,
\]
welche aus dem Begriffe des Potentials folgt, so erhält man, indem man erstere mit $\alpha$ multiplicirt und für $\xi$, $\eta$, $\zeta$ ihre Werthe setzt, die Gleichung
\[
\alpha \frac{\partial \varrho}{\partial t} + \varrho - \beta^2 \left(\frac{\partial^2 \varrho}{\partial x^2} + \frac{\partial^2 \varrho}{\partial y^2} + \frac{\partial^2 \varrho}{\partial z^2}\right) = 0.
\]
Diese giebt für $u$ eine partielle Differentialgleichung, welche in Bezug auf $t$ vom ersten, in Bezug auf die Raumcoordinaten vom vierten Grade ist, und um von einem bestimmten Zeitpunkte an $u$ innerhalb des ponderabeln Körpers allenthalben vollständig zu bestimmen, werden ausser dieser Gleichung in jedem Punkte desselben Eine Bedingung für die Anfangszeit und für die Folge in jedem Oberflächenpunkte zwei Bedingungen erforderlich sein.

Ich werde nun die Consequenzen dieser Gesetze in einigen besonderen Fällen mit der Erfahrung vergleichen.

Für das Gleichgewicht (in einem System isolirter Leiter) ist
\[
\frac{\partial u}{\partial x} - \beta^2 \frac{\partial \varrho}{\partial x} = 0, \quad
\frac{\partial u}{\partial y} - \beta^2 \frac{\partial \varrho}{\partial y} = 0, \quad
\frac{\partial u}{\partial z} - \beta^2 \frac{\partial \varrho}{\partial z} = 0
\]
oder
\[
u + \beta^2 \varrho = \text{Const.},
\]
oder, da $\frac{\partial^2 u}{\partial x^2} + \frac{\partial^2 u}{\partial y^2} + \frac{\partial^2 u}{\partial z^2} = -4\pi \varrho$,
\[
\frac{\partial^2 u}{\partial x^2} + \frac{\partial^2 u}{\partial y^2} + \frac{\partial^2 u}{\partial z^2} - \frac{4\pi}{\beta^2} (u - \text{Const.}) = 0.
\]

Für die Stromausgleichung oder den Beharrungszustand der Vertheilung (im Schliessungsbogen constanter Ketten) ist
\[
\frac{\partial \varrho}{\partial t} = 0,
\]
oder
\[
\varrho - \beta^2 \left(\frac{\partial^2 \varrho}{\partial x^2} + \frac{\partial^2 \varrho}{\partial y^2} + \frac{\partial^2 \varrho}{\partial z^2}\right) = 0.
\]

Wenn nun die Länge $\beta$ gegen die Dimensionen des ponderabeln Körpers sehr klein ist, so nimmt $u - \text{Const.}$ im erstem Falle und $\varrho$ im zweiten von der Oberfläche ab sehr schnell ab und ist im Innern überall sehr klein, und zwar ändern diese Grössen mit dem Abstande $p$ von der Oberfläche nahe wie $e^{-\frac{p}{\beta}}$. Dieser Fall wird bei den metallischen Leitern angenommen werden müssen; wird $\beta = 0$ gesetzt, so erhält man die bekannten Formeln für vollkommene Leiter.

Bei der Anwendung dieser Gesetze auf die Rückstandsbildung in der Leydener Flasche musste ich, da Angaben über die Dimensionen der Apparate fehlten, annehmen, dass die Dimensionen derselben gegen den Abstand der Belegungen als unendlich gross betrachtet werden dürften. Mit der Ausführung der Rechnung wage ich die verehrten Anwesenden nicht zu ermüden und begnüge mich das Resultat derselben anzugeben.

Aus den Messungen des Herrn Prof.~Kohlrausch hatte sich ergeben, dass die disponible Ladung, als Function der Zeit betrachtet, nahe durch eine Parabel dargestellt wird, dass jedoch der Parameter der Parabel, welche sich der Ladungscurve am nächsten anschliesst, langsam abnimmt, so dass wenn man die anfängliche Ladung durch $L_0$, die zur Zeit $t$ durch $L_t$ bezeichnet, $-\frac{dL_t}{\sqrt{t}}$ eine Grösse ist, welche mit wachsendem $t$ allmählich abnimmt.

Dasselbe ergab sich auch aus der Rechnung, wenn angenommen wurde, dass sowohl $\alpha$ als $\beta^2$ beim Glase, wie dies von vorn herein zu erwarten war, sehr gross sei und als unendlich gross betrachtet werden dürfe, während ihr Quotient endlich bleibt. Eine schärfere Vergleichung der Rechnung mit den Beobachtungen habe ich nicht angestellt, namentlich aus dem Grunde, weil mir Angaben über die Dimensionen der Apparate und überhaupt alle Mittel fehlten, die wegen der Abweichungen von den Voraussetzungen der Rechnung nöthigen Correctionen zu bestimmen. Es wäre eine solche namentlich zur Bestimmung der electrischen Constanten des Glases zu wünschen. Doch halte ich das hier aufgestellte Gesetz für die Vertheilung der Spannungselectricität für vollkommen durch die Messungen des Herrn Prof.~Kohlrausch bestätigt.

Ich darf wohl noch in der Kürze die Anwendung dieses Gesetzes auf einen andern Gegenstand besprechen.

Bekanntlich wird die Fortpflanzung der galvanischen Ströme in metallischen Leitern und die in Folge derselben stattfindende Stromausgleichung bei Constanten oder langsam sich ändernden electromotorischen Kräften durch die dabei auftretende Spannungselectricität bewirkt. Dieser Vorgang ist wegen seiner ungemein kurzen Dauer und der hinzukommenden thermischen und magnetischen Wirkungen nur in seinen Resultaten der experimentalen Forschung zugänglich, und die einzigen experimentellen Bestimmungen, welche wir darüber haben, sind die Messungen der Fortpflanzungsgeschwindigkeit in Telegraphendrähten und die Ohm'schen Gesetze der Stromausgleichung. Eine genauere Analyse der Ohm'schen Gesetze führt indess ebenfalls zu der hier gemachten Annahme, und ich wurde in der That dadurch zuerst auf sie geführt.

Ohm bestimmt die Stromvertheilung bei der Stromausgleichung durch folgende zwei Bedingungen:
\begin{enumerate}
\item Um die den wirklich erfolgten Stromintensitäten proportionalen electromotorischen Kräfte zu erhalten, muss man zu den äussern electromotorischen Kräften Kräfte hinzufügen, welche die Differentialquotienten Einer Function des Orts, der \textit{Spannung}, sind.
\item Bei der Stromausgleichung strömt in jeden Theil des ponderabeln Leiters eben so viel Electricität ein als aus.
\end{enumerate}

Ohm glaubte nun, dass die Spannung, diese Function des Orts, von welcher die inneren electromotorischen Kräfte die Differentialquotienten sind, von der Spannungselectricität so abhinge, dass sie ihrer Dichtigkeit proportional sei, welche Annahme in der That das Zustandekommen beider Bedingungen erklärt. Aber es haben schon, fast gleichzeitig, Herr Prof.~Weber\footnote{Abhandlungen d.\ k.\ sächs.\ Ges.\ d.\ W. 1852, I. S.~293.} und Kirchhoff\footnote{Poggendorffs Annalen. Bd.~79, S.~506.} darauf aufmerksam gemacht, dass dann die Electricität im Gleichgewicht sein müsste, wenn sie den ponderabeln Körper mit gleichmässiger Dichtigkeit erfüllte, während sie doch der Erfahrung nach beim Gleichgewicht auf der Oberfläche vertheilt ist. Die Spannung muss eine Function sein, welche beim Gleichgewicht im ganzen Leiter constant ist, und also vielmehr dem Potential der Spannungselectricität proportional sein, und diese innern electromotorischen Kräfte sind mit den dem Coulomb'schen Gesetz gemässen identisch.

Diese Ansicht über die Spannung wurde auch von den meisten Forschern angenommen. Dabei aber blieb es ununtersucht, durch welche Ursachen bei der Stromausgleichung die zweite Bedingung hergestellt wurde, dass in jedem ponderabeln Körpertheil die Electricitätsmenge constant bleibe.

Nach der dualistischen Auffassung muss sowohl die positive als die negative Electricitätsmenge constant bleiben; dass kein merklicher Ueberschuss Einer Electricität sich bilde, scheint man, wenigstens so lange man auf die Grössenverhältnisse nicht näher eingeht, aus der Anziehung der entgegengesetzten Electricitäten nach dem Coulomb'schen Gesetz erklären zu können, und man muss dann noch eine Ursache, dass die neutrale Electricität in jedem Körpertheil constant bleibe, also einen Druck des Ponderabile auf sie, annehmen. Diese Annahme habe ich auf Anregung des Herrn Prof.~Weber schon vor mehreren Jahren der Rechnung zu unterwerfen gesucht, ohne zu einem befriedigenden Resultat zu gelangen.

Nach unitarischer Auffassung bedarf es nur einer Ursache, welche die in einem ponderabeln Körpertheil enthaltene Electricitätsmenge constant zu erhalten strebt. Man wird so geradeswegs zu der obigen Annahme geführt, dass jeder ponderabele Körper Electricität von bestimmter Dichtigkeit zu besitzen strebt und sowohl einem grösseren als einem geringeren erfüllt Sein widerstrebt. Das Gesetz dieses Widerstrebens kann man so annehmen, wie es sich für das Glas durch die Erfahrung bestätigt hat.

Diese Betrachtungen führen also dazu, die ursprüngliche Franklin'sche Auffassung der electrischen Erscheinungen als diejenige anzunehmen, welche man für das tiefere Eindringen in den Zusammenhang dieser Erscheinungen unter sich und mit andern Erscheinungen zu Grunde zu legen und der weitern Aus- und Umbildung nach den Geboten und Winken der Erfahrung zu unterwerfen hat.

Möchten sie in dem Kreise bewährter Forscher, vor denen ich sie zu entwickeln die Ehre hatte, einer nähern Prüfung werth gefunden werden.

\vspace{2em}
\begin{center}
\rule{0.5\textwidth}{0.4pt}
\end{center}
\vspace{1em}

% ============================================================
% III. Zur Theorie der Nobili'schen Farbenringe
% ============================================================

\section*{III. Zur Theorie der Nobili'schen Farbenringe.}
\addcontentsline{toc}{section}{III. Zur Theorie der Nobili'schen Farbenringe}

\begin{center}
\small
(Aus Poggendorff's Annalen der Physik und Chemie. Bd.~95, 28.~März 1855.)
\end{center}
\vspace{0.5em}

Die Nobili'schen Farbenringe bilden ein schätzbares Mittel, die Gesetze der Stromverzweigung in einem durch Zersetzung leitenden Körper experimentell zu studiren. Die Erzeugungsweise dieser Ringe ist folgende. Man übergiesst eine Platte von Platin, vergoldetem Silber oder Neusilber mit einer Auflösung von Bleioxyd in concentrirter Kalilauge und lässt den Strom einer starken galvanischen Batterie durch die Spitze eines feinen in eine Glasröhre eingeschmolzenen Platindrahts in die Flüssigkeitsschicht ein- und durch die Platte austreten. Das Anion, Bleisuperoxyd nach Beetz, lagert sich dann auf der Metallplatte in einer zarten durchsichtigen Schicht ab, welche je nach der Entfernung vom Eintrittspunkte des Stroms verschiedene Dicke besitzt, so dass die Platte nach Entfernung der Flüssigkeit Newtonische Farbenringe zeigt. Aus diesen Farbenringen lässt sich dann die relative Dicke der Schicht in verschiedenen Entfernungen bestimmen und hieraus mittelst des Faraday'schen Gesetzes, nach welchem die Menge der abgeschiedenen Substanz der durchgegangenen Electricitätsmenge allenthalben proportional sein muss, die Stromvertheilung beim Austritt aus der Flüssigkeit ableiten.

Der erste Versuch, die Stromvertheilung durch Rechnung zu bestimmen und das gefundene Resultat mit der Erfahrung zu vergleichen, ist von E.~Becquerel gemacht worden. Derselbe hat vorausgesetzt, dass die Ausdehnung der Flüssigkeitsschicht gegen ihre Dicke als unendlich gross betrachtet werden dürfe, der Strom durch einen Punkt ihrer Oberfläche eintrete und sich nach den Ohm'schen Gesetzen in derselben ausbreite. Er glaubt nun bei diesen Voraussetzungen ohne merklichen Fehler die Strömungscurven als gerade Linien betrachten zu können und leitet aus dieser Annahme das Gesetz ab, dass die Dicke der niedergeschlagenen Schicht dem Abstände vom Eintrittspunkte umgekehrt proportional sein müsste, welches Gesetz er experimentell bestätigt habe.

Herr Du-Bois-Reymond hat dagegen in einem vor der physikalischen Gesellschaft zu Berlin gehaltenen Vortrage gezeigt, dass bei Voraussetzung gerader Strömungslinien die Dicke der in ihrem Endpunkte abgeschiedenen Substanz vielmehr dem Cubus ihrer Länge umgekehrt proportional sich ergiebt und dadurch Herrn Beetz zu einer Reihe von dem Anschein nach bestätigenden Versuchen veranlasst, welche in Poggendorff's Annalen Bd.~71, S.~71 beschrieben sind und viel Vertrauen erwecken.

Die genaue Rechnung indessen lehrt, dass die Voraussetzung gerader Strömungslinien unzulässig ist und ein ganz falsches Resultat liefert. Allerdings sind die Strömungslinien, wenigstens bei grösserer Entfernung ihres Austrittspunktes (da sie zwischen zwei sehr nahen Parallel-Linien liegen und höchstens einen Wendepunkt besitzen), in dem mittleren Theile ihres Laufes in beträchtlicher Ausdehnung sehr wenig gekrümmt; hieraus aber darf man keineswegs schliessen, dass sie ohne merklichen Fehler durch gerade von ihrem Eintrittspunkte nach ihrem Austrittspunkte gehende Linien ersetzt werden können.

Ich werde zunächst die bei genauer Rechnung aus den Voraussetzungen der Herren E.~Becquerel und Du-Bois-Reymond fliessenden Folgerungen entwickeln und schliesslich auf die Versuche des Herrn Beetz zurückzukommen mir erlauben.

Ich nehme an, dass der Eintritt des Stromes in die durch zwei horizontale Ebenen begrenzte Flüssigkeitsschicht in einem Punkte stattfinde, und bezeichne für einen Punkt derselben den Horizontalabstand vom Einströmungspunkt durch $r$, die Höhe über der unteren Grenzfläche durch $z$, die Erhebung seiner Spannung über die Spannung an der oberen Seite dieser Grenzfläche durch $u$.(\footnote{Ferner sei die Stärke des ganzen Stromes $S$, der specifische Leitungswiderstand der Flüssigkeit $W$, im Einströmungspunkt $z = \alpha$, an der Oberfläche $0 = \beta$.}) Es muss nun $u$ als Function von $r$ und $z$ bestimmt werden; die Stromintensität im Punkte $(r, z)$, welcher nach dem Faraday'schen Gesetz die gesuchte Dicke der dort niedergeschlagenen Schicht proportional sein muss, ist dann gleich dem Werthe von $-\frac{1}{W}\frac{\partial u}{\partial z}$ in diesem Punkte.

Wird zunächst vorausgesetzt, dass die Ausdehnung der Flüssigkeitsschicht gegen ihre Dicke als unendlich gross betrachtet werden dürfe, so sind die Bedingungen zur Bestimmung von $u$:

\medskip
\begin{tabular}{@{}l@{\quad}p{0.7\textwidth}@{}}
(1) & für $-\infty < r < \infty$, $0 < z < \beta$:
\[\frac{\partial^2 u}{\partial r^2} + \frac{1}{r}\frac{\partial u}{\partial r} + \frac{\partial^2 u}{\partial z^2} = 0;\] \\[0.5em]
(2) & für $-\infty < r < \infty$, $z = 0$: $u = 0$; \\[0.3em]
(3) & für $-\infty < r < \infty$, $z = \beta$: $\dfrac{\partial u}{\partial z} = 0$; \\[0.5em]
(4) & für $r = \pm\infty$, $0 < z < \beta$: $u$ endlich; \\[0.3em]
(5) & für $r = 0$, $z = \alpha$: $u = \dfrac{1}{WS}\dfrac{1}{\sqrt{rr + (z - \alpha)^2}}$ oder $= \dfrac{1}{WS}\dfrac{1}{\sqrt{rr + (z - \alpha)^2}} + \text{einer stetigen Function von } r, z$,
\end{tabular}
\medskip

\noindent je nachdem der Einströmungspunkt im Innern oder in der Oberfläche liegt.

Diesen Bedingungen genügt
\[
u = \frac{1}{WS} \sum_{m=-\infty}^{+\infty} (-1)^m \left( \frac{1}{\sqrt{rr + (z + 2m\beta - \alpha)^2}} + \frac{1}{\sqrt{rr + (z + 2m\beta + \alpha)^2}} \right)
\]
oder wenn man zur Vereinfachung $S = \dfrac{1}{W}$ annimmt:
\[
u = \sum_{m=-\infty}^{+\infty} (-1)^m \left( \frac{1}{\sqrt{rr + (z + 2m\beta - \alpha)^2}} + \frac{1}{\sqrt{rr + (z + 2m\beta + \alpha)^2}} \right).
\]

Setzt man $u = a_1 \sin \frac{\pi z}{2\beta} + a_2 \sin 2\frac{\pi z}{2\beta} + a_3 \sin 3\frac{\pi z}{2\beta} + \cdots$, so wird für ein gerades $m$ der Coefficient $a_m = 0$ und für ein ungerades
\[
a_m = \frac{2}{\beta} \int_0^\beta u \sin m\frac{\pi z}{2\beta}\, dz
\]
\[
= \frac{2}{\beta} \int_0^\beta \sum_{-\infty}^{+\infty} (-1)^m \left( \frac{1}{\sqrt{rr + (z + 2m\beta - \alpha)^2}} + \frac{1}{\sqrt{rr + (z + 2m\beta + \alpha)^2}} \right) \sin m\frac{\pi z}{2\beta}\, dz
\]
\[
= 2 \sin m\frac{\pi \alpha}{2\beta} \int_{-\infty}^{+\infty} \frac{\sin m\frac{\pi t}{2\beta}\, dt}{\sqrt{rr + tt}}.
\]

In letzterem Integral kann statt $\int_{-\infty}^{+\infty}$ auch $2\int_0^{+\infty}$ geschrieben werden.

Führt man für $t$ als Veränderliche $t\tau$ ein, so erhält man

\[
\frac{\beta}{2} a_m = 2 \sin m\frac{\pi \alpha}{2\beta} \int_0^\infty \frac{\sin m\frac{\pi \tau}{2\beta}\, d\tau}{\sqrt{rr + \tau^2}},
\]
also
\[
u = \frac{4}{\beta} \sum_n \sin n\frac{\pi \alpha}{2\beta} \sin n\frac{\pi z}{2\beta} \int_0^\infty \frac{\sin n\frac{\pi t}{2\beta}\, dt}{\sqrt{rr + tt}},
\]
über alle positiven ungeraden Werthe von $n$ ausgedehnt.

Nimmt man an, dass die Flüssigkeit bei $r = c$ begrenzt sei und zwar beispielshalber durch einen Nichtleiter, so muss für $r = c$ $\frac{\partial u}{\partial r} = 0$ werden und also zu dem oben erhaltenen Werth von $u$, der durch $u'$ bezeichnet werden möge, noch eine Function $u''$ hinzugefügt werden, welche folgenden Bedingungen genügt:

\medskip
\begin{tabular}{@{}l@{\quad}p{0.7\textwidth}@{}}
(1) & für $-c < r < c$, $0 < z < \beta$:
\[\frac{\partial^2 u''}{\partial r^2} + \frac{1}{r}\frac{\partial u''}{\partial r} + \frac{\partial^2 u''}{\partial z^2} = 0;\] \\[0.5em]
(2) & für $-c < r < c$, $z = 0$: $u'' = 0$; \\[0.3em]
(3) & für $-c < r < c$, $z = \beta$: $\dfrac{\partial u''}{\partial z} = 0$; \\[0.5em]
(4) & für $r = \pm c$, $0 < z < \beta$: $\dfrac{\partial u''}{\partial r} = -\dfrac{\partial u'}{\partial r}$,
\end{tabular}
\medskip

\noindent und überall stetig ist.

Den Bedingungen (1) bis (3) zufolge muss $u''$ ebenfalls in der Form
\[
u'' = b_1 \sin \frac{\pi z}{2\beta} + b_3 \sin 3\frac{\pi z}{2\beta} + b_5 \sin 5\frac{\pi z}{2\beta} + \cdots
\]
darstellbar sein, und zwar fliesst aus (1) für $b_n$ die Bedingung
\[
\frac{\partial^2 b_n}{\partial r^2} + \frac{1}{r}\frac{\partial b_n}{\partial r} - \frac{nn\pi\pi}{4\beta\beta} b_n = 0.
\]
Eine particuläre Lösung dieser Gleichung ist, wie schon bekannt, $\dfrac{1}{\sqrt{rr - 1}}$; eine andere erhält man, wenn man dasselbe Integral zwischen $-1$ und $1$ nimmt; die allgemeinste ist also, wenn $c_n$ und $\gamma_n$ Constanten bedeuten,
\[
b_n = c_n \int_1^\infty \frac{e^{-\frac{n\pi r q}{2\beta}}\, dq}{\sqrt{qq - 1}} + \gamma_n \int_{-1}^{+1} \frac{e^{-\frac{n\pi r q}{2\beta}}\, dq}{\sqrt{1 - qq}}.
\]
oder wenn man
\[
\int_1^\infty \frac{e^{-q^2 t}\, dt}{\sqrt{t-1}} \text{ durch } f(q), \quad \int_{-1}^{+1} \frac{e^{-q^2 t}\, dt}{\sqrt{1-tt}} \text{ durch } \varphi(q)
\]
bezeichnet:
\[
b_n = c_n f\left(\frac{n\pi r}{2\beta}\right) + \gamma_n \varphi\left(\frac{n\pi r}{2\beta}\right).
\]

Die Entwicklung nach steigenden Potenzen von $q$ giebt
\[
\varphi(q) = \pi \sum_{0,\infty} \frac{q^{2m}}{\Pi(m)\Pi(m)},
\]
\[
f(q) = -\varphi(q) \log q + \pi \sum_{0,\infty} W(m) \frac{q^{2m}}{\Pi(m)\Pi(m)};
\]
es wird also $f(q)$ für $q = 0$ unendlich und damit $u''$ für $r = 0$ stetig bleibe, muss $c_n = 0$ sein; $\gamma_n$ ergiebt sich dann aus (4) gleich
\[
\gamma_n = \frac{4 \sin n\frac{\pi \alpha}{2\beta} \cdot a}{\varphi\left(\frac{n\pi c}{2\beta}\right)},
\]
mithin
\[
u'' = \frac{4}{\beta} \sum_n \sin n\frac{\pi \alpha}{2\beta} \sin n\frac{\pi z}{2\beta} \cdot a \cdot \frac{\varphi\left(\frac{n\pi r}{2\beta}\right)}{\varphi\left(\frac{n\pi c}{2\beta}\right)}
\]
über alle positiven ungeraden Werthe von $n$ ausgedehnt.

Zur Berechnung von $f(q)$ und $\varphi(q)$ können für grosse Werthe von $q$ die halbconvergenten Reihen
\[
f(q) = \frac{e^{-2q}}{\sqrt{2q}} \sum_{0,\infty} \frac{(1 \cdot 3 \cdots 2m - 1)^2}{m! (16q)^m},
\]
\[
\varphi(q) = \frac{e^{2q}}{\sqrt{2q}} \sum_{0,\infty} (-1)^m \frac{(1 \cdot 3 \cdots 2m - 1)^2}{m! (16q)^m}
\]
benutzt werden, welche indess ihren Werth nur bis auf Bruchtheile von der Ordnung der Grösse $e^{-4q}$ geben; genügt diese Genauigkeit nicht, so ist es wohl am zweckmässigsten die Entwicklungen nach steigenden Potenzen von $q$ anzuwenden.

Für hinreichend grosse Werthe von $\dfrac{c}{r}$ erhält man also mit Vernachlässigung von Grössen von der Ordnung der Grösse $e^{-\frac{4\pi r}{2\beta}}$:

\begin{multline*}
-\frac{\partial u}{\partial z} = \frac{4}{\beta} \sum_n n \sin n\frac{\pi \alpha}{2\beta} \cos n\frac{\pi z}{2\beta} \cdot a \cdot \frac{e^{-\frac{n\pi(c-r)}{2\beta}}}{e^{-\frac{n\pi c}{2\beta}} \sqrt{\frac{c}{r}}}\\
\times \sum_{0,\infty} \frac{(1 \cdot 3 \cdots 2m - 1)^2}{m! (2m - 1)^2} \left(\frac{\beta}{4\pi r}\right)^m \left(\frac{n}{2}\right)^{2m-1},
\end{multline*}
und die Dicke der Schicht proportional $\left(\dfrac{\partial u}{\partial z}\right)_0$ oder proportional
\[
\frac{1}{\sqrt{r}} \sum_{0,\infty} \frac{(1 \cdot 3 \cdots 2m - 1)^2}{m! (47rc)^m} \left(\frac{\beta}{r}\right)^m \sum_n n^{2m} \sin n\frac{\pi \alpha}{2\beta} \cdot e^{-\frac{n\pi(c-r)}{2\beta}}.
\]

Dieses Resultat bleibt im Allgemeinen auch richtig, wenn statt des Einströmungspunktes eine beliebige Umdrehungsfläche als Kathode angenommen wird; denn für Werthe von $r$ zwischen $c$ und demjenigen Werthe, bis zu welchem die Bedingungen (1) bis (3) gültig bleiben, muss $u$ auch dann durch eine Reihe von der Form
\[
u = \sum_n b_n \sin n\frac{\pi z}{2\beta} \cdot \frac{\varphi\left(\frac{n\pi r}{2\beta}\right)}{\varphi\left(\frac{n\pi c}{2\beta}\right) - f\left(\frac{n\pi c}{2\beta}\right) \cdot \dfrac{f\left(\frac{n\pi a}{2\beta}\right)}{\varphi\left(\frac{n\pi a}{2\beta}\right)}}
\]
dargestellt werden. Eine Ausnahme würde nur eintreten, wenn $\dfrac{f\left(\frac{n\pi a}{2\beta}\right)}{\varphi\left(\frac{n\pi a}{2\beta}\right)} = 0$ würde.

Die von Herrn E.~Becquerel gemachte und von Herrn Du-Bois-Reymond im Wesentlichen beibehaltene specielle Voraussetzung ist die, dass die Kathode ein Punkt der Oberfläche, also $\alpha = \beta$ sei; in diesem Falle ist, wie die geführte Rechnung zeigt, die Dicke der Schicht für grosse Werthe von $\dfrac{c}{r}$ weder der Entfernung vom Einströmungspunkte, wie Herr Becquerel, noch ihrem Cubus, wie Herr Du-Bois-Reymond gefunden hat, umgekehrt proportional, sondern sie nimmt mit wachsenden $\dfrac{c}{r}$ vielmehr ab, wie eine Potenz mit dem Exponenten $-\dfrac{1}{2}$, so dass $\left(\dfrac{\partial u}{\partial z}\right)_0 \sqrt{r}$ sich einem festen Grenzwerthe schliesslich bis zu jedem Grade nähert. Dagegen ist das Gesetz des Herrn Du-Bois-Reymond nicht bloss näherungsweise für grosse Werthe von $\dfrac{c}{r}$, sondern streng richtig, wenn $\beta = \infty$ ist, da sich alsdann
\[
u = \sum_{-\infty}^{+\infty} (-1)^m \left( \frac{1}{\sqrt{rr + (z - 2m\beta - \alpha)^2}} - \frac{1}{\sqrt{rr + (z - 2m\beta + \alpha)^2}} \right)
\]
auf
\[
u = \frac{1}{\sqrt{rr + (z - \alpha)^2}}
\]
und folglich
\[
\left(\frac{\partial u}{\partial z}\right)_0 = \frac{\alpha}{\sqrt{rr + \alpha\alpha}}
\]
reducirt. Die Vermuthung aber, aus welcher derselbe dieses Resultat abgeleitet hat, dass nämlich die Strömungslinien als gerade betrachtet werden dürften, bestätigt sich keineswegs. Die Gleichung der Strömungslinien ist
\[
\int \left(\frac{z\,dr - r\,dz}{z^2 + r^2}\right) = v = \text{const.},
\]
und zwar ist die Constante, multiplicirt mit $\dfrac{1}{2\pi}$, wenn man das Integral so nimmt, dass es für $r = 0$ verschwindet, gleich dem innerhalb der Umdrehungsfläche ($v = \text{const.}$) fliessenden Theile des Stromes. In unserm Falle also sind die Strömungslinien die in der Gleichung
\[
v = \sum_m \frac{z + \alpha}{\sqrt{rr + (z + \alpha)^2}} - \frac{z - \alpha}{\sqrt{rr + (z - \alpha)^2}} = \text{const.}
\]
enthaltenen Linien, welche Linien für alle grösseren Werthe der const. beträchtlich von einer geraden abweichen. Da Herr Du-Bois-Reymond zwar die Annahme macht, dass der Einströmungspunkt in der Oberfläche liege, seine ferneren Schlüsse aber nicht wesentlich auf diese Annahme stützt, so liegt wohl die Vermuthung nahe, dass bei den Versuchen des Herrn Beetz, welche eine nicht zu verkennende Annäherung an das Gesetz der Cuben ergeben, die Forderung des Herrn Du-Bois-Reymond, dass die Oberfläche der Flüssigkeit durch den Einströmungspunkt gehe, nicht berücksichtigt worden ist, sondern dass Herr Beetz, was zweckmässiger sein dürfte, grössere Flüssigkeitsmengen anwandte, so dass in der Reihe für $\left(\dfrac{\partial u}{\partial z}\right)_0$
\[
\sum_m (-1)^m \left( \frac{2m\beta + \alpha}{\sqrt{rr + (2m\beta + \alpha)^2}} - \frac{2m\beta - \alpha}{\sqrt{rr + (2m\beta - \alpha)^2}} \right)
\]
die späteren Glieder oder doch ihre Summe gegen das erste vernachlässigt werden konnten. In diesem Falle würden die hübschen Versuche des Herrn Beetz wirklich als ein Beweis anzusehen sein, dass die Stromvertheilung nahezu nach den vorausgesetzten Gesetzen erfolgt.

Sollte aber diese Vermuthung irrig sein, so wäre aus Herrn Beetz's Versuchen zu schliessen, dass noch andere Umstände bei der Berechnung der Stromvertheilung in Betracht zu ziehen sind, deren Ermittlung einer neuen experimentellen Untersuchung obliegen würde.\footnote{In einer späteren Abhandlung (Poggendorff's Annalen Bd.~95, p.~22) ist Beetz auf diesen Gegenstand zurückgekommen. Es ergiebt sich daraus zunächst, dass bei den Versuchen von Beetz die Einströmungsstelle immer unmittelbar an der Oberfläche der Flüssigkeit lag und mithin die Vermuthung von Riemann irrig ist. Es ist aber gleichwohl nicht nothwendig, nach anderen Umständen zu suchen, welche die Gesetze der Stromvertheilung beeinflussen könnten, da das theoretische Resultat von Riemann mit den Versuchen in noch vollständigerer Uebereinstimmung steht als das von Du-Bois-Reymond, wie aus den in der erwähnten Abhandlung enthaltenen Zusammenstellungen zu ersehen ist.

Am Anfang der achtziger Jahre sind solche Versuche in ausgedehntem Maasse und unter mannigfach veränderten Umständen von Guebhard angestellt worden. Die theoretische Deutung, die Guebhard seinen Versuchen giebt, stimmt mit der von Riemann nicht überein. Es scheint, dass bei diesen Vorgängen ausser der Leitung der Electricität noch andere Einflüsse, besonders die electromotorische Kraft der Polarisation zu berücksichtigen sind. (Compt.~rend. 90, 93, 94. Journal de Physique (2) I, u.~II. und in der Zeitschrift L'Electricien.) W.}

\vspace{2em}
\begin{center}
\rule{0.5\textwidth}{0.4pt}
\end{center}
\vspace{1em}

% ============================================================
% Anmerkungen (zu den Papieren II und III)
% ============================================================

\section*{Anmerkungen.}
\addcontentsline{toc}{section}{Anmerkungen (zu den Papieren II und III)}

\textbf{(1)} (Zu Seite 56.) Es ist hier die Voraussetzung gemacht, dass die Spannung (das Potential) in der begrenzenden Platte constant sei, oder, was damit gleichbedeutend ist, dass die Strömung überall senkrecht gegen die Grenzfläche erfolgt. Da die begrenzende Platte in den Versuchen von Metall ist, dessen Leitungsfähigkeit gegen die der Flüssigkeit ausserordentlich gross ist, so ist diese Annahme statthaft. Ein Bedenken dagegen würde nur dann zu erheben sein, wenn die Metallplatte sehr dünn wäre. (Vgl. eine Mittheilung des Herausgebers ``Ueber stationäre Strömung der Electricität in Platten'' in den Nachrichten der Göttinger Gesellschaft der Wissenschaften 1889 Nr.~6.)

\textbf{(2)} (Zu Seite 58.) Die Umformung
\[
\int_{-\infty}^{+\infty} \frac{e^{-\frac{n\pi t}{2\beta}}\, dt}{\sqrt{rr + tt}} = 2 \int_0^{\infty} \frac{\sin n\frac{\pi t}{2\beta}\, dt}{\sqrt{rr + tt}}
\]
beruht auf dem Satze, dass bei der Integration einer Function complexen Arguments der Integrationsweg geändert werden darf, wenn dabei kein singulärer Punkt überschritten wird.

\textbf{(3)} (Zu Seite 59.) Die Reihenentwicklungen für die Functionen $f(q)$ und $\varphi(q)$ lassen sich auf elementarem Wege folgendermaassen ableiten.

Die beiden Functionen
\[
\int_1^{\infty} \frac{e^{-2qt}\, dt}{\sqrt{tt-1}}, \qquad \int_{-1}^{+1} \frac{e^{-2qt}\, dt}{\sqrt{1-tt}}
\]
sind particuläre Lösungen der Differentialgleichung (auf Seite 58)
\[
\frac{d^2 y}{dq^2} + \frac{1}{q}\frac{dy}{dq} - 4y = 0.
\]
Integrirt man diese Differentialgleichung durch eine nach steigenden Potenzen von $q$ fortschreitende Reihe, so erhält man, wenn man einen constanten Factor durch $\varphi(0) = \pi$ bestimmt,
\[
\varphi(q) = \pi \sum_{0,\infty} \frac{q^{2m}}{\Pi(m)\Pi(m)},
\]
 wenn man sich des Gauss'schen Zeichens $\Pi(m)$ bedient, das für ein ganzzahliges $m$ die Bedeutung $1 \cdot 2 \cdot 3 \cdots m$ oder $m!$ hat.

Eine zweite particuläre Lösung findet sich nun, wenn man
\[
y = u - \frac{1}{\pi} \varphi(q) \log q
\]
setzt und für $u$ eine Potenzreihe von der Form
\[
u = \sum_{0,\infty} a_m \frac{q^{2m}}{\Pi(m)\Pi(m)}
\]
aus der Differentialgleichung ableitet.

Für die Coefficienten $a_m$ erhält man die Recursionsformel
\[
a_m - a_{m-1} = -\frac{1}{m}
\]
und folglich, wenn man den ersten unbestimmt bleibenden Coefficienten $a_0 = W(0)$ setzt,
\[
a_m = W(0) + 1 + \frac{1}{2} + \cdots + \frac{1}{m} = W(m),
\]
worin $W(m)$ die von Gauss eingeführte Function $\dfrac{d\log\Pi(m)}{dm}$ ist. (Gauss.\ Disq.\ circa ser.\ inf.\ Werke Bd.~III, Seite 153.)

Setzt man also
\[
F(q) = \sum_{0,\infty} W(m) \frac{q^{2m}}{\Pi(m)\Pi(m)},
\]
so muss, wenn $A$, $B$ noch zu bestimmende Constanten sind,
\[
f(q) = AF(q) + B\varphi(q)
\]
sein. Die Werthe dieser Constanten erhält man aus $q = 0$. Es ist nämlich
\[
\lim_{q=0} (F(q) + \log q) = W(0),
\]
und durch partielle Integration:
\[
f(q) = \int_1^{\infty} \frac{e^{-2qt}\, dt}{\sqrt{tt-1}} = \int_1^{\infty} e^{-2qt}\, d\log(\sqrt{tt-1} + t)
\]
\[
= -e^{-2q}\log q + 2\int_1^{\infty} e^{-2qt} \log(\sqrt{tt-1} + t)\, dt,
\]
also
\[
\lim_{q=0} (f(q) + \log q) = 2\int_1^{\infty} e^{-2qt} \log 2t\, dt = \int_0^{\infty} e^{-t} \log t\, dt = W(0).
\]
Also erhält man $A = 1$, $B = 0$ und in Uebereinstimmung mit der Formel des Textes $f(q) = F(q)$.

\textbf{(4)} (Zu Seite 59.) Von den beiden nach fallenden Potenzen von $q$ fortschreitenden (halbconvergenten) Reihen erhält man die erste, einschliesslich der Genauigkeitsgrenze, aus dem Integralausdruck
\[
f(q) = \int_1^{\infty} \frac{e^{-2qt}\, dt}{\sqrt{tt-1}} = \frac{e^{-2q}}{\sqrt{2q}} \int_0^{\infty} \frac{e^{-s}\, ds}{\sqrt{s(1+\frac{s}{4q})}}
\]
durch Entwicklung von $\left(1 + \dfrac{s}{4q}\right)^{-\frac{1}{2}}$ nach steigenden Potenzen von $s$ nach dem Taylor'schen Lehrsatz mit Rücksicht auf das Restglied. Die zweite dieser Formeln aber bot Schwierigkeiten, die bereits H.~Hankel in einer Arbeit über die Cylinderfunctionen (Bd.~I der Mathem.\ Annalen) hervorhob. Zur Aufklärung sollen mit Bezugnahme auf eine Abhandlung des Herausgebers (Zur Theorie der Bessel'schen Functionen, Mathem.\ Annalen Bd.~XXXVII, S.~404) die folgenden Bemerkungen hier Platz finden.

Definirt man mit Benutzung einer jetzt üblichen Bezeichnung zwei Functionen des complexen Argumentes $x$
\[
J(x) = \sum_{0,\infty} \frac{(-1)^m x^{2m}}{\Pi(m)\Pi(m)}, \qquad
Y(x) = 2\sum_{0,\infty} \frac{W(m)(-1)^m x^{2m}}{\Pi(m)\Pi(m)},
\]
wobei, um $Y(x)$ eindeutig zu bestimmen, $\log \dfrac{x}{2}$ für reelle positive Werthe von $x$ reell und die $x$-Ebene durch einen längs des negativen Theils der reellen Axe verlaufenden Schnitt begrenzt angenommen sei, so wird, indem wir $q$ als reell und positiv voraussetzen,
\[
\varphi(q) = \pi J(2iq), \qquad f(q) = -\frac{\pi}{2} Y(2iq) + \frac{\pi W(0)}{2} J(2iq).
\]
In der oben erwähnten Arbeit sind zwei Functionen
\[
\sqrt{\frac{\pi}{2}} e^{-q} f(q) = S_2(2iq), \qquad
\frac{1}{\sqrt{2\pi}} e^{q} \varphi(q) = S_1(2iq)
\]
definirt, durch die sich $J(x)$, $Y(x)$ und folglich auch $\varphi(q)$, $f(q)$ ausdrücken lassen:
\[
\sqrt{\frac{\pi}{2}} e^{-q} f(q) = S_2(2iq) - iS_1(2iq) e^{-2q}.
\]
Nun folgt aus Art.~X der genannten Abhandlung, dass, wenn $S_1(2iq)$, $S_2(2iq)$ durch die endlichen Reihen
\[
S_1(2iq) = \sum_{0,n-1} \frac{(1 \cdot 3 \cdots 2m - 1)^2}{m! (16q)^m},
\]
\[
S_2(2iq) = \sum_{0,n-1} (-1)^m \frac{(1 \cdot 3 \cdots 2m - 1)^2}{m! (16q)^m}
\]
ersetzt werden, der Fehler für ein reelles positives $q$, sofern $n < 4q$ und nicht zu klein ist, etwa von der Grösse
\[
\sqrt{\frac{\pi}{2}} \frac{(1 \cdot 3 \cdots 2n - 1)^2}{n! (16q)^n}
\]
oder, wenn man mit der willkürlichen Zahl $n$ möglichst nahe an $4q$ herangeht,
\[
\frac{e^{-4q}}{\sqrt{q}}
\]
ist. Mit gleicher Genauigkeit kann man also auch
\[
\sqrt{\frac{\pi}{2}} e^{-q} f(q), \quad \sqrt{\frac{\pi}{2}} e^{q} \varphi(q)
\]
durch dieselben Ausdrücke darstellen, in Uebereinstimmung mit den Formeln des Textes.

Es darf hier natürlich der Begriff der halbconvergenten Reihe nicht in dem beschränkten Sinne gebraucht werden, wonach die Summe der $n$ ersten Glieder einem bestimmten Werthe so sich nähert, dass der Unterschied immer kleiner ist als das zuletzt hinzugefügte Glied. Eine divergente Reihe mit nur positiven Gliedern kann selbstverständlich jeden beliebigen positiven Werth in dem Sinne genähert darstellen, dass der Werth der Summe den darzustellenden Werth durch Hinzüfügung eines weiteren Gliedes überschreitet, und dass also, wenn man unmittelbar vor oder nach diesem Gliede abbricht, der Fehler kleiner ist als dies Glied, und dann wieder fortwährend wächst. Es kommt also nur darauf an, die am zweckmässigsten anzuwendende Gliederzahl und die ungefähre Grösse des letzten dieser Glieder zu ermitteln. In dem vorliegenden Fall ist dieser Zweck erreicht, wenn man die Gliederzahl $n$ möglichst nahe an $4q$ heranrückt.

\vspace{2em}
\begin{center}
\rule{0.5\textwidth}{0.4pt}
\end{center}
\vspace{1em}

% ============================================================
% IV. Beiträge zur Theorie der durch die Gauss'sche Reihe
%     F(α, β, γ, x) darstellbaren Functionen
% ============================================================

\section*{IV. Beiträge zur Theorie der durch die Gauss'sche Reihe\\
$F(\alpha, \beta, \gamma, x)$ darstellbaren Functionen.}
\addcontentsline{toc}{section}{IV. Beiträge zur Theorie der durch die Gauss'sche Reihe $F(\alpha, \beta, \gamma, x)$ darstellbaren Functionen}

\begin{center}
\small
(Aus dem siebenten Bande der Abhandlungen der Königlichen Gesellschaft\\
der Wissenschaften zu Göttingen. 1857.)
\end{center}
\vspace{0.5em}

Die Gauss'sche Reihe $F(\alpha, \beta, \gamma, x)$, als Function ihres vierten Elements $x$ betrachtet, stellt diese Function nur dar, so lange der Modul von $x$ die Einheit nicht überschreitet. Um diese Function in ihrem ganzen Umfange, bei unbeschränkter Veränderlichkeit dieses ihres Arguments, zu untersuchen, bieten die bisherigen Arbeiten über dieselbe zwei Wege dar. Man kann nämlich entweder von einer linearen Differentialgleichung, welcher sie genügt, ausgehen, oder von ihrem Ausdrucke durch bestimmte Integrale. Jeder dieser Wege gewährt eigenthümliche Vortheile; jedoch ist bis jetzt, in der reichhaltigen Abhandlung von Kummer im 15.\ Bande des mathematischen Journals von Grelle und auch in den noch unveröffentlichten Untersuchungen von Gauss\footnote{Gauss Werke. Bd.~III 1886. S.~207. W.}, nur der erste betreten, wohl hauptsächlich deshalb, weil die Rechnung mit bestimmten Integralen zwischen complexen Grenzen noch zu wenig ausgebildet war, oder doch nicht als einem grossen Leserkreise geläufig vorausgesetzt werden konnte.

In der folgenden Abhandlung habe ich diese Transcendente nach einer neuen Methode behandelt, welche im Wesentlichen auf jede Function, die einer linearen Differentialgleichung mit algebraischen Coefficienten genügt, anwendbar bleibt. Nach derselben lassen sich die früher zum Theil durch ziemlich mühsame Rechnungen gefundenen Resultate fast unmittelbar aus der Definition ableiten, und dies ist in dem hier vorliegenden Theile dieser Abhandlung geschehen, hauptsächlich in der Absicht für die vielfachen Anwendungen dieser Function in physikalischen und astronomischen Untersuchungen eine bequeme Uebersicht über ihre möglichen Darstellungen zu geben. Es ist nöthig, einige allgemeine Vorbemerkungen über die Betrachtung einer Function bei unbeschränkter Veränderlichkeit ihres Arguments voraufzuschicken.

Betrachtet man den Werth der unabhängig veränderlichen Grösse $x = y + zi$ zur leichteren Auffassung ihrer Veränderlichkeit als vertreten durch einen Punkt einer unendlichen Ebene, dessen rechtwinklige Coordination $y$, $z$ sind, und denkt sich die Function $w$ in einem Theile dieser Ebene gegeben, so kann sie von dort aus nach einem leicht zu beweisenden Satze nur auf eine Weise der Gleichung $\frac{\partial w}{\partial x} = \frac{1}{i}\frac{\partial w}{\partial y}$ gemäss stetig fortgesetzt werden. Diese Fortsetzung muss selbstredend nicht in blossen Linien geschehen, worauf eine partielle Differentialgleichung nicht angewandt werden könnte, sondern in Flächenstreifen von endlicher Breite. Bei Functionen, welche, wie die hier zu untersuchende, ``mehrwerthig'' sind oder für denselben Werth von $x$ je nach dem Wege, auf welchem die Fortsetzung geschehen ist, mehrere Werthe annehmen können, giebt es gewisse Punkte der $x$-Ebene, um welche herum sich die Function in eine andere fortsetzt, wie z.~B. bei $\sqrt{(x-a)}$, $\log(x-a)$, $(x-a)^\mu$, wenn $\mu$ keine ganze Zahl ist, der Punkt $a$.

Wenn man von diesem Punkte $a$ aus sich eine beliebige Linie gezogen denkt, so kann der Werth der Function in der Umgebung von $a$ so gewählt werden, dass er sich ausserhalb dieser Linie überall stetig ändert; sie nimmt aber dann zu beiden Seiten dieser Linie verschiedene Werthe an, so dass die Fortsetzung der Function über diese Linie hinüber eine von der jenseits schon vorhandenen verschiedene Function giebt.

Zur Erleichterung des Ausdrucks sollen die verschiedenen Fortsetzungen Einer Function für denselben Theil der $x$-Ebene ``Zweige'' dieser Function genannt werden und ein Werth von $x$, um welchen herum sich ein Zweig einer Function in einen andern fortsetzt, ein ``Verzweigungswerth''; für einen Werth, in welchem keine Verzweigung stattfindet, heisst die Function ``einändrig oder monodrom''.

\subsection*{1.}
\addcontentsline{toc}{subsection}{Art. 1}

Ich bezeichne durch
\[
P\left\{\begin{array}{ccc|c}
a & b & c & \\
\alpha & \beta & \gamma & x \\
\alpha' & \beta' & \gamma' &
\end{array}\right\}
\]
eine Function von $x$, welche folgende Bedingungen erfüllt:
\begin{enumerate}
\item Sie ist für alle Werthe von $x$ ausser $a$, $b$, $c$ einändrig und endlich.
\item Zwischen je drei Zweigen dieser Function $P'$, $P''$, $P'''$ findet eine lineare homogene Gleichung mit constanten Coefficienten Statt,
\[
c'P' + c''P'' + c'''P''' = 0.
\]
\item Die Function lässt sich in die Formen
\[
C_{\alpha} P^{(\alpha)} + C_{\alpha'} P^{(\alpha')}, \quad
C_{\beta} P^{(\beta)} + C_{\beta'} P^{(\beta')}, \quad
C_{\gamma} P^{(\gamma)} + C_{\gamma'} P^{(\gamma')}
\]
mit Constanten $C_{\alpha}$, $C_{\alpha'}$, \dots\ setzen, so dass
\[
P^{(\alpha)}(x-a)^{-\alpha}, \quad P^{(\alpha')}(x-a)^{-\alpha'}
\]
für $x = a$ einändrig bleiben und weder Null noch unendlich werden, und ebenso $P^{(\beta)}(x-b)^{-\beta}$, $P^{(\beta')}(x-b)^{-\beta'}$ für $x = b$ und $P^{(\gamma)}(x-c)^{-\gamma}$, $P^{(\gamma')}(x-c)^{-\gamma'}$ für $x = c$. In Betreff der sechs Grössen $\alpha$, $\alpha'$, \dots, $\gamma'$ wird vorausgesetzt, dass keine der Differenzen $\alpha - \alpha'$, $\beta - \beta'$, $\gamma - \gamma'$ eine ganze Zahl und die Summe aller, $\alpha + \alpha' + \beta + \beta' + \gamma + \gamma' = 1$ sei.
\end{enumerate}

Wie mannigfaltig die Functionen seien, welche diesen Bedingungen genügen, bleibt vorläufig unentschieden und wird sich im Laufe der Untersuchung (Art.~4) ergeben. Zu grösserer Bequemlichkeit des Ausdrucks werde ich $x$ die \textit{Veränderliche}, $a$, $b$, $c$ den \textit{ersten, zweiten, dritten Verzweigungswerth} und $\alpha$, $\alpha'$; $\beta$, $\beta'$; $\gamma$, $\gamma'$ das \textit{erste, zweite, dritte Exponentenpaar} der $P$-function nennen.

\subsection*{2.}
\addcontentsline{toc}{subsection}{Art. 2}

Zunächst einige unmittelbare Folgerungen aus der Definition.

In der Function
\[
P\left\{\begin{array}{ccc|c}
a & b & c & \\
\alpha & \beta & \gamma & x \\
\alpha' & \beta' & \gamma' &
\end{array}\right\}
\]
können die drei ersten Verticalreihen beliebig unter einander vertauscht werden, sowie auch $\alpha$ mit $\alpha'$, $\beta$ mit $\beta'$, $\gamma$ mit $\gamma'$. Es ist ferner
\[
P\left\{\begin{array}{ccc|c}
a & b & c & \\
\alpha & \beta & \gamma & x \\
\alpha' & \beta' & \gamma' &
\end{array}\right\}
= P\left\{\begin{array}{ccc|c}
a_1 & b_1 & c_1 & \\
\alpha & \beta & \gamma & x_1 \\
\alpha' & \beta' & \gamma' &
\end{array}\right\},
\]
wenn man für $x_1$ einen rationalen Ausdruck ersten Grades von $x$ setzt, der für $x = a$, $b$, $c$ die Werthe $a_1$, $b_1$, $c_1$ annimmt.

Für
\[
P\left\{\begin{array}{ccc|c}
0 & \infty & 1 & \\
\alpha & \beta & \gamma & x \\
\alpha' & \beta' & \gamma' &
\end{array}\right\},
\]
auf welche Function sich demzufolge alle $P$-functionen mit denselben $\alpha$, $\alpha'$, \dots, $\gamma'$ zurückführen lassen, werde ich zur Abkürzung auch blos $P\left(\begin{array}{ccc|} \alpha & \beta & \gamma \\ \alpha' & \beta' & \gamma' \end{array}\, x\right)$ setzen.

Durch diese Umformung können zwei Exponenten verschiedener Paare beliebig gegebene Werthe erhalten und als Werthe der Exponenten, da zwischen ihnen die Bedingung $\alpha + \alpha' + \beta + \beta' + \gamma + \gamma' = 1$ stattfindet, jedwede andere eingeführt werden, für welche die drei Differenzen $\alpha - \alpha'$, $\beta - \beta'$, $\gamma - \gamma'$ dieselben sind. Aus diesem Grunde werde ich später zur Erleichterung der Uebersicht durch
\[
P(\alpha - \alpha', \beta - \beta', \gamma - \gamma', x)
\]
sämmtliche in der Form $x^{\delta} (1-x)^{\varepsilon} P\left(\begin{array}{ccc|} \alpha & \beta & \gamma \\ \alpha' & \beta' & \gamma' \end{array}\, x\right)$ enthaltenen Functionen bezeichnen.

\subsection*{3.}
\addcontentsline{toc}{subsection}{Art. 3}

Es ist jetzt vor allen Dingen nöthig, den Verlauf der Function etwas genauer zu untersuchen. Zu diesem Ende denke man sich durch sämmtliche Verzweigungspunkte der Function eine in sich zurücklaufende Linie $l$ gezogen, welche die Gesammtheit der complexen Werthe in zwei Grössengebiete scheidet. Innerhalb jedes von ihnen wird alsdann jeder Zweig der Function stetig und von den übrigen gesondert verlaufen; längs der gemeinschaftlichen Grenzlinie aber werden zwischen den Zweigen des einen und des andern Gebiets in verschiedenen Begrenzungstheilen verschiedene Relationen stattfinden. Zu ihrer bequemeren Darstellung werde ich die mittels des Coefficienten-Systems $S = \begin{pmatrix} p & q \ r & s \end{pmatrix}$ aus den Grössen $t$, $u$ gebildeten linearen Ausdrücke $pt + qu$, $rt + su$ durch $(S)(t, u)$ bezeichnen. Es möge ferner nach Analogie der von Gauss vorgeschlagenen Benennung ``positiv laterale Einheit'' für $+i$ als ``positive'' Seitenrichtung zu einer gegebenen Richtung diejenige bezeichnet werden, welche zu ihr ebenso liegt, wie $+i$ zu $1$ (also bei der üblichen Darstellungsweise der complexen Grössen die linke). Demgemäss macht $x$ einen ``positiven Umlauf um einen Verzweigungswerth $a$'', wenn es sich durch die ganze Begrenzung eines nur diesen und keinen andern Verzweigungswerth enthaltenden Grössengebiets in einer gegen die Richtung von Innen nach Aussen positiv liegenden Richtung bewegt. Es gehe nun die Linie $l$ der Reihe nach durch die Punkte $x = c$, $x = b$, $x = a$, und in dem auf ihrer positiven Seite liegenden Gebiete seien $P'$, $P''$ zwei in keinem constanten Verhältnisse stehende Zweige der Function $P$.

Jeder andere Zweig $P'''$ lässt sich dann, da in der Vorausgesetztermassen stattfindenden Gleichung $c'P' + c''P'' + c'''P''' = 0$ $c'''$ nicht verschwinden kann, linear und mit constanten Coefficienten in $P'$ und $P''$ ausdrücken. Nimmt man nun an, dass $P'$, $P''$ durch einen positiven Umlauf der Grösse $x$ um $a$ in $(A)(P', P'')$, um $b$ in $(B)(P', P'')$, um $c$ in $(C)(P', P'')$ übergehe, so wird durch die Coefficienten der Systeme $(A)$, $(B)$, $(C)$ die Periodicität der Function völlig bestimmt sein. Zwischen diesen finden aber noch Relationen Statt.

Wenn nämlich $x$ das negative Ufer der Linie $l$ durchläuft, so müssen die Functionen $P'$, $P''$ die vorigen Werthe wieder annehmen, da der durchlaufene Weg negativerseits die ganze Begrenzung eines Grössengebiets bildet, innerhalb dessen diese Functionen allenthalben einändrig sind. Es ist dies aber dasselbe, als ob der Werth $x$ sich von einem der Werthe $c$, $b$, $a$ bis zum folgenden auf der positiven Seite fortbewegt, dann aber jedesmal um diesen Werth positiv herum, wobei $(P', P'')$ der Reihe nach in $(C)(P', P'')$, $(C)(B)(P', P'')$, schliesslich in $(C)(B)(A)(P', P'')$ übergeht. Es ist daher
\begin{equation}
(C)(B)(A) = (E), \tag{1}
\end{equation}
welche Gleichung vier Bedingungsgleichungen zwischen den zwölf Coefficienten von $A$, $B$, $C$ liefert.

Bei der Discussion dieser Bedingungsgleichungen beschränke ich mich, zur Fixirung der Vorstellungen, auf die Function $P\left(\begin{array}{ccc|} \alpha & \beta & \gamma \\ \alpha' & \beta' & \gamma' \end{array}\, x\right)$, also auf den Fall, wenn $a = 0$, $b = \infty$, $c = 1$, was die Allgemeinheit der Resultate nicht wesentlich beeinträchtigt, und wähle für die durch $1$, $\infty$, $0$ zu ziehende Linie $l$ die Linie der reellen Werthe, welche, um der Reihe nach durch $c$, $b$, $a$ zu gehen, von $-\infty$ nach $+\infty$ gerichtet sein muss. Innerhalb des auf der positiven Seite dieser Linie liegenden Gebiets, welches die complexen Werthe mit positiv imaginärem Gliede enthält, sind dann die oben charakterisirten Bestandtheile der Function $P$, die Grössen $P^{\alpha}$, $P^{\alpha'}$, $P^{\beta}$, $P^{\beta'}$, $P^{\gamma}$, $P^{\gamma'}$, einändrige Functionen von $x$ und sind bis auf constante Factoren, welche von der Wahl der Grössen $c_{\alpha}$, $c_{\alpha'}$, \dots, $c_{\gamma'}$ abhängen, völlig bestimmt, wenn die Function $P$ gegeben ist.

Die Functionen $P^{\alpha}$, $P^{\alpha'}$ gehen durch einen positiven Umlauf der Grösse $x$ um $0$ in $P^{\alpha} e^{\alpha 2\pi i}$, $P^{\alpha'} e^{\alpha' 2\pi i}$ über und ebenso durch einen positiven Umlauf dieser Grösse um $\infty$ die Functionen $P^{\beta}$, $P^{\beta'}$ in $P^{\beta} e^{\beta 2\pi i}$, $P^{\beta'} e^{\beta' 2\pi i}$ und durch einen positiven Umlauf um $1$ die Functionen $P^{\gamma}$, $P^{\gamma'}$ in $P^{\gamma} e^{\gamma 2\pi i}$, $P^{\gamma'} e^{\gamma' 2\pi i}$. Bezeichnet man den Werth, in welchen $P$ durch einen positiven Umlauf von $x$ um $0$ übergeht, durch $P'$, so ist, wenn
\[
P = c_{\alpha} P^{\alpha} + c_{\alpha'} P^{\alpha')}, \quad P' = c_{\alpha} e^{\alpha 2\pi i} P^{\alpha} + c_{\alpha'} e^{\alpha' 2\pi i} P^{\alpha'}.
\]
Diese Ausdrücke haben eine von Null verschiedene Determinante, da n.\ V.\ $\alpha - \alpha'$ keine ganze Zahl ist, und folglich können $P^{\alpha}$, $P^{\alpha'}$ auch umgekehrt in $P$, $P'$ also auch in $P^{\beta}$, $P^{\beta'}$; $P^{\gamma}$, $P^{\gamma'}$ linear mit constanten Coefficienten ausgedrückt werden. Setzt man nun
\[
P^{\alpha} = a_{\beta} P^{\beta} + a_{\beta'} P^{\beta')} = a_{\gamma} P^{\gamma} + a_{\gamma'} P^{\gamma')},
\]
\[
P^{\alpha'} = a'_{\beta} P^{\beta} + a'_{\beta'} P^{\beta')} = a'_{\gamma} P^{\gamma} + a'_{\gamma'} P^{\gamma')},
\]
und zur Abkürzung
\[
\begin{pmatrix} a_{\beta} & a_{\beta'} \\ a'_{\beta} & a'_{\beta'} \end{pmatrix} = (b), \quad
\begin{pmatrix} a_{\gamma} & a_{\gamma'} \\ a'_{\gamma} & a'_{\gamma'} \end{pmatrix} = (c)
\]
und die inversen Substitutionen von $(b)$ und $(c)$ bezw.\ $= (b)^{-1}$ und $(c)^{-1}$, so ergeben sich für die Functionen $(P^{\alpha}, P^{\alpha'})$ die Substitutionen:

um $0$: $(A) = \begin{pmatrix} e^{\alpha 2\pi i} & 0 \\ 0 & e^{\alpha' 2\pi i} \end{pmatrix}$,
um $\infty$: $(b) \begin{pmatrix} e^{\beta 2\pi i} & 0 \\ 0 & e^{\beta' 2\pi i} \end{pmatrix} (b)^{-1}$,
um $1$: $(c) \begin{pmatrix} e^{\gamma 2\pi i} & 0 \\ 0 & e^{\gamma' 2\pi i} \end{pmatrix} (c)^{-1}$.

Aus der Gleichung $(C)(B)(A) = (E)$ folgt nun zunächst, da die Determinante einer zusammengesetzten Substitution dem Producte aus den Determinanten ihrer Componenten gleich ist,
\[
1 = \det(A) \det(B) \det(C)
\]
\[
= e^{(\alpha+\alpha'+\beta+\beta'+\gamma+\gamma')2\pi i} \det(b) \det(b)^{-1} \det(c) \det(c)^{-1}
\]
oder, da $\det(b) \det(b)^{-1} = 1$, $\det(c) \det(c)^{-1} = 1$,
\begin{equation}
\alpha + \alpha' + \beta + \beta' + \gamma + \gamma' = \text{einer ganzen Zahl}, \tag{2}
\end{equation}
womit die obige Annahme, dass diese Exponentensumme $= 1$ sei, vereinbar ist.

Die übrigen drei in $(C)(B)(A) = (E)$ enthaltenen Relationen geben drei Bedingungen für $(b)$ und $(c)$, welche indess leichter auf folgendem Wege gefunden werden.

Wenn $x$ erst um $0$ und dann um $\infty$ negativ herumgeht, so bildet der durchlaufene Weg zugleich einen positiven Umlauf um $1$. Der Werth, in welchen $P^{\alpha}$ dadurch übergeht, ist daher
\[
= a_{\gamma} e^{\gamma 2\pi i} P^{\gamma} + a_{\gamma'} e^{\gamma' 2\pi i} P^{\gamma'}
\]
\[
= (c)(b) e^{-\beta 2\pi i} P^{\beta} + a_{\beta'} e^{\beta' 2\pi i} P^{\beta'} e^{-\alpha 2\pi i}.
\]

\subsection*{4.}
\addcontentsline{toc}{subsection}{Art. 4}

Multiplicirt man diese Gleichung mit einem willk{"u}rlichen Factor $e^{-(\alpha+\beta+\gamma')\pi i}$ und die Gleichung
\[
a_{\gamma} P^{\gamma} + a_{\gamma'} P^{\gamma'} = a_{\beta} P^{\beta} + a_{\beta'} P^{\beta'} \text{ mit } e^{-(\alpha'+\beta+\gamma')\pi i}
\]
und subtrahirt, so ergiebt sich nach Abwerfung eines allgemeinen Factors
\begin{multline*}
a_{\gamma} \sin(\sigma - \gamma)\pi \cdot e^{\gamma'\pi i} P^{\gamma} + a_{\gamma'} \sin(\sigma - \gamma')\pi \cdot e^{\gamma\pi i} P^{\gamma'} = \\
a_{\beta} \sin(\sigma + \alpha + \beta)\pi \cdot e^{-(\alpha+\beta')\pi i} P^{\beta} + a_{\beta'} \sin(\sigma + \alpha + \beta')\pi \cdot e^{-(\alpha+\beta')\pi i} P^{\beta'}.
\end{multline*}

Aus ganz {"a}hnlichen Gr{"u}nden hat man auch, wenn man {"u}berall $\alpha'$ f{"u}r $\alpha$ setzt, die Gleichung
\begin{multline*}
a'_{\gamma} \sin(\sigma - \gamma)\pi \cdot e^{\gamma'\pi i} P^{\gamma} + a'_{\gamma'} \sin(\sigma - \gamma')\pi \cdot e^{\gamma\pi i} P^{\gamma'} = \\
a'_{\beta} \sin(\sigma + \alpha' + \beta)\pi \cdot e^{-(\alpha'+\beta')\pi i} P^{\beta} + a'_{\beta'} \sin(\sigma + \alpha' + \beta')\pi \cdot e^{-(\alpha'+\beta')\pi i} P^{\beta'}
\end{multline*}
mit der willk{"u}rlichen Gr{"o}sse $\sigma$. Befreit man beide Gleichungen von einer der Functionen, z.~B.\ $P^{\gamma'}$, indem man $\sigma$ demgem{"a}ss bestimmt, so k{"o}nnen sich die resultirenden Gleichungen nur durch einen allgemeinen constanten Factor unterscheiden, da $\frac{P^{\beta}}{P^{\beta'}}$ nicht constant ist.

Diese Elimination von $P^{\gamma'}$ giebt daher:
\begin{equation}
\frac{a_{\gamma}}{a'_{\gamma}} = \frac{a_{\beta} \sin(\alpha+\beta+\gamma')\pi \cdot e^{-\alpha\pi i} - a_{\beta'} \sin(\alpha+\beta'+\gamma')\pi \cdot e^{-\alpha'\pi i}}{a'_{\beta} \sin(\alpha'+\beta+\gamma')\pi \cdot e^{-\alpha'\pi i} - a'_{\beta'} \sin(\alpha'+\beta'+\gamma')\pi \cdot e^{-\alpha'\pi i}}, \tag{3}
\end{equation}
und die {"a}hnliche Elimination von $P^{\gamma}$
\begin{equation}
\frac{a_{\gamma'}}{a'_{\gamma'}} = \frac{a_{\beta} \sin(\alpha+\beta+\gamma)\pi \cdot e^{-\alpha\pi i} - a_{\beta'} \sin(\alpha+\beta'+\gamma)\pi \cdot e^{-\alpha'\pi i}}{a'_{\beta} \sin(\alpha'+\beta+\gamma)\pi \cdot e^{-\alpha'\pi i} - a'_{\beta'} \sin(\alpha'+\beta'+\gamma)\pi \cdot e^{-\alpha'\pi i}}, \tag{4}
\end{equation}
welches die vier gesuchten Relationen sind. Aus ihnen ergeben sich die Verh{"a}ltnisse der Quotienten $\frac{a_{\beta}}{a'_{\beta}}$, $\frac{a_{\beta'}}{a'_{\beta'}}$, $\frac{a_{\gamma}}{a'_{\gamma}}$, $\frac{a_{\gamma'}}{a'_{\gamma'}}$. Die Gleichheit der beiden aus der zweiten und vierten fliessenden Werthe von $\frac{a_{\beta}}{a'_{\beta}} : \frac{a_{\gamma}}{a'_{\gamma}}$ erhellt leicht als eine Folge aus $\alpha + \alpha' + \beta + \beta' + \gamma + \gamma' = 1$ mittelst der Identit{"a}t $\sin s\pi = \sin(1-s)\pi$.

Demnach sind von den Gr{"o}ssen $\frac{a_{\beta}}{a'_{\beta}}$, $\frac{a_{\beta'}}{a'_{\beta'}}$, $\frac{a_{\gamma}}{a'_{\gamma}}$, $\frac{a_{\gamma'}}{a'_{\gamma'}}$ durch eine von ihnen, z.~B.\ $\frac{a_{\beta}}{a'_{\beta}}$, die {"u}brigen bestimmt und die drei Gr{"o}ssen $a'_{\beta}$, $a'_{\gamma}$, $a_{\gamma'}$ durch die f{"u}nf Gr{"o}ssen $a_{\beta}$, $a_{\beta'}$, $a'_{\beta'}$, $a_{\gamma}$, $a_{\gamma'}$. Diese f{"u}nf Gr{"o}ssen aber h{"a}ngen von den in $P^{\alpha}$, $P^{\alpha'}$, $P^{\beta}$, $P^{\beta'}$, $P^{\gamma}$, $P^{\gamma'}$, wenn die Function $P$ gegeben ist, noch willk{"u}rlichen Factoren oder vielmehr von deren Verh{"a}ltnissen ab, und k{"o}nnen durch geeignete Bestimmung derselben jedwede endliche Werthe erhalten.\footnote{Diese Bemerkung zeigt, dass man jene Factoren auch so bestimmen kann, dass $a'_{\beta} = a_{\gamma} = a_{\gamma'} = a'_{\gamma} = a'_{\gamma'} = 1$, $a_{\beta'} = a_{\beta}$ wird.}

Die soeben gemachte Bemerkung bahnt den Weg zu dem Satze, dass in zwei $P$-functionen mit gleichen Exponenten die denselben Exponenten entsprechenden Bestandtheile sich nur durch einen constanten Factor unterscheiden.

In der That, ist $P_1$ eine Function mit denselben Exponenten wie $P$, so kann man die f{"u}nf Gr{"o}ssen $a_{\beta}$, $a_{\beta'}$, $a_{\gamma}$, $a_{\gamma'}$ und $a'_{\beta}$ bei beiden gleich annehmen und dann m{"u}ssen auch die Gr{"o}ssen $a'_{\beta'}$, $a'_{\gamma}$, $a'_{\gamma'}$ bei beiden {"u}bereinstimmen. Man hat also gleichzeitig:
\[
(P^{\alpha}, P^{\alpha'}) = (b)(P^{\beta}, P^{\beta'}) = (c)(P^{\gamma}, P^{\gamma'})
\]
und
\[
(P_1^{\alpha}, P_1^{\alpha'}) = (b)(P_1^{\beta}, P_1^{\beta'}) = (c)(P_1^{\gamma}, P_1^{\gamma'}),
\]
folglich
\[
(P^{\alpha} P_1^{\alpha'} - P^{\alpha'} P_1^{\alpha}) = \det(b)(P^{\beta} P_1^{\beta'} - P^{\beta'} P_1^{\beta}) = \det(c)(P^{\gamma} P_1^{\gamma'} - P^{\gamma'} P_1^{\gamma}).
\]

Von diesen drei Ausdr{"u}cken bleibt der erste, mit $x^{\alpha+\alpha'-1}$ multiplicirt, offenbar f{"u}r $x = 0$ ein{"a}ndrig und endlich; ebenso der zweite, mit $x^{\beta+\beta'+1} = x^{\alpha-\alpha'-\gamma-\gamma'+1}$ multiplicirt, f{"u}r $x = \infty$, der dritte, mit $(1-x)^{-\gamma-\gamma'}$ multiplicirt, f{"u}r $x = 1$, und dasselbe gilt von allen drei Ausdr{"u}cken f{"u}r alle von $0$, $\infty$, $1$ verschiedenen Werthe von $x$; es ist daher
\[
(P^{\alpha} P_1^{\alpha'} - P^{\alpha'} P_1^{\alpha}) \cdot x^{\alpha+\alpha'-1} (1-x)^{-\gamma-\gamma'}
\]
eine allenthalben stetige und ein{"a}ndrige Function, also eine Constante. Sie ist ferner $= 0$ f{"u}r $x = \infty$ und muss folglich allenthalben $= 0$ sein. Hieraus folgt
\[
\frac{P^{\alpha'}}{P^{\alpha}} = \frac{P_1^{\alpha'}}{P_1^{\alpha}}.
\]
Ebenso
\[
\frac{P^{\beta}}{P^{\beta'}} = \frac{P^{\alpha}}{P^{\alpha'}} = \frac{a_{\beta} P^{\beta} + a_{\beta'} P^{\beta'}}{P^{\beta'}} = \frac{P^{\alpha}}{P^{\alpha'}} = \frac{P^{\gamma}}{P^{\gamma'}}.
\]
Die Function $\frac{P^{\alpha'}}{P^{\alpha}}$ ist demnach einwerthig und muss {"u}berdies allenthalben endlich, also, w.~z.~b.~w., constant sein, wenn noch bewiesen wird, dass $P^{\alpha}$ und $P^{\alpha'}$ nicht zugleich f{"u}r einen von $0$, $1$, $\infty$ verschiedenen Werth von $x$ verschwinden k{"o}nnen.

Zu diesem Ende bemerke man, dass
\begin{multline*}
\frac{dP^{\alpha}}{dx} P^{\alpha'} - P^{\alpha} \frac{dP^{\alpha'}}{dx} = \det(b) \left(\frac{dP^{\beta}}{dx} P^{\beta'} - P^{\beta} \frac{dP^{\beta'}}{dx}\right) \\
= \det(c) \left(\frac{dP^{\gamma}}{dx} P^{\gamma'} - P^{\gamma} \frac{dP^{\gamma'}}{dx}\right)
\end{multline*}
und folglich f{"u}r $x = 0$, $\infty$, $1$ unendlich klein von den Ordnungen $\alpha + \alpha' - 1$, $\beta + \beta' + 1 = 2 - \alpha - \alpha' - \gamma - \gamma'$, $\gamma + \gamma' - 1$ wird, {"u}brigens aber stetig und ein{"a}ndrig bleibt, so dass
\[
\left(\frac{dP^{\alpha}}{dx} P^{\alpha'} - P^{\alpha} \frac{dP^{\alpha'}}{dx}\right) x^{\alpha+\alpha'-1} (1-x)^{-\gamma-\gamma'+1}
\]
eine allenthalben stetige und ein{"a}ndrige Function bildet, folglich einen Constanten Werth hat. Dieser constante Werth dieser Function ist nothwendig von Null verschieden, weil sonst $\log P^{\alpha} - \log P^{\alpha'} = \text{const.}$, folglich $\alpha = \alpha'$ sein w{"u}rde gegen die Voraussetzung; offenbar m{"u}sste sie gleich Null werden, wenn f{"u}r einen von $0$, $1$, $\infty$ verschiedenen Werth von $x$ $P^{\alpha}$ und $P^{\alpha'}$ gleichzeitig verschw{"a}nden, da $\frac{dP^{\alpha}}{dx}$, $\frac{dP^{\alpha'}}{dx}$ als Derivirte ein{"a}ndrig und stetig bleibender Functionen nicht unendlich werden k{"o}nnen.

Es werden daher $P^{\alpha}$ und $P^{\alpha'}$ f{"u}r keinen von $0$, $1$, $\infty$ verschiedenen Werth von $x$ gleichzeitig $= 0$, und es bleibt die einwerthige Function
\[
\frac{P^{\alpha}}{P^{\alpha'}} = \frac{P^{\beta}}{P^{\beta'}} = \frac{P^{\gamma}}{P^{\gamma'}}
\]
allenthalben endlich, mithin constant, w.~z.~b.~w.

Aus dem eben bewiesenen Satze folgt, dass in zwei Zweige Einer $P$-function, deren Quotient nicht constant ist, jede andere $P$-function mit gleichen Exponenten sich linear mit constanten Coefficienten ausdr{"u}cken l{"a}sst und dass durch die im Art.~1 geforderten Eigenschaften die zu definirende Function bis auf zwei linear in ihr enthaltene Constanten v{"o}llig bestimmt ist. Diese werden in jedem Falle leicht aus den Werthen der Function f{"u}r specielle Werthe der Ver{"a}nderlichen gefunden, am bequemsten, indem man die Ver{"a}nderliche einem der Verzweigungswerthe gleich setzt.

Ob es immer eine jenen Bedingungen gen{"u}gende Function gebe, bleibt freilich noch unentschieden, wird sich aber sp{"a}ter durch die wirkliche Darstellung der Function mittelst bestimmter Integrale und hypergeometrischer Reihen erledigen und bedarf daher keiner besondern Untersuchung.

\subsection*{5.}
\addcontentsline{toc}{subsection}{Art. 5}

Ausser den f{"u}r jedwede Werthe der Exponenten m{"o}glichen Transformationen des Art.~2 ergeben sich aus der Definition noch leicht die beiden Transformationen:

\medskip
\noindent (A)
\[
P\left\{\begin{array}{ccc|c}
0 & \infty & 1 \\
\alpha & \beta & \gamma & x^2 \\
\alpha' & \beta' & \gamma'
\end{array}\right\}
= P\left\{\begin{array}{ccc|c}
-1 & \infty & 1 \\
\gamma & 2\beta & \gamma & x \\
\gamma' & \beta' & \gamma'
\end{array}\right\},
\]
wo nach dem Fr{"u}heren $\beta + \beta' + \gamma + \gamma' = \frac{1}{2}$ sein muss, und

\medskip
\noindent (B)
\[
P\left\{\begin{array}{ccc|c}
0 & \infty & 1 \\
\alpha & \beta & \gamma & x^3 \\
\alpha' & \beta' & \gamma'
\end{array}\right\}
= P\left\{\begin{array}{ccc|c}
1 & \varrho & \varrho^2 \\
\gamma & \gamma & \gamma & x \\
\gamma' & \gamma' & \gamma'
\end{array}\right\},
\]
wo $\gamma + \gamma' = \frac{1}{3}$ und $\varrho$ eine imagin{"a}re dritte Wurzel der Einheit bezeichnet. Um s{"a}mmtliche Functionen, welche sich mit H{"u}lfe dieser Transformationen auf einander zur{"u}ckf{"u}hren lassen, bequem zu {"u}bersehen, ist es zweckm{"a}ssig, statt der Exponenten ihre Differenzen einzuf{"u}hren und, wie oben vorgeschlagen, durch $P(\alpha - \alpha', \beta - \beta', \gamma - \gamma', x)$ s{"a}mmtliche in der Form $x^{\delta} (1-x)^{\varepsilon} P\left(\begin{array}{ccc|} \alpha & \beta & \gamma \\ \alpha' & \beta' & \gamma' \end{array}\, x\right)$ enthaltenen Functionen zu bezeichnen, wobei $\alpha - \alpha'$, $\beta - \beta'$, $\gamma - \gamma'$ die erste, zweite, dritte Exponentendifferenz genannt werden mag.

Aus den Formeln im Art.~2 folgt dann, dass in der Function $P(\lambda, \mu, \nu, x)$ die Gr{"o}ssen $\lambda$, $\mu$, $\nu$ beliebig in's Entgegengesetzte verwandelt und beliebig unter einander vertauscht werden k{"o}nnen. Die Ver{"a}nderliche nimmt dabei einen der 6 Werthe $x$, $1 - x$, $\frac{1}{x}$, $\frac{1}{1-x}$, $\frac{x}{x-1}$, $\frac{x-1}{x}$ an, und zwar haben von den 48 auf diese Weise sich ergebenden $P$-functionen je acht, welche durch blosse Zeichen{"a}nderung der Gr{"o}ssen $\lambda$, $\mu$, $\nu$ aus einander hervorgehen, dieselbe Ver{"a}nderliche.

Von den in diesem Art.\ angegebenen Transformationen A und B ist die erste anwendbar, wenn von den Exponentendifferenzen entweder eine gleich $\frac{1}{2}$ oder zwei einander gleich sind, die zweite, wenn von ihnen entweder zwei $= \frac{1}{3}$ oder alle drei einander gleich sind. Durch successive Anwendung dieser Transformationen erh{"a}lt man daher durch einander ausgedr{"u}ckt:

I.\ $P(\mu, \nu, \frac{1}{2}, x_1)$, $P(\mu, 2\nu, \nu, x_2)$ und $P(\nu, 2\mu, \nu, x_2)$, wobei $\sqrt{1-x_2} = 1 - 2x_1$, $\sqrt{\frac{1}{x_2}} = 1 - 2x_1$, also $x_2 = 4x_1(1-x_1) = 4x_2(1-x_2)$ sich ergiebt.

II.\ $P(\nu, \nu, \nu, x_3)$, $P(\nu, \frac{1}{2}, \frac{1}{2}, x_4)$, $P(\frac{1}{2}, \nu, \frac{1}{2}, x_4)$, $P(\frac{1}{2}, \frac{1}{2}, \nu, x_4)$, $P(\frac{1}{2}, \nu, \frac{1}{3}, x_5)$, $P(\frac{1}{3}, \frac{1}{2}, \nu, x_5)$, $P(\nu, \frac{1}{3}, \frac{1}{2}, x_5)$, wenn
\[
\frac{1-x_4}{x_4} = \frac{x_3^3}{27} \text{ und folglich } x_4 = \frac{27}{27+x_3^3};
\]
\[
x_4(1-x_4) = \frac{x_3^3}{(27+x_3^3)^2} = \frac{x_5}{27}, \text{ ferner nach I.}
\]
\[
4x_5(1-x_5) = x_6 = \frac{4x_3^3}{(27+x_3^3)^2}, \quad 4x_5(1-x_5) = x_6 = 4x_7(1-x_7).
\]

III.\ $P(\nu, \nu, \frac{1}{2}, x_8)$, $P(\nu, 2\nu, \frac{1}{2}, x_9)$, $P(\frac{1}{2}, \nu, \frac{1}{2}, x_8)$, $P(\frac{1}{2}, 2\nu, \frac{1}{2}, x_9)$, wenn $x_8 = 4x_9(1-x_9) = \frac{4}{x_3^3} = 4x_4(1-x_4)$, $x_9 = 4x_8(1-x_8)$.

Alle diese Functionen k{"o}nnen noch mittelst der allgemeinen Transformationen umgeformt und dadurch ihre Exponentendifferenzen beliebig vertauscht und mit beliebigen Vorzeichen versehen werden. Ausser den beiden Transcendenten II.\ und III.\ l{"a}sst, wenn eine Exponentendifferenz willk{"u}rlich bleiben soll, nur noch die Function $P(\nu, \frac{1}{2}, \frac{1}{2}) = P(\nu, \frac{1}{2}, \nu)$ eine h{"a}ufigere Wiederholung der Transformationen A und B zu, welche indess, da
\[
P\left(\begin{array}{ccc|} \nu & \frac{1}{2} & \frac{1}{2} \\ -\nu & \frac{1}{2} & \frac{1}{2} \end{array}\, x\right) = \text{const.} \cdot x^{\nu} + \text{const.}'
\]
auf ganz elementare Formeln f{"u}hrt.

In der That ist die Transformation B nur anwendbar auf $P(\nu, \nu, \nu)$ oder $P(\frac{1}{3}, \nu, \frac{1}{3})$, also nur auf die Transcendente II.; die Transformation A aber l{"a}sst sich h{"a}ufiger als in I.\ nur wiederholen, wenn entweder von den Gr{"o}ssen $\mu$, $\nu$, $2\mu$, $2\nu$ eine gleich $\frac{1}{2}$ gesetzt oder eine der Gleichungen $\mu = \nu$, $\mu = 2\nu$, $\nu = 2\mu$ angenommen wird. Von diesen Annahmen f{"u}hrt $\mu = 2\nu$ oder $\nu = 2\mu$ auf die Transcendente II., $\mu = \nu$, sowie $2\mu$ oder $2\nu = \frac{1}{2}$ auf die Transcendente III., endlich $\mu$ oder $\nu = \frac{1}{2}$ auf die Function $P(\nu, \frac{1}{2}, \frac{1}{2})$.

Die Anzahl der verschiedenen Ausdr{"u}cke, welche man durch diese Transformationen f{"u}r jede der Transcendenten I---III.\ erh{"a}lt, ergiebt sich, wenn man ber{"u}cksichtigt, dass in den obigen $P$-functionen als Ver{"a}nderliche alle Wurzeln der Gleichungen, durch welche sie bestimmt werden, zul{"a}ssig sind und jede Wurzel zu einem Systeme von 6 Werthen geh{"o}rt, welche mittelst der allgemeinen Transformation f{"u}r einander als Ver{"a}nderliche eingef{"u}hrt werden k{"o}nnen.

Es f{"u}hren aber im Falle I.\ die beiden Werthe von $x_1$ und $x_2$, welche zu einem gegebenen $x_3$ geh{"o}ren, auf dasselbe System von 6 Werthen, so dass jede der Functionen I.\ durch $P$-functionen mit $6 \cdot 3 = 18$ verschiedenen Ver{"a}nderlichen ausgedr{"u}ckt werden kann.

Im Falle II.\ f{"u}hren von den zu einem gegebenen Werthe von $x_3$ geh{"o}rigen Werthen die beiden Werthe von $x_4$ und $x_5$, die 6 Werthe von $x_6$ und von den 6 Werthen von $x_7$ je zwei zu demselben Systeme von 6 Werthen, w{"a}hrend die drei Werthe von $x_8$ zu drei verschiedenen Systemen von je 6 Werthen f{"u}hren. Es liefern also $x_4$ und $x_5$ je drei und $x_6$, $x_7$, $x_8$, $x_9$ je ein System von 6 Werthen, also alle zusammen $6 \cdot 10 = 60$ Werthe, durch deren $P$-functionen sich jede der Functionen II.\ ausdr{"u}cken l{"a}sst.

Im Falle III.\ endlich liefern $x_8$ die beiden Werthe von $x_9$, die beiden Werthe von $x_{10}$ und von den vier Werthen von $x_7$ je zwei ein System von 6 Werthen, so dass jede der Functionen III.\ durch $P$-functionen von $6 \cdot 5 = 30$ verschiedenen Ver{"a}nderlichen darstellbar ist.

In jeder $P$-function k{"o}nnen nun ohne Aenderung der Ver{"a}nderlichen mittelst der allgemeinen Transformationen die Exponentendifferenzen beliebige Vorzeichen erhalten, und also kann, da keine dieser Exponentendifferenzen $= 0$ ist, eine und dieselbe Function auf 8 verschiedene Arten als $P$-function derselben Ver{"a}nderlichen dargestellt werden. Die Anzahl s{"a}mmtlicher Ausdr{"u}cke betr{"a}gt also im Falle I.\ $8 \cdot 6 \cdot 3 = 144$, im Falle II.\ $8 \cdot 6 \cdot 10 = 480$, im Falle III.\ $8 \cdot 6 \cdot 5 = 240$.

\subsection*{6.}
\addcontentsline{toc}{subsection}{Art. 6}

Wenn man s{"a}mmtliche Exponenten einer $P$-function um ganze Zahlen {"a}ndert, so bleiben in den Gleichungen (3), (4) Art.~3 die Gr{"o}ssen
\[
\frac{\sin(\alpha + \beta + \gamma')\pi \cdot e^{-\alpha\pi i}}{\sin(\alpha' + \beta + \gamma')\pi \cdot e^{-\alpha'\pi i}}, \quad
\frac{\sin(\alpha + \beta' + \gamma')\pi \cdot e^{-\alpha\pi i}}{\sin(\alpha' + \beta' + \gamma')\pi \cdot e^{-\alpha'\pi i}},
\]
\[
\frac{\sin(\alpha + \beta + \gamma)\pi \cdot e^{-\alpha\pi i}}{\sin(\alpha' + \beta + \gamma)\pi \cdot e^{-\alpha'\pi i}}, \quad
\frac{\sin(\alpha + \beta' + \gamma)\pi \cdot e^{-\alpha\pi i}}{\sin(\alpha' + \beta' + \gamma)\pi \cdot e^{-\alpha'\pi i}}
\]
unge{"a}ndert.

Sind daher in den Functionen $P\left(\begin{array}{ccc|} \alpha & \beta & \gamma \\ \alpha' & \beta' & \gamma' \end{array}\, x\right)$, $P_1\left(\begin{array}{ccc|} \alpha_1 & \beta_1 & \gamma_1 \\ \alpha_1' & \beta_1' & \gamma_1' \end{array}\, x\right)$ die entsprechenden Exponenten $\alpha_1$ und $\alpha$, etc., um ganze Zahlen verschieden, so kann man die acht Gr{"o}ssen $(a_{\beta})_1$, $(a'_{\beta})_1$, $(a_{\beta'})_1$, \dots\ und die acht Gr{"o}ssen $a_{\beta}$, $a'_{\beta}$, $a_{\beta'}$, \dots\ gleich annehmen, da aus der Gleichheit der f{"u}nf willk{"u}rlichen die Gleichheit der drei {"u}brigen folgt.

Nach der im Art.~4 angewandten Schlussweise folgt hieraus:
\begin{multline*}
P^{\alpha} P_1^{\alpha'} - P^{\alpha'} P_1^{\alpha} = \det(b)(P^{\beta} P_1^{\beta'} - P^{\beta'} P_1^{\beta}) \\
= \det(c)(P^{\gamma} P_1^{\gamma'} - P^{\gamma'} P_1^{\gamma});
\end{multline*}
und wenn man von den Gr{"o}ssen $\alpha + \alpha_1'$ und $\alpha_1 + \alpha'$, $\beta + \beta_1'$ und $\beta_1 + \beta'$, $\gamma + \gamma_1'$ und $\gamma_1 + \gamma'$ diejenigen Gr{"o}ssen jedes Paars, welche um eine positive ganze Zahl kleiner sind, als die andern, durch $\bar{\alpha}$, $\bar{\beta}$, $\bar{\gamma}$ bezeichnet, so ist
\[
(P^{\alpha} P_1^{\alpha'} - P^{\alpha'} P_1^{\alpha}) \cdot x^{\alpha+\alpha'-1} (1-x)^{-\gamma-\gamma'}
\]
eine Function von $x$, welche ein{"a}ndrig und endlich bleibt f{"u}r $x = 0$, $x = 1$ und alle {"u}brigen endlichen Werthe von $x$, f{"u}r $x = \infty$ aber unendlich wird von der Ordnung $-\bar{\alpha} - \bar{\gamma} - \bar{\beta}$, folglich eine ganze Function $F$ vom Grade $-\bar{\alpha} - \bar{\beta} - \bar{\gamma}$.

Man bezeichne nun, wie fr{"u}her, die Exponentendifferenzen $\alpha - \alpha'$, $\beta - \beta'$, $\gamma - \gamma'$ durch $\lambda$, $\mu$, $\nu$. In Betreff dieser ergiebt sich zun{"a}chst: ihre Summe {"a}ndert sich um eine gerade Zahl, wenn sich s{"a}mmtliche Exponenten um ganze Zahlen {"a}ndern; denn sie {"u}bertrifft die Summe s{"a}mmtlicher Exponenten, welche unver{"a}ndert $= 1$ bleibt, um
\[
(\alpha - \alpha') + (\beta - \beta') + (\gamma - \gamma'),
\]
welche Gr{"o}sse sich dabei um eine gerade Zahl {"a}ndert. Sie k{"o}nnen sich aber dabei um jedwede ganze Zahlen {"a}ndern, deren Summe gerade ist. Bezeichnet man ferner $\alpha_1 - \alpha_1'$, $\beta_1 - \beta_1'$, $\gamma_1 - \gamma_1'$ durch $\lambda_1$, $\mu_1$, $\nu_1$ und durch $\varkappa_{\lambda}$, $\varkappa_{\mu}$, $\varkappa_{\nu}$ die absoluten Werthe der Differenzen $\lambda - \lambda_1$, $\mu - \mu_1$, $\nu - \nu_1$, so ist von den Gr{"o}ssen $\alpha + \alpha_1'$ und $\alpha_1 + \alpha'$ diejenige, welche um die positive Zahl $\varkappa_{\lambda}$ kleiner ist als die andere
\[
\bar{\alpha} = -\frac{\alpha + \alpha' + \alpha_1 + \alpha_1'}{2} - \frac{\varkappa_{\lambda}}{2}
\]
und ebenso
\[
\bar{\beta} = -\frac{\beta + \beta' + \beta_1 + \beta_1'}{2} - \frac{\varkappa_{\mu}}{2}, \quad
\bar{\gamma} = -\frac{\gamma + \gamma' + \gamma_1 + \gamma_1'}{2} - \frac{\varkappa_{\nu}}{2}.
\]
Der Grad der ganzen Function $F$, welcher gleich der Summe dieser Gr{"o}ssen ist, ergiebt sich daher
\[
-\bar{\alpha} - \bar{\beta} - \bar{\gamma} = 1 + \frac{\varkappa_{\lambda} + \varkappa_{\mu} + \varkappa_{\nu}}{2}.
\]

\subsection*{7.}
\addcontentsline{toc}{subsection}{Art. 7}

Sind jetzt
\[
P\left(\begin{array}{ccc|} \alpha & \beta & \gamma \\ \alpha' & \beta' & \gamma' \end{array}\, x\right), \quad
P_1\left(\begin{array}{ccc|} \alpha_1 & \beta_1 & \gamma_1 \\ \alpha_1' & \beta_1' & \gamma_1' \end{array}\, x\right), \quad
P_2\left(\begin{array}{ccc|} \alpha_2 & \beta_2 & \gamma_2 \\ \alpha_2' & \beta_2' & \gamma_2' \end{array}\, x\right)
\]
drei Functionen, in welchen sich die entsprechenden Exponenten um ganze Zahlen unterscheiden, so fliesst aus diesem Satze mittelst der identischen Gleichung
\begin{multline*}
P^{\alpha} (P_1^{\alpha'} P_2^{\alpha'} - P_1^{\alpha'} P_2^{\alpha}) + P_1^{\alpha} (P_2^{\alpha} P^{\alpha'} - P_2^{\alpha'} P^{\alpha}) \\
+ P_2^{\alpha} (P^{\alpha} P_1^{\alpha'} - P^{\alpha'} P_1^{\alpha}) = 0
\end{multline*}
der wichtige Satz, dass zwischen ihren entsprechenden Gliedern eine lineare homogene Gleichung stattfindet, deren Coefficienten ganze Functionen von $x$ sind, und dass also

\medskip
\textit{``s{"a}mmtliche $P$-functionen, deren entsprechende Exponenten sich um ganze Zahlen unterscheiden, sich in zwei beliebige von ihnen linear mit rationalen Functionen von $x$ als Coefficienten ausdr{"u}cken lassen''.}
\medskip

Eine specielle Folge aus den Beweisgr{"u}nden dieses Satzes ist, dass sich der zweite Differentialquotient einer $P$-function linear mit rationalen Functionen als Coefficienten in den ersten und die Function selbst ausdr{"u}cken l{"a}sst, und also die Function einer linearen homogenen Differentialgleichung zweiter Ordnung gen{"u}gt.

Beschr{"a}nkt man sich, um ihre Ableitung m{"o}glichst zu vereinfachen, auf den Fall $\gamma = 0$, auf welchen der allgemeine nach Art.~2 leicht zur{"u}ckgef{"u}hrt wird, und setzt $P = y$, $P^{\alpha} = y'$, $P^{\alpha'} = y''$, so ergiebt sich, dass die Functionen
\[
y \frac{\partial^2 y''}{\partial \log x^2} - y'' \frac{\partial^2 y}{\partial \log x^2}, \quad
y' \frac{\partial^2 y}{\partial \log x^2} - y \frac{\partial^2 y'}{\partial \log x^2},
\]
\[
\frac{\partial y'}{\partial \log x} \frac{\partial^2 y''}{\partial \log x^2} - \frac{\partial y''}{\partial \log x} \frac{\partial^2 y'}{\partial \log x^2}
\]
mit $x^{-\alpha-\alpha'+1} (1-x)^{-\gamma+\gamma'+1}$ multiplicirt, endlich und ein{"a}ndrig bleiben f{"u}r endliche Werthe von $x$ und unendlich von der ersten Ordnung werden f{"u}r $x = \infty$, und dass {"u}berdies das erste dieser Producte f{"u}r $x = 1$ unendlich klein von der ersten Ordnung wird. F{"u}r
\[
y = \text{const.}' \cdot y' + \text{const.}'' \cdot y''
\]
findet daher eine Gleichung von der Form statt
\[
(1-x) \frac{\partial^2 y}{\partial \log x^2} + (A + A'x) \frac{\partial y}{\partial \log x} + (B' - Bx) \cdot y = 0,
\]
in welcher $A$, $A'$, $B$, $B'$ noch zu bestimmende Constanten bezeichnen.

Nach der Methode der unbestimmten Coefficienten l{"a}sst sich eine L{"o}sung dieser Differentialgleichung nach um $1$ steigenden oder fallenden Potenzen in eine Reihe
\[
y = \sum_n a_n x^n
\]
entwickeln, und zwar wird der Exponent $\mu$ des Anfangsgliedes im ersten Falle, wo er der niedrigste ist, durch die Gleichung
\[
\mu\mu - A'\mu + B' = 0,
\]
und im zweiten, wo er der h{"o}chste ist, durch die Gleichung
\[
\mu\mu + B\mu + B' = 0
\]
bestimmt. Die Wurzeln der ersteren Gleichung m{"u}ssen $\alpha$ und $\alpha'$, die der letztern $-\beta$ und $-\beta'$ sein und folglich ist
\[
A = \alpha + \alpha', \quad A' = \alpha\alpha', \\quad B = \beta + \beta', \quad B' = \beta\beta'.
\]

\subsection*{8.}
\addcontentsline{toc}{subsection}{Art. 8}

Es bestimmen sich ferner die Coefficienten aus einem von ihnen mittelst der Recursionsformel
\[
\frac{a_{n+1}}{a_n} = \frac{(n + \beta)(n + \beta')}{(n + 1 - \alpha)(n + 1 - \alpha')},
\]
welcher $a = \frac{\Pi(n - \alpha) \Pi(n - \alpha')}{\Pi(n) \Pi(n - \alpha) \Pi(-n - \beta) \Pi(-n - \beta')} \cdot \text{const.}$ gen{"u}gt.

Demnach bildet die Reihe
\[
y = \text{Const.} \sum_n \frac{\Pi(n - \alpha) \Pi(n - \alpha')}{\Pi(n) \Pi(n) \Pi(-n - \beta) \Pi(-n - \beta')} x^n
\]
sowohl wenn die Exponenten von $\alpha$ oder $\alpha'$ an um die Einheit steigen, als auch wenn sie von $-\beta$ oder $-\beta'$ an um die Einheit fallen, eine L{"o}sung der Differentialgleichung und zwar bezw.\ diejenigen particul{"a}ren L{"o}sungen, welche oben durch $P^{\alpha}$, $P^{\alpha'}$, $P^{\beta}$, $P^{\beta'}$ bezeichnet worden sind.

Nach Gauss, welcher durch $F(a, b, c, x)$ eine Reihe bezeichnet, in welcher der Quotient des $(n + 1)$ten Gliedes in das folgende $= \frac{(n + a)(n + b)}{(n + 1)(n + c)}$ und das erste Glied $= 1$ ist, l{"a}sst sich dieses Resultat f{"u}r den einfachsten Fall, f{"u}r $\gamma = 0$, so ausdr{"u}cken:
\[
F(a, b, c, x) = P^{\alpha} \left(\begin{array}{ccc|} a & 0 & c - a - b \\ a' & \beta' & \gamma' \end{array}\, x\right)
\]
oder
\[
F(a, b, c, x) = P^{\alpha} \left(\begin{array}{ccc|} a & 0 & c - a - b \\ 1 - c & b - c & 0 \end{array}\, x\right).
\]
Aus demselben erh{"a}lt man auch leicht einen Ausdruck der $P$-function durch ein bestimmtes Integral, indem man in dem allgemeinen Gliede der Reihe f{"u}r die $\Pi$-functionen ein Euler'sches Integral zweiter Gattung einf{"u}hrt und dann die Ordnung der Summation und Integration vertauscht. Auf diese Weise findet man, dass das Integral
\[
x^{\alpha} (1 - x)^{\gamma} \int s^{\alpha' + \beta + \gamma - 1} (1 - s)^{-\alpha' - \beta' - \gamma} (1 - xs)^{-\alpha - \beta' - \gamma} \, ds
\]
von einem der vier Werthe $0$, $1$, $\frac{1}{x}$, $\infty$ bis zu einem dieser vier Werthe auf beliebigem Wege erstreckt eine Function $P\left(\begin{array}{ccc|} \alpha & \beta & \gamma \\ \alpha' & \beta' & \gamma' \end{array}\, x\right)$ bildet und bei passender Wahl dieser Grenzwerthe und des Weges von einem zum andern jede der sechs Functionen $P^{\alpha}$, $P^{\beta}$, \dots, $P^{\gamma'}$ darstellt.\footnote{Es l{"a}sst sich aber auch direct zeigen, dass das Integral die charakteristischen Eigenschaften einer solchen Function besitzt. Es wird dies in der Folge geschehen, wo dieser Ausdruck der $P$-function durch ein bestimmtes Integral zur Bestimmung der in $P^{\alpha}$, $P^{\alpha'}$, \dots\ noch willk{"u}rlich gebliebenen Factoren benutzt werden soll; und ich bemerke hier nur noch, dass es, um diesen Ausdruck allgemein anwendbar zu machen, einer Modification des Weges der Integration bedarf, wenn die Function unter dem Integralzeichen f{"u}r einen der Werthe $0$, $1$, $\frac{1}{x}$, $\infty$ so unendlich wird, dass sie die Integration bis an denselben nicht zul{"a}sst.}

Zufolge der im Art.~2 und dem vorigen erhaltenen Gleichungen
\[
P^{\alpha} \left(\begin{array}{ccc|} \alpha & \beta & \gamma \\ \alpha' & \beta' & \gamma' \end{array}\, x\right) = x^{\alpha} (1 - x)^{\gamma} P^{\alpha} \left(\begin{array}{ccc|} \alpha - \alpha & \beta + \alpha + \gamma & \gamma - \gamma \\ 0 & \beta' + \alpha + \gamma & 0 \end{array}\, x\right)
\]
\[
= \text{Const.} \cdot x^{\alpha} (1 - x)^{\gamma} F(\beta + \alpha + \gamma, \beta' + \alpha + \gamma, \alpha - \alpha' + 1, x)
\]
fliesst aus jedem Ausdrucke einer Function durch eine $P$-function eine Entwicklung derselben in eine hypergeometrische Reihe, welche nach steigenden Potenzen der Ver{"a}nderlichen in dieser $P$-function fortschreitet. Nach Art.~5 giebt es 8 Darstellungen einer Function durch $P$-functionen mit derselben Ver{"a}nderlichen, welche durch Vertauschung zusammengeh{"o}riger Exponenten aus einander erhalten werden, also z.~B.\ 8 Darstellungen mit der Ver{"a}nderlichen $x$. Von diesen liefern aber je zwei, welche durch Vertauschung ihres zweiten Paares, $\beta$ und $\beta'$, aus einander entstehen, dieselbe Entwicklung; man erh{"a}lt also vier Entwicklungen nach steigenden Potenzen von $x$, von denen zwei, welche durch Vertauschung von $\gamma$ und $\gamma'$ aus einander erhalten werden, die Function $P^{\alpha}$, die beiden andern die Function $P^{\alpha'}$ darstellen. Diese vier Entwicklungen convergiren, so lange der Modul von $x < 1$, und divergiren, wenn er gr{"o}sser als $1$ ist, w{"a}hrend die vier Reihen nach fallenden Potenzen von $x$, welche $P^{\beta}$ und $P^{\beta'}$ darstellen, sich umgekehrt verhalten. F{"u}r den Fall, wenn der Modul von $x$ gleich $1$ ist, folgt aus der Fourier'schen Reihe, dass die Reihen zu convergiren aufh{"o}ren, wenn die Function f{"u}r $x = 1$ unendlich von einer hohem Ordnung als der ersten wird, aber convergent bleiben, wenn sie nur unendlich von einer niedrigem Ordnung als $1$ wird oder endlich bleibt.\footnote{Dies folgt aus der bekannten Formel f{"u}r die Summe einer Reihe complexer Gr{"o}ssen: $S_n = S + \varrho_n e^{\psi_n i}$, worin $S$ die Summe der unendlichen Reihe, $\varrho_n$ das letzte Glied bedeutet, und aus der Bemerkung, dass $\varrho_n$ unendlich klein wird, wenn das letzte Glied der Reihe verschwindet, was bei der hier betrachteten hypergeometrischen Reihe eintritt, wenn der reelle Theil von $\gamma - \gamma'$ negativ ist.}

Es convergirt also auch in diesem Falle nur die H{"a}lfte der 8 Entwicklungen nach Potenzen von $x$, so lange der reelle Theil von $\gamma - \gamma'$ nicht zwischen $-1$ und $+1$ liegt, und sie convergiren s{"a}mmtlich, sobald dieses stattfindet.

\subsection*{9.}
\addcontentsline{toc}{subsection}{Art. 9}

Demnach hat man zur Darstellung einer $P$-function im Allgemeinen 24 verschiedene hypergeometrische Reihen, welche nach steigenden oder fallenden Potenzen von drei verschiedenen Gr{"o}ssen fortschreiten, und von denen f{"u}r einen gegebenen Werth von $x$ jedenfalls die H{"a}lfte, also zw{"o}lf convergiren. Im Falle I.\ Art.~5 sind alle diese Anzahlen mit $3$, im Falle II.\ mit $10$, im Falle III.\ mit $5$ zu multipliciren. Am geeignetsten zur numerischen Rechnung werden von diesen Reihen meistens diejenigen sein, deren viertes Element den kleinsten Modul hat.

Was die Ausdr{"u}cke einer $P$-function durch bestimmte Integrale betrifft, die sich durch die am Schlüsse des vorigen Art.\ aus den Transformationen des Art.~5 ableiten lassen, so sind diese Ausdr{"u}cke s{"a}mmtlich von einander verschieden. Man erh{"a}lt also im Allgemeinen $48$, im Falle I.\ $144$, im Falle II.\ $480$, im Falle III.\ $240$ bestimmte Integrale, welche dasselbe Glied einer $P$-function darstellen und also zu einander ein von $x$ unabh{"a}ngiges Verh{"a}ltniss haben. Von diesen lassen sich je $24$, welche durch eine gerade Anzahl von Vertauschungen der Exponenten aus einander hervorgehen, auch in einander transformiren durch eine solche Substitution ersten Grades, dass f{"u}r irgend drei von den Werthen $0$, $1$, $\infty$, $\frac{1}{x}$ der Integrationsver{"a}nderlichen $s$ die neue Ver{"a}nderliche die Werthe $0$, $1$, $\infty$ annimmt. Die {"u}brigen Gleichungen erfordern, soweit ich sie untersucht habe, zu ihrer Best{"a}tigung durch Methoden der Integralrechnung die Transformation von vielfachen Integralen.

\vspace{2em}
\begin{center}
\rule{0.5\textwidth}{0.4pt}
\end{center}
\vspace{1em}

% ============================================================
% V. Selbstanzeige der vorstehenden Abhandlung
% ============================================================

\section*{V. Selbstanzeige der vorstehenden Abhandlung.}
\addcontentsline{toc}{section}{V. Selbstanzeige der vorstehenden Abhandlung}

\begin{center}
\small
(G{"o}ttinger Nachrichten, 1857, Nr.~1.)
\end{center}
\vspace{0.5em}

Am 6.~November 1856 wurde der k{"o}niglichen Societ{"a}t eine von ihrem Assessor, Herrn Doctor Riemann, eingereichte mathematische Abhandlung vorgelegt, welche ``Beitr{"a}ge zur Theorie der durch die Gauss'sche Reihe $F(\alpha, \beta, \gamma, x)$ darstellbaren Functionen'' enth{"a}lt.

Diese Abhandlung ist einer Classe von Functionen gewidmet, welche bei der L{"o}sung mancher Aufgaben der mathematischen Physik gebraucht werden. Aus ihnen gebildete Reihen leisten bei schwierigeren Problemen dieselben Dienste, wie in den einfacheren F{"a}llen die jetzt so vielfach angewandten Reihen, welche nach Cosinus und Sinus der Vielfachen einer ver{"a}nderlichen Gr{"o}sse fortschreiten. Diese Anwendungen, namentlich astronomische, scheinen, nachdem schon Euler sich aus theoretischem Interesse mehrfach mit diesen Functionen besch{"a}ftigt hatte, Gauss zu seinen Untersuchungen {"u}ber dieselben veranlasst zu haben, von denen er einen Theil in seiner der K{"o}n.\ Soc.\ im J.\ 1812 {"u}bergebenen Abhandlung {"u}ber die Reihe, welche er durch $F(\alpha, \beta, \gamma, x)$ bezeichnet, ver{"o}ffentlicht hat.

Diese Reihe ist eine Reihe, in welcher der Quotient des $(n + 1)$ten Gliedes in das folgende
\[
= \frac{(n + \alpha)(n + \beta)}{(n + 1)(n + \gamma)}
\]
und das erste Glied $= 1$ ist. Die f{"u}r sie jetzt gew{"o}hnliche Benennung hypergeometrische Reihe ist schon fr{"u}her von Johann Friedrich Pfaff f{"u}r die allgemeineren Reihen vorgeschlagen worden, in denen der Quotient eines Gliedes in das folgende eine rationale Function des Stellenzeigers ist; w{"a}hrend Euler nach Wallis darunter eine Reihe verstand, in welcher dieser Quotient eine ganze Function ersten Grades des Stellenzeigers ist.

Der unver{"o}ffentlichte Theil der Gauss'schen Untersuchungen {"u}ber diese Reihe, welcher sich in seinem Nachlasse vorgefunden hat, ist 

\vspace{2em}
\begin{center}
\rule{0.5\textwidth}{0.4pt}
\end{center}

\vspace{1em}
\tableofcontents

\clearpage
\chapter*{Source segment 3: riemann pages 101 150}
\addcontentsline{toc}{chapter}{Source segment 3: riemann pages 101 150}
\selectlanguage{ngerman}
\maketitle
\thispagestyle{empty}
\tableofcontents
\newpage

%%%%%%%%%%%%%%%%%%%%%%%%%%%%%%%%%%%%%%%%%%%%%%%%%%%%%%%%%%%%%%%%%%%%%%%%
% V. Selbstanzeige der vorstehenden Abhandlung
%%%%%%%%%%%%%%%%%%%%%%%%%%%%%%%%%%%%%%%%%%%%%%%%%%%%%%%%%%%%%%%%%%%%%%%%
\section*{V. Selbstanzeige der vorstehenden Abhandlung.}
\addcontentsline{toc}{section}{V. Selbstanzeige der vorstehenden Abhandlung}
\label{sec:selbstanzeige}

unterdessen schon im J.\ 1835 durch die im 15.\ Bande des Journals
von Crelle enthaltenen Arbeiten Kummer's erg"anzt worden. Sie betreffen
die Ausdr"ucke der Reihe durch "ahnliche Reihen, in denen statt
des Elements $x$ eine algebraische Function dieser Gr"osse vorkommt.
Einen speciellen Fall dieser Umformungen hatte schon Euler aufgefunden
und in seiner Integralrechnung, so wie in mehren Abhandlungen behandelt
(in der einfachsten Gestalt in den N.\ Acta Acad.\ Petr.\ T.\ XII.\ p.\ 58);
und diese Relation ward sp"ater von Pfaff (Disquis.\ anal.\ Helmstadii
1797), Gudermann (Crelle J.\ Bd.\ 7.\ S.\ 306) und Jacobi auf
verschiedenen Wegen bewiesen. Kummer gelang es, die Methode Euler's
zu einem Verfahren auszubilden, durch welches s"ammtliche Transformationen
gefunden werden konnten; die wirkliche Ausf"uhrung desselben erforderte
aber so weitl"aufige Discussionen, da\ss{} er f"ur die Transformationen
dritten Grades von der Durchf"uhrung derselben abstand und sich begn"ugte,
die Transformationen ersten und zweiten Grades und die aus ihnen
zusammengesetzten vollst"andig abzuleiten.

In der anzuzeigenden Abhandlung wird auf diese Transcendenten eine
Methode angewandt, deren Princip in der Inaug.\ Diss.\ des Verfassers
(Art.\ 20) ausgesprochen worden ist und durch die sich s"ammtliche
fr"uher gefundenen Resultate fast ohne Rechnung ergeben. Einige weitere
mittelst derselben Methode gewonnenen Ergebnisse hofft der Verf.\ demn"achst
der K"oniglichen Societ"at vorlegen zu k"onnen.

\newpage
%%%%%%%%%%%%%%%%%%%%%%%%%%%%%%%%%%%%%%%%%%%%%%%%%%%%%%%%%%%%%%%%%%%%%%%%
% Anmerkungen zur Abhandlung IV
%%%%%%%%%%%%%%%%%%%%%%%%%%%%%%%%%%%%%%%%%%%%%%%%%%%%%%%%%%%%%%%%%%%%%%%%
\section*{Anmerkungen zur Abhandlung IV.}
\addcontentsline{toc}{section}{Anmerkungen zur Abhandlung IV}
\label{sec:anmerkungen}

\textbf{(1) (Zu Seite 73.)} In einer handschriftlichen Notiz von Riemann
(vom Juli 1866) finden sich die folgenden Formeln, die man aus~(3) erh"alt,
wenn man "uber die 5 willk"urlichen Coefficienten angemessen verf"ugt:
\[
\begin{aligned}
\alpha &= \frac{\sin(\alpha'+\beta'+\gamma')\pi}{\sin(\beta'-\beta)\pi},\\
\beta &= \frac{\sin(\alpha+\beta''+\gamma)\pi}{\sin(\beta'-\beta)\pi},\\
\gamma &= \frac{\sin(\alpha'+\beta+\gamma)\pi}{\sin(\gamma'-\gamma)\pi},\\
\alpha' &= \frac{\sin(\alpha+\beta+\gamma)\pi}{\sin(\beta'-\beta)\pi},\\
\beta' &= \frac{\sin(\alpha+\beta'+\gamma)\pi \cdot e^{(\alpha'+\gamma')\pi i}}{\sin(\gamma'-\gamma)\pi},\\
\gamma' &= \frac{\sin(\alpha+\beta+\gamma)\pi \cdot e^{(\alpha+\gamma)\pi i}}{\sin(\gamma'-\gamma)\pi}.
\end{aligned}
\]

\textbf{(2) (Zu Seite 81.)} Man erh"alt, wenn man zur Abk"urzung
\[
S = s^{-\alpha'-\beta'-\gamma'}(1-s)^{-\alpha'-\beta'-\gamma''}(1-xs)^{-\alpha''-\beta'-\gamma'}
\]
setzt, von Constanten Factoren abgesehen,
\[
\small
\begin{aligned}
p^\alpha &= x^\alpha(1-x)^\beta \int_0^1 S\,ds,\qquad
&&p^\beta = x^\alpha(1-x)^\beta \int_{-\infty}^0 S\,ds,\qquad
&&p^\gamma = x^{\alpha''}(1-x)^\beta \int_1^{\frac{1}{x}} S\,ds,\\
p^{\alpha'} &= x^{\alpha'}(1-x)^{\beta'} \int_0^1 S\,ds,\qquad
&&p^{\beta'} = x^{\alpha'}(1-x)^{\beta'} \int_{-\infty}^0 S\,ds,\qquad
&&p^{\gamma'} = x^{\alpha'}(1-x)^{\beta'} \int_1^{\frac{1}{x}} S\,ds.
\end{aligned}
\]
In jedem dieser Integrale kann die Bedeutung der mehrwerthigen
Function~$S$ in beliebiger Weise festgelegt werden. Verf"ugt man in
bestimmter Weise dar"uber, so erh"alt man zur Bestimmung constanter
Factoren:
\begingroup
\footnotesize
\begin{align*}
&\frac{(p^\alpha x^{-\alpha})'}{p^\alpha x^{-\alpha}} =
\frac{\Pi(-\alpha'-\beta'-\gamma')\,\Pi(-\alpha'-\beta''-\gamma')}
     {\Pi(\alpha' - \alpha)}
\cdot \frac{\Pi(\gamma' - \gamma)}
     {\Pi(-\alpha - \beta - \gamma')\,\Pi(-\alpha - \beta' - \gamma)},\\
&\frac{(p^\beta x^{-\beta})'}{p^\beta x^{-\beta}} =
\frac{\Pi(-\alpha'-\beta'-\gamma')\,\Pi(-\alpha - \beta' - \gamma')}
     {\Pi(\beta' - \beta)}
\cdot \frac{e^{(\alpha+\gamma)\pi i}\,\Pi(\gamma' - \gamma)}
     {\Pi(-\alpha - \beta - \gamma')\,\Pi(-\alpha' - \beta - \gamma)}
\cdot e^{(\alpha'+\beta'+\gamma')\pi i},\\
&\frac{(p^\gamma(1-x)^{-\gamma})'}{p^\gamma(1-x)^{-\gamma}} =
\frac{\Pi(-\alpha'-\beta'-\gamma')\,\Pi(-\alpha - \beta - \gamma')}
     {\Pi(\gamma' - \gamma)}
\cdot \frac{e^{(\alpha'+\beta'+\gamma')\pi i}}
     {\Pi(-\alpha - \beta' - \gamma')\,\Pi(-\alpha' - \beta - \gamma)}.
\end{align*}
\endgroup

Auch diese Formeln finden sich in Riemann's Papieren an verschiedenen
Stellen.

Man kann die Constanten $\alpha$,\ \ldots\ auch auf folgendem Wege bestimmen.
Nimmt man die Function~$S$ in dem Viereck $0$, $\infty$, $1$, $\frac{1}{x}$ der
nebenstehenden Figur an, so k"onnen die Zweige $P^\alpha$, $P^{\alpha'}$,
$P^\beta$, $P^{\beta'}$, $P^\gamma$, $P^{\gamma'}$ durch die in der Figur durch
die Pfeile angedeuteten Integrationen definirt werden.

Man liest dann unmittelbar aus der Figur die Relationen ab
\[
p^\alpha = p^\beta + p^\gamma,\qquad
p^{\alpha'} = p^{\beta'} + p^{\gamma'},\qquad
p^\beta = -p^{\beta'} - p^{\gamma'} = p^\alpha - p^{\gamma'},
\]
die zusammen mit den Formeln~(3) auf S.~73 zur Bestimmung der
Coefficienten $a_\alpha^\alpha$, $a_\beta^\alpha$, $a_\gamma^\alpha$, $a_\alpha^{\alpha'}$,
$a_\beta^{\alpha'}$, $a_\gamma^{\alpha'}$, $a_\alpha^\beta$, \ldots\ ausreichen.

\textbf{(3) (Zu Seite 82.)} Einen Integrationsweg, der in allen F"allen
statthaft ist, erh"alt man nach Pochhammer (Math.\ Annalen Bd.~35)
durch einen Doppelumlauf um zwei Verzweigungspunkte, wie ihn die
Figur zeigt. Ist die Integration bis~$a$ und~$b$ gestattet, so kann
dieser Integrationsweg so zusammengezogen werden, da\ss{} er aus vier
zwischen~$a$ und~$b$ verlaufenden Linien besteht. Bezeichnet man mit~$P$
das Integral "uber eine dieser Linien, so ist das "uber den Doppelumlauf
erstreckte Integral
\[
(1 - e^{2\alpha\pi i})(1 - e^{-2\beta\pi i})P.
\]
F.\ Klein hat diesen Darstellungen der $P$-functionen durch Einf"uhrung
homogener Variablen eine noch elegantere Fassung gegeben.
(Math.\ Annalen Bd.~38.)

\textbf{(4) (Zu Seite 82.)} Nach der Erg"anzung, die Dirichlet seinem
Convergenzbeweis der Fourier'schen Reihe in dem Zusatz zu der
Abhandlung "uber Kugelfunctionen gegeben hat (Crelle's Journal
Bd.~4, Dove's Repertorium Bd.~I, Crelle's Journal Bd.~17, Dirichlet's
Werke S.~117, 133, 305) kann eine periodische Function eines reellen
Argumentes, die in einem Punkt von niedrigerer als der ersten Ordnung
unendlich wird, in eine Fourier'sche Reihe entwickelt werden. Wendet
man diesen Satz an auf die Werthe, die eine in eine hypergeometrische
Reihe entwickelbare $P$-Function auf dem Einheitskreis hat, so ergiebt
sich eine Reihe, die von der nicht verschieden sein kann, die man erh"alt,
wenn man in der hypergeometrischen Reihe den Modul (absoluten Werth)
von~$x = 1$ setzt.

\newpage
%%%%%%%%%%%%%%%%%%%%%%%%%%%%%%%%%%%%%%%%%%%%%%%%%%%%%%%%%%%%%%%%%%%%%%%%
% VI. Theorie der Abel'schen Functionen
%%%%%%%%%%%%%%%%%%%%%%%%%%%%%%%%%%%%%%%%%%%%%%%%%%%%%%%%%%%%%%%%%%%%%%%%
\section*{VI. Theorie der Abel'schen Functionen.}
\addcontentsline{toc}{section}{VI. Theorie der Abel'schen Functionen}
\label{sec:abel}

\begin{center}
(Aus Borchardt's Journal f"ur reine und angewandte Mathematik, Bd.~54. 1857.)
\end{center}

\subsection*{I. Allgemeine Voraussetzungen und H"ulfsmittel f"ur die
Untersuchung von Functionen unbeschr"ankt ver"anderlicher Gr"ossen.}
\addcontentsline{toc}{subsection}{I. Allgemeine Voraussetzungen und H"ulfsmittel}

Die Absicht, den Lesern des Journals f"ur Mathematik Untersuchungen
"uber verschiedene Transcendenten, insbesondere auch "uber Abel'sche
Functionen vorzulegen, macht es mir w"unschenswerth, um Wiederholungen
zu vermeiden, eine Zusammenstellung der allgemeinen Voraussetzungen,
von denen ich bei ihrer Behandlung ausgehen werde, in einem besonderen
Aufsatze voraufschicken.

F"ur die unabh"angig ver"anderliche Gr"osse setze ich stets die jetzt
allgemein bekannte Gauss'sche geometrische Repr"asentation voraus,
nach welcher eine complexe Gr"osse~$z = x + yi$ vertreten wird durch
einen Punkt einer unendlichen Ebene, dessen rechtwinklige Coordinaten
$x$,~$y$ sind; ich werde dabei die complexen Gr"ossen und die sie
repr"asentirenden Punkte durch dieselben Buchstaben bezeichnen. Als
Function von~$x + yi$ betrachte ich jede Gr"osse~$w$, die sich mit ihr
der Gleichung
\begin{equation}
\label{eq:diff}
\frac{\partial w}{\partial x} = \frac{1}{i}\frac{\partial w}{\partial y}
\end{equation}
gem"ass "andert, ohne einen Ausdruck von~$w$ durch~$x$ und~$y$
vorauszusetzen.

Aus dieser Differentialgleichung folgt nach einem bekannten Satze, da\ss{}
die Gr"osse~$w$ durch eine nach ganzen Potenzen von~$z - a$
forschreitende Reihe von der Form $\sum a_n(z-a)^n$ darstellbar ist,
sobald sie in der Umgebung von~$a$ allenthalben einen bestimmten
mit~$z$ stetig sich "andernden Werth hat, und da\ss{} diese Darstellbarkeit
stattfindet bis zu einem Abst"ande von~$a$ oder Modul von~$z - a$, f"ur
welchen eine Unstetigkeit eintritt. Es ergiebt sich aber aus den
Betrachtungen, welche der Methode der unbestimmten Coefficienten zu
Grunde liegen, da\ss{} die Coefficienten~$a_n$ v"ollig bestimmt sind, wenn~$w$
in einer endlichen "ubrigens beliebig kleinen von~$a$ ausgehenden Linie
gegeben ist.

Beide Ueberlegungen verbindend, wird man sich leicht von der Richtigkeit
des Satzes "uberzeugen:

\begin{quote}
\textit{Eine Function von $x + yi$, die in einem Theile der $(x,y)$-Ebene
gegeben ist, kann dar"uber hinaus nur auf \emph{Eine} Weise stetig
fortgesetzt werden.}
\end{quote}

Man denke sich nun die zu untersuchende Function nicht durch irgend
welche $z$ enthaltende analytische Ausdr"ucke oder Gleichungen bestimmt,
sondern dadurch, da\ss{} der Werth der Function in einem beliebig begrenzten
Theile der $z$-Ebene gegeben ist und sie von dort aus stetig (der partiellen
Differentialgleichung \eqref{eq:diff} gem"ass) fortgesetzt wird. Diese
Fortsetzung ist nach den obigen S"atzen eine v"ollig bestimmte,
vorausgesetzt, da\ss{} sie nicht in blossen Linien geschieht, wobei eine
partielle Differentialgleichung nicht zur Anwendung kommen k"onnte,
sondern durch Fl"achenstreifen von endlicher Breite. Je nach der
Beschaffenheit der fortzusetzenden Function wird nun entweder die
Function f"ur denselben Werth von~$z$ immer wieder denselben Werth
annehmen, auf welchem Wege auch die Fortsetzung geschehen sein m"oge,
oder nicht.

Im ersteren Falle nenne ich sie \emph{einwerthig}, sie bildet dann eine
f"ur jeden Werth von~$z$ v"ollig bestimmte und nicht l"angs einer Linie
unstetige Function. Im letzteren Falle, wo sie \emph{mehrwerthig} heissen
soll, hat man, um ihren Verlauf aufzufassen, vor Allem seine Aufmerksamkeit
auf gewisse Punkte der $z$-Ebene zu richten, um welche herum sich die
Function in eine andere fortsetzt.

Ein solcher Punkt ist z.\,B.\ bei der Function~$\log(z-a)$ der Punkt~$a$.
Denkt man sich von diesem Punkte~$a$ aus eine beliebige Linie gezogen,
so wird man in der Umgebung von~$a$ den Werth der Function so w"ahlen
k"onnen, da\ss{} sie sich ausser dieser Linie "uberall stetig "andert; zu beiden
Seiten dieser Linie nimmt sie aber dann verschiedene Werthe an, auf der
negativen\footnote{Im Anschlusse an die von Gauss vorgeschlagene Benennung
\emph{positiv laterale Einheit} f"ur~$+i$ werde ich als positive
Seitenrichtung zu einer gegebenen Richtung diejenige bezeichnen, welche
zu ihr ebenso liegt, wie~$+i$ zu~$1$.} einen um~$2\pi i$ gr"osseren, als auf der
positiven. Die Fortsetzung der Function von einer Seite dieser Linie aus,
z.\,B.\ von der negativen, "uber sie hin"uber in das jenseitige Gebiet
giebt dann offenbar eine von der dort schon vorhandenen verschiedene
Function und zwar im hier betrachteten Falle eine allenthalben um~$2\pi i$
gr"ossere.

\medskip
\textbf{2.}

Zur bequemeren Bezeichnung dieser Verh"altnisse sollen die verschiedenen
Fortsetzungen einer Function f"ur denselben Theil der $z$-Ebene
\emph{Zweige} dieser Function genannt werden und ein Punkt, um welchen
sich ein Zweig einer Function in einen andern fortsetzt, eine
\emph{Verzweigungsstelle} dieser Function; wo keine Verzweigung stattfindet,
heisst die Function \emph{ein"andrig} oder \emph{monodrom}.

Ein Zweig einer Function von mehreren unabh"angig ver"anderlichen
Gr"ossen, $z$, $s$, $t$,\ \ldots\ ist ein"andrig in der Umgebung eines bestimmten
Werthensystems $z = a$, $s = b$, $t = c$,\ \ldots, wenn allen Werthencombinationen dieser Gr"ossen, die diesem Werthensysteme hinreichend benachbart
sind, nur ein Werth dieses Zweiges entspricht. Dies vorausgesetzt
ersch"opft die weitere Untersuchung dieser Function den in dem Obigen
f"ur den Fall einer unabh"angig ver"anderlichen Gr"osse behandelten
Gegenstand.

\medskip
\textbf{2.} (Fortsetzung: Zerschneidung einer mehrfach zusammenh"angenden
Fl"ache in eine einfach zusammenh"angende.)

Es m"ogen nun in einer Fl"ache, in welcher die zu untersuchende
Function urspr"unglich gegeben war, verschiedene Fl"achenst"ucke mit der
Begrenzung der Fl"ache zusammenh"angen. Eine solche Fl"ache werde
ich \emph{mehrfach zusammenh"angend} nennen. Um darin eine
\emph{einfach zusammenh"angende Fl"ache} zu bilden, m"ogen in derselben
von einem Begrenzungspunkte durch das Innere nach einem Begrenzungspunkte
Querschnitte, d.\,h.\ Linien, welche keinen Punkt mehrfach durchlaufen,
gezogen werden. Diese Querschnitte sollen so gew"ahlt werden, da\ss{}
ein zusammenh"angendes begrenztes Fl"achenst"uck entsteht, dessen
Begrenzung aus der urspr"unglichen Begrenzung der Fl"ache und aus
den beiden Seiten der Querschnitte besteht. Durch jeden solchen
Querschnitt wird offenbar die Anzahl der Zusammenh"ange der Fl"ache
um Eins vermindert, da man um ihn herum nicht mehr von einem
Theile der Fl"ache, der vorher mit dem andern zusammenh"angte, in
den andern gelangen kann, ohne die Begrenzung zu "uberschreiten, und
daher nach einer geraden Anzahl von Querschnitten eine einfach
zusammenh"angende Fl"ache entsteht.

Die hierdurch zerschnittene Fl"ache werde ich im Folgenden durch~$T'$
bezeichnen.

\medskip
\textbf{3.}

\textit{Bestimmung einer Function einer ver"anderlichen complexen Gr"osse
durch Grenz- und Unstetigkeitsbedingungen.}

Wenn in einer Ebene, in welcher die rechtwinkligen Coordinaten eines
Punkts~$x$,~$y$ sind, der Werth einer Function von~$x + yi$ in einer
endlichen Linie gegeben ist, so kann diese von dort aus nur auf eine
Weise stetig fortgesetzt werden und ist also dadurch v"ollig bestimmt
(siehe oben S.~89). Sie kann aber auch in dieser Linie nicht willk"urlich
angenommen werden, wenn sie von ihr aus einer stetigen Fortsetzung in
die anstossenden Fl"achentheile nach beiden Seiten hin f"ahig sein soll,
da sie durch ihren Verlauf in einem noch so kleinen endlichen Theile
dieser Linie schon f"ur den "ubrigen Theil bestimmt ist. Bei dieser
Bestimmungsweise einer Function sind also die zu ihrer Bestimmung
dienenden Bedingungen nicht von einander unabh"angig.

Als Grundlage f"ur die Untersuchung einer Transcendenten ist es vor
allen Dingen n"othig, ein System zu ihrer Bestimmung hinreichender von
einander unabh"angiger Bedingungen aufzustellen. Hierzu kann in vielen
F"allen, namentlich bei den Integralen algebraischer Functionen und ihren
inversen Functionen, ein Princip dienen, welches Dirichlet zur L"osung
dieser Aufgabe f"ur eine der Laplace'schen partiellen Differentialgleichung
gen"ugende Function von drei Ver"anderlichen, --- wohl durch einen "ahnlichen
gedanken von Gauss veranla\ss{}t --- in seinen Vorlesungen "uber die dem
umgekehrten Quadrat der Entfernung proportional wirkenden Kr"afte seit
einer Reihe von Jahren zu geben pflegt.

F"ur diese Anwendung auf die Theorie von Transcendenten ist jedoch
gerade ein Fall besonders wichtig, auf welchen dies Princip in seiner
dortigen einfachsten Form nicht anwendbar ist, und welcher dort als von
ganz untergeordneter Bedeutung unber"ucksichtigt bleiben kann. Dieser
Fall ist der, wenn die Function an gewissen Stellen des Gebiets, wo sie
zu bestimmen ist, vorgeschriebene Unstetigkeiten annehmen soll; was so
zu verstehen ist, da\ss{} sie an jeder solchen Stelle der Bedingung unterworfen
ist, unstetig zu werden, wie eine dort gegebene unstetige Function, oder
sich nur um eine dort stetige Function von ihr zu unterscheiden.

Ich werde hier das Princip in der f"ur die beabsichtigte Anwendung
erforderlichen Form darstellen und erlaube mir dabei in Betreff einiger
Nebenuntersuchungen auf die in meiner Doctordissertation (\textit{Grundlagen
f"ur eine allgemeine Theorie der Functionen einer ver"anderlichen complexen
Gr"osse. G"ottingen 1851}) gegebene Darstellung desselben zu verweisen.

Man nehme an, da\ss{} eine die $(x,y)$-Ebene einfach oder mehrfach
bedeckende beliebig begrenzte Fl"ache~$T$ und in derselben zwei f"ur jeden
ihrer Punkte eindeutig bestimmte reelle Functionen von~$x$,~$y$, die
Functionen $\alpha$ und~$\beta$ gegeben seien, und bezeichne das durch die Fl"ache~$T$
ausgedehnte Integral
\begin{equation}
\label{eq:omega}
\Omega(\alpha) = \int_T\Biggl[\Bigl(\frac{\partial\alpha}{\partial x}
- \frac{\partial\beta}{\partial y}\Bigr)^2 + \Bigl(\frac{\partial\alpha}{\partial y}
+ \frac{\partial\beta}{\partial x}\Bigr)^2\Biggr]\,dT
\end{equation}
durch~$\Omega(\alpha)$, wobei die Functionen $\alpha$ und~$\beta$ beliebige
Unstetigkeiten besitzen k"onnen, wenn nur das Integral dadurch nicht
unendlich wird. Es bleibt dann auch $\Omega(\alpha - \lambda)$ endlich, wenn~$\lambda$
allenthalben stetig ist und endliche Differentialquotienten hat. Wird diese
stetige Function~$\lambda$ der Bedingung unterworfen, nur in einem unendlich
kleinen Theile der Fl"ache~$T$ von einer unstetigen Function~$\gamma$
verschieden zu sein, so wird $\Omega(\alpha - \gamma)$ unendlich gro\ss{}, wenn~$\gamma$
l"angs einer Linie unstetig ist oder in einem Punkte so unstetig ist, da\ss{}

$\displaystyle \frac{\partial\gamma}{\partial x} + \frac{\partial\gamma}{\partial y}$

unendlich wird (Meine Inaug.\ Diss.\ Art.~17); es bleibt aber
$\Omega(\alpha - \lambda)$ endlich, wenn~$\gamma$ nur in einzelnen Punkten und
nur so unstetig ist, da\ss{}
\[
\int\Bigl[\Bigl(\frac{\partial\gamma}{\partial x}\Bigr)^2
+ \Bigl(\frac{\partial\gamma}{\partial y}\Bigr)^2\Bigr]\,dT
\]
durch die Fl"ache~$T$ erstreckt endlich bleibt, wie z.\,B.\ wenn~$\gamma$ in
der Umgebung eines Punktes im Abst"ande~$r$ von demselben
$= (-\log r)^\varepsilon$ und $0 < \varepsilon < \frac{1}{2}$ ist. Zur Abk"urzung m"ogen
hier die Functionen, in welche~$\lambda$ unbeschadet der Endlichkeit
von~$\Omega(\alpha - \lambda)$ "ubergehen kann, \emph{unstetig von der ersten Art},
die Functionen, f"ur welche dies nicht m"oglich ist, \emph{unstetig von der
zweiten Art} genannt werden.

Denkt man sich nun in $\Omega(\alpha - \mu)$ f"ur~$\mu$ alle stetigen oder von der
ersten Art unstetigen Functionen gesetzt, welche an der Grenze verschwinden,
so erh"alt dies Integral immer einen endlichen, aber seiner Natur nach nie
einen negativen Werth, und es mu\ss{} daher wenigstens einmal,
f"ur $\alpha - \mu = u$, ein Minimumwerth eintreten, so da\ss{}~$\Omega$ f"ur jede
Function $\alpha - \mu$, die unendlich wenig von~$u$ verschieden ist, gr"o\ss{}er
als~$\Omega(u)$ wird.

Bezeichnet daher~$\theta$ eine beliebige stetige oder von erster Art unstetige
Function des Orts in der Fl"ache~$T$, die an der Grenze allenthalben
gleich~$0$ ist, und~$h$ eine von~$x$,~$y$ unabh"angige Gr"osse, so mu\ss{}
$\Omega(u + h\theta)$ sowohl f"ur ein positives, als f"ur ein negatives
hinreichend kleines~$h$ gr"o\ss{}er als~$\Omega(u)$ werden, und daher in der
Entwicklung dieses Ausdrucks nach Potenzen von~$h$ der Coefficient von~$h$
verschwinden. Ist dieser~$0$, so ist
\[
\Omega(u + h\theta) = \Omega(u) + h^2 \Omega(\theta)
\]
und folglich~$\Omega$ immer ein Minimum.

Das Minimum tritt nur f"ur eine einzige Function~$u$ ein; denn f"ande
auch ein Minimum f"ur $u + \delta$ statt, so k"onnte $\Omega(u + \delta)$
nicht~$> \Omega(u)$ sein, weil sonst $\Omega(u + h\delta) < \Omega(u + \delta)$
f"ur $h < 1$ w"urde; also k"onnte $\Omega(u + \delta)$ nicht kleiner als die
anliegenden Werthe sein. Ist aber $\Omega(u + \delta) = \Omega(u)$, so mu\ss{}
$\delta$ constant, also da es in der Begrenzung~$0$ ist, "uberall~$0$ sein.
Es wird daher nur f"ur eine einzige Function~$u$ das Integral~$\Omega$ ein
Minimum und die Variation erster Ordnung oder das $h$ proportionale Glied
in $\Omega(u + h\theta)$:
\begin{equation}
\label{eq:var1}
\int\Biggl[\Bigl(\frac{\partial u}{\partial x} - \frac{\partial\beta}{\partial y}\Bigr)
\frac{\partial\theta}{\partial x}
+ \Bigl(\frac{\partial u}{\partial y} + \frac{\partial\beta}{\partial x}\Bigr)
\frac{\partial\theta}{\partial y}\Biggr]\,dT = 0.
\end{equation}

Aus dieser Gleichung folgt, da\ss{} das Integral
\[
\int\Bigl(\frac{\partial u}{\partial x} - \frac{\partial\beta}{\partial y}\Bigr)\,dy
- \Bigl(\frac{\partial u}{\partial y} + \frac{\partial\beta}{\partial x}\Bigr)\,dx
\]
durch die ganze Begrenzung eines Theils der Fl"ache~$T$ erstreckt stets
$= 0$ ist. Zerlegt man nun (nach der vorhergehenden Abhandlung) die
Fl"ache~$T$, wenn sie eine mehrfach zusammenh"angende ist, in eine einfach
zusammenh"angende~$T'$, so liefert die Integration durch das Innere von~$T'$
von einem festen Anfangspunkte bis zum Punkte~$(x,y)$ eine Function
von~$x$,~$y$,
\begin{equation}
v = \int\Biggl[\Bigl(\frac{\partial u}{\partial x} - \frac{\partial\beta}{\partial y}\Bigr)\,dx
+ \Bigl(\frac{\partial u}{\partial y} + \frac{\partial\beta}{\partial x}\Bigr)\,dy\Biggr]
+ \text{const.},
\end{equation}
welche in~$T'$ "uberall stetig oder unstetig von der ersten Art ist und
sich beim "Uberschreiten der Querschnitte um endliche von einem
Knotenpunkte des Schnittnetzes zum andern constante Gr"ossen "andert.

Es gen"ugt dann $v = \beta - v$ den Gleichungen
\[
\frac{\partial v}{\partial x} = -\frac{\partial u}{\partial y},\qquad
\frac{\partial v}{\partial y} = \frac{\partial u}{\partial x},
\]
und folglich ist $u + vi$ eine L"osung der Differentialgleichung
\[
i\,\frac{\partial(u + vi)}{\partial y} = \frac{\partial(u + vi)}{\partial x}
\]
oder eine Function von~$x + yi$.

Man erh"alt auf diesem Wege den in der erw"ahnten Abhandlung Art.~18
ausgesprochenen Satz:

\begin{quote}
\textit{Ist in einer zusammenh"angenden durch Querschnitte in eine einfach
zusammenh"angende~$T'$ zerlegten Fl"ache~$T$ eine complexe Function
$\alpha + \beta i$ von~$x$,~$y$ gegeben, f"ur welche}
\[
\int\Biggl[\Bigl(\frac{\partial\alpha}{\partial x}
- \frac{\partial\beta}{\partial y}\Bigr)^2 + \Bigl(\frac{\partial\alpha}{\partial y}
+ \frac{\partial\beta}{\partial x}\Bigr)^2\Biggr]\,dT
\]
\textit{durch die ganze Fl"ache ausgedehnt einen endlichen Werth hat, so kann
sie immer und nur auf Eine Art in eine Function von~$x + yi$ verwandelt
werden durch Subtraction einer Function~$\mu + \nu i$ von~$x$,~$y$, welche
folgenden Bedingungen gen"ugt:
\begin{enumerate}
\item[$1)$] $\mu$ ist am Rande~$= 0$ oder doch nur in einzelnen Punkten
davon verschieden, $\nu$ in Einem Punkte beliebig gegeben.
\item[$2)$] Die "Aenderungen von~$\mu$ sind in~$T$, von~$\nu$ in~$T'$ nur in
einzelnen Punkten und nur so unstetig, da\ss{}
\[
\int\Bigl[\Bigl(\frac{\partial\mu}{\partial x}\Bigr)^2
+ \Bigl(\frac{\partial\mu}{\partial y}\Bigr)^2\Bigr]\,dT
\quad\text{und}\quad
\int\Bigl[\Bigl(\frac{\partial\nu}{\partial x}\Bigr)^2
+ \Bigl(\frac{\partial\nu}{\partial y}\Bigr)^2\Bigr]\,dT
\]
durch die ganze Fl"ache erstreckt, endlich bleiben, und letztere l"angs der
Querschnitte beiderseits gleich.
\end{enumerate}}
\end{quote}

Wenn die Function~$\alpha + \beta i$, wo ihre Differentialquotienten unendlich
werden, unstetig wird, wie eine gegebene dort unstetige Function von~$x + yi$,
und keine durch eine Ab"anderung ihres Werthes in einem einzelnen Punkte
hebbare Unstetigkeit besitzt, so bleibt~$\Omega(\alpha)$ endlich, und es wird
$\mu + \nu i$ in~$T'$ allenthalben stetig. Denn da eine Function von~$x + yi$
gewisse Unstetigkeiten, wie z.\,B.\ Unstetigkeiten erster Art, gar nicht
annehmen kann (Meine Diss.\ Art.~12), so mu\ss{} die Differenz zweier solcher
Functionen stetig sein, sobald sie nicht von der zweiten Art unstetig ist.

Nach dem eben bewiesenen Satze l"asst sich daher eine Function von~$x + yi$
so bestimmen, da\ss{} sie im Innern von~$T$, von der Unstetigkeit des
imagin"aren Theils in den Querschnitten abgesehen, gegebene Unstetigkeiten
annimmt, und ihr reeller Theil an der Grenze einen dort allenthalben
beliebig gegebenen Werth erh"alt; wenn nur f"ur jeden Punkt, wo ihre
Differentialquotienten unendlich werden sollen, die vorgeschriebene
Unstetigkeit die einer gegebenen dort unstetigen Function von~$x + yi$ ist.
Die Bedingung an der Grenze kann man, wie leicht zu sehen, ohne eine
wesentliche Aenderung der gemachten Schl"usse durch manche andere ersetzen.

\newpage
\subsection*{4. Theorie der Abel'schen Functionen.}
\addcontentsline{toc}{subsection}{4.\ Theorie der Abel'schen Functionen}

In der folgenden Abhandlung habe ich die Abel'schen Functionen nach
einer Methode behandelt, deren Principien in meiner Inauguraldissertation\footnote{\textit{Grundlagen f"ur eine allgemeine Theorie der Functionen einer
ver"anderlichen complexen Gr"osse. G"ottingen 1851.}} aufgestellt und in einer
etwas ver"anderten Form in den drei vorhergehenden Aufs"atzen dargestellt
worden sind. Zur Erleichterung der Uebersicht schicke ich eine kurze
Inhaltsangabe vorauf.

Die erste Abtheilung enth"alt die Theorie eines Systems von
gleichverzweigten algebraischen Functionen und ihren Integralen, soweit
f"ur dieselbe nicht die Betrachtung von $\vartheta$-Reihen ma\ss{}gebend ist, und
handelt im \S\,1--5 von der Bestimmung dieser Functionen durch ihre
Verzweigungsart und ihre Unstetigkeiten, im \S\,6--10 von den rationalen
Ausdr"ucken derselben in zwei durch eine algebraische Gleichung verkn"upfte
ver"anderliche Gr"ossen, und im \S\,11--13 von der Transformation dieser
Ausdr"ucke durch rationale Substitutionen. Der bei dieser Untersuchung
sich darbietende Begriff einer \emph{Klasse} von algebraischen Gleichungen,
welche sich durch rationale Substitutionen in einander transformiren lassen,
d"urfte auch f"ur andere Untersuchungen wichtig und die Transformation
einer solchen Gleichung in Gleichungen niedrigsten Grades ihrer Klasse
(\S\,13) auch bei anderen Gelegenheiten von Nutzen sein. Diese Abtheilung
behandelt endlich im \S\,14--16 zur Vorbereitung der folgenden die
Anwendung des Abel'schen Additionstheorems f"ur ein beliebiges System
allenthalben endlicher Integrale von gleichverzweigten algebraischen
Functionen zur Integration eines Systems von Differentialgleichungen.

In der \emph{zweiten Abtheilung} werden f"ur ein beliebiges System von
immer endlichen Integralen gleichverzweigter, algebraischer, $(2p+1)$fach
zusammenh"angender Functionen die \emph{Jacobi'schen Umkehrungsfunctionen}
von~$p$ ver"anderlichen Gr"ossen durch $p$fach unendliche $\vartheta$-Reihen
ausgedr"uckt, d.\,h.\ durch Reihen von der Form
\[
\vartheta\bigl(v_1, v_2, \ldots, v_p\bigr) =
\sum_{-\infty}^{+\infty} e^{\bigl(\sum_{\mu} \sum_{\mu'}
  a_{\mu,\mu'} m_\mu m_{\mu'} + 2\sum_\mu v_\mu m_\mu\bigr)}\!,
\]
worin die Summationen im Exponenten sich auf~$\mu$ und~$\mu'$, die "ausseren
Summationen auf $m_1$, $m_2$, \ldots, $m_p$ beziehen. Es ergiebt sich, da\ss{}
zur allgemeinen L"osung dieser Aufgabe \emph{eine} Gattung von
$\vartheta$-Reihen ausreicht, in denen wenn $p > 3$ specielle --- zwischen
den $\frac{p(p+1)}{2}$ Gr"ossen~$a$ --- Relationen stattfinden, so da\ss{} nur
$3p - 3$ willk"urlich bleiben. Dieser Theil der Abhandlung bildet zugleich
eine Theorie dieser speciellen Gattung von $\vartheta$-Functionen; die allgemeinen
$\vartheta$-Functionen bleiben hier ausgeschlossen, lassen sich jedoch nach einer
ganz "ahnlichen Methode behandeln.

Das hier erledigte Jacobi'sche Umkehrungsproblem ist f"ur die
hyperelliptischen Integrale schon auf mehreren Wegen durch die beharrlichen
mit so sch"onem Erfolge gekr"onten Arbeiten von Weierstrass gel"ost worden,
von denen eine Uebersicht im 47.\ Bande des Journ.\ f"ur Math.\ (S.~289)
mitgetheilt worden ist. Es ist jedoch bis jetzt nur von dem Theile dieser
Arbeiten, welcher in den \S\S\,1 und~2 und der ersten die elliptischen
Functionen betreffenden H"alfte des \S\,3 der angef"uhrten Abhandlung
skizzirt wird, die wirkliche Ausf"uhrung ver"offentlicht (Bd.~52, S.~285
d.\ Journ.\ f.\ Math.); in wie weit zwischen den sp"ateren Theilen dieser
Arbeiten und meinen hier dargestellten eine Uebereinstimmung nicht bloss
in Resultaten, sondern auch in den zu ihnen f"uhrenden Methoden stattfindet,
wird gr"osstentheils erst die versprochene ausf"uhrliche Darstellung derselben
ergeben k"onnen.

Die gegenw"artige Abhandlung bildet mit Ausnahme der beiden letzten
\S\S.\ 26 und~27, deren Gegenstand damals nur kurz angedeutet werden
konnte, einen Auszug aus einem Theile meiner von Michaelis 1855 bis
Michaelis 1856 zu G"ottingen gehaltenen Vorlesungen. Was die Auffindung
der einzelnen Resultate betrifft, so wurde ich auf das im \S\,1--5, 9 und~12
Mitgetheilte und die dazu n"othigen vorbereitenden S"atze, welche sp"ater
Behufs der Vorlesungen so, wie es in dieser Abhandlung geschehen ist,
weiter ausgef"uhrt wurden, im Herbste 1851 und zu Anfang 1852 durch
Untersuchungen "uber die conforme Abbildung mehrfach zusammenh"angender
Fl"achen gef"uhrt, ward aber dann durch einen andern Gegenstand von dieser
Untersuchung abgezogen. Erst um Ostern 1855 wurde sie wieder aufgenommen
und in den Oster- und Michaelisferien jenes Jahres bis zu \S\,21 incl.\ fortgef"uhrt; das "Ubrige wurde bis Michaelis 1856 hinzugef"ugt. Einzelne
erg"anzende Zus"atze sind an manchen Stellen w"ahrend der Ausarbeitung
hinzugekommen.

\newpage
\subsection*{Erste Abtheilung.}
\addcontentsline{toc}{subsection}{Erste Abtheilung}

\medskip
\textbf{\S\,1.}

Ist~$s$ die Wurzel einer irreductibeln Gleichung $n$ten Grades, deren
Coefficienten ganze Functionen $m$ten Grades von~$z$ sind, so entsprechen
jedem Werthe von~$z$ $n$ Werthe von~$s$, die sich mit~$z$ "uberall, wo sie
nicht unendlich werden, stetig "andern. Stellt man daher (nach Art.\ 4
der Diss.) jeden Werth von~$z$ durch einen Punkt einer mit~$z$-Ebene
parallelen unendlichen Ebene, die $s$-Ebene, ab und theilt man die letztere
so in~$n$ Fl"achenst"ucke, da\ss{} in jedem derselben der Punkt einen ganz
bestimmten, sich stetig "andernden Werth von~$s$ repr"asentirt, so erh"alt
man eine die $s$-Ebene allenthalben $n$fach bedeckende und die $z$-Ebene
(in den kleinsten Theilen "ahnlich) abbildende Fl"ache~$T$. Diese Fl"ache
wird durch jeden Punkt, um welchen die Punkte von~$s$ mit einander
vertauscht werden, in eine solche Verzweigungspunkt genannte Punkt,
um welchen sich die Fl"ache in eine windende heissen.

Die Function~$s$ ist algebraisch und h"angt daher von einer Gleichung
$F(s,z) = 0$ ab. Die Fl"ache~$T$ ist durch ihre Verzweigungsart bestimmt,
und wenn man "uber die Anfangs- und Endpunkte der in ihr vorhandenen
Windsungen frei verf"ugen kann, so ist dadurch die Fl"ache v"ollig
gegeben, da man sie dann aus lauter einfach zusammenh"angenden
Bl"attern, welche "uber der $z$-Ebene ausgebreitet sind, zusammensetzen
kann. Die Verzweigungsart einer solchen Function~$s$ wird daher
vollst"andig dargestellt durch eine die $z$-Ebene bedeckende und in
Verzweigungspunkten windende Fl"ache~$T$.

\medskip
\textbf{\S\,2.}

Umgekehrt bilden f"ur jede beliebig gegebenen Verzweigungsart der
Fl"ache die Functionen, die sich mit~$z$ stetig "andern und ausser in den
Verzweigungspunkten allenthalben einwerthig sind, ein System gleich
verzweigter algebraischer Functionen. Dieses System enth"alt insbesondere
eine Function, die f"ur jeden Punkt der Fl"ache unendlich von der ersten
Ordnung wird und deren reeller Theil am Rande einen gegebenen Werth
annimmt.

\medskip
\textbf{\S\,3.}

Es werde nun die Fl"ache~$T$ in eine einfach zusammenh"angende
Fl"ache~$T'$ durch~$2p$ Querschnitte zerlegt. Durch jeden Punkt der
Fl"ache gehen dann $2p$ von einander unabh"angige Integrale erster
Gattung, d.\,h.\ allenthalben endliche Integrale rationaler Functionen
von~$s$ und~$z$. Die Periodicit"atsmoduln dieser Integrale an den
Querschnitten bilden ein System von $2p^2$ Gr"ossen, welche die fundamentale
Bedeutung f"ur die Theorie der Abel'schen Functionen haben.

\medskip
\textbf{\S\,4.}

Eine wie~$T$ verzweigte Function $\zeta$ von~$z$, die f"ur~$\mu$ Punkte dieser
Fl"ache unendlich von der ersten Ordnung wird, ist nach dem Fr"uheren
die Wurzel einer Gleichung von der Form
\[
G(\zeta, z) = 0,
\]
und nimmt daher jeden Werth f"ur~$\mu$ Punkte der Fl"ache~$T$ an.
Wenn man sich also jeden Punkt von~$T$ durch einen den Werth
von~$\zeta$ in diesem Punkte geometrisch repr"asentirenden Punkt einer
Ebene abgebildet denkt, so bildet die Gesammtheit dieser Punkte eine
in der $\zeta$-Ebene allenthalben $\mu$fach ausgebreitete und die Fl"ache~$T$
--- bekanntlich in den kleinsten Theilen "ahnlich --- abbildende
Fl"ache~$T_1$.

Jedem Punkt in der einen Fl"ache entspricht dann ein Punkt in der
andern. Die Functionen~$\omega$ oder die Integrale wie~$T$ verzweigter
Functionen von~$z$ gehen daher, wenn man f"ur~$z$ als unabh"angig
ver"anderliche Gr"osse~$z_1$ einf"uhrt, in Functionen "uber, welche in der
Fl"ache~$T_1$ allenthalben einen bestimmten Werth und dieselben
Unstetigkeiten haben, wie die Functionen~$\omega$ in den entsprechenden
Punkten von~$T$, und welche folglich Integrale wie~$T_1$ verzweigter
Functionen von~$z_1$ sind.

\medskip
\textbf{\S\,5.}

Es seien $\omega_1$, $\omega_2$, \ldots, $\omega_p$ von einander unabh"angige allenthalben
endliche Integrale rationaler Functionen von~$s$ und~$z$. Die allgemeinste
Function, die in der Fl"ache~$T$ nur f"ur $p$~Punkte unendlich von der
ersten Ordnung wird, l"asst sich dann als Quotient zweier solcher
Integrale darstellen. Die Anzahl der Bedingungsgleichungen, welche
eine Function bestimmen, die f"ur~$\mu$ Punkte unendlich von der ersten
Ordnung wird, betr"agt $p + \mu - 1$; denn es m"ussen $2\mu - p - 1$
von den~$2\mu$ reellen Constanten der Function willk"urlich bleiben.

F"ur die Bestimmung einer solchen Function sind also $2\mu - p - 1$
Bedingungen erforderlich. Soll die Function f"ur m"oglichst wenig Punkte
unendlich werden, so mu\ss{} $\mu$ den kleinstm"oglichen Werth haben, und
es zeigt sich, da\ss{} dieser Werth~$\mu$ nicht kleiner als $p + 1$ sein kann.

\medskip
\textbf{\S\,6.}

Es seien $s_1$, $s_2$, \ldots, $s_n$ die $n$ Werthe der Function~$s$ f"ur dasselbe~$z$,
und bezeichnet~$\vartheta$ eine beliebige Gr"osse, so ist
$(\vartheta - s_1)(\vartheta - s_2)\cdots(\vartheta - s_n)$ eine einwerthige Function von~$z$,
die nur f"ur einen Punkt der $z$-Ebene, der mit einem Punkte~$\varepsilon$
zusammenf"allt, unendlich wird und unendlich von einer so hohen
Ordnung, als Punkte~$\varepsilon$ auf ihn fallen.

Bezeichnet man nun die Werthe von~$z$ in den Punkten~$\varepsilon$, wo~$s$ nicht
unendlich ist, durch $\varepsilon_1$, $\varepsilon_2$, \ldots, $\varepsilon_\nu$ und
$(z - \varepsilon_1)(z - \varepsilon_2)\cdots(z - \varepsilon_\nu)$ durch~$g_0$, so ist
$g_0(\vartheta - s_1)\cdots(\vartheta - s_n)$ eine einwerthige Function von~$z$, die
f"ur alle endlichen Werthe von~$z$ endlich ist und f"ur $z = \infty$ unendlich
von der $m$ten Ordnung wird, also eine ganze Function $m$ten Grades
von~$z$. Sie ist zugleich eine ganze Function $n$ten Grades von~$\vartheta$,
die f"ur $\vartheta = s$ verschwindet. Bezeichnet man sie durch~$F$ und, wie
wir in der Folge thun wollen, eine ganze Function $n$ten Grades von~$\vartheta$
und $m$ten Grades von~$z$ durch $F(\vartheta, z)$, so ist~$s$ die Wurzel der
Gleichung $F(s,z) = 0$.

Die Function~$F$ ist eine Potenz einer \emph{unzerf"allbaren} --- d.\,h.\ nicht
als ein Product aus ganzen Functionen von~$\vartheta$ und~$z$ darstellbaren ---
Function. Denn jeder ganze rationale Factor von~$F(\vartheta,z)$ bildet, da er
f"ur einige der Wurzeln $s_1$, $s_2$, \ldots, $s_n$ verschwinden mu\ss{}, f"ur
$\vartheta = s$ eine Function von~$z$, die in einem Theile der Fl"ache~$T$
verschwindet und folglich, da diese Fl"ache zusammenh"angend ist, in
der ganzen Fl"ache~$0$ sein mu\ss{}. Zwei unzerf"allbare Factoren
von~$F(\vartheta,z)$ k"onnten aber nur f"ur eine endliche Anzahl von Werthenpaaren
zugleich verschwinden, wenn die eine nicht durch Multiplication mit einer
Constanten aus der andern erhalten werden k"onnte. Folglich mu\ss{} $F$
eine Potenz einer unzerf"allbaren Function sein.

Wenn der Exponent~$\nu$ dieser Potenz~$> 1$ ist, so wird die
Verzweigungsart der Function~$s$ nicht dargestellt durch die Fl"ache~$T$,
sondern durch eine in der $z$-Ebene allenthalben $\nu$fach ausgebreitete
Fl"ache~$t$, in welcher die Fl"ache~$T$ allenthalben $\nu$fach ausgebreitet ist.
Es kann dann zwar~$s$ als eine wie~$T$ verzweigte Function betrachtet
werden, nicht aber umgekehrt $T$ als verzweigt, wie~$s$.

\medskip
\textbf{\S\,7.}

Eine solche nur in einzelnen Punkten von~$T$ unstetige Function,
wie~$s$, ist auch $\displaystyle\frac{ds}{dz}$. Denn diese Function nimmt zu beiden
Seiten der Querschnitte und der Linien~$l$ denselben Werth an, da die
Differenz der beiden Werthe von~$\omega$ in diesen Linien l"angs denselben
constant ist; sie kann nur unendlich werden, wo~$\omega$ unendlich wird,
und in den Verzweigungspunkten der Fl"ache und ist sonst allenthalben
stetig, da die Derivirte einer ein"andrig und endlich bleibenden Function
ebenfalls ein"andrig und endlich bleibt.

Es sind daher s"ammtliche Functionen~$\omega$ algebraische wie~$T$
verzweigte Functionen von~$z$ oder Integrale solcher Functionen.
Dieses System von Functionen ist bestimmt, wenn die Fl"ache~$T$ gegeben
ist und h"angt nur von der Lage ihrer Verzweigungspunkte ab.

\medskip
\textbf{\S\,8.}

Es sei jetzt die irreductible Gleichung $F(s,z) = 0$ gegeben und die
Art der Verzweigung der Function~$s$ oder der sie darstellenden Fl"ache~$T$
zu bestimmen. Wenn f"ur einen Werth~$\beta$ von~$z$ $\mu$~Zweige der Function
zusammenh"angen, so da\ss{} einer dieser Zweige sich erst nach~$\mu$ Uml"aufen
des~$z$ um~$\beta$ wieder in sich selbst fortsetzt, so k"onnen diese~$\mu$ Zweige
der Function (wie nach Cauchy oder durch die Fourier'sche Reihe leicht
bewiesen werden kann) dargestellt werden durch eine Reihe nach steigenden
rationalen Potenzen von~$z - \beta$ mit Exponenten vom kleinsten
gemeinschaftlichen Nenner~$\mu$, und umgekehrt.

Ein Punkt der Fl"ache~$T$, in welchem nur zwei Zweige einer Function
zusammenh"angen, so da\ss{} sich um diesen Punkt der erste in den zweiten
und dieser in jenen fortsetzt, heisse ein \emph{einfacher Verzweigungspunkt}.

Ein Punkt der Fl"ache, um welchen sie sich $(\mu+1)$mal windet,
kann dann angesehen werden als~$\mu$ zusammengefallene (oder unendlich
nahe) einfache Verzweigungspunkte.

Um dies zu zeigen, seien in einem diesen Punkt umgebenden St"ucke der
$z$-Ebene $s_1$, $s_2$, \ldots, $s_{\mu+1}$ ein"andrige Zweige der Function~$s$
und in der Begrenzung desselben, bei positiver Umschreibung auf einander
folgend, $a_1$, $a_2$, \ldots, $a_\mu$ einfache Verzweigungspunkte. Durch einen
positiven Umlauf um~$a_1$ werde $s_1$ mit~$s_2$, um~$a_2$ $s_2$ mit~$s_3$, \ldots,
um~$a_\mu$, $s_\mu$ mit~$s_{\mu+1}$ vertauscht. Es gehen dann nach einem
positiven Umlaufe um ein alle diese Punkte (und keinen andern
Verzweigungspunkt) enthaltendes Gebiet
\[
s_1,\; s_2,\; \ldots,\; s_\mu,\; s_{\mu+1}
\]
in
\[
s_2,\; s_3,\; \ldots,\; s_{\mu+1},\; s_1
\]
"uber, und es entsteht daher, wenn sie zusammenfallen, ein $\mu$facher
Windungspunkt.

\medskip
\textbf{\S\,9.}

Die Eigenschaften der Functionen~$\omega$ h"angen wesentlich davon ab,
wie vielfach zusammenh"angend die Fl"ache~$T$ ist. Um dies zu entscheiden,
wollen wir zun"achst die Anzahl der einfachen Verzweigungspunkte der
Function~$s$ bestimmen.

In einem Verzweigungspunkte nehmen die dort zusammenh"angenden Zweige
der Function denselben Werth an, und es werden daher zwei oder mehrere
Wurzeln der Gleichung
\[
F(s) = a_0 s^n + a_1 s^{n-1} + \cdots + a_n = 0
\]
einander gleich. Dies kann nur geschehen, wenn
\[
F'(s) = a_0 n s^{n-1} + a_1 (n-1) s^{n-2} + \cdots + a_{n-1} = 0
\]
oder die einwerthige Function von~$z$,
\[
F'(s_1)\, F'(s_2) \cdots F'(s_n),
\]
verschwindet. Diese Function wird f"ur endliche Werthe von~$z$ nur unendlich,
 wenn~$s = \infty$, also $a_0 = 0$ ist und mu\ss{}, um endlich zu bleiben,
mit $a_0^{n-2}$ multiplicirt werden. Sie wird dann eine einwerthige, f"ur
ein endliches~$z$ endliche Function von~$z$, welche f"ur $z = \infty$ unendlich
von der $2m(n-1)$ten Ordnung wird, also eine ganze Function
$2m(n-1)$ten Grades. Die Werthe von~$z$, f"ur welche $F'(s)$ und~$F(s)$
gleichzeitig verschwinden, sind also die Wurzeln der Gleichung
\begin{equation}
\label{eq:Q}
Q(z) = a_0^{n-2} \prod_{i=1}^{n} F'(s_i) = 0
\end{equation}
oder auch, da $\displaystyle F'(s_i) = a_0 \prod_{i' \neq i}(s_i - s_{i'})$, die aus
$F'(s) = 0$ und $F(s) = 0$ durch Elimination von~$s$ gebildete werden kann.

Wird $F(s,z) = 0$ f"ur $s = \alpha$, $z = \beta$, so ist
\[
F(s,z) = F(\alpha,\beta) + \frac{\partial F}{\partial s}(s-\alpha)
+ \frac{\partial F}{\partial z}(z-\beta) + \cdots.
\]
Ist also f"ur $(s = \alpha, z = \beta)$ gleichzeitig $\displaystyle\frac{\partial F}{\partial s} = 0$
und $\displaystyle\frac{\partial F}{\partial z}$ nicht $= 0$, so wird $s - \alpha$ unendlich klein, wie
$(z - \beta)^{\frac{1}{2}}$, und findet also ein einfacher Verzweigungspunkt statt.
Es werden zugleich in dem Producte $Q(z) = \prod F'(s_i)$ zwei Factoren
dadurch den Factor $(z - \beta)$ unendlich klein.

In dem Falle, da\ss{} $\displaystyle\frac{\partial F}{\partial s}$ und $\displaystyle\frac{\partial F}{\partial z}$ nie
verschwinden, wenn gleichzeitig $F = 0$ und $\displaystyle\frac{\partial F}{\partial s} = 0$ werden,
entspricht demnach jedem linearen Factor von~$Q(z)$ ein einfacher
Verzweigungspunkt, und die Anzahl dieser Punkte ist also $= 2m(n-1)$.

Die Lage der Verzweigungspunkte h"angt von den Coefficienten der
Potenzen von~$z$ in den Functionen~$a$ ab und "andert sich stetig mit denselben.

Wenn diese Coefficienten solche Werthe annehmen, da\ss{} zwei demselben
Zweigepaare angeh"orige einfache Verzweigungspunkte zusammenfallen, so
heben diese sich auf, und es werden zwei Wurzeln von~$F(s)$ einander
gleich, ohne da\ss{} eine Verzweigung stattfindet. Setzt sich um jeden von
ihnen $s_1$ in~$s_2$ und $s_2$ in~$s_1$ fort, so geht durch einen Umlauf um ein
beide enthaltendes St"uck der $z$-Ebene $s_1$ in~$s_2$ und $s_2$ in~$s_1$ "uber,
und beide Zweige werden ein"andrig, wenn sie zusammenfallen.

Es bleibt dann also auch ihre Derivirte~$\dfrac{ds}{dz}$ ein"andrig und endlich,
und folglich wird $\dfrac{\partial F}{\partial z} = -\dfrac{\partial F}{\partial s} \cdot \dfrac{ds}{dz} = 0$.

Wird $F = \dfrac{\partial F}{\partial s} = \dfrac{\partial F}{\partial z} = 0$ f"ur $s = \alpha$, $z = \beta$,
so ergeben sich aus den drei folgenden Gliedern der Entwicklung von
$F(s,z)$ zwei Werthe f"ur $\dfrac{ds}{dz}$ $(s = \alpha, z = \beta)$. Sind diese
Werthe ungleich und endlich, so k"onnen die beiden Zweige der Function~$s$,
denen sie angeh"oren, dort nicht zusammenh"angen und sich nicht verzweigen.
Es wird dann $\dfrac{ds}{dz}$ f"ur beide unendlich klein wie $z - \beta$, und $Q(\beta)$
erh"alt dadurch den Factor $(z - \beta)^2$; es fallen also nur zwei einfache
Verzweigungspunkte zusammen.

Um in jedem Falle, wenn f"ur $z = \beta$ mehrere Wurzeln der Gleichung
$F(s) = 0$ gleich~$\alpha$ werden, zu entscheiden, wie viele einfache
Verzweigungspunkte f"ur $(s = \alpha, z = \beta)$ zusammenfallen, und wie viele
von diesen sich aufheben, mu\ss{} man diese Wurzeln (nach dem Verfahren
von Lagrange\footnote{Nach einer handschriftlichen Bemerkung Riemann's
ist dieses Verfahren von Puiseux in dessen Abhandlung \selectlanguage{latin}\textit{Recherches
sur les fonctions alg\'{e}briques}, Journ.\ de Math.\ 1850, durchgef"uhrt.
Riemann hatte offenbar die Puiseux'sche Arbeit erst nach Abfassung des
Textes kennen gelernt.\selectlanguage{ngerman}}) soweit nach steigenden Potenzen von $z - \beta$
entwickeln, bis diese Entwicklungen s"ammtlich von einander verschieden
werden, wodurch sich die wirklich noch stattfindenden Verzweigungen ergeben.
Und man mu\ss{} dann untersuchen, von welcher Ordnung $F'(s)$ f"ur jede
dieser Wurzeln unendlich klein wird, um die Anzahl der ihnen zugeh"origen
linearen Factoren von~$Q(z)$ oder der f"ur $(s = \alpha, z = \beta)$
zusammengefallenen einfachen Verzweigungspunkte zu bestimmen.

Bezeichnet die Zahl~$q$, wie oft sich die Fl"ache~$T$ um den Punkt
$(s,z)$ windet, so wird f"ur diesen Punkt der Ausdruck
$\dfrac{Q(z)}{(z-\beta)^{q-1}}$ weder verschwinden noch unendlich.

\medskip
\textbf{\S\,10.}

F"ur die Functionen, die nur in einzelnen Punkten der Fl"ache~$T$
unendlich werden, gilt folgender wichtiger Satz:

\begin{quote}
\textit{Eine wie~$T$ verzweigte Function~$\zeta$ von~$z$, die f"ur~$\mu$ Punkte
dieser Fl"ache unendlich von der ersten Ordnung wird, enth"alt in ihrem
allgemeinen Ausdrucke $2\mu - p + 1$ willk"urliche Constanten.}
\end{quote}

In der That, die Differenz zweier solcher Functionen ist eine Function,
die f"ur dieselben~$\mu$ Punkte unendlich wird und in den $2\mu$ reellen
Theilen dieser Unstetigkeiten gegebene Gr"ossen hat. Damit diese Differenz
allenthalben endlich sei, mu\ss{} sie identisch verschwinden, und dazu sind
$2\mu$ reelle Bedingungen erforderlich. Da nun eine allgemeine solche
Function $4\mu$ reelle Constanten enth"alt, so bleiben $4\mu - 2\mu = 2\mu$
"ubrig, die sich jedoch nicht unabh"angig voneinander "andern, da die
reellen und imagin"aren Theile durch die Cauchy-Riemann'schen
Differentialgleichungen mit einander verkn"upft sind. Es bleiben daher
$2\mu - p + 1$ willk"urliche Constanten.

Eine solche Function ist die Wurzel einer Gleichung $n$ten Grades,
deren Coefficienten ganze Functionen $m$ten Grades von~$z$ sind. Die
Anzahl der in einer Function~$s$, die nur f"ur~$m$ Punkte der Fl"ache~$T$
unendlich von der ersten Ordnung wird und "ubrigens stetig bleibt,
enthaltenen willk"urlichen Constanten ist in allen F"allen $= 2m - p + 1$.

\medskip
\textbf{\S\,11.}

Eine wie~$T$ verzweigte Function~$\zeta$ von~$z$, die f"ur $\mu_1$~Punkte
dieser Fl"ache unendlich von der ersten Ordnung wird, ist nach dem
Fr"uheren die Wurzel einer $\mu_1$ten Gleichung von der Form
\[
G(\zeta_1, z) = 0
\]
und nimmt daher jeden Werth f"ur $\mu_1$~Punkte der Fl"ache~$T$ an. Wenn
man sich also jeden Punkt von~$T$ durch einen den Werth von~$\zeta_1$ in
diesem Punkte geometrisch repr"asentirenden Punkt einer Ebene abgebildet
denkt, so bildet die Gesammtheit dieser Punkte eine in der
$\zeta_1$-Ebene allenthalben $\mu_1$fach ausgebreitete und die Fl"ache~$T$
--- bekanntlich in den kleinsten Theilen "ahnlich --- abbildende Fl"ache~$T_1$.

Bezeichnet~$s_1$ irgend eine andere wie~$T$ verzweigte Function von~$z$,
die f"ur $m_1$~Punkte von~$T$ und also auch von~$T_1$ unendlich von der
ersten Ordnung wird, so findet (\S\,5) zwischen~$s_1$ und~$\zeta_1$ eine
Gleichung von der Form
\[
F_1(s_1, \zeta_1) = 0
\]
statt, worin $F_1$ eine Potenz einer unzerf"allbaren ganzen Function von~$s_1$
und~$\zeta_1$ ist, und es lassen sich, wenn diese Potenz die erste ist, alle
wie~$T_1$ verzweigten Functionen von~$\zeta_1$, folglich alle rationale Functionen
von~$s$ und~$z$ rational in~$s_1$ und~$\zeta_1$ ausdr"ucken (\S\,8).

Die Gleichung $F(s,z) = 0$ kann also durch eine rationale Substitution
in $F_1(s_1, \zeta_1) = 0$ und diese in jene transformirt werden.

\medskip
\textbf{\S\,12.}

Man betrachte nun als zu \emph{Einer Klasse} geh"orend alle irreductiblen
algebraischen Gleichungen zwischen zwei ver"anderlichen Gr"ossen, welche
sich durch rationale Substitutionen in einander transformiren lassen, so
da\ss{} $F(s,z) = 0$ und $F_1(s_1, z_1) = 0$ zu derselben Klasse geh"oren,
wenn sich f"ur~$s$ und~$z$ solche rationale Functionen von~$s_1$ und~$z_1$
setzen lassen, da\ss{} $F(s,z) = 0$ in $F_1(s_1, z_1) = 0$ "ubergeht und
zugleich $s_1$ und~$z_1$ rationale Functionen von~$s$ und~$z$ sind.

Die rationalen Functionen von~$s$ und~$z$ bilden, als Functionen von
irgend einer von ihnen~$\zeta$ betrachtet, ein System gleich verzweigter
algebraischer Functionen. Auf diese Weise f"uhrt jede Gleichung offenbar
zu einer Klasse von Systemen gleichverzweigter algebraischer Functionen,
welche sich durch Einf"uhrung einer Function des Systems als unabh"angig
ver"anderlicher Gr"osse in einander transformiren lassen und zwar alle
Gleichungen \emph{Einer} Klasse zu derselben Klasse von Systemen algebraischer
Functionen, und umgekehrt f"uhrt (\S\,11) jede Klasse von solchen Systemen
zu \emph{Einer} Klasse von Gleichungen.

Ist das Gr"ossengebiet $(s,z)$ $(2p+1)$fach zusammenh"angend und die
Function~$\zeta$ in~$\mu$ Punkten desselben unendlich von der ersten Ordnung,
so ist die Anzahl der Verzweigungswerthe der gleichverzweigten Functionen
von~$z$, welche durch die "ubrigen rationalen Functionen von~$s$ und~$z$
gebildet werden, $2(\mu + p - 1)$, und die Anzahl der willk"urlichen Constanten
in der Function~$\zeta$ $2\mu - p + 1$ (\S\,5). Diese lassen sich so bestimmen,
da\ss{} $2\mu - p + 1$ Verzweigungswerthe gegebene Werthe annehmen, wenn
diese Verzweigungswerthe von einander unabh"angige Functionen von ihnen
sind, und zwar nur auf eine endliche Anzahl Arten, da die Bedingungsgleichungen
algebraisch sind.

In jeder Klasse von Systemen gleichverzweigter $(2p+1)$fach
zusammenh"angender Functionen giebt es daher eine endliche Anzahl von
Systemen $\mu$werthiger Functionen, in welchen $2\mu - p + 1$
Verzweigungswerthe gegebene Werthe annehmen. Wenn andererseits die
$2(\mu + p - 1)$ Verzweigungspunkte einer die $z$-Ebene allenthalben
$\mu$fach bedeckenden $(2p+1)$fach zusammenh"angenden Fl"ache beliebig
gegeben sind, so giebt es (\S\S\,3--5) immer ein System wie diese Fl"ache
verzweigter algebraischer Functionen von~$z$. Die $3p - 3$ "ubrigen
Verzweigungswerthe in jenen Systemen gleichverzweigter $\mu$werthiger
Functionen k"onnen daher beliebige Werthe annehmen; und es h"angt also
eine \emph{Klasse} von Systemen gleichverzweigter $(2p+1)$fach zusammenh"angender
Functionen und die zu ihr geh"orende Klasse algebraischer Gleichungen
von $3p - 3$ stetig ver"anderlichen Gr"ossen ab, welche die \emph{Moduln}
dieser Klasse genannt werden sollen.

\medskip
\textbf{\S\,13.}

Diese Bestimmung der Anzahl der Moduln einer Klasse $(2p+1)$fach
zusammenh"angender algebraischer Functionen gilt jedoch nur unter der
Voraussetzung, da\ss{} es $2\mu - p + 1$ Verzweigungswerthe giebt, welche von
einander unabh"angige Functionen der willk"urlichen Constanten in der
Function~$\zeta$ sind. Diese Voraussetzung trifft nur zu, wenn $p > 1$, und
die Anzahl der Moduln ist nur dann $= 3p - 3$, f"ur $p = 1$ aber~$= 1$.
Die directe Untersuchung derselben wird indess schwierig durch die Art und
Weise, wie die willk"urlichen Constanten in~$\zeta$ enthalten sind. Man f"uhre
deshalb in einem Systeme gleich verzweigter $(2p+1)$fach zusammenh"angender
Functionen, um die Anzahl der Moduln zu bestimmen, als unabh"angig
ver"anderliche Gr"osse nicht eine dieser Functionen, sondern ein allenthalben
endliches Integral einer solchen Function ein.

Die Werthe, welche die Function~$\omega$ von~$z$ innerhalb der Fl"ache~$T'$
annimmt, werden geometrisch repr"asentirt durch eine einen endlichen Theil
der $w$-Ebene einfach oder mehrfach bedeckende und die Fl"ache~$T'$ (in
den kleinsten Theilen "ahnlich) abbildende Fl"ache, welche durch~$S$
bezeichnet werden soll. Da~$\omega$ auf der positiven Seite des $\nu$ten
Querschnitts um die Constante~$\Pi^{(\nu)}$ gr"osser ist, als auf der negativen,
so besteht die Begrenzung von~$S$ aus Paaren von parallelen Curven,
welche denselben Theil des $T'$ begrenzenden Schnittsystems abbilden, und
es wird die Ortsverschiedenheit der entsprechenden Punkte in den parallelen,
den $\nu$ten Querschnitt abbildenden Begrenzungstheilen von~$S$ durch die
complexe Gr"osse~$\Pi^{(\nu)}$ ausgedr"uckt. Die Anzahl der einfachen
Verzweigungspunkte der Fl"ache~$S$ ist~$2p - 2$, da $d\omega$ in $2p - 2$
Punkten der Fl"ache~$T$ unendlich klein von der zweiten Ordnung wird.
Die rationalen Functionen von~$s$ und~$z$ sind dann Functionen von~$\omega$,
welche f"ur jeden Punkt von~$S$ \emph{Einen} bestimmten Werth haben und
"uberall, mit Ausnahme der Verzweigungspunkte und der Unstetigkeitspunkte
der Fl"ache~$T$, stetig sind.

\medskip
\textbf{\S\,14.}

Das Abel'sche Additionstheorem f"ur ein System allenthalben endlicher
Integrale von gleichverzweigten algebraischen Functionen l"asst sich
folgenderma\ss{}en aussprechen. Es seien $\omega_1$, $\omega_2$, \ldots, $\omega_p$ von
einander unabh"angige allenthalben endliche Integrale; es seien ferner
$(s_1,z_1)$, $(s_2,z_2)$, \ldots, $(s_m,z_m)$ $m$~beliebige Werthenpaare, und
$(\sigma_1,\zeta_1)$, $(\sigma_2,\zeta_2)$, \ldots, $(\sigma_m,\zeta_m)$ seien durch die
$p$~Gleichungen
\begin{equation}
\label{eq:abel}
\sum_{\mu=1}^{m} \omega_\mu^{(s_\mu,z_\mu)} =
\sum_{\mu=1}^{m} \omega_\mu^{(\sigma_\mu,\zeta_\mu)}
\quad (\mu = 1, 2, \ldots, p)
\end{equation}
verbunden, worin die Integrale von einem festen Anfangspunkte bis zu den
entsprechenden Endpunkten erstreckt sind. Es wird behauptet, da\ss{} dann
die Gr"ossenpaare $(\sigma,\zeta)$ algebraische Functionen der Gr"ossenpaare
$(s,z)$ sind.

\medskip
\textbf{\S\,15.}

F"ur ein beliebiges System von immer endlichen Integralen
 gleichverzweigter algebraischer Functionen erh"alt man den allgemeinsten
Ausdruck einer Function, die nur f"ur $p$~Punkte der Fl"ache~$T$ unendlich
von der ersten Ordnung wird, wenn man zwei allenthalben endliche
Integrale rationaler Functionen von~$s$ und~$z$ in der Form
\[
\frac{\varphi^{(1)}}{\varphi^{(2)}} \quad \text{oder} \quad \frac{dw^{(1)}}{dw^{(2)}}
\]
im Verh"altniss setzt. Es l"asst sich daher jede Function, die nur f"ur
weniger als $2p + 1$ Punkte der Fl"ache~$T$ unendlich von der ersten
Ordnung wird, als Quotient zweier solcher Integrale darstellen.

\medskip
\textbf{\S\,16.}

Dr"uckt man~$\xi$ als Quotienten zweier ganzen Functionen von~$s$
und~$z$, $\eta$ und~$\zeta$ aus, so sind die Gr"ossenpaare $(s_1,z_1)$, \ldots,
$(s_m,z_m)$ die gemeinschaftlichen Wurzeln der Gleichungen $F = 0$ und
$\xi = \eta/\zeta$. Da die ganze Function~$\eta$ f"ur alle Werthenpaare, f"ur
welche~$\xi$ und~$\psi$ gleichzeitig verschwinden ebenfalls, was auch~$\psi$ sei,
verschwindet, so k"onnen die Gr"ossenpaare $(s_1,z_1)$, \ldots, $(s_m,z_m)$
auch \emph{definirt} werden als gemeinschaftliche Wurzeln der Gleichung $F = 0$
und einer Gleichung $f(s,z) = 0$, deren Coefficienten so sich "andern,
da\ss{} alle "ubrigen gemeinschaftlichen Wurzeln constant bleiben. Wenn
$m < p + 1$, kann~$\xi$ in der Form~$\dfrac{\varphi^{(1)}}{\varphi^{(2)}}$ dargestellt werden
(\S\,10) und~$f$ in der Form
\[
f(s,z) = \varphi^{(1)} - \xi\, \varphi^{(2)} = 0.
\]

Die allgemeinsten Werthe der den $p$~Gleichungen
\[
\sum_{\mu=1}^{m} d\omega_\mu^{(s_\mu,z_\mu)} = 0
\quad (\mu = 1, 2, \ldots, p)
\]
gen"ugenden Functionenpaare $(s_1,z_1)$, \ldots, $(s_p,z_p)$ werden daher
gebildet durch $p$ gemeinschaftliche Wurzeln der Gleichungen $F = 0$ und
$\xi = \eta/\zeta$, welche so sich "andern, da\ss{} die "ubrigen gemeinschaftlichen
Wurzeln constant bleiben. Hieraus folgt leicht der sp"ater n"othige Satz,
da\ss{} die Aufgabe, $p - 1$ von den $2p - 2$ Gr"ossenpaaren $(s_1,z_1)$,
\ldots, $(s_{2p-2}, z_{2p-2})$ als Functionen der $p - 1$ "ubrigen so zu bestimmen,
da\ss{} die $p$~Gleichungen
\[
\sum_{\mu=1}^{2p-2} d\omega_\mu^{(s_\mu,z_\mu)} = 0
\quad (\mu = 1, 2, \ldots, p)
\]
erf"ullt werden, v"ollig allgemein gel"ost wird, wenn man f"ur diese
$2p - 2$ Gr"ossenpaare die von den $r$~Wurzeln $s = \gamma_\mu$, $z = \delta_\mu$
(\S\,6) verschiedenen gemeinschaftlichen Wurzeln der Gleichungen $F = 0$
und $\varphi = 0$ oder die $2p - 2$ Werthenpaare nimmt, f"ur welche $d\omega$
unendlich klein von der zweiten Ordnung wird, und da\ss{} diese Aufgabe
daher nur eine L"osung zul"asst. Solche Gr"ossenpaare sollen durch die
Gleichung $\varphi = 0$ \emph{verkn"upft} heissen.

Infolge der Gleichungen
\[
\sum_{\mu=1}^{2p-2} d\omega_\mu^{(s_\mu,z_\mu)} = 0
\quad\text{wird}\quad
\sum_{\mu=1}^{2p-2} \omega_\mu^{(s_\mu,z_\mu)} \equiv (c_1, c_2, \ldots, c_p),
\]
die Summe "uber solche Gr"ossenpaare ausgedehnt, congruent einem
constanten Gr"ossensysteme $(c_1, c_2, \ldots, c_p)$, worin~$c_\mu$ nur von der
additiven Constanten in der Function~$\omega_\mu$ oder dem Anfangswerthe des
sie ausdr"uckenden Integrals abh"angt.

\newpage
\subsection*{Zweite Abtheilung.}
\addcontentsline{toc}{subsection}{Zweite Abtheilung}

\medskip
\textbf{\S\,17.}

F"ur die ferneren Untersuchungen "uber Integrale von algebraischen,
$(2p+1)$fach zusammenh"angenden Functionen ist die Betrachtung einer
$p$fach unendlichen $\vartheta$-Reihe von grossem Nutzen, d.\,h.\ einer $p$fach
unendlichen Reihe, in welcher der Logarithmus des allgemeinen Gliedes
eine ganze Function zweiten Grades der Stellenzeiger ist. Es sei in
dieser Function f"ur ein Glied, dessen Stellenzeiger $m_1$, $m_2$, \ldots, $m_p$
sind, der Coefficient des Quadrats~$m_\mu^2$ gleich~$a_{\mu,\mu}$, des doppelten
Products $m_\mu m_{\mu'}$ gleich $a_{\mu,\mu'} = a_{\mu',\mu}$, der doppelten Gr"osse~$m_\mu$
gleich~$v_\mu$, und das constante Glied~$= 0$. Die Summe der Reihe, "uber
alle ganze positive oder negative Werthe der Gr"ossen~$m$ ausgedehnt,
werde als Function der $p$~Gr"ossen~$v$ betrachtet und durch
$\vartheta(v_1, v_2, \ldots, v_p)$ bezeichnet, so da\ss{}
\begin{equation}
\label{eq:theta}
\vartheta(v_1, v_2, \ldots, v_p) =
\sum_{-\infty}^{+\infty} e^{\Bigl(\sum_\mu \sum_{\mu'}
  a_{\mu,\mu'} m_\mu m_{\mu'} + 2\sum_\mu v_\mu m_\mu\Bigr)},
\end{equation}
worin die Summationen im Exponenten sich auf~$\mu$ und~$\mu'$, die "ausseren
Summationen auf $m_1$, $m_2$, \ldots, $m_p$ beziehen.

Damit diese Reihe convergirt, mu\ss{} der reelle Theil von
$\sum_\mu \sum_{\mu'} a_{\mu,\mu'} m_\mu m_{\mu'}$ wesentlich negativ sein oder, als eine
Summe von positiven oder negativen Quadraten reeller linearer von einander
unabh"angiger Functionen der Gr"ossen~$m$ dargestellt, aus $p$ negativen
Quadraten zusammengesetzt sein.

Die Function~$\vartheta$ hat die Eigenschaft, da\ss{} es Systeme von
gleichzeitigen "Aenderungen der $p$~Gr"ossen~$v$ giebt, durch welche
$\log\vartheta$ nur um eine lineare Function der Gr"ossen~$v$ ge"andert wird,
und zwar $2p$ von einander unabh"angige Systeme (d.\,h.\ von denen
keins eine Folge der "ubrigen ist). Denn man hat, die unge"andert
bleibenden Gr"ossen~$v$ unter dem Functionszeichen~$\vartheta$ weglassend,
f"ur $\mu = 1, 2, \ldots, p$
\begin{align}
\vartheta &= e^{a_{\mu,\mu} + 2v_\mu}\,
  \vartheta(v_1, \ldots, v_\mu + \pi i, \ldots, v_p),
  \label{eq:theta2}\\
\vartheta &= e^{a_{\mu,\nu}m_\nu + v_\mu}\,
  \vartheta(v_1 + a_{1,\nu}, v_2 + a_{2,\nu}, \ldots, v_p + a_{p,\nu}),
  \label{eq:theta3}
\end{align}
wie sich sofort ergiebt, wenn man in der Reihe f"ur~$\vartheta$ den
Stellenzeiger~$m_\mu$ in $m_\mu + 1$ verwandelt, wodurch sie, w"ahrend ihr
Werth unge"andert bleibt, in den Ausdruck zur Rechten "ubergeht.

Die Function~$\vartheta$ ist durch diese Relationen und durch die Eigenschaft,
allenthalben endlich zu bleiben, bis auf einen constanten Factor bestimmt.
Denn in Folge der letzteren Eigenschaft und der Relationen~(\ref{eq:theta2})
ist sie eine einwerthige, f"ur endliche~$v$ endliche Function von
$e^{2v_1}, e^{2v_2}, \ldots, e^{2v_p}$ und folglich in eine $p$fach unendliche Reihe
von der Form
\[
\sum_{-\infty}^{+\infty} A_{m_1,\ldots,m_p}\,
  e^{2(m_1 v_1 + m_2 v_2 + \cdots + m_p v_p)}
\]
mit den constanten Coefficienten~$A$ entwickelbar. Aus den Relationen
(\ref{eq:theta3}) ergiebt sich aber
\[
A_{m_1,\ldots,m_\nu+1,\ldots,m_p} =
A_{m_1,\ldots,m_\nu,\ldots,m_p}\,
  e^{\sum_{\mu'} a_{\nu,\mu'} m_{\mu'}},
\]
folglich $A_m = \text{const.} \cdot e^{\sum \sum a_{\mu,\mu'} m_\mu m_{\mu'}}$, w.\,z.\,b.\,w.

Man kann daher diese Eigenschaften der Function zu ihrer Definition
verwenden. Die Systeme gleichzeitiger "Aenderungen der Gr"ossen~$v$,
durch welche sich $\log\vartheta$ nur um eine lineare Function von ihnen "andert,
sollen \emph{Systeme zusammengeh"origer Periodicit"atsmoduln} der unabh"angig
ver"anderlichen Gr"ossen in dieser $\vartheta$-Function genannt werden.

\medskip
\textbf{\S\,18.}

Ich substituire nun f"ur die $p$~Gr"ossen $v_1, v_2, \ldots, v_p$ $p$~immer
endlich bleibende Integrale $u_1, u_2, \ldots, u_p$ rationaler Functionen einer
ver"anderlichen Gr"osse~$z$ und einer $(2p+1)$fach zusammenh"angenden
algebraischen Function~$s$ dieser Gr"osse, und f"ur die zusammengeh"origen
Periodicit"atsmoduln der Gr"ossen~$v$ zusammengeh"orige (d.\,h.\ an demselben
Querschnitte stattfindende) Periodicit"atsmoduln dieser Integrale, so da\ss{}
$\log\vartheta$ in eine Function einer Ver"anderlichen~$z$ "ubergeht, welche sich,
wenn~$s$ und~$z$ nach beliebiger stetiger "Aenderung von~$z$ den vorigen
Werth wieder annehmen, um lineare Functionen der Gr"ossen~$u$ "andert.

Es soll zun"achst gezeigt werden, da\ss{} eine solche Substitution f"ur jede
$(2p+1)$fach zusammenh"angende Function~$s$ m"oglich ist. Die Zerschneidung
der Fl"ache~$T$ mu\ss{} zu diesem Zwecke so durch $2p$ in sich zur"ucklaufende
Schnitte $a_1, a_2, \ldots, a_p, b_1, b_2, \ldots, b_p$ geschehen, da\ss{} folgende
Bedingungen erf"ullt werden. Wenn man $u_1, u_2, \ldots, u_p$ so w"ahlt, da\ss{}
der Periodicit"atsmodul von~$u_\mu$ an dem Schnitte~$a_\nu$ gleich~$\pi i$ (f"ur
$\nu = \mu$), an den "ubrigen Schnitten~$a$ gleich~$0$ ist, und man den
Periodicit"atsmodul von~$u_\mu$ an dem Schnitte~$b_\nu$ durch $a_{\mu,\nu}$
bezeichnet, so mu\ss{} $a_{\mu,\nu} = a_{\nu,\mu}$ und der reelle Theil von
$\sum_\mu \sum_{\mu'} a_{\mu,\mu'} m_\mu m_{\mu'}$ f"ur alle reellen (ganzen) Werthe der
$p$~Gr"ossen~$m$ negativ sein.

\medskip
\textbf{\S\,19.}

Die Zerlegung der Fl"ache~$T$ werde nicht wie bisher nur durch in
sich zur"ucklaufende Querschnitte, sondern folgenderma\ss{}en ausgef"uhrt.
Man mache zuerst einen in sich zur"ucklaufenden die Fl"ache nicht
zerst"uckelnden Schnitt~$a_1$ und f"uhre dann einen Querschnitt~$b_1$ von der
positiven Seite von~$a_1$ auf die negative zum Anfangspunkte zur"uck,
worauf die Begrenzung aus einem St"ucke bestehen wird. Einen dritten
die Fl"ache nicht zerst"uckelnden Querschnitt kann man demzufolge (wenn
die Fl"ache noch nicht einfach zusammenh"angend ist) von einem beliebigen
Punkte dieser Begrenzung bis zu einem beliebigen Begrenzungspunkte,
also auch zu einem fr"uheren Punkte dieses Querschnitts f"uhren. Man
thue das Letztere, so da\ss{} dieser Querschnitt aus einer in sich
zur"ucklaufenden Linie~$a_2$ und einem dieser Linie voraufgehenden Theile~$c_1$
besteht, welcher das fr"uhere Schnittsystem mit ihr verbindet. Den
folgenden Querschnitt~$b_2$ ziehe man von der positiven Seite von~$a_2$ auf
die negative zum Anfangspunkte zur"uck, worauf die Begrenzung wieder
aus einem St"ucke besteht. Die weitere Zerschneidung kann daher, wenn
n"othig, wieder durch zwei in demselben Punkte anfangende und endende
Schnitte $a_3$ und~$b_3$ und eine das System der Linien $a_2$ und~$b_2$ mit
ihnen verbindende Linie~$c_2$ geschehen.

Wird dieses Verfahren fortgesetzt, bis die Fl"ache einfach zusammenh"angend
ist, so erh"alt man ein Schnittnetz, welches aus $p$~Paaren von zwei in
einem und demselben Punkte anfangenden und endenden Linien
$a_1$ und~$b_1$, $a_2$ und~$b_2$, \ldots, $a_p$ und~$b_p$ besteht und aus $p - 1$
Linien $c_1, c_2, \ldots, c_{p-1}$, welche jedes Paar mit dem folgenden verbinden.
Es m"oge $c_\nu$ von einem Punkte von~$b_\nu$ nach einem Punkte von~$a_{\nu+1}$
gehen.

Das Schnittnetz wird als so entstanden betrachtet, da\ss{} der
$(2\nu - 1)$te Querschnitt aus $c_{\nu-1}$ und der von dem Endpunkte von
$c_{\nu-1}$ zu diesem zur"uckgezogenen Linie~$a_\nu$ besteht, und der $2\nu$te
durch die von der positiven auf die negative Seite von~$a_\nu$ gezogene
Linie~$b_\nu$ gebildet wird. Die Begrenzung der Fl"ache besteht bei dieser
Zerschneidung nach einer geraden Anzahl von Schnitten aus einem, nach
einer ungeraden aus zwei St"ucken.

Ein allenthalben endliches Integral~$\omega$ einer rationalen Function von~$s$
und~$z$ nimmt dann zu beiden Seiten einer Linie~$c$ denselben Werth an.
Denn die ganze fr"uher entstandene Begrenzung besteht aus einem St"ucke
und bei der Integration l"angs derselben von der einen Seite der Linie~$c$
bis auf die andere wird $\int d\omega$ durch jedes fr"uher entstandene
Schnittelement zweimal, in entgegengesetzter Richtung, erstreckt.
Eine solche Function ist daher in~$T$ allenthalben ausser den Linien~$a$
und~$b$ stetig. Die durch diese Linien zerschnittene Fl"ache~$T$ m"oge
durch~$T''$ bezeichnet werden.

\medskip
\textbf{\S\,20.}

Es seien nun $\omega_1, \omega_2, \ldots, \omega_p$ von einander unabh"angige solche
Functionen, und der Periodicit"atsmodul von~$\omega_\mu$ an dem Querschnitte~$a_\nu$
gleich $A_\mu^{(\nu)}$ und an dem Querschnitte~$b_\nu$ gleich~$B_\mu^{(\nu)}$.
Es ist dann das Integral $\int \omega_\mu\, d\omega_\mu'$, um die Fl"ache~$T''$
positiv herum ausgedehnt, $= 0$, da die Function unter dem Integralzeichen
allenthalben endlich ist. Bei dieser Integration wird jede der Linien~$a$
und~$b$ zweimal, einmal in positiver und einmal in negativer Richtung
durchlaufen, und es mu\ss{} w"ahrend jener Integration, wo sie als Begrenzung
des positiverseits gelegenen Gebiets dient, f"ur $\omega_\mu$ der Werth auf der
positiven Seite oder~$\omega_\mu^+$, w"ahrend dieser der Werth auf der negativen
oder~$\omega_\mu^-$ genommen werden.

Es ist also dies Integral gleich der Summe
\[
\sum \int (\omega_\mu^+ - \omega_\mu^-)\, d\omega_\mu'
\]
aller Integrale $\int (\omega_\mu^+ - \omega_\mu^-)\, d\omega_\mu'$ durch die Linien~$a$
und~$b$. Die Linien~$b$ f"uhren von der positiven zur negativen Seite
der Linien~$a$, und folglich die Linien~$a$ von der negativen zur positiven
Seite der Linien~$b$. Das Integral durch die Linie~$a_\nu$ ist daher
\[
\int A_\mu^{(\nu)}\, d\omega_\mu' = A_\mu^{(\nu)} B_{\mu'}^{(\nu)},
\]
und das Integral durch die Linie~$b_\nu$
\[
\int B_\mu^{(\nu)}\, d\omega_\mu' = -B_\mu^{(\nu)} A_{\mu'}^{(\nu)}.
\]
Das Integral $\int \omega_\mu\, d\omega_\mu'$, um die Fl"ache~$T''$ positiv herum
erstreckt, ist also
\[
\sum_{\nu=1}^{p} \bigl(A_\mu^{(\nu)} B_{\mu'}^{(\nu)}
  - B_\mu^{(\nu)} A_{\mu'}^{(\nu)}\bigr) = 0,
\]
und diese Summe folglich $= 0$. Diese Gleichung gilt f"ur je zwei von
den Functionen $\omega_1, \omega_2, \ldots, \omega_p$ und liefert also
$\dfrac{p(p-1)}{2}$ Relationen zwischen deren Periodicit"atsmoduln.

Nimmt man f"ur die Functionen~$\omega$ die Functionen~$u$ oder w"ahlt
man sie so, da\ss{} $A_\mu^{(\nu)}$ f"ur ein von~$\mu$ verschiedenes~$\nu$ gleich~$0$ und
$A_\mu^{(\mu)} = \pi i$ ist, so gehen diese Relationen "uber in
$B_{\mu'}^{(\mu)} \pi i - B_\mu^{(\mu')} \pi i = 0$ oder in $a_{\mu',\mu} = a_{\mu,\mu'}$.

\medskip
\textbf{\S\,21.}

Es bleibt noch zu zeigen, da\ss{} die Gr"ossen~$a$ die zweite oben
n"othig gefundene Eigenschaft besitzen.

Man setze $\omega = u + vi$ und den Periodicit"atsmodul dieser Function
an dem Schnitte~$a_\nu$ gleich $A^{(\nu)} = \alpha_\nu + \gamma_\nu i$ und an dem
Schnitte~$b_\nu$ gleich $B^{(\nu)} = \beta_\nu + \delta_\nu i$.
Es ist dann das Integral
\[
\int\Bigl[\Bigl(\frac{\partial u}{\partial x}\Bigr)^2
  + \Bigl(\frac{\partial u}{\partial y}\Bigr)^2\Bigr]\, dx\,dy
\quad\text{oder}\quad
\int\Bigl(\frac{\partial u}{\partial x}\frac{\partial v}{\partial y}
  - \frac{\partial u}{\partial y}\frac{\partial v}{\partial x}\Bigr)\, dx\,dy,
\]
durch die Fl"ache~$T''$ gleich dem Begrenzungsintegral $\int u\,dv$
um~$T''$ positiv herum erstreckt, also gleich der Summe der Integrale
$\int (u^+ - u^-)\, dv$ durch die Linien~$a$ und~$b$. Das Integral durch
die Linie~$a_\nu$ ist $= \alpha_\nu \int dv = \alpha_\nu \gamma_\nu$, das Integral durch
die Linie~$b_\nu$ gleich $\beta_\nu \int dv = -\beta_\nu \delta_\nu$.
und folglich
\[
\sum_{\nu=1}^{p} (\alpha_\nu \gamma_\nu - \beta_\nu \delta_\nu)
  = \int\Bigl[\Bigl(\frac{\partial u}{\partial x}\Bigr)^2
    + \Bigl(\frac{\partial u}{\partial y}\Bigr)^2\Bigr]\, dx\,dy.
\]
Diese Summe ist daher stets \emph{positiv}.

Hieraus ergiebt sich die zu beweisende Eigenschaft der Gr"ossen~$a$,
wenn man f"ur~$\omega$ setzt $u_1 m_1 + u_2 m_2 + \cdots + u_p m_p$. Denn es ist
dann $A^{(\nu)} = m_\nu \pi i$, $B^{(\nu)} = \sum_\mu a_{\nu,\mu} m_\mu$, folglich
$\alpha_\nu$ stets $= 0$ und
\[
\int\Bigl[\Bigl(\frac{\partial u}{\partial x}\Bigr)^2
  + \Bigl(\frac{\partial u}{\partial y}\Bigr)^2\Bigr]\, dx\,dy
= -\sum_\nu \sum_\mu m_\nu a_{\nu,\mu} m_\mu
\]
oder gleich dem reellen Theile von $-\sum_\nu \sum_\mu \pi i\, a_{\nu,\mu} m_\nu m_\mu$,
welcher also f"ur alle reellen Werthe der Gr"ossen~$m$ positiv ist.

\medskip
\textbf{\S\,22.}

Setzt man nun in der $\vartheta$-Reihe~(\ref{eq:theta}) \S\,17 f"ur
$a_{\mu,\mu'}$ den Periodicit"atsmodul der Function~$u_\mu$ an dem Schnitte~$b_{\mu'}$
und, durch $e_1, e_2, \ldots, e_p$ beliebige Constanten bezeichnend,
$u_\mu - e_\mu$ f"ur~$v_\mu$, so erh"alt man eine in jedem Punkte von~$T$
eindeutig bestimmte Function von~$z$, welche ausser den Linien~$b$ stetig
und endlich und auf der positiven Seite der Linie~$b_\nu$
$e^{2u_\nu - 2e_\nu - a_{\nu,\nu}}$ mal so gro\ss{} als auf der negativen ist, wenn man
den Functionen~$u$ in den Linien~$b$ selbst den Mittelwerth von den
Werthen zu beiden Seiten beilegt.

F"ur wie viele Punkte von~$T'$ oder Werthenpaare von~$s$ und~$z$
diese Function unendlich klein von der ersten Ordnung wird, kann durch
Betrachtung des Begrenzungsintegrals $\int d\log\vartheta$, um~$T'$ positiv herum
erstreckt, gefunden werden; denn dieses Integral ist gleich der Anzahl
dieser Punkte multiplicirt mit $2\pi i$. Andererseits ist dies Integral
gleich der Summe der Integrale $\int (d\log\vartheta^+ - d\log\vartheta^-)$ durch
s"ammtliche Schnittlinien $a, b$ und~$c$. Die Integrale durch die Linien~$a$
und~$c$ sind~$= 0$, das Integral durch~$b_\nu$ aber gleich $-2 \int du_\nu = 2\pi i$,
die Summe aller also $= p \cdot 2\pi i$. Die Function~$\vartheta$ wird daher
unendlich klein von der ersten Ordnung in $p$~Punkten der Fl"ache~$T'$,
welche durch $\eta_1, \eta_2, \ldots, \eta_p$ bezeichnet werden m"ogen.

\medskip
\textbf{\S\,23.}

Durch einen positiven Umlauf des Punktes $(s,z)$ um einen dieser
Punkte w"achst $\log\vartheta$ um $2\pi i$, durch einen positiven Umlauf um das
Schnittepaar $a_\nu$ und~$b_\nu$ um $-2\pi i$. Um daher die Function
$\log\vartheta$ allenthalben eindeutig zu bestimmen, f"uhre man von jedem
Punkte~$\eta_\nu$ einen Schnitt~$l_\nu$ durch das Innere nach je einem
Linienpaare, von~$\eta_\nu$ den Schnitt~$l_\nu$ nach $a_\nu$ und~$b_\nu$, und
zwar nach ihrem gemeinschaftlichen Anfangs- und Endpunkte, und nehme
in der dadurch entstandenen Fl"ache~$T^*$ die Function allenthalben stetig
an. Sie ist dann auf der positiven Seite der Linien~$l$ um $-2\pi i$,
auf der positiven Seite der Linie~$a_\nu$ um $g_\nu \cdot 2\pi i$ und auf der
positiven Seite der Linie~$b_\nu$ um $-2(\alpha_\nu - e_\nu) - h_\nu \cdot 2\pi i$
gr"osser, als auf der negativen, wenn $g_\nu$ und~$h_\nu$ ganze Zahlen bezeichnen.

Die Lage der Punkte~$\eta$ und die Werthe der Zahlen~$g$ und~$h$
h"angen von den Gr"ossen~$e$ ab, und diese Abh"angigkeit l"asst sich auf
folgendem Wege n"aher bestimmen. Das Integral $\int \log\vartheta\, du_\mu$,
um~$T^*$ positiv herum erstreckt, ist~$= 0$, da die Function~$\log\vartheta$
in~$T^*$ stetig bleibt. Dieses Integral ist aber auch gleich der Summe
der Integrale $\int (\log\vartheta^+ - \log\vartheta^-)\, du_\mu$ durch s"ammtliche
Schnittlinien $l, a, h$ und~$c$ und findet sich, wenn man den Werth von~$u_\mu$
im Punkte~$\eta_\nu$ durch~$\alpha_\mu^{(\nu)}$ bezeichnet,
\[
2\pi i \Bigl(\sum_\nu \alpha_\mu^{(\nu)} + h_\mu \pi i + \sum_\nu g_\nu a_{\nu,\mu}
  - e_\mu + k_\mu\Bigr) = 0,
\]
worin $k_\mu$ von den Gr"ossen $e, g, h$ und der Lage der Punkte~$\eta$
unabh"angig ist.

Dieser Ausdruck ist also~$= 0$.

\bigskip
\begin{center}
\rule{0.5\textwidth}{0.4pt}
\end{center}

\vfill
\begin{center}
\textit{Ende der Theorie der Abel'schen Functionen f"ur die hier behandelten
Paragraphen.}
\end{center}

\clearpage
\chapter*{Source segment 4: riemann pages 151 200}
\addcontentsline{toc}{chapter}{Source segment 4: riemann pages 151 200}
\selectlanguage{german}

\begin{center}
\Large\textbf{Riemann --- Gesammelte Mathematische Werke}\\[0.5em]
\large\textbf{Seiten 151--200}
\end{center}

\vspace{1em}

%===================================================================
% ENDE VON PAPER VII: Ueber die Anzahl der Primzahlen
%===================================================================

\noindent\textbf{VII. Ueber die Anzahl der Primzahlen unter einer gegebenen Gr\"osse.}

\bigskip

\noindent oder
\[
\int\limits_{\beta-\infty i}^{x}\frac{dx}{\log x} + \text{const.}
\]
je nachdem der reelle Theil von $\beta$ negativ oder positiv ist. Man hat daher
\[
2\pi i\log\eta = \int\limits_{\alpha-\infty i}^{\alpha+\infty i} \frac{d}{ds}\log\xi(s)\,x^{s}\,\frac{ds}{s}
\]
und
\[
\int\limits_{\beta}^{x}\frac{dx}{\log x} + \text{const. im ersten}
\]
und
\[
\int\limits_{0}^{x}\frac{dx}{\log x} + \text{const. im zweiten Falle.}
\]

Im ersten Falle bestimmt sich die Integrationsconstante, wenn man den reellen Theil von $\beta$ negativ unendlich werden l\"asst; im zweiten Falle erh\"alt das Integral von $0$ bis $x$ um $2\pi i$ verschiedene Werthe, je nachdem die Integration durch complexe Werthe mit positivem oder negativem Arcus geschieht, und wird, auf jenem Wege genommen, unendlich klein, wenn der Coefficient von $i$ in dem Werthe von $\beta$ positiv unendlich wird, auf letzterem aber, wenn dieser Coefficient negativ unendlich wird. Hieraus ergiebt sich, wie auf der linken Seite $\log\bigl(1-\frac{s}{\beta}\bigr)$ zu bestimmen ist, damit die Integrationsconstante wegf\"allt.

Durch Einsetzung dieser Werthe in den Ausdruck von $f(x)$ erh\"alt man
\[
f(x) = \Li(x) - \sum_{\alpha}\bigl(\Li(x^{\frac{1}{2}+\alpha i}) + \Li(x^{\frac{1}{2}-\alpha i})\bigr) + \int\limits_{x^{2}}^{\infty}\frac{dt}{t(t^{2}-1)\log t} + \log\xi(0),
\]
 wenn in $\sum_{\alpha}$ f\"ur $\alpha$ s\"ammtliche positive (oder einen positiven reellen Theil enthaltenden) Wurzeln der Gleichung $\xi(\alpha) = 0$, ihrer Gr\"osse nach geordnet, gesetzt werden. Es l\"asst sich, mit H\"ulfe einer genaueren Discussion der Function $\xi$, leicht zeigen, dass bei dieser Anordnung der Werth der Reihe

% --- page 152 ---
\[
\sum_{\alpha}\bigl(\Li(x^{\frac{1}{2}+\alpha i}) + \Li(x^{\frac{1}{2}-\alpha i})\bigr)\log x
\]
mit dem Grenzwerth, gegen welchen
\[
\frac{1}{2\pi i\log x} \int\limits_{\alpha-bi}^{\alpha+bi} \frac{d\frac{1}{s}\log\xi(s)}{ds}\,x^{s}\,ds
\]
bei unaufh\"orlichem Wachsen der Gr\"osse $b$ convergirt, \"ubereinstimmt; durch ver\"anderte Anordnung aber w\"urde sie jeden beliebigen reellen Werth erhalten k\"onnen.

Aus $f(x)$ findet sich $F(x)$ mittelst der durch Umkehrung der Relation
\[
f(x) = \sum\frac{1}{m}F(x^{\frac{1}{m}})
\]
sich ergebenden Gleichung
\[
F(x) = \sum(-1)^{\mu}\frac{1}{m}f(x^{\frac{1}{m}}),
\]
worin f\"ur $m$ der Reihe nach die durch kein Quadrat ausser $1$ theilbaren Zahlen zu setzen sind und $\mu$ die Anzahl der Primfactoren von $m$ bezeichnet.

Beschr\"ankt man $\Li$ auf eine endliche Zahl von Gliedern, so giebt die Derivirte des Ausdrucks f\"ur $f(x)$ oder, bis auf einen mit wachsendem $x$ sehr schnell abnehmenden Theil,
\[
\frac{1}{\log x} - 2\sum_{\alpha}\frac{\cos(\alpha\log x)x^{-\frac{1}{2}}}{\log x}
\]
einen angen\"aherten Ausdruck f\"ur die Dichtigkeit der Primzahlen $+$ der halben Dichtigkeit der Primzahlquadrate $+\frac{1}{3}$ von der Dichtigkeit der Primzahlcuben u.~s.~w.~von der Gr\"osse $x$.

Die bekannte N\"aherungsformel $F(x) = \Li(x)$ ist also nur bis auf Gr\"ossen von der Ordnung $x^{\frac{1}{2}}$ richtig und giebt einen etwas zu grossen Werth; denn die nicht periodischen Glieder in dem Ausdrucke von $F(x)$ sind, von Gr\"ossen, die mit $x$ nicht in's Unendliche wachsen, abgesehen:
\[
\Li(x) - \tfrac{1}{2}\Li(x^{\frac{1}{2}}) - \tfrac{1}{3}\Li(x^{\frac{1}{3}}) - \tfrac{1}{5}\Li(x^{\frac{1}{5}}) + \tfrac{1}{6}\Li(x^{\frac{1}{6}}) - \tfrac{1}{7}\Li(x^{\frac{1}{7}}) + \cdots
\]

In der That hat sich bei der von Gauss und Goldschmidt vorgenommenen und bis zu $x =$ drei Millionen fortgesetzten Vergleichung von $\Li(x)$ mit der Anzahl der Primzahlen unter $x$ diese Anzahl schon vom ersten Hunderttausend an stets kleiner als $\Li(x)$ ergeben, und

% --- page 153 ---
zwar w\"achst die Differenz unter manchen Schwankungen allm\"ahlich mit $x$. Aber auch die von den periodischen Gliedern abh\"angige stellenweise Verdichtung und Verd\"unnung der Primzahlen hat schon bei den Z\"ahlungen die Aufmerksamkeit erregt, ohne dass jedoch hierin eine Gesetzm\"assigkeit bemerkt worden w\"are. Bei einer etwaigen neuen Z\"ahlung w\"urde es interessant sein, den Einfluss der einzelnen in dem Ausdrucke f\"ur die Dichtigkeit der Primzahlen enthaltenen periodischen Glieder zu verfolgen. Einen regelm\"assigeren Gang als $F(x)$ w\"urde die Function $f(x)$ zeigen, welche sich schon im ersten Hundert sehr deutlich als mit $\Li(x) + \log\xi(0)$ im Mittel \"ubereinstimmend erkennen l\"asst.

\bigskip
% --- page 154 ---
\noindent\textbf{Anmerkungen.}

\smallskip
\noindent In einem Briefe, dessen Entwurf im Nachlass vorliegt, findet sich, nachdem das Resultat der Arbeit mitgetheilt ist, folgende Bemerkung:

\smallskip
\noindent\textit{,,Den Beweis habe ich noch nicht v\"ollig ausgef\"uhrt, und ich m\"ochte in Betreff desselben \dots\ noch die Bemerkung beif\"ugen, dass die beiden S\"atze, welche ich dort nur angef\"uhrt habe,
\[
\text{dass zwischen }0\text{ und }T\text{ etwa }\frac{T}{2\pi}\log\frac{T}{2\pi}-\frac{T}{2\pi}\text{ reelle Wurzeln der Gleichung }\xi(t) = 0\text{ liegen,}
\]
und
\[
\text{dass die Reihe }\sum_{\alpha}\bigl(\Li(x^{\frac{1}{2}+\alpha i}) + \Li(x^{\frac{1}{2}-\alpha i})\bigr),\text{ wenn die Glieder nach wachsenden }\alpha\text{ geordnet werden,}
\]
gegen denselben Grenzwerth convergirt, wie
\[
\frac{1}{2\pi i\log x}\int\limits_{a-bi}^{a+bi}\frac{d\frac{1}{s}\log\xi(s)}{ds}\,x^{s}\,ds
\]
bei unaufh\"orlichem Wachsen der Gr\"osse $b$,
\ aus einer neuen Entwicklung der Function $\xi$ folgen, welche ich aber noch nicht genug vereinfacht hatte, um sie mittheilen zu k\"onnen.``}

\smallskip
\noindent Trotz mancher sp\"aterer Untersuchungen (Scheibner, Pilz, Stieltjes) sind die Dunkelheiten dieser Arbeit noch nicht v\"ollig aufgehellt.

\smallskip
\noindent (1) (Zu Seite~146.) Dies Verhalten der Function $J(s)$ ergiebt sich mit Benutzung der zweiten Form dieser Function
\[
J(s) = \frac{\Gamma(\tfrac{1}{2}s)}{\Gamma(\tfrac{1}{2})}\pi^{-\frac{1}{2}s}\sum_{n=1}^{\infty}\frac{1}{n^{s}} - \frac{1}{s(1-s)}
\]
und mit R\"ucksicht darauf, dass $\frac{1}{e^{x}-1}$ in der Entwicklung nach steigenden Potenzen von $x$ nur ungerade Potenzen enth\"alt.

\smallskip
\noindent (2) (Zu Seite~149.) Der Ausdruck dieses Satzes ist nicht ganz genau. Die beiden Gleichungen, einzeln in der angegebenen Weise behandelt, ergeben, wenn die Integrationsgrenzen $0$, $\infty$ auf $\log x$ bezogen werden,
\[
x^{\frac{1}{2}+\alpha i}\bigl(\vartheta(\alpha) + \vartheta(-\alpha)\bigr)
\]
und also erst in ihrer Summe die Formel des Textes.

\smallskip
\noindent (3) (Zu Seite~151.) Die Function $\Li(x)$ ist f\"ur reelle Werthe von $x$, die gr\"osser als $1$ sind,
\[
\int\limits_{0}^{x}\frac{dx}{\log x}
\]
 wenn das obere oder
\[
\int\limits_{-\infty}^{x}\frac{dx}{\log x}
\]
 das untere Zeichen zu nehmen ist, je nachdem die Integration durch complexe

% --- page 155 ---
Werthe mit positivem oder negativem Arcus geschieht. Daraus leitet man leicht die von Scheibner (Schl\"omilch's Zeitschrift Bd.~V) gegebene Entwicklung her
\[
\Li(x) = \log\log x - \Gamma'(1) + \sum_{n=1}^{\infty}\frac{(\log x)^{n}}{n\cdot n!},
\]
die f\"ur alle Werthe von $x$ gilt, und f\"ur negative reelle Werthe eine Unstetigkeit ergiebt. (Vgl. Gauss--Bessel Briefwechsel.)

Befolgt man die von Riemann angedeutete Rechnung, so findet man in der Formel $\log\xi(0)$ anstatt $\log\xi(0)$. M\"oglicherweise liegt nur ein Schreib- oder Druckfehler vor, $\log\xi(0)$ an Stelle von $\log\xi(0)$, da $\xi(0) = -\frac{1}{2}$ ist.


%===================================================================
% PAPER VIII: Ueber die Fortpflanzung ebener Luftwellen
%===================================================================

\newpage
\begin{center}
\Large\textbf{VIII.}
\end{center}
\vspace{-0.5em}
\begin{center}
\Large\textbf{Ueber die Fortpflanzung ebener Luftwellen}\\[0.3em]
\Large\textbf{von endlicher Schwingungsweite.}
\end{center}
\vspace{0.3em}
\begin{center}
\small (Aus dem achten Bande der Abhandlungen der K\"oniglichen Gesellschaft der\\
\small Wissenschaften zu G\"ottingen. 1860.)
\end{center}

\bigskip

Obwohl die Differentialgleichungen, nach welchen sich die Bewegung der Gase bestimmt, l\"angst aufgestellt worden sind, so ist doch ihre Integration fast nur f\"ur den Fall ausgef\"uhrt worden, wenn die Druckverschiedenheiten als unendlich kleine Bruchtheile des ganzen Drucks betrachtet werden k\"onnen, und man hat sich bis auf die neueste Zeit begn\"ugt, nur die ersten Potenzen dieser Bruchtheile zu ber\"ucksichtigen. Erst ganz vor Kurzem hat Helmholtz auch die Glieder zweiter Ordnung mit in die Rechnung gezogen und daraus die objective Entstehung von Combinationst\"onen erkl\"art. Es lassen sich indess f\"ur den Fall, dass die anf\"angliche Bewegung allenthalben in gleicher Richtung stattfindet und in jeder auf dieser Richtung senkrechten Ebene Geschwindigkeit und Druck constant sind, die exacten Differentialgleichungen vollst\"andig integriren; und wenn auch zur Erkl\"arung der bis jetzt experimentell festgestellten Erscheinungen die bisherige Behandlung vollkommen ausreicht, so k\"onnten doch, bei den grossen Fortschritten, welche in neuester Zeit durch Helmholtz auch in der experimentellen Behandlung akustischer Fragen gemacht worden sind, die Resultate dieser genaueren Rechnung in nicht allzu ferner Zeit vielleicht der experimentellen Forschung einige Anhaltspunkte gew\"ahren; und dies mag, abgesehen von dem theoretischen Interesse, welches die Behandlung nicht linearer partieller Differentialgleichungen hat, die Mittheilung derselben rechtfertigen.

F\"ur die Abh\"angigkeit des Drucks von der Dichtigkeit w\"urde das Boyle'sche Gesetz vorauszusetzen sein, wenn die durch die Druckver\"anderungen bewirkten Temperaturverschiedenheiten sich so schnell ausglichen, dass die Temperatur des Gases als constant betrachtet werden d\"urfte. Es ist aber wahrscheinlich der W\"armeaustausch ganz

% --- page 157 ---
zu vernachl\"assigen, und man muss daher f\"ur diese Abh\"angigkeit das Gesetz zu Grunde legen, nach welchem sich der Druck des Gases mit der Dichtigkeit \"andert, wenn es keine W\"arme aufnimmt oder abgiebt.

Nach dem Boyle'schen und Gay-Lussac'schen Gesetze ist, wenn $V$ das Volumen der Gewichtseinheit, $p$ den Druck und $T$ die Temperatur von $-273^{\circ}$ C an gerechnet bezeichnet,
\[
\log p + \log V = \log T + \text{const.}
\]
Betrachten wir hier $T$ als Function von $p$ und $v$ und nennen die specifische W\"arme bei constantem Drucke $c$, bei constantem Volumen $c'$, beide auf die Gewichtseinheit bezogen, so wird von dieser Gewichtseinheit, wenn $p$ und $v$ sich um $dp$ und $dv$ \"andern, die W\"armemenge
\[
c\frac{dT}{dv}dv + c'\frac{dT}{dp}dp
\]
oder, da $\frac{d\log T}{d\log v} = \frac{c'}{c-c'}$, $\frac{d\log T}{d\log p} = 1$,
\[
T(c\,d\log v + c'\,d\log p)
\]
aufgenommen. Wenn daher keine W\"armeaufnahme stattfindet, so ist $d\log p = -\frac{c}{c'}d\log v$, und also, wenn man mit Poisson annimmt, dass das Verh\"altniss der beiden specifischen W\"armen $\frac{c}{c'} = k$ von Temperatur und Druck unabh\"angig ist,
\[
\log p = -k\log v + \text{const.}
\]

Nach neueren Versuchen von Regnault, Joule und W.~Thomson sind diese S\"atze f\"ur Sauerstoff, Stickstoff und Wasserstoff und deren Gemenge unter allen darstellbaren Drucken und Temperaturen wahrscheinlich sehr nahe g\"ultig.

Durch Regnault ist f\"ur diese Gase eine sehr nahe Anschmiegung an das Boyle'sche und Gay-Lussac'sche Gesetz und die Unabh\"angigkeit der specifischen W\"arme $c$ von Temperatur und Druck festgestellt worden.

F\"ur atmosph\"arische Luft fand Regnault
\begin{align*}
\text{zwischen } -30^{\circ}\text{C und } +100^{\circ}\text{C} \quad c &= 0{,}2377\\
\text{``}\quad +10^{\circ}\text{C}\quad\text{``}\quad +100^{\circ}\text{C} \quad c &= 0{,}2379\\
\text{``}\quad +100^{\circ}\text{C}\quad\text{``}\quad +215^{\circ}\text{C} \quad c &= 0{,}2376.
\end{align*}
Ebenso ergab sich f\"ur Drucke von $1$ bis $10$ Atmosph\"aren kein merklicher Unterschied der specifischen W\"arme.

Nach Versuchen von Regnault und Joule scheint ferner f\"ur diese Gase die von Clausius adoptirte Annahme Mayer's sehr nahe richtig zu sein, dass ein bei constanter Temperatur sich ausdehnendes Gas nur so viel W\"arme aufnimmt, als zur Erzeugung der \"ausseren Arbeit erforderlich ist. Wenn das Volumen des Gases sich um $dv$ \"andert, w\"ahrend die Temperatur constant bleibt, so ist $d\log p = -d\log v$, die aufgenommene W\"armemenge $T(c-c')d\log v$, die geleistete Arbeit $p\,dv$. Diese Hypothese giebt daher, wenn $A$ das mechanische Aequivalent der W\"arme bezeichnet,
\[
AT(c-c')d\log v = p\,dv
\]
oder
\[
\frac{pv}{AT} = c - c',
\]
also von Druck und Temperatur unabh\"angig.

Hiernach ist auch $k = \frac{c}{c'}$ von Druck und Temperatur unabh\"angig und ergiebt sich, wenn $c = 0{,}237733$, $A$ nach Joule $= 424{,}55$ Kilogrammmet. und, f\"ur die Temperatur $0^{\circ}$C oder $T = \frac{100^{\circ}}{273}$ mit Regnault $= 7990^{m},267$ angenommen wird, gleich $1{,}4101$. Die Schallgeschwindigkeit in trockner Luft von $0^{\circ}$C betr\"agt in der Secunde
\[
\sqrt{7990^{m},267 \cdot 9^{m},8088/k},
\]
und w\"urde also mit diesem Werthe von $k$ gleich $332^{m},440$ gefunden werden, w\"ahrend die beiden vollst\"andigsten Versuchsreihen von Moll und van Beck daf\"ur, einzeln berechnet, $332^{m},528$ und $331^{m},867$, vereinigt $332^{m},271$ geben und die Versuche von Martins und A.~Bravais nach ihrer eignen Berechnung $332^{m},37$.

F\"ur's erste ist es nicht n\"othig, \"uber die Abh\"angigkeit des Drucks von der Dichtigkeit eine bestimmte Voraussetzung zu machen; wir nehmen daher an, dass bei der Dichtigkeit $\varrho$ der Druck $\varphi(\varrho)$ sei, und lassen die Function $\varphi$ vorl\"aufig noch unbestimmt.

Man denke sich nun rechtwinklige Coordinaten $x$, $y$, $z$ eingef\"uhrt, die $x$-Axe in der Richtung der Bewegung, und bezeichne durch $\varrho$ die Dichtigkeit, durch $p$ den Druck, durch $u$ die Geschwindigkeit f\"ur die Coordinate $x$ zur Zeit $t$ und durch $\omega$ ein Element der Ebene, deren Coordinate $x$ ist.

Der Inhalt des auf dem Element $\omega$ stehenden geraden Cylinders von der H\"ohe $dx$ ist dann $\omega\,dx$, die in ihm enthaltene Masse $\varrho\omega\,dx$. Die Aenderung dieser Masse w\"ahrend des Zeitelements $dt$ oder die Gr\"osse $\varrho\omega\,dt\,dx$ bestimmt sich durch die in ihn einstr\"omende Masse, welche $= -\omega\frac{\partial(\varrho u)}{\partial x}dx\,dt$ gefunden wird. Ihre Beschleunigung ist

% --- page 159 ---
$+\varrho\omega\frac{\partial u}{\partial t}$ und die Kraft, welche sie in der Richtung der positiven $x$-Axe forttreibt, $= -\frac{\partial p}{\partial x}\omega\,dx = -\varphi'(\varrho)\frac{\partial\varrho}{\partial x}\omega\,dx$, wenn $\varphi'(\varrho)$ die Derivirte von $\varphi(\varrho)$ bezeichnet. Man hat daher f\"ur $\varrho$ und $u$ die beiden Differentialgleichungen
\begin{align*}
\frac{\partial\varrho}{\partial t} + \frac{\partial(\varrho u)}{\partial x} &= 0,\tag{1}\\
\frac{\partial u}{\partial t} + u\frac{\partial u}{\partial x} &= -\varphi'(\varrho)\frac{\partial\varrho}{\partial x}.\tag{2}
\end{align*}

Wenn man die zweite Gleichung, mit $\pm\sqrt{\varphi'(\varrho)}$ multiplicirt, zur ersteren addirt und zur Abk\"urzung
\begin{align*}
\int\sqrt{\varphi'(\varrho)}\,d\log\varrho &= f(\varrho)\quad\text{und}\\
f(\varrho) + u &= 2r,\quad f(\varrho) - u = 2s\tag{2}
\end{align*}
setzt, so erhalten diese Gleichungen die einfachere Gestalt
\begin{align*}
\frac{\partial r}{\partial t} &= -\bigl(u + \sqrt{\varphi'(\varrho)}\bigr)\frac{\partial r}{\partial x},\\
\frac{\partial s}{\partial t} &= -\bigl(u - \sqrt{\varphi'(\varrho)}\bigr)\frac{\partial s}{\partial x},\tag{3}
\end{align*}
worin $u$ und $\varrho$ durch die Gleichungen~(2) bestimmte Functionen von $r$ und $s$ sind. Aus ihnen folgt
\begin{align*}
dr &= \frac{\partial r}{\partial x}\bigl(dx - (u + \sqrt{\varphi'(\varrho)})\,dt\bigr),\\
ds &= \frac{\partial s}{\partial x}\bigl(dx - (u - \sqrt{\varphi'(\varrho)})\,dt\bigr).\tag{4,5}
\end{align*}

Unter der in der Wirklichkeit immer zutreffenden Voraussetzung, dass $\varphi'(\varrho)$ positiv ist, besagen diese Gleichungen, dass $r$ constant bleibt, wenn $x$ sich mit $t$ so \"andert, dass $dx = (u + \sqrt{\varphi'(\varrho)})\,dt$, und $s$ constant bleibt, wenn $x$ sich mit $t$ so \"andert, dass $dx = (u - \sqrt{\varphi'(\varrho)})\,dt$ ist.

Ein bestimmter Werth von $r$ oder von $f(\varrho) + u$ r\"uckt daher zu gr\"osseren Werthen von $x$ mit der Geschwindigkeit $\sqrt{\varphi'(\varrho)} + u$ fort, ein bestimmter Werth von $s$ oder von $f(\varrho) - u$ zu kleineren Werthen von $x$ mit der Geschwindigkeit $\sqrt{\varphi'(\varrho)} - u$.

Ein bestimmter Werth von $r$ wird also nach und nach mit jedem vor ihm stattfindenden Werthe von $s$ zusammentreffen, und die Geschwindigkeit seines Fortr\"uckens wird in jedem Augenblicke von dem Werthe von $s$ abh\"angen, mit welchem er zusammentrifft.

\newpage
% --- page 160 ---
\begin{center}
\textbf{2.}
\end{center}

Die Analysis bietet nun zun\"achst die Mittel, die Frage zu beantworten, wo und wann ein Werth $r$ von $r$ einem vor ihm befindlichen Werthe $s'$ von $s$ begegnet, d.~h.~$x$ und $t$ als Functionen von $r$ und $s$ zu bestimmen. In der That, wenn man in den Gleichungen~(3) des vor.~Art. $r$ und $s$ als unabh\"angige Variable einf\"uhrt, so gehen diese Gleichungen in lineare Differentialgleichungen f\"ur $x$ und $t$ \"uber und lassen sich also nach bekannten Methoden integriren. Um die Zur\"uckf\"uhrung der Differentialgleichungen auf eine lineare zu bewirken, ist es am zweckm\"assigsten, die Gleichungen (4) und (5) des vorigen Art. in die Form zu setzen:
\begin{align*}
d(x - (u + \sqrt{\varphi'(\varrho)})t) &= \frac{\partial(x - (u + \sqrt{\varphi'(\varrho)})t)}{\partial s}\,ds\\
&\quad + \Bigl(\frac{d\log\sqrt{\varphi'(\varrho)}}{d\log\varrho} + 1\Bigr)t\,dr,\tag{1}\\
d(x - (u - \sqrt{\varphi'(\varrho)})t) &= \frac{\partial(x - (u - \sqrt{\varphi'(\varrho)})t)}{\partial r}\,dr\\
&\quad + \Bigl(\frac{d\log\sqrt{\varphi'(\varrho)}}{d\log\varrho} + 1\Bigr)t\,ds.\tag{2}
\end{align*}

Man erh\"alt dann, wenn man $s$ und $r$ als unabh\"angige Variable betrachtet, f\"ur $x$ und $t$ die beiden linearen Differentialgleichungen:
\begin{align*}
\frac{\partial(x - (u + \sqrt{\varphi'(\varrho)})t)}{\partial s} &= -\Bigl(\frac{d\log\sqrt{\varphi'(\varrho)}}{d\log\varrho} + 1\Bigr)t,\\
\frac{\partial(x - (u - \sqrt{\varphi'(\varrho)})t)}{\partial r} &= -\Bigl(\frac{d\log\sqrt{\varphi'(\varrho)}}{d\log\varrho} + 1\Bigr)t.
\end{align*}

In Folge derselben ist
\begin{equation*}
(x - (u + \sqrt{\varphi'(\varrho)})t)\,dr - (x - (u - \sqrt{\varphi'(\varrho)})t)\,ds\tag{3}
\end{equation*}
ein vollst\"andiges Differential, dessen Integral, $w$, der Gleichung
\[
\frac{\partial^{2}w}{\partial r\,\partial s} + m\Bigl(\frac{\partial w}{\partial r} + \frac{\partial w}{\partial s}\Bigr) = 0
\]
gen\"ugt, worin $m = \frac{1}{2\sqrt{\varphi'(\varrho)}}\frac{d\log\sqrt{\varphi'(\varrho)}}{d\log\varrho}$, also eine Function von $r + s$ ist. Setzt man $f(\varrho) = r + s = \vartheta$, so wird $\sqrt{\varphi'(\varrho)} = \vartheta'$, $\frac{d\log\vartheta'}{d\vartheta} = \frac{d\log\sqrt{\varphi'(\varrho)}}{d\log\varrho} \cdot \frac{d\log\varrho}{d\vartheta}$, folglich $m = \frac{d\log\vartheta'}{d\vartheta}$.

Bei der Poisson'schen Annahme $\varphi(\varrho) = a^{2}\varrho^{k}$ wird $f(\varrho) = \frac{2a\sqrt{k}}{k-1}\varrho^{\frac{k-1}{2}}$, und, wenn man f\"ur die willk\"urliche Constante den Werth Null w\"ahlt,
\begin{align*}
f(\varrho) + u = \frac{k+1}{k-1}r + s,\quad f(\varrho) - u = r + \frac{k+1}{k-1}s,\\
m = \frac{k+1}{2(k-1)(r+s)}.\qquad\qquad\qquad\qquad
\end{align*}

Unter Voraussetzung des Boyle'schen Gesetzes $\varphi(\varrho) = a^{2}\varrho$ erh\"alt man $f(\varrho) = a\log\varrho$ und
\begin{align*}
f(\varrho) + u &= r - s + a,\\
f(\varrho) - u &= s - r + a,\\
m &= -\frac{1}{2a},
\end{align*}
Werthe, die aus den obigen fliessen, wenn man $f(\varrho)$ um die Constante $a$, also $r$ und $s$ um $\frac{a}{2}$, vermindert und dann $k = 1$ setzt.

Die Einf\"uhrung von $r$ und $s$ als unabh\"angig ver\"anderlichen Gr\"ossen ist indess nur m\"oglich, wenn die Determinante dieser Functionen von $x$ und $t$, welche $= 2\sqrt{\varphi'(\varrho)}\frac{\partial(r,s)}{\partial(x,t)}$ nicht verschwindet, also nur wenn $\frac{\partial r}{\partial x}$ und $\frac{\partial s}{\partial x}$ beide von Null verschieden sind.

Wenn $\frac{\partial r}{\partial x} = 0$ ist, ergiebt sich aus (1) $dr = 0$ und aus (2) $x - (u - \sqrt{\varphi'(\varrho)})t =$ einer Function von $s$. Es ist folglich auch dann der Ausdruck (3) ein vollst\"andiges Differential, und es wird $w$ eine blosse Function von $s$.

Aus \"ahnlichen Gr\"unden werden, wenn $\frac{\partial s}{\partial x} = 0$ ist, $s$ auch in Bezug auf $t$ constant, $x - (u + \sqrt{\varphi'(\varrho)})t$ und $w$ Functionen von $r$.

Wenn endlich $\frac{\partial r}{\partial x}$ und $\frac{\partial s}{\partial x}$ beide $= 0$ sind, so werden in Folge der Differentialgleichungen $r$, $s$ und $w$ Constanten.

\newpage
% --- page 161 --- (Art. 3)
\begin{center}
\textbf{3.}
\end{center}

Um die Aufgabe zu l\"osen, muss nun zun\"achst $w$ als Function von $r$ und $s$ so bestimmt werden, dass sie der Differentialgleichung
\begin{equation*}
\frac{\partial^{2}w}{\partial r\,\partial s} + m\Bigl(\frac{\partial w}{\partial r} + \frac{\partial w}{\partial s}\Bigr) = 0\tag{1}
\end{equation*}
und den Anfangsbedingungen gen\"ugt, wodurch sie bis auf eine Constante, die ihr offenbar willk\"urlich hinzugef\"ugt werden kann, bestimmt ist.

Wo und wann ein bestimmter Werth von $r$ mit einem bestimmten Werthe von $s$ zusammentrifft, ergiebt sich dann aus der Gleichung
\begin{equation*}
(x - (u + \sqrt{\varphi'(\varrho)})t)\,dr - (x - (u - \sqrt{\varphi'(\varrho)})t)\,ds = dw,\tag{2}
\end{equation*}
und hierauf findet man schliesslich $u$ und $\varrho$ als Functionen von $x$ und $t$ durch Hinzuziehung der Gleichungen
\begin{equation*}
f(\varrho) + u = 2r,\quad f(\varrho) - u = 2s.\tag{3}
\end{equation*}

In der That folgen, wenn nicht etwa in einer endlichen Strecke $dr$ oder $ds$ Null und folglich $r$ oder $s$ constant ist, aus (2) die Gleichungen
\begin{align*}
x - (u + \sqrt{\varphi'(\varrho)})t &= \frac{\partial w}{\partial r},\\
x - (u - \sqrt{\varphi'(\varrho)})t &= -\frac{\partial w}{\partial s},\tag{4,5}
\end{align*}
durch deren Verbindung mit (3) man $u$ und $\varrho$ in $x$ und $t$ ausgedr\"uckt erh\"alt.

Wenn aber $r$ anfangs in einer endlichen Strecke denselben Werth $r'$ hat, so r\"uckt diese Strecke allm\"ahlich zu gr\"osseren Werthen von $x$ fort. Innerhalb dieses Gebietes, wo $r = r'$, kann man dann aus der Gleichung (2) den Werth von $x - (u + \sqrt{\varphi'(\varrho)})t$ nicht ableiten, da $dr = 0$, und in der That l\"asst die Frage, wo und wann dieser Werth $r'$ einem bestimmten Werthe von $s$ begegnet, dann keine bestimmte Antwort zu. Die Gleichung (4) gilt dann nur an den Grenzen dieses Gebietes und giebt an, zwischen welchen Werthen von $x$ zu einer bestimmten Zeit der constante Werth $r'$ von $r$ stattfindet, oder auch, w\"ahrend welches Zeitraums $r$ an einer bestimmten Stelle diesen Werth beh\"alt. Zwischen diesen Grenzen bestimmen sich $u$ und $\varrho$ als Functionen von $x$ und $t$ aus den Gleichungen (3) und (5). Auf \"ahnlichem Wege findet man diese Functionen, wenn $s$ den Werth $s'$ in einem endlichen Gebiete besitzt, w\"ahrend $r$ ver\"anderlich ist, sowie auch wenn $r$ und $s$ beide constant sind. In letzterem Falle nehmen sie zwischen gewissen durch (4) und (5) bestimmten Grenzen constante aus (3) fliessende Werthe an.

Bevor wir die Integration der Gleichung (1) des vor.~Art. in Angriff nehmen, scheint es zweckm\"assig, einige Er\"orterungen voraufzuschicken, welche die Ausf\"uhrung dieser Integration nicht voraussetzen. \"Uber die Function $\varphi(\varrho)$ ist dabei nur die Annahme n\"othig, dass ihre Derivirte bei wachsendem $\varrho$ nicht abnimmt, was in der Wirklichkeit gewiss immer der Fall ist; und wir bemerken gleich hier, was im folgenden Art. mehrfach angewandt werden wird, dass dann
\[
\frac{\varphi(\varrho_{1}) - \varphi(\varrho_{2})}{\varrho_{1} - \varrho_{2}}\Bigl(\frac{1}{\varrho_{1}} + \frac{1}{\varrho_{2}}\Bigr) - \frac{\varphi(\varrho_{1})}{\varrho_{1}^{2}} - \frac{\varphi(\varrho_{2})}{\varrho_{2}^{2}},
\]
 wenn nur eine der Gr\"ossen $\varrho_{1}$ und $\varrho_{2}$ sich \"andert, entweder constant bleibt oder mit dieser Gr\"osse zugleich w\"achst und abnimmt, woraus zugleich folgt, dass der Werth dieses Ausdrucks stets zwischen $\varphi'(\varrho_{1})$ und $\varphi'(\varrho_{2})$ liegt.

Wir betrachten zun\"achst den Fall, wo die anf\"angliche Gleichgewichtsst\"orung auf ein endliches durch die Ungleichheiten $a < x < b$ begrenztes Gebiet beschr\"ankt ist, so dass ausserhalb desselben $u$ und $\varrho$ und folglich auch $r$ und $s$ constant sind; die Werthe dieser Gr\"ossen f\"ur $x < a$ m\"ogen durch Anh\"angung des Index 1, f\"ur $x > b$ durch den Index 2 bezeichnet werden. Das Gebiet, in welchem $r$ ver\"anderlich ist, bewegt sich nach Art.~1 allm\"ahlich vorw\"arts und zwar seine hintere Grenze mit der Geschwindigkeit $\sqrt{\varphi'(\varrho_{1})} + u_{1}$, w\"ahrend die vordere Grenze des Gebiets, in welchem $s$ ver\"anderlich ist, mit der Geschwindigkeit $\sqrt{\varphi'(\varrho_{2})} - u_{2}$ r\"uckw\"arts geht. Nach Verlauf der Zeit $\frac{b-a}{\sqrt{\varphi'(\varrho_{1})} + \sqrt{\varphi'(\varrho_{2})} + u_{1} - u_{2}}$ fallen daher beide Gebiete auseinander, und zwischen ihnen bildet sich ein Raum, in welchem $s = s_{2}$ und $r = r_{1}$ ist und folglich die Gas-theilchen wieder im Gleichgewicht sind. Von der anfangs ersch\"utterten Stelle gehen also zwei nach entgegengesetzten Richtungen fortschreitende Wellen aus. In der vorw\"artsgehenden ist $s = s_{2}$; es ist daher mit einem bestimmten Werthe $\varrho$ der Dichtigkeit stets die Geschwindigkeit $u = f(\varrho) - 2s_{2}$ verbunden, und beide Werthe r\"ucken mit der constanten Geschwindigkeit $\sqrt{\varphi'(\varrho)} + u = \sqrt{\varphi'(\varrho)} + f(\varrho) - 2s_{2}$ vorw\"arts. In der r\"uckw\"artslaufenden ist dagegen mit der Dichtigkeit $\varrho$ die Geschwindigkeit $-f(\varrho) + 2r_{1}$ verbunden, und diese beiden Werthe bewegen sich mit der Geschwindigkeit $\sqrt{\varphi'(\varrho)} + f(\varrho) - 2r_{1}$ r\"uckw\"arts. Die Fortpflanzungsgeschwindigkeit ist f\"ur gr\"ossere Dichtigkeiten eine gr\"ossere, da sowohl $\sqrt{\varphi'(\varrho)}$, als $f(\varrho)$ mit $\varrho$ zugleich w\"achst.

Denkt man sich $\varrho$ als Ordinate einer Curve f\"ur die Abscisse $x$, so bewegt sich jeder Punkt dieser Curve parallel der Abscissenaxe mit constanter Geschwindigkeit fort und zwar mit desto gr\"osserer, je gr\"osser seine Ordinate ist. Man bemerkt leicht, dass bei diesem Gesetze Punkte mit gr\"osseren Ordinaten schliesslich voraufgehende Punkte mit kleineren Ordinaten \"uberholen w\"urden, so dass zu einem Werthe von $x$ mehr als ein Werth von $\varrho$ geh\"oren w\"urde. Da nun dieses in Wirklichkeit

% --- page 164 ---
nicht stattfinden kann, so muss ein Umstand eintreten, wodurch dieses Gesetz ung\"ultig wird. In der That liegt nun der Herleitung der Differentialgleichungen die Voraussetzung zu Grunde, dass $u$ und $\varrho$ stetige Functionen von $x$ sind und endliche Derivirten haben; diese Voraussetzung h\"ort aber auf erf\"ullt zu sein, sobald in irgend einem Punkte die Dichtigkeitscurve senkrecht zur Abscissenaxe wird, und von diesem Augenblicke an tritt in dieser Curve eine Discontinuit\"at ein, so dass ein gr\"osserer Werth von $\varrho$ einem kleineren unmittelbar nachfolgt; ein Fall, der im n\"achsten Art. er\"ortert werden wird.

Die Verdichtungswellen, d.~h.~die Theile der Welle, in welchen die Dichtigkeit in der Fortpflanzungsrichtung abnimmt, werden demnach bei ihrem Fortschreiten immer schm\"aler und gehen schliesslich in Verdichtungsst\"osse \"uber; die Breite der Verd\"unnungswellen aber w\"achst best\"andig der Zeit proportional.

Es l\"asst sich, wenigstens unter Voraussetzung des Poisson'schen (oder Boyle'schen) Gesetzes, leicht zeigen, dass auch dann, wenn die anf\"angliche Gleichgewichtsst\"orung nicht auf ein endliches Gebiet beschr\"ankt ist, sich stets, von ganz besonderen F\"allen abgesehen, im Laufe der Bewegung Verdichtungsst\"osse bilden m\"ussen. Die Geschwindigkeit, mit welcher ein Werth von $r$ vorw\"arts r\"uckt, ist bei dieser Annahme $\frac{k+1}{2}r + \frac{k-3}{2}s$; gr\"ossere Werthe werden sich also durchschnittlich mit gr\"osserer Geschwindigkeit bewegen, und ein gr\"osserer Werth $r$ wird einen voraufgehenden kleineren Werth $r'$ schliesslich einholen m\"ussen, wenn nicht der mit $r'$ zusammentreffende Werth von $s$ durchschnittlich um $\frac{k+1}{2}(r - r')$ kleiner ist, als der gleichzeitig mit $r$ zusammentreffende. In diesem Falle w\"urde $s$ f\"ur ein positiv unendliches $x$ negativ unendlich werden, und also f\"ur $u = +\infty$ die Geschwindigkeit $u = +\infty$ (oder auch statt dessen beim Boyle'schen Gesetz die Dichtigkeit unendlich klein) werden. Von speciellen F\"allen abgesehen, wird also immer der Fall eintreten m\"ussen, dass ein um eine endliche Gr\"osse gr\"osserer Werth von $r$ einem kleineren unmittelbar nachfolgt; es werden folglich, durch ein Unendlichwerden von $\frac{\partial\varrho}{\partial x}$, die Differentialgleichungen ihre G\"ultigkeit verlieren und vorw\"artslaufende Verdichtungsst\"osse entstehen m\"ussen. Ebenso werden fast immer, indem $\frac{\partial s}{\partial x}$ unendlich wird, r\"uckw\"artslaufende Verdichtungsst\"osse sich bilden.

\newpage
% --- page 165 --- (Art. 5)
\begin{center}
\textbf{5.}
\end{center}

Zur Bestimmung der Zeiten und Orte, f\"ur welche $\frac{\partial\varrho}{\partial x}$ oder $\frac{\partial r}{\partial x}$ unendlich wird und pl\"otzliche Verdichtungen ihren Anfang nehmen, erh\"alt man aus den Gleichungen (1) und (2) des Art.~2, wenn man darin die Function $w$ einf\"uhrt.

Wir m\"ussen nun, da sich pl\"otzliche Verdichtungen fast immer einstellen, auch wenn sich Dichtigkeit und Geschwindigkeit anfangs allenthalben stetig \"andern, die Gesetze f\"ur das Fortschreiten von Verdichtungsst\"ossen aufsuchen.

Wir nehmen an, dass zur Zeit $t$ f\"ur $x = \xi$ eine sprungweise Aenderung von $u$ und $\varrho$ stattfinde, und bezeichnen die Werthe dieser und der von ihnen abh\"angigen Gr\"ossen f\"ur $x = \xi - 0$ durch Anh\"angung des Index 1 und f\"ur $x = \xi + 0$ durch den Index 2; die relativen Geschwindigkeiten, mit welchen das Gas sich gegen die Unstetigkeitsstelle bewegt, $u_{1} - \frac{d\xi}{dt}$, $u_{2} - \frac{d\xi}{dt}$, m\"ogen durch $v_{1}$ und $v_{2}$ bezeichnet werden. Die Masse, welche durch ein Element $\omega$ der Ebene, wo $x = \xi$, im Zeitelement $dt$ in positiver Richtung hindurchgeht, ist dann $= v_{1}\varrho_{1}\omega\,dt = v_{2}\varrho_{2}\omega\,dt$, und die eingedr\"uckte Kraft $(\varphi(\varrho_{1}) - \varphi(\varrho_{2}))\omega\,dt$ und der dadurch bewirkte Zuwachs an Geschwindigkeit $v_{2} - v_{1}$; man hat daher
\[
(\varphi(\varrho_{1}) - \varphi(\varrho_{2}))\omega\,dt = (v_{2} - v_{1})v_{1}\varrho_{1}\omega\,dt\quad\text{und}\quad v_{1}\varrho_{1} = v_{2}\varrho_{2},
\]
woraus folgt $v_{1} = \pm\sqrt{\frac{\varrho_{2}}{\varrho_{1}}}\sqrt{\frac{\varphi(\varrho_{1}) - \varphi(\varrho_{2})}{\varrho_{1} - \varrho_{2}}}$, also
\[
\frac{d\xi}{dt} = \frac{v_{1}\varrho_{1} - v_{2}\varrho_{2}}{\varrho_{1} - \varrho_{2}} = \mp\sqrt{\frac{\varphi(\varrho_{1}) - \varphi(\varrho_{2})}{\varrho_{1} - \varrho_{2}}}\sqrt{\varrho_{1}\varrho_{2}}.\tag{1}
\]

F\"ur einen Verdichtungsstoss muss $\varrho_{2} - \varrho_{1}$ dasselbe Zeichen, wie $v_{1}$ und $v_{2}$, haben, und zwar f\"ur einen vorw\"artslaufenden das negative, f\"ur einen r\"uckw\"artslaufenden das positive. Im erstem Falle gelten die oberen Zeichen und $\varrho_{1}$ ist gr\"osser als $\varrho_{2}$; es ist daher, bei der zu Anfang des vorigen Artikels gemachten Annahme \"uber die Function $\varphi(\varrho)$

% --- page 166 ---
\[
\sqrt{\frac{\varphi(\varrho_{1}) - \varphi(\varrho_{2})}{\varrho_{1} - \varrho_{2}}} > \sqrt{\varphi'(\varrho_{1})}\quad\text{und}\quad < \sqrt{\varphi'(\varrho_{2})},
\]
und folglich r\"uckt die Unstetigkeitsstelle langsamer fort als die nachfolgenden und schneller als die voraufgehenden Werthe von $\varrho$; $r_{1}$ und $r_{2}$ sind also in jedem Augenblicke durch die zu beiden Seiten der Unstetigkeitsstelle geltenden Differentialgleichungen bestimmt. Dasselbe gilt, da die Werthe von $s$ sich mit der Geschwindigkeit $\sqrt{\varphi'(\varrho)} - u$ r\"uckw\"arts bewegen, auch f\"ur $s_{2}$ und folglich f\"ur $\varrho_{2}$ und $u_{2}$, aber nicht f\"ur $s_{1}$. Die Werthe von $s_{2}$ und $\frac{d\xi}{dt}$ bestimmen sich aus $r_{1}$, $\varrho_{2}$ und $u_{2}$ eindeutig durch die Gleichungen (1). In der That gen\"ugt der Gleichung
\begin{equation*}
2(r_{2} - s_{2}) = f(\varrho_{1}) - f(\varrho_{2}) + \sqrt{\frac{\varphi(\varrho_{1}) - \varphi(\varrho_{2})}{\varrho_{1} - \varrho_{2}}}\sqrt{\frac{\varrho_{1} - \varrho_{2}}{\varrho_{2}}}\tag{3}
\end{equation*}
nur ein Werth von $\varrho_{2}$; denn die rechte Seite nimmt, wenn $\varrho_{2}$ von $\varrho_{1}$ an in's Unendliche w\"achst, jeden positiven Werth nur einmal an, da sowohl $f(\varrho_{1})$ als auch die beiden Factoren
\[
\sqrt{\varrho_{1}} - \sqrt{\varrho_{2}}\quad\text{und}\quad\sqrt{\varrho_{1}\varphi(\varrho_{1})} - \sqrt{\varrho_{2}\varphi(\varrho_{2})},
\]
in welche sich das letzte Glied zerlegen l\"asst, best\"andig wachsen oder doch nur der letztere Factor constant bleibt. Wenn aber $\varrho_{2}$ bestimmt ist, erh\"alt man durch die Gleichungen (1) offenbar v\"ollig bestimmte Werthe f\"ur $u_{2}$ und $\frac{d\xi}{dt}$.

Ganz Aehnliches gilt f\"ur einen r\"uckw\"artslaufenden Verdichtungsstoss.

\newpage
% --- page 167 --- (Art. 6)
\begin{center}
\textbf{6.}
\end{center}

Wir haben eben gefunden, dass in einem fortschreitenden Verdichtungsstosse zwischen den Werthen von $u$ und $\varrho$ zu beiden Seiten desselben stets die Gleichung
\[
(u_{2} - u_{1})^{2} = \frac{(\varrho_{1} - \varrho_{2})^{2}}{\varrho_{1}\varrho_{2}}\bigl(\varphi(\varrho_{1}) - \varphi(\varrho_{2})\bigr)\bigl(\frac{1}{\varrho_{1}} - \frac{1}{\varrho_{2}}\bigr)
\]
stattfindet. Es fragt sich nun, was eintritt, wenn zu einer gegebenen Zeit an einer gegebenen Stelle beliebig gegebene Unstetigkeiten vorhanden sind. Es k\"onnen dann von dieser Stelle, je nach den Werthen von $u_{1}$, $\varrho_{1}$, $u_{2}$, $\varrho_{2}$, entweder zwei nach entgegengesetzten Seiten laufende Verdichtungsst\"osse ausgehen, oder ein vorw\"artslaufender, oder ein r\"uckw\"artslaufender, oder endlich kein Verdichtungsstoss, so dass die Bewegung nach den Differentialgleichungen erfolgt.

Bezeichnet man die Werthe, welche $u$ und $\varrho$ hinter oder zwischen den Verdichtungsst\"ossen im ersten Augenblicke ihres Fortschreitens annehmen, durch Hinzuf\"ugung eines Accents, so ist im ersten Falle $\varrho' > \varrho_{1}$ und $> \varrho_{2}$; und man hat
\begin{align*}
u' &= \sqrt{\frac{(\varrho' - \varrho_{1})(\varrho' - \varrho_{2})}{\varrho'\varrho_{1}\varrho'\varrho_{2}}\bigl(\varphi(\varrho') - \varphi(\varrho_{1})\bigr)\bigl(\varrho' - \varrho_{2}\bigr)},\tag{1}\\
u' - u_{1} &= \sqrt{\frac{(\varrho' - \varrho_{1})(\varrho' - \varrho_{2})}{\varrho'\varrho_{1}\varrho'\varrho_{2}}\bigl(\varphi(\varrho') - \varphi(\varrho_{1})\bigr)(\varrho' - \varrho_{2})},\\
u' - u_{2} &= \sqrt{\frac{(\varrho' - \varrho_{1})(\varrho' - \varrho_{2})}{\varrho'\varrho_{1}\varrho'\varrho_{2}}\bigl(\varphi(\varrho') - \varphi(\varrho_{2})\bigr)(\varrho' - \varrho_{1})}.\tag{2}
\end{align*}

Es muss also, da beide Glieder der rechten Seite von (2) mit $\varrho'$ zugleich wachsen, $u_{1} - u_{2}$ positiv sein und
\[
(u_{1} - u_{2})^{2} > \frac{(\varrho_{1} - \varrho_{2})\bigl(\varphi(\varrho_{1}) - \varphi(\varrho_{2})\bigr)}{\varrho_{1}\varrho_{2}}
\]
und umgekehrt giebt es, wenn diese Bedingungen erf\"ullt sind, stets ein und nur ein den Gleichungen (1) gen\"ugendes Werthenpaar von $u$ und $\varrho'$.

Damit der letzte Fall eintritt und also die Bewegung sich den Differentialgleichungen gem\"ass bestimmen l\"asst, ist es nothwendig und hinreichend, dass $r_{1} < r_{2}$ und $s_{1} < s_{2}$ sei, also $u_{1} - u_{2}$ negativ und $(u_{1} - u_{2})^{2} > (f(\varrho_{1}) - f(\varrho_{2}))^{2}$. Die Werthe $r_{1}$ und $r_{2}$, $s_{1}$ und $s_{2}$ treten dann, da der voraufgehende Werth mit gr\"osserer Geschwindigkeit fortr\"uckt, im Fortschreiten auseinander, so dass die Unstetigkeit verschwindet.

Wenn weder die ersteren, noch die letzteren Bedingungen erf\"ullt sind, so gen\"ugt den Anfangswerthen Ein Verdichtungsstoss, und zwar ein vorw\"arts oder r\"uckw\"arts laufender, je nachdem $\varrho_{1}$ gr\"osser oder kleiner als $\varrho_{2}$ ist.

In der That ist dann, wenn $\varrho_{1} > \varrho_{2}$,
\[
2(r_{1} - r_{2}\text{ oder }f(\varrho_{1}) - f(\varrho_{2}) + u_{1} - u_{2})
\]
positiv, --- weil $(u_{1} - u_{2})^{2} < (f(\varrho_{1}) - f(\varrho_{2}))^{2}$ --- und zugleich
\[
(u_{1} - u_{2})^{2} < \frac{(\varrho_{1} - \varrho_{2})\bigl(\varphi(\varrho_{1}) - \varphi(\varrho_{2})\bigr)}{\varrho_{1}\varrho_{2}}
\]
weil
\[
(u_{1} - u_{2})^{2} < (\varrho_{1} - \varrho_{2})\bigl(\varphi(\varrho_{1}) - \varphi(\varrho_{2})\bigr).
\]
Es l\"asst sich also f\"ur die Dichtigkeit $\varrho'$ hinter dem Verdichtungsstoss ein der Bedingung (3) des vor.~Art. gen\"ugender Werth finden und dieser ist $< \varrho_{1}$. Folglich wird, da $s = f(\varrho') - u'$, $s_{2} = f(\varrho_{1}) - u_{1}$, auch $s' < s_{1}$, so dass die Bewegung hinter dem Verdichtungsstosse nach den Differentialgleichungen erfolgen kann.

Der andere Fall, wenn $\varrho_{1} < \varrho_{2}$, ist offenbar von diesem nicht wesentlich verschieden.

\newpage
% --- page 168 --- (Art. 7)
\begin{center}
\textbf{7.}
\end{center}

Um das Bisherige durch ein einfaches Beispiel zu erl\"autern, wo sich die Bewegung mit den bis jetzt gewonnenen Mitteln bestimmen l\"asst, wollen wir annehmen, dass Druck und Dichtigkeit von einander nach dem Boyle'schen Gesetz abh\"angen und anfangs Dichtigkeit und Geschwindigkeit sich bei $x = 0$ sprungweise \"andern, aber zu beiden Seiten dieser Stelle constant sind.

Es sind dann nach dem Obigen vier F\"alle zu unterscheiden.

\smallskip
\noindent\textbf{I.} Wenn $u_{1} - u_{2} > 0$, also die beiden Gasmassen sich einander entgegen bewegen und $\bigl(\frac{u_{1} - u_{2}}{a}\bigr)^{2} > \frac{(\varrho_{1} - \varrho_{2})^{2}}{\varrho_{1}\varrho_{2}}$, so bilden sich zwei entgegengesetzt laufende Verdichtungsst\"osse. Nach Art.~6~(1) ist, wenn $\sigma$ durch $a$ und durch $\varrho'$ die positive Wurzel der Gleichung
\[
\frac{(\varrho' - \varrho_{1})(\varrho' - \varrho_{2})}{\varrho'\varrho_{1}\varrho'\varrho_{2}} = \frac{(\varrho_{1} - \varrho_{2})^{2}}{\varrho_{1}\varrho_{2}}
\]
bezeichnet wird, die Dichtigkeit zwischen den Verdichtungsst\"ossen $\varrho' = \varrho_{1}\varrho_{2}\sigma^{2}$ und nach Art.~5~(1) hat man f\"ur den vorw\"artslaufenden Verdichtungsstoss
\[
\frac{d\xi}{dt} = u_{1} + a\sigma = u' - a\sigma,
\]
f\"ur den r\"uckw\"artslaufenden
\[
\frac{d\xi}{dt} = u_{2} - a\sigma = u' + a\sigma;
\]
die Werthe der Geschwindigkeit und Dichtigkeit sind also nach Verlauf der Zeit $t$, wenn $(u_{1} - a\sigma)t < x < (u_{2} + a\sigma)t$, $u'$ und $\varrho'$; f\"ur ein kleineres $x$ $u_{1}$ und $\varrho_{1}$ und f\"ur ein gr\"osseres $u_{2}$ und $\varrho_{2}$.

\smallskip
\noindent\textbf{II.} Wenn $u_{1} - u_{2} < 0$, folglich die Gasmassen sich aus einander bewegen, und zugleich
\[
\frac{(u_{1} - u_{2})^{2}}{a^{2}} < \frac{(\varrho_{1} - \varrho_{2})^{2}}{\varrho_{1}\varrho_{2}},
\]
so gehen von der Grenze nach entgegengesetzten Richtungen zwei allm\"ahlich breiter werdende Verd\"unnungswellen aus. Nach Art.~4 ist zwischen ihnen $r = r_{1}$, $s = s_{2}$, $u = r_{1} - s_{2}$. In der vorw\"artslaufenden ist $s = s_{2}$ und $x - (u + a)t$ eine Function von $r$, deren Werth, aus den Anfangswerthen $t = 0$, $x = 0$, sich $= 0$ findet; f\"ur die r\"uckw\"artslaufende dagegen hat man $r = r_{1}$ und $x - (u - a)t = 0$. Die eine Gleichung zur Bestimmung von $u$ und $\varrho$ ist also, wenn

% --- page 169 ---
\[
(u_{1} - s_{2} + a)t < x < (u_{2} + a)t,\quad u = -a + \frac{x}{t};
\]
f\"ur kleinere Werthe von $x$ $r = r_{1}$ und f\"ur gr\"ossere $r = r_{2}$; die andere Gleichung ist, wenn
\[
(u_{1} - a)t < x < (u_{1} - s_{2} - a)t,\quad u = a + \frac{x}{t},
\]
f\"ur ein kleineres $x$ $s = s_{1}$ und f\"ur ein gr\"osseres $s = s_{2}$.

\smallskip
\noindent\textbf{III.} Wenn keiner dieser beiden F\"alle stattfindet und $\varrho_{1} > \varrho_{2}$, so entsteht eine r\"uckw\"artslaufende Verd\"unnungswelle und ein vorw\"artschreitender Verdichtungsstoss. F\"ur letzteren findet sich aus Art.~5~(3), wenn $\varepsilon$ die Wurzel der Gleichung
\[
\frac{\varrho_{1}\varepsilon^{2} - \varrho_{2}}{\varrho_{1}\varepsilon^{2}} = 2\log\varepsilon + \varepsilon^{-2} - 1
\]
bezeichnet, $\varrho' = \varrho_{1}\varepsilon^{2}\varrho_{2}$ und aus Art.~5~(1)
\[
\frac{d\xi}{dt} = u_{2} + a\varepsilon = u' + a\varepsilon.
\]
Nach Verlauf der Zeit $t$ ist demnach vor dem Verdichtungsstosse, also wenn $x > (u_{2} + a\varepsilon)t$, $u = u_{2}$, $\varrho = \varrho_{2}$; hinter dem Verdichtungsstosse aber hat man $r = r_{1}$ und ausserdem, wenn
\[
(u_{1} - a)t < x < (u' - a)t,\quad u = a + \frac{x}{t},
\]
f\"ur ein kleineres $x$ $u = u_{1}$ und f\"ur ein gr\"osseres $u = u'$.

\smallskip
\noindent\textbf{IV.} Wenn endlich die beiden ersten F\"alle nicht stattfinden und $\varrho_{1} < \varrho_{2}$, so ist der Verlauf ganz wie in III., nur der Richtung nach entgegengesetzt.

\newpage
% --- page 170 --- (Art. 8)
\begin{center}
\textbf{8.}
\end{center}

Um unsere Aufgabe allgemein zu l\"osen, muss nach Art.~3 die Function $w$ so bestimmt werden, dass sie der Differentialgleichung
\begin{equation*}
\frac{\partial^{2}w}{\partial r\,\partial s} + m\Bigl(\frac{\partial w}{\partial r} + \frac{\partial w}{\partial s}\Bigr) = 0\tag{1}
\end{equation*}
und den Anfangsbedingungen gen\"ugt.

Schliessen wir den Fall aus, dass Unstetigkeiten eintreten, so sind offenbar nach Art.~1 Ort und Zeit oder die Werthe von $x$ und $t$, f\"ur welche ein bestimmter Werth $r'$ von $r$ mit einem bestimmten Werthe $s'$ von $s$ zusammentrifft, v\"ollig bestimmt, wenn die Anfangswerthe von $r$ und $s$ f\"ur die Strecke zwischen den beiden Werthen $r'$ von $r$ und $s'$ von $s$ gegeben sind und \"uberall in dem Gr\"ossengebiet $(S)$, welches f\"ur jeden Werth von $t$ die zwischen den beiden Werthen, wo $r = r'$ und $s = s'$, liegenden Werthe von $x$ umfasst, die Differentialgleichungen (3) des Art.~1 erf\"ullt sind. Es ist also auch der Werth

% --- page 171 ---
von $w$ f\"ur $r = r'$, $s = s'$ v\"ollig bestimmt, wenn $w$ \"uberall in dem Gr\"ossengebiet $(S)$ der Differentialgleichung (1) gen\"ugt und f\"ur die Anfangswerthe von $r$ und $s$ die Werthe von $\frac{\partial w}{\partial r}$ und $\frac{\partial w}{\partial s}$, also, bis auf eine additive Constante, auch von $w$ gegeben sind und diese Constante beliebig gew\"ahlt worden ist. Denn diese Bedingungen sind mit den obigen gleichbedeutend. Auch folgt aus Art.~3 noch, dass $\frac{\partial w}{\partial r}$ zwar zu beiden Seiten eines Werthes $r'$ von $r$, wenn dieser Werth in einer endlichen Strecke stattfindet, verschiedene Werthe annimmt, sich aber allenthalben stetig mit $s$ \"andert; ebenso \"andert sich $\frac{\partial w}{\partial s}$ mit $r$, die Function $w$ selbst aber sowohl mit $r$, als mit $s$ allenthalben stetig.

Nach diesen Vorbereitungen k\"onnen wir nun an die L\"osung unserer Aufgabe gehen, an die Bestimmung des Werthes von $w$ f\"ur zwei beliebige Werthe, $r'$ und $s'$, von $r$ und $s$.

Zur Veranschaulichung denke man sich $x$ und $t$ als Abscisse und Ordinate eines Punkts in einer Ebene und in dieser Ebene die Curven gezogen, wo $r$ und wo $s$ constante Werthe hat. Von diesen Curven m\"ogen die ersteren durch $(r)$, die letzteren durch $(s)$ bezeichnet und in ihnen die Richtung, in welcher $t$ w\"achst, als die positive betrachtet werden. Das Gr\"ossengebiet $(S)$ wird dann repr\"asentirt durch ein St\"uck der Ebene, welches begrenzt ist durch die Curve $(r')$, die Curve $(s')$ und das zwischen beiden liegende St\"uck der Abscissenaxe, und es handelt sich darum, den Werth von $w$ in dem Durchschnittspunkte der beiden ersteren aus den in letzterer Linie gegebenen Werthen zu bestimmen. Wir wollen die Aufgabe noch etwas verallgemeinern und annehmen, dass das Gr\"ossengebiet $(S)$, statt durch diese letztere Linie, durch eine beliebige Curve $c$ begrenzt werde, welche keine der Curven $(r)$ und $(s)$ mehr als einmal schneidet, und dass f\"ur die dieser Curve angeh\"origen Werthenpaare von $r$ und $s$ die Werthe von $\frac{\partial w}{\partial r}$ und $\frac{\partial w}{\partial s}$ gegeben seien. Wie sich aus der Aufl\"osung der Aufgabe ergeben wird, unterliegen auch dann diese Werthe von $\frac{\partial w}{\partial r}$ und $\frac{\partial w}{\partial s}$ nur der Bedingung, sich stetig mit dem Ort in der Curve zu \"andern, k\"onnen aber \"ubrigens willk\"urlich angenommen werden, w\"ahrend diese Werthe nicht von einander unabh\"angig sein w\"urden, wenn die Curve $c$ eine der Curven $(r)$ oder $(s)$ mehr als einmal schnitte.

Um Functionen zu bestimmen, welche linearen partiellen Differentialgleichungen und linearen Grenzbedingungen gen\"ugen sollen, kann man ein ganz \"ahnliches Verfahren anwenden, wie wenn man zur Aufl\"osung

% --- page 172 ---
eines Systems von linearen Gleichungen s\"ammtliche Gleichungen, mit unbestimmten Factoren multiplicirt, addirt und diese Factoren dann so bestimmt, dass aus der Summe alle unbekannten Gr\"ossen bis auf eine herausfallen.

Man denke sich das St\"uck $(S)$ der Ebene durch die Curven $(r)$ und $(s)$ in unendlich kleine Parallelogramme zerschnitten und bezeichne durch $dr$ und $ds$ die Aenderungen, welche die Gr\"ossen $r$ und $s$ erleiden, wenn die Curvenelemente, welche die Seiten dieser Parallelogramme bilden, in positiver Richtung durchlaufen werden; man bezeichne ferner durch $v$ eine beliebige Function von $r$ und $s$, welche allenthalben stetig ist und stetige Derivirten hat. In Folge der Gleichung (1) hat man dann
\[
\iint_{S}\Bigl(\frac{\partial^{2}v}{\partial r\,\partial s}w + \frac{\partial v}{\partial r}m\frac{\partial w}{\partial s} + \frac{\partial v}{\partial s}m\frac{\partial w}{\partial r}\Bigr)\,dr\,ds = 0.
\]
\"Uber das ganze Gr\"ossengebiet $(S)$ ausgedehnt. Es muss nun die rechte Seite dieser Gleichung nach den Unbekannten geordnet, d.~h.~hier, das Integral durch partielle Integration so umgeformt werden, dass es ausser bekannten Gr\"ossen nur die gesuchte Function, nicht ihre Derivirten enth\"alt. Bei Ausf\"uhrung dieser Operation geht das Integral zun\"achst \"uber in das \"uber $(S)$ ausgedehnte Integral
\[
\iint\Bigl(w\frac{\partial^{2}v}{\partial r\,\partial s} - \frac{\partial v}{\partial r}m\frac{\partial w}{\partial s} - \frac{\partial v}{\partial s}m\frac{\partial w}{\partial r}\Bigr)\,dr\,ds
\]
und ein einfaches Integral, welches sich, weil sich $\frac{\partial w}{\partial r}$ mit $s$, $\frac{\partial w}{\partial s}$ mit $r$ und $w$ mit beiden Gr\"ossen stetig \"andert, nur \"uber die Begrenzung von $(S)$ erstrecken wird. Bedeuten $dr$ und $ds$ die Aenderungen von $r$ und $s$ in einem Begrenzungselemente, wenn die Begrenzung in der Richtung durchlaufen wird, welche gegen die Richtung nach Innen ebenso liegt, wie die positive Richtung in den Curven $(r)$ gegen die positive Richtung in den Curven $(s)$, so ist dies Begrenzungsintegral
\[
\int\Bigl(v\frac{\partial w}{\partial s} - mvw\Bigr)ds + \Bigl(v\frac{\partial w}{\partial r} + mvw\Bigr)dr.
\]
Das Integral durch die ganze Begrenzung von $S$ ist gleich der Summe der Integrale durch die Curven $c$, $(s')$, $(r')$, welche diese Begrenzung bilden, also, wenn ihre Durchschnittspunkte durch $(c,r')$, $(c,s')$, $(r',s')$ bezeichnet werden,
\[
\int_{c,r'}^{c,s'} + \int_{c,s'}^{r',s'} + \int_{r',s'}^{c,r'}.
\]
Von diesen drei Bestandtheilen enth\"alt der erste ausser der Function $v$ nur bekannte Gr\"ossen, der zweite enth\"alt, da in ihm $ds = 0$ ist, nur die unbekannte Function $w$ selbst, nicht ihre Derivirten; der dritte Bestandtheil aber kann durch partielle Integration in
\[
(vw)_{r',s'} - (vw)_{c,r'} + \int_{r',s'}^{c,r'}w\Bigl(\frac{\partial v}{\partial r} + mv\Bigr)dr
\]
verwandelt werden, so dass in ihm ebenfalls nur die gesuchte Function $w$ selbst vorkommt.

Nach diesen Umformungen liefert die Gleichung (2) offenbar den Werth der Function $w$ im Punkte $(r', s')$, durch bekannte Gr\"ossen ausgedr\"uckt, wenn man die Function $v$ den folgenden Bedingungen gem\"ass bestimmt:

\begin{enumerate}
\item allenthalben in $S$: $\displaystyle\frac{\partial^{2}v}{\partial r\,\partial s} + m\Bigl(\frac{\partial v}{\partial r} + \frac{\partial v}{\partial s}\Bigr) = 0$;
\item f\"ur $r = r'$: $\displaystyle\frac{\partial v}{\partial r} + mv = 0$;
\item f\"ur $s = s'$: $\displaystyle\frac{\partial v}{\partial s} + mv = 0$;
\item f\"ur $r = r'$, $s = s'$: $v = 1$.
\end{enumerate}

Man hat dann
\begin{equation*}
w_{r',s'} = (vw)_{c,r'} + \int_{c,r'}^{c,s'}\Bigl(v\frac{\partial w}{\partial s} - mvw\Bigr)ds + \Bigl(v\frac{\partial w}{\partial r} + mvw\Bigr)dr.\tag{4}
\end{equation*}

\newpage
% --- page 173 --- (Art. 9)
\begin{center}
\textbf{9.}
\end{center}

Durch das eben angewandte Verfahren wird die Aufgabe, eine Function $w$ einer linearen Differentialgleichung und linearen Grenzbedingungen gem\"ass zu bestimmen, auf die L\"osung einer \"ahnlichen, aber viel einfacheren Aufgabe f\"ur eine andere Function $v$ zur\"uckgef\"uhrt; die Bestimmung dieser Function erreicht man meistens am Leichtesten durch Behandlung eines speciellen Falls jener Aufgabe nach der Fourier'schen Methode. Wir m\"ussen uns hier begn\"ugen, diese Rechnung nur anzudeuten und das Resultat auf anderem Wege zu beweisen.\footnote{Weitere Ausf\"uhrungen sind in dem Nachlass vorhanden.}

F\"uhrt man in der Gleichung (1) des vor.~Art. f\"ur $r$ und $s$ als unabh\"angig ver\"anderliche Gr\"ossen $\vartheta = r + s$ und $u = r - s$ ein und w\"ahlt man f\"ur die Curve $c$ eine Curve, in welcher $\vartheta$ constant ist, so l\"asst sich die Aufgabe nach den Regeln Fourier's behandeln, und man erh\"alt durch Vergleichung des Resultats mit der Gleichung (4) des vor.~Art., wenn $r + s = \vartheta$, $r - s' = u'$ gesetzt wird,
\[
v = \int_{0}^{\infty}\cos\mu(u - u')\frac{\psi_{1}(\vartheta')\psi_{2}(\vartheta) - \psi_{2}(\vartheta')\psi_{1}(\vartheta)}{\psi_{1}'(\vartheta) - \psi_{2}'(\vartheta)}\,d\mu,
\]
worin $\psi_{1}(\vartheta)$ und $\psi_{2}(\vartheta)$ zwei solche particul\"are L\"osungen der Differentialgleichung $\psi'' - 2m\psi' + \mu^{2}\psi = 0$ bezeichnen, dass
\[
\psi_{1}(\vartheta_{0}) = 1,\quad \psi_{2}(\vartheta_{0}) = 0,\quad \psi_{1}'(\vartheta_{0}) = 0,\quad \psi_{2}'(\vartheta_{0}) = 1.
\]

Bei Voraussetzung des Poisson'schen Gesetzes, nach welchem $m = \frac{k+1}{2(k-1)(r+s)}$, kann man $\psi_{1}$ und $\psi_{2}$ durch bestimmte Integrale ausdr\"ucken, so dass man f\"ur $v$ ein dreifaches Integral erh\"alt, durch dessen Reduction sich ergiebt
\[
v = \Bigl(\frac{r+s}{r'+s'}\Bigr)^{\frac{k-1}{2}}P\Bigl(\frac{k-1}{2}, \frac{1}{2}, \frac{(r-r')(s-s')}{(r+s)(r'+s')}\Bigr),
\]
worin $P(\alpha, \beta, x)$ die Gauss'sche Reihe bezeichnet.

Man kann nun die Richtigkeit dieses Ausdrucks leicht beweisen, indem man zeigt, dass er wirklich den Bedingungen (3) des vor.~Art. gen\"ugt.

Setzt man $v = e^{-\int md\vartheta}y$, so gehen diese f\"ur $y$ \"uber in
\[
\frac{\partial^{2}y}{\partial r\,\partial s} + \Bigl(\frac{dm}{d\vartheta} - m^{2}\Bigr)y = 0
\]
und $y = 1$ sowohl f\"ur $r = r'$, als f\"ur $s = s'$. Bei der Poisson'schen Annahme kann man aber diesen Bedingungen gen\"ugen, wenn man annimmt, dass $y$ eine Function von $z = \frac{(r-r')(s-s')}{(r+s)(r'+s')}$ sei. Denn es wird dann, wenn man $\frac{d\log\sqrt{\varphi'(\varrho)}}{d\log\varrho}$ durch $\lambda$ bezeichnet, $m = \frac{\lambda}{\vartheta}$, also $3mm = \frac{3\lambda^{2}}{\vartheta^{2}}$ und
\[
\frac{dm}{d\vartheta} = -\frac{\lambda}{\vartheta^{2}},\quad\frac{dm}{d\vartheta} - m^{2} = -\frac{\lambda(\lambda+1)}{\vartheta^{2}}.
\]

Da
\[
\frac{\partial^{2}y}{\partial r\,\partial s} = \frac{d^{2}y}{dz^{2}}\frac{\partial z}{\partial r}\frac{\partial z}{\partial s} + \frac{dy}{dz}\frac{\partial^{2}z}{\partial r\,\partial s},
\]
\[
\frac{\partial z}{\partial r} = \frac{s+s'}{(r+s)^{2}(r'+s')}(s-s'),\quad \frac{\partial z}{\partial s} = \frac{r+r'}{(r+s)^{2}(r'+s')}(r-r'),
\]
und $\frac{\partial^{2}z}{\partial r\,\partial s} = \frac{1+z}{(r+s)^{2}}$, so wird
\[
\frac{\partial^{2}y}{\partial r\,\partial s} = \frac{(r-r')(s-s')}{(r+s)^{4}(r'+s')^{2}}\Bigl[(s+s')(r+r')\frac{d^{2}y}{dz^{2}} + (r+s)(r'+s')(1+z)\frac{dy}{dz}\Bigr].
\]

% --- page 174 ---
Setzt man nun
\[
v = e^{-\int md\vartheta}y = \Bigl(\frac{r+s}{r'+s'}\Bigr)^{-\lambda}y,
\]
so gen\"ugt diese Function allen Bedingungen (3) des Art.~8, und es wird $y$ eine L\"osung der Differentialgleichung
\[
z(1-z)\frac{d^{2}y}{dz^{2}} + (1 - (2\lambda+2)z)\frac{dy}{dz} - \lambda(\lambda+1)y = 0
\]
oder nach der in meiner Abhandlung \"uber die Gauss'sche Reihe eingef\"uhrten Bezeichnung eine Function
\[
P\Bigl(\lambda, \frac{1}{2}, \lambda + \frac{1}{2}, z\Bigr)
\]
und zwar diejenige particul\"are L\"osung, welche f\"ur $z = 0$ gleich $1$ wird.

Nach den in jener Abhandlung entwickelten Transformationsprincipien l\"asst sich $y$ nicht bloss durch die Functionen $P(0, 2\lambda+1, 0, z)$, sondern auch durch die Functionen $P(-\lambda, 0, \lambda+1, z)$, $P(0, \lambda+\frac{1}{2}, \lambda+1, z)$ ausdr\"ucken; man erh\"alt daher f\"ur $y$ eine grosse Menge von Darstellungen durch hypergeometrische Reihen und bestimmte Integrale, von denen wir hier nur die folgenden
\[
y = F(\lambda+1, -\lambda, 1, z) = (1-z)^{\lambda}F\Bigl(\tfrac{1}{2}, \lambda, 1, \frac{z}{z-1}\Bigr) = (1-z)^{-\lambda}F\Bigl(\tfrac{1}{2}+\lambda, 1+\lambda, 1, \frac{z}{z-1}\Bigr)
\]
bemerken, mit denen man in allen F\"allen ausreicht.

Um aus diesen f\"ur das Poisson'sche Gesetz gefundenen Resultaten die f\"ur das Boyle'sche geltenden abzuleiten, muss man nach Art.~2 die Gr\"ossen $r$, $s$, $r'$, $s'$ um $\frac{a}{2}$ vermindern und dann $k = 1$ werden lassen, wodurch man erh\"alt $m = -\frac{1}{2a}$ und
\[
v = e^{\frac{r-r'+s-s'}{2a}}F\Bigl(\tfrac{1}{2}, \tfrac{1}{2}, 1, \frac{(r'-r)(s'-s)}{4a^{2}}\Bigr).
\]

\newpage
% --- page 175 --- (Art. 10)
\begin{center}
\textbf{10.}
\end{center}

Wenn man den im vor.~Art. gefundenen Ausdruck f\"ur $v$ in die Gleichung (4) des Art.~8 einsetzt, erh\"alt man den Werth von $w$ f\"ur $r = r'$, $s = s'$ durch die Werthe von $w$, $\frac{\partial w}{\partial r}$ und $\frac{\partial w}{\partial s}$ in der Curve $c$ ausgedr\"uckt; da aber bei unserm Problem in dieser Curve immer nur $\frac{\partial w}{\partial r}$ und $\frac{\partial w}{\partial s}$ unmittelbar gegeben sind und $w$ erst durch eine Quadratur aus ihnen gefunden werden m\"usste, so ist es zweckm\"assig, den Ausdruck f\"ur $w_{r',s'}$ so umzuformen, dass unter dem Integralzeichen nur die Derivirten von $w$ vorkommen.

Man bezeichne die Integrale der Ausdr\"ucke $-mv\,ds + (\frac{\partial v}{\partial r} + mv)\,dr$ und $(\frac{\partial v}{\partial s} + mv)\,ds - mv\,dr$, welche in Folge der Gleichung
\[
\frac{\partial^{2}v}{\partial r\,\partial s} + m\frac{\partial v}{\partial r} + m\frac{\partial v}{\partial s} = 0
\]
vollst\"andige Differentiale sind, durch $P$ und $U$ und das Integral von $P\,dr + U\,ds$, welcher Ausdruck wegen $\frac{\partial P}{\partial s} = -mv = \frac{\partial U}{\partial r}$ ebenfalls ein vollst\"andiges Differential ist, durch $\omega$.

Bestimmt man nun die Integrationsconstanten in diesen Integralen so, dass $\omega$, $P$ und $U$ f\"ur $r = r'$, $s = s'$ verschwinden, so gen\"ugt $\omega$ den Gleichungen $\frac{\partial\omega}{\partial r} = P$, $\frac{\partial\omega}{\partial s} = U$ und sowol f\"ur $r = r'$, als f\"ur $s = s'$ der Gleichung $\omega = 0$ und ist, beil\"aufig bemerkt, durch diese Grenzbedingung und die Differentialgleichung
\[
\frac{\partial^{2}\omega}{\partial r\,\partial s} + m\Bigl(\frac{\partial\omega}{\partial r} + \frac{\partial\omega}{\partial s}\Bigr) = 0
\]
v\"ollig bestimmt.

F\"uhrt man nun in dem Ausdrucke von $w_{r',s'}$ f\"ur $v$ die Function $\omega$ ein, so kann man ihn durch partielle Integration in
\begin{equation*}
w_{r',s'} = w_{c,r'} + \int_{c,r'}^{c,s'}\Bigl(\frac{\partial\omega}{\partial s}\frac{\partial w}{\partial r}\,ds - \frac{\partial\omega}{\partial r}\frac{\partial w}{\partial s}\,dr\Bigr)\tag{1}
\end{equation*}
umwandeln.

Um die Bewegung des Gases aus dem Anfangszustande zu bestimmen, muss man f\"ur $c$ die Curve, in welcher $t = 0$ ist, nehmen; in dieser Curve hat man dann $\frac{\partial w}{\partial r} = x$, $\frac{\partial w}{\partial s} = -x$, und man erh\"alt durch abermalige partielle Integration
\[
w_{r',s'} = w_{c,r'} + \int(\omega\,dx - x\,d\omega),
\]
folglich nach Art.~3, (4) und (5)
\begin{align*}
(x - (\sqrt{\varphi'(\varrho)} + u)t)_{r',s'} &= x_{r'} + \int\omega\,dx,\\
(x - (\sqrt{\varphi'(\varrho)} - u)t)_{r',s'} &= x_{s'} - \int\omega\,dx.\tag{2}
\end{align*}

Diese Gleichungen (2) dr\"ucken aber die Bewegung nur aus, so lange $\frac{\partial r}{\partial x} + (\frac{d\log\sqrt{\varphi'(\varrho)}}{d\log\varrho} + 1)\frac{\partial r}{\partial t}$ und $\frac{\partial s}{\partial x} + (\frac{d\log\sqrt{\varphi'(\varrho)}}{d\log\varrho} + 1)\frac{\partial s}{\partial t}$ von Null verschieden bleiben. Sobald eine dieser Gr\"ossen verschwindet, entsteht ein Verdichtungsstoss, und die Gleichung (1) gilt dann nur innerhalb solcher Gr\"ossengebiete, welche ganz auf einer und derselben Seite dieses Verdichtungsstosses liegen. Die hier entwickelten Principien reichen dann, wenigstens im Allgemeinen, nicht aus, um aus dem Anfangszustande die Bewegung zu bestimmen; wohl aber kann man mit H\"ulfe der Gleichung (1) und der Gleichungen, welche nach Art.~5 f\"ur den Verdichtungsstoss gelten, die Bewegung bestimmen, wenn der Ort des Verdichtungsstosses zur Zeit $t$, also $\xi$ als Function von $t$, gegeben ist. Wir wollen indess dies nicht weiter verfolgen und verzichten auch auf die Behandlung des Falles, wenn die Luft durch eine feste Wand begrenzt ist, da die Rechnung keine Schwierigkeiten hat und eine Vergleichung der Resultate mit der Erfahrung gegenw\"artig noch nicht m\"oglich ist.

\newpage
%===================================================================
% PAPER IX: Selbstanzeige
%===================================================================
\begin{center}
\Large\textbf{IX.}
\end{center}
\vspace{-0.5em}
\begin{center}
\Large\textbf{Selbstanzeige der vorstehenden Abhandlung.}
\end{center}
\vspace{0.3em}
\begin{center}
\small (G\"ottinger Nachrichten, 1859, Nr.~19.)
\end{center}

\bigskip

Diese Untersuchung macht nicht darauf Anspruch, der experimentellen Forschung n\"utzliche Ergebnisse zu liefern; der Verfasser w\"unscht sie nur als einen Beitrag zur Theorie der nicht linearen partiellen Differentialgleichungen betrachtet zu sehen. Wie f\"ur die Integration der linearen partiellen Differentialgleichungen die fruchtbarsten Methoden nicht durch Entwicklung des allgemeinen Begriffs dieser Aufgabe gefunden worden, sondern vielmehr aus der Behandung specieller physikalischer Probleme hervorgegangen sind, so scheint auch die Theorie der nichtlinearen partiellen Differentialgleichungen durch eine eingehende, alle Nebenbedingungen ber\"ucksichtigende, Behandlung specieller physikalischer Probleme am meisten gef\"ordert zu werden, und in der That hat die L\"osung der ganz speciellen Aufgabe, welche den Gegenstand dieser Abhandlung bildet, neue Methoden und Auffassungen erfordert, und zu Ergebnissen gef\"uhrt, welche wahrscheinlich auch bei allgemeineren Aufgaben eine Rolle spielen werden.

Durch die vollst\"andige L\"osung dieser Aufgabe d\"urften die vor einiger Zeit zwischen den englischen Mathematikern Challis, Airy und Stokes lebhaft verhandelten Fragen\footnote{Phil.~mag.~voll.~33, 34, 36.}), soweit dies nicht schon durch Stokes\footnote{Phil.~mag.~vol.~33, p.~349.} geschehen ist, zu klarer Entscheidung gebracht worden sein, so wie auch der Streit, welcher \"uber eine andere denselben Gegenstand betreffende Frage in der K.~K.~Ges.~d.~W.~zu Wien zwischen den Herrn Petzval, Doppler und A.~von Ettinghausen\footnote{Sitzungsberichte der K.~K.~Ges.~d.~W. vom 15.~Jan., 21.~Mai und 1.~Juni 1852.} gef\"uhrt wurde.

Das einzige empirische Gesetz, welches ausser den allgemeinen Bewegungsgesetzen bei dieser Untersuchung vorausgesetzt werden

% --- page 177 ---
musste, ist das Gesetz, nach welchem der Druck eines Gases sich mit der Dichtigkeit \"andert, wenn es keine W\"arme aufnimmt oder abgiebt. Die schon von Poisson gemachte, aber damals auf sehr unsicherer Grundlage ruhende Annahme, dass der Druck bei der Dichtigkeit $\varrho$ proportional $\varrho^{k}$ sich \"andere, wenn $k$ das Verh\"altniss der specifischen W\"arme bei constantem Druck zu der bei constantem Volumen bedeutet, kann jetzt durch die Versuche von Regnault \"uber die specifischen W\"armen der Gase und ein Princip der mechanischen W\"armetheorie begr\"undet werden, und es schien n\"othig diese Begr\"undung des Poisson'schen Gesetzes, da sie noch wenig bekannt zu sein scheint, in der Einleitung voranzuschicken. Der Werth von $k$ findet sich dabei $= 1{,}4101$, w\"ahrend die Schallgeschwindigkeit bei $0^{\circ}$C. und trockner Luft nach den Versuchen von Martins und A.~Bravais\footnote{Ann.~de chim.~et de phys. S\"er.~III, T.~XIII, p.~5.} $= \frac{332^{m},37}{\sqrt{1,4101}}$ sich ergeben und f\"ur $k$ den Werth $1{,}4095$ liefern w\"urde.

Obwohl die Vergleichung der Resultate unserer Untersuchung mit der Erfahrung durch Versuche und Beobachtungen grosse Schwierigkeiten hat und gegenw\"artig kaum ausf\"uhrbar sein wird, so m\"ogen diese doch, soweit es ohne Weitl\"aufigkeit m\"oglich ist, hier mitgetheilt werden.

Die Abhandlung behandelt die Bewegung der Luft oder eines Gases nur f\"ur den Fall, wenn anfangs und also auch in der Folge die Bewegung allenthalben gleich gerichtet ist, und in jeder auf ihrer Richtung senkrechten Ebene Geschwindigkeit und Dichtigkeit constant sind.

F\"ur den Fall, wo die anf\"angliche Gleichgewichtsst\"orung auf eine endliche Strecke beschr\"ankt ist, ergiebt sich bekanntlich bei der gew\"ohnlichen Voraussetzung, dass die Druckverschiedenheiten unendlich kleine Bruchtheile des ganzen Drucks sind, das Resultat, dass von der ersch\"utterten Stelle zwei Wellen, in deren jeder die Geschwindigkeit eine bestimmte Function der Dichtigkeit ist, ausgehen und in entgegengesetzten Richtungen mit der bei dieser Voraussetzung constanten Geschwindigkeit $\sqrt{\varphi'(\varrho)}$ fortschreiten, wenn $\varphi(\varrho)$ den Druck bei der Dichtigkeit $\varrho$ und $\varphi'(\varrho)$ die Derivirte dieser Function bezeichnet. Etwas ganz \"ahnliches gilt nun f\"ur diesen Fall auch, wenn die Druckverschiedenheiten endlich sind. Die Stelle, wo das Gleichgewicht gest\"ort ist, zerf\"allt sich ebenfalls nach Verlauf einer endlichen Zeit in zwei nach entgegengesetzten Richtungen fortschreitende Wellen. In diesen ist die Geschwindigkeit, in der Fortpflanzungsrichtung gemessen, eine bestimmte Function $\int\sqrt{\varphi'(\varrho)}\,d\log\varrho$ der Dichtigkeit, wobei die Integrationsconstante in beiden verschieden sein kann; in jeder ist also mit einem und demselben Werthe der Dichtigkeit stets derselbe Werth der Geschwindigkeit verbunden, und zwar mit einem gr\"osseren Werthe ein algebraisch gr\"osserer Werth der Geschwindigkeit. Beide Werthe r\"ucken mit constanter Geschwindigkeit fort. Ihre Fortpflanzungsgeschwindigkeit im Gase ist $\sqrt{\varphi'(\varrho)}$, im Raume aber um die in der Fortpflanzungsrichtung gemessene Geschwindigkeit des Gases gr\"osser.

Unter der in der Wirklichkeit zutreffenden Voraussetzung, dass $\varphi'(\varrho)$ bei wachsendem $\varrho$ nicht abnimmt, r\"ucken daher gr\"ossere Dichtigkeiten mit gr\"osserer Geschwindigkeit fort, und hieraus folgt, dass die Verd\"unnungswellen, d.~h.~die Theile der Welle, in denen die Dichtigkeit in der Fortpflanzungsrichtung w\"achst, der Zeit proportional an Breite zunehmen, die Verdichtungswellen aber ebenso an Breite abnehmen, und schliesslich in Verdichtungsst\"osse \"ubergehen m\"ussen. Die Gesetze, welche vor der Scheidung beider Wellen oder bei einer \"uber den ganzen Raum sich erstreckenden Gleichgewichtsst\"orung gelten, so wie die Gesetze f\"ur das Fortschreiten von Verdichtungsst\"ossen, k\"onnen hier, weil dazu gr\"ossere Formeln erforderlich w\"aren, nicht angegeben werden.

In akustischer Beziehung liefert demnach diese Untersuchung das Resultat, dass in den F\"allen, wo die Druckverschiedenheiten nicht als unendlich klein betrachtet werden k\"onnen, eine Aenderung der Form der Schallwellen, also des Klanges, w\"ahrend der Fortpflanzung eintritt. Eine Pr\"ufung dieses Resultats durch Versuche scheint aber trotz der Fortschritte, welche in der Analyse des Klanges in neuester Zeit durch Helmholtz u.~A. gemacht worden sind, sehr schwer zu sein; denn in geringeren Entfernungen ist eine Aenderung des Klanges nicht merklich, und bei gr\"osseren Entfernungen wird es schwer sein, die mannigfachen Ursachen, welche den Klang modificiren k\"onnen, zu sondern. An eine Anwendung auf die Meteorologie ist wohl nicht zu denken, da die hier untersuchten Bewegungen der Luft solche Bewegungen sind, die sich mit der Schallgeschwindigkeit fortpflanzen, die Str\"omungen in der Atmosph\"are aber allem Anschein nach mit viel geringerer Geschwindigkeit fortschreiten.

\newpage
%===================================================================
% ANMERKUNGEN ZU PAPER VIII
%===================================================================
\begin{center}
\Large\textbf{Anmerkungen.}
\end{center}

\smallskip
\noindent (1) (Zu Seite~172.) Zum leichteren Verst\"andniss dieser sehr kurzen Andeutungen sollen folgende Bemerkungen dienen.

Die durch die Bedingungen (3) des Art.~8 definirte Function $v$ enth\"alt nichts mehr, was von den besonderen f\"ur die Function $w$ geltenden Grenzbedingungen abh\"angt. Kann man aber unter einer speciellen Voraussetzung $w$ bestimmen, so ergiebt umgekehrt die Formel (4) eine auch f\"ur alle anderen F\"alle g\"ultige Bestimmung von $v$. Es kommt dabei nicht darauf an, dass die gew\"ahlten Grenzbedingungen von $w$ auf das mechanische Problem passen; man kann $w$ und seine beiden Differentialquotienten auf der Curve $c$ beliebig w\"ahlen.

Wir nehmen also f\"ur $c$ eine Curve, in der $\vartheta$ einen constanten Werth hat, und nehmen auf dieser Linie $w = 0$, $\frac{\partial w}{\partial r}$ aber gleich einer beliebigen Function von $u$.

Die Differentialgleichung (1) des Art.~8 wird nun durch Einf\"uhrung von $u$ und $\vartheta$ in folgende transformirt:
\begin{equation*}
\frac{\partial^{2}w}{\partial\vartheta^{2}} - \frac{\partial^{2}w}{\partial u^{2}} - 2m\frac{\partial w}{\partial\vartheta} = 0\tag{1}
\end{equation*}
und ergiebt als particul\"are L\"osungen
\[
\psi\cos\mu u,\qquad \psi\sin\mu u,
\]
wenn $\mu$ ein willk\"urlicher Parameter ist, und $\psi$ als Function von $\vartheta$ durch die Differentialgleichung
\begin{equation*}
\psi'' - 2m\psi' + \mu^{2}\psi = 0\tag{2}
\end{equation*}
bestimmt wird. Sind $\psi_{1}$ und $\psi_{2}$ zwei particul\"are L\"osungen dieser Gleichung, so ist nach einem bekannten Satz
\[
\psi_{1}\psi_{2}' - \psi_{2}\psi_{1}' = \text{const.},
\]
und da (nach Art.~2) $2m = \frac{d\log a}{d\vartheta}$ ist, so kann man die particul\"aren L\"osungen $\psi_{1}$, $\psi_{2}$ so bestimmen, dass
\[
\psi_{1}\psi_{2}' - \psi_{2}\psi_{1}' = \frac{da}{d\vartheta},
\]
wie im Text verlangt wird.

Bezeichnen wir den constanten Werth, den $\vartheta$ auf der Linie $c$ hat, ohne Index, mit $\vartheta'$, einen beliebigen Werth von $\vartheta$, so erhalten wir ein f\"ur $\vartheta' = \vartheta$ der Grenzbedingung $w = 0$ gen\"ugendes $w$ in dem Ausdruck
\begin{equation*}
w_{\vartheta',u'} = \frac{1}{2}\int_{0}^{\infty}(A\cos\mu u' + B\sin\mu u')(\psi_{1}(\vartheta)\psi_{2}(\vartheta') - \psi_{2}(\vartheta)\psi_{1}(\vartheta'))\,d\mu,\tag{4}
\end{equation*}
wenn $A$, $B$ willk\"urliche Functionen von $\mu$ sind. Differentiirt man nach $\vartheta'$ und setzt dann $\vartheta' = \vartheta$, so folgt
\[
\Bigl(\frac{\partial w}{\partial\vartheta'}\Bigr)_{\vartheta'=\vartheta} = \frac{1}{2}\int_{0}^{\infty}(A\cos\mu u' + B\sin\mu u')\,d\mu.
\]

Nun ist nach dem Fourier'schen Lehrsatz, wenn $\varphi(u)$ eine willk\"urliche Function von $u$ ist,
\[
\varphi(u') = \frac{1}{\pi}\int_{0}^{\infty}d\mu\int_{-\infty}^{\infty}du\,\varphi(u)\cos\mu(u - u'),
\]
und man erh\"alt also
\[
A\cos\mu u' + B\sin\mu u' = \frac{1}{\pi}\frac{d\vartheta}{da}\int du\,\frac{\partial w}{\partial\vartheta}\cos\mu(u - u'),
\]
\[
w_{\vartheta',u'} = \frac{1}{2\pi}\int_{-\infty}^{+\infty}du\int_{0}^{\infty}d\mu\,\cos\mu(u - u')\frac{\psi_{1}(\vartheta)\psi_{2}(\vartheta') - \psi_{2}(\vartheta)\psi_{1}(\vartheta')}{da/d\vartheta}.
\]

Setzt man dies in (4) ein, so erh\"alt man
\begin{equation*}
w_{\vartheta',u'} = \frac{1}{2\pi}\int_{-\infty}^{+\infty}du\int_{0}^{\infty}d\mu\,\cos\mu(u - u')\frac{\psi_{1}(\vartheta)\psi_{2}(\vartheta') - \psi_{2}(\vartheta)\psi_{1}(\vartheta')}{da/d\vartheta}.\tag{5}
\end{equation*}

Nun giebt die Formel (4) des Art.~8 unter den gegenw\"artigen Voraussetzungen
\begin{equation*}
w_{\vartheta',u'} = \frac{w_{\vartheta'+u'-\vartheta} + w_{\vartheta'-u'+\vartheta}}{2}.\tag{6}
\end{equation*}

Nimmt man, was freisteht, $\frac{\partial w}{\partial\vartheta}$ nur in einer innerhalb des Intervalles $u' - \vartheta + \vartheta'$ und $u' + \vartheta - \vartheta'$ gelegenen Strecke von $0$ verschieden, ausserhalb dieser Strecke $= 0$ an, so ergiebt die Vergleichung von (5) und (6) unmittelbar die Formel des Textes
\[
v = \int_{0}^{\infty}\cos\mu(u - u')\frac{\psi_{1}(\vartheta')\psi_{2}(\vartheta) - \psi_{2}(\vartheta')\psi_{1}(\vartheta)}{\psi_{1}'(\vartheta) - \psi_{2}'(\vartheta)}\,d\mu.
\]

Die Differentialgleichung (2) wird unter Voraussetzung des Poisson'schen Gesetzes
\[
\psi'' - \frac{2\lambda}{\vartheta}\psi' + \mu^{2}\psi = 0.
\]
Integrirt man sie durch Potenzreihen, so erh\"alt man zwei particul\"are Integrale in der Form
\[
\psi = \vartheta^{n}\sum\frac{(-\mu^{2}\vartheta^{2})^{n}}{n(n)(n + \frac{2\lambda}{k-1})},
\]
wenn man $n$ das eine Mal von Null, das andere Mal von $1 - \frac{2\lambda}{k-1}$ in Stufen von einer Einheit ins Unendliche wachsen l\"asst.

% --- page 181 ---
Durch Anwendung der Formel
\[
\frac{1}{\Gamma(n)(n + \frac{2\lambda}{k-1})} = \frac{1}{2\pi i}\int x^{-n-1}(1 - x)^{n + \frac{2\lambda}{k-1} - 1}\,dx
\]
kann man diese Reihen summiren und erh\"alt unter Weglassung des Factors $-\frac{1}{2\mu}$: f\"ur die beiden particul\"aren Integrale
\[
\psi = \int e^{\mu\vartheta i(t - \frac{1}{t})}t^{\frac{2\lambda}{k-1} - 1}\,dt,
\]
worin die Integrationen auf complexem Wege von $-\infty$ nach $-\infty$ um den Nullpunkt herum zu nehmen sind. Diese Formeln finden sich wenigstens angedeutet in Riemann's Papieren.

\newpage
%===================================================================
% PAPER X: Ein Beitrag zu den Untersuchungen über die Bewegung
%===================================================================
\begin{center}
\Large\textbf{X.}
\end{center}
\vspace{-0.5em}
\begin{center}
\Large\textbf{Ein Beitrag zu den Untersuchungen \"uber die Bewegung eines}\\[0.3em]
\Large\textbf{fl\"ussigen gleichartigen Ellipsoides.}
\end{center}
\vspace{0.3em}
\begin{center}
\small (Aus dem neunten Bande der Abhandlungen der K\"oniglichen Gesellschaft der\\
\small Wissenschaften zu G\"ottingen. 1861.)
\end{center}

\bigskip

F\"ur die Untersuchungen \"uber die Bewegung eines gleichartigen fl\"ussigen Ellipsoides, dessen Elemente sich nach dem Gesetze der Schwere anziehen, hat Dirichlet durch seine letzte von Dedekind herausgegebene Arbeit auf \"uberraschende Weise eine neue Bahn gebrochen. Die Verfolgung dieser sch\"onen Entdeckung hat f\"ur den Mathematiker ihren besondern Reiz, ganz abgesehen von der Frage nach den Gr\"unden der Gestalt der Himmelsk\"orper, durch welche diese Untersuchungen veranlasst worden sind. Dirichlet selbst hat die L\"osung der von ihm behandelten Aufgabe nur in den einfachsten F\"allen vollst\"andig durchgef\"uhrt. F\"ur die weitere Ausf\"uhrung der Untersuchung ist es zweckm\"assig, den Differentialgleichungen f\"ur die Bewegung der fl\"ussigen Masse eine von dem gew\"ahlten Anfangszeitpunkte unabh\"angige Form zu geben, was z.~B.~dadurch geschehen kann, dass man die Gesetze aufsucht, nach welchen die Gr\"osse der Hauptaxen des Ellipsoides und die relative Bewegung der fl\"ussigen Masse gegen dieselben sich \"andert. Indem wir hier die Aufgabe in dieser Weise behandeln, werden wir zwar die Dirichlet'sche Abhandlung voraussetzen, m\"ussen aber dabei zur Vermeidung von Irrungen gleich bevorworten, dass es nicht m\"oglich gewesen ist, die dort gebrauchten Zeichen unver\"andert beizubehalten.

\begin{center}
\textbf{1.}
\end{center}

Wir bezeichnen durch $a$, $b$, $c$ die Hauptaxen des Ellipsoides zur Zeit $t$, ferner durch $x$, $y$, $z$ die Coordinaten eines Elements der fl\"ussigen Masse zur Zeit $t$ und die Anfangswerthe dieser Gr\"ossen durch Anh\"angung des Index $0$ und nehmen an, dass f\"ur die Anfangszeit die Hauptaxen des Ellipsoides mit den Coordinatenaxen zusammenfallen.

Den Ausgangspunkt f\"ur die Untersuchung Dirichlet's bildet bekanntlich die Bemerkung, dass man den Differentialgleichungen f\"ur die Bewegung der Fl\"ussigkeitstheile gen\"ugen kann, wenn man die Coordinaten $x$, $y$, $z$ linearen Ausdr\"ucken von ihren Anfangswerthen gleichsetzt, in denen die Coefficienten blosse Functionen der Zeit sind. Diese Ausdr\"ucke setzen wir in die Form

% --- page 183 ---
\begin{align*}
x &= \frac{a}{a_{0}}x_{0} + \frac{a_{1}}{b_{0}}y_{0} + \frac{a_{2}}{c_{0}}z_{0}\\
y &= \frac{b}{a_{0}}x_{0} + \frac{b_{1}}{b_{0}}y_{0} + \frac{b_{2}}{c_{0}}z_{0}\\
z &= \frac{c}{a_{0}}x_{0} + \frac{c_{1}}{b_{0}}y_{0} + \frac{c_{2}}{c_{0}}z_{0}\tag{1}
\end{align*}

Bezeichnet man nun durch $\xi$, $\eta$, $\zeta$ die Coordinaten des Punktes $(x,y,z)$ in Bezug auf ein bewegliches Coordinatensystem, dessen Axen in jedem Augenblicke mit den Hauptaxen des Ellipsoides zusammenfallen, so sind bekanntlich $\xi$, $\eta$, $\zeta$ gleich linearen Ausdr\"ucken von $x$, $y$, $z$
\begin{align*}
\xi &= \alpha x + \beta y + \gamma z\\
\eta &= \alpha' x + \beta' y + \gamma' z\\
\zeta &= \alpha'' x + \beta'' y + \gamma'' z,\tag{2}
\end{align*}
worin die Coefficienten die Cosinus der Winkel sind, welche die Axen des einen Systems mit den Axen des andern bilden, $\alpha = \cos\xi x$, $\beta = \cos\xi y$ etc., und zwischen diesen Coefficienten finden sechs Bedingungsgleichungen statt, welche sich daraus herleiten lassen, dass durch die Substitution dieser Ausdr\"ucke
\[
\xi^{2} + \eta^{2} + \zeta^{2} = x^{2} + y^{2} + z^{2}
\]
werden muss.

Da die Oberfl\"ache stets von denselben Fl\"ussigkeitstheilchen gebildet wird, so muss
\[
\frac{\xi^{2}}{a^{2}} + \frac{\eta^{2}}{b^{2}} + \frac{\zeta^{2}}{c^{2}} = \frac{x_{0}^{2}}{a_{0}^{2}} + \frac{y_{0}^{2}}{b_{0}^{2}} + \frac{z_{0}^{2}}{c_{0}^{2}}
\]
sein; setzt man also
\begin{align*}
\frac{\xi}{a} &= \alpha_{1}\frac{x_{0}}{a_{0}} + \beta_{1}\frac{y_{0}}{b_{0}} + \gamma_{1}\frac{z_{0}}{c_{0}}\\
\frac{\eta}{b} &= \alpha_{1}'\frac{x_{0}}{a_{0}} + \beta_{1}'\frac{y_{0}}{b_{0}} + \gamma_{1}'\frac{z_{0}}{c_{0}}\\
\frac{\zeta}{c} &= \alpha_{1}''\frac{x_{0}}{a_{0}} + \beta_{1}''\frac{y_{0}}{b_{0}} + \gamma_{1}''\frac{z_{0}}{c_{0}},\tag{3}
\end{align*}
d.~h.~bezeichnet man in den Ausdr\"ucken von $\frac{\xi}{a}$, $\frac{\eta}{b}$, $\frac{\zeta}{c}$ durch $\frac{x_{0}}{a_{0}}$, $\frac{y_{0}}{b_{0}}$, $\frac{z_{0}}{c_{0}}$, welche man durch Einsetzung der Werthe (1) in die Gleichungen (2) erh\"alt, die Coefficienten durch $\alpha_{1}$, $\beta_{1}$, \dots, $\gamma_{1}''$, so bilden diese Gr\"ossen $\alpha_{1}$, $\beta_{1}$, \dots, $\gamma_{1}''$ ebenfalls die Coefficienten einer orthogonalen Coordinatentransformation: sie k\"onnen betrachtet werden als die Cosinus der Winkel, welche die Axen eines beweglichen Coordinatensystems der $\xi_{1}$, $\eta_{1}$, $\zeta_{1}$ mit den Axen des festen Coordinatensystems der $x$, $y$, $z$ bilden. Dr\"uckt man die Gr\"ossen $x$, $y$, $z$ mit H\"ulfe der Gleichungen (2) und (3) in $\frac{x_{0}}{a_{0}}$, $\frac{y_{0}}{b_{0}}$, $\frac{z_{0}}{c_{0}}$ aus, so ergiebt sich
\begin{align*}
l &= a\alpha\alpha_{1} + b\alpha'\alpha_{1}' + c\alpha''\alpha_{1}''\\
m &= a\alpha\beta_{1} + b\alpha'\beta_{1}' + c\alpha''\beta_{1}''\\
n &= a\alpha\gamma_{1} + b\alpha'\gamma_{1}' + c\alpha''\gamma_{1}''\\
l' &= a\beta\alpha_{1} + b\beta'\alpha_{1}' + c\beta''\alpha_{1}''\\
m' &= a\beta\beta_{1} + b\beta'\beta_{1}' + c\beta''\beta_{1}''\\
n' &= a\beta\gamma_{1} + b\beta'\gamma_{1}' + c\beta''\gamma_{1}''\\
l'' &= a\gamma\alpha_{1} + b\gamma'\alpha_{1}' + c\gamma''\alpha_{1}''\\
m'' &= a\gamma\beta_{1} + b\gamma'\beta_{1}' + c\gamma''\beta_{1}''\\
n'' &= a\gamma\gamma_{1} + b\gamma'\gamma_{1}' + c\gamma''\gamma_{1}''.\tag{4}
\end{align*}

Wir k\"onnen daher die Lage der Fl\"ussigkeitstheilchen oder die Werthe der Gr\"ossen $l$, $m$, \dots, $n''$ zur Zeit $t$ als abh\"angig betrachten von den Gr\"ossen $a$, $b$, $c$ und der Lage zweier beweglichen Coordinatensysteme und k\"onnen zugleich bemerken, dass durch Vertauschung dieser beiden Coordinatensysteme in dem Systeme der Gr\"ossen $l$, $m$, $n$ die Horizontalreihen mit den Verticalreihen vertauscht werden, also $l$, $m'$, $n''$ unge\"andert bleiben, w\"ahrend von den Gr\"ossen $m$ und $l'$, $n$ und $l''$, $n'$ und $m''$ jede in die andere \"ubergeht. Es wird nun unser n\"achstes Gesch\"aft sein, die Differentialgleichungen f\"ur die Ver\"anderungen der Hauptaxen und die Bewegung dieser beiden Coordinatensysteme aus den in der Dirichlet'schen Abhandlung (\S.~1, 1) angegebenen Grundgleichungen f\"ur die Bewegung der Fl\"ussigkeitstheilchen abzuleiten.

\begin{center}
\textbf{2.}
\end{center}

Offenbar ist es erlaubt, in jenen Gleichungen statt der Derivirten nach den Anfangswerthen der Gr\"ossen $x$, $y$, $z$, welche dort durch $a$, $b$, $c$ bezeichnet sind, die Derivirten nach den Gr\"ossen $\xi$, $\eta$, $\zeta$ zu setzen; denn die hierdurch gebildeten Gleichungen lassen sich als Aggregate von jenen darstellen und umgekehrt. Wir erhalten dadurch, wenn wir f\"ur $\frac{\partial x}{\partial\xi}$, $\frac{\partial y}{\partial\eta}$, \dots, $\frac{\partial z}{\partial\zeta}$ ihre Werthe einsetzen,

% --- page 185 ---
\begin{align*}
\frac{\partial^{2}x}{\partial t^{2}}\xi + \frac{\partial^{2}y}{\partial t^{2}}\eta + \frac{\partial^{2}z}{\partial t^{2}}\zeta &= \varepsilon\frac{\partial V}{\partial x}\xi + \varepsilon\frac{\partial V}{\partial y}\eta + \varepsilon\frac{\partial V}{\partial z}\zeta - \frac{1}{\varrho}\frac{\partial p}{\partial\xi},\\
\frac{\partial^{2}x}{\partial t^{2}}\xi' + \frac{\partial^{2}y}{\partial t^{2}}\eta' + \frac{\partial^{2}z}{\partial t^{2}}\zeta' &= \varepsilon\frac{\partial V}{\partial x}\xi' + \varepsilon\frac{\partial V}{\partial y}\eta' + \varepsilon\frac{\partial V}{\partial z}\zeta' - \frac{1}{\varrho}\frac{\partial p}{\partial\eta},\\
\frac{\partial^{2}x}{\partial t^{2}}\xi'' + \frac{\partial^{2}y}{\partial t^{2}}\eta'' + \frac{\partial^{2}z}{\partial t^{2}}\zeta'' &= \varepsilon\frac{\partial V}{\partial x}\xi'' + \varepsilon\frac{\partial V}{\partial y}\eta'' + \varepsilon\frac{\partial V}{\partial z}\zeta'' - \frac{1}{\varrho}\frac{\partial p}{\partial\zeta},\tag{1}
\end{align*}
worin $V$ das Potential, $p$ den Druck im Punkte $x$, $y$, $z$ zur Zeit $t$ und $\varepsilon$ die Constante bezeichnet, welche die Anziehung zwischen zwei Masseneinheiten in der Entfernungseinheit ausdr\"uckt.

Es handelt sich nun zun\"achst darum, die Gr\"ossen links vom Gleichheitszeichen in die Form linearer Functionen von den Gr\"ossen $\xi$, $\eta$, $\zeta$ zu setzen, wozu einige Vorbereitungen n\"othig sind.

Durch Differentiation der Gleichungen (2) erh\"alt man, wenn man zur Abk\"urzung
\begin{align*}
\frac{dx}{dt} &= p\,y + q\,z,\quad \frac{dy}{dt} = r\,x + s\,z,\quad \frac{dz}{dt} = t\,x + u\,y\quad\text{setzt,}\\
\frac{d\xi}{dt} &= \frac{d\alpha}{dt}x + \frac{d\beta}{dt}y + \frac{d\gamma}{dt}z,\\
\frac{d\eta}{dt} &= \frac{d\alpha'}{dt}x + \frac{d\beta'}{dt}y + \frac{d\gamma'}{dt}z,\\
\frac{d\zeta}{dt} &= \frac{d\alpha''}{dt}x + \frac{d\beta''}{dt}y + \frac{d\gamma''}{dt}z,\tag{2}
\end{align*}
und wenn man hierin $x$, $y$, $z$ wieder durch $\xi$, $\eta$, $\zeta$ ausdr\"uckt
\begin{align*}
\frac{d\xi}{dt} &= (q\beta\gamma + r\beta'\gamma' + s\beta''\gamma'')\xi + (q\gamma\alpha + r\gamma'\alpha' + s\gamma''\alpha'')\eta\\
&\quad + (q\alpha\beta + r\alpha'\beta' + s\alpha''\beta'')\zeta + p',\tag{3}
\end{align*}
und \"ahnliche Gleichungen f\"ur $\frac{d\eta}{dt}$, $\frac{d\zeta}{dt}$.

Nun giebt aber die Differentiation der bekannten Gleichungen $\alpha^{2} + \beta^{2} + \gamma^{2} = 1$, $\alpha\alpha' + \beta\beta' + \gamma\gamma' = 0$, etc.

% --- page 186 ---
\begin{align*}
\frac{d\alpha}{dt} &= \alpha''\frac{d\xi}{dt} + \beta''\frac{d\eta}{dt} + \gamma''\frac{d\zeta}{dt},\\
\frac{d\alpha'}{dt} &= \alpha'''\frac{d\xi}{dt} + \beta'''\frac{d\eta}{dt} + \gamma'''\frac{d\zeta}{dt},\\
\frac{d\alpha''}{dt} &= \alpha''''\frac{d\xi}{dt} + \beta''''\frac{d\eta}{dt} + \gamma''''\frac{d\zeta}{dt}\tag{3}
\end{align*}
und es wird folglich, wenn man diese letzteren drei Gr\"ossen durch $p$, $q$, $r$ bezeichnet,
\begin{align*}
\frac{d\xi}{dt} &= r\eta - q\zeta,\\
\frac{d\eta}{dt} &= p\zeta - r\xi,\\
\frac{d\zeta}{dt} &= q\xi - p\eta.\tag{4}
\end{align*}

Durch ein ganz \"ahnliches Verfahren ergiebt sich aus den Gleichungen (2)
\begin{align*}
\xi' \frac{d\xi}{dt} + \eta' \frac{d\eta}{dt} + \zeta' \frac{d\zeta}{dt} &= r'\xi - q'\eta,\quad\text{etc.}\tag{5}
\end{align*}
und aus den Gleichungen Art.~1, (3), wenn $p_{1}$, $q_{1}$, $r_{1}$ die Gr\"ossen bezeichnen, welche von den Functionen $\alpha_{1}$, $\beta_{1}$, \dots, $\gamma_{1}''$ ebenso abh\"angen, wie die Gr\"ossen $p$, $q$, $r$ von den Functionen $\alpha$, $\beta$, \dots, $\gamma''$,
\begin{align*}
\frac{1}{a}\frac{da}{dt} &= p_{1},\quad \frac{1}{b}\frac{db}{dt} = q_{1},\quad \frac{1}{c}\frac{dc}{dt} = r_{1}.\tag{6}
\end{align*}

Setzt man die Werthe $\xi'$, $\eta'$, $\zeta'$ aus (6) in (4) ein, so erh\"alt man

Was die geometrische Bedeutung dieser Gr\"ossen betrifft, so sind, wie leicht ersichtlich ist, $\xi'$, $\eta'$, $\zeta'$ die Geschwindigkeitscomponenten des Punktes $x$, $y$, $z$ der fl\"ussigen Masse parallel den Axen $\xi$, $\eta$, $\zeta$; $\frac{d\xi}{dt}$, $\frac{d\eta}{dt}$, $\frac{d\zeta}{dt}$ die ebenso zerlegten relativen Geschwindigkeiten gegen das Coordinatensystem der $\xi$, $\eta$, $\zeta$; ferner in den Gleichungen (1) die Gr\"ossen auf der linken Seite die Beschleunigungen und die auf der rechten die beschleunigenden Kr\"afte parallel diesen Axen; endlich sind $p$, $q$, $r$ die augenblicklichen Rotationen des Coordinatensystems der $\xi$, $\eta$, $\zeta$ um seine Axen und $p_{1}$, $q_{1}$, $r_{1}$ haben dieselbe Bedeutung f\"ur das Coordinatensystem der $\xi_{1}$, $\eta_{1}$, $\zeta_{1}$.

\begin{center}
\textbf{3.}
\end{center}

Wenn man nun die Werthe der Gr\"ossen $\xi'$, $\eta'$, $\zeta'$ aus (7) in die Gleichungen (5) substituirt und mit H\"ulfe der Gleichungen (6) die Derivirten von $\frac{\xi}{a}$, $\frac{\eta}{b}$, $\frac{\zeta}{c}$ wieder durch die Gr\"ossen $\xi$, $\eta$, $\zeta$ ausdr\"uckt, so nehmen die Gr\"ossen auf der linken Seite der Gleichungen (1) die Form linearer Ausdr\"ucke von den Gr\"ossen $\xi$, $\eta$, $\zeta$ an. Auf der rechten Seite hat $V$ die Form
\[
V = H - A\xi^{2} - B\eta^{2} - C\zeta^{2},
\]
worin $H$, $A$, $B$, $C$ auf bekannte Weise von den Gr\"ossen $a$, $b$, $c$ abh\"angen; und man gen\"ugt ihnen daher, wenn an der Oberfl\"ache der Druck den constanten Werth $0$ hat, indem man
\[
p = \varrho\bigl(\Theta - 2\varepsilon(A\xi^{2} + B\eta^{2} + C\zeta^{2})\bigr)
\]
setzt und die zehn Functionen der Zeit $a$, $b$, $c$; $p$, $q$, $r$; $p_{1}$, $q_{1}$, $r_{1}$ und $\Theta$ so bestimmt, dass die neun Coefficienten der Gr\"ossen $\xi$, $\eta$, $\zeta$ auf beiden Seiten einander gleich werden und zugleich die aus der Incompressibilit\"at folgende Bedingungsgleichung $abc = a_{0}b_{0}c_{0}$ befriedigt wird.

% --- page 188 ---
Durch Gleichsetzung der Coefficienten von $\frac{\xi^{2}}{a^{2}}$, $\frac{\eta^{2}}{b^{2}}$ in der ersten und von $\frac{\zeta^{2}}{c^{2}}$ in der zweiten Gleichung ergiebt sich
\begin{align*}
\frac{d^{2}a}{dt^{2}} + 2brr_{1} + 2cqq_{1} - a(r^{2} + r_{1}^{2} + q^{2} + q_{1}^{2}) &= 2\varepsilon aA,\\
\frac{d^{2}b}{dt^{2}} + 2cpr_{1} + 2app_{1} - b(p^{2} + p_{1}^{2} + r^{2} + r_{1}^{2}) &= 2\varepsilon bB,\\
\frac{d^{2}c}{dt^{2}} + 2aqp_{1} + 2bqq_{1} - c(q^{2} + q_{1}^{2} + p^{2} + p_{1}^{2}) &= 2\varepsilon cC.\tag{a}
\end{align*}
Aus diesen Gleichungen erh\"alt man die sechs \"ubrigen durch cyclische Versetzung der Axen, oder auch durch beliebige Vertauschungen, wenn man nur dabei beachtet, dass durch Vertauschung zweier Axen nicht bloss die ihnen entsprechenden Gr\"ossen vertauscht werden, sondern zugleich die sechs Gr\"ossen $p$, $q$, \dots, $r_{1}$ ihr Zeichen \"andern.

Man kann diesen Gleichungen eine f\"ur die weitere Untersuchung bequemere Form geben, wenn man statt der Gr\"ossen $p$, $p_{1}$, $q$, $q_{1}$, $r$, $r_{1}$ ihre halben Summen und Differenzen
\[
\frac{p + p_{1}}{2},\quad \frac{q + q_{1}}{2},\quad \frac{r + r_{1}}{2},\quad \frac{p - p_{1}}{2},\quad \frac{q - q_{1}}{2},\quad \frac{r - r_{1}}{2}
\]
als unbekannte Functionen einf\"uhrt.

Dadurch wird das System von Gleichungen, welchen die zehn unbekannten Functionen der Zeit gen\"ugen m\"ussen
\begin{align*}
&(b - c)u^{2} + (b + c)u'^{2} + (a - b)w^{2} + (a + b)w'^{2} + \cdots = \varepsilon aA,\\
&(b - a)w^{2} + (b + a)w'^{2} + (b - c)u^{2} + (b + c)u'^{2} + \cdots = \varepsilon bB - \cdots,\\
&(b - c)\frac{dh}{dt} + 2\frac{db}{dt}\frac{dc}{dt}w + (b + c - 2a)vw + (b + c + 2a)v'w' = 0,\\
&(b + c)\frac{dh}{dt} + 2\frac{db}{dt}\frac{dc}{dt}w' + (b - c + 2a)vw' + (b - c - 2a)v'w = 0,\\
&(c - a)\frac{dg}{dt} + 2\frac{dc}{dt}\frac{da}{dt}v + (c + a - 2b)wu + (c + a + 2b)w'u' = 0,\\
&(c + a)\frac{dg}{dt} + 2\frac{dc}{dt}\frac{da}{dt}v' + (c - a + 2b)wu' + (c - a - 2b)w'u = 0,\\
&(a - b)\frac{df}{dt} + 2\frac{da}{dt}\frac{db}{dt}w + (a + b - 2c)uv + (a + b + 2c)u'v' = 0,\\
&(a + b)\frac{df}{dt} + 2\frac{da}{dt}\frac{db}{dt}w' + (a - b + 2c)uv' + (a - b - 2c)u'v = 0,\\
&abc = a_{0}b_{0}c_{0}.
\end{align*}

Die Werthe von $A$, $B$, $C$ ergeben sich aus dem bekannten Ausdrucke f\"ur $V$
\begin{equation*}
V = \pi abc\int_{0}^{\infty}\frac{ds}{\Delta}\Bigl(1 - \frac{\xi^{2}}{a^{2} + s} - \frac{\eta^{2}}{b^{2} + s} - \frac{\zeta^{2}}{c^{2} + s}\Bigr),\tag{$\alpha$}
\end{equation*}
worin
\[
\Delta = \sqrt{(a^{2} + s)(b^{2} + s)(c^{2} + s)}.
\]

Nach ausgef\"uhrter Integration dieser Differentialgleichungen hat man noch, um die Functionen $\alpha$, $\beta$, \dots, $\gamma''$ zu bestimmen, die allgemeine L\"osung $\theta$, $\theta'$, $\theta''$ der Differentialgleichungen
\[
\frac{d\theta}{dt} = -r\theta' + q\theta'',\quad \frac{d\theta'}{dt} = -p\theta'' + r\theta,\quad \frac{d\theta''}{dt} = -q\theta + p\theta'\tag{$\beta$}
\]
zu suchen, --- von welchen, wie aus Art.~2, (3) hervorgeht, $\alpha$, $\alpha'$, $\alpha''$; $\beta$, $\beta'$, $\beta''$; $\gamma$, $\gamma'$, $\gamma''$ die drei particul\"aren Aufl\"osungen sind, die f\"ur $t = 0$ die Werthe $1$, $0$, $0$; $0$, $1$, $0$; $0$, $0$, $1$ annehmen, --- und zur Bestimmung der Functionen $\alpha_{1}$, $\beta_{1}$, \dots, $\gamma_{1}''$ die allgemeine L\"osung der simultanen Differentialgleichungen
\[
\frac{d\theta}{dt} = -r_{1}\theta' + q_{1}\theta'',\quad \frac{d\theta'}{dt} = -p_{1}\theta'' + r_{1}\theta,\quad \frac{d\theta''}{dt} = -q_{1}\theta + p_{1}\theta'.\tag{$\gamma$}
\]

Es fragt sich nun, welche H\"ulfsmittel f\"ur die Integration dieser Differentialgleichungen ($\alpha$), ($\beta$), ($\gamma$) die allgemeinen hydrodynamischen Principien darbieten, aus denen Dirichlet sieben Integrale erster Ordnung der durch die Functionen $l$, $m$, \dots, $n''$ zu erf\"ullenden Differentialgleichungen (\S.~1, ($\alpha$)) sch\"opfte. Die aus ihnen fliessenden Gleichungen lassen sich mit H\"ulfe der oben f\"ur $\xi'$, $\eta'$, $\zeta'$ gegebenen Ausdr\"ucke leicht herleiten.

Der Satz von der Erhaltung der Fl\"achen giebt
\begin{align*}
(b - c)fu + (b + c)f'u' &= g = \alpha g^{0} + \beta h^{0} + \gamma k^{0}\\
(c - a)gv + (c + a)g'v' &= h = \alpha' g^{0} + \beta' h^{0} + \gamma' k^{0}\\
(a - b)hw + (a + b)h'w' &= k = \alpha'' g^{0} + \beta'' h^{0} + \gamma'' k^{0},\tag{1}
\end{align*}
worin die Constanten $g^{0}$, $h^{0}$, $k^{0}$, die Anfangswerthe von $g$, $h$, $k$, mit den Constanten $G$, $H$, $K$ in der Abhandlung von Dirichlet \"ubereinkommen; er liefert also das aus den sechs letzten Differentialgleichungen ($\alpha$) leicht zu best\"atigende Resultat, dass $\theta = g$, $\theta' = h$, $\theta'' = k$ eine L\"osung der Differentialgleichungen ($\beta$) ist.

Aus dem Helmholtz'schen Princip der Erhaltung der Rotation folgen die Gleichungen

% --- page 190 ---
\begin{align*}
(b - c)fu - (b + c)f'u' &= g_{1} = \alpha_{1} g_{1}^{0} + \beta_{1} h_{1}^{0} + \gamma_{1} k_{1}^{0}\\
(c - a)gv - (c + a)g'v' &= h_{1} = \alpha_{1}' g_{1}^{0} + \beta_{1}' h_{1}^{0} + \gamma_{1}' k_{1}^{0}\\
(a - b)hw - (a + b)h'w' &= k_{1} = \alpha_{1}'' g_{1}^{0} + \beta_{1}'' h_{1}^{0} + \gamma_{1}'' k_{1}^{0},\tag{2}
\end{align*}
in welchen die Constanten $g_{1}^{0}$, $h_{1}^{0}$, $k_{1}^{0}$ den Gr\"ossen $BC^{0}$, $CA^{0}$, $AB^{0}$ der genannten Abhandlung gleich sind.

Der Satz von der Erhaltung der lebendigen Kraft endlich giebt ein Integral erster Ordnung der Differentialgleichungen ($\alpha$)
\begin{multline*}
\Bigl(\frac{da}{dt}\Bigr)^{2} + \Bigl(\frac{db}{dt}\Bigr)^{2} + \Bigl(\frac{dc}{dt}\Bigr)^{2} + (b - c)^{2}u^{2} + (c - a)^{2}v^{2} + (a - b)^{2}w^{2}\\
+ (b + c)^{2}u'^{2} + (c + a)^{2}v'^{2} + (a + b)^{2}w'^{2} = 2\varepsilon H + \text{const.}\tag{I}
\end{multline*}

Aus den Gleichungen (1) und (2) folgen zun\"achst noch zwei Integrale der Gleichungen ($\alpha$)
\begin{align*}
g^{2} + h^{2} + k^{2} &= \text{const.} = C^{2},\tag{II}\\
g_{1}^{2} + h_{1}^{2} + k_{1}^{2} &= \text{const.} = C'^{2}.\tag{III}
\end{align*}

Ferner lassen sich von den Gleichungen ($\beta$) zwei Integrale
\begin{align*}
\theta^{2} + \theta'^{2} + \theta''^{2} &= \text{const.},\tag{IV}\\
\theta g + \theta' h + \theta'' k &= \text{const.}\tag{V}
\end{align*}
angeben, wodurch ihre Integration allgemein auf eine Quadratur zur\"uckgef\"uhrt wird. Zur Aufstellung ihrer allgemeinen L\"osung ist es jedoch, da sie linear und homogen sind, nur n\"othig, noch zwei von der L\"osung $g$, $h$, $k$ verschiedene particul\"are L\"osungen zu suchen, f\"ur welchen Zweck man die willk\"urlichen Constanten in diesen beiden Integralgleichungen so w\"ahlen kann, dass sich die Rechnung vereinfacht. Giebt man beiden den Werth Null, so hat man
\begin{equation*}
\theta' h + \theta'' k = -\theta g,\tag{3}
\end{equation*}
und ferner erh\"alt man, wenn man diese Gleichung quadrirt und dazu die Gleichung
\[
\theta'^{2} + \theta''^{2} = -\theta^{2}
\]
multiplicirt mit $h^{2} + k^{2}$, addirt
\[
-(\theta' k - \theta'' h)^{2} = \theta^{2}C^{2},
\]
folglich
\begin{equation*}
\theta' k - \theta'' h = \theta_{1} C.\tag{4}
\end{equation*}

Durch Aufl\"osung dieser beiden linearen Gleichungen (3) und (4) findet sich

% --- page 191 ---
\begin{align*}
\theta' &= \frac{-h\theta + k\theta_{1}}{C},\tag{5}\\
\theta'' &= \frac{-k\theta - h\theta_{1}}{C},\tag{6}
\end{align*}
und durch Einsetzung dieser Werthe in die erste der Gleichungen ($\beta$)
\[
C\frac{d\theta}{dt} = \frac{-gk + gh}{C}\theta + \frac{rk + gh}{C}\theta_{1},
\]
\begin{equation*}
\log\theta = \tfrac{1}{2}\log(h^{2} + k^{2}) + \theta_{1}\int\frac{rh - gk}{h^{2} + k^{2}}\,dt + \text{const.}\tag{7}
\end{equation*}

Aus dieser in (5), (6) und (7) enthaltenen L\"osung der Differentialgleichungen ($\beta$) erh\"alt man eine dritte, indem man f\"ur $\theta_{1}$ \"uberall $-\theta_{1}$ setzt, und es ist dann leicht aus den gefundenen drei particul\"aren L\"osungen die Ausdr\"ucke f\"ur die Functionen $\alpha$, $\beta$, \dots, $\gamma''$ zu bilden.

Die geometrische Bedeutung jeder reellen L\"osung der Differentialgleichungen ($\beta$) besteht darin, dass sie, mit einem geeigneten constanten Factor multiplicirt, die Cosinus der Winkel ausdr\"uckt, welche die Axen der $\xi$, $\eta$, $\zeta$ zur Zeit $t$ mit einer festen Linie machen. Diese feste Linie wird f\"ur die erste der drei eben gefundenen L\"osungen durch die Normale auf der unver\"anderlichen Ebene der ganzen bewegten Masse gebildet, f\"ur den reellen und den imagin\"aren Bestandtheil der beiden andern durch zwei in dieser Ebene enthaltene und auf einander senkrechte Linien. Die Cosinus der Winkel zwischen den Axen und jener Normalen sind demnach $\frac{g}{C}$, $\frac{h}{C}$, $\frac{k}{C}$; die Lage der Axen gegen diese Normale ergiebt sich also nach Aufl\"osung der Gleichungen ($\alpha$) ohne weitere Integration und zur vollst\"andigen Bestimmung ihrer Lage gen\"ugt eine einzige Quadratur, z.~B.~die Integration
\begin{equation*}
\int\frac{rh - gk}{h^{2} + k^{2}}\,dt,\tag{$\delta$}
\end{equation*}
welche die Drehung der durch die Normale und die Axe der $\xi$ gehenden Ebene um die Normale giebt.

Ganz Aehnliches gilt von den Differentialgleichungen ($\gamma$). Man kann auf demselben Wege aus den beiden Integralen
\begin{align*}
\theta_{1}^{2} + \theta_{1}'^{2} + \theta_{1}''^{2} &= \text{const.},\tag{VI}\\
\theta_{1}g_{1} + \theta_{1}'h_{1} + \theta_{1}''k_{1} &= \text{const.}\tag{VII}
\end{align*}
ihre allgemeine L\"osung und folglich auch die Werthe der Gr\"ossen $\alpha_{1}$, $\beta_{1}$, \dots, $\gamma_{1}''$ zur Zeit $t$ ableiten, und es wird dabei nur eine Quadratur erforderlich sein. Es ergiebt sich dann schliesslich der Ort eines

% --- page 192 ---
beliebigen Fl\"ussigkeitstheilchens zur Zeit $t$ aus den oben (Art.~1, (1) und (4)) f\"ur die Gr\"ossen $x$, $y$, $z$ und die Functionen $l$, $m$, \dots, $n''$ gegebenen Ausdr\"ucken.

Wir wollen uns jetzt Rechenschaft dar\"uber geben, was durch die Zur\"uckf\"uhrung der Differentialgleichungen zwischen den Functionen $l$, $m$, \dots, $n$ (der Differentialgleichungen ($\alpha$) \S.~1 bei Dirichlet) auf unsere Differentialgleichungen f\"ur das Gesch\"aft der Integration gewonnen ist. Das System der Differentialgleichungen ($\alpha$) ist von der sechszehnten Ordnung, und man kennt von denselben sieben Integrale erster Ordnung, wodurch es auf ein System der neunten Ordnung zur\"uckgef\"uhrt wird. Das System ($\alpha$) ist nur von der zehnten Ordnung, und man kennt von demselben noch drei Integrale erster Ordnung. Durch die hier bewirkte Umformung jener Differentialgleichungen ist also die Ordnung des noch zu integrirenden Systems von Differentialgleichungen um zwei Einheiten erniedrigt, und man hat statt dessen nur schliesslich noch zwei Quadraturen auszuf\"uhren. Diese Umformung leistet also dasselbe, wie die Auffindung von zwei Integralen erster Ordnung.

Wir bemerken indess ausdr\"ucklich, dass hierdurch unsere Form der Differentialgleichungen nur f\"ur die Integration und die wirkliche Bestimmung der Bewegung einen Vorzug erh\"alt. F\"ur die allgemeinsten Untersuchungen \"uber diese Bewegung ist dagegen diese Form der Differentialgleichungen weniger geeignet, nicht bloss, weil ihre Herleitung weniger einfach ist, sondern auch deshalb, weil der Fall der Gleichheit zweier Axen eine besondere Betrachtung erfordert. Bei Gleichheit zweier Axen tritt n\"amlich der besondere Umstand ein, dass die ihnen zu gebende Lage durch die Gestalt der fl\"ussigen Masse nicht v\"ollig bestimmt ist; sie h\"angt dann im Allgemeinen auch von der augenblicklichen Bewegung ab und bleibt nur dann willk\"urlich, wenn diese Bewegung so beschaffen ist, dass die Axen fortw\"ahrend einander gleich bleiben. Die Untersuchung dieses Falles ist zwar immer leicht und bedarf daher keiner weiteren Ausf\"uhrung, kann aber in speciellen F\"allen noch wieder besondere Formen annehmen, und die allgemeinen Untersuchungen, wie z.~B.~der allgemeine Nachweis der M\"oglichkeit der Bewegung (\S.~2 bei Dirichlet), w\"urden daher wegen der Menge von besonders zu behandelnden F\"allen ziemlich weitl\"aufig werden.

Ehe wir zur Behandlung von speciellen F\"allen schreiten, in welchen sich die Differentialgleichungen ($\alpha$) integriren lassen, ist es zweckm\"assig, zu bemerken, dass in einer L\"osung dieser Differentialgleichungen, wie unmittelbar aus der Form dieser Gleichungen hervorgeht, jede Zeichen\"anderung der Functionen $u$, $v$, \dots, $w'$ zul\"assig ist, bei welcher $uv'w'$, $uv'w$, $uvw'$, $u'v'w$ unge\"andert bleiben. Es k\"onnen also erstens die Zeichen der Functionen $u'$, $v$, $w$ gleichzeitig ge\"andert werden, und dadurch werden die Gr\"ossen $\alpha$, $\beta$, \dots, $\gamma''$ mit den Gr\"ossen $\alpha_{1}$, $\beta_{1}$, \dots, $\gamma_{1}''$, also in dem Systeme der Gr\"ossen $l$, $m$, \dots, $n''$ die Horizontalreihen mit den Verticalreihen vertauscht. Zweitens k\"onnen gleichzeitig zwei der Gr\"ossenpaare $u$, $u'$; $v$, $v'$; $w$, $w'$ mit den entgegengesetzten Zeichen versehen werden, und diese Aenderung l\"asst sich auf eine Aenderung in dem Zeichen einer Coordinatenaxe zur\"uckf\"uhren, wobei die Bewegung in eine ihr symmetrisch gleiche \"ubergeht. In dieser Bemerkung ist der von Dedekind gefundene Reciprocit\"atssatz enthalten.

\begin{center}
\textbf{6.}
\end{center}

Wir wollen nun den Fall untersuchen, in welchem eins der Gr\"ossenpaare $u$, $u'$; $v$, $v'$; $w$, $w'$ fortw\"ahrend gleich Null ist, also z.~B. $u = u' = 0$; die geometrische Bedeutung dieser Voraussetzung ist diese, dass die Hauptaxe $a$ stets in der unver\"anderlichen Ebene der ganzen bewegten Masse liegt und die augenblickliche Rotationsaxe auf dieser Hauptaxe senkrecht steht.

Aus den sechs letzten Differentialgleichungen ($\alpha$) folgt sogleich, dass in diesem Falle die Gr\"ossen
\begin{equation*}
(b - c)^{2}v,\quad (b + c)^{2}v',\quad (a - b)^{2}w,\quad (a + b)^{2}w'\tag{$\beta$}
\end{equation*}
constant sind und die Gleichungen
\begin{align*}
(b + c - 2a)vw + (b + c + 2a)v'w' &= 0,\\
(b - c + 2a)vw' + (b - c - 2a)v'w &= 0\tag{$\gamma$}
\end{align*}
stattfinden m\"ussen.

Bei der weiteren Untersuchung ist zu unterscheiden, ob noch ein zweites der drei Gr\"ossenpaare Null ist oder nicht, und wir k\"onnen im Allgemeinen nur noch bemerken, dass in Folge der Gleichungen ($\beta$) die Gr\"ossen $h$, $k$, $h_{1}$, $k_{1}$ constant sind und folglich auch die Winkel zwischen den Hauptaxen und der unver\"anderlichen Ebene der ganzen bewegten Masse, und dass dann ferner aus den Differentialgleichungen ($\beta$) und ($\gamma$) die Verh\"altnissgleichungen
\[
g : h : k = p : q : r,\qquad g_{1} : h_{1} : k_{1} = p_{1} : q_{1} : r_{1}
\]
folgen, wodurch die L\"osungen dieser Gleichungen sich vereinfachen.

\begin{center}
\textbf{Erster Fall.} Nur eins der drei Gr\"ossenpaare $u$, $u'$; $v$, $v'$; $w$, $w'$ ist gleich Null.
\end{center}

Wenn weder zugleich $v$ und $v'$, noch zugleich $w$ und $w'$ Null sind, folgt aus den Gleichungen ($\beta$) und ($\gamma$)
\[
\frac{v'^{2}}{v^{2}} = \frac{(2a - b - c)(2a + b - c)}{(2a + b + c)(2a - b + c)} = \Bigl(\frac{a - c}{a + c}\Bigr)^{2},\quad\text{u.~s.~w.},
\]
\[
\frac{w'^{2}}{w^{2}} = \frac{(2a - b - c)(2a - b + c)}{(2a + b + c)(2a + b - c)} = \Bigl(\frac{a - b}{a + b}\Bigr)^{2},
\]
woraus sich mit Hinzuziehung von
\[
abc = \text{const.}
\]
ergiebt, dass $a$, $b$, $c$ und folglich auch $v$, $v'$, $w$, $w'$ constant sind.

Setzen wir nun
\begin{align*}
v' &= v\sqrt{\frac{(2a + b + c)(2a - b + c)}{(2a - b - c)(2a + b - c)}} = vS,\\
w' &= w\sqrt{\frac{(2a + b + c)(2a + b - c)}{(2a - b - c)(2a - b + c)}} = wT,\tag{2}
\end{align*}
so erhalten wir aus den drei ersten Differentialgleichungen ($\alpha$) die drei Gleichungen
\begin{align*}
(4a^{2} - b^{2} - 3c^{2})S + (4a^{2} - 3b^{2} - c^{2})T &= \frac{\varepsilon A}{a^{2}} - \frac{\Theta}{a^{2}},\\
(c^{2} - a^{2})S &= \frac{\varepsilon B}{b^{2}} - \frac{\Theta}{b^{2}},\\
(a^{2} - b^{2})T &= \frac{\varepsilon C}{c^{2}} - \frac{\Theta}{c^{2}}.\tag{3,4}
\end{align*}

Um hieraus die Werthe von $S$, $T$ und $\Theta$ abzuleiten, bilde man aus den Gleichungen (4) die Gleichungen
\[
\int_{0}^{\infty}\frac{s\,ds}{\Delta(b^{2} + s)(c^{2} + s)},\qquad \int_{0}^{\infty}\frac{s\,ds}{\Delta(c^{2} + s)(a^{2} + s)},
\]
und substituire diese Werthe in der Gleichung (3)
\[
(4a^{2} - b^{2} - c^{2})(T + S) - 2(b^{2}T + c^{2}S) = \frac{\varepsilon A}{a^{2}} - \frac{\Theta}{a^{2}},
\]
wodurch man
\begin{equation*}
\frac{\varepsilon B - \Theta}{2a^{2}b^{2}c^{2}} = \frac{\varepsilon}{2}\int_{0}^{\infty}\frac{ds}{\Delta}\Bigl(\frac{2s + 4a^{2} - b^{2} - c^{2}}{(b^{2} + s)(c^{2} + s)} + \frac{1}{a^{2} + s}\Bigr)\tag{5}
\end{equation*}
erh\"alt, wenn zur Abk\"urzung
\begin{equation*}
4a^{4} - a^{2}(b^{2} + c^{2}) + b^{2}c^{2} = D\tag{6}
\end{equation*}
gesetzt wird.

% --- page 195 ---
Durch Einsetzung des Werthes von $\Theta$ in die Gleichungen (4) findet sich dann
\begin{align*}
S &= \frac{\varepsilon b^{2}c^{2}}{D}\int_{0}^{\infty}\frac{s\,ds}{\Delta(b^{2} + s)}\Bigl(\frac{4a^{2} - c^{2} + b^{2}}{c^{2} + s} - \frac{b^{2}}{a^{2} + s}\Bigr),\\
T &= \frac{\varepsilon b^{2}c^{2}}{D}\int_{0}^{\infty}\frac{s\,ds}{\Delta(c^{2} + s)}\Bigl(\frac{4a^{2} - b^{2} + c^{2}}{b^{2} + s} - \frac{c^{2}}{a^{2} + s}\Bigr).\tag{7,8}
\end{align*}

Es bleibt nun noch zu untersuchen, welchen Bedingungen $a$, $b$, $c$ gen\"ugen m\"ussen, damit sich aus den Gleichungen (7) und (8) und den Gleichungen (2) f\"ur $v$, $v'$, $w$, $w'$ reelle Werthe ergeben.

Damit $\bigl(\frac{v'}{v}\bigr)^{2}$ und $\bigl(\frac{w'}{w}\bigr)^{2}$ nicht negativ werden, ist es nothwendig und hinreichend, dass die Gr\"osse
\[
(4a^{2} - (b + c)^{2})(4a^{2} - (b - c)^{2}) \geq 0
\]
sei. Es muss also $a^{2}$ entweder $\geq \bigl(\frac{b+c}{2}\bigr)^{2}$ oder $\leq \bigl(\frac{b-c}{2}\bigr)^{2}$ sein.

Wenn $a > \frac{b+c}{2}$, m\"ussen die Gr\"ossen $S$ und $T$ beide $> 0$ sein, damit die Gleichungen (2) f\"ur $v$, $v'$, $w$, $w'$ reelle Werthe liefern. Man kann nun aber leicht zeigen, dass, wenn $a^{2} > \frac{(b+c)^{2}}{4}$, $D$ und die beiden Integrale auf der rechten Seite der Gleichungen (7) und (8) immer positiv sind. Man hat dazu nur n\"othig, $D$ in die Form zu setzen
\[
a^{2}(4a^{2} - (b + c)^{2}) + bc(2a^{2} + bc)
\]
und das in (7) enthaltene Integral in die Form
\[
\int_{0}^{\infty}\frac{s\,ds}{\Delta(c^{2} + s)}((4a^{2} - c^{2})s + a^{2}(4a^{2} + b^{2} - c^{2}) - b^{2}c^{2}),
\]
und dann zu bemerken, dass aus $a \geq \frac{b+c}{2}$ die folgenden Ungleichheiten fliessen: $4a^{2} - (b+c)^{2} \geq 0$, $4a^{2} - c^{2} > 0$, ferner
\[
4a^{2} + b^{2} - c^{2} \geq (b+c)^{2} + b^{2} - c^{2} = 2b(b+c),
\]
und folglich
\[
a^{2}(4a^{2} + b^{2} - c^{2}) \geq 2b(b+c)a^{2} \geq \tfrac{1}{2}b(b+c)^{3} > b^{2}c^{2}.
\]
Aus diesen Ungleichheiten folgt, dass sowohl $D$, als das betrachtete Integral nur positive Bestandtheile hat, und dasselbe gilt auch von dem Integral auf der rechten Seite der Gleichung (8), welches aus diesem durch Vertauschung von $b$ und $c$ erhalten wird. Lassen wir

% --- page 196 ---
nun $a$ die Werthe von $\frac{b+c}{2}$ bis $\infty$ durchlaufen, so wird, wenn $b > c$, $T$ immer positiv bleiben, $S$ aber nur so lange $a < b$. Die Bedingungen f\"ur diesen Fall sind also, wenn $b$ die gr\"ossere der beiden Axen $b$ und $c$ bezeichnet,
\[
\text{(I)}\qquad \frac{b+c}{2} \leq a \leq b.
\]

F\"ur die Untersuchung des zweiten Falles, wenn $a^{2} \leq \bigl(\frac{b-c}{2}\bigr)^{2}$, wollen wir annehmen, dass $b$ die gr\"ossere der beiden Axen $b$ und $c$ sei, so dass $a \leq \frac{b-c}{2}$. Es muss dann, damit $v$, $v'$, $w$, $w'$ reell werden, $S \geq 0$ und $T \geq 0$ sein. Da aus den Ungleichheiten
\[
b^{2} > (2a+c)^{2} > 4a^{2} + c^{2}
\]
hervorgeht, dass das Integral auf der rechten Seite der Gleichung (8) in unserm Falle stets negativ ist, so wird die letztere Bedingung $T \geq 0$ nur erf\"ullt werden, wenn $D(c^{2} - a^{2}) \geq 0$, also $c^{2}$ entweder $< \frac{D}{\cdots}$ oder $> a^{2}$ ist. Dieser Fall spaltet sich also wieder in zwei F\"alle, und diese sind, da $\frac{D}{\cdots} < a^{2} < c^{2}$ durch einen endlichen Zwischenraum getrennt, so dass von einem zum andern kein stetiger Uebergang stattfindet. Da das Integral in der Gleichung (7), so lange $c^{2} < a^{2}$ ist, wegen der beiden Ungleichheiten $c^{2} + s < a^{2} + s$, $4a^{2} - c^{2} + b^{2} > b^{2}$ nur positiv sein kann, so reduciren sich die zu erf\"ullenden Bedingungen im ersten dieser F\"alle auf $a \leq \frac{b-c}{2}$ oder
\[
\text{(II)}\qquad c \leq b - 2a\quad\text{und}\quad c < \sqrt{a^{2} + \sqrt{D}}
\]
und im zweiten auf
\[
\text{(III)}\qquad a < \frac{b-c}{2}\quad\text{und}\quad\int_{0}^{\infty}\frac{s\,ds}{\Delta(b^{2}+s)}\Bigl(\frac{4a^{2}-c^{2}+b^{2}}{c^{2}+s} - \frac{b^{2}}{a^{2}+s}\Bigr) < 0.
\]

Es ist leicht zu sehen, dass das Integral auf der linken Seite der letzten Ungleichheit, wenn $a$ die Werthe von $0$ bis $c$ durchl\"auft, negativ bleibt, so lange $a \leq \frac{c}{2}$ ist, w\"ahrend es f\"ur $a = c$ einen positiven Werth annimmt; die genaue Bestimmung der Grenzen aber, innerhalb welcher diese Ungleichheit erf\"ullt ist, h\"angt, wie man sieht, von der Aufl\"osung einer transcendenten Gleichung ab.

In Bezug auf das Zeichen von $\Theta$, welches bekanntlich entscheidet, ob die Bewegung ohne \"aussern Druck m\"oglich ist, k\"onnen wir bemerken, dass sich der oben gefundene Werth dieser Gr\"osse in die Form

% --- page 197 ---
\[
\Theta = \varepsilon\int_{0}^{\infty}\frac{ds}{\Delta}\frac{3s^{3} + 6a^{2}s^{2} + D}{s + a^{2}}
\]
setzen l\"asst, und also in den F\"allen I und III, wo $D > 0$, jedenfalls positiv ist, f\"ur einen negativen Werth von $D$ aber, wenigstens so lange dieser Werth absolut genommen unter einer gewissen Grenze liegt, negativ wird.

\begin{center}
\textbf{7.}
\end{center}
\begin{center}
\textbf{Zweiter Fall.} Zwei der Gr\"ossenpaare $u$, $u'$; $v$, $v'$; $w$, $w'$ sind gleich Null.
\end{center}

Wir haben nun noch den Fall zu behandeln, wenn zwei der Gr\"ossenpaare $u$, $u'$; $v$, $v'$; $w$, $w'$ fortw\"ahrend Null sind, und also nur um eine Hauptaxe eine Rotation stattfindet.

Wenn ausser $u$ und $u'$ auch $v$ und $v'$ fortw\"ahrend Null sind, so reduciren sich die Gleichungen ($\beta$) und ($\gamma$) auf
\[
(a - b)^{2}w = \text{const.} = \varkappa,\quad (a + b)^{2}w' = \text{const.} = \varkappa'
\]
und die ersten drei Differentialgleichungen ($\alpha$) liefern daher die Gleichungen
\begin{align*}
\frac{\varkappa^{2}}{(a-b)^{4}} + \frac{\varkappa'^{2}}{(a+b)^{4}} + \frac{d^{2}a}{dt^{2}} &= \frac{2\varepsilon aA - \Theta}{a},\\
\frac{\varkappa^{2}}{(b-a)^{4}} + \frac{\varkappa'^{2}}{(b+a)^{4}} + \frac{d^{2}b}{dt^{2}} &= \frac{2\varepsilon bB - \Theta}{b},\\
\frac{d^{2}c}{dt^{2}} &= \frac{2\varepsilon cC - \Theta}{c},\tag{1}
\end{align*}
welche verbunden mit
\[
abc = a_{0}b_{0}c_{0}
\]
die Gr\"ossen $a$, $b$, $c$ und $\Theta$ als Functionen der Zeit bestimmen. Das Princip der Erhaltung der lebendigen Kraft giebt f\"ur diese Differentialgleichungen das Integral erster Ordnung
\[
\Bigl(\frac{da}{dt}\Bigr)^{2} + \Bigl(\frac{db}{dt}\Bigr)^{2} + \Bigl(\frac{dc}{dt}\Bigr)^{2} = 2\varepsilon H + \text{const.},
\]
woraus unmittelbar hervorgeht, dass wenn $\varkappa$ nicht Null ist, die Hauptaxen $a$ und $b$ nie einander gleich werden k\"onnen.

Ausser den schon von Mac Laurin und Dirichlet untersuchten F\"allen, wenn $a = b$, l\"asst noch der Fall, wenn die Gr\"ossen $a$, $b$, $c$ constant sind, eine Bestimmung der Bewegung in geschlossenen Ausdr\"ucken zu. In diesem Falle erh\"alt man aus (1) durch Elimination von $\Theta$ die beiden Gleichungen

% --- page 198 ---
\begin{align*}
\frac{\varkappa^{2}}{(a-b)^{4}} - \frac{\varkappa'^{2}}{(b+a)^{4}} &= \frac{2\varepsilon(A-B)}{a^{2}-b^{2}}\\
&= 2\varepsilon abc\int_{0}^{\infty}\frac{s\,ds}{\Delta}\frac{(b^{2}-c^{2})s}{(a^{2}+s)(b^{2}+s)},\\
\frac{\varkappa^{2}}{(b+a)^{4}} + \frac{\varkappa'^{2}}{(b-a)^{4}} &= \frac{2\varepsilon B}{b} - \frac{2\varepsilon C}{c}\\
&= 2\varepsilon abc\int_{0}^{\infty}\frac{s\,ds}{\Delta}\frac{(a^{2}-c^{2})s}{(b^{2}+s)(c^{2}+s)},\tag{3}
\end{align*}
worin die Integrale auf der rechten Seite durch $K$ und $L$ bezeichnet werden m\"ogen; sie lassen sich auch in die Form setzen
\begin{align*}
\frac{\varkappa^{2}}{(a+b)^{4}} - \frac{\varkappa'^{2}}{(b-a)^{4}} &= K = \frac{1}{2}\int_{0}^{\infty}\frac{ds}{\Delta}\Bigl(\frac{s - ab}{(a^{2}+s)(b^{2}+s)} + \frac{1}{ab(c^{2}+s)}\Bigr),\\
\frac{\varkappa^{2}}{(b-a)^{4}} &= L = \frac{1}{2}\int_{0}^{\infty}\frac{ds}{\Delta}\Bigl(\frac{s + ab}{(a^{2}+s)(b^{2}+s)} - \frac{1}{ab(c^{2}+s)}\Bigr).\tag{5}
\end{align*}

Nehmen wir an, dass $b$, wie in den fr\"uher betrachteten F\"allen, die gr\"ossere der beiden Axen $a$ und $b$ bezeichne, so liefern diese beiden Gleichungen dann und auch nur dann f\"ur $\varkappa^{2}$ und $\varkappa'^{2}$ positive Werthe, wenn $K$ positiv und abgesehen vom Zeichen gr\"osser als $L$ ist; und es ist klar, dass die erste Bedingung erf\"ullt ist, so lange $c < b$. Der zweiten Bedingung wird gen\"ugt, wenn $c = a$, also $\varkappa = 0$ ist, und folglich auch, da $K$ und $L$ sich mit $c$ stetig \"andern, innerhalb eines endlichen Gebiets zu beiden Seiten dieses Werthes. Dieses erstreckt sich aber nicht bis zu den Werthen $b$ und $0$; denn f\"ur $c = b$ w\"urde $\varkappa'^{2}$ negativ werden, f\"ur ein unendlich kleines $c$ aber $\varkappa^{2}$, da dann
\[
K = \frac{1}{2ab}\int_{0}^{\infty}\frac{ds}{\Delta}\frac{s - ab}{a^{2} + s},\qquad L = \frac{1}{2ab}\int_{0}^{\infty}\frac{ds}{\Delta}\frac{s + ab}{b^{2} + s}
\]
und folglich $L > K$ wird. W\"achst $b$, w\"ahrend $a$ und $c$ endlich bleiben, in's Unendliche, so kann $L$ nur dann kleiner als $K$ bleiben, wenn zugleich $a^{2} - c^{2}$ in's Unendliche abnimmt; beide Grenzen f\"ur $c$ sind also dann nur unendlich wenig von $a$ verschieden. Wenn dagegen $b$ seiner unteren Grenze $a$ unendlich nahe kommt, so convergirt die obere Grenze f\"ur $c$, wo $\varkappa'^{2} = 0$ wird, gegen $a$, die untere Grenze aber gegen einen Werth, f\"ur welchen das Integral auf der rechten Seite von (5) verschwindet. Zur Bestimmung dieses Werthes erh\"alt man, wenn man $\frac{c}{a} = \sin\psi$ setzt, die Gleichung
\[
(-5 + 2\cos 2\psi) + \cos 4\psi)(\pi - 2\psi) + 10\sin 2\psi + 2\sin 4\psi = 0,
\]
und diese hat zwischen $\psi = 0$ und $\psi = \frac{\pi}{2}$ nur eine Wurzel, welche

% --- page 199 ---
\[
\psi = 0{,}303327\ldots
\]
giebt. F\"ur $b = a$ kann freilich $c$ jeden Werth zwischen $0$ und $b$ annehmen, da dann $\varkappa^{2}$ wegen des Factors $b - a$ immer Null wird. Man erh\"alt dann den von Mac Laurin untersuchten Fall, w\"ahrend sich f\"ur $\varkappa^{2} = \varkappa'^{2}$ die beiden von Jacobi und Dedekind gefundenen F\"alle ergeben.

Der eben behandelte Fall f\"allt f\"ur $b = a$ mit dem Falle (I) des vorigen Artikels zusammen und, wenn
\[
\frac{w^{2}}{(b+c+2a)(b-c+2a)} = \frac{w'^{2}}{(b+c-2a)(b-c-2a)},
\]
mit dem Falle (III). Von den bisher gefundenen vier F\"allen, in denen das fl\"ussige Ellipsoid w\"ahrend der Bewegung seine Form nicht \"andert, h\"angen also diese drei F\"alle stetig unter einander zusammen, w\"ahrend der Fall (II) isolirt bleibt.

\begin{center}
\textbf{8.}
\end{center}

Die Untersuchung, ob ausser diesen vier F\"allen noch andere vorhanden sind, in denen die Hauptaxen w\"ahrend der Bewegung constant bleiben, f\"uhrt auf eine ziemlich weitl\"aufige Rechnung, welche wir nur kurz andeuten wollen, da sie nur ein negatives Resultat liefert.

Aus der Voraussetzung, dass $a$, $b$, $c$ constant sind, kann man zun\"achst leicht folgern, dass $\Theta$ constant ist, indem man die drei ersten Differentialgleichungen ($\alpha$), multiplicirt mit $a$, $b$, $c$, zu einander addirt und dann die Integralgleichung I, Art.~4, also den Satz von der Erhaltung der lebendigen Kraft, benutzt.

Durch Differentiation dieser drei Gleichungen erh\"alt man dann ferner, wenn man die Werthe von $\frac{du}{dt}$, $\frac{dv}{dt}$, \dots, $\frac{dw'}{dt}$ aus den sechs letzten Differentialgleichungen ($\alpha$) einsetzt, die drei Gleichungen
\begin{align*}
(b - c)u(vw - v'w') + (b + c)u'(v'w - vw') &= 0\\
(c - a)v(wu - w'u') + (c + a)v'(w'u - wu') &= 0\\
(a - b)w(uv - u'v') + (a + b)w'(uv' - u'v) &= 0,\tag{1}
\end{align*}
von denen eine eine Folge der \"ubrigen ist.

\smallskip
\noindent\textbf{I.} Wenn nun keine von den sechs Gr\"ossen $u$, $u'$, \dots, $w'$ Null ist, folgt aus diesen Gleichungen die Gleichheit der folgenden drei Gr\"ossenpaare, deren Werthe wir durch $2a'$, $2b'$, $2c'$ bezeichnen wollen:
\begin{align*}
(b-c)u + (b+c)u' &= (c-a)v + (c+a)v' = (a-b)w + (a+b)w' = 2a'\\
(b-c)u - (b+c)u' &= (c-a)v - (c+a)v' = (a-b)w - (a+b)w' = 2b'\\
&= 2c'.\tag{2}
\end{align*}

% --- page 200 ---
Es ergiebt sich dann $a'^{2} - b'^{2} = a^{2} - b^{2}$, $b'^{2} - c'^{2} = b^{2} - c^{2}$, so dass wir
\[
aa' - aa = bb' - b'b = cc' - c'c = \delta
\]
setzen k\"onnen, und aus den drei ersten Differentialgleichungen ($\alpha$)
\[
2a'd' = \text{const.},\quad 2b'b' = \text{const.},\quad 2c'c' = \text{const.},
\]
wenn wir $vv' + ww'$, $ww' + uu'$, $uu' + vv'$ zur Abk\"urzung durch $\pi$, $\varkappa$, $\varrho$ bezeichnen. Aus diesen Gleichungen und der aus den Integralgleichungen II und III leicht herzuleitenden Gleichung
\[
(a^{2} - b^{2})(a'^{2} - c'^{2})\pi + (b^{2} - a^{2})(b'^{2} - c'^{2})\varkappa + (c^{2} - a^{2})(c'^{2} - b'^{2})\varrho
\]
folgt, wenn nicht $a = b = c$, dass $\delta$ und folglich $u$, $u'$, \dots, $w'$ constant sein m\"ussen. Es ergiebt sich aber leicht, dass dann die sechs letzten Differentialgleichungen ($\alpha$) nicht erf\"ullt werden k\"onnen; und hierdurch ist, wenn nicht alle drei Axen einander gleich sind, die Unzul\"assigkeit der Annahme, dass $u$, $u'$, \dots, $w'$ s\"ammtlich von Null verschieden sind, erwiesen.

Die Annahme $a = b = c$ w\"urde auf den Fall einer ruhenden Kugel f\"uhren; $u'$, $v'$, $w'$ ergeben sich $= 0$, $u$, $v$, $w$ aber bleiben ganz willk\"urlich, was davon herr\"uhrt, dass die

\clearpage
\chapter*{Source segment 5: riemann pages 201 250}
\addcontentsline{toc}{chapter}{Source segment 5: riemann pages 201 250}
\begin{center}
{\LARGE\bfseries Riemann --- Gesammelte Mathematische Werke}\\[1em]
{\Large (Seiten 201--250 des vorliegenden Bandes)}\\[2em]
\end{center}

\tableofcontents
\newpage

% ============================================================
% PAPER X: Liquid Ellipsoid
% ============================================================
\selectlanguage{ngerman}

\section*{X. Ein Beitrag zu den Untersuchungen \\ über die Bewegung eines flüssigen gleichartigen Ellipsoides.}
\addcontentsline{toc}{section}{X. Ein Beitrag zu den Untersuchungen über die Bewegung eines flüssigen gleichartigen Ellipsoides}

\subsection*{2.}
\addcontentsline{toc}{subsection}{Art. 2}

Die Grundgleichungen der Bewegung lauten:
\begin{equation}
\tag{1}
\begin{aligned}
\frac{\partial^2 x}{\partial t^2}\alpha + \frac{\partial^2 y}{\partial t^2}\beta + \frac{\partial^2 z}{\partial t^2}\gamma &= \varepsilon\,\frac{\partial V}{\partial \xi} - \frac{\partial P}{\partial \xi},\\
\frac{\partial^2 x}{\partial t^2}\alpha' + \frac{\partial^2 y}{\partial t^2}\beta' + \frac{\partial^2 z}{\partial t^2}\gamma' &= \varepsilon\,\frac{\partial V}{\partial \eta} - \frac{\partial P}{\partial \eta},\\
\frac{\partial^2 x}{\partial t^2}\alpha'' + \frac{\partial^2 y}{\partial t^2}\beta'' + \frac{\partial^2 z}{\partial t^2}\gamma'' &= \varepsilon\,\frac{\partial V}{\partial \zeta} - \frac{\partial P}{\partial \zeta},
\end{aligned}
\end{equation}
worin $V$ das Potential, $P$ den Druck im Punkte $x$, $y$, $z$ zur Zeit $t$ und $\varepsilon$ die Constante bezeichnet, welche die Anziehung zwischen zwei Masseneinheiten in der Entfernungseinheit ausdrückt.

Es handelt sich nun zunächst darum, die Grössen links vom Gleichheitszeichen in die Form linearer Functionen von den Grössen $\xi$, $\eta$, $\zeta$ zu setzen, wozu einige Vorbereitungen nöthig sind.

Durch Differentiation der Gleichungen (2) erhält man, wenn man zur Abkürzung
\begin{equation}
\tag{2}
\begin{aligned}
\frac{\partial x}{\partial t}\alpha + \frac{\partial y}{\partial t}\beta + \frac{\partial z}{\partial t}\gamma &= \xi',\\
\frac{\partial x}{\partial t}\alpha' + \frac{\partial y}{\partial t}\beta' + \frac{\partial z}{\partial t}\gamma' &= \eta',\\
\frac{\partial x}{\partial t}\alpha'' + \frac{\partial y}{\partial t}\beta'' + \frac{\partial z}{\partial t}\gamma'' &= \zeta'
\end{aligned}
\end{equation}
setzt,
\begin{align*}
\frac{\partial \xi}{\partial t} &= \frac{d\alpha}{dt}x + \frac{d\beta}{dt}y + \frac{d\gamma}{dt}z + \xi',\\
\frac{\partial \eta}{\partial t} &= \frac{d\alpha'}{dt}x + \frac{d\beta'}{dt}y + \frac{d\gamma'}{dt}z + \eta',\\
\frac{\partial \zeta}{\partial t} &= \frac{d\alpha''}{dt}x + \frac{d\beta''}{dt}y + \frac{d\gamma''}{dt}z + \zeta',
\end{align*}
und wenn man hierin $x$, $y$, $z$ wieder durch $\xi$, $\eta$, $\zeta$ ausdrückt
\begin{multline*}
\frac{\partial \xi}{\partial t} = \Bigl(\frac{d\alpha}{dt}\alpha + \frac{d\beta}{dt}\beta + \frac{d\gamma}{dt}\gamma\Bigr)\xi + \Bigl(\frac{d\alpha}{dt}\alpha' + \frac{d\beta}{dt}\beta' + \frac{d\gamma}{dt}\gamma'\Bigr)\eta\\
+ \Bigl(\frac{d\alpha}{dt}\alpha'' + \frac{d\beta}{dt}\beta'' + \frac{d\gamma}{dt}\gamma''\Bigr)\zeta + \xi'
\end{multline*}
\begin{multline*}
\frac{\partial \eta}{\partial t} = \Bigl(\frac{d\alpha'}{dt}\alpha + \frac{d\beta'}{dt}\beta + \frac{d\gamma'}{dt}\gamma\Bigr)\xi + \Bigl(\frac{d\alpha'}{dt}\alpha' + \frac{d\beta'}{dt}\beta' + \frac{d\gamma'}{dt}\gamma'\Bigr)\eta\\
+ \Bigl(\frac{d\alpha'}{dt}\alpha'' + \frac{d\beta'}{dt}\beta'' + \frac{d\gamma'}{dt}\gamma''\Bigr)\zeta + \eta'
\end{multline*}
\begin{multline*}
\frac{\partial \zeta}{\partial t} = \Bigl(\frac{d\alpha''}{dt}\alpha + \frac{d\beta''}{dt}\beta + \frac{d\gamma''}{dt}\gamma\Bigr)\xi + \Bigl(\frac{d\alpha''}{dt}\alpha' + \frac{d\beta''}{dt}\beta' + \frac{d\gamma''}{dt}\gamma'\Bigr)\eta\\
+ \Bigl(\frac{d\alpha''}{dt}\alpha'' + \frac{d\beta''}{dt}\beta'' + \frac{d\gamma''}{dt}\gamma''\Bigr)\zeta + \zeta'.
\end{multline*}

Nun giebt aber die Differentiation der bekannten Gleichungen
\[
\alpha^2 + \beta^2 + \gamma^2 = 1,\quad \alpha\alpha' + \beta\beta' + \gamma\gamma' = 0,\ \text{etc.}
\]

\newpage
% ============================================================
% Continue Paper X
% ============================================================

\begin{align*}
\alpha\frac{d\alpha}{dt} + \beta\frac{d\beta}{dt} + \gamma\frac{d\gamma}{dt} &= 0, &
\alpha'\frac{d\alpha'}{dt} + \beta'\frac{d\beta'}{dt} + \gamma'\frac{d\gamma'}{dt} &= 0,\\
&\alpha''\frac{d\alpha''}{dt} + \beta''\frac{d\beta''}{dt} + \gamma''\frac{d\gamma''}{dt} = 0 &\\
\frac{d\alpha'}{dt}\alpha'' + \frac{d\beta'}{dt}\beta'' + \frac{d\gamma'}{dt}\gamma'' &= -\Bigl(\frac{d\alpha''}{dt}\alpha' + \frac{d\beta''}{dt}\beta' + \frac{d\gamma''}{dt}\gamma'\Bigr) &\\
\frac{d\alpha''}{dt}\alpha + \frac{d\beta''}{dt}\beta + \frac{d\gamma''}{dt}\gamma &= -\Bigl(\frac{d\alpha}{dt}\alpha'' + \frac{d\beta}{dt}\beta'' + \frac{d\gamma}{dt}\gamma''\Bigr) &\\
\frac{d\alpha}{dt}\alpha' + \frac{d\beta}{dt}\beta' + \frac{d\gamma}{dt}\gamma' &= -\Bigl(\frac{d\alpha'}{dt}\alpha + \frac{d\beta'}{dt}\beta + \frac{d\gamma'}{dt}\gamma\Bigr) &
\end{align*}

\begin{equation}
\tag{3}
\end{equation}
und es wird folglich, wenn man diese letzteren drei Grössen durch $p$, $q$, $r$ bezeichnet,
\begin{equation}
\tag{4}
\begin{aligned}
\xi' &= \quad\ \ \frac{\partial \xi}{\partial t} - r\eta + q\zeta,\\
\eta' &= \quad r\xi + \frac{\partial \eta}{\partial t} - p\zeta,\\
\zeta' &= -q\xi + p\eta + \frac{\partial \zeta}{\partial t}.
\end{aligned}
\end{equation}

Durch ein ganz ähnliches Verfahren ergiebt sich aus den Gleichungen (2)
\begin{equation}
\tag{5}
\begin{aligned}
\frac{\partial^2 x}{\partial t^2}\alpha + \frac{\partial^2 y}{\partial t^2}\beta + \frac{\partial^2 z}{\partial t^2}\gamma &= \quad\ \ \frac{\partial \xi'}{\partial t} - r\eta' + q\zeta',\\
\frac{\partial^2 x}{\partial t^2}\alpha' + \frac{\partial^2 y}{\partial t^2}\beta' + \frac{\partial^2 z}{\partial t^2}\gamma' &= \quad r\xi' + \frac{\partial \eta'}{\partial t} - p\zeta',\\
\frac{\partial^2 x}{\partial t^2}\alpha'' + \frac{\partial^2 y}{\partial t^2}\beta'' + \frac{\partial^2 z}{\partial t^2}\gamma'' &= -q\xi' + p\eta' + \frac{\partial \zeta'}{\partial t},
\end{aligned}
\end{equation}
und aus den Gleichungen Art.~1, (3), wenn $p_\iota$, $q_\iota$, $r_\iota$ die Grössen bezeichnen, welche von den Functionen $\alpha_\iota$, $\beta_\iota$, $\ldots$, $\gamma''_{\iota}$ ebenso abhängen, wie die Grössen $p$, $q$, $r$ von den Functionen $\alpha$, $\beta$, $\ldots$, $\gamma''$
\begin{equation}
\tag{6}
\begin{aligned}
\frac{\partial \frac{\xi}{a}}{\partial t} &= r_\iota\frac{\eta}{b} - q_\iota\frac{\zeta}{c},\\
\frac{\partial \frac{\eta}{b}}{\partial t} &= p_\iota\frac{\zeta}{c} - r_\iota\frac{\xi}{a},\\
\frac{\partial \frac{\zeta}{c}}{\partial t} &= q_\iota\frac{\xi}{a} - p_\iota\frac{\eta}{b}.
\end{aligned}
\end{equation}

Setzt man die Werthe $\frac{\partial \xi}{\partial t}$, $\frac{\partial \eta}{\partial t}$, $\frac{\partial \zeta}{\partial t}$ aus (6) in (4) ein, so erhält man
\begin{equation}
\tag{7}
\begin{aligned}
\xi' &= \frac{da}{dt}\frac{\xi}{a} + (ar_\iota - br)\frac{\eta}{b} + (cq - aq_\iota)\frac{\zeta}{c},\\
\eta' &= (ar_\iota - br_\iota)\frac{\xi}{a} + \frac{db}{dt}\frac{\eta}{b} + (bp_\iota - cp)\frac{\zeta}{c},\\
\zeta' &= (cq_\iota - aq)\frac{\xi}{a} + (bp - cp_\iota)\frac{\eta}{b} + \frac{dc}{dt}\frac{\zeta}{c}.
\end{aligned}
\end{equation}

Was die geometrische Bedeutung dieser Grössen betrifft, so sind, wie leicht ersichtlich ist, $\xi'$, $\eta'$, $\zeta'$ die Geschwindigkeitscomponenten des Punktes $x$, $y$, $z$ der flüssigen Masse parallel den Axen $\xi$, $\eta$, $\zeta$; $\frac{\partial \xi}{\partial t}$, $\frac{\partial \eta}{\partial t}$, $\frac{\partial \zeta}{\partial t}$ die ebenso zerlegten relativen Geschwindigkeiten gegen das Coordinatensystem der $\xi$, $\eta$, $\zeta$; ferner in den Gleichungen (1) die Grössen auf der linken Seite die Beschleunigungen und die auf der rechten die beschleunigenden Kräfte parallel diesen Axen; endlich sind $p$, $q$, $r$ die augenblicklichen Rotationen des Coordinatensystems der $\xi$, $\eta$, $\zeta$ um seine Axen und $p_\iota$, $q_\iota$, $r_\iota$ haben dieselbe Bedeutung für das Coordinatensystem der $\xi_\iota$, $\eta_\iota$, $\zeta_\iota$.

\subsection*{3.}
\addcontentsline{toc}{subsection}{Art. 3}

Wenn man nun die Werthe der Grössen $\xi'$, $\eta'$, $\zeta'$ aus (7) in die Gleichungen (5) substituirt und mit Hülfe der Gleichungen (6) die Derivirten von $\frac{\xi}{a}$, $\frac{\eta}{b}$, $\frac{\zeta}{c}$ wieder durch die Grössen $\xi$, $\eta$, $\zeta$ ausdrückt, so nehmen die Grössen auf der linken Seite der Gleichungen (1) die Form linearer Ausdrücke von den Grössen $\xi$, $\eta$, $\zeta$ an. Auf der rechten Seite hat $V$ die Form
\[
H - A\xi^2 - B\eta^2 - C\zeta^2,
\]
worin $H$, $A$, $B$, $C$ auf bekannte Weise von den Grössen $a$, $b$, $c$ abhängen; und man genügt ihnen daher, wenn an der Oberfläche der Druck den constanten Werth $Q$ hat, indem man
\[
P = Q + \sigma\Bigl(1 - \frac{\xi^2}{a^2} - \frac{\eta^2}{b^2} - \frac{\zeta^2}{c^2}\Bigr)
\]
setzt und die zehn Functionen der Zeit $a$, $b$, $c$; $p$, $q$, $r$; $p_\iota$, $q_\iota$, $r_\iota$ und $\sigma$ so bestimmt, dass die neun Coefficienten der Grössen $\xi$, $\eta$, $\zeta$ auf beiden Seiten einander gleich werden und zugleich die aus der Incompressibilität folgende Bedingungsgleichung $abc = a_0b_0c_0$ befriedigt wird.

\newpage
% ============================================================
% Continue with more equations
% ============================================================

Durch Gleichsetzung der Coefficienten von $\frac{\xi}{a}$, $\frac{\eta}{b}$ in der ersten und von $\frac{\xi}{a}$ in der zweiten Gleichung ergiebt sich
\begin{align*}
\frac{d^2a}{dt^2} + 2brr_\iota + 2cqq_\iota - a(r^2 + r_\iota^2 + q^2 + q_\iota^2) &= 2\frac{\sigma}{a} - 2\varepsilon aA,\\
a\frac{dr}{dt} - b\frac{dr_\iota}{dt} + 2\frac{da}{dt}r - 2\frac{db}{dt}r_\iota + ap_\iota q + bp_\iota q_\iota - 2cpq_\iota &= 0,\\
a\frac{dr_\iota}{dt} - b\frac{dr}{dt} + 2\frac{da}{dt}r_\iota - 2\frac{db}{dt}r + ap_\iota q_\iota + bpq - 2cp_\iota q &= 0.
\end{align*}

Aus diesen Gleichungen erhält man die sechs übrigen durch cyclische Versetzung der Axen, oder auch durch beliebige Vertauschungen, wenn man nur dabei beachtet, dass durch Vertauschung zweier Axen nicht bloss die ihnen entsprechenden Grössen vertauscht werden, sondern zugleich die sechs Grössen $p$, $q$, $\ldots$, $r_\iota$ ihr Zeichen ändern.

Man kann diesen Gleichungen eine für die weitere Untersuchung bequemere Form geben, wenn man statt der Grössen $p$, $p_\iota$; $q$, $q_\iota$; $r$, $r_\iota$ ihre halben Summen und Differenzen
\[
u = \frac{p+p_\iota}{2},\quad v = \frac{q+q_\iota}{2},\quad w = \frac{r+r_\iota}{2},
\]
\[
u' = \frac{p-p_\iota}{2},\quad v' = \frac{q-q_\iota}{2},\quad w' = \frac{r-r_\iota}{2}
\]
als unbekannte Functionen einführt.

Dadurch wird das System von Gleichungen, welchen die zehn unbekannten Functionen der Zeit genügen müssen:
\begin{equation}
\tag{$\alpha$}
\begin{cases}
\begin{aligned}
&(a-c)v^2 + (a+c)v'^2 + (a-b)w^2 + (a+b)w'^2 - \tfrac{1}{2}\frac{d^2a}{dt^2} = \varepsilon aA - \frac{\sigma}{a}\\
&(b-a)w^2 + (b+a)w'^2 + (b-c)u^2 + (b+c)u'^2 - \tfrac{1}{2}\frac{d^2b}{dt^2} = \varepsilon bB - \frac{\sigma}{b}\\
&(c-b)u^2 + (c+b)u'^2 + (c-a)v^2 + (c+a)v'^2 - \tfrac{1}{2}\frac{d^2c}{dt^2} = \varepsilon cC - \frac{\sigma}{c}\\
&(b-c)\frac{du}{dt} + 2\frac{d(b-c)}{dt}u + (b+c-2a)vw + (b+c+2a)v'w' = 0\\
&(b+c)\frac{du'}{dt} + 2\frac{d(b+c)}{dt}u' + (b-c+2a)vw' + (b-c-2a)v'w = 0\\
&(c-a)\frac{dv}{dt} + 2\frac{d(c-a)}{dt}v + (c+a-2b)wu + (c+a+2b)w'u' = 0\\
&(c+a)\frac{dv'}{dt} + 2\frac{d(c+a)}{dt}v' + (c-a+2b)wu' + (c-a-2b)w'u = 0\\
&(a-b)\frac{dw}{dt} + 2\frac{d(a-b)}{dt}w + (a+b-2c)uv + (a+b+2c)u'v' = 0\\
&(a+b)\frac{dw'}{dt} + 2\frac{d(a+b)}{dt}w' + (a-b+2c)uv' + (a-b-2c)u'v = 0\\
&\qquad\qquad\qquad\qquad abc = a_0b_0c_0.
\end{aligned}
\end{cases}
\end{equation}

Die Werthe von $A$, $B$, $C$ ergeben sich aus dem bekannten Ausdrucke für $V$
\[
V = H - A\xi^2 - B\eta^2 - C\zeta^2 = \pi\int_0^{\infty}\!\frac{ds}{\Delta}\Bigl(1 - \frac{\xi^2}{a^2+s} - \frac{\eta^2}{b^2+s} - \frac{\zeta^2}{c^2+s}\Bigr),
\]
worin
\[
\Delta = \sqrt{\Bigl(1+\frac{s}{a^2}\Bigr)\Bigl(1+\frac{s}{b^2}\Bigr)\Bigl(1+\frac{s}{c^2}\Bigr)}.
\]

Nach ausgeführter Integration dieser Differentialgleichungen hat man noch, um die Functionen $\alpha$, $\beta$, $\ldots$, $\gamma''$ zu bestimmen, die allgemeine Lösung $\theta$, $\theta'$, $\theta''$ der Differentialgleichungen
\begin{equation}
\tag{$\beta$}
\frac{d\theta}{dt} = r\theta' - q\theta'',\quad \frac{d\theta'}{dt} = -r\theta + p\theta'',\quad \frac{d\theta''}{dt} = q\theta - p\theta'
\end{equation}
zu suchen, --- von welchen, wie aus Art.~2, (3) hervorgeht, $\alpha$, $\alpha'$, $\alpha''$; $\beta$, $\beta'$, $\beta''$; $\gamma$, $\gamma'$, $\gamma''$ die drei particularen Auflösungen sind, die für $t=0$ die Werthe $1$, $0$, $0$; $0$, $1$, $0$; $0$, $0$, $1$ annehmen, --- und zur Bestimmung der Functionen $\alpha_\iota$, $\beta_\iota$, $\ldots$, $\gamma''_{\iota}$ die allgemeine Lösung der simultanen Differentialgleichungen
\begin{equation}
\tag{$\gamma$}
\frac{d\theta_\iota}{dt} = r_\iota\theta'_\iota - q_\iota\theta''_{\iota},\quad \frac{d\theta'_\iota}{dt} = -r_\iota\theta_\iota + p_\iota\theta''_{\iota},\quad \frac{d\theta''_{\iota}}{dt} = q_\iota\theta_\iota - p_\iota\theta'_\iota.
\end{equation}

\subsection*{4.}
\addcontentsline{toc}{subsection}{Art. 4}

Es fragt sich nun, welche Hülfsmittel für die Integration dieser Differentialgleichungen ($\alpha$), ($\beta$), ($\gamma$) die allgemeinen hydrodynamischen Principien darbieten, aus denen Dirichlet sieben Integrale erster Ordnung der durch die Functionen $l$, $m$, $\ldots$, $n''$ zu erfüllenden Differentialgleichungen (\S.~1.~(a)) schöpfte. Die aus ihnen fliessenden Gleichungen lassen sich mit Hülfe der oben für $\xi'$, $\eta'$, $\zeta'$ gegebenen Ausdrücke leicht herleiten.

Der Satz von der Erhaltung der Flächen giebt
\begin{equation}
\tag{1}
\begin{aligned}
(b-c)^2u + (b+c)^2u' &= g = \alpha\,g^0 + \beta\,h^0 + \gamma\,k^0,\\
(c-a)^2v + (c+a)^2v' &= h = \alpha'\,g^0 + \beta'\,h^0 + \gamma'\,k^0,\\
(a-b)^2w + (a+b)^2w' &= k = \alpha''\,g^0 + \beta''\,h^0 + \gamma''\,k^0,
\end{aligned}
\end{equation}
worin die Constanten $g^0$, $h^0$, $k^0$, die Anfangswerthe von $g$, $h$, $k$, mit den Constanten $\mathfrak{K}$, $\mathfrak{K}'$, $\mathfrak{K}''$ in der Abhandlung von Dirichlet übereinkommen; er liefert also das aus den sechs letzten Differentialgleichungen ($\alpha$) leicht zu bestätigende Resultat, dass $\theta = g$, $\theta' = h$, $\theta'' = k$ eine Lösung der Differentialgleichungen ($\beta$) ist.

Aus dem Helmholtz'schen Princip der Erhaltung der Rotation folgen die Gleichungen
\begin{equation}
\tag{2}
\begin{aligned}
(b-c)^2u - (b+c)^2u' &= g_\iota = \alpha_\iota\,g^0_\iota + \beta_\iota\,h^0_\iota + \gamma_\iota\,k^0_\iota,\\
(c-a)^2v - (c+a)^2v' &= h_\iota = \alpha'_\iota\,g^0_\iota + \beta'_\iota\,h^0_\iota + \gamma'_\iota\,k^0_\iota,\\
(a-b)^2w - (a+b)^2w' &= k_\iota = \alpha''_{\iota}\,g^0_\iota + \beta''_{\iota}\,h^0_\iota + \gamma''_{\iota}\,k^0_\iota,
\end{aligned}
\end{equation}
in welchen die Constanten $g^0_\iota$, $h^0_\iota$, $k^0_\iota$ den Grössen $BC\mathfrak{K}$, $CA\mathfrak{B}$, $AB\mathfrak{A}$ der genannten Abhandlung gleich sind.

Der Satz von der Erhaltung der lebendigen Kraft endlich giebt ein Integral erster Ordnung der Differentialgleichungen ($\alpha$)
\begin{equation}
\tag{I}
\begin{cases}
\tfrac{1}{2}\Bigl(\bigl(\frac{da}{dt}\bigr)^2 + \bigl(\frac{db}{dt}\bigr)^2 + \bigl(\frac{dc}{dt}\bigr)^2\Bigr)\\
+ (b-c)^2u^2 + (c-a)^2v^2 + (a-b)^2w^2\\
+ (b+c)^2u'^2 + (c+a)^2v'^2 + (a+b)^2w'^2
\end{cases}
= 2\varepsilon H + \text{const.}
\end{equation}

Aus den Gleichungen (1) und (2) folgen zunächst noch zwei Integrale der Gleichungen ($\alpha$)
\begin{align*}
\tag{II} g^2 + h^2 + k^2 &= \text{const.} = \omega^2\\
\tag{III} g_\iota^2 + h_\iota^2 + k_\iota^2 &= \text{const.} = \omega_\iota^2.
\end{align*}

Ferner lassen sich von den Gleichungen ($\beta$) zwei Integrale
\begin{align*}
\tag{IV} \theta^2 + \theta'^2 + \theta''^2 &= \text{const.}\\
\tag{V} \theta g + \theta'h + \theta''k &= \text{const.}
\end{align*}
angeben, wodurch ihre Integration \textit{allgemein} auf eine Quadratur zurückgeführt wird. Zur Aufstellung ihrer allgemeinen Lösung ist es jedoch, da sie linear und homogen sind, nur nöthig, noch zwei von der Lösung $g$, $h$, $k$ verschiedene \textit{particulare} Lösungen zu suchen, für welchen Zweck man die willkürlichen Constanten in diesen beiden Integralgleichungen so wählen kann, dass sich die Rechnung vereinfacht. Giebt man beiden den Werth Null, so hat man
\begin{equation}
\tag{3} \theta'h + \theta''k = -g\theta,
\end{equation}
und ferner erhält man, wenn man diese Gleichung quadrirt und dazu die Gleichung
\[
-\theta'^2 - \theta''^2 = \theta^2
\]
multiplicirt mit $h^2 + k^2$, addirt
\[
-(\theta'k - \theta''h)^2 = \omega^2\theta^2,
\]
folglich
\begin{equation}
\tag{4} \theta'k - \theta''h = \omega i\theta.
\end{equation}

Durch Auflösung dieser beiden linearen Gleichungen (3) und (4) findet sich
\begin{equation}
\tag{5} \theta' = \frac{-gh + k\omega i}{h^2 + k^2}\theta
\end{equation}
\begin{equation}
\tag{6} \theta'' = \frac{-gk - h\omega i}{h^2 + k^2}\theta
\end{equation}
und durch Einsetzung dieser Werthe in die erste der Gleichungen ($\beta$)
\[
\frac{1}{\theta}\frac{d\theta}{dt} = \frac{-g\frac{dg}{dt} + r\frac{dk}{dt} + q\frac{dh}{dt}}{h^2 + k^2} + \omega i
\]
\begin{equation}
\tag{7} \log\theta = \tfrac{1}{2}\log(h^2 + k^2) + \omega i\int\!\frac{qh + rk}{h^2 + k^2}\,dt + \text{const.}
\end{equation}

Aus dieser in (5), (6) und (7) enthaltenen Lösung der Differentialgleichungen ($\beta$) erhält man eine dritte, indem man für $\sqrt{-1}$ überall $-\sqrt{-1}$ setzt, und es ist dann leicht aus den gefundenen drei particularen Lösungen die Ausdrücke für die Functionen $\alpha$, $\beta$, $\ldots$, $\gamma''$ zu bilden.

Die geometrische Bedeutung jeder reellen Lösung der Differentialgleichungen ($\beta$) besteht darin, dass sie, mit einem geeigneten constanten Factor multiplicirt, die Cosinus der Winkel ausdrückt, welche die Axen der $\xi$, $\eta$, $\zeta$ zur Zeit $t$ mit einer festen Linie machen. Diese feste Linie wird für die erste der drei eben gefundenen Lösungen durch die Normale auf der unveränderlichen Ebene der ganzen bewegten Masse gebildet, für den reellen und den imaginären Bestandtheil der beiden andern durch zwei in dieser Ebene enthaltene und auf einander senkrechte Linien. Die Cosinus der Winkel zwischen den Axen und jener Normalen sind demnach $\frac{g}{\omega}$, $\frac{h}{\omega}$, $\frac{k}{\omega}$; die Lage der Axen gegen diese Normale ergiebt sich also nach Auflösung der Gleichungen ($\alpha$) ohne weitere Integration und zur vollständigen Bestimmung ihrer Lage genügt eine einzige Quadratur, z.~B. die Integration $\omega\int\limits_0^t\!\frac{qh + rk}{h^2 + k^2}\,dt$, welche die Drehung der durch die Normale und die Axe der $\xi$ gehenden Ebene um die Normale giebt.

Ganz Aehnliches gilt von den Differentialgleichungen ($\gamma$). Man kann auf demselben Wege aus den beiden Integralen
\begin{align*}
\tag{VI} \theta_\iota^2 + \theta'_\iota{}^2 + \theta''_{\iota}{}^2 &= \text{const.}\\
\tag{VII} \theta_\iota g_\iota + \theta'_\iota h_\iota + \theta''_{\iota} k_\iota &= \text{const.}
\end{align*}
ihre allgemeine Lösung und folglich auch die Werthe der Grössen $\alpha_\iota$, $\beta_\iota$, $\ldots$, $\gamma''_{\iota}$ zur Zeit $t$ ableiten, und es wird dabei nur eine Quadratur erforderlich sein. Es ergiebt sich dann schliesslich der Ort eines beliebigen Flüssigkeitstheilchens zur Zeit $t$ aus den oben (Art.~1, (1) und (4)) für die Grössen $x$, $y$, $z$ und die Functionen $l$, $m$, $\ldots$, $n''$ gegebenen Ausdrücken.

\subsection*{5.}
\addcontentsline{toc}{subsection}{Art. 5}

Wir wollen uns jetzt Rechenschaft darüber geben, was durch die Zurückführung der Differentialgleichungen zwischen den Functionen $l$, $m$, $\ldots$, $n''$ (der Differentialgleichungen (a) \S.~1 bei Dirichlet) auf unsere Differentialgleichungen für das Geschäft der Integration gewonnen ist. Das System der Differentialgleichungen (a) ist von der sechszehnten Ordnung, und man kennt von demselben sieben Integrale erster Ordnung, wodurch es auf ein System der neunten Ordnung zurückgeführt wird. Das System ($\alpha$) ist nur von der zehnten Ordnung, und man kennt von demselben noch drei Integrale erster Ordnung. Durch die hier bewirkte Umformung jener Differentialgleichungen ist also die Ordnung des noch zu integrirenden Systems von Differentialgleichungen um zwei Einheiten erniedrigt, und man hat statt dessen nur schliesslich noch zwei Quadraturen auszuführen. Diese Umformung leistet also dasselbe, wie die Auffindung von zwei Integralen erster Ordnung.

Wir bemerken indess ausdrücklich, dass hierdurch unsere Form der Differentialgleichungen nur für die Integration und die wirkliche Bestimmung der Bewegung einen Vorzug erhält. Für die allgemeinsten Untersuchungen über diese Bewegung ist dagegen diese Form der Differentialgleichungen weniger geeignet, nicht bloss, weil ihre Herleitung weniger einfach ist, sondern auch deshalb, weil der Fall der Gleichheit zweier Axen eine besondere Betrachtung erfordert. Bei Gleichheit zweier Axen tritt nämlich der besondere Umstand ein, dass die ihnen zu gebende Lage durch die Gestalt der flüssigen Masse nicht völlig bestimmt ist; sie hängt dann im Allgemeinen auch von der augenblicklichen Bewegung ab und bleibt nur dann willkürlich, wenn diese Bewegung so beschaffen ist, dass die Axen fortwährend einander gleich bleiben. Die Untersuchung dieses Falles ist zwar immer leicht und bedarf daher keiner weiteren Ausführung, kann aber in speciellen Fällen noch wieder besondere Formen annehmen, und die allgemeinen Untersuchungen, wie z.~B. der allgemeine Nachweis der Möglichkeit der Bewegung (\S.~2 bei Dirichlet), würden daher wegen der Menge von besonders zu behandelnden Fällen ziemlich weitläufig werden.

Ehe wir zur Behandlung von speciellen Fällen schreiten, in welchen sich die Differentialgleichungen ($\alpha$) integriren lassen, ist es zweckmässig, zu bemerken, dass in einer Lösung dieser Differentialgleichungen, wie unmittelbar aus der Form dieser Gleichungen hervorgeht, jede Zeichenänderung der Functionen $u$, $v$, $\ldots$, $w'$ zulässig ist, bei welcher $uvw$, $uv'w'$, $u'vw'$, $u'v'w$ ungeändert bleiben. Es können also erstens die Zeichen der Functionen $u'$, $v'$, $w'$ gleichzeitig geändert werden, und dadurch werden die Grössen $\alpha$, $\beta$, $\ldots$, $\gamma''$ mit den Grössen $\alpha_\iota$, $\beta_\iota$, $\ldots$, $\gamma''_{\iota}$, also in dem System der Grössen $l$, $m$, $\ldots$, $n''$ die Horizontalreihen mit den Verticalreihen vertauscht. Zweitens können gleichzeitig zwei der Grössenpaare $u$, $u'$; $v$, $v'$; $w$, $w'$ mit den entgegengesetzten Zeichen versehen werden, und diese Aenderung lässt sich auf eine Aenderung in dem Zeichen einer Coordinataxe zurückführen, wobei die Bewegung in eine ihr symmetrisch gleiche übergeht. In dieser Bemerkung ist der von Dedekind gefundene Reciprocitätssatz enthalten.

\subsection*{6.}
\addcontentsline{toc}{subsection}{Art. 6}

Wir wollen nun den Fall untersuchen, in welchem eins der Grössenpaare $u$, $u'$; $v$, $v'$; $w$, $w'$ fortwährend gleich Null ist, also z.~B. $u = u' = 0$; die geometrische Bedeutung dieser Voraussetzung ist diese, dass die Hauptaxe stets in der unveränderlichen Ebene der ganzen bewegten Masse liegt und die augenblickliche Rotationsaxe auf dieser Hauptaxe senkrecht steht.

Aus den sechs letzten Differentialgleichungen ($\alpha$) folgt sogleich, dass in diesem Falle die Grössen
\begin{equation}
\tag{$\mu$} (c-a)^2v,\ (c+a)^2v',\ (a-b)^2w,\ (a+b)^2w' \end{equation}
constant sind und die Gleichungen
\begin{equation}
\tag{$\nu$}
\begin{aligned} (b+c-2a)vw + (b+c+2a)v'w' &= 0\\ (b-c+2a)vw' + (b-c-2a)v'w &= 0 \end{aligned}
\end{equation}
stattfinden müssen.

Bei der weiteren Untersuchung ist zu unterscheiden, ob noch ein zweites der drei Grössenpaare Null ist oder nicht, und wir können im Allgemeinen nur noch bemerken, dass in Folge der Gleichungen ($\mu$) die Grössen $h$, $k$, $h_\iota$, $k_\iota$ constant sind und folglich auch die Winkel zwischen den Hauptaxen und der unveränderlichen Ebene der ganzen bewegten Masse, und dass dann ferner aus den Differentialgleichungen ($\beta$) und ($\gamma$) die Verhältnissgleichungen
\begin{align*} g : h : k &= p : q : r\\ g_\iota : h_\iota : k_\iota &= p_\iota : q_\iota : r_\iota \end{align*}
folgen, wodurch die Lösungen dieser Gleichungen sich vereinfachen.

\textit{Erster Fall. Nur eins der drei Grössenpaare $u$, $u'$; $v$, $v'$; $w$, $w'$ ist gleich Null.}

Wenn weder zugleich $v$ und $v'$, noch zugleich $w$ und $w'$ Null sind, folgt aus den Gleichungen ($\mu$) und ($\nu$)
\begin{equation}
\tag{1}
\begin{aligned} \frac{v'^2}{v^2} &= \frac{(2a-b-c)(2a+b-c)}{(2a+b+c)(2a-b+c)} = \Bigl(\frac{a-c}{a+c}\Bigr)^4\text{const.,}\\ \frac{w'^2}{w^2} &= \frac{(2a-b-c)(2a-b+c)}{(2a+b+c)(2a+b-c)} = \Bigl(\frac{a-b}{a+b}\Bigr)^4\text{const.,} \end{aligned}
\end{equation}
woraus sich mit Hinzuziehung von
\[ abc = \text{const.} \]
ergibt, dass $a$, $b$, $c$ und folglich auch $v$, $v'$, $w$, $w'$ constant sind.

Setzen wir nun
\begin{equation}
\tag{2}
\begin{aligned} \frac{v^2}{(2a+b+c)(2a-b+c)} &= \frac{v'^2}{(2a-b-c)(2a+b-c)} = S\\ \frac{w^2}{(2a+b+c)(2a+b-c)} &= \frac{w'^2}{(2a-b-c)(2a-b+c)} = T, \end{aligned}
\end{equation}
so erhalten wir aus den drei ersten Differentialgleichungen ($\alpha$) die drei Gleichungen
\begin{equation}
\tag{3} (4a^2 - b^2 - 3c^2)S + (4a^2 - 3b^2 - c^2)T = \frac{\varepsilon A}{2} - \frac{\sigma}{2a^2} \end{equation}
\begin{equation}
\tag{4}
\begin{cases} (b^2 - c^2)\,T = \dfrac{\varepsilon B}{2} - \dfrac{\sigma}{2b^2}\\[1ex] (c^2 - b^2)\,S = \dfrac{\varepsilon C}{2} - \dfrac{\sigma}{2c^2}. \end{cases}
\end{equation}

Um hieraus die Werthe von $S$, $T$ und $\sigma$ abzuleiten, bilde man aus den Gleichungen (4) die Gleichungen
\begin{align*} b^2T + c^2S &= \frac{\varepsilon\pi}{2}\int_0^{\infty}\!\frac{s\,ds}{\Delta\,(b^2+s)(c^2+s)}\\ T + S &= \frac{\sigma}{2b^2c^2} - \frac{\varepsilon\pi}{2}\int_0^{\infty}\!\frac{ds}{\Delta\,(b^2+s)(c^2+s)}, \end{align*}
und substituire diese Werthe in der Gleichung (3)
\[
(4a^2 - b^2 - c^2)(T+S) - 2(b^2T + c^2S) = \frac{\varepsilon A}{2} - \frac{\sigma}{2a^2},
\]
wodurch man
\begin{equation}
\tag{5} \frac{D\sigma}{2a^2b^2c^2} = \frac{\varepsilon\pi}{2}\int_0^{\infty}\!\frac{ds}{\Delta}\Bigl(\frac{2s + 4a^2 - b^2 - c^2}{(b^2+s)(c^2+s)} + \frac{1}{a^2+s}\Bigr) \end{equation}
erhält, wenn zur Abkürzung
\begin{equation}
\tag{6} 4a^4 - a^2(b^2 + c^2) + b^2c^2 = D \end{equation}
gesetzt wird.

Durch Einsetzung des Werthes von $\sigma$ in die Gleichungen (4) findet sich dann
\begin{equation}
\tag{7} \frac{b^2 - c^2}{b^2 - a^2}\,DS = \frac{\varepsilon\pi}{2}\int_0^{\infty}\!\frac{s\,ds}{\Delta\,(b^2+s)}\Bigl(\frac{4a^2 - c^2 + b^2}{c^2+s} - \frac{b^2}{a^2+s}\Bigr) \end{equation}
\begin{equation}
\tag{8} \frac{c^2 - b^2}{c^2 - a^2}\,DT = \frac{\varepsilon\pi}{2}\int_0^{\infty}\!\frac{s\,ds}{\Delta\,(c^2+s)}\Bigl(\frac{4a^2 - b^2 + c^2}{b^2+s} - \frac{c^2}{a^2+s}\Bigr). \end{equation}

Es bleibt nun noch zu untersuchen, welchen Bedingungen $a$, $b$, $c$ genügen müssen, damit sich aus den Gleichungen (7) und (8) und den Gleichungen (2) für $v$, $v'$, $w$, $w'$ reelle Werthe ergeben.

Damit $\bigl(\frac{v'}{v}\bigr)^2$ und $\bigl(\frac{w'}{w}\bigr)^2$ nicht negativ werden, ist es nothwendig und hinreichend, dass die Grösse
\[ (4a^2 - (b+c)^2)\,(4a^2 - (b-c)^2) \geq 0 \]
sei. Es muss also $a^2$ entweder $\geq\bigl(\frac{b+c}{2}\bigr)^2$ oder $\leq\bigl(\frac{b-c}{2}\bigr)^2$ sein.

Wenn $a > \frac{b+c}{2}$, müssen die Grössen $S$ und $T$ beide $> 0$ sein, damit die Gleichungen (2) für $v$, $v'$, $w$, $w'$ reelle Werthe liefern. Man kann nun aber leicht zeigen, dass, wenn $a > \frac{b+c}{2}$, $D$ und die beiden Integrale auf der rechten Seite der Gleichungen (7) und (8) immer positiv sind. Man hat dazu nur nöthig, $D$ in die Form zu setzen
\[ a^2(4a^2 - (b+c)^2) + bc(2a^2 + bc) \]
und das in (7) enthaltene Integral in die Form
\[ \frac{\varepsilon\pi}{2a^2b^2c^2}\int_0^{\infty}\!\frac{s\,ds}{\Delta^3}\bigl((4a^2-c^2)s + a^2(4a^2+b^2-c^2) - b^2c^2\bigr), \]
und dann zu bemerken, dass aus $a \geq \frac{b+c}{2}$ die folgenden Ungleichheiten fliessen: $4a^2 - (b+c)^2 \geq 0$, $4a^2 - c^2 > 0$, ferner
\[ 4a^2 + b^2 - c^2 \geq (b+c)^2 + b^2 - c^2 = 2b(b+c), \]
und folglich
\[ a^2(4a^2 + b^2 - c^2) \geq 2b(b+c)a^2 > \tfrac{1}{2}b(b+c)^3 > b^2c^2. \]
Aus diesen Ungleichheiten folgt, dass sowohl $D$, als das betrachtete Integral nur positive Bestandtheile hat, und dasselbe gilt auch von dem Integral auf der rechten Seite der Gleichung (8), welches aus diesem durch Vertauschung von $b$ und $c$ erhalten wird. Lassen wir nun $a$ die Werthe von $\frac{b+c}{2}$ bis $\infty$ durchlaufen, so wird, wenn $b > c$, $T$ immer positiv bleiben, $S$ aber nur so lange $a < b$. Die Bedingungen für diesen Fall sind also, wenn $b$ die grössere der beiden Axen $b$ und $c$ bezeichnet,
\begin{equation}
\tag{I} \frac{b+c}{2} \leq a \leq b. \end{equation}

Für die Untersuchung des zweiten Falles, wenn $a^2 \leq \bigl(\frac{b-c}{2}\bigr)^2$, wollen wir annehmen, dass $b$ die grössere der beiden Axen $b$ und $c$ sei, so dass $a \leq \frac{b-c}{2}$. Es muss dann, damit $v$, $v'$, $w$, $w'$ reell werden, $S \leq 0$ und $T \geq 0$ sein. Da aus den Ungleichheiten
\[ b^2 \geq (2a+c)^2 > 4a^2 + c^2 \]
hervorgeht, dass das Integral auf der rechten Seite der Gleichung (8) in unserm Falle stets negativ ist, so wird die letztere Bedingung $T > 0$ nur erfüllt werden, wenn $D(c^2 - a^2) \geq 0$, also $c^2$ entweder $< \frac{a^2(b^2 - 4a^2)}{b^2 - a^2}$, oder $\geq a^2$ ist. Dieser Fall spaltet sich also wieder in zwei Fälle, und diese sind, da $\frac{a^2(b^2-4a^2)}{b^2-a^2} < a^2$, durch einen endlichen Zwischenraum getrennt, so dass von einem zum andern kein stetiger Uebergang stattfindet. Da das Integral in der Gleichung (7), so lange $c^2 \leq a^2$ ist, wegen der beiden Ungleichheiten $c^2 + s \leq a^2 + s$, $4a^2 - c^2 + b^2 > b^2$ nur positiv sein kann, so reduciren sich die zu erfüllenden Bedingungen im ersten dieser Fälle auf $a \leq \frac{b-c}{2}$ oder
\begin{equation}
\tag{II} c \leq b - 2a \quad\text{und}\quad c^2 < \frac{a^2(b^2 - 4a^2)}{b^2 - a^2} \end{equation}
und im zweiten auf
\begin{equation}
\tag{III} a \leq \frac{b-c}{2} \quad\text{und}\quad \int_0^{\infty}\!\frac{s\,ds}{\Delta\,(b^2+s)}\Bigl(\frac{4a^2 - c^2 + b^2}{c^2+s} - \frac{b^2}{a^2+s}\Bigr) \leq 0. \end{equation}

Es ist leicht zu sehen, dass das Integral auf der linken Seite der letzten Ungleichheit, wenn $a$ die Werthe von $0$ bis $c$ durchläuft, negativ bleibt, so lange $a \leq \frac{c}{2}$ ist, während es für $a = c$ einen positiven Werth annimmt; die genaue Bestimmung der Grenzen aber, innerhalb deren diese Ungleichheit erfüllt ist, hängt, wie man sieht, von der Auflösung einer transcendenten Gleichung ab.

In Bezug auf das Zeichen von $\sigma$, welches bekanntlich entscheidet, ob die Bewegung ohne äussern Druck möglich ist, können wir bemerken, dass sich der oben gefundene Werth dieser Grösse in die Form
\[ \frac{\varepsilon\pi}{D}\int_0^{\infty}\!\frac{3s^2 + 6a^2s + D}{\Delta^3}\,ds \]
setzen lässt, und also in den Fällen I und III, wo $D > 0$, jedenfalls positiv ist, für einen negativen Werth von $D$ aber, wenigstens so lange dieser Werth absolut genommen unter einer gewissen Grenze liegt, negativ wird.

\subsection*{7.}
\addcontentsline{toc}{subsection}{Art. 7}

\textit{Zweiter Fall. Zwei der Grössenpaare $u$, $u'$; $v$, $v'$; $w$, $w'$ sind gleich Null.}

Wir haben nun noch den Fall zu behandeln, wenn zwei der Grössenpaare $u$, $u'$; $v$, $v'$; $w$, $w'$ fortwährend Null sind, und also nur um eine Hauptaxe eine Rotation stattfindet.

Wenn ausser $u$ und $u'$ auch $v$ und $v'$ fortwährend Null sind, so reduciren sich die Gleichungen ($\mu$) und ($\nu$) auf
\[ (a-b)^2w = \text{const.} = \tau \qquad (a+b)^2w' = \text{const.} = \tau' \]
und die ersten drei Differentialgleichungen ($\alpha$) liefern daher die Gleichungen
\begin{equation}
\tag{1}
\begin{aligned}
\frac{\tau^2}{(a-b)^3} + \frac{\tau'^2}{(a+b)^3} - \tfrac{1}{2}\frac{d^2a}{dt^2} &= \varepsilon aA - \frac{\sigma}{a}\\
\frac{\tau^2}{(b-a)^3} + \frac{\tau'^2}{(b+a)^3} - \tfrac{1}{2}\frac{d^2b}{dt^2} &= \varepsilon bB - \frac{\sigma}{b}\\
-\tfrac{1}{2}\frac{d^2c}{dt^2} &= \varepsilon cC - \frac{\sigma}{c},
\end{aligned}
\end{equation}
welche verbunden mit
\[ abc = a_0b_0c_0 \]
die Grössen $a$, $b$, $c$ und $\sigma$ als Functionen der Zeit bestimmen. Das Princip der Erhaltung der lebendigen Kraft giebt für diese Differentialgleichungen das Integral erster Ordnung
\begin{equation}
\tag{2} \tfrac{1}{2}\Bigl(\bigl(\tfrac{da}{dt}\bigr)^2 + \bigl(\tfrac{db}{dt}\bigr)^2 + \bigl(\tfrac{dc}{dt}\bigr)^2\Bigr) + \frac{\tau^2}{(a-b)^2} + \frac{\tau'^2}{(a+b)^2} = 2\varepsilon H + \text{const.}, \end{equation}
woraus unmittelbar hervorgeht, dass wenn $\tau$ nicht Null ist, die Hauptaxen $a$ und $b$ nie einander gleich werden können.

Ausser den schon von Mac Laurin und Dirichlet untersuchten Fällen, wenn $a = b$, lässt noch der Fall, wenn die Grössen $a$, $b$, $c$ constant sind, eine Bestimmung der Bewegung in geschlossenen Ausdrücken zu. In diesem Falle erhält man aus (1) durch Elimination von $\sigma$ die beiden Gleichungen
\begin{equation}
\tag{3}
\begin{aligned}
\frac{\tau'^2}{(b+a)^3} + \frac{\tau^2}{(b-a)^3} &= \frac{\varepsilon\pi}{b}\int_0^{\infty}\!\frac{ds}{\Delta}\frac{(b^2-c^2)s}{(b^2+s)(c^2+s)} = K\\
\frac{\tau'^2}{(b+a)^3} - \frac{\tau^2}{(b-a)^3} &= \frac{\varepsilon\pi}{a}\int_0^{\infty}\!\frac{ds}{\Delta}\frac{(a^2-c^2)s}{(a^2+s)(c^2+s)} = L,
\end{aligned}
\end{equation}
worin die Integrale auf der rechten Seite durch $K$ und $L$ bezeichnet werden mögen; sie lassen sich auch in die Form setzen
\begin{equation}
\tag{4} w'^2 = \frac{\tau'^2}{(b+a)^4} = \frac{\varepsilon\pi}{2}\int_0^{\infty}\!\frac{ds}{\Delta}\Bigl(\frac{s+ab}{(a^2+s)(b^2+s)} - \frac{c^2}{ab(c^2+s)}\Bigr) \end{equation}
\begin{equation}
\tag{5} w^2 = \frac{\tau^2}{(b-a)^4} = \frac{\varepsilon\pi}{2}\int_0^{\infty}\!\frac{ds}{\Delta}\Bigl(\frac{s-ab}{(a^2+s)(b^2+s)} + \frac{c^2}{ab(c^2+s)}\Bigr). \end{equation}

Nehmen wir an, dass $b$, wie in den früher betrachteten Fällen, die grössere der beiden Axen $a$ und $b$ bezeichne, so liefern diese beiden Gleichungen dann und auch nur dann für $\tau^2$ und $\tau'^2$ positive Werthe, wenn $K$ positiv und abgesehen vom Zeichen grösser als $L$ ist; und es ist klar, dass die erste Bedingung erfüllt ist, so lange $c < b$. Der zweiten Bedingung wird genügt, wenn $c = a$, also $L = 0$ ist, und folglich auch, da $K$ und $L$ sich mit $c$ stetig ändern, innerhalb eines endlichen Gebiets zu beiden Seiten dieses Werthes. Dieses erstreckt sich aber nicht bis zu den Werthen $b$ und $0$; denn für $c = b$ würde $\tau'^2$ negativ werden, für ein unendlich kleines $c$ aber $\tau^2$, da dann
\[
\frac{K}{c} = \varepsilon\pi\int_0^{\infty}\!\frac{ds}{s^{\frac{1}{2}}(1+s)^{\frac{3}{2}}(1+\frac{b^2}{a^2}s)^{\frac{1}{2}}},\quad
\frac{L}{c} = \varepsilon\pi\int_0^{\infty}\!\frac{ds}{s^{\frac{1}{2}}(1+s)^{\frac{3}{2}}(1+\frac{a^2}{b^2}s)^{\frac{1}{2}}}
\]
und folglich $L > K$ wird. Wächst $b$, während $a$ und $c$ endlich bleiben, in's Unendliche, so kann $L$ nur dann kleiner als $K$ bleiben, wenn zugleich $a^2 - c^2$ in's Unendliche abnimmt; beide Grenzen für $c$ sind also dann nur unendlich wenig von $a$ verschieden. Wenn dagegen $b$ seiner unteren Grenze $a$ unendlich nahe kommt, so convergirt die obere Grenze für $c$, wo $\tau'^2 = 0$ wird, gegen $a$, die untere aber gegen einen Werth, für welchen das Integral auf der rechten Seite von (5) verschwindet. Zur Bestimmung dieses Werthes erhält man, wenn man $\frac{c}{a} = \sin\psi$ setzt, die Gleichung
\[ (-5 + 2\cos 2\psi + \cos 4\psi)(\pi - 2\psi) + 10\sin 2\psi + 2\sin 4\psi = 0, \]
und diese hat zwischen $\psi = 0$ und $\psi = \frac{\pi}{2}$ nur eine Wurzel, welche
\[ \frac{c}{a} = 0{,}303327\ldots \]
giebt. Für $b = a$ kann freilich $c$ jeden Werth zwischen $0$ und $b$ annehmen, da dann $\tau^2$ wegen des Factors $b-a$ immer Null wird. Man erhält dann den von Mac Laurin untersuchten Fall, während sich für $w^2 = w'^2$ die beiden von Jacobi und Dedekind gefundenen Fälle ergeben.

Der eben behandelte Fall fällt für $b = a$ mit dem Falle (I) des vorigen Artikels zusammen und, wenn
\[ \frac{w'^2}{(b+c+2a)(b-c+2a)} = \frac{w^2}{(b+c-2a)(b-c-2a)}, \]
mit dem Falle (III). Von den bisher gefundenen vier Fällen, in denen das flüssige Ellipsoid während der Bewegung seine Form nicht ändert, hängen also diese drei Fälle stetig unter einander zusammen, während der Fall (II) isolirt bleibt.

\subsection*{8.}
\addcontentsline{toc}{subsection}{Art. 8}

Die Untersuchung, ob ausser diesen vier Fällen noch andere vorhanden sind, in denen die Hauptaxen während der Bewegung constant bleiben, führt auf eine ziemlich weitläufige Rechnung, welche wir nur kurz andeuten wollen, da sie nur ein negatives Resultat liefert.

Aus der Voraussetzung, dass $a$, $b$, $c$ constant sind, kann man zunächst leicht folgern, dass $\sigma$ constant ist, indem man die drei ersten Differentialgleichungen ($\alpha$), multiplicirt mit $a$, $b$, $c$, zu einander addirt und dann die Integralgleichung I, Art.~4, also den Satz von der Erhaltung der lebendigen Kraft, benutzt.

Durch Differentiation dieser drei Gleichungen erhält man dann ferner, wenn man die Werthe von $\frac{du}{dt}$, $\frac{du'}{dt}$, $\ldots$, $\frac{dw'}{dt}$ aus den sechs letzten Differentialgleichungen ($\alpha$) einsetzt, die drei Gleichungen
\begin{equation}
\tag{1}
\begin{aligned}
(b-c)\,u\,(vw - v'w') + (b+c)\,u'\,(v'w - vw') &= 0\\
(c-a)\,v\,(wu - w'u') + (c+a)\,v'\,(w'u - wu') &= 0\\
(a-b)\,w\,(uv - u'v') + (a+b)\,w'\,(u'v - uv') &= 0,
\end{aligned}
\end{equation}
von denen eine eine Folge der übrigen ist.

I. Wenn nun keine von den sechs Grössen $u$, $u'$, $\ldots$, $w'$ Null ist, folgt aus diesen Gleichungen die Gleichheit der folgenden drei Grössenpaare, deren Werthe wir durch $2a'$, $2b'$, $2c'$ bezeichnen wollen:
\begin{align*}
(a-c)\frac{v}{v'} + (a+c)\frac{v'}{v} &= (a-b)\frac{w}{w'} + (a+b)\frac{w'}{w} = 2a'\\
(b-a)\frac{w}{w'} + (b+a)\frac{w'}{w} &= (b-c)\frac{u}{u'} + (b+c)\frac{u'}{u} = 2b'\\
(c-b)\frac{u}{u'} + (c+b)\frac{u'}{u} &= (c-a)\frac{v}{v'} + (c+a)\frac{v'}{v} = 2c'.
\end{align*}

Es ergiebt sich dann $a'^2 - b'^2 = a^2 - b^2$, $b'^2 - c'^2 = b^2 - c^2$, so dass wir
\[ aa - a'a' = bb - b'b' = cc - c'c' = \theta \]
setzen können, und aus den drei ersten Differentialgleichungen ($\alpha$)
\[ 2\pi a' = \text{const.},\quad 2\chi b' = \text{const.},\quad 2\varrho c' = \text{const.}; \]
wenn wir $vv' + ww'$, $ww' + uu'$, $uu' + vv'$ zur Abkürzung durch $\pi$, $\chi$, $\varrho$ bezeichnen. Aus diesen Gleichungen und der aus den Integralgleichungen II und III leicht herzuleitenden Gleichung
\begin{multline*}
(a^2-b^2)(a^2-c^2)\pi + (b^2-a^2)(b^2-c^2)\chi + (c^2-a^2)(c^2-b^2)\varrho\\
= \tfrac{1}{4}(\omega^2 - \omega_\iota^2)
\end{multline*}
folgt, wenn nicht $a = b = c$, dass $\theta$ und folglich $u$, $u'$, $\ldots$, $w'$ constant sein müssen. Es ergiebt sich aber leicht, dass dann die sechs letzten Differentialgleichungen ($\alpha$) nicht erfüllt werden können; und hierdurch ist, wenn nicht alle drei Axen einander gleich sind, die Unzulässigkeit der Annahme, dass $u$, $u'$, $\ldots$, $w'$ sämmtlich von Null verschieden sind, erwiesen.

Die Annahme $a = b = c$ würde auf den Fall einer ruhenden Kugel führen; $u'$, $v'$, $w'$ ergeben sich $= 0$, $u$, $v$, $w$ aber bleiben ganz willkürlich, was davon herrührt, dass die Lage der Axen in jedem Augenblicke willkürlich geändert werden kann.

II. Es bleibt also nur die Annahme übrig, dass eine der Grössen $u$, $u'$, $\ldots$, $w'$ Null ist, und diese zieht, wie wir gleich sehen werden, immer die früher untersuchte Voraussetzung nach sich, dass eins der drei Grössenpaare $u$, $u'$; $v$, $v'$; $w$, $w'$ verschwinde.

1. Wenn eine der Grössen $u'$, $v'$, $w'$, z.~B. $u' = 0$ ist, folgen aus (1) die Gleichungen
\[ (b-c)uvw = 0,\quad (b-c)uv'w' = 0 \]
und diese lassen nur eine von den folgenden Annahmen zu: erstens die früher untersuchte Voraussetzung, zweitens $b = c$, drittens $v = 0$ und $w = 0$ oder $v' = 0$ und $w = 0$, was nicht wesentlich verschieden ist.

Wenn $b = c$, bleibt $u$ ganz willkürlich und kann also auch $= 0$ gesetzt werden, wodurch der früher untersuchte Fall eintritt.

Wenn $v = 0$ und $w' = 0$, erhält man aus den Differentialgleichungen ($\alpha$)
\[ (b-c-2a)uv'w = 0,\ (c+a-2b)uv'w = 0,\ (a-b+2c)uv'w = 0, \]
und, wenn man die erste dieser Gleichungen zur zweiten addirt,
\[ -(a+b)uv'w = 0; \]
es muss also ausser den Grössen $u'$, $v$, $w'$ noch eine der Grössen $u$, $v'$, $w$ Null sein, wodurch wieder der früher untersuchte Fall eintritt.

2. Wenn endlich eine der Grössen $u$, $v$, $w$, z.~B. $u = 0$ ist, folgt aus den Gleichungen (1)
\[ u'v'w = 0,\quad u'vw' = 0 \]
und diese Gleichungen führen entweder zu unserer früheren Voraussetzung, oder zu der Annahme, $u = v' = w' = 0$, welche von der eben untersuchten $u' = v = w' = 0$ nicht wesentlich verschieden ist, oder endlich zu der Annahme $u = v = w = 0$. Unter dieser Voraussetzung aber geben die Differentialgleichungen ($\alpha$) $v'w' = w'u' = u'v' = 0$, und es müssen also noch zwei von den Grössen $u'$, $v'$, $w'$ Null sein, was wieder den früher behandelten Fall liefert.

Es hat sich also ergeben, dass mit der Beständigkeit der Gestalt nothwendig eine Beständigkeit des Bewegungszustandes verbunden ist d.~h., dass allemal, wenn die flüssige Masse fortwährend denselben Körper bildet, auch die relative Bewegung aller Theile dieses Körpers immerfort dieselbe bleibt. Die absolute Bewegung im Raume kann man sich in diesem Falle aus zwei einfacheren zusammengesetzt denken, indem man sich zuerst der flüssigen Masse eine innere Bewegung ertheilt denkt, bei welcher sich die Flüssigkeitstheilchen in ähnlichen, parallelen und auf einem Hauptschnitte senkrechten Ellipsen bewegen, und dann dem ganzen System eine gleichförmige Rotation um eine in diesem Hauptschnitte liegende Axe. Wenn dieser Hauptschnitt, wie oben angenommen, senkrecht zur Hauptaxe $a$ ist, so sind die Cosinus der Winkel zwischen der Umdrehungsaxe und den Hauptaxen $0$, $\frac{h}{\omega}$, $\frac{k}{\omega}$ und die Umdrehungszeit $\frac{2\pi}{\sqrt{q^2+r^2}}$. Ferner sind $0$, $b\frac{h}{\omega_\iota}$, $c\frac{k}{\omega_\iota}$ die auf die Hauptaxen bezogenen Coordinaten des Endpunkts der augenblicklichen Rotationsaxe, und bei der innern Bewegung sind die elliptischen Bahnen der Flüssigkeitstheilchen der in diesem Punkte an das Ellipsoid gelegten Tangentialebene parallel, so dass ihre Mittelpunkte in dieser Rotationsaxe liegen. Die Theilchen bewegen sich in diesen Bahnen so, dass die nach den Mittelpunkten gezogenen Radienvectoren in gleichen Zeiten gleiche Flächen durchstreichen, und durchlaufen sie in der Zeit $\frac{2\pi}{\sqrt{q_\iota^2 + r_\iota^2}}$.

\subsection*{9.}
\addcontentsline{toc}{subsection}{Art. 9}

Wir kehren jetzt zurück zur Betrachtung der Bewegung der flüssigen Masse in dem Falle, wenn $u$, $u'$; $v$, $v'$ fortwährend Null sind und also nur um eine Hauptaxe eine Rotation stattfindet, und bemerken zunächst, dass sich den Gleichungen (1) Art.~7, nach welchen sich die Hauptaxen in diesem Falle ändern, noch eine andere anschaulichere mechanische Bedeutung geben lässt. Man kann sie nämlich betrachten als die Gleichungen für die Bewegung eines materiellen Punktes $(a, b, c)$ von der Masse $1$, der gezwungen ist auf einer durch die Gleichung $abc = \text{const.}$ bestimmten Fläche zu bleiben und von Kräften getrieben wird, deren Potentialfunction der Grösse
\[ \frac{\tau^2}{(a-b)^2} + \frac{\tau'^2}{(a+b)^2} - 2\varepsilon H \]
dem Werthe nach gleich und dem Zeichen nach entgegengesetzt ist.

Bezeichnen wir diese Grösse mit $G$, so lassen sich die Gleichungen für beide Bewegungen in die Form setzen:
\begin{equation}
\tag{1} \frac{d^2a}{dt^2}\,\delta a + \frac{d^2b}{dt^2}\,\delta b + \frac{d^2c}{dt^2}\,\delta c + \delta G = 0 \end{equation}
für alle unendlich kleinen Werthe von $\delta a$, $\delta b$, $\delta c$, welche der Bedingung $abc = \text{const.}$ genügen; und der Satz von der Erhaltung der mechanischen Kraft giebt
\[ \tfrac{1}{2}\Bigl(\bigl(\tfrac{da}{dt}\bigr)^2 + \bigl(\tfrac{db}{dt}\bigr)^2 + \bigl(\tfrac{dc}{dt}\bigr)^2\Bigr) + G = \text{const.}, \]
wonach der von der Formänderung der flüssigen Masse unabhängige Theil der mechanischen Kraft $= G$ ist.

Damit $a$, $b$, $c$ und folglich Form und Bewegungszustand des flüssigen Ellipsoids constant bleiben, wenn $\frac{da}{dt}$, $\frac{db}{dt}$, $\frac{dc}{dt}$ Null sind, ist es offenbar nothwendig und hinreichend, dass die Variation erster Ordnung der Function $G$ von den veränderlichen Grössen $a$, $b$, $c$, zwischen welchen die Bedingung $abc = \text{const.}$ stattfindet, verschwinde, was auf die Gleichungen (3) oder (4) und (5) des Art.~7 führt. Diese Beständigkeit des Bewegungszustandes wird aber nur eine labile sein, wenn der Werth der Function kein Minimumwerth ist; es lassen sich dann immer beliebig kleine Aenderungen des Zustandes der flüssigen Masse angeben, welche eine völlige Aenderung desselben zur Folge haben.

Die directe Untersuchung der Variation zweiter Ordnung für den Fall, wenn die Variation erster Ordnung der Function $G$ verschwindet, würde sehr verwickelt werden; es lässt sich jedoch die Frage, ob die Function für diesen Fall einen Minimumwerth habe, auf folgendem Wege entscheiden.

Zunächst lässt sich leicht zeigen, dass die Function immer, welche Werthe auch $\tau^2$, $\tau'^2$ und $a$, $b$, $c$ haben mögen, für ein System von Werthen der unabhängig veränderlichen Grössen ein Minimum haben müsse; es folgt dies offenbar aus den drei Umständen, dass erstens die Function $G$ für den Grenzfall, wenn die Axen unendlich klein oder

\newpage
% ============================================================
% Paper XI: Theta Functions
% ============================================================
\section*{XI. Ueber das Verschwinden der Theta-Functionen.}
\addcontentsline{toc}{section}{XI. Ueber das Verschwinden der Theta-Functionen}

\begin{center}
(Aus Borchardt's Journal für reine und angewandte Mathematik, Bd.~65. 1865.)
\end{center}

Die zweite Abtheilung meiner im 54.~Bande des mathematischen Journals erschienenen Theorie der Abel'schen Functionen enthält den Beweis eines Satzes über das Verschwinden der $\vartheta$-Functionen, welchen ich sogleich wieder anführen werde, indem ich dabei die in jener Abhandlung angewandten Bezeichnungen als dem Leser bekannt voraussetze. Alles in der Abhandlung noch Folgende enthält kurze Andeutungen über die Anwendung dieses Satzes, welcher bei unserer Methode, die sich auf die Bestimmung der Functionen durch ihre Unstetigkeiten und ihr Unendlichwerden stützt, wie man leicht sieht, die Grundlage der Theorie der Abel'schen Functionen bilden muss. Bei dem Satze selbst und dessen Beweis ist jedoch der Umstand nicht gehörig berücksichtigt worden, dass die $\vartheta$-Function durch die Substitution der Integrale algebraischer Functionen Einer Veränderlichen identisch, d.~h. für jeden Werth dieser Veränderlichen, verschwinden kann. Diesem Mangel abzuhelfen ist die folgende kleine Abhandlung bestimmt.

Bei der Darstellung der Untersuchungen über $\vartheta$-Functionen mit einer unbestimmten Anzahl von Variablen macht sich das Bedürfniss einer abkürzenden Bezeichnung einer Reihe, wie
\[ v_1,\ v_2,\ \ldots,\ v_m \]
geltend, sobald der Ausdruck von $v_\nu$ durch $\nu$ complicirt ist. Man könnte dieses Zeichen ganz analog den Summen- und Productenzeichen bilden; eine solche Bezeichnung würde aber zu viel Raum wegnehmen und innerhalb der Functionszeichen unbequem für den Druck sein; ich ziehe es daher vor
\[ v_1,\ v_2,\ \ldots,\ v_m \quad\text{durch}\quad \left(\begin{array}{c} m\\ \nu\ (v_\nu)\\ 1 \end{array}\right) \]
zu bezeichnen, also
\[ \vartheta(v_1,\ v_2,\ \ldots,\ v_p) \quad\text{durch}\quad \vartheta\left(\begin{array}{c} p\\ \nu\ (v_\nu)\\ 1 \end{array}\right). \]

\subsection*{1.}
\addcontentsline{toc}{subsection}{Art. 1}

Wenn man in der Function $\vartheta(v_1,\ v_2,\ \ldots,\ v_p)$ für die $p$ Veränderlichen $v$ die $p$ Integrale $u_1 - e_1$, $u_2 - e_2$, $\ldots$, $u_p - e_p$ algebraischer wie die Fläche $T$ verzweigter Functionen von $z$ substituirt, so erhält man eine Function von $z$, welche in der ganzen Fläche $T$ ausser den Linien $b$ sich stetig ändert, beim Uebertritt von der negativen auf die positive Seite der Linie $b_\nu$ aber den Factor $e^{-u_\nu^+ + u_\nu^- + 2e_\nu}$ erlangt. Wie in \S.~22 bewiesen worden ist, wird diese Function, wenn sie nicht für alle Werthe von $z$ verschwindet, nur für $p$ Punkte der Fläche $T$ unendlich klein von der ersten Ordnung. Diese Punkte wurden durch $\eta_1$, $\eta_2$, $\ldots$, $\eta_p$ bezeichnet, und der Werth der Function $u_\nu$ im Punkte $\eta_\mu$ durch $\alpha_\nu^{(\mu)}$. Es ergab sich dann nach den $2p$ Modulsystemen der $\vartheta$-Function die Congruenz
\begin{equation}
\tag{1} (e_1,\ e_2,\ \ldots,\ e_p) \equiv \Bigl(\sum_1^p \alpha_1^{(\mu)} + K_1,\ \sum_1^p \alpha_2^{(\mu)} + K_2,\ \ldots,\ \sum_1^p \alpha_p^{(\mu)} + K_p\Bigr), \end{equation}
worin die Grössen $K$ von den bis dahin noch willkürlichen additiven Constanten in den Functionen $u$ abhingen, aber von den Grössen $e$ und den Punkten $\eta$ unabhängig waren.

Führt man die dort angegebene Rechnung aus, so findet sich
\begin{equation}
\tag{2} 2K_\nu = \sum \frac{1}{\pi i}\int (u_\nu^+ + u_\nu^-)\,du_\nu - \varepsilon_\nu\pi i - \sum_{\mu=1}^{\mu=p} \varepsilon'_\mu a_{\mu,\nu}. \end{equation}
In diesem Ausdrucke ist das Integral $\int(u_\nu^+ + u_\nu^-)\,du_\nu$ positiv durch $b_\nu$ auszudehnen, und in der Summe sind für $\nu'$ alle Zahlen von $1$ bis $p$ ausser $\nu$ zu setzen; $\varepsilon_\nu = \pm 1$, je nachdem das Ende von $b_\nu$ auf der positiven oder negativen Seite von $a_\nu$ liegt, und $\varepsilon'_\nu = \pm 1$, je nachdem dasselbe auf der positiven oder negativen Seite von $b_\nu$ liegt. Die Bestimmung der Vorzeichen ist übrigens nur nöthig, wenn die Grössen $e$ nach den in \S.~22 gegebenen Gleichungen aus den Unstetigkeiten von $\log\vartheta$ völlig bestimmt werden sollen; die obige Congruenz (1) bleibt richtig, welche Vorzeichen man wählen mag.

Wir behalten zunächst die dort gemachte vereinfachende Voraussetzung bei, dass die additiven Constanten in den Functionen $u$ so bestimmt werden, dass die Grössen $K$ sämmtlich gleich Null sind. Um die so gewonnenen Resultate schliesslich von dieser beschränkenden Voraussetzung zu befreien, hat man offenbar nur nöthig, überall in den $\vartheta$-Functionen zu den Argumenten $-K_1$, $-K_2$, $\ldots$, $-K_p$ hinzuzufügen.

Wenn also die Function $\vartheta(u_1-e_1,\ u_2-e_2,\ \ldots,\ u_p-e_p)$ für die $p$ Punkte $\eta_1$, $\eta_2$, $\ldots$, $\eta_p$ verschwindet und \textit{nicht identisch für jeden Werth von $z$ verschwindet}, so ist
\[ (e_1,\ e_2,\ \ldots,\ e_p) \equiv \Bigl(\sum_1^p \alpha_1^{(\mu)},\ \sum_1^p \alpha_2^{(\mu)},\ \ldots,\ \sum_1^p \alpha_p^{(\mu)}\Bigr). \]

Dieser Satz gilt für ganz beliebige Werthe der Grössen $e$, und wir haben hieraus, indem wir den Punkt $(s, z)$ mit dem Punkte $\eta_p$ zusammenfallen liessen, geschlossen, dass
\[ \vartheta\Bigl(-\sum_1^{p-1} \alpha_1^{(\mu)},\ -\sum_1^{p-1} \alpha_2^{(\mu)},\ \ldots,\ -\sum_1^{p-1} \alpha_p^{(\mu)}\Bigr) = 0, \]
oder da die $\vartheta$-Function gerade ist,
\[ \vartheta\Bigl(\sum_1^{p-1} \alpha_1^{(\mu)},\ \sum_1^{p-1} \alpha_2^{(\mu)},\ \ldots,\ \sum_1^{p-1} \alpha_p^{(\mu)}\Bigr) = 0, \]
welches auch die Punkte $\eta_1$, $\eta_2$, $\ldots$, $\eta_{p-1}$ seien.

\subsection*{2.}
\addcontentsline{toc}{subsection}{Art. 2}

Der Beweis dieses Satzes bedarf jedoch einer Vervollständigung wegen des Umstandes, dass die Function
\[ \vartheta(u_1-e_1,\ u_2-e_2,\ \ldots,\ u_p-e_p) \]
identisch verschwinden kann (was in der That bei jedem System von gleich verzweigten algebraischen Functionen für gewisse Werthe der Grössen $e$ eintritt).

Wegen dieses Umstandes muss man sich begnügen, zunächst zu zeigen, dass der Satz richtig bleibt, während die Punkte $\eta$ unabhängig von einander innerhalb endlicher Grenzen ihre Lage ändern. Hieraus folgt dann die allgemeine Richtigkeit des Satzes nach dem Principe, dass eine Function einer complexen Grösse nicht innerhalb eines endlichen Gebiets gleich Null sein kann, ohne überall gleich Null zu sein.

Wenn $z$ gegeben ist, so können die Grössen $e_1$, $e_2$, $\ldots$, $e_p$ immer so gewählt werden, dass
\[ \vartheta(u_1-e_1,\ u_2-e_2,\ \ldots,\ u_p-e_p) \]
nicht verschwindet; denn sonst müsste die Function $\vartheta(v_1,\ v_2,\ \ldots,\ v_p)$ für jedwede Werthe der Grössen $v$ verschwinden, und folglich müssten in ihrer Entwicklung nach ganzen Potenzen von $e^{2v_1}$, $e^{2v_2}$, $\ldots$, $e^{2v_p}$ sämmtliche Coefficienten gleich Null sein, was nicht der Fall ist. Die Grössen $e$ können sich dann von einander unabhängig innerhalb endlicher Grössengebiete ändern, ohne dass die Function
\[ \vartheta(u_1-e_1,\ u_2-e_2,\ \ldots,\ u_p-e_p) \]
für diesen Werth von $z$ verschwindet. Oder mit andern Worten: man kann immer ein Grössengebiet $E$ von $2p$ Dimensionen angeben, innerhalb dessen sich das System der Grössen $e$ bewegen kann, ohne dass die Function
\[ \vartheta(u_1-e_1,\ u_2-e_2,\ \ldots,\ u_p-e_p) \]
für diesen Werth von $z$ verschwindet. Sie wird also nur für $p$ Lagen von $(s, z)$ unendlich klein von der ersten Ordnung, und bezeichnet man diese Punkte durch $\eta_1$, $\eta_2$, $\ldots$, $\eta_p$, so ist
\begin{equation}
\tag{1} (e_1,\ e_2,\ \ldots,\ e_p) \equiv \Bigl(\sum_1^p \alpha_1^{(\mu)},\ \sum_1^p \alpha_2^{(\mu)},\ \ldots,\ \sum_1^p \alpha_p^{(\mu)}\Bigr). \end{equation}

Jeder Bestimmungsweise des Systems der Grössen $e$ innerhalb $E$ oder jedem Punkte von $E$ entspricht dann eine Bestimmungsweise der Punkte $\eta$, deren Gesammtheit ein dem Grössengebiete $E$ entsprechendes Grössengebiet $H$ bildet. In Folge der Gleichung (1) entspricht jedem Punkte von $H$ aber auch nur ein Punkt von $E$; hätte also $H$ nur $2p-1$, oder weniger Dimensionen, so würde $E$ nicht $2p$ Dimensionen haben können. Es hat folglich $H$ $2p$ Dimensionen. Die Schlüsse, auf welche sich unser Satz stützt, bleiben daher anwendbar für beliebige Lagen der Punkte $\eta$ innerhalb endlicher Gebiete, und die Gleichung
\[ \vartheta\Bigl(-\sum_1^{p-1} \alpha_1^{(\mu)},\ -\sum_1^{p-1} \alpha_2^{(\mu)},\ \ldots,\ -\sum_1^{p-1} \alpha_p^{(\mu)}\Bigr) = 0 \]
gilt für beliebige Lagen der Punkte $\eta_1$, $\eta_2$, $\ldots$, $\eta_{p-1}$ innerhalb endlicher Gebiete und folglich allgemein.

\subsection*{3.}
\addcontentsline{toc}{subsection}{Art. 3}

Hieraus folgt, dass sich das Grössensystem $(e_1,\ e_2,\ \ldots,\ e_p)$ immer und nur auf eine Weise congruent einem Ausdrucke von der Form $\left(\begin{array}{c} p\\ \nu\ \bigl(\sum\limits_1^p \alpha_\nu^{(\mu)}\bigr)\\ 1 \end{array}\right)$ setzen lässt, wenn $\vartheta\left(\begin{array}{c} p\\ \nu\ (u_\nu - e_\nu)\\ 1 \end{array}\right)$ nicht für jeden Werth von $z$ verschwindet; denn liessen sich die Punkte $\eta_1$, $\eta_2$, $\ldots$, $\eta_p$ auf mehr als eine Weise so bestimmen, dass der Congruenz
\[ \left(\begin{array}{c} p\\ \nu\ (e_\nu)\\ 1 \end{array}\right) \equiv \left(\begin{array}{c} p\\ \nu\ \bigl(\sum\limits_1^p \alpha_\nu^{(\mu)}\bigr)\\ 1 \end{array}\right) \]
genügt wäre, so würde nach dem eben bewiesenen Satze die Function $\vartheta\left(\begin{array}{c} p\\ \nu\ (u_\nu - e_\nu)\\ 1 \end{array}\right)$ für mehr als $p$ Punkte verschwinden, ohne identisch gleich Null zu sein, was unmöglich ist.

Wenn $\vartheta\left(\begin{array}{c} p\\ \nu\ (u_\nu - e_\nu)\\ 1 \end{array}\right)$ identisch verschwindet, muss man, um $\left(\begin{array}{c} p\\ \nu\ (e_\nu)\\ 1 \end{array}\right)$ in die obige Form zu setzen,
\[ \vartheta\left(\begin{array}{c} p\\ \nu\ (u_\nu + \alpha_\nu^{(1)} - u_\nu^{(1)} - e_\nu)\\ 1 \end{array}\right) \]
betrachten, und wenn diese Function identisch für jeden Werth $z$, $\xi_1$, $z_1$ verschwindet, die Function
\[ \vartheta\left(\begin{array}{c} p\\ \nu\ \bigl(u_\nu + \sum\limits_1^2 \alpha_\nu^{(\mu)} - \sum\limits_1^2 u_\nu^{(\mu)} - e_\nu\bigr)\\ 1 \end{array}\right). \]

Wir nehmen an, dass
\begin{equation}
\tag{1}
\left\{
\begin{aligned}
&\vartheta\left(\begin{array}{c} p\\ \nu\ \bigl(\sum\limits_1^m \alpha_\nu^{(p+2-\mu)} - \sum\limits_1^{m-1} u_\nu^{(p-\mu)} - e_\nu\bigr)\\ 1 \end{array}\right)\\
&\quad\text{identisch verschwindet,}\\
&\vartheta\left(\begin{array}{c} p\\ \nu\ \bigl(\sum\limits_1^{m+1} \alpha_\nu^{(p+2-\mu)} - \sum\limits_1^m u_\nu^{(p-\mu)} - e_\nu\bigr)\\ 1 \end{array}\right)\\
&\quad\text{aber nicht identisch verschwindet.}
\end{aligned}
\right.
\end{equation}

Diese letztere Function verschwindet dann, als Function von $\xi_{p+1}$ betrachtet, für $\xi_{p-1}$, $\xi_{p-2}$, $\ldots$, $\xi_{p-m}$, ausserdem also noch für $p-m$ Punkte, und bezeichnet man diese mit $\eta_1$, $\eta_2$, $\ldots$, $\eta_{p-m}$, so ist
\[ \left(\begin{array}{c} p\\ \nu\ \bigl(-\sum\limits_{p-m+1}^p \alpha_\nu^{(\mu)} + e_\nu\bigr)\\ 1 \end{array}\right) \equiv \left(\begin{array}{c} p\\ \nu\ \bigl(\sum\limits_1^{p-m} \alpha_\nu^{(\mu)}\bigr)\\ 1 \end{array}\right) \]
und diese Punkte $\eta_1$, $\eta_2$, $\ldots$, $\eta_{p-m}$ können nur auf eine Weise so bestimmt werden, dass diese Congruenz erfüllt wird, weil sonst die Function für mehr als $p$ Punkte verschwinden würde. Dieselbe Function verschwindet, als Function von $\xi_{p-1}$ betrachtet, ausser für
\[ \eta_{p+1},\ \eta_p,\ \ldots,\ \eta_{p-m+1} \]
noch für $p-m-1$ Punkte, und bezeichnet man diese durch
\[ \varepsilon_1,\ \varepsilon_2,\ \ldots,\ \varepsilon_{p-m-1}, \]
so ist
\[ \left(\begin{array}{c} p\\ \nu\ \bigl(-\sum\limits_{p-m}^{p-2} u_\nu^{(\mu)} - e_\nu\bigr)\\ 1 \end{array}\right) \equiv \left(\begin{array}{c} p\\ \nu\ \bigl(\sum\limits_1^{p-m-1} u_\nu^{(\mu)}\bigr)\\ 1 \end{array}\right), \]
und die Punkte $\varepsilon_1$, $\varepsilon_2$, $\ldots$, $\varepsilon_{p-m-1}$ sind durch diese Congruenz völlig bestimmt.

\subsection*{4.}
\addcontentsline{toc}{subsection}{Art. 4}

Bezeichnen wir die Derivirte von
\begin{equation}
\tag{1}
\left\{
\begin{aligned}
&\vartheta(v_1,\ v_2,\ \ldots,\ v_p)\\
&\text{nach } v_\nu \text{ mit } \vartheta'_\nu, \text{ die zweite Derivirte nach } v_\nu \text{ und } v_\mu \text{ mit } \vartheta'_{\nu,\mu} \text{ u.~s.~f.,}
\end{aligned}
\right.
\end{equation}
so sind, wenn
\[ \vartheta\left(\begin{array}{c} p\\ \nu\ (u_\nu^{(1)} - \alpha_\nu^{(1)} + r_\nu)\\ 1 \end{array}\right) \]
identisch für jeden Werth von $z_1$ und $\xi_1$ verschwindet, sämmtliche Functionen $\vartheta'\left(\begin{array}{c} p\\ \nu\ (r_\nu)\\ 1 \end{array}\right)$ gleich Null. In der That geht die Gleichung
\[ \vartheta\left(\begin{array}{c} p\\ \nu\ (u_\nu^{(1)} - \alpha_\nu^{(1)} + r_\nu)\\ 1 \end{array}\right) = 0, \]
 wenn $s_1$ und $z_1$ unendlich wenig von $\sigma_1$ und $\xi_1$ verschieden sind, über in die Gleichung
\[ \sum_1^p \vartheta'_\mu\left(\begin{array}{c} p\\ \nu\ (r_\nu)\\ 1 \end{array}\right) d\alpha_\mu^{(1)} = 0. \]
Nehmen wir an, dass
\[ du_\mu = \frac{\varphi_\mu(s,\ z)\,dz}{\frac{\partial F}{\partial s}} \]
sei, so verwandelt sich diese Gleichung nach Weglassung des Factors $\frac{\partial F(\sigma_1,\ \xi_1)}{\partial \sigma_1}$ in
\[ \sum_1^p \vartheta'_\mu\left(\begin{array}{c} p\\ \nu\ (r_\nu)\\ 1 \end{array}\right) \varphi_\mu(\sigma_1,\ \xi_1) = 0; \]
und da zwischen den Functionen $\varphi$ keine lineare Gleichung mit constanten Coefficienten stattfindet, so folgt hieraus, dass sämmtliche erste Derivirten von $\vartheta(v_1,\ v_2,\ \ldots,\ v_p)$ für $\left(\begin{array}{c} p\\ \nu\ (v_\nu = r_\nu)\\ 1 \end{array}\right)$ verschwinden müssen.

Um den umgekehrten Satz zu beweisen, nehmen wir an, dass $\left(\begin{array}{c} p\\ \nu\ (v_\nu = r_\nu)\\ 1 \end{array}\right)$ und $\left(\begin{array}{c} p\\ \nu\ (v_\nu = t_\nu)\\ 1 \end{array}\right)$ zwei Werthsysteme seien, für welche die Function $\vartheta$ verschwindet, ohne für $\left(\begin{array}{c} p\\ \nu\ (v_\nu = u_\nu^{(1)} - \alpha_\nu^{(1)} + r_\nu)\\ 1 \end{array}\right)$ und $\left(\begin{array}{c} p\\ \nu\ (v_\nu = u_\nu^{(1)} - \alpha_\nu^{(1)} + t_\nu)\\ 1 \end{array}\right)$ identisch zu verschwinden, und bilden den Ausdruck
\begin{equation}
\tag{2}
\frac{\vartheta\left(\begin{array}{c} p\\ \nu\ (u_\nu^{(1)} - \alpha_\nu^{(1)} + r_\nu)\\ 1 \end{array}\right) \vartheta\left(\begin{array}{c} p\\ \nu\ (\alpha_\nu^{(1)} - u_\nu^{(1)} + r_\nu)\\ 1 \end{array}\right)}{\vartheta\left(\begin{array}{c} p\\ \nu\ (u_\nu^{(1)} - \alpha_\nu^{(1)} + t_\nu)\\ 1 \end{array}\right) \vartheta\left(\begin{array}{c} p\\ \nu\ (\alpha_\nu^{(1)} - u_\nu^{(1)} + t_\nu)\\ 1 \end{array}\right)}. \end{equation}

Betrachten wir diesen Ausdruck als Function von $z_1$, so ergiebt sich, dass er eine algebraische Function von $z_1$ und zwar eine rationale Function von $s_1$ und $z_1$ ist, da Nenner und Zähler in $T''$ stetig sind und an den Querschnitten dieselben Factoren erlangen. Für $z_1 = \xi_1$ und $s_1 = \sigma_1$ werden Nenner und Zähler unendlich klein von der zweiten Ordnung, so dass die Function endlich bleibt; die übrigen Werthe aber, für welche Nenner oder Zähler verschwinden, sind, wie oben bewiesen, durch die Werthe der Grössen $r$ und der Grössen $t$ völlig bestimmt, also von $\xi_1$ ganz unabhängig. Da nun eine algebraische Function durch die Werthe, für welche sie Null und unendlich wird, bis auf einen constanten Factor bestimmt ist, so ist der Ausdruck gleich einer rationalen von $\xi_1$ unabhängigen Function von $s_1$ und $z_1$, $\chi(s_1,\ z_1)$, multiplicirt in eine Constante, d.~h. eine von $z_1$ unabhängige Grösse. Da der Ausdruck symmetrisch in Bezug auf die Grössensysteme $(s_1,\ z_1)$ und $(\sigma_1,\ \xi_1)$ ist, so ist diese Constante gleich $\chi(\sigma_1,\ \xi_1)$, multiplicirt in eine auch von $\xi_1$ unabhängige Grösse $A$. Setzt man nun
\[ \sqrt{A}\,\chi(s,\ z) = \varrho(s,\ z), \]
so erhält man für unsern Ausdruck (2) den Werth
\begin{equation}
\tag{3} \varrho(s_1,\ z_1)\,\varrho(\sigma_1,\ \xi_1), \end{equation}
wo $\varrho(s,\ z)$ eine rationale Function von $s$ und $z$ ist.

Um diese zu bestimmen, hat man nur nöthig $\xi_1 = z_1$ und $\sigma_1 = s_1$ werden zu lassen; es ergiebt sich dann
\begin{equation}
\tag{4} \bigl(\varrho(s_1,\ z_1)\bigr)^2 = \frac{\sum\limits_\mu \vartheta'_\mu\left(\begin{array}{c} p\\ \nu\ (r_\nu)\\ 1 \end{array}\right) \varphi_\mu(s_1,\ z_1)}{\sum\limits_\mu \vartheta'_\mu\left(\begin{array}{c} p\\ \nu\ (t_\nu)\\ 1 \end{array}\right) \varphi_\mu(s_1,\ z_1)}. \end{equation}

Man hat daher aus (3) und (4) die Gleichung
\begin{equation}
\tag{5}
\left\{
\begin{aligned}
&\frac{\vartheta\left(\begin{array}{c} p\\ \nu\ (u_\nu^{(1)} - \alpha_\nu^{(1)} + r_\nu)\\ 1 \end{array}\right) \vartheta\left(\begin{array}{c} p\\ \nu\ (\alpha_\nu^{(1)} - u_\nu^{(1)} + r_\nu)\\ 1 \end{array}\right)}{\vartheta\left(\begin{array}{c} p\\ \nu\ (u_\nu^{(1)} - \alpha_\nu^{(1)} + t_\nu)\\ 1 \end{array}\right) \vartheta\left(\begin{array}{c} p\\ \nu\ (\alpha_\nu^{(1)} - u_\nu^{(1)} + t_\nu)\\ 1 \end{array}\right)}\\[1ex]
&= \frac{\sum\limits_\mu \vartheta'_\mu\left(\begin{array}{c} p\\ \nu\ (r_\nu)\\ 1 \end{array}\right) \varphi_\mu(s_1,\ z_1)}{\sum\limits_\mu \vartheta'_\mu\left(\begin{array}{c} p\\ \nu\ (t_\nu)\\ 1 \end{array}\right) \varphi_\mu(s_1,\ z_1)} \cdot \frac{\sum\limits_\mu \vartheta'_\mu\left(\begin{array}{c} p\\ \nu\ (r_\nu)\\ 1 \end{array}\right) \varphi_\mu(\sigma_1,\ \xi_1)}{\sum\limits_\mu \vartheta'_\mu\left(\begin{array}{c} p\\ \nu\ (t_\nu)\\ 1 \end{array}\right) \varphi_\mu(\sigma_1,\ \xi_1)}. \end{aligned}
\right.
\end{equation}

Aus dieser Gleichung folgt, dass
\[ \vartheta\left(\begin{array}{c} p\\ \nu\ (u_\nu^{(1)} - \alpha_\nu^{(1)} + r_\nu)\\ 1 \end{array}\right) \]
für jeden Werth von $z_1$ und $\xi_1$ gleich Null sein muss, wenn die ersten Derivirten der Function $\vartheta(v_1,\ v_2,\ \ldots,\ v_p)$ für $\left(\begin{array}{c} p\\ \nu\ (v_\nu = r_\nu)\\ 1 \end{array}\right)$ sämmtlich verschwinden.

\subsection*{5.}
\addcontentsline{toc}{subsection}{Art. 5}

Wenn
\begin{equation}
\tag{1} \vartheta\left(\begin{array}{c} p\\ \nu\ \bigl(\sum\limits_1^m \alpha_\nu^{(\mu)} - \sum\limits_1^m u_\nu^{(\mu)} + r_\nu\bigr)\\ 1 \end{array}\right) \end{equation}
identisch, d.~h. für jedwede Werthe von $\left(\begin{array}{c} m\\ \mu\ (\sigma_\mu,\ \xi_\mu)\\ 1 \end{array}\right)$ und $\left(\begin{array}{c} m\\ \mu\ (s_\mu,\ z_\mu)\\ 1 \end{array}\right)$, verschwindet, so findet man auf dem oben angegebenen Wege zunächst, indem man $\xi_m = z_m$, $\sigma_m = s_m$ werden lässt, dass die ersten Derivirten der Function
\[ \vartheta(v_1,\ v_2,\ \ldots,\ v_p) \text{ für } \left(\begin{array}{c} p\\ \nu\ \bigl(v_\nu = \sum\limits_1^{m-1} \alpha_\nu^{(\mu)} - \sum\limits_1^{m-1} u_\nu^{(\mu)} + r_\nu\bigr)\\ 1 \end{array}\right) \]
sämmtlich verschwinden, dann, indem man $\xi_{m-1} - z_{m-1}$, $\sigma_{m-1} - s_{m-1}$ unendlich klein werden lässt, dass für
\[ \left(\begin{array}{c} p\\ \nu\ \bigl(v_\nu = \sum\limits_1^{m-2} \alpha_\nu^{(\mu)} - \sum\limits_1^{m-2} u_\nu^{(\mu)} + r_\nu\bigr)\\ 1 \end{array}\right) \]
auch die zweiten Derivirten sämmtlich verschwinden; und offenbar ergiebt sich allgemein, dass die Derivirten $n$ter Ordnung sämmtlich verschwinden für
\[ \left(\begin{array}{c} p\\ \nu\ \bigl(v_\nu = \sum\limits_1^{m-n} \alpha_\nu^{(\mu)} - \sum\limits_1^{m-n} u_\nu^{(\mu)} + r_\nu\bigr)\\ 1 \end{array}\right), \]
welche Werthe auch die Grössen $z$ und die Grössen $\xi$ haben mögen.

Es folgt hieraus, dass unter der gegenwärtigen Voraussetzung (1) für $\left(\begin{array}{c} p\\ \nu\ (v_\nu = r_\nu)\\ 1 \end{array}\right)$ die ersten bis $m$ten Derivirten der Function $\vartheta(v_1,\ v_2,\ \ldots,\ v_p)$ sämmtlich gleich Null sind.

Um zu zeigen, dass dieser Satz auch umgekehrt gilt, beweisen wir zunächst, dass wenn
\[ \vartheta\left(\begin{array}{c} p\\ \nu\ \bigl(\sum\limits_1^{m-1} \alpha_\nu^{(\mu)} - \sum\limits_1^{m-1} u_\nu^{(\mu)} + r_\nu\bigr)\\ 1 \end{array}\right) \]
identisch verschwindet und die Grössen $\vartheta^{(m)}\left(\begin{array}{c} p\\ \nu\ (r_\nu)\\ 1 \end{array}\right)$ sämmtlich gleich Null sind, auch
\[ \vartheta\left(\begin{array}{c} p\\ \nu\ \bigl(\sum\limits_1^m \alpha_\nu^{(\mu)} - \sum\limits_1^m u_\nu^{(\mu)} + r_\nu\bigr)\\ 1 \end{array}\right) \]
identisch verschwinden muss und verallgemeinern zu diesem Zwecke die Gleichung \S.~4, (5).

\newpage
% ============================================================
% End of Paper XI and transition to Paper XII
% ============================================================

Wir nehmen an, dass
\[ \vartheta\left(\begin{array}{c} p\\ \nu\ \bigl(\sum\limits_1^{m-1} u_\nu^{(\mu)} - \sum\limits_1^{m-1} \alpha_\nu^{(\mu)} + r_\nu\bigr)\\ 1 \end{array}\right) \]
identisch verschwinde,
\[ \vartheta\left(\begin{array}{c} p\\ \nu\ \bigl(\sum\limits_1^m u_\nu^{(\mu)} - \sum\limits_1^m \alpha_\nu^{(\mu)} + r_\nu\bigr)\\ 1 \end{array}\right) \]
aber nicht identisch verschwinde, behalten in Bezug auf die Grössen $t$ die frühere Voraussetzung bei und betrachten den Ausdruck
\begin{multline}
\tag{2}
\frac{\vartheta\left(\begin{smallmatrix} p\\ \nu\ \bigl(\sum_1^m u_\nu^{(\mu)} - \sum_1^m \alpha_\nu^{(\mu)} + r_\nu\bigr)\\ 1 \end{smallmatrix}\right)
\vartheta\left(\begin{smallmatrix} p\\ \nu\ \bigl(\sum_1^m \alpha_\nu^{(\mu)} - \sum_1^m u_\nu^{(\mu)} + r_\nu\bigr)\\ 1 \end{smallmatrix}\right)}{\prod_{\varrho,\varrho'}
\vartheta\left(\begin{smallmatrix} p\\ \nu\ (u_\nu^{(\varrho)} - u_\nu^{(\varrho')} + t_\nu)\\ 1 \end{smallmatrix}\right)
\vartheta\left(\begin{smallmatrix} p\\ \nu\ (\alpha_\nu^{(\varrho)} - \alpha_\nu^{(\varrho')} + t_\nu)\\ 1 \end{smallmatrix}\right)}\\
\cdot\Bigl(\prod_{\varrho=1}^m
\vartheta\left(\begin{smallmatrix} p\\ \nu\ (u_\nu^{(\varrho)} - \alpha_\nu^{(\varrho)} + t_\nu)\\ 1 \end{smallmatrix}\right)
\vartheta\left(\begin{smallmatrix} p\\ \nu\ (\alpha_\nu^{(\varrho)} - u_\nu^{(\varrho)} + t_\nu)\\ 1 \end{smallmatrix}\right)\Bigr)^{-2}.
\end{multline}

In diesem Ausdrucke sind unter den Productzeichen sowohl für $\varrho$, als für $\varrho'$ sämmtliche Werthe von $1$ bis $m$ zu setzen, im Zähler aber die Fälle, wo $\varrho = \varrho'$ würde, wegzulassen.

Betrachten wir diesen Ausdruck als Function von $z_1$, so ergiebt sich, dass er an den Querschnitten den Factor $1$ erlangt und folglich eine algebraische Function von $z_1$ ist. Für $z_1 = \xi_\varrho$ und $s_1 = \sigma_\varrho$ werden Nenner und Zähler unendlich klein von der zweiten Ordnung, der Bruch bleibt also endlich; die übrigen Werthe aber, für welche Zähler und Nenner verschwinden, sind durch die Grössen $\left(\begin{array}{c} m\\ \mu\ (s_\mu,\ z_\mu)\\ 2 \end{array}\right)$, die Grössen $r$ und die Grössen $t$, wie oben (\S.~3) bewiesen, völlig bestimmt, und folglich von den Grössen $\xi$ ganz unabhängig. Da der Ausdruck nun eine symmetrische Function von den Grössen $z$ ist, so gilt dasselbe für jedes beliebige $z_\mu$: er ist eine algebraische Function von $z_\mu$, und die Werthe dieser Grösse, für welche er unendlich gross oder unendlich klein wird, sind von den Grössen $\xi$ unabhängig. Er ist daher gleich einer von den Grössen $\xi$ unabhängigen algebraischen Function der Grössen $z$, $\chi(z_1,\ z_2,\ \ldots,\ z_m)$, multiplicirt in einen von den Grössen $z$ unabhängigen Factor. Da er aber ungeändert bleibt, wenn man die Grössen $z$ mit den Grössen $\xi$ vertauscht, so ist dieser Factor gleich $\chi(\xi_1,\ \xi_2,\ \ldots,\ \xi_m)$, multiplicirt mit einer von den Grössen $z$ und den Grössen $\xi$ unabhängigen Constanten $A$; und wir können daher, wenn wir $\sqrt{A}\,\chi(z_1,\ z_2,\ \ldots,\ z_m) = \psi(z_1,\ z_2,\ \ldots,\ z_m)$ setzen, unserm Ausdrucke (2) die Form
\begin{equation}
\tag{3} \psi(z_1,\ z_2,\ \ldots,\ z_m)\,\psi(\xi_1,\ \xi_2,\ \ldots,\ \xi_m) \end{equation}
geben, wo $\psi(z_1,\ z_2,\ \ldots,\ z_m)$ eine algebraische von den Grössen $\xi$ unabhängige Function der Grössen $z$ ist, welche in Folge ihrer Verzweigungsart sich rational in $\left(\begin{array}{c} m\\ \mu\ (s_\mu,\ z_\mu)\\ 1 \end{array}\right)$ ausdrücken lassen muss. Lässt man nun die Punkte $\eta$ mit den Punkten $\varepsilon$ zusammenfallen, so dass die Grössen $\xi_\mu - z_\mu$ und die Grössen $\sigma_\mu - s_\mu$ sämmtlich unendlich

\newpage
% ============================================================
% Division Title: Zweite Abtheilung
% ============================================================
\thispagestyle{empty}
\vspace*{\fill}
\begin{center}
{\Huge\scshape Zweite Abtheilung.}
\end{center}
\vspace*{\fill}
\clearpage
\thispagestyle{empty}
\clearpage

% ============================================================
% Paper XII: Trigonometric Series
% ============================================================
\section*{XII. Ueber die Darstellbarkeit einer Function durch eine trigonometrische Reihe.}
\addcontentsline{toc}{section}{XII. Ueber die Darstellbarkeit einer Function durch eine trigonometrische Reihe}

\begin{center}
(Aus dem dreizehnten Bande der Abhandlungen der Königlichen Gesellschaft der\\
Wissenschaften zu Göttingen.\textsuperscript{*})
\end{center}

Der folgende Aufsatz über die trigonometrischen Reihen besteht aus zwei wesentlich verschiedenen Theilen. Der erste Theil enthält eine Geschichte der Untersuchungen und Ansichten über die willkürlichen (graphisch gegebenen) Functionen und ihre Darstellbarkeit durch trigonometrische Reihen. Bei ihrer Zusammenstellung war es mir vergönnt, einige Wink e des berühmten Mathematikers zu benutzen, welchem man die erste gründliche Arbeit über diesen Gegenstand verdankt. Im zweiten Theile liefere ich über die Darstellbarkeit einer Function durch eine trigonometrische Reihe eine Untersuchung, welche auch die bis jetzt noch unerledigten Fälle umfasst. Es war nöthig, ihr einen kurzen Aufsatz über den Begriff eines bestimmten Integrales und den Umfang seiner Gültigkeit vorauszuschicken.

\begingroup
\renewcommand{\thefootnote}{}
\footnotetext{\textsuperscript{*} Diese Abhandlung ist im Jahre 1854 von dem Verfasser behufs seiner Habilitation an der Universität zu Göttingen der philosophischen Facultät eingereicht. Wiewohl der Verfasser ihre Veröffentlichung, wie es scheint, nicht beabsichtigt hat, so wird doch die hiermit erfolgende Herausgabe derselben in gänzlich ungeänderter Form sowohl durch das hohe Interesse des Gegenstandes an sich als durch die in ihr niedergelegte Behandlungsweise der wichtigsten Principien der Infinitesimal-Analysis wohl hinlänglich gerechtfertigt erscheinen.

Braunschweig, im Juli 1867.\hfill R.\ Dedekind.}
\endgroup

\begin{center}
\textbf{Geschichte der Frage über die Darstellbarkeit einer willkürlich\\gegebenen Function durch eine trigonometrische Reihe.}
\end{center}

\subsection*{1.}
\addcontentsline{toc}{subsection}{Art. 1 (Geschichte)}

Die von \textit{Fourier} so genannten trigonometrischen Reihen, d.~h. die Reihen von der Form
\begin{multline*}
a_1\sin x + a_2\sin 2x + a_3\sin 3x + \cdots\\
+ \tfrac{1}{2}b_0 + b_1\cos x + b_2\cos 2x + b_3\cos 3x + \cdots
\end{multline*}
spielen in demjenigen Theile der Mathematik, wo ganz willkürliche Functionen vorkommen, eine bedeutende Rolle; ja, es lässt sich mit Grund behaupten, dass die wesentlichsten Fortschritte in diesem für die Physik so wichtigen Theile der Mathematik von der klareren Einsicht in die Natur dieser Reihen abhängig gewesen sind. Schon gleich bei den ersten mathematischen Untersuchungen, die auf die Betrachtung willkürlicher Functionen führten, kam die Frage zur Sprache, ob sich eine solche ganz willkürliche Function durch eine Reihe von obiger Form ausdrücken lasse.

Es geschah dies in der Mitte des vorigen Jahrhunderts bei Gelegenheit der Untersuchungen über die schwingenden Saiten, mit welchen sich damals die berühmtesten Mathematiker beschäftigten. Ihre Ansichten über unsern Gegenstand lassen sich nicht wohl darstellen, ohne auf dieses Problem einzugehen.

Unter gewissen Voraussetzungen, die in der Wirklichkeit näherungsweise zutreffen, wird bekanntlich die Form einer gespannten in einer Ebene schwingenden Saite, wenn $x$ die Entfernung eines unbestimmten ihrer Punkte von ihrem Anfangspunkte, $y$ seine Entfernung aus der Ruhelage zur Zeit $t$ bedeutet, durch die partielle Differentialgleichung
\[ \frac{\partial^2 y}{\partial t^2} = \alpha\alpha\,\frac{\partial^2 y}{\partial x^2} \]
bestimmt, wo $\alpha$ von $t$ und bei einer überall gleich dicken Saite von $x$ unabhängig ist.

Der erste, welcher eine allgemeine Lösung dieser Differentialgleichung gab, war d'\textit{Alembert}.

Er zeigte\footnote{\textit{Mémoires de l'académie de Berlin}. 1747. pag.~214.}), dass jede Function von $x$ und $t$, welche für $y$ gesetzt, die Gleichung zu einer identischen macht, in der Form
\[ f(x + \alpha t) + \varphi(x - \alpha t) \]
enthalten sein müsse, wie sich dies durch Einführung der unabhängig veränderlichen Grössen $x + \alpha t$, $x - \alpha t$ anstatt $x$, $t$ ergiebt, wodurch
\[ \frac{\partial^2 y}{\partial x^2} - \frac{1}{\alpha\alpha}\frac{\partial^2 y}{\partial t^2} \quad\text{in}\quad 4\frac{\partial\frac{\partial y}{\partial(x+\alpha t)}}{\partial(x-\alpha t)} \]
übergeht.

Ausser dieser partiellen Differentialgleichung, welche sich aus den allgemeinen Bewegungsgesetzen ergiebt, muss nun $y$ noch die Bedingung erfüllen, in den Befestigungspunkten der Saite stets $= 0$ zu sein; man hat also, wenn in dem einen dieser Punkte $x = 0$, in dem andern $x = l$ ist,
\[ f(\alpha t) = -\varphi(-\alpha t),\quad f(l + \alpha t) = -\varphi(l - \alpha t) \]
und folglich
\begin{align*}
f(z) &= -\varphi(-z) = -\varphi(l - (l+z)) = f(2l + z),\\
y &= f(\alpha t + x) - f(\alpha t - x).
\end{align*}

Nachdem d'\textit{Alembert} dies für die allgemeine Lösung des Problems geleistet hatte, beschäftigte er sich in einer Fortsetzung\footnote{Ibid. pag.~220.} seiner Abhandlung mit der Gleichung $f(z) = f(2l+z)$; d.~h. er sucht analytische Ausdrücke, welche unverändert bleiben, wenn $z$ um $2l$ wächst.

Es war ein wesentliches Verdienst \textit{Euler}'s, der im folgenden Jahrgange der Berliner Abhandlungen\footnote{Mémoires de l'académie de Berlin. 1748. pag.~69.} eine neue Darstellung dieser d'\textit{Alembert}'schen Arbeiten gab, dass er das Wesen der Bedingungen, welchen die Function $f(z)$ genügen muss, richtiger erkannte. Er bemerkte, dass der Natur des Problems nach die Bewegung der Saite vollständig bestimmt sei, wenn für irgend einen Zeitpunkt die Form der Saite und die Geschwindigkeit jedes Punktes (also $y$ und $\frac{\partial y}{\partial t}$) gegeben seien, und zeigte, dass sich, wenn man diese beiden Functionen sich durch willkürlich gezogene Curven bestimmt denkt, daraus stets durch eine einfache geometrische Construction die d'\textit{Alembert}'sche Function $f(z)$ finden lässt. In der That, nimmt man an, dass für
\[ t = 0,\ y = g(x)\ \text{und}\ \frac{\partial y}{\partial t} = h(x) \]
sei, so erhält man für die Werthe von $x$ zwischen $0$ und $l$
\[ f(x) - f(-x) = g(x),\quad f(x) + f(-x) = \frac{1}{\alpha}\int h(x)\,dx \]
und folglich die Function $f(z)$ zwischen $-l$ und $l$; hieraus aber ergiebt sich ihr Werth für jeden andern Werth von $z$ vermittelst der Gleichung
\[ f(z) = f(2l + z). \]
Dies ist in abstracten, aber jetzt allgemein geläufigen Begriffen dargestellt, die \textit{Euler}'sche Bestimmung der Function $f(z)$.

Gegen diese Ausdehnung seiner Methode durch \textit{Euler} verwahrte sich indess d'\textit{Alembert} sofort\footnote{Mémoires de l'académie de Berlin. 1750. pag.~358. En effet on ne peut ce me semble exprimer $y$ analytiquement d'une manière plus générale, qu'en la supposant une fonction de $t$ et de $x$. Mais dans cette supposition on ne trouve la solution du problème que pour les cas où les différentes figures de la corde vibrante peuvent être renfermées dans une seule et même équation.}, weil seine Methode nothwendig voraussetze, dass $y$ sich in $t$ und $x$ analytisch ausdrücken lasse.

Ehe eine Antwort \textit{Euler}'s hierauf erfolgte, erschien eine dritte von diesen beiden ganz verschiedene Behandlung dieses Gegenstandes von \textit{Daniel Bernoulli}\footnote{Mémoires de l'académie de Berlin. 1753. p.~147.}. Schon vor d'\textit{Alembert} hatte \textit{Taylor}\footnote{Taylor de methodo incrementorum.} gesehen, dass $\frac{\partial^2 y}{\partial t^2} = \alpha\alpha\frac{\partial^2 y}{\partial x^2}$ und zugleich $y$ für $x = 0$ und für $x = l$ stets gleich $0$ sei, wenn man $y = \sin\frac{n\pi x}{l}\cos\frac{n\pi\alpha t}{l}$ und hierin für $n$ eine ganze Zahl setze. Er erklärte hieraus die physikalische Thatsache, dass eine Saite ausser ihrem Grundtone auch den Grundton einer $\frac{1}{2}$, $\frac{1}{3}$, $\frac{1}{4}$, $\ldots$ so langen (übrigens ebenso beschaffenen) Saite geben könne, und hielt seine particuläre Lösung für allgemein, d.~h. er glaubte, die Schwingung der Saite würde stets, wenn die ganze Zahl $n$ der Höhe des Tons gemäss bestimmt würde, wenigstens sehr nahe durch die Gleichung ausgedrückt. Die Beobachtung, dass eine Saite ihre verschiedenen Töne gleichzeitig geben könne, führte nun \textit{Bernoulli} zu der Bemerkung, dass die Saite (der Theorie nach) auch der Gleichung
\[ y = \Sigma a_n\sin\frac{n\pi x}{l}\cos\frac{n\pi\alpha}{l}(t - \beta_n) \]
gemäss schwingen könne, und weil sich aus dieser Gleichung alle beobachteten Modificationen der Erscheinung erklären liessen, so hielt er sie für die allgemeinste\footnote{l.~c. p.~157. art.~XIII.}. Um diese Ansicht zu stützen, untersuchte er die Schwingungen eines masselosen gespannten Fadens, der in einzelnen Punkten mit endlichen Massen beschwert ist, und zeigte, dass die Schwingungen desselben stets in eine der Zahl der Punkte gleiche Anzahl von solchen Schwingungen zerlegt werden kann, deren jede für alle Massen gleich lange dauert.

Diese Arbeiten \textit{Bernoulli}'s veranlassten einen neuen Aufsatz \textit{Euler}'s, welcher unmittelbar nach ihnen unter den Abhandlungen der Berliner Akademie abgedruckt ist\footnote{Mémoires de l'académie de Berlin. 1753. p.~196.}. Er hält darin d'\textit{Alembert} gegenüber fest\footnote{l.~c. p.~214.}, dass die Function $f(z)$ eine zwischen den Grenzen $-l$ und $l$ ganz willkürliche sein könne, und bemerkt\footnote{l.~c. art.~III---X.}, dass \textit{Bernoulli}'s Lösung (welche er schon früher als eine besondere aufgestellt hatte) dann allgemein sei und zwar nur dann allgemein sei, wenn die Reihe
\begin{multline*}
a_1\sin\frac{x\pi}{l} + a_2\sin\frac{2x\pi}{l} + \cdots\\
+ \tfrac{1}{2}b_0 + b_1\cos\frac{x\pi}{l} + b_2\cos\frac{2x\pi}{l} + \cdots
\end{multline*}
für die Abscisse $x$ die Ordinate einer zwischen den Abscissen $0$ und $l$ ganz willkürlichen Curve darstellen könne. Nun wurde es damals von Niemand bezweifelt, dass alle Umformungen, welche man mit einem analytischen Ausdrucke --- er sei endlich oder unendlich --- vornehmen könne, für jedwede Werthe der unbestimmten Grössen gültig seien oder doch nur in ganz speciellen Fällen unanwendbar würden. Es schien daher unmöglich, eine algebraische Curve oder überhaupt eine analytisch gegebene nicht periodische Curve durch obigen Ausdruck darzustellen, und \textit{Euler} glaubte daher, die Frage gegen \textit{Bernoulli} entscheiden zu müssen.

Der Streit zwischen \textit{Euler} und d'\textit{Alembert} war indess noch immer unerledigt. Dies veranlasste einen jungen, damals noch wenig bekannten Mathematiker, \textit{Lagrange}, die Lösung der Aufgabe auf einem ganz neuen Wege zu versuchen, auf welchem er zu \textit{Euler}'s Resultaten gelangte. Er unternahm es\footnote{Miscellanea Taurinensia. Tom.~I. Recherches sur la nature et la propagation du son.}, die Schwingungen eines masselosen Fadens zu bestimmen, welcher mit einer endlichen unbestimmten Anzahl gleich grosser Massen in gleich grossen Abständen beschwert ist, und untersuchte dann, wie sich diese Schwingungen ändern, wenn die Anzahl der Massen in's Unendliche wächst. Mit welcher Gewandtheit, mit welchem Aufwande analytischer Kunstgriffe er aber auch den ersten Theil dieser Untersuchung durchführte, so liess der Uebergang vom Endlichen zum Unendlichen doch viel zu wünschen übrig, so dass d'\textit{Alembert} in einer Schrift, welche er an die Spitze seiner opuscules mathématiques stellte, fortfahren konnte, seiner Lösung den Ruhm der grössten Allgemeinheit zu vindiciren. Die Ansichten der damaligen berühmten Mathematiker waren und blieben daher in dieser Sache getheilt; denn auch in spätern Arbeiten behielt jeder im Wesentlichen seinen Standpunkt bei.

Um also schliesslich ihre bei Gelegenheit dieses Problems entwickelten Ansichten über die willkürlichen Functionen und über die Darstellbarkeit derselben durch eine trigonometrische Reihe zusammenzustellen, so hatte \textit{Euler} zuerst diese Functionen in die Analysis eingeführt und, auf geometrische Anschauung gestützt, die Infinitesimalrechnung auf sie angewandt. \textit{Lagrange}\footnote{Miscellanea Taurinensia. Tom.~II. Pars math. pag.~18.} hielt \textit{Euler}'s Resultate (eine geometrische Construction des Schwingungsverlaufs) für richtig; aber ihm genügte die \textit{Euler}'sche geometrische Behandlung dieser

\newpage
% ============================================================
% Continue Paper XII
% ============================================================

Functionen nicht. D'\textit{Alembert}\footnote{Opuscules mathématiques p.\ d'Alembert. Tome premier. 1761. pag.~16. art.~VII---XX.} dagegen ging auf die \textit{Euler}'sche Auffassungsweise der Differentialgleichung ein und beschränkte sich, die Richtigkeit seiner Resultate anzufechten, weil man bei ganz willkürlichen Functionen nicht wissen könne, ob ihre Differentialquotienten stetig seien. Was die \textit{Bernoulli}'sche Lösung betraf, so kamen alle drei darin überein, sie nicht für allgemein zu halten; aber während d'\textit{Alembert}\footnote{Opuscules mathématiques. Tome~I. pag.~42. art.~XXIV.}, um \textit{Bernoulli}'s Lösung für minder allgemein, als die seinige, erklären zu können, behaupten musste, dass auch eine analytisch gegebene periodische Function sich nicht immer durch eine trigonometrische Reihe darstellen lasse, glaubte \textit{Lagrange}\footnote{Misc.\ Taur.\ Tom.~III. Pars math. pag.~221. art.~XXV.} diese Möglichkeit beweisen zu können.

\subsection*{2.}
\addcontentsline{toc}{subsection}{Art. 2 (Geschichte, Forts.)}

Fast fünfzig Jahre vergingen, ohne dass in der Frage über die analytische Darstellbarkeit willkürlicher Functionen ein wesentlicher Fortschritt gemacht wurde. Da warf eine Bemerkung \textit{Fourier}'s ein neues Licht auf diesen Gegenstand; eine neue Epoche in der Entwicklung dieses Theils der Mathematik begann, die sich bald auch äusserlich in grossartigen Erweiterungen der mathematischen Physik kund that. \textit{Fourier} bemerkte, dass in der trigonometrischen Reihe
\[ f(x) = \begin{cases} &a_1\sin x + a_2\sin 2x + \cdots\\ + &\tfrac{1}{2}b_0 + b_1\cos x + b_2\cos 2x + \cdots \end{cases} \]
die Coefficienten sich durch die Formeln
\[ a_n = \frac{1}{\pi}\int\limits_{-\pi}^{\pi} f(x)\sin nx\,dx,\quad b_n = \frac{1}{\pi}\int\limits_{-\pi}^{\pi} f(x)\cos nx\,dx \]
bestimmen lassen. Er sah, dass diese Bestimmungsweise auch anwendbar bleibe, wenn die Function $f(x)$ ganz willkürlich gegeben sei; er setzte für $f(x)$ eine so genannte discontinuirliche Function (die Ordinate einer gebrochenen Linie für die Abscisse $x$) und erhielt so eine Reihe, welche in der That stets den Werth der Function gab.

Als \textit{Fourier} in einer seiner ersten Arbeiten über die Wärme, welche er der französischen Akademie vorlegte\footnote{Bulletin des sciences p.~la soc.~philomatique. Tome~I. p.~112.}, (21.~Dec.\ 1807) zuerst den Satz aussprach, dass eine ganz willkürlich (graphisch) gegebene Function sich durch eine trigonometrische Reihe ausdrücken lasse,

\newpage

war diese Behauptung dem greisen \textit{Lagrange} so unerwartet, dass er ihr auf das Entschiedenste entgegentrat. Es soll\footnote{Nach einer mündlichen Mittheilung des Herrn Professor Dirichlet.} sich hierüber noch ein Schriftstück im Archiv der Pariser Akademie befinden. Dessennungeachtet verweist\footnote{Unter Andern in dem verbreiteten Traité de mécanique Nro.~323, p.~638.} \textit{Poisson} überall, wo er sich der trigonometrischen Reihen zur Darstellung willkürlicher Functionen bedient, auf eine Stelle in \textit{Lagrange}'s Arbeiten über die schwingenden Saiten, wo sich diese Darstellungsweise finden soll. Um diese Behauptung, die sich nur aus der bekannten Rivalität zwischen \textit{Fourier} und \textit{Poisson} erklären lässt\footnote{Der Bericht im bulletin des sciences über die von Fourier der Akademie vorgelegte Abhandlung ist von Poisson.}, zu widerlegen, sehen wir uns genöthigt, noch einmal auf die Abhandlung \textit{Lagrange}'s zurückzukommen; denn über jenen Vorgang in der Akademie findet sich nichts veröffentlicht.

Man findet in der That an der von \textit{Poisson} citirten Stelle\footnote{Misc.\ Taur.\ Tom.~III.\ Pars math.\ pag.~261.} die Formel:
\begin{multline*}
y = 2\!\int\! Y\sin X\pi\,dX \times \sin x\pi + 2\!\int\! Y\sin 2X\pi\,dX \times \sin 2x\pi\\
+\, 2\!\int\! Y\sin 3X\pi\,dX \times \sin 3x\pi + \text{etc.} + 2\!\int\! Y\sin nX\pi\,dX \times \sin nx\pi,
\end{multline*}
\textit{de sorte que, lorsque $x = X$, on aura $y = Y$, $Y$ étant l'ordonnée qui répond à l'abscisse $X$.}

Diese Formel sieht nun allerdings ganz so aus wie die \textit{Fourier}'sche Reihe, so dass bei flüchtiger Ansicht eine Verwechslung leicht möglich ist; aber dieser Schein rührt bloss daher, weil \textit{Lagrange} das Zeichen $\int\!dX$ anwandte, wo er heute das Zeichen $\Sigma\Delta X$ angewandt haben würde. Sie giebt die Lösung der Aufgabe, die endliche Sinusreihe
\[ a_1\sin x\pi + a_2\sin 2x\pi + \cdots + a_n\sin nx\pi \]
so zu bestimmen, dass sie für die Werthe
\[ \frac{1}{n+1},\ \frac{2}{n+1},\ \ldots,\ \frac{n}{n+1} \]
von $x$, welche \textit{Lagrange} unbestimmt durch $X$ bezeichnet, gegebene Werthe erhält. Hätte \textit{Lagrange} in dieser Formel $n$ unendlich gross werden lassen, so wäre er allerdings zu dem \textit{Fourier}'schen Resultat gelangt. Wenn man aber seine Abhandlung durchliest, so sieht man, dass er weit davon entfernt ist zu glauben, eine ganz willkürliche Function lasse sich wirklich durch eine unendliche Sinusreihe darstellen. Er hatte vielmehr die ganze Arbeit gerade unternommen, weil er glaubte, diese willkürlichen Functionen liessen sich nicht durch eine

\newpage

Formel ausdrücken, und von der trigonometrischen Reihe glaubte er, dass sie jede analytisch gegebene periodische Function darstellen könne. Freilich erscheint es uns jetzt kaum denkbar, dass \textit{Lagrange} von seiner Summenformel nicht zur \textit{Fourier}'schen Reihe gelangt sein sollte; aber dies erklärt sich daraus, dass durch den Streit zwischen \textit{Euler} und d'\textit{Alembert} sich bei ihm im Voraus eine bestimmte Ansicht über den einzuschlagenden Weg gebildet hatte. Er glaubte das Schwingungsproblem für eine unbestimmte endliche Anzahl von Massen erst vollständig absolviren zu müssen, bevor er seine Grenzbetrachtungen anwandte. Diese erfordern eine ziemlich ausgedehnte Untersuchung\footnote{Misc.\ Taur.\ Tom.~III.\ Pars math.\ S.~251.}, welche unnöthig war, wenn er die \textit{Fourier}'sche Reihe kannte.

Durch \textit{Fourier} war nun zwar die Natur der trigonometrischen Reihen vollkommen richtig erkannt\footnote{Bulletin d.\ sc.\ Tom.~I.\ p.~115. Les coefficients $a$, $a'$, $a''$, $\ldots$, étant ainsi déterminés etc.}; sie wurden seitdem in der mathematischen Physik zur Darstellung willkürlicher Functionen vielfach angewandt, und in jedem einzelnen Falle überzeugte man sich leicht, dass die \textit{Fourier}'sche Reihe wirklich gegen den Werth der Function convergire; aber es dauerte lange, ehe dieser wichtige Satz allgemein bewiesen wurde.

Der Beweis, welchen \textit{Cauchy} in einer der Pariser Akademie am 27.~Febr.\ 1826 vorgelesenen Abhandlung gab\footnote{Mémoires de l'ac.\ d.\ sc.\ de Paris.\ Tom.~VI.\ p.~603.}, ist unzureichend, wie \textit{Dirichlet} gezeigt hat\footnote{Crelle Journal für die Mathematik. Bd.~IV. p.~157 \& 158.}. \textit{Cauchy} setzt voraus, dass, wenn man in der willkürlich gegebenen periodischen Function $f(x)$ für $x$ ein complexes Argument $x + yi$ setzt, diese Function für jeden Werth von $y$ endlich sei. Dies findet aber nur statt, wenn die Function gleich einer constanten Grösse ist. Man sieht indess leicht, dass diese Voraussetzung für die ferneren Schlüsse nicht nothwendig ist. Es reicht hin, wenn eine Function $\varphi(x + yi)$ vorhanden ist, welche für alle positiven Werthe von $y$ endlich ist und deren reeller Theil für $y = 0$ der gegebenen periodischen Function $f(x)$ gleich wird. Will man diesen Satz, der in der That richtig ist\footnote{Der Beweis findet sich in der Inauguraldissertation des Verfassers.}, voraussetzen, so führt allerdings der von \textit{Cauchy} eingeschlagene Weg zum Ziele, wie umgekehrt dieser Satz sich aus der \textit{Fourier}'schen Reihe ableiten lässt.

\clearpage
\chapter*{Source segment 6: riemann pages 251 300}
\addcontentsline{toc}{chapter}{Source segment 6: riemann pages 251 300}
\selectlanguage{ngerman}
\maketitle
\thispagestyle{fancy}
\vspace{-0.5cm}
\hrule
\vspace{0.8cm}






% ===== MAIN CONTENT =====

% --- Page 251 ---
Erst  im  Januar  1829  erschien  im  Journal  von  Grelle$^{*)}$  eine  Ab-
handlung von  Dirichlet,  worin  für  Functionen,  die  durchgehends  eine
Integration  zulassen  und  nicht  unendlich  viele  Maxima  und  Minima
haben,  die  Frage  ihrer  Darstellbarkeit  durch  trigonometrische  Reihen
in  aller  Strenge  entschieden  wurde.
Die  Erkenntniss  des  zur  Lösung  dieser  Aufgabe  einzuschlagenden
Weges  ergab  sich  ihm  aus  der  Einsicht,  dass  die  unendlichen  Reihen
in  zwei  wesentlich  verschiedene  Klassen  zerfallen,  je  nachdem  sie,
wenn  man  sämmtliche  Glieder  positiv  macht,  convergent  bleiben  oder
nicht.  In  den  ersteren  können  die  Glieder  beliebig  versetzt  werden,
der  Werth  der  letzteren  dagegen  ist  von  der  Ordnung  der  Glieder  ab-
hängig. In  der  That,  bezeichnet  man  in  einer  Reihe  zweiter  Klasse
die  positiven  Glieder  der  Reihe  nach  durch
die  negativen  durch
so  ist  klar,  dass  sowohl  Ua,  als  Üb  unendlich  sein  müssen;  denn
wären  beide  endlich,  so  würde  die  Reihe  auch  nach  Gleichmachung
der  Zeichen  convergiren;  wäre  aber  eine  unendlich,  so  würde  die  Reihe
divergiren.  Offenbar  kann  nun  die  Reihe  durch  geeignete  Anordnung
der  Glieder  einen  beliebig  gegebenen  Werth  C  erhalten.  Denn  nimmt
man  abwechselnd  so  lange  positive  Glieder  der  Reihe,  bis  ihr  Werth
grösser  als  C  wird,  und  so  lange  negative,  bis  ihr  Werth  kleiner  als
G  wird,  so  wird  die  Abweichung  von  C  nie  mehr  betragen,  als  der
Werth  des  dem  letzten  Zeichenwechsel  voraufgehenden  Gliedes.  Da
nun  sowohl  die  Grössen  a,  als  die  Grössen  h  mit  wachsendem  Index
zuletzt  unendlich  klein  werden,  so  werden  auch  die  Abweichungen
von  Cf  wenn  man  in  der  Reihe  nur  hinreichend  weit  fortgeht,  be-
liebig klein  werden,  d.  h.  die  Reihe  wird  gegen  C  convergiren.
Nur  auf  die  Reihen  erster  Klasse  sind  die  Gesetze  endlicher  Sum-
men anwendbar;  nur  sie  können  wirklich  als  Inbegriff  ihrer  Glieder
betrachtet  werden,  die  Reihen  der  zweiten  Klasse  nicht;  ein  Umstand,
welcher  von  den  Mathematikern  des  vorigen  Jahrhunderts  übersehen
wurde,  hauptsächlich  wohl  aus  dem  Grunde,  weil  die  Reihen,  welche
nach  steigenden  Potenzen  einer  veränderlichen  Grösse  fortschreiten,  all-
gemein zu  reden  (d.  h.  einzelne  Werthe  dieser  Grösse  ausgenommen),
zur  ersten  Klasse  gehören.
\textsuperscript{$\diamondsuit$})  Bd.  IV.  pag.  157.
\newpage

% --- Page 252 ---
Die  Fourier'sclie  Reihe  gehört  nun  offenbar  nicht  noth wendig
zur  ersten  Klasse;  ihre  Convergenz  konnte  also  gar  nicht,  wie  Cauchy
vergeblich$^{*)}$  versucht  hatte,  aus  dem  Gesetze,  nach  welchem  die  Glieder
abnehmen,  abgeleitet  werden.  Es  musste  vielmehr  gezeigt  werden,  dass
die  endliche  Reihe
\[
\frac{1}{\pi} \int_{-\pi}^{\pi} f(\alpha)\, d\alpha \sin x + \frac{1}{\pi} \int_{-\pi}^{\pi} f(\alpha) \sin 2\alpha\, d\alpha \sin 2x + \cdots
\]
\[
+ \frac{1}{\pi} \int_{-\pi}^{\pi} f(\alpha) \sin n\alpha\, d\alpha \sin nx
\]
\[
+ \frac{1}{\pi} \int_{-\pi}^{\pi} f(\alpha)\, d\alpha \cos x + \frac{1}{\pi} \int_{-\pi}^{\pi} f(\alpha) \cos 2\alpha\, d\alpha \cos 2x + \cdots
\]
\[
+ \frac{1}{\pi} \int_{-\pi}^{\pi} f(\alpha) \cos n\alpha\, d\alpha \cos nx,
\]
oder,  was  dasselbe  ist,  das  Integral
\[
\frac{1}{\pi} \int_{-\pi}^{\pi} f(\alpha) \frac{\sin \frac{2n+1}{2}(x-\alpha)}{\sin \frac{x-\alpha}{2}}\, d\alpha
\]
sich,  wenn  n  in's  Unendliche  wächst,  dem  Werthe  $$f(x)$$  unendlich  an-
nähert.
Dirichlet  stützt  diesen  Beweis  auf  die  beiden  Sätze:
\begin{enumerate}
\item Wenn $0 < c \leq \frac{\pi}{2}$, nähert sich $\int_0^c \varphi(\beta) \frac{\sin(2n+1)\beta}{\sin\beta}\, d\beta$ mit wachsendem $n$ zuletzt unendlich dem Werth $\frac{\pi}{2}\varphi(0)$;
\item wenn $0 < \delta < c \leq \frac{\pi}{2}$, nähert sich $\int_\delta^c \varphi(\beta) \frac{\sin(2n+1)\beta}{\sin\beta}\, d\beta$ mit wachsendem $n$ zuletzt unendlich dem Werth $0$,
\end{enumerate}
vorausgesetzt, dass die Function $\varphi(\beta)$ zwischen den Grenzen dieser Integrale entweder immer abnimmt, oder immer zunimmt.
Mit  Hülfe  dieser  beiden  Sätze  lässt  sich,  wenn  die  Function  f
nicht  unendlich  oft  vom  Zunehmen  zum  Abnehmen  oder  vom  Ab-
nehmen zum  Zunehmen  übergeht,  das  Integral
$^{*)}$  Dirichlet   in   Crelle's  Journal.     Bd.  IV.  pag.  158.     Quoi  qu'il  en  soit  de
cette  premiöre  Observation,  .  .  .  ä  mesure  que  n  croit.
\newpage

% --- Page 253 ---
\[
\frac{1}{\pi} \int_{-\pi}^{\pi} f(\alpha) \frac{\sin \frac{2n+1}{2}(x-\alpha)}{\sin \frac{x-\alpha}{2}}\, d\alpha
\]
offenbar  in  eine  endliche  Anzahl  von  Gliedern  zerlegen,  von  denen
eins$^{*)}$  gegen $\frac{1}{2}f(x+0)$,  ein  anderes  gegen  $\frac{1}{2}f(x-0)$,  die  übrigen
aber  gegen  0  convergiren,  wenn  n  ins  Unendliche  wächst.
Hieraus  folgt,  dass  durch  eine  trigonometrische  Reihe  jede  sich
nach  dem  Intervall  27t  periodisch  wiederholende  Function  darstellbar
ist,  welche
1)  durchgehends  eine  Integration  zulässt,
2)  nicht  unendlich  viele  Maxima  und  Minima  hat  und
3)  wo  ihr  Werth  sich  sprungweise  ändert,  den  Mittel werth  zwi-
schen den  beiderseitigen  Grenzwerthen  annimmt.
Eine  Function,  welche  die  ersten  beiden  Eigenschaften  hat,  die
dritte  aber  nicht,  kann  durch  eine  trigonometrische  Reihe  offenbar
nicht  dargestellt  werden;  denn  die  trigonometrische  Reihe,  die  sie
ausser  den  Unstetigkeiten  darstellt,  würde  in  den  Unstetigkeitspunkten
selbst  von  ihr  abweichen.  Ob  und  wann  aber  eine  Function,  welche
die  ersten  beiden  Bedingungen  nicht  erfüllt,  durch  eine  trigonometrische
Reihe  darstellbar  sei,  bleibt  durch  diese  Untersuchung  unentschieden.
Durch  die  Arbeit  Dirichlet's  ward  einer  grossen  Menge  wich-
tiger analytischer  Untersuchungen  eine  feste  Grundlage  gegeben.  Es
war  ihm  gelungen,  indem  er  den  Punkt,  wo  Euler  irrte,  in  volles
Licht  brachte,  eine  Frage  zu  erledigen,  die  so  viele  ausgezeichnete
Mathematiker  seit  mehr  als  siebzig  Jahren  (seit  dem  Jahre  1753)  be-
schäftigt hatte.  In  der  That  für  alle  Fälle  der  Natur,  um  welche  es
sich  allein  handelte,  war  sie  vollkommen  erledigt,  denn  so  gross  auch
unsere  Unwissenheit  darüber  ist,  wie  sich  die  Kräfte  und  Zustände  der
Materie  nach  Ort  und  Zeit  im  Unendlichkleinen  ändern,  so  können
wir  doch  sicher  annehmen,  dass  die  Functionen,  auf  welche  sich  die
Dirichlet'sQhe  Untersuchung  nicht  erstreckt,  in  der  Natur  nicht  vor-
kommen.
Dessenungeachtet  scheinen  diese  von  Dirichlet  unerledigten  Fälle
aus  einem  zweifachen  Grunde  Beachtung  zu  verdienen.
$^{*)}$  Es  ist  nicht  schwer  zu  beweisen,  dass  der  Werth  einer  Function  f,  welche
nicht  unendlich  viele  Maxima  und  Minima  hat,  stets,  sowohl  wenn  der  Argument-
werth  abnehmend,  als  wenn  er  zunehmend  gleich  x  wird,  entweder  festen  Grenz-
werthen f{x  -j-  0)  und  f{x  ---  0)  (nach  Dirichlet's  Bezeichnung  in  Dove's  Reper-
torium  der  Physik.  Bd.  1.  pag.  170)  sich  nähern,  oder  unendlich  gross  werden
müsse.  (')
\newpage

% --- Page 254 ---
238  Xir.    Ueber  die  Darstellbarkeit  einer  Function
Erstlich  steht,  wie  Dirichlet  selbst  am  Schlüsse  seiner  Abhand-
lung bemerkt,  dieser  Gegenstand  mit  den  Principien  der  Infinitesimal-
rechnung in  der  engsten  Verbindung  und  kann  dazu  dienen,  diese
Principien  zu  grösserer  Klarheit  und  Bestimmtheit;  zu  bringen.  In
dieser  Beziehung  hat  die  Behandlung  desselben  ein  unmittelbares
Interesse.
Zweitens  aber  ist  die  Anwendbarkeit  der  Fourier'schen  Reihen
nicht  auf  physikalische  Untersuchungen  beschränkt;  sie  ist  jetzt  auch
in  einem  Gebiete  der  reinen  Mathematik,  der  Zahlentheorie,  mit  Er-
folg angewandt,  und  hier  scheinen  gerade  diejenigen  Functionen,  deren
Darstellbarkeit  durch  eine  trigonometrische  Reihe  Dirichlet  nicht
untersucht  hat,  von  Wichtigkeit  zu  sein.
Am  Schlüsse  seiner  Abhandlung  verspricht  freilich  Dirichlet,
später  auf  diese  Fälle  zurückzukommen,  aber  dieses  Versprechen  ist
bis  jetzt  unerfüllt  geblieben.  Auch  die  Arbeiten  von  Dirksen  und
Bessel  über  die  Cosinus-  und  Sinusreihen  leisten  diese  Ergänzung
nicht;  sie  stehen  vielmehr  der  Dirichlet'schen  an  Strenge  und  All-
gemeinheit nach.  Der  mit  ihr  fast  ganz  gleichzeitige  Aufsatz  D  i rk  s  e n'  s  $^{*)}$ ,
welcher  offenbar  ohne  Kenntniss  derselben  geschrieben  ist,  schlägt
zwar  im  Allgemeinen  einen  richtigen  Weg  ein,  enthält  aber  im  Ein-
zelnen einige  Ungenauigkeiten.  Denn  abgesehen  davon,  dass  er  in
einem  speciellen  Falle$^{\dagger)}$  für  die  Summe  der  Reihe ein  falsches  Re-
sultat findet,  stützt  er  sich  in  einer  Nebenbetrachtung  auf  eine  nur  in
besonderen  Fällen  mögliche  Reihenentwicklung$^{\ddagger)}$,  so  dass  sein  Be-
weis nur  für  Functionen  mit  überall  endlichen  ersten  Differential-
quotienten vollständig  ist.  Bessel$^{\S)}$  sucht  den  Dirichlet'schen  Be-
weis zu  vereinfachen.  Aber  die  Aenderungen  in  diesem  Beweise  ge-
währen keine  wesentliche  Vereinfachung  in  den  Schlüssen,  sondern
dienen  höchstens  dazu,  ihn  in  geläufigere  Begriffe  zu  kleiden,  während
seine  Strenge  und  Allgemeinheit  beträchtlich  darunter  leidet.
Die  Frage  über  die  Darstellbarkeit  einer  Function  durch  eine  tri-
gonometrische Reihe  ist  also  bis  jetzt  nur  unter  den  beiden  Voraus-
setzungen entschieden,  dass  die  Function  durchgehends  eine  Integration
zulässt  und  nicht  unendlich  viele  Maxima  und  Minima  hat.  Wenn  die
letztere  Voraussetzung  nicht  gemacht  wird,  so  sind  die  beiden  Integral-
theoreme Dirichlet's  zur  Entscheidung  der  Frage  unzulänglich;  wenn
aber  die  erstere  wegfällt,  so  ist  schon  die  Fourier'sche  Coefficienten-
$^{*)}$  Crelle's  Journal.     Bd.  IV.  p.  170.
$^{\dagger)}$  1.  c.  Formel  22.
$^{\ddagger)}$  1.  c.  Art.  3.
$^{\S)}$  Schumacher.    Astronomische  Nachrichten.     Nro.  374  (Bd    16.  p.  229)
\newpage

% --- Page 255 ---
bestimniung  nicht  anwendbar.  Der  im  Folgenden,  wo  diese  Frage
ohne  besondere  Voraussetzungen  über  die  Natur  der  Function  unter-
sucht werden  soll,  eingeschlagene  Weg  ist  hierdurch,  wie  man  sehen
wird,  bedingt;  ein  so  directer  Weg,  wie  der  Diriclilet's,  ist  der  Natur
der  Sache  nach  nicht  möglich.
Ueber  den  Begriff  eines  bestimmten  Integrals  und  den  Umfang
seiner  Gültigkeit.
4.
Die  Unbestimmtheit,  welche  noch  in  einigen  Fundamentalpunkten
der  Lehre  von  den  bestimmten  Integralen  herrscht,  nöthigt  uns.
Einiges  voraufzuschicken  über  den  Begriff  eines  bestimmten  Integrals
und  den  Umfang  seiner  Gültigkeit.
$b$
Also  zuerst:  Was  hat  man  unter  $\int_a^b f(x)\, dx$  zu  verstehen?

Um  dieses  festzusetzen,  nehmen  wir  zwischen  $a$  und  $b$  der  Grösse
nach  auf  einander  folgend,  eine  Reihe  von  Werthen  $x_1, x_2, \ldots, x_{n-1}$
an  und  bezeichnen  der  Kürze  wegen  $x_1 - a$  durch  $\delta_1$,  $x_2 - x_1$  durch
$\delta_2$,  $\ldots$,  $b - x_{n-1}$  durch  $\delta_n$  und  durch  $\varepsilon$  einen  positiven  ächten  Bruch.
Es  wird  alsdann  der  Werth  der  Summe
\[ S = \delta_1 f(a + \varepsilon_1 \delta_1) + \delta_2 f(x_1 + \varepsilon_2 \delta_2) + \delta_3 f(x_2 + \varepsilon_3 \delta_3) + \cdots + \delta_n f(x_{n-1} + \varepsilon_n \delta_n) \]
von  der  Wahl  der  Intervalle  $\delta$  und  der  Grössen  $\varepsilon$  abhängen.  Hat  sie
nun  die  Eigenschaft,  wie  auch  $\delta$  und  $\varepsilon$  gewählt  werden  mögen,  sich
einer  festen  Grenze  $A$  unendlich  zu  nähern,  sobald  sämmtliche  $\delta$  un-
endlich  klein  werden,  so  heisst  dieser  Werth  $\int_a^b f(x)\, dx$.

Hat  sie  diese  Eigenschaft  nicht,  so  hat  $\int_a^b f(x)\, dx$  keine  Bedeutung.
Man  hat  jedoch  in  mehreren  Fällen  versucht,  diesem  Zeichen  auch
dann  eine  Bedeutung  beizulegen,  und  unter  diesen  Erweiterungen  des
Begriffs  eines  bestimmten  Integrals  ist  eine  von  allen  Mathematikern
angenommen.  Wenn  nämlich  die  Function  fix)  bei  Annäherung  des
Arguments  an  einen  einzelnen  Werth  c  in  dem  Intervalle  (a,  \&)  un-
endlich gross  wird,  so  kann  offenbar  die  Summe  S,  welchen  Grad
von  Kleinheit  man  auch   den   d  vorschreiben   möge,  jeden  beliebigen
\newpage

% --- Page 256 ---
Werth  erhalten;  sie  hat  also  keinen  Grenzwerth,  und  $\int_a^b f(x)\, dx$  würde
nach  dem  Obigen  keine  Bedeutung  haben.  Wenn  aber  alsdann
\[ \int_a^{c-\alpha_1} f(x)\, dx + \int_{c+\alpha_2}^b f(x)\, dx \]
sich,  wenn  $\alpha_1$  und  $\alpha_2$  unendlich  klein  werden,  einer  festen  Grenze
nähert,  so  versteht  man  unter  $\int_a^b f(x)\, dx$  diesen  Grenzwerth.
Andere  Festsetzungen  von  Cauchy  über  den  Begriff  des  bestimmten
Integrales  in  den  Fällen,  wo  es  dem  Grundbegriffe  nach  ein  solches
nicht  giebt,  mögen  für  einzelne  Klassen  von  Untersuchungen  zweck-
mässig sein;  sie  sind  indess  nicht  allgemein  eingeführt  und  dazu,  schon
wegen  ihrer  grossen  Willkürlichkeit,  wohl  kaum  geeignet.
Untersuchen  wir  jetzt  zweitens  den  Umfang  der  Gültigkeit  dieses
Begriffs  oder  die  Frage:  in  welchen  Fällen  lässt  eine  Function  eine
Integration  zu  und  in  welchen  nicht?
Wir  betrachten  zunächst  den  Integralbegriff  im  engern  Sinne,
d,  h.  wir  setzen  voraus,  dass  die  Summe  S,  wenn  sämmtliche  d  un-
endlich klein  werden,  convergirt.  Bezeichnen  wir  also  die  grösste
Schwankung  der  Function  zwischen  $a$  und  $x_1$,  d.~h.  den  Unterschied
ihres  grössten  und  kleinsten  Werthes  in  diesem  Intervalle,  durch  $D_1$,
zwischen  $x_1$  und  $x_2$  durch  $D_2$,  $\ldots$,  zwischen  $x_{n-1}$  und  $b$  durch  $D_n$,
so  muss
\[ \delta_1 D_1 + \delta_2 D_2 + \cdots + \delta_n D_n \]
mit  den  Grössen  $\delta$  unendlich  klein  werden.  Wir  nehmen  ferner  an,
dass,  so  lange  sämmtliche  $\delta$  kleiner  als  $d$  bleiben,  der  grösste  Werth,
den  diese  Summe  erhalten  kann,  $\Delta$  sei;  $\Delta$  wird  alsdann  eine  Function
von  $d$  sein,  welche  mit  $d$  immer  abnimmt  und  mit  dieser  Grösse  un-
endlich  klein  wird.  Ist  nun  die  Gesammtgrösse  der  Intervalle,  in
welchen  die  Schwankungen  grösser  als  $\sigma$  sind,  $= s$,  so  wird  der  Bei-
trag  dieser  Intervalle  zur  Summe  $\delta_1 D_1 + \delta_2 D_2 + \cdots + \delta_n D_n$  offenbar
$\leq \sigma s$.  Man  hat  daher
\[ \sigma s \leq \delta_1 D_1 + \delta_2 D_2 + \cdots + \delta_n D_n \leq \Delta; \quad \text{folglich} \quad s \leq \frac{\Delta}{\sigma} \]
---  kann  nun,  wenn  tf  gegeben  ist,  immer  durch  geeignete  Wahl  von
\newpage

% --- Page 257 ---
d  beliebig  klein  gemacht   werden j  dasselbe  gilt  daher  von  s,  und  es
ergiebt  sich  also:
Damit  die  Summe  S,  wenn  säramtliche  d  unendlich  klein  werden,
convergirt,  ist  ausser  der  Endlichkeit  der  Function  $$f(x)$$  noch  erfor-
derlich, dass  die  Gesammtgrösse  der  Intervalle,  in  welchen  die
Schwankungen  >  6  sind,  was  auch  6  sei,  durch  geeignete  Wahl  von
d  beliebig  klein  gemacht  werden  kann.
Dieser  Satz  lässt  sich  auch  umkehren:
Wenn  die  Function  $f(x)$  immer  endlich  ist,  und  bei  unendlichem
Abnehmen  sämmtlicher  Grössen  d  die  Gesammtgrösse  s  der  Intervalle,
in  welchen  die  Schwankungen  der  Function  $f(x)$  grösser,  als  eine  ge-
gebene Grösse  6,  sind,  stets  zuletzt  unendlich  klein  wird,  so  conver-
girt  die  Summe  S,  wenn  sämmtliche  d  unendlich  klein  werden.
Denn  diejenigen  Intervalle,  in  welchen  die  Schwankungen  $> \sigma$
sind,  liefern  zur  Summe  $\delta_1 D_1 + \delta_2 D_2 + \cdots + \delta_n D_n$  einen  Beitrag,
kleiner  als  $s$,  multiplicirt  in  die  grösste  Schwankung  der  Function
zwischen  $a$  und  $b$,  welche  (n.~V.)  endlich  ist;  die  übrigen  Intervalle
einen  Beitrag  $< \sigma(b-a)$.  Offenbar  kann  man  nun  erst  $\sigma$  beliebig
klein  annehmen  und  dann  immer  noch  die  Grösse  der  Intervalle  (n.~V.)
so  bestimmen,  dass  auch  $s$  beliebig  klein  wird,  wodurch  der  Summe
$\delta_1 D_1 + \cdots + \delta_n D_n$  jede  beliebige  Kleinheit  gegeben,  und  folglich
der  Werth  der  Summe  $S$  in  beliebig  enge  Grenzen  eingeschlossen
werden  kann.
Wir  haben  also  Bedingungen  gefunden,  welche  nothwendig  und
hinreichend  sind,  damit  die  Summe  $S$  bei  unendlichem  Abnehmen  der
Grössen  $\delta$  convergire  und  also  im  engern  Sinne  von  einem  Integrale
der  Function  $f(x)$  zwischen  $a$  und  $b$  die  Rede  sein  könne.
Wird  nun  der  Integralbegriff  wie  oben  erweitert,  so  ist  offenbar,
damit  die  Integration  durchgehends  möglich  sei,  die  letzte  der  beiden
gefundenen  Bedingungen  auch  dann  noch  nothwendig;  an  die  Stelle
der  Bedingung,  dass  die  Function  immer  endlich  sei,  aber  tritt  die
Bedingung,  dass  die  Function  nur  bei  Annäherung  des  Arguments  an
einzelne  Werthe  unendlich  werde,  und  dass  sich  ein  bestimmter
Grenzwerth  ergebe,  wenn  die  Grenzen  der  Integration  diesen  Werthen
unendlich  genähert  werden.
6.
Nachdem  wir  die  Bedingungen  für  die  Möglichkeit  eines  bestimmten
Integrals  im  Allgemeinen,  d.  h.  ohne  besondere  Voraussetzungen  über
die  Natur  der  zu  integrirenden  Function,  untersucht  haben,  soll  nun
diese    Untersuchung    in    besonderen    Fällen    theils    angewandt,    theils
iliKHANN's  gesammelte  mathematiacbe  Werke.  16
\newpage

% --- Page 258 ---
242  XII.    üeber  die  Darstellbarkeit  einer  Function
weiter  ausgeführt  werden,  und  zwar  zunächst  für  die  Functionen,
welche  zwischen  je  zwei  uoch  so  engen  Grenzen  unendlich  oft  un-
stetig sind.
Da  diese  Functionen  noch  nirgends  betrachtet  sind,  wird  es  gut
sein,  von  einem  bestimmten  Beispiele  auszugehen.  Man  bezeichne  der
Kürze  wegen  durch  $(x)$  den  Ueberschuss  von  $x$  über  die  nächste  ganze
Zahl,  oder,  wenn  $x$  zwischen  zweien  in  der  Mitte  liegt  und  diese  Be-
stimmung zweideutig  wird,  den  Mittelwerth  aus  den  beiden  Werthen
$\frac{1}{2}$  und  $-\frac{1}{2}$,  also  die  Null,  ferner  durch  $n$  eine  ganze,  durch  $p$  eine
ungerade  Zahl  und  bilde  alsdann  die  Reihe
\[ f(x) = \sum_{n=1}^{\infty} \frac{(nx)}{n} + \frac{1}{n^p} = \cdots \]
so  convergirt,  wie  leicht  zu  sehen,  diese  Reihe  für  jeden  Werth  von  $x$;
ihr  Werth  nähert  sich,  sowohl,  wenn  der  Argumentwerth  stetig  ab-
nehmend, als  wenn  er  stetig  zunehmend  gleich  $x$  wird,  stets einem
festen  Grenzwerth,  und  zwar  ist,  wenn  $x = \frac{p}{2n}$  (wo  $p$,  $n$  relative  Prim-
zahlen)
\[ f(x+0) = f(x) - \frac{1}{2n}\left(1 + \frac{1}{2} + \frac{1}{3} + \cdots\right) = f(x) - \frac{\infty}{2n}, \]
\[ f(x-0) = f(x) + \frac{1}{2n}\left(1 + \frac{1}{2} + \frac{1}{3} + \cdots\right) = f(x) + \frac{\infty}{2n}, \]
sonst  aber  überall  $f(x+0) = f(x)$,  $f(x-0) = f(x)$.
sonst  aber  überall  f(x  -j-  0)  =  $f(x)$,  f(x  ---  0)  =  $$f(x)$$ .
Diese  Function  ist  also  für  jeden  rationalen  Werth  von  x,  der  in
den  kleinsten  Zahlen  ausgedrückt  ein  Bruch  mit  geradem  Nenner  ist,
unstetig,  also  zwischen  je  zwei  noch  so  engen  Grenzen  unendlich  oft,
so  jedoch,  dass  die  Zahl  der  Sprünge,  welche  grösser  als  eine  ge-
gebene Grösse  sind,  immer  endlich  ist.  Sie  lässt  durchgehends  eine
Integration  zu.  In  der  That  genügen  hierzu  neben  ihrer  Endlichkeit
die  beiden  Eigenschaften,  dass  sie  für  jeden  Werth  von  x  beiderseits
einen  Grenzwerth  f{x  -f-  0)  und  f{x  ---  0)  hat,  und  dass  die  Zahl  der
Sprünge,  welche  grösser  oder  gleich  einer  gegebenen  Grösse  6  sind,
stets  endlich  ist.  Denn  wenden  wir  unsere  obige  Untersuchung  an,
so  lässt  sich  offenbar  in  Folge  dieser  beiden  Umstände  d  stets  so  klein
annehmen,  dass  in  sämmtlichen  Intervallen,  welche  diese  Sprünge  nicht
enthalten,  die  Schwankungen  kleiner  als  6  sind,  und  dass  die  Ge-
sammtgrösse  der  Intervalle,  welche  diese  Sprünge  enthalten,  beliebig
klein  wird.
Es  verdient  bemerkt  zu  werden,  dass  die  Functionen,  welche  nicht
unendlich  viele  Maxima  und  Miuima  haben  (zu  welchen  übrigens  die
eben  betrachtete  nicht  gehört),  wo  sie  nicht  unendlich  werden,   stets
\newpage

% --- Page 259 ---
diese  beiden  Eigenschaften  besitzen  und  daher  allenthalben,  wo  sie
nicht  unendlich  werden,  eine  Integration  zulassen,  wie  sich  auch  leicht
direct  zeigen  lässt.
Um  jetzt  den  Fall,  wo  die  zu  integrirende  Function  $f(x)$  für  einen
einzelnen  Werth  unendlich  gross  wird,  näher  in  Betracht  zu  ziehen,
nehmen  wir  an,  dass  dies  für  ä;  =  0  stattfinde,  so  dass  bei  abnehmen-
dem positiven  x  ihr  Werth  zuletzt  über  jede  gegebene  Grenze  wächst.
Es  lässt  sich  dann  leicht  zeigen,  dass  x$f(x)$  bei  abnehmendem  x
von  einer  endlichen  Grenze  a  an,  nicht  fortwährend  grösser  als  eine
endliche  Grösse  c  bleiben  könne.     Denn  dann  wäre
a  a
I
$$f(x)$$  dx>  c  I
dx
---  j
X
also  grösser  als  c  (log   log  --- j ,  welche  Grösse  mit  abnehmendem
x  zuletzt  in's  Unendliche  wächst.  Es  muss  also  x$$f(x)$$,  wenn  diese
Function  nicht  in  der  Nähe  von  x  =  0  unendlich  viele  Maxima  und
Minima  hat,  noth wendig  mit  x  unendlich  klein  werden,  damit  $$f(x)$$
einer  Integration  fähig  sein  könne.     Wenn  andererseits
bei  einem  Werth  von  a  <  1  mit  x  unendlich  klein  wird,  so  ist  klar,
dass  das  Integral  bei  unendlichem  Abnehmen  der  unteren  Grenze  con-
vergirt.
Ebenso  findet  man,  dass  im  Falle  der  Convergenz  des  Integrals
die  Functionen
\[ f(x) \cdot x \log \frac{1}{x} = \frac{d}{d\log\frac{1}{x}}, \quad f(x) \cdot x \log\frac{1}{x} \log\log\frac{1}{x} = \frac{d}{d\log\log\frac{1}{x}} \]
\[ f(x) \cdot x \log\frac{1}{x} \log\log\frac{1}{x} \cdots \log^{m-1}\frac{1}{x} \log^m\frac{1}{x} = \frac{d}{d\log^m\frac{1}{x}} \]
nicht  bei  abnehmendem  $x$  von  einer  endlichen  Grenze  an  fortwährend
grösser  als  eine  endliche  Grösse  bleiben  können,  und  also,  wenn  sie
nicht  unendlich  viele  Maxima  und  Minima  haben,  mit  $x$  unendlich
klein  werden  müssen;  dass  dagegen  das  Integral  $\int f(x)\, dx$  bei  un-
endlichem Abnehmen  der  unteren  Grenze  convergire,  sobald
\[ f(x) \cdot x \log\frac{1}{x} \cdots \log^{m-1}\frac{1}{x} \left(\log^m\frac{1}{x}\right)^\alpha = \frac{d}{dx} \cdot \frac{1}{\cdots} \]
für  $\alpha > 1$  mit  $x$  unendlich  klein  wird.
\newpage

% --- Page 260 ---
244  ^11-    Ueber  die  Darstellbarkeit  einer  Function
Hat  aber  die  Function  $$f(x)$$  unendlich  viele  Maxima  und  Minima
so  lässt  sich  über  die  Ordnung  ihres  ünendlichwerdens  nichts  be-
stimmen. In  der  That,  nehmen  wir  an,  die  Function  sei  ihrem  ab-
soluten Werthe  nach,  wovon  die  Ordnung  des  Unendlichwerdens  allein
abhängt,  gegeben,  so  wird  man  immer  durch  geeignete  Bestimmung
des  Zeichens  bewirken  können,  dass  das  IniQgrdl  ff (x)dx  bei  unend-
lichem Abnehmen  der  unteren  Grenze  convergire.  Als  Beispiel  einer
solchen  Function,  welche  unendlich  wird  und  zwar  so,  dass  ihre  Ordnung
(die  Ordnung  von  ---  als  Einheit  genommen)  unendlich  gross  ist,  mag
die  Function
(  X  cos  e  ^  )
i         1     -i  >-
dx
dienen.
Das  möge  über  diesen  im  Grunde  in  ein  anderes  Gebiet  gehörigen
Gegenstand  genügen;  wir  gehen  jetzt  au  unsere  eigentliche  Aufgabe,
eine  allgemeine  Untersuchung  über  die  Darstellbarkeit  einer  Function
TJntersuehTing  der  Darstellbarkeit  einer  Function  durch  eine  trigo-
nometrische Beihe  ohne  besondere  Voraussetzungen  über  die  Natur
der  Function.
Die  bisherigen  Arbeiten  über  diesen  Gegenstand  hatten  den  Zweck,
die  Fourier'sche  Reihe  für  die  in  der  Natur  vorkommenden  Fälle  zu
beweisen;  es  konnte  daher  der  Beweis  für  eine  ganz  willkürlich  an-
genommene Function  begonnen,  und  später  der  Gang  der  Function  behufs
des  Beweises  willkürlichen  Beschränkungen  unterworfen  werden,  wenn
sie  nur  jenen  Zweck  nicht  beeinträchtigten.  Für  unsern  Zweck  darf
derselbe  nur  den  zur  Darstellbarkeit  der  Function  nothwendigen  Be-
dingungen unterworfen  werden;  es  müssen  daher  zunächst  zur  Dar-
stellbarkeit nothwendige  Bedingungen  aufgesucht  und  aus  diesen  dann
zur  Darstellbarkeit  hinreichende  ausgewählt  werden.  Während  also
die  bisherigen  Arbeiten  zeigten:  wenn  eine  Function  diese  und  jene
Eigenschaften  hat,  so  ist  sie  durch  die  Fourier'sche  Reihe  darstell-
bar; müssen  wir  von  der  umgekehrten  Frage  ausgehen:  Wenn  eine
Function  durch  eine  trigonometrische  Reihe  darstellbar  ist^  was  folgt
daraus  über  ihren  Gang,  über  die  Aenderung  ihres  Werthes  bei  stetiger
Aenderung  des  Arguments?
\newpage

% --- Page 261 ---
Demnach  betrachten  wir  die  Reihe
«1  sin  a;  +  ^z  sin  2a;  +  *  \textbullet{}  \textbullet{}
-\textbackslash{}-  ^  Iq -^  h^  cos  X  -^  \&2  cos  2x  -\textbackslash{}-  $\blacksquare$  \textbullet{}  \textbullet{}
oder,  wenn  wir  der  Kürze  wegen
^\textbackslash{}  =  ^0,  a^  sin X  -\textbackslash{}-hi  cos x  =  A^,  a^  sin 2x  -\textbackslash{}-  h^  cos 2x  =  Ä,,,  ...
setzen,  die  Reihe
A,  +  A,-\textbackslash{}-A,  +  ---
als  gegeben.  Wir  bezeichnen  diesen  Ausdruck  durch  Sl  und  seinen
Werth  durch  $$f(x)$$j  so  dass  diese  Function  nur  für  diejenigen  Werthe
von  X  vorhanden  ist,  wo  die  Reihe  convergirt.
Zur  Convergenz  einer  Reihe  ist  nothwendig,  dass  ihre  .Glieder
zuletzt  unendlich  klein  werden.  Wenn  die  Coefficienten  a„,  \&«  mit
wachsendem  n  iu's  Unendliche  abnehmen,  so  werden  die  Glieder  der
Reihe  Sl  für  jeden  Werth  von  x  zuletzt  unendlich  klein;  andernfalls
kann  dies  nur  für  besondere  Werthe  von  x  stattfinden.  Es  ist  nöthig,
beide  Fälle  getrennt  zu  behandeln.
Wir  setzen  also  zunächst  voraus,  dass  die  Glieder  der  Reihe  Sl
für  jeden  Werth  von  x  zuletzt  unendlich  klein  werden.
Unter  dieser  Voraussetzung  convergirt  die  Reihe
C  +  Ca;  +  A  f  -  -4i  -  4"  -  ^
\textbullet{}  \textbullet{}  \textbullet{}  ==  F{x),
welche  man  aus  Sl  durch  zweimalige  Integration  jedes  Gliedes  nach  x
erhält,  für  jeden  Werth  von  x.     Ihr  Werth  F{x)  ändert   sich  mit  x
stetig,  und  diese  Function  F  von  x  lässt  folglich  allenthalben  eine
Integration  zu.
Um   Beides  ---  die  Convergenz  der  Reihe  und  die  Stetigkeit  der
Function  Fix)  ---  einzusehen,  bezeichne  man  die  Summe  der  Glieder
K    \textbullet{}         \textbullet{}       .
bis   einschliesslich  durch  N,  den  Rest  der  Reihe,  d.  h.  die  Reihe
durch  jR  und  den  grössten  Werth  von  Am  für  m>  n  durch  s.  Als-
dann bleibt  der  Werth  von  B,  wie  weit  man  diese  Reihe  fortsetzen
möge,  offenbar  abgesehen  vom  Zeichen
<  *  Vn"\_}-  r,«  +  (n  -4-  2)"  "l  /  ^
{{n  +  1)*    '    (n  +  2)"    '  /  ^  n
und  kann  also  in  beliebig  kleine  Grenzen  eingeschlossen  werden,  wenn
man  n  nur  hinreichend  gross  annimmt;  folglich  convergirt  die  Reihe.
\newpage

% --- Page 262 ---
246  XII.   üeber  die  Darstellbarkeit  einer  Function
Ferner  ist  die  Function  F{x)  stetig;  d.h.  ihrer  Aenderung  kann  jede
Kleinheit  gegeben  werden,  wenn  man  der  entsprechenden  Aenderung
von  X  eine  hinreichende  Kleinheit  vorschreibt.  Denn  die  Aenderung
von  F{x)  setzt  sich  zusammen  aus  der  Aenderung  von  H  und  von  N]
offenbar  kann  man  nun  erst  n  so  gross  annehmen,  dass  M,  was  auch
X  sei,  und  folglich  auch  die  Aenderung  von  H  für  jede  Aenderung
von  X  beliebig  klein  wird,  und  dann  die  Aenderung  von  x  so  klein
annehmen,  dass  auch  die  Aenderung  von  N  beliebig  klein  wird.
Es  wird  gut  sein,  einige  Sätze  über  diese  Function  F(x),  deren
Beweise  den  Faden  der  Untersuchung  unterbrechen  würden,  vorauf-
zuschicken.
Lehrsatz  1.     Falls  die  Reihe  Sl  convergirt,  convergirt
F(x  -f-  «  -f  (3)  ---  F(x  +  K  ---  (3)  ---  F{x  ---  «  4-  (3)  -I-  F{x  ---  a  ---  ß)
wenn  cc  und  ß  so  unendlich  klein  werden,  dass  ihr  Verhältniss  endlich
bleibt,  gegen  denselben  Werth  wie  die  Reihe.
In  der  That  wird
F{x  -\textbackslash{}-cc-\textbackslash{}-ß)  ---  F{x  -\textbackslash{}-  a  ---  ß)  ---  F{x  ---  «  +  ^)  -{-  F {x  -  cc  ---  ß)
laß
j     j^    j    sin  a  sin  ß  j^     .    sin  2  a  sin  2  ß  ^^     .    sin  3  a  sin  3  ß
---  ^0  i-  ^1  -^  -ß-  i-  ^2  -2^  -jf  -f-  ^3  -3^  -Yß-  H
oder,  um  den  einfacheren  Fall,  wo  ß  =  a,  zuerst  zu  erledigen,
F  {X  -}$\blacksquare$  2a)  ---  2F  (x)  +  F  (x  -  2a)  ^  ^      1    ^    /sin  aV  /sin2ay
Ist  die  unendliche  Reihe
A  +  A  +  -12  +  -  ---/"W,
die  Reihe
A  +  A-\textbackslash{} --- ^+  A-i  =  $$f(x)$$  +  \pounds„,
so  muss  sich  für  eine  beliebig  gegebene  Grösse  d  ein  Werth  m  von  n
angeben  lassen,  so  dass,  wenn  n'>  m,  f „  <  d  wird.  Nehmen  wir  nun
a  so  klein  an,  dass  m«  <  ä,  setzen  wir  ferner  mittelst  der  Substitution
^(^)'^»"  die  Form
1,  cc  '
und  theilen  wir  diese  letztere  unendliche  Reihe  in  drei  Theile,  indem  wir
1)  die  Glieder  vom  Index  1  bis  m  einschliesslich,
2)  vom  Index  m  -\textbackslash{}-  1   bis  zur  grössten  unter  ---  liegenden  ganzen
Zahl,  welche  s  sei,
3)  von  s  +  1  bis  unendlich,
\newpage

% --- Page 263 ---
zusammenfassen,  so  besteht  der  erste  Theil  aus  einer  endlichen  Anzahl
stetig  sich  ändernder  Glieder  und  kann  daher  seinem  Grenzwerth  0
beliebig  genähert  werden,  wenn  man  a  hinreichend  klein  werden  lässt;
der  zweite  Theil  ist,  da  der  Factor  von  e„  beständig  positiv  ist,  offen-
bar abgesehen  vom  Zeichen
^   f/8inwa\textbackslash{}"''         /sinsa\textbackslash{}$^{*)}$
<^|(    m«    j   -{-^.\textasciitilde{})  p
um  endlich  den  dritten  Theil  in  Grenzen  einzuschliessen,  zerlege  man
das  allgemeine  Glied  in
Kein
(n  ---  1)  a\textbackslash{}  /sin  (n  ---  1) «  \textbackslash{}  *  |
und
K8in(n  ---  l)a\textbackslash{}8         /sin  na\textbackslash{} 2 1     8in(2w ---  1)  a  sin  a
^        )  \textasciitilde{}  (   nir)  j
---  ---  *«  (^^p  '$\blacksquare$>
so  leuchtet  ein,  dass  es
<*(^^   '-!  +  «-L
|(n ---  lycccc        nnaal     '        nna
und  folglich  die  Summe  von  n  =  s  -\textbackslash{}-  1  bis  n  =  oo
welcher  Werth  für  ein  unendlich  kleines  a  in
^{;;^  +  ;^}     übergeht.
Die  Reihe
^7        [  /sin  (n  ---  1)  a\textbackslash{}  2
Ksin
(n ---  1)  «\textbackslash{}*         /sinwa\textbackslash{}2|
{n-l)a    )   \textasciitilde{}  [\textasciitilde{}^i^)  )
nähert  sich  daher  mit  abnehmendem  a  einem  Grenzwerth,  der  nicht
grösser  als
1  '       TT        '       7C7t\textbackslash{}
sein  kann,  also  Null  sein  muss,  und  folglich  convergirt
F{x  +  2a)  ---  2F(x)  +  F  {x  ---  2a)
iua  '
welches
mit  in's  Unendliche  abnehmendem  «  gegen  $f(x)$,   wodurch  unser  Satz
für  den  Fall  ß  =  a  bewiesen  isi
Um  ihn  allgemein  zu  beweisen,  sei
F{x-\textbackslash{}-a-\textbackslash{}-ß)-  2F(x)  -{-F{x-a-ß)  =  {a-\textbackslash{}-  ß)'  (fix)  +  d,)
F{x-\textbackslash{}-a-ß)-  2F{x)  -\textbackslash{}.F{x-a-{-ß)=^{a-  ßfifix)  +  d,),
woraus
\newpage

% --- Page 264 ---
248  XII.    üeber  die  Darstellbarkeit  einer  Function
F{x-i-a-\textbackslash{}-ß)  ---  Fix-j-a  ---  ß)  ---  F(x  ---  a-j-ß)-{-  Fix  ---  cc  ---  ß)
=  Aaß  fix)  +  («  +  ßy  d,  -  («  -  ßf  d, .
In  Folge  des  eben  Bewiesenen  werden  nun  d^  und  d^  unendlich  klein,
sobald  a  und  ß  unendlich  klein  werden;  es  wird  also  auch
unendlich  klein,  wenn  dabei  die  Coefficienten  von  ^^  und  ^2  nicht  un-
endlich gross  werden,  was  nicht  stattfindet,  wenn  zugleich  ---  endlich
bleibt;  und  folglich  convergirt  alsdann
F{x  +  K  +  (3)  ---  F{x  -f  «  ---  p)  ---  Fjx  ---  K  +  ^)  +  Fix.-  a  ---  ß)
iaß
gegen  fix),  w.  z.  b.  w.
Lehrsatz  2.
F(x  +  2«)  4-  Fjx  ---  2«)  ---  2Fix)
2«
wird  stets  mit  a  unendlich  klein.
Um  dieses  zu  beweisen,  theile  man  die  Reihe
in  drei  Gruppen,  von  welchen  die  erste  alle  Glieder  bis  zu  einem  festen
Index  m  enthält,  von  dem  an  A»  immer  kleiner  als  s  bleibt,  die  zweite
alle  folgenden  Glieder,  für  welche  wa  <  als  eine  feste  Grösse  c  ist,  die
dritte  den  Rest  der  Reihe  umfasst.  Es  ist  dann  leicht  zu  sehen,  dass,
wenn  «  in's  unendliche  abnimmt,  die  Summe  der  ersten  endlichen  Gruppe
endlich  bleibt,  d.  h.  <  eine  feste  Grösse  Q ;  die  der  zweiten  <  «  --- ,  die
der  dritten
^^  nnacc        ccc
c  <  na
Folglich  bleibt
2a  '  X I      "\textbackslash{}    na    /   '
<2(^«  +  .(c  +  i)),
woraus  der  z.  b.  Satz  folgt.
Lehrsatz  3.  Bezeichnet  man  durch  h  und  c  zwei  beliebige  Con-
stanten, die  grössere  durch  c,  und  durch  A(a;)  eine  Function,  welche
nebst  ihrem  ersten  Differentialquotienten  zwischen  b  und  c  immer
stetig  ist  und  an  den  Grenzen  gleich  Null  wird,  und  von  welcher  der
zweite  Differentialquotient  nicht  unendlich  viele  Maxima  und  Minima
hat,  so  wird  das  Integral
\newpage

% --- Page 265 ---
$c$
fi^  I  F(x)  cos  ^{x  ---  d)k  (x)  dx ,
$b$
wenn  ^  in's  Unendliche  wächst,  zuletzt  kleiner  als  jede  gegebene  Grösse.
Setzt  man  für  F{x)  seinen  Ausdruck  durch  die  Reihe,   so  erhält
man  für
$c$
|it/i  1  F(x)  cos n,{x  ---  ä)X {x) dx
$b$
die  Reihe  {O)
$c$
jtift  /  ( 0  +  C'x  +  Äq  \textasciitilde{})  cos  (i(x  ---  a)X  (x)  dx
$l$
$c$
\textasciitilde{}2^J  ^«cos/^Ca;  --- a)A(a;)<^a;.
1,00  6
Nun   lässt   sich  An  cos  [i(x  ---  a)    offenbar  als  ein  Aggregat  von
cos{}i-\textbackslash{}-n)(x --- a),  sin(|ü4"*^)(^\textasciitilde{}^);  cos(|!* --- n)(x --- d),  sin(ft --- n)(x  ---  a)'
ausdrücken,  und  bezeichnet  man  in  demselben  die  Summe  der  beiden
ersten  Glieder  durch  Sfi^n,  die  Summe  der  beiden  letzten  Glieder  durch
Bu---n,  so  hat  man  cos|ü(ä;  ---  a)  A,i  =  -B^,+n  +  Bf,---n,
^^  =  -  (ft  +  nyB,  +  .,     ^^"  =  -(.«-  n)2\pounds,\_„,
und  es  werden  J5«+„  und  J5^\_«  mit  wachsendem  n,  was  auch  x  sei,
zuletzt  unendlich  klein.
Das  allgemeine  Glied  der  Reihe  (0)
---  I
An  cos  (i(x  ---  a)X  (x)  dx
wird  daher
=  -T7-^^,  f^^'P-  A  (x)  dx  -f    ,/'     ,,  r^^^q^  A  (o;)  (?a;
oder  durch  zweimalige  partielle  Integration,  indem  man  zuerst  k{x),
dann  k'  {x)  als  constant  betrachtet,
c  c
6  6
da    A  (a;)    und    A'  (a;)    und    daher    auch    die    aus    dem    Integralzeichen
tretenden  Glieder  an  den  Grenzen  =  0  werden.
$c$
Man  überzeugt  sich  nun  leicht,  dass    /  B/^i  +  n  k"{x)dx<,   wenn  fi
\newpage

% --- Page 266 ---
250  \textbullet{}XII.   Ueber  die  Darstellbarkeit  einer  Function
in's  Unendliclie  wächst,  was  auch  n  sei,  unendlich  klein  wird;  denn
dieser  Ausdruck  ist  gleich  einem  Aggregat  der  Integrale
/  cos  (ft  +  w)  (ic  ---  a)  X'  (ic)  dx ,       /  sin  (ft  +  w)  (^  ---  a)  X'  {x)  dx ,
b  l
und  wenn  ft  +  ^  unendlich  gross  wird,  so  werden  diese  Integrale,
wenn  aber  nicht,  weil  dann  n  unendlich  gross  wird,  ihre  Coefficienten
in  diesem  Ausdrucke  unendlich  klein.
Zum  Beweise  unseres  Satzes  genügt  es  daher  offenbar,  wenn  von
der  Summe
über  alle  ganzen  Werthe  von  n  ausgedehnt,  welche  den  Bedingungen
w  <  ---  c',  c"<Cn  <C  fi  ---  c",  n  -f-  c^^  <  n  genügen,  für  irgend  welche
positive  Werthe  der  Grössen  c  gezeigt  wird,  dass  sie,  wenn  fi  unend-
lich gross  wird,  endlich  bleibt.  Denn  abgesehen  von  den  Gliedern,  für
welche  ---  c  <n  <Cc',  (i  ---  c"'<  n  <  [i  -{-  c^^,  welche  offenbar  un-
endlich klein  werden  und  von  endlicher  Anzahl  sind,  bleibt  die  Reihe  (0)
offenbar  kleiner  als  diese  Summe,  multiplicirt  mit  dem  grössten  Werthe
0
von    /  B^$\pm$nX' {x)dXy  welcher  unendlich  klein  wird,
ö
Nun  ist  aber,  wenn  die  Grössen  c  >  1  sind,  die  Summe
^
^'
=
i  'V
^
in  den  obigen  Grenzen,  kleiner  als
dx
f^
J   (1
ausgedehnt  von
---  cx)  bis   ,   bis  1
denn  zerlegt  man  das  ganze  Intervall  von  ---  oo  bis  -j-  c»  von  Null
anfangend  in  Intervalle  von  der  Grösse  --- ,  und  ersetzt  man  überall  die
Function  unter  dem  Integralzeichen  durch  den  kleinsten  Werth  in  jedem
Intervall,  so  erhält  man,  da  diese  Function  zwischen  den  Integrations-
grenzen nirgends  ein  Maximum  hat,  sämmtliche  Glieder  der  Reihe.
Führt  man  die  Integration  aus,  so  erhält  man
jj
^^(^t-xY  =  j{-\textbackslash{}  +  ih^  +  21oga;-  21og(l  -a;))  +  const.
\newpage

% --- Page 267 ---
und  folglich  zwischen  den  obigen  Grenzen  einen  Werth,  der  mit  fi
nicht  unendlich  gross  wird.  ($^{*)}$
9.
Mit  Hülfe  dieser  Sätze  lässt  sich  über  die  Darstellbarkeit  einer
Function  durch  eine  trigonometrische  Reihe,  deren  Glieder  für  jeden
Argumentwerth  zuletzt  unendlich  klein  werden,  Folgendes  feststellen:
\section*{I.  Wenn  eine  nach  dem  Intervall  27t  periodisch  sich  wieder-}
holende Function  $f(x)$  durch  eine  trigonometrische  Reihe,  deren  Glieder
für  jeden  Werth  von  x  zuletzt  unendlich  klein  werden,  darstellbar  sein
soll,  so  muss  es  eine  stetige  Function  F(x)  geben,  von  welcher  $f(x)$
so  abhängt,  dass
F(x-\textbackslash{}-a-]-ß)  ---  Fix-\textbackslash{}-cc  ---  ß)  ---  F{x  ---  a-\textbackslash{}-ß)-\textbackslash{}-Fix---  a  ---  ß)
iaß  '
wenn  a  und  ß  unendlich  klein  werden  und  dabei  ihr  Verhältniss  end-
lich bleibt,  gegen  $f(x)$  convergirt.
Es  muss  ferner
$c$
fi^  I
F(x)  cos  ^{x  ---  a)X  (x)  dx ,
6
wenn  X  (x)  und  A'  {x)  an  den  Grenzen  des  Integrals  =  0  und  zwischen
denselben  immer  stetig  sind,  und  X' {x)  nicht  unendlich  viele  Maxima
und  Minima  hat,  mit  wachsendem  ^  zuletzt  unendlich  klein  werden.
IL  Wenn  umgekehrt  diese  beiden  Bedingungen  erfüllt  sind,  so
giebt  es  eine  trigonometrische  Reihe,  in  welcher  die  Coefficienten  zu-
letzt unendlich  klein  werden,  und  welche  überall,  wo  sie  convergirt,
die  Function  darstellt.
Denn  bestimmt  man  die  Grössen  C,  Äq  so,  dass
F{x)-C'x-A,'^
eine  nach  dem  Intervall  2n  periodisch  wiederkehrende  Function  ist
und  entwickelt  diese  nach  Fourier's  Methode  in  die  trigonometrische
Reihe
^      A-i        A^        A^
indem  man
7t
---  7t
^f(F{t)  -et-A^-Pj  cos  n(x-t)dt==-f^
\newpage

% --- Page 268 ---
setzt,  so  muss  (n.  V.)
$n$
^„  ==  \_  ^  C(f{S)  ---  C't  ---  A^  y)  cosnix  ---  t)dt
---  n
mit  wachsendem  n  zuletzt  unendlich  klein  werden;  woraus  nach  Satz  1
des  vorigen  Art.  folgt,  dass  die  Reihe
Aq  -]-  Ä^  -\textbackslash{}-  Ä2  -{-  \textbullet{}  \textbullet{}  \textbullet{}
überall,  wo  sie  convergirt,  gegen  $f(x)$  convergirt.
\section*{III.  Es  sei  h<.x<.c,  und  Q(t)  eine  solche  Function,  dass  Q{t)}
und  Q  (t)  inr  t  =b  und  t  =  c  den  Werth  0  haben  und  zwischen  diesen
Werthen  stetig  sich  ändern,  (»"(0  nicht  unendlich  viele  Maxima  und
Minima  hat,  und  dass  ferner  fürt  =  x  q  (t)  =  1,  q'  {t)  =  0,  Q'(t)  =  0,
Q"{t)  und  Q^^(t)  aber  endlich  und  stetig  sind;  so  wird  der  Unterschied
zwischen  der  Reihe
A^-Jr  Ai-\textbackslash{}   \textbackslash{}-  An
und  dem  Integral
sin   ^
---  (x  ---  t)
dd   -.
Tv
c  .     {X  ---  t)
8in
U'^(0   äj^  9«''*
2
$b$
mit  wachsendem  n  zuletzt  unendlich  klein.     Die  Reihe
Jo  +  A  +  ^2  +  \textbullet{}  \textbullet{}  \textbullet{}
wird  daher  convergiren  oder  nicht  convergiren,  je  nachdem
.    2«  4-  1 ,        .,
sin  \textbullet{}  -! ---  {X  ---  t)
dd   ---   i
c  .     X  ---  t
sm-
^/^W   äi^ --- «>w*
2
6
sich  mit  wachsendem  n  zuletzt  einer  festen  Grenze  nähert  oder  dies
nicht  stattfindet.
In  der  That  wird
$n$
^  +  ^  +  \textbullet{}\textbullet{}\textbullet{}^„  =  1  f(F(t)-C't-A,\textasciitilde{}\textbackslash{}^-nnco8n(x-t)dt,
oder,  da  2w  +  1,       ,,
sin   -^ ---  {x  ---  t)
dd-
n  "STT  f  4.\textbackslash{}        n  ^1^^  cos  n{x  ---  t)
ist.
.    X  ---  t
sin  --- - ---
dt""
1, »
\newpage

% --- Page 269 ---
\textbullet{}    2n  4-  1  .
sin   -' ---  {x  ---  t)
dd
.    X  ---  t
sm
---  n
Nun  wird  aber  nach  Satz  3  des  vorigen  Art.
dd
.    2«+  1  .         ^.
8in   J ---  {X  ---  t)
.    X  ---  t
Bin
---  n
bei  unendlichem  Zunehmen  von  n  unendlich  klein,  wenn  A  nebst
ihrem  ersten  Differentialquotienten  stetig  ist,  A"  {f)  nicht  unendlich  viele
Maxima  und  Minima  hat,  und  für  t  =  x  X(t)  =  0,  k'(t)  =  0,  A"(^)  =  0,
X"'{t)  und  X^^  (t)  aber  endlich  und  stetig  sind.
Setzt  man  hierin  A  (f)  ausserhalb  der  Grenzen  \&,  c  gleich  1  und
zwischen  diesen  Grenzen  =1  ---  (>  (0  >  ^^\textregistered  offenbar  verstattet  ist,  so
folgt,  dass  die  Differenz  zwischen  der  Reihe  Äj^  -}-\textbullet{}\textbullet{}'-{-  A„  und  dem
Integral
.    2n  4-  1
sm   {X  ---  t)
dd-
.    X
sin-
i-„f{m  -c-i-A,  i')  ^ --- ,  (t)  dt
$b$
mit  wachsendem  n  zuletzt  unendlich  klein  wird.     Man  überzeugt  sich
aber  leicht  durch  partielle  Integration,  dass
sm  --- - ---  {x  ---  t)
dd  ---
X  ---  *
sin
hj-i^C^t  +  A«'-!)   ^   ,{t)dt,
2
$b$
wenn   n    unendlich    gross    wird,    gegen  Ä^  convergirt,    wodurch  man
obigen  Satz  erhält.
10.
Aus  dieser  Untersuchung  hat  sich  also  ergeben,  dass,  wenn  die
Coefficienten  der  Reihe  Sl  zuletzt  unendlich  klein  werden,  dann  die
Convergenz  der  Reihe  für  einen  bestimmten  Werth  von  x  nur  abhängt
von  dem  Verhalten  der  Function  $$f(x)$$  in  unmittelbarer  Nähe  dieses
Werthes.
Ob  nun  die  Coefficienten  der  Reihe  zuletzt  unendlich  klein  werden,
\newpage

% --- Page 270 ---
wird  in  vielen  Fällen  nicht  aus  ihrem  Ausdrucke  durch  bestimmte
Integrale,  sondern  auf  anderm  Wege  entschieden  werden  müssen.  Es
verdient  indess  ein  Fall  hervorgehoben  zu  werden,  wo  sich  dies  un-
mittelbar aus  der  Natur  der  Function  entscheiden  lässt,  wenn  nämlich
die  Function  $$f(x)$$  durchgehends  endlich  bleibt  und  eine  Integration
zulässt.
In  diesem  Falle  muss,  wenn  man  das   ganze  Intervall  von  ---  n
bis  7t  der  Reihe  nach  in  Stücke  von  der  Grösse
zerlegt,  und  durch  D^  die  grösste  Schwankung  der  Function  im  ersten,
durch  Dg  im  zweiten,  u.  s.  w.  bezeichnet,
unendlich  klein  werden,  sobald  sämmtliche  d  unendlich  klein  werden.
7t
Zerlegt  man  aber  das  Integral    i  $$f(x)$$  9,mn(x  ---  a)dx,  in  welcher
7t
Form  von  dem  Factor  abgesehen  die  Coefficienten  der  Reihe  ent-
halten  sind,  oder  was  dasselbe  ist,    /  $f(x)$  smn{x  ---  a)dx  von  x  =  a
$a$
anfangend  in  Integrale  vom  Umfange  --- ,  so  liefert  jedes  derselben  zur
2
Summe  einen  Beitrag  kleiner  als  --- ,  multiplicirt  mit  der  grössten
Schwankung  in  seinem  Intervall,  und  ihre  Summe  ist  also  kleiner  als
eine  Grosse,  welche  n.  V.  mit  ---  unendlich  klein  werden  muss.
'  n
In  der  That:  diese  Integrale  haben  die  Form
a-U ---I\_2rt
$n$
I
$f(x)$  sin  w  (ic  ---  a)  dx .
aA   2«
$n$
Der  Sinus  wird  in  der  ersten  Hälfte  positiv,  in  der  zweiten  negativ.
Bezeichnet  man  also  den  grössten  Werth  von  $$f(x)$$  in  dem  Intervall
des  Integrals  durch  M,  den  kleinsten  durch  m,  so  ist  einleuchtend,
dass  man  das  Integral  vergrössert,  wenn  man  in  der  ersten  Hälfte
fix)  durch  Jf,  in  der  zweiten  durch  m  ersetzt,  dass  man  aber  das
Integral  verkleinert,  wenn  man  in  der  ersten  Hälfte  $$f(x)$$  durch  m  und
in  der  zweiten  durch  M.  ersetzt.     Im  ersteren  Falle  aber  erhält  man
--- „  2  2
den  Werth  ---{M.  ---  w),  im  letzteren  ---  (w  ---  ütf).    Es  ist  daher  dies
\newpage

% --- Page 271 ---
Integral    abgesehen    vom    Zeichen    kleiner    als      -  (M --- m)    und    das
Integral
a+2«
/
$f(x)$  sin  n  {x  ---  a)  dx
kleiner  als
1  (M,  -  m,)  +  1  {M,  -  m,)  +  |
{M,  -  m,)  +  \textbullet{}  \textbullet{}  \textbullet{} ,
wenn  man  durch  M^  den  grössten,  durch  m^  den  kleinsten  Werth  von
$$f(x)$$  im  sten  Intervall  bezeichnet;  diese  Summe  aber  muss,  wenn  $f(x)$
einer  Integration  fähig  ist,  unendlich  klein  werden,  sobald  n  unendlich
gross  und  also  der  Umfang  der  Intervalle  ---  unendlich  klein  wird.
In  dem  vorausgesetzten  Falle  werden  daher  die  Coefficienten  der
Reihe  unendlich  klein.
11.
Es  bleibt  nun  noch  der  Fall  zu  untersuchen,  wo  die  Glieder  der
Reihe  Sl  für  den  Argumentwerth  x  zuletzt  unendlich  klein  werden,
ohne  dass  dies  für  jeden  Argumentwerth  stattfindet.  Dieser  Fall  l'ässt
sich  auf  den  vorigen  zurückführen.
Wenn  man  nämlich  in  den  Reihen  für  den  Argumentwerth  x  -}-t
und  X  ---  t  die  Glieder  gleichen  Ranges  addirt,  so  erhält  man  die  Reihe
2^0  +  2^1  cos^  +  2^2  cos 2^  H  ,
in  welcher  die  Glieder  für  jeden  Werth  von  t  zuletzt  unendlich  klein
werden  und  auf  welche  also  die  vorige  Untersuchung  angewandt  wer-
den kann.
Bezeichnet  man  zu  diesem  Ende  den  Werth  der  unendlichen  Reihe
^,    ,     ,,,       ,      .    XX    ,      .    tt  .    cOBt  .    cos2f  .    cos3f
C+  C  x-{-  Äq^  +  A^-^  ---  A^---   A^  ---^   A^  ---
durch    G(t)f   so   dass  ---  ^ --- ^
überall,  wo  die  Reihen  für
F(x-{-t)  und  F(x --- t)  convergiren,  =  G  (f)  ist,  so  ergiebt  sich
Folgendes:
V.  Wenn  die  Glieder  der  Reihe  Sl  für  den  Argumentwerth  x  zu-
letzt unendlich  klein  werden,  so  muss
$c$
[ifi  I  G(t)  cos  fi(t  ---  o)  A (t)  dt,
$b$
wenn  A  eine  Function  wie  oben  ---  Art.  9  ---  bezeichnet,  mit  wachsen-
\newpage

% --- Page 272 ---
dem  fi  zuletzt  unendlich  klein  werden.  Der  Werth  dieses  Integrals
setzt  sich  zusammen  aus  den  beiden  Bestandtheilen
c  c
^^  I        2        ^^^  ^{t---a)k  (t)  dt  und  f*^  /  ---^ --- -  cos  ii(t  ---  a)  A  (t)  dt ,
b  h
wofern  diese  Ausdrücke  einen  Werth  haben.  Das  Unendlichkleinwerden
desselben  yrird  daher  bewirkt  durch  das  Verhalten  der  Function  F  an
zwei  symmetrisch  zu  beiden  Seiten  von  x  gelegenen  Stellen.  Es  ist
aber  zu  bemerken,  dass  hier  Stellen  vorkommen  müssen,  wo  jeder  Be-
staüdtheil  für  sich  nicht  unendlich  klein  wird;  denn  sonst  würden  die
Glieder  der  Reihe  für  jeden  Argumentwerth  zuletzt  unendlich  klein
werden.  Es  müssen  also  dann  die  Beiträge  der  symmetrisch  zu  beiden
Seiten  von  x  gelegenen  Stellen  einander  aufheben,  so  dass  ihre  Summe
für  ein  unendliches  (i  unendlich  klein  wird.  Hieraus  folgt,  dass  die
Reihe  il  nur  für  solche  Werthe  der  Grösse  x  convergiren  kann,  zu
welchen  die  Stellen,  wo  nicht
c  \textbullet{}
fift  /  F(x)  cos (i{x  ---  d)  k  (x) dx
für  ein  unendliches  ft  unendlich  klein  wird,  symmetrisch  liegen.  Offenbar
kann  daher  nur  dann,  wenn  die  Anzahl  dieser  Stellen  unendlich  gross
ist,  die  trigonometrische  Reihe  mit  nicht  in's  Unendliche  abnehmenden
Coefficienten  für  eine  unendliche  Anzahl  von  Argumentwerthen  con-
vergiren.
Umgekehrt  ist
$n$
An  =  ---  nn---  j
iG{t)  -\textasciitilde{}  Aq  y j
cos ntdt
0
und  wird  also  mit  wachsendem  n  zuletzt  unendlich  klein,  wenn
$e$
^fi  /  G  (f)  cos  ^(t  ---  a)  -l  (t)  dt
0
für  ein  unendliches  ft  immer  unendlich  klein  wird.
\section*{II.  Wenn  die  Glieder  der  Reihe  Sl  für  den  Argumentwerth  x}
zuletzt  unendlich  klein  werden,  so  hängt  es  nur  von  dem  Gange  der
Function  G  (t)  für  ein  unendlich  kleines  t  ab,  ob  die  Reihe  convergirt
oder  nicht,  und  zwar  wird  der  Unterschied  zwischen
Jo  +  -4,  H  ]r  A„
und  dem  Integrale
\newpage

% --- Page 273 ---
.    2«  4-1,
am  --- \textasciitilde{} ---  t
dd  ---
-'.h
sin  ---
-di^\textasciitilde{}QQyji
mit  wachsendem  n  zuletzt  unendlich  klein,  wenn  h  eine  zwischen  0
und  n  enthaltene  noch  so  kleine  Constante  und  q  {t)  eine  solche  Function
bezeichnet,  dass  Q(t)  und  Q'(t)  immer  stetig  und  für  t  =  h  gleich  Null
sind,  (>"(0  ^icht  unendlich  viele  Maxima  und  Minima  hat,  und  für
t^O,  Q(t)=  1,  Q(t)  =  0,  Q"{t)  =  0,  ()'"(0  und  P""(0  aber  endlich
und  stetig  sind.
12.
Die  Bedingungen  für  die  Darstellbarkeit  einer  Function  durch
eine  trigonometrische  Reihe  können  freilich  noch  etwas  beschränkt
und  dadurch  unsere  Untersuchungen  ohne  besondere  Voraussetzungen
über  die  Natur  der  Function  noch  etwas  weiter  geführt  werden.  So
z.  B.  kann  in  dem  zuletzt  erhaltenen  Satze  die  Bedingung,  dass
q"  (0)  =  0  sei,  weggelassen  werden,  wenn  man  in  dem  Integrale
.    2wH-l,
sin- --- ---   1
dd---
l/«(0
.    t
8in----
2
d«^
Q{t)dt
G{t)  durch  G{t)  ---  6r(0)  ersetzt.  Es  wird  aber  dadurch  nichts  Wesent-
liches gewonnen.
Indem  wir  uns  daher  zur  Betrachtung  besonderer  Fälle  wenden,
wollen  wir  zunächst  der  Untersuchung  für  eine  Function,  welche  nicht
unendlich  viele  Maxima  und  Minima  hat,  diejenige  Vervollständigung
zu  geben  suchen,  deren  sie  nach  den  Arbeiten  Dirichlet's  noch
fähig  ist.
Es  ist  oben  bemerkt,  dass  eine  solche  Function  allenthalben  in-
tegrirt  werden  kann,  wo  sie  nicht  unendlich  wird,  und  es  ist  offenbar,
dass  dies  nur  für  eine  endliche  Anzahl  von  Argumentwerthen  eintreten
kann.  Auch  lässt  der  Beweis  Dirichlet's,  dass  in  dem  Integralaus-
drucke für  das  wte  Glied  der  Reihe  und  für  die  Summe  ihrer  n  ersten
Glieder  der  Beitrag  aller  Strecken  mit  Ausnahme  derer,  wo  die  Function
unendlich  wird,  und  der  dem  Argumentwerth  der  Reihe  unendlich  nahe
liegenden  mit  wachsendem  n  zuletzt  unendlich  klein  wird,  imd  dass
Bikmank's  gesammelte  mathematische  Werke.  17
\newpage

% --- Page 274 ---
=H-6      sm  --- ^---^  {x  ---  t)
fm
dt
sin
2
wenn  0  <  \&  <  ä  und  f{ii)  zwischen  den  Grenzen  des  Integrals  nicht
unendlich  wird,  für  ein  unendliches  n  gegen  nf(x  -f-  0)  convergirt,  in
der  That  nichts  zu  wünschen  übrig,  wenn  man  die  unnöthige  Voraus-
setzung, dass  die  Function  stetig  sei,  weglässt.  Es  bleibt  also  nur  noch
zu  untersuchen,  in  welchen  Fällen  in  diesen  Integralausdrücken  der
Beitrag  der  Stellen,  wo  die  Function  unendlich  wird,  mit  wachsendem
n  zuletzt  unendlich  klein  wird.  Diese  Untersuchung  ist  noch  nicht
erledigt;  sondern  es  ist  nur  gelegentlich  von  Dirichlet  gezeigt,  dass
dies  stattfindet  unter  der  Voraussetzung,  dass  die  darzustellende  Function
eine  Integration  zulässt,  was  nicht  nothwendig  ist.
Wir  haben  oben  gesehen,  dass,  wenn  die  Glieder  der  Reihe  ß  für
jeden  Werth  von  x  zuletzt  unendlich  klein  werden,  die  Function  F(x),
deren  zweiter  Differentialquotient  $$f(x)$$  ist,  endlich  und  stetig  sein  muss,
und  dass
F{x  +  cc)  ---  2  Fix)  +  F{x  ---  cc)
mit  «  stets  unendlich  klein  wird.  Wenn  nun  F'  {x  +  t)  ---  F'-  {x  ---  t)
nicht  unendlich  viele  Maxima  und  Minima  hat,  so  muss  es,  wenn  t
Null  wird,  gegen  einen  festen  Grenzwerth  L  convergiren  oder  unend-
lich gross  werden,  und  es  ist  offenbar,  dass
^j\textbackslash{}F'{x^t)-F'{x-t))dt  =
F{x  +  a)  -  2  F{x)  +  F{x  ---  et)
0
dann  ebenfalls  gegen  L  oder  gegen  oo  convergiren  muss  und  daher
nur  unendlich  klein  werden  kann,  wenn  F'  {x-^t)---  F'  {x  ---  i)  gegen
Null  convergirt.  Es  muss  daher,  wenn  $$f(x)$$  ivx  x=a  unendlich  gross
wird,  doch  immer  f{a  +  0  +  f^^  ---  0  ^i^  an  i(  =  0  integrirt  werden
können.     Dies  reicht  hin,  damit
a --- 8        c
I    I  -\textbackslash{}-  I  ]dx  [fix)  cos  n{x  ---  a))
6        a-\textbackslash{}-B
mit  abnehmendem  s  convergire  und  mit  wachsendem  n  unendlich  klein
werde.  Weil  ferner  die  Function  F{x)  endlich  und  stetig  ist,  so  muss
F'  (x)  bis  an  ic  =  a  eine  Integration  zulassen  und  (x  ---  d)  F'  (x)  mit
{x  ---  d)  unendlich  klein  werden,  wenn  diese  Function  nicht  unendlich
viele  Maxima  und  Minima  hat;  woraus  folgt,  dass
'"\%''r""=(-^-")Aw+j"w
\newpage

% --- Page 275 ---
und  also  aucli  (x  ---  d)$$f(x)$$  bis  an  x  ==  a  integrirt  werden  kann.  Es
kann  daher  auch  f$$f(x)$$  sin  w  {x^---  a)  dx  bis  an  x  =  a  integrirt  werden,
und  damit  die  Coefficienten  der  Reihe  zuletzt  unendlich  klein  werden,
ist  oflfenbar  nur  noch  nöthig,  dass
$c$
I
$$f(x)$$  smn{x  ---  d)dx,  wo  6  <  a  <  c,
$b$
mit  wachsendem  n  zuletzt  unendlich  klein  werde.     Setzt  man
$$f(x)$$  {x  ---  d)  =  (p  {x),
so  ist,  wenn  diese  Function  nicht  unendlich  viele  Maxima  und  Mioinia
hat,  für  ein  unendliches  n
c  c
I
$f(x)$  sin« {x  ---  a)dx=  1  \textasciitilde{}^-  smn{x  ---  d)dx  =  it
9!  y
wie  Dirichlet  gezeigt  hat.     Es  muss  daher
gj (a  +  0  +  9?  («  ---  0  =  /(«  -\textbackslash{}-t)t  ---  f{a---t)t
mit  t  unendlich  klein  werden,  und  da
bis  an  t  =f=  0  integrirt  werden  kann  und  folglich  auch
f{a-\textbackslash{}rt)t  +  f{a-t)t
mit  t  unendlich  klein  wird,  so  muss  sowohl  f(a  -\textbackslash{}-  t)  t,  als  f(a  ---  t)t
mit  abnehmendem  t  zuletzt  unendlich  klein  werden.  Von  Functionen,
welche  unendlich  viele  Maxima  und  Minima  haben,  abgesehen,  ist  es
also  zur  Darstellbarkeit  der  Function  f{cc)  durch  eine  trigonometrische
Reihe  mit  in's  Unendliche  abnehmenden  Coefficienten  hinreichend  und
nothwendig,  dass,  wenn  sie  für  x  =  a  unendlich  wird,  f(a  -\textbackslash{}-  f)t  und
f(a  ---  t)t  mit  t  unendlich  klein  werden  und  f{a  +  0  +  f{<^  \textasciitilde{}  0  ^^^
an  ^  «=  0  integrirt  werden  kann.
unendlich  klein  werden,  kann  eine  Function  $$f(x)$$,  welche  nicht  unend-
lich viele  Maxima  und  M'inima  hat,  da
0
^^  I  Fix)  cosft  (x  ---  a)  X  (x) dx
nur  für  eine  endliche  Anzahl  von  Stellen  für  ein  unendliches  /t  nicht
unendlich  klein  wird,  auch  nur  für  eine  endliche  Anzahl  von  Argument-
werthen  dargestellt  werden,  wobei  es  unnöthig  ist  länger  zu  verweilen.
\newpage

% --- Page 276 ---
260  XII.    üeber  die  Darstellbarkeit  einer  Function
13.
Was  die  Functionen  betrifft,  welche  unendlich  viele  Maxima  und
Minima  haben,  so  ist  es  wohl  nicht  überflüssig  zu  bemerken,  dass  eine
Function  $f(x)$,  welche  unendlich  viele  Maxima  und  Minima  hat,  einer
Integration  durchgehends  fähig  sein  kann,  ohne  durch  die  Fourier'sche
Reihe  darstellbar  zu  sein.  Dies  findet  z.  B.  statt,  wenn  $f(x)$  zwischen
0  und  2n  gleich
d  ( a;*  cos  --- )
ist.     Denn  es  wird    in  dem  Integral  /  $$f(x)$$cosn{x  ---  a)dx  mit  wach-
0
sendem  n  der  Beitrag    derjenigen    Stelle,    wo  x  nahe  =y  ---  ist,  all-
gemein zu  reden,  zuletzt  unendlich  gross,  so  dass  das  Verhältniss  dieses
Integrals  zu
^ sin  1 2]/w  ---  wa -|- -^  I  ]/jr«
gegen  1  convergirt,  wie  man  auf  dem  gleich  anzugebenden  Wege  finden
wird.  Um  dabei  das  Beispiel  zu  verallgemeinern,  wodurch  das  Wesen
der  Sache  mehr  hervortritt,  setze  man
/  $f(x)$  dx  =  q)  {x)  cos  Tp  {x)
und  nehme  an,  dass  (p{x)  für  ein  unendlich  kleines  x  unendlich  klein
und  tl^ix)  unendlich  gross  werde,  übrigens  aber  diese  Functionen  nebst
ihren  Differentialquotienten  stetig  seien  und  nicht  unendlich  viele
Maxima  und  Minima  haben.     Es  wird  dann
f(^x)  =  (p{x)  tos  il}  (x)  ---  (p{x)  tjj'ix)  ^mip{x)
und
/  fix)  cos  n(x  ---  a)  dx
gleich  der  Summe  der  vier  Integrale  *
\textbackslash{}  I
g)\textbackslash{}x)  cos(i/'(ic)  +  n(a;  ---  a))  dx,
---  \textbackslash{}  i
(p{x)  tlf'(x)  sin  (^  (x)  +  w  (a;  ---  a))  dx.
Man  betrachte  nun,  il^ix)  positiv  genommen,  das  Glied
---  i  /
9  ^'  sin  (ip  (x)  -\textbackslash{}-  n{x  ---  a))  dx
\newpage

% --- Page 277 ---
und  untersuche  in  diesem  Integrale  die  Stelle,  wo  die  Zeichenwechsel
des  Sinus  sich  am  langsamsten  folgen.     Setzt  mau
ip(x)  -{-n(x  ---  a)  =  y,
so  geschieht  dies,  wo  ^^  =  0  ist,  und  also,
i)'(cc)  -f-  w  =  0
gesetzt,  für  x  =  cc.     Man  untersuche  also  das  Verhalten  des  Integrals
---  ^  I
'p{^)t'  sin  y  dx
für  den  Fall,  dass  a  für  ein  unendliches  n  unendlich  klein  wird,  und
führe  hiezu  y  als  Variable  ein.     Setzt  man
tp  (cc)  -\textbackslash{}-  n{cc  ---  a)  =  ß,
so  wird  für  ein  hinreichend  kleines  \pounds
und  zwar  ist  \textbullet{}^"{a)  positiv,  da  tp{x)  für  ein  unendlich  kleines  x  positiv
unendlich  wird;  es  wird  ferner
$f$
I  =  ^"(«)  {x-a)^$\pm$  y2rp'\textbackslash{}a){y-ß) ,
je  nachdem  x  ---  «  ^  0,  und
a-\textbackslash{}-t
---  i  /
9j(^)  ^'(^)  sin  y  dx
ß     /*+<(«) ,
.'^  +  V"{«)2
2
.jsin(,  +  ^)^|
dy   cp{ce)tp'{a)
0
Lässt  man  also  mit  wachsendem  n  die  Grösse  e  so   abnehmen,  dass
4)"{a)ee  unendlich  gross  wird,  so  wird,  falls
0
welches   bekanntlich  gleich   ist  sin  [ß  -\textbackslash{}-  -j)  j/jr,   nicht  Null   ist,   von
Grössen  niederer  Ordnung  abgesehen
\newpage

% --- Page 278 ---
262  XII.    üeber  die  Darstellbarkeit  einer  Function
«4"*
-  ij
V(^)^'(^)  ^m(t{x)  +  n{x  -  a))dx  =  -  sin  [ß  +  y)  ^\textasciitilde{}|=^-
\textbullet{}
a  ---  s
Es  wird  daher,  wenn  diese  Grösse  nicht  unendlich  klein  wird,  das  Ver-
hältniss  von
271
fix)  coBn{x  ---  a)  dx
0
zu  dieser  Grösse,  da  dessen  übrige  Bestandtheile  unendlich  klein  wer-
den, bei  unendlichem  Zunehmen  von  n  gegen  1  convergiren.
Nimmt  man  an,  dass  (p  (x)  und  ^'  (x)  für  ein  unendlich  kleines  x
mit  Potenzen  von  x  von  gleicher  Ordnung  sind  und  zwar  cp{x)  mit  'x^
und  il^'(x)  mit  x\textasciitilde{}f*'\textasciitilde{}^,  so  dass  v  >  0  und  ft^  0  äein  mUss,  so  wird
für  ein  unendliches  n
cp  (a)  tp'  (a)
y2ip"{a)
von  gleicher  Ordnung  mit  «     ^   und  daher  nicht  unendlich  klein,  wenn
fi^2v.     Ueberhaupt    aber  wird,    wenn  it;^'(a;)  oder,  was  damit  iden-
tisch ist,  wenn  ^\textasciitilde{}^  für  ein  unendlich  kleines  x  unendlich  gross  ist,
'  loga;  $^\circ$  '
sich  rp  (x)  immer  so  annehmen  lassen,  dass  für  ein  unendlich  kleines  x
tp  (x)  unendlich  klein,
ip'(x)  (p{x)  (fix)
(p{x)
]/iip"{x)        -,/      „   d
V      ^dx^'(x)         V'
2  lim
dx  ip'{x)         Y  xil)'{x)
aber  unendlich   gross   wird,   und  folglich    /  $f(x)$  dx  bis  an  a;  =  0  er-
streckt werden  kann,  während  *
2n
fix)  cos  nix  ---  d)  dx
I
0
für  ein  unendliches  n  nicht  unendlich   klein  wird.     Wie   man    sieht,
heben  sich  in  dem  Integrale  1  $$f(x)$$  dx  bei  unendlichem  Abnehmen  von
X
X  die  Zuwachse  des  Integrals,  obwohl  ihr  Verhältniss  zu  den  Aen-
derungen  von  x  sehr  rasch  wächst,  wegen  des  raschen  Zeichenwechsels
der  Function  fix)  einander  auf;  durch  das  Hinzutreten  des  Factors
cosM(a; ---  d)  aber  wird  hier  bewirkt,  dass  diese  Zuwachse  sich  summiren.
Ebenso  wohl  aber,  wie  hienach  für  eine  Function  trotz  der  durch-
gängigen Möglichkeit  der  Integration  die  Fourier'sche  Reihe  nicht
convergiren  und  selbst  ihr  Glied  zuletzt  unendlich  gross  werden  kann,
\newpage

% --- Page 279 ---
---  ebenso  wohl  können  trotz  der  durchgängigen  Unmöglichkeit  der
Integration  von  $$f(x)$$  zwischen  je  zwei  noch  so  nahen  Werthen  un-
endlich viele  Werthe  von  x  liegen,  für  welche  die  Reihe  iß  convergirt.
Ein  Beispiel  liefert,  {nx)  in  der  Bedeutung,  wie  oben  (Art.  6.)
genommen,  die  durch  die  Reihe
^j     n
1,00
gegebene  Function,  welche  für  jeden  rationalen  Werth  von  x  vorhanden
ist  und  sich  durch  die  trigonometrische  Reihe
1,  OC
wo  für  6  alle  Theiler  von  n  zu  setzen   sind,    darstellen  lässt,   welche
aber  in  keinem  noch    so  kleinen  Grössenintervall  zwischen  endlichen
Grenzen  enthalten   ist  und   folglich  nirgends   eine  Integration  zulässt.
Ein  anderes  Beispiel  erhält  man,  wenn  mau  in  den  Reihen
^CnCosnwa;,     2^  Cns\textbackslash{}.nnnx
0,  00  1,  CO
für  Cq,  Ci,  c.y,  . .  .  positive  Grössen  setzt,  welche  immer  abnehmen  und
$s$
zuletzt  unendlich  klein  werden,   während  2JCs    mit  n   unendlich  gross
1,»
wird.  Denn  wenn  das  Verhältniss  von  ij;  zu  2;r  rational  und  in  den
kleinsten  Zahlen  ausgedrückt,  ein  Bruch  mit  dem  Nenner  m  ist,  so
werden  ofifenbar  diese  Reihen  convergiren  oder  in's  Unendliche  wachsen,
je  nachdem
^  cos  nnx,      ^  sin  nnx
0,m --- 1  0,m  --- 1
gleich  Null  oder  nicht  gleich  Null  sind.  Beide  Fälle  aber  treten  nach
einem  bekannten  Theoreme  der  Kreistheilung$^{*)}$  zwischen  je  zwei  noch
so  engen  Grenzen  für  unendlich  viele  Werthe  von  x  ein.
In  einem  eben  so  grossen  Umfange  kann  die  Reihe  Sl  auch  con-
vergiren, ohne  dass  der  Werth  der  Reihe
,c+Ä,x-y^
1    dAn
y^i  nn  dx  '
welche  man  durch  Integration  jedes  Gliedes   aus  Sl  erhält,   durch  ein
noch  so  kleines  Grössenintervall  integrirt  werden  könnte.
Wenn  mau  z.  B.  den  Ausdruck
")  Disquis,  ar.  pag.  636  art.  356.     (Gauss   Werke  Bd.  I.  pag.  442.)
\newpage

% --- Page 280 ---
1,  00
WO  die  Logarithmen  so  zu  nehmen  sind,  dass  sie  für  g  =  0  ver-
schwinden ,  nach  steigenden  Potenzen  von  q  entwickelt  und  darin
qz=^i  setzt,  so  bildet  der  imaginäre  Theil  eine  trigonometrische  Reihe,
welche  zweimal  nach  x  differentiirt  in  jedem  Grössenintervall  unend-
lich oft  convergirt,  während  ihr  erster  DifiFerentialquotient  unendlich
oft  unendlich  wird.
In  demselben  Umfange,  d.  h.  zwischen  je  zwei  noch  so  nahen
Argumentwerthen  unendlich  oft,  kann  die  trigonometrische  Reihe  auch
selbst  dann  convergiren,  wenn  ihre  Coefficienten  nicht  zuletzt  unend-
lich klein  werden.  Ein  einfaches  Beispiel  einer  solchen  Reihe  bildet
die  unendliche  Reihe  E  sm(n\textbackslash{}xn),  wo  n\textbackslash{},  wie  gebräuchlich,
1,00
=  1.2.3...w,
welche  nicht  bloss  für  jeden  rationalen  Werth  von  x  convergirt,  indem
sie  sich  in  eine  endliche  verwandelt,  sondern  auch  für  eine  unendliche
Anzahl  von  irrationalen ,  von    denen  die  einfachsten  sind  sin  1 ,  cos  1,
1
$e$
---  und  deren  Vielfache,  ungerade  Vielfache  von  e,  --- r --- ,  u.  s.  w.  (")
\newpage

% --- Page 281 ---
Geschichte   der  Frage   über   die   Darstellbarkeit   einer  Function    durch   eine
trigonometrische  Reihe.
Ursprung  der  Frage  in  dem  Streite  über  die  Tragweite  der  d'Alem-
b er t' sehen  und  Bernoulli' sehen  Lösung  des  Problems  der  schwin-      »,
genden  Saiten  im  Jahre  1753.    Ansichten  von  Euler,  d'Alembert,
Lagrange   227
Richtige    Ansicht    Fourier's,    bekämpft    von   Lagrange.     1807.
-Cauchy.     1826
.232
Erledigung  der  Frage  durch  Dirichlet  für  die  in  der  Natur  vor-
kommenden Functionen.     1829.     Dirksen.     Bessel.     1839.    .     .     .  235
üeber  den  Begriff  eines  bestimmten  Integrals  und  den  Umfang  seiner  Gültigkeit.
\subsection*{\S.  5.     Bedingungen  der  Möglichkeit  eines  bestimmten  Integrals     .     .     .     .  240}
Untersuchung  der  Darstellbarkeit  einer  Function  durch  eine  trigonometrische
Reihe,  ohne  besondere  Voraussetzungen  über  die  Natur  der  Function.
\section*{I.     Ueber  die  Darstellbarkeit  einer  Function  durch  eine  trigonometrische  Reihe,}
deren  Coefficienten  zuletzt  unendlich  klein  werden.
\subsection*{\S.  8.     Beweise  einiger  für  diese  Untersuchung  wichtigen  Sätze      ....  245}
\subsection*{\S.  9.     Bedingungen  für   die  Darstellbarkeit  einer  Function  durch  eine  tri-}
gonometrische Reihe  mit  in's  Unendliche  abnehmenden  Coefficienten  251
\subsection*{\S.  10.  Die  Coefficienten  der  Fourier' sehen  Reihe  werden  zuletzt  unend-}
lich klein,  wenn  die  darzustellende  Function  durchgehends  endlich
bleibt  und  eine  Integration  zulässt   253
\section*{II.     Ueber  die  Darstellbarkeit  einer  Function  durch  eine  trigonometrische  Reihe}
mit  nicht  in's  Unendliche  abnehmenden  Coefficienten.
\subsection*{\S.  11.  Zurückführung  dieses  Falles  auf  den  vorigen   255}
Betrachtung  besonderer  Fälle.
\subsection*{\S.  12.  Functionen,  welche  nicht  unendlich  viele  Maxima  und  Minima  haben  257}
\subsection*{\S.  13.  Functionen,  welche  unendlich  viele  Maxima  und  Minima  haben  .    .  260}
\newpage

% --- Page 282 ---
(4)  (Zu  Seite  237.  Anmerkung.)  Nehmen  wir  an,  dass  die  Function  /"(a;)  in  dem
Intervalle  A  zwischen  a;  und  aji  >\textbullet{}  x  nicht  wachse,  und  bezeichnen  wir  mit  g
die  obere  Grenze  der  Werthe,  die  f{x-\textbackslash{}-^  für  0<[KA  annimmt,  d.  h.  einen
Werth,  der  von  keinem  dieser  Functionswerthe  überschritten,  aber  mit  jedem
beliebigen  Grad  von  Annäherung  erreicht  wird,  so  wird  g  ---  /"(x -}- |)  mit
wachsendem  \%  niemals  abnehmen,  während  es  doch  beliebig  klein  werden  soll,
d.  h.  es  ist  Lim  {g  ---  f{x  -)-  |))  =  Ö,  flr  =  f{x  -f  0).    Der  Satz,  dass  ein  System
1=0     :  $\blacksquare$    .'
reeller  Zahlen  \textcopyright,  dessen  in  endlicher,  oder  unendlicher  Zahl  vorhandene  In-
dividuen s  einen  endlichen  Zahlwerth  nicht  übersteigen,  eine  obere  Grenze  hat,
ist  wohl  zuerst  von  Weierstrass  präcis  ausgesprochen  und  bewiesen  (vgl.
0.  Biermann,  Theorie  der  analytischen  Functionen  \S.  16.     Leipzig,   Teubner,
. .  1884).  Sehr  einfach  ist  der  Beweis  auf  Grund  der  Anschauungen  von  Dedekind
über  Irrationalzahlen  (Stetigkeit  und  irrationale  Zahlen,  Braunschweig,  Vieweg,
1872).  Theilt  man  nämlich  die  reelle  Zahlenreihe  in  zwei  Theile  A  und  B
.  ein,  so  dass  jede  Zahl  a  in  A  von  Zahlen  des  Systems  \textcopyright  überschritten,  jede
Zahl  \&:  in  .B  nicht  überschritten  wird,  so  werden  diese  beiden  Theile  A  und  Iß
durch  eine  existirende  Zahl  gf. getrennt,  der  offenbar  die  charäkterischen  Merk-
\_   male  der  oberen  Grenze. von  (3  zukommen.
(2)  (Zu  Seite.  241.)  ;  Es .  findet  sich  hierzu  eine  fragmentarische  handschriftliche
Bemerkung  vonRiemann,  die;  wir  folgendermassen  herziistellen  versuchen,
da  sie  ,  zur  Vervollständigung  des  Beweises  nothwendig  ist,  dass  das  .Ver-
schwinden von  A  mit  d  auch  die  für  die  Convergenz  von  /?  ausreichende  Be-
dingung ist.  Es  könnte  scheinen  als  ob,  wenn  auch  bei  zwei  verschiedenen
Eintheiluhgen,  in  denen  die  Intervalle  d',  ö"  kleiner  als  <Z  und ;  folglich  der
Unterschied  zwischen  dem  grössten  und  kleinsten  Werthe  (oberer  uüd  unterer
Grenze)  der  Summe  ä,  die  für  die  beiden  Eintheilungen  mit  :iS','S^"  gezeichnet
sei,  kleiner  als  eine  gegebene  Grösse  \pounds  ist,  doch  die  Summen  S',  S"  selbst  um
ein  endlichesStück  auseinander  liegen  könnten.  Um  die  Unmöglichkeit  hier-
von einzusehen,  bilde  man  eine  dritte  Eintheilung  d,  der  die  Summe  ä  ent-
spreche, indem  man  d'  und  8"  gleichzeitig  ausführt.  Da  nun  jedes  Element  8'
aus  einer  ganzen  Anzahl  von  Elementen  8  besteht,  so  wird,  wenn  ein  beliebiger
Werth  von  S  betrachtet  wird,  die  Summe  der  diesen  Elementen  8  entsprechenden
Glieder  von  /S'  zwischen  dem  grössten  und  kleinsten  Werth  des  dem  Element  8'
entsprechenden  Gliedes  von  >S"  liegen,  und  folglich  auch  die  ganze  Summe  S
zwischen  dem  grössten  und  kleinsten  Werth  von  S'  und  ebenso  auch  zwischen
dem  grössten  und  kleinsten  Werth  von  S"  \textbackslash{}  folglich  können  /S,  B\textbackslash{}  S"  nicht
um  mehr  als  s  von  einander  verschieden  sein.
\newpage

% --- Page 283 ---
(3)  (Zu  Seite  243.)    Dass  jode  endliche  Function  $$f(x)$$,  die  zwischen  den  Grenzen
a  und  h  nicht  wächst,  und  folglich  jede  Function,  die  nicht  unendlich  viele
Maxima  und  Minima  hat,  einer  Integration  fähig  ist,  beweist  man  so.
Sei
a  =  a;,  <  a;,  <  a^a  \textbullet{}  \textbullet{}  \textbullet{}  <  aj^  =  \& ,
"i   "^  '^'2         '*'*l  >     "2  =  ^3  ^2>  $\blacksquare$  $\blacksquare$  $\blacksquare$  \%--- 1  "^  ^n         ^»---1,
A  +  A  +  -  \textbullet{}  \textbullet{}  +  i>«\_i  =  /\textbullet{}(«)  -  /$\blacksquare$(\&)\textbullet{}
Da  nach  Voraussetzung  $$f(x)$$  nicht  wächst,  so  sind  die  Grössen  Dj ,  D^ ,  \textbullet{}  \textbullet{}  -  -D„\_i
die  grössten  Schwankungen  in  den  Intervallen  5^ ,  d,^  . . . ,  ^„\_i  und  alle  positiv
oder  wenigstens  keines  von  ihnen  negativ.  Ist  m  die  Anzahl  derjenigen  Inter-
valle, in  denen  D  >  ff,  so  ist  ma  <.  f{a)  ---  f{b)  oder
m  <  \textasciitilde{}^-^   ---^-^  $\blacksquare$
$a$
Sind  also  die  sämmtlichen  Intervalle  8  kleiner  als  d,   so  ist  die  Gesammt-
grösse  der  Intervalle,   von  denen  die  grösste  Schwankung  grösser  als  a  ist,
<<   ---^d  und  wird  also,  w.  z.  b.  w. ,  mit  d  zugleich  unendlich  klein.
(4)  (Zu  Seite  251.)    Die  hier  gebrauchten  Grössen  B  ^^  haben  den  Ausdruck
$\blacksquare$^ft  +  n  =  i  COS  (jti  -\textbackslash{}-  n){x  ---  ä)  (a^  sin  na  -f  \&„  cos  wa)
-|-  \textbackslash{}  sin  (jit  -f"  '^)  (^  --- '  «)  (<"«  cos  na  ---  b^  sin  wa^
-B„\_^  =  i  cos  (fi  ---  n)  {x  ---  a)  (a^  sin  nä  +  \&^  cos  na\textbackslash{}
---  \textbackslash{}  sin  (m  ---  n)  {x  ---  a)  (a^  cos  na  ---  \&^  sin  wo") .
Zur  Ergänzung  ist  noch  zu  beweisen,  dass  auch
$c$
f^f*  /
\textbackslash{}G  -{-  G'x  -f  A  ---  1  cos  fi  (a;  ---  a)  i  (a;)  dx
$b$
den  Grenzwerth  0  hat.    Dies  erreicht  man  wohl  am  einfachsten ,  wenn  man
B
(ifi  dx'  '
-,        //^  ;    3^0    I >w     I.  j   *^\textbackslash{}            /            \textbackslash{}        «  //T    1     j     \textbackslash{}sinu(a;  ---  a)
^  \textbackslash{}         "ii^  +  Ca;'+ J[,  --j  cosft  (a;  -  a)  -  2  (^G  +  ^,a;j
- --- '^   ^
setzt  und  eine  zweimalige  partielle  Integration  anwendet.    Dass  Integrale  wie
c  c
I  cos  ^i  (x  ---  a)  i"  (a;)  dx  ,         l  sin  ft  {x  ---  a)  X."  (x)  dx
ö  b
mit  unendlich  wachsendem  (i.  verschwinden,  kann  man  entweder  nach  der
Dirichlet'schen  Methode,  oder  einfacher  mittels  des  du  Bois-Reymond'schen
Mittelwerthsatzes  beweisen,  wonach,  wenn  qo  (a;)  eine  zwischen  den  Grenzen  b
und  c  nicht  wachsende  oder  nicht  abnehmende  Function  und  |  ein  zwischen
b  und  c  gelegener  Werth  ist,
/  $$f(x)$$  cp  (x)  dx  =  cp  ib)   I
fix)  dx  -\textbackslash{}-  (p  (c)   I  $$f(x)$$  dx  .
\textbackslash{}C  -\textbackslash{}-  C x-\textbackslash{}-  A^  \textasciitilde{}\textasciitilde{}-j  cos  (i(x  ---  a)  =
\newpage

% --- Page 284 ---
(5)  (Zu  Seite  262.)     Die  unter  II.   aufgestellten  Sätze  bedürfen  einer  Erläuterung:
Da  die  Function  $$f(x)$$  um  2«  periodisch  angenommen  ist,  so  muss
F {x -{- In)  ---  F {x)  =  (f  {x)
die  Eigenschaft  haben,  dass
cp{x-\textbackslash{}-a-\textbackslash{}-ß)  ---  cp{x-[-a  ---  ß)  ---  cp  {x  ---  a  -\textbackslash{}-  ^)  -^  (f  {x  ---  «  ---  ß)
laß       \textasciitilde{}\textasciitilde{}
unter   der  im  Text  gemachten
Voraussetzung
sich  mit  a  und  ß  der  Grenze  0
nähert.
Es  ist  daher  g?  {x)  eine  lineare  Function
von  aj,  und  folglich
lassen
sich  die  Constanten
C ,  A^,  so  bestimmen,
dass
^{x)^F{x)-C'x-A,'^
eine  um  2n  periodische  Function  von  x  ist.
Nun  ist  über  die  Function  F {x)  weiter  die  Voraussetzung  gemacht,  dass
für  beliebige  Grenzen  h,  c
0
HH  j  F{x)  cos  fi  (x  ---  a)X  (x)  dx
$b$
mit  unendlich  wachsendem  ft  sich  der  Grenze  0  nähere ,  wenn  X  {x)  den  im
Text  angegebenen  Bedingungen  genügt,  woraus  folgt,  dass  unter  den  gleichen
Voraussetzungen
^ft   /
<
*  {x)  cos  ^  {x  ---  a)  l  {x)  dx
$b$
sich  der  Grenze  0  nähert.
Es  sei  nun  h  <i  ---  n,  c  '^  n,  und  man  nehme,  was  zulässig  ist,  X  (x)  im
Intervall  von  ---  tt  bis  -\textbackslash{}-n  =  1,  so  folgt,  dass  auch:
---  7t  G
Uli   I  ^  (a;)  cos  (i{x  ---  a)X  (x)  dx  -{-  (ifi   1   ^  (x)  cos  {i{x  ---  a)X  {x)  dx
b  n
-|-  /Aft   I   $  (ic)  cos  fi  (ic  ---  a)  dx
--- n
Null  zur  Grenze  hat.    Nun  kann  man,  wenn  ft  eine  ganze  Zahl  n   ist,  mit
Rücksicht  auf  die  Periodicität  von  *  (a;)  für  diese  Summe  setzen:
c  +«
nn  I
$  (x)  cos  n(x  ---  a)  X^  (x)  dx  -\textbackslash{}-  nn  j
^  {x)  cos  n(x  ---  a)dx,
6  +  2«  ---7t
wenn  in  dem  Intervall  von  \&  -f  2«  bis  n  Xi{x)  =  X{x  ---  2n)  und  in  dem
Intervall  von  n  bis  c  X^{x)  =  X{x)  ist,  so  dass  X^ix)  zwischen  den  Grenzen
6  4-  271!  und  c  den  Voraussetzungen  über  die  Function  X{x)  genügt.  Demnach
hat  das  erste  Glied  der  obigen  Summe  für  sich  den  Grenzwerth  0,  und  folglich
ist  auch  der  Grenzwerth  von
$  (x)  cos  ii{x  ---  a)  dx
(ifi  1
<
gleich  Null.
\newpage

% --- Page 285 ---
(6)  (Zu  Seite  253.)  Hier  scheint  für  die  Function  X  (x)  die  Bedingung  hinzugefügt
werden  zu  müssen,  dass  sie  sich  nach  dem  Interrall  2«  periodisch  wiederholt
(die  mit  der  nachher  gemachten  Annahme  verträglich  ist).  In  der  That  wurde
z.  B.   das  in  Rede  stehende   Integral  nicht  sich  der  Grenze  0  nähern,  wenn
F{t)  ---  C't  ---  Af,  --- -  =  const.   und   X  {t)  =  (x  ---  ty  gesetzt  würde.     Dagegen
lässt  sich  unter  der  Voraussetzung  der  Periodicität  von  X  (x)  das  Verschwinden
dieses  Integrals  durch  Ausführung  der  Differentiation
dd
sin   \_L\_(a:\_<)
.    X  ---  *
sin.
dt'
durch   Anwendung  des  Satzes  3,  Art,  8   und   eines  ähnlichen  Verfahrens  wie
in  der  Anmerkung  (5)  leicht  darthun.
Gingen  die  Anmerkungen  (5)  (6),  die  in  der  ersten  Auflage  dieser  Gesammt-
ausgabe  die  Nummern  (1)  (2)  hatten,  hat  Ascoli  in  einer  Abhandlung  über  trigo-
nometrische Reihen  (Accademia  dei  Lincei  1880)  verschiedene  Bedenken  erhoben.
Sie  sollen  gleichwohl  hier  unverändert  bleiben;  nur  möge  Folgendes  beigefügt
werden.
Der  Beweis  des  Satzes ,  nach  dem  die  in  der  Anmerkung  (5)  mit  cp  (x)  be-
zeichnete Function  eine  lineare  sein  muss  (für  den  ich  auf  eine  Abhandlung  von
G.  Cantor  in  Crelle's  Journal  Bd.  72,  S.  141  verweise),  setzt  allerdings  voraus,
dass  die  Function  $$f(x)$$  für  jeden  Werth  von  x  existirt  (also  auch  endlich  ist).
Mir  scheinen  aber  die  Nummern  I,  II  des  Art.  9  überhaupt  nur  dann  ganz  ver-
ständlich,  wenn  diese  Existenz  vorausgesetzt  wird,    selbst  dann,   wenn  man  wie
XX
Ascoli  will,  die  Forderung,  durch  Hinzufügung  eines  Ausdrucks  ---  C'x  ---  -^o  "^
in  eine  periodische  Function  verwandelt  zu  werden,  unter  die  Bedingungen  für
die  Function  F  (x)  mit  aufnimmt.  Lässt  man  die  Voraussetzung  der  durchgängigen
Existenz  von  $$f(x)$$  fallen,  so  kann  es  unendlieh  viele  verschiedene  Functionen  F(x)
geben,  die  sich  nicht  blos  um  lineare  Ausdrücke  von  einander  unterscheiden.
Art.  9,  III  behält  allerdings  seine  Bedeutung  auch  dann  noch,  wenn  die  durch-
gehende Existenz  von  $f(x)$  nicht  vorausgesetzt  wird,  sondern  wie  in  Art.  8  F (x)
durch  die  Reihe  C   \textasciitilde{}-  ---  "---    $\blacksquare$     \textbullet{}  definirt  ist.
14  9
Bei  der  Anmerkung  (6)  ist  zuzugeben,  dass  es  genügt,   da  die  Function  X{t)
in  den  Formeln  des  Textes  nur  in  dem  Intervall  --- n  bis  -\textbackslash{}- n  vorkommt,  nicht
die  Periodicität   von   X  (t)  und  X' {t)^    sondern  nur  die  Formeln  X{jt)  =  X( --- «),
X' {n)  =  X' ( --- n)    vorauszusetzen,    also  nicht  eigentlich  die  Periodicität,    sondern
die    Möglichkeit   der    stetigen    periodischen   Fortsetzung.     Da  aber  die  Function
F{t)  -  C't  ---  ^- 1^  auf  Seite  261  durch  die  Reihe  C  ---  ^
--- -^- --- '^^
$\blacksquare$  $\blacksquare$  $\blacksquare$  und
Ä         A        A,
nicht  wie  Ascoli  annimmt,   durch   -'   ---  -^  \textbullet{}  \textbullet{}  \textbullet{}  definirt  ist,  so  ist  die
14  9
Annahme  des  Beispiels,  dass  F(t)  ---  et   ^  t*  eine  von  Null  verschiedene  Con-
staute  sei,  sehr  wohl  zulässig.
\newpage

% --- Page 286 ---
Auch  das  Verfahren,  das  ich  nach  Analogie  der  Note  (5)   zum  Beweis  des
Verschwindens  des  Integrals
sm  ---^ ---  {x  ---  t)
dd   '---^
,  „  .    X  ---  t
^Jt(t)\textasciitilde{}\textasciitilde{} --- ^.$\blacksquare$- --- \textasciitilde{}m<u
---  n
angewandt
habe,  muss  ich  etwas  genauer
darlegen.
Führt  man  die  Differentiation  unter  dem  Integral  aus,  so  erhält  man  einen
2n-\textbackslash{}-  1
mehrgliedrigen  Ausdruck,   dessen   einer  Term,  wenn  zur  Abkürzung      -- ---  =  fi
gesetzt  wird,  lautet
+  7t
d'
^(t)      '"^^^^       sin  (1  (x  ---  t)  dt ,
\textbullet{}    X  ---  t
sin  ---:
oder  wenn  l  (t)  =  l^  (f)  sin  --- - ---  und  x  =  a\textasciitilde{}{-7t.  gesetzt  wird,
,(---  1)"  (i^  1
*(<)  l,  (t)  cos  (i(a  ---  t)dt.
---n
Man  wähle  nun  \&,  c  so,  dass  das  Intervall  von  \&  bis  c  das  Intervall  von  ---  jr
bis  4"  '^   einschliesst  und  bestimme   im   ersteren  Intervall   eine  Function  X  (<)   so,
dass  zwischen  ---  n  und  \textbullet{}\textbackslash{}-n  i  (t)  =  i^  (f) ,  an  den  Grenzen  aber  X  {t)  und  X'  (t)
verischwinden,  ferner  eine  Function  X^  (t)  in  dem  Intervall  von  b  -\textbackslash{}-  2n  his  e  so,
dass    zwischen    \&  -|-  2«    und   «   X^  (t)  =  ---  X(t  ---  27r)   und    zwischen  -n   und    c
Xj  (t)  =  X  {t)  wird,  was  also  voraussetzt,  dass  l^  (n)  =  ---  X^  {--- ti)  ,  X'^  («)  =  X'j^  (---  n)
sei.    Dann  ergiebt  sich  wie  in  der  Anmerkung  (5)
4-7r  c
(i^  1  *(t)  X^  (t)  cos  n(a  ---  t)dt  =  (i^   /  * (t)  X  (t)  cos  (i{a  ---  t)dt  ,
^(t)  X^  (t)  cos  [i(a  ---  t)dt
"'/'
und  die  beiden  Bestandtheile  rechts  verschwinden  nach  Satz  3,  Art.  8  mit  un-
endlich wachsendem  (x.  Ebenso  verfährt  man  mit  den  übrigen  Bestandtheileu  des
vorliegenden  Integrals.
(7)  (Zu  Seite  260.)  Es  möge  hier  auf  die  Arbeiten  von  P.  du  Bois-Eeymond  hin-
gewiesen werden,  die  nach  Riemann  die  Theorie  der  trigonometrischen  Reihen
noch  wesentlich  gefördert  haben.  Es  ist  dort  durch  Beispiele  nachgewiesen,
dass  es  selbst  durchaus  endliche  und  stetige  Functionen  mit  unendlich  vielen
Maxima  und  Minima  giebt,  die  eine  Darstellung  durch  trigonometrische  Reihen
nicht  gestatten.
(8)  (Zu  Seite  263.)  Das  Zeichen  Z*^  ---  (---  1)**  ist  als  eine  Summe  von  positiven
und  negativen  Einheiten  zu  verstehen,  so  dass  jedem  geraden  Theiler  von  w
ein    negatives,   jedem  ungeraden  ein  positives   Glied  entspricht.  .  Man  findet
\newpage

% --- Page 287 ---
diese    Entwicklung  (wenn    auch  auf  einem  nicht  ganz    eiuwurfafreien  Wege),
wenn  man  die  Function  {x)  durch  die  bekannte  Formel
$m$
\_  ^  c\_  lY^  sin  2  m  nx
1,00
ausdrückt,  dies  in  die  Summe  S-^-  einsetzt  und  die  Ordnung  der  Summa-
tionen  vertauscht.
(9)  (Zu  Seite  264.)  Der  Werth  x  =  l  \textbackslash{}e  ---  -j  gehört,  wie  Genochi  in  einem
diese  Beispiele  betreffenden  Aufsatz  bemerkt  (Intorno  ad  alcune  serie,  Torino
1876)  nicht  zu  den  Werthen  von  x,  für  welche  die  Reihe  ^  sin(w!a;jr)  con-
1,00
vergirt.     Aber    auch  für  x  =  \textbackslash{}  \textbackslash{}e   )   ist  die  Reihe  nicht,  wie  Genochi
angiebt,  convergent.
\newpage

% --- Page 288 ---
(Aus  dem  dreizehnten  Bande  der  Abhandlungen   der  Königlichen  Gesellschaft  der
Wissenschaften  zu  Göttingen.)$^{*)}$
Plan  der  Untersuch.ung.
Bekanntlich  setzt  die  Geometrie  sowohl  den  Begriff  des  Raumes,
als  die  ersten  Grundbegriffe  für  die  Constructionen  im  Räume  als  etwas
Gegebenes  voraus.  Sie  giebt  von  ihnen  nur  Nominaldefinitionen,  wäh-
rend die  wesentlichen  Bestimmungen  in  Form  von  Axiomen  auftreten.
Das  Verhältniss  dieser  Voraussetzungen  bleibt  dabei  im  Dunkeln;  man
sieht  weder  ein,  ob  und  in  wie  weit  ihre  Verbindung  nothwendig,
noch  a  priori,  ob  sie  möglich  ist.
Diese  Dunkelheit  wurde  auch  von  Euklid  bis  auf  Legendr e,  um
den  berühmtesten  neueren  Bearbeiter  der  Geometrie  zu  nennen,  weder
von  den  Mathematikern,  noch  von  den  Philosophen,  welche  sich  da-
mit beschäftigten,  gehoben.  Es  hatte  dies  seinen  Grund  wohl  darin,
dass  der  allgemeine  Begriff  mehrfach  ausgedehnter  Grössen,  unter
welchem  die  Raumgrössen  enthalten  sind,  ganz  unbearbeitet  blieb.
Ich  habe  mir  daher  zunächst  die  Aufgabe  gestellt,  den  Begriff  einer
mehrfach  ausgedehnten  Grosse  aus  allgemeinen  Grössenbegriffen  zu
construiren.  Es  wird  daraus  hervorgehen,  dass  eine  mehrfach  aus-
gedehnte Grösse  verschiedener  Massverhältnisse  fähig  ist  und  der  Raum
also  nur  einen  besonderen  Fall  einer  dreifach  ausgedehnten  Grösse
bildet.     Hiervon  aber  ist  eine  nothwendige  Folge,    dass  die  Sätze  der
$^{*)}$  Diese  Abhandlung  ist  am  10.  Juni  1854  von  dem  Verfasser  bei  dem  zum
Zweck  seiner  Habilitation  veranstalteten  Colloquium  mit  der  philosophischen
Facultät  zu  Göttingen  vorgelesen  worden.  Hieraus  erklärt  sich  die  Form  der  Dar-
stellung, in  welcher  die  analytischen  Untersuchungen  nur  angedeutet  werden
konnten;  einige  Ausführungen  derselben  findet  man  in  der  Beantwortung  der
Pariser  Preisaufgabe  nebst  den  Anmerkungen  zu  derselben.
\newpage

% --- Page 289 ---
XIIT.    Ueber  die  Hypothesen,  welche  der  Geometrie  zu  Grunde  liegen.    273
Geometrie  sich  nicht  aus  allgemeinen  Grössenbegriffen  ableiten  lassen,
sondern  dass  diejenigen  Eigenschaften,  durch  welche  sich  der  Raum
von  anderen  denkbaren  dreifach  ausgedehnten  Grössen  unterscheidet,
nur  aus  der  Erfahrung  entnommen  werden  können.  Hieraus  entsteht
die  Aufgabe,  die  einfachsten  Thatsachen  aufzusuchen,  aus  denen  sich
die  Massverhältnisse  des  Raumes  bestimmen  lassen  ---  eine  Aufgabe,
die  der  Natur  der  Sache  nach  nicht  völlig  bestimmt  ist;  denn  es  lassen
sich  mehrere  Systeme  einfacher  Thatsachen  angeben,  welche  zur  Be-
stimmung der  Massverhältnisse  des  Raumes  hinreichen;  am  wichtigsten
ist  für  den  gegenwärtigen  Zweck  das  von  Euklid  zu  Grunde  gelegte.
Diese  Thatsachen  sind  wie  alle  Thatsachen  nicht  nothwendig,  sondern
nur  von  empirischer  Gewissheit,  sie  sind  Hypothesen;  man  kann  also
ihre  Wahrscheinlichkeit,  welche  innerhalb  der  Grenzen  der  Beobachtuug
allerdings  sehr  gross  ist,  untersuchen  und  hienach  über  die  Zulässig-
keit  ihrer  Ausdehnung  jenseits  der  Grenzen  der  Beobachtung,  sowohl
nach  der  Seite  des  Unmessbargrossen,  als  nach  der  Seite  des  Un-
messbarkleinen  urtheilen.
\section*{I.    Begriff  einer  nfaoh  ausgedehnten  Grösse.}
Indem  ich  nun  von  diesen  Aufgaben  zunächst  die  erste,  die  Ent-
wicklung des  Begriffs  mehrfach  ausgedehnter  Grössen,  zu  lösen  ver-
suche, glaube  ich  um  so  mehr  auf  eine  nachsichtige  Beurtheilung  An-
spruch machen  zu  dürfen,  da  ich  in  dergleichen  Arbeiten  philosophischer
Natur,  wo  die  Schwierigkeiten  mehr  in  den  Begriffen,  als  in  der  Con-
struction  liegen,  wenig  geübt  bin  und  ich  ausser  einigen  ganz  kurzen
Andeutungen,  welche  Herr  Geheimer  Hofrath  Gauss  in  der  zweiten
Abhandlung  über  die  biquadratischen  Reste,  in  den  Göttingenschen
gelehrten  Anzeigen  und  in  seiner  Jubiläumsschrift  darüber  gegeben
hat,  und  einigen  philosophischen  Untersuchungen  Herbart's,  durchaus
keine  Vorarbeiten  benutzen  konnte.
Grössenbegriffe  sind  nur  da  möglich,  wo  sich  ein  allgemeiner  Be-
griff vorfindet,  der  verschiedene  Bestimmungsweisen  zulässt.  Je  nach-
dem unter  diesen  Bestimmungsweisen  von  einer  zu  einer  andern  ^iu
stetiger  üebergang  stattfindet  oder  nicht,  bilden  sie  eine  stetige  oder
discrete  Mannigfaltigkeit;  die  einzelnen  Bestimmungsweisen  heissen  im
er\&tern  Falle  Punkte,  im  letztern  Elemente  dieser  Mannigfaltigkeit.
Begriffe,  deren  Bestimmungsweisen  eine  discrete  Mannigfaltigkeit  bil-
den, sind  so  häufig,  dass  sich  für  beliebig  gegebene  Dinge  wenigstens
Kiemann's  gesammelte  matliemutisclie  Werke.  18
\newpage

% --- Page 290 ---
in  den  gebildeteren  Sprachen  immer  ein  Begriff  auffinden  lässt,  unter
welchem  sie  enthalten  sind  (und  die  Mathematiker  konnten  daher  in
der  Lehre  von  den  discreten  Grössen  unbedenklich  von  der  Forderung
ausgehen,  gegebene  Dinge  als  gleichartig  zu  betrachten),  dagegen  sind
die  Veranlassungen  zur  Bildung  von  Begriffen,  deren  Bestimmungs-
weisen eine  stetige  Mannigfaltigkeit  bilden,  im  gemeinen  Leben  so
selten,  dass  die  Orte  der  Sinnengegenstände  und  die  Farben  wohl  die
einzigen  einfachen  Begriffe  sind,  deren  Bestimmungsweisen  eine  mehr-
fach ausgedehnte  Mannigfaltigkeit  bilden.  Häufigere  Veranlassung  zur
Erzeugung  und  Ausbildung  dieser  Begriffe  findet  sich  erst  in  der
höhern  Mathematik.
Bestimmte,  durch  ein  Merkmal  oder  eine  Grenze  unterschiedene
Theile  einer  Mannigfaltigkeit  heissen  Quanta.  Ihre  Vergleichung  der
Quantität  nach  geschieht  bei  den  discreten  Grössen  durch  Zählung,  bei
den  stetigen  durch  Messung.  Das  Messen  besteht  in  einem  Aufeinander-
legen der  zu  vergleichenden  Grössen;  zum  Messen  wird  also  ein  Mittel
erfordert,  die  eine  Grösse  als  Massstab  für  die  andere  fortzutragen.
Fehlt  dieses,  so  kann  man  zwei  Grössen  nur  vergleichen,  wenn  die
eine  ein  Theil  der  andern  ist,  und  auch  dann  nur  das  Mehr  oder  Min-
der, nicht  das  Wieviel  entscheiden.  Die  Untersuchungen,  welche  sich
in  diesem  Falle  über  sie  anstellen  lassen,  bilden  einen  allgemeinen  von
Massbestimmungen  unabhängigen  Theil  der  Grössenlehre,  wo  die  Grössen
nicht  als  unabhängig  von  der  Lage  existirend  und  nicht  als  durch  eine
Einheit  ausdrückbar,  sondern  als  Gebiete  in  einer  Mannigfaltigkeit  be-
trachtet werden.  Solche  Untersuchungen  sind  für  mehrere  Theile  der
Mathematik,  namentlich  für  die  Behandlung  der  mehrwerthigen  ana-
lytischen Functionen  ein  Bedürfniss  geworden,  und  der  Mangel  der-
selben ist  wohl  eine  Hauptursache,  dass  der  berühmte  Abel 'sehe  Satz
und  die  Leistungen  von  Lagrange,  P/aff,  Jacobi  für  die  allgemeine
Theorie  der  Differentialgleichungen  so  lange  unfruchtbar  geblieben  sind.
Für  den  gegenwärtigen  Zweck  genügt  es,  aus  diesem  allgemeinen  Theile
der  Lehre  von  den  ausgedehnten  Grössen,  wo  weiter  nichts  vorausgesetzt
wird,  als  was  in  dem  Begriffe  derselben  schon  enthalten  ist,  zwei
Punkte  hervorzuheben,  wovon  der  erste  die  Erzeugung  des  Begriffs
einer  mehrfach  ausgedehnten  Mannigfaltigkeit,  der  zweite  die  Zurück-
führung  der  Ortsbestimmungen  in  einer  gegebenen  Mannigfaltigkeit
auf  Quantitätsbestimmungen  betrifft  und  das  wesentliche  Kennzeichen
einer  w fachen  Ausdehnung  deutlich  machen  wird.
\newpage

% --- Page 291 ---
XIII.   Ueber  die  Hypothesen,  welche  der  Geometrie  zu  Grnnde  liegen.    275
2.
Geht  man  bei  einem  Begriffe,  dessen  Bestimmungsweisen  eine
stetige  Mannigfaltigkeit  bilden,  von  einer  Bestimmungsweise  auf  eine
bestimmte  Art  zu  einer  andern  über,  so  bilden  die  durchlaufenen  Be-
stimmungsweisen eine  einfach  ausgedehnte  Mannigfaltigkeit,  deren
wesentliches  Kennzeichen  ist,  dass  in  ihr  von  einem  Punkte  nur  nach
zwei  Seiten,  vorwärts  oder  rückwärts,  ein  stetiger  Fortgang  möglich
ist.  Denkt  man  sich  nun,  dass  diese  Mannigfaltigkeit  wieder  in  eine
andere,  völlig  verschiedene,  übergeht,  und  zwar  wieder  auf  bestimmte
Art,  d.  h.  so,  dass  jeder  Punkt  in  einen  bestimmten  Punkt  der  andern
übergeht,  so  bilden  sämmtliche  so  erhaltene  Besfimmungsweisen  eine
zweifach  ausgedehnte  Mannigfaltigkeit.  In  ähnlicher  Weise  erhält  man
eine  dreifach  ausgedehnte  Mannigfaltigkeit,  wenn  man  sich  vorstellt,
dass  eine  zweifach  ausgedehnte  in  eine  völlig  verschiedene  auf  be-
stimmte Art  übergeht,  und  es  ist  leicht  zu  sehen,  wie  man  diese  Con-
struction  fortsetzen  kann.  Wenn  man,  anstatt  den  Begriff  als  be-
stimmbar, seinen  Gegenstand  als  veränderlich  betrachtet,  so  kann  diese
Construction  bezeichnet  werden  als  eine  Zusammensetzung  einer  Ver-
änderlichkeit von  n  -{-  1  Dimensionen  aus  einer  Veränderlichkeit  von
n  Dimensionen  und  aus  einer  Veränderlichkeit  von  Einer  Dimension.
Ich  werde  nun  zeigen,  wie  man  umgekehrt  eine  Veränderlichkeit,
deren  Gebiet  gegeben  ist,  in  eine  Veränderlichkeit  von  einer  Dimension
und  eine  Veränderlichkeit  von  weniger  Dimensionen  zerlegen  kann.
Zu  diesem  Ende  denke  man  sich  ein  veränderliches  Stück  einer  Mannig-
faltigkeit von  Einer  Dimension  ---  von  einem  festen  Anfangspunkte  an
gerechnet,  so  dass  die  Werthe  desselben  unter  einander  vergleichbar
sind  --- ,  welches  für  jeden  Punkt  der  gegebenen  Mannigfaltigkeit  einen
bestimmten  mit  ihm  stetig  sich  ändernden  Werth  hat,  oder  mit  andern
Worten,  man  nehme  innerhalb  der  gegebenen  Mannigfaltigkeit  eine
stetige  Function  des  Orts  an,  und  zwar  eine  solche  Function,  welche
nicht  längs  eines  Theils  dieser  Mannigfaltigkeit  constant  ist.  Jedes
System  von  Punkten,  wo  die  Function  einen  constanten  Werth  hat,
bildet  dann  eine  stetige  Mannigfaltigkeit  von  weniger  Dimensionen,
als  die  gegebene.  Diese  Mannigfaltigkeiten  gehen  bei  Aenderung  der
Function  stetig  in  einander  über;  man  wird  daher  annehmen  können,
dass  aus  einer  von  ihnen  die  übrigen  hervorgehen,  und  es  wird  dies,
allgemein  zu  reden,  so  geschehen  können,  dass  jeder  Punkt  in  einen
bestimmten    Punkt    der    andern    übergeht;    die    Ausnahmsfalle,    deren
\newpage

% --- Page 292 ---
276   XIII.    üeber  die  Hypothesen,  welche  der  Geometrie  zu  Grunde  liegen.
Untersuchung  wichtig  ist,  können  hier  unberücksichtigt  bleiben.  Hier-
durch wird  die  Ortsbestimmung  in  der  gegebeneu  Mannigfaltigkeit  zurück-
geführt auf  eine  Grössenbestimmung  und  auf  eine  Ortsbestimmung  in
einer  minderfach  ausgedehnten  Mannigfaltigkeit.  Es  ist  nun  leicht  zu
zeigen,  dass  diese  Mannigfaltigkeit  n  ---  1  Dimensionen  hat,  wenn  die
gegebene  Mannigfaltigkeit  eine  wfach  ausgedehnte  ist.  Durch  »^malige
Wiederholung  dieses  Verfahrens  wird  daher  die  Ortsbestimmung  in
einer  n  fach  ausgedehnten  Mannigfaltigkeit  auf  n  Grössenbestimmungen,
und  also  die  Ortsbestimmung  in  einer  gegebenen  Mannigfaltigkeit,
wenn  dieses  möglich  ist,  auf  eine  endliche  Anzahl  von  Quantitäts-
bestimmungen  zurückgeführt.  Es  giebt  indess  auch  Mannigfaltigkeiten,
in  welchen  die  Ortsbestimmung  nicht  eine  endliche  Zahl,  sondern  ent-
weder eine  unendliche  Reihe  oder  eine  stetige  Mannigfaltigkeit  von
Grössenbestimmungen  erfordert.  Solche  Mannigfaltigkeiten  bilden  z.  B.
die  möglichen  Bestimmungen  einer  Function  für  ein  gegebenes  Gebiet,
die  möglichen  Gestalten  ein  erräumlichen  Figur  u.  s.  w.
\section*{II.    Massverhältnisse,  deren  eine  Mannigfaltigkeit  von  n  Dimensionen}
fähig  ist,    unter    der   Voraussetzung,    dass    die    Linien    unabhängig
von  der  Lage  eine  Länge  besitzen,  also  jede  Linie  durch  jede
messbar  ist.
Es  folgt  nun,  nachdem  der  Begriff  einer  wfach  ausgedehnten  Mannig-
faltigkeit construirt  und  als  wesentliches  Kennzeichen  derselben  ge-
funden worden  ist,  dass  sich  die  Ortsbestimmung  in  derselben  auf
n  Grössenbestimmungen  zurückführen  lässt,  als  zweite  der  oben  ge-
stellten Aufgaben  eine  Untersuchung  über  die  Massverhältnisse,  deren
eine  solche  Mannigfaltigkeit  fähig  ist,  und  über  die  Bedingungeu,  welclie
zur  Bestimmung  dieser  Massverhältnisse  hinreichen.  Diese  Massver-
hältnisse lassen  sich  nur  in  abstracten  Grössenbegriffen  untersuchen
und  im  Zusammenhange  nur  durch  Formeln  darstellen;  unter  gewissen
Voraussetzungen  kann  man  sie  indess  in  Verhältnisse  zerlegen,  welche
einzeln  genommen  einer  geometrischen  Darstellung  fähig  sind,  und
hiedurch  wird  es  möglich,  die  Resultate  der  Rechnung  geometrisch
auszudrücken.  Es  wird  daher,  um  festen  Boden  zu  gewinnen,  zwar
eine  abstracte  Untersuchung  in  Formeln  nicht  zu  vermeiden  sein,  die
Resultate  derselben  aber  werden  sich  im  geometrischen  Gewände  dar-
stellen lassen.  Zu  Beidem  sind  die  Grundlagen  enthalten  in  der  be-
rühmten Abhandlung  des  Herrn  Geheimen  Hofraths  Gauss  über  die
krummen  Flächen.
\newpage

% --- Page 293 ---
XIII.    lieber  die  Hypothesen,  welche  der  Geometrie  zu  Grunde  «liegen.    277
1.
Massbestimmungen  erfordern  eine  Unabhängigkeit  der  Grossen
vom  Ort,  die  in  mehr  als  einer  Weise  stattfinden  kann:  die  zunächst
sich  darbietende  Annahme,  welche  ich  hier  verfolgen  will,  ist  wohl
die,  dass  die  Länge  der  Linien  unabhängig  von  der  Lage  sei,  also
jede  Linie  durch  jede  messbar  sei.  Wird  die  Ortsbestimmung  auf
Grössenbestimmungen  zurückgeführt,  also  die  Lage  eines  Punktes  in
der  gegebenen  «fach  ausgedehnten  Mannigfaltigkeit  durch  n  veränder-
liche Grössen  x^,  x.j.y  x-^,  und  so  fort  bis  Xn  ausgedrückt,  so  wird  die
Bestimmung  einer  Linie  darauf  hinauskommen,  dass  die  Grössen  x  als
Functionen  Einer  Veränderlichen  gegeben  werden.  Die  Aufgabe  ist
dann,  für  die  Länge  der  Linien  einen  mathematischen  Ausdruck  auf-
zustellen, zu  welchem  Zwecke  die  Grössen  x  als  in  Einheiten  ausdrück-
bar betrachtet  werden  müssen.  Ich  werde  diese  Aufgabe  nur  unter
gewissen  Beschränkimgen  behandeln  und  beschränke  mich  erstlich  auf
solche  Linien,  in  welchen  die  Verhältnisse  zwischen  den  Grössen  dx
---  den  zusammengehörigen  Aeuderungen  der  Grössen  x  ---  sich  stetig
ändern;  man  kann  dann  die  Linien  in  Elemente  zerlegt  denken,  inner-
halb deren  die  Verhältnisse  der  Grössen  dx  als  constant  betrachtet
werden  dürfen,  und  die  Aufgabe  kommt  dann  darauf  zurück,  für  jeden
Punkt  einen  allgemeinen  Ausdruck  des  von  ihm  ausgehenden  Linien-
elements ds  aufzustellen,  welcher  also  die  Grössen  x  und  die  Grössen
dx  enthalten  wird.  Ich  nehme  nun  zweitens  an,  dass  die  Länge  des
Linienelements,  von  Grössen  zweiter  Ordnung  abgesehen,  ungeändert
bleibt,  wenn  sämmtliche  Punkte  desselben  dieselbe  unendlich  kleine
Ortsänderung  erleiden,  worin  zugleich  enthalten  ist,  dass,  wenn  sämmt-
liche Grössen  dx  in  demselben  Verhältnisse  wachsen,  das  Linienelement
sich  ebenfalls  in  diesem  Verhältnisse  ändert.  Unter  diesen  Annahmen
wird  das  Linienelement  eine  beliebige  homogene  Function  ersten  Grades
der  Grössen  dx  sein  können,  welche  ungeändert  bleibt,  wenn  sämmt-
liche Grössen  dx  ihr  Zeichen  ändern,  und  worin  die  willkürlichen
Constanten  stetige  Functionen  der  Grössen  x  sind.  Um  die  einfachsten
Fälle  zu  finden,  suche  ich  zunächst  einen  Ausdruck  für  die  (w  ---  1)  fach
ausgedehnten  Mannigfaltigkeiten,  welche  vom  Anfangspunkte  des  Linien-
elements überall  gleich  weit  abstehen,  d.  h.  ich  suche  eine  stetige
Function  des  Orts,  welche  sie  von  einander  unterscheidet.  Diese  wird
vom  Anfangspunkt  aus  nach  allen  Seiten  entweder  ab-  oder  zunehmen
müssen;  ich  will  annehmen,  dass  sie  nach  allen  Seiten  zunimmt  und
also  in  dem  Punkte  ein  Minimum  hat.  Es  mus.5  dann,  wenn  ihre
ersten  und  zweiten  Differentialquotienten  endlich  sind,  das  Differential
\newpage

% --- Page 294 ---
278    XIII.   üeber  die  Hypothesen,  welche  der  Geometrie  zu  Grunde  liegen.
erster  Ordnung  verschwinden  und  das  zweiter  Ordnung  darf  nie  negativ
werden;  ich  nehme  an,  dass  es  immer  positiv  bleibt.  Dieser  Differential-
ausdruck zweiter  Ordnung  bleibt  alsdann  constant,  wenn  ds  constant
bleibt,  und  wächst  im  quadratischen  Verhältnisse,  wenn  die  Grössen
dx  und  also  auch  ds  sich  sämmtlich  in  demselben  Verhältnisse  ändern;
er  ist  also  =  const.  ds^  und  folglich  ist  ds  =  der  Quadratwurzel  aus
einer  immer  positiven  ganzen  homogenen  Function  zweiten  Grades  der
Grössen  dx,  in  welcher  die  Coefficienten  stetige  Functionen  der  Grössen
X  sind.  Für  den  Raum  wird,  wenn  man  die  Lage  der  Punkte  durch
rechtwinklige  Coordinaten  ausdrückt,  ds  =  YuJdxY]  der  Raum  ist
also  unter  diesem  einfachsten  Falle  enthalten.  Der  nächst  einfache
Fall  würde  wohl  die  Mannigfaltigkeiten  umfassen,  in  welchen  sich  das
Linienelement  durch  die  vierte  Wurzel  aus  einem  Differentialausdrucke
vierten  Grades  ausdrücken  lässt.  Die  Untersuchung  dieser  allgemeinern
Gattung  würde  zwar  keine  wesentlich  andere  Principien  erfordern,  aber
ziemlich  zeitraubend  sein  und  verhältnissmässig  auf  die  Lehre  vom
Räume  wenig  neues  Licht  werfen,  zumal  da  sich  die  Resultate  nicht
geometrisch  ausdrücken  lassen;  ich  beschränke  mich  daher  auf  die
Mannigfaltigkeiten,  wo  das  Linienelement  durch  die  Quadratwurzel  aus
einem  Differentialausdruck  zweiten  Grades  ausgedrückt  wird.  Man  kann
einen  solchen  Ausdruck  in  einen  andern  ähnlichen  transformiren,  in-
dem man  für  die  n  unabhängigen  Veränderlichen  Functionen  von  n
neuen  unabhängigen  Veränderlichen  setzt.  Auf  diesem  Wege  wird
man  aber  nicht  jeden  Ausdruck  in  jeden   transformiren  können;  denn
n-\textbackslash{}-l
der  Ausdruck  enthält  n-^-   Coefficienten,  welche  willkürliche  Functionen
der  unabhängigen  Veränderlichen  sind;  durch  Einführung  neuer  Ver-
änderlicher wird  man  aber  nur  «Relationen  genügen  und  also  nur  n
der  Coefficienten  gegebenen  Grössen  gleich   machen  können.     Es  sind
dann  die  übrigen  n  --- - ---  durch  die  Natur  der  darzustellenden  Mannig-
faltigkeit schon  völlig  bestimmt,  und  zur  Bestimmung  ihrer  Massver-
hältnisse also  n --- - ---  Functionen  des  Orts  erforderlich.  Die  Mannig-
faltigkeiten, in  welchen  sich,  wie  in  der  Ebene  und  im  Räume,  das
Linienelement  auf  die  Form  ^ Edx^  bringen  lässt,  bilden  daher  nur
einen  besondern  Fall  der  hier  zu  untersuchenden  Mannigfaltigkeiten;
sie  verdienen  wohl  einen  besonderen  Namen,  und  ich  will  also  diese
Mannigfaltigkeiten,  in  welchen  sich  das  Quadrat  des  Linienelements
auf  die  Summe  der  Quadrate  von  selbständigen  Differentialien  bringen
lässt,  eben  nennen.  Um  nun  die  wesentlichen  Verschiedenheiten  sämmt-
licher  in  der  vorausgesetzten  Form  darstellbarer  Mannigfaltigkeiten  über-
\newpage

% --- Page 295 ---
Xlir.   Ueber  die  Hypothesen,  welche  der  Geometrie  zu  Grande  liegen.    279
sehen  zu  können,  ist  es  nöthig,  die  von  der  Darstellungsweise  her-
rührenden zu  beseitigen ,  was  durch  Wahl  der  veränderlichen  Grössen
nach  einem  bestimmten  Princip  erreicht  wird.
2.
Zu  diesem  Ende  denke  man  sich  von  einem  beliebigen  Punkte  aus
das  System  der  von  ihm  ausgehenden  kürzesten  Linien  construirt;  die
Lage  eines  unbestimmten  Punktes  wird  dann  bestimmt  werden  können
durch  die  Anfangsrichtung  der  kürzesten  Linie,  in  welcher  er  liegt,
und  durch  seine  Entfernung  in  derselben  vom  Anfangspunkte  und  kann
daher  durch  die  Verhältnisse  der  Grössen  dx^,  d.  h.  der  Grössen  dx  im
Anfang  dieser  kürzesten  Linie  und  durch  die  Länge  s  dieser  Linie  aus-
gedrückt werden.  Man  führe  nun  statt  dx^  solche  aus  ihnen  gebildete
lineare  Ausdrücke  da  ein,  dass  der  Anfangswerth  des  Quadrats  des  Linien-
elements gleich  der  Summe  der  Quadrate  dieser  Ausdrücke  wird,  so  dass
die  unabhängigen  Variabein  sind:  die  Grösse  s  und  die  Verhältnisse  der
Grössen  da-^  und  setze  schliesslich  statt  da  solche  ihnen  proportionale
Grössen  x^^  x^,  . . .,  Xn,  dass  die  Quadratsumme  =  s^  wird.  Führt  man
diese  Grössen  ein,  so  wird  für  unendlich  kleine  Werthe  von  x  das
Quadrat  des  Linienelements  ==  2Jdx^,  das  Glied  der  nächsten  Ordnung  in
demselben  aber  gleich  einem  homogenen  Ausdruck  zweiten  Grades  der
n     -      Grössen   (x^  dx^  ---  x^  dx^) ,    (x^  dx^  ---  x.^  ^^\textbackslash{}))  \textbullet{}  \textbullet{}  \textbullet{} ;     also     eine
unendlich  kleine  Grösse  von  der  vierten  Dimension,  so  dass  man  eine
endliche  Grösse  erhält,  wenn  man  sie  durch  das  Quadrat  des  unendlich
kleinen  Dreiecks  dividirt,  in  dessen  Eckpunkten  die  Werthe  der  Ver-
änderlichen sind  (0,  0,  0,  .  . .),  {x^,  x^,  x.^,  . .  .),  (dx^,  dx.^,  dx^,  ....).
Diese  Grösse  behält  denselben  Werth,  so  lange  die  Grössen  x  und  dx
in  denselben  binären  Linearformen  enthalten  sind,  oder  so  lange  die
beiden  kürzesten  Linien  von  den  Werthen  0  bis  zu  den  Werthen  x
und  von  den  Werthen  0  bis  zu  den  Werthen  dx  in  demselben  Flächen-
element bleiben,  und  hängt  also  nur  von  Ort  und  Richtung  desselben
ab.  Sie  wird  offenbar  =  0,  wenn  die  dargestellte  Mannigfaltigkeit
eben,  d.  h.  das  Quadrat  des  Linienelements  auf  Udx^  reducirbar  ist,
und  kann  daher  als  das  Mass  der  in  diesem  Punkte  in  dieser  Flächen-
richtung stattfindenden  Abweichung  der  Mannigfaltigkeit  von  der  Eben-
heit angesehen  werden.  Multiplicirt  mit  ---  f  wird  sie  der  Grösse
gleich,  welche  Herr  Geheimer  Hofrath  Gauss  das  Krümmungsmass
einer  Fläche  genannt  hat.  Zur  Bestimmung  der  Massverhältnisse  einer
nfach  ausgedehnten  in  der  vorausgesetzten  Form  darstellbaren  Mannig-
faltigkeit   wurden   vorhin  n --- ^---  Functionen  des  Orts  nöthig  gefunden;
\newpage

% --- Page 296 ---
wenn  also  das  Krüramungsmass  in  jedem  Punkte  in  n  --- - ---  Flächen-
richtungen gegeben  wird,  so  werden  daraus  die  Massverhältnisse  der
Mannigfaltigkeit  sich  bestimmen  lassen,  wofern  nur  zwischen  diesen
Werthen  keine  identischen  Relationen  stattfinden,  was  in  der  That,
allgemein  zu  reden,  nicht  der  Fall  ist.  Die  Massverhältnisse  dieser
Mannigfaltigkeiten,  wo  das  Linienelemeut  durch  die  Quadratwurzel  aus
einem  Differentialausdruck  zweiten  Grades  dargestellt  wird,  lassen  sich
so  auf  eine  von  der  Wahl  der  veränderlichen  Grössen  völlig  unab-
hängige Weise  ausdrücken.  Ein  ganz  ähnlicher  Weg  lässt  sich  zu
diesem  Ziele  auch  bei  den  Mannigfaltigkeiten  einschlagen,  in  welchen
das  Linienelement  durch  einen  weniger  einfachen  Ausdruck,  z.  B.  durch
die  vierte  Wurzel  aus  einem  Differentialausdruck  vierten  Grades,  aus-
gedrückt wird.  Es  würde  sich  dann  das  Linienelement,  allgemein  zu
reden,  nicht  mehr  auf  die  Form  der  Quadratwurzel  aus  einer  Quadrat-
summe von  Differential  ausdrücken  bringen  lassen  und  also  in  dem
Ausdrucke  für  das  Quadrat  des  Linienelements  die  Abweichung  von
der  Ebenheit  eine  unendlich  kleine  Grösse  von  der  zweiten  Dimension
sein,  während  sie  bei  jenen  Mannigfaltigkeiten  eine  unendlich  kleine
Grösse  von  der  vierten  Dimension  war.  Diese  EigenthOmlichkeit  der
letztern  Mannigfaltigkeiten  kann  daher  wohl  Ebenheit  in  den  kleinsten
Theilen  genannt  werden.  Die  für  den  jetzigen  Zweck  wichtigste  Eigen-
thümlichkeit  dieser  Mannigfaltigkeiten,  derentwegen  sie  hier  allein
untersucht  worden  sind,  ist  aber  die,  dass  sich  die  Verhältnisse  der
zweifach  ausgedehnten  geometrisch  durch  Flächen  darstellen  und  die
der  mehrfach  a,usgedehnten  auf  die  der  in  ihnen  enthaltenen  Flächen
zurückführen  lassen,  was  jetzt  noch   einer    kurzen  Erörterung  bedarf.
3.
In  die  Auffassung  der  Flächen  mischt  sich  neben  den  inneren
Massverhältnissen,  bei  welchen  nur  die  Länge  der  Wege  in  ihnen  in
Betracht  kommt,  immer  auch  ihre  Lage  zu  ausser  ihnen  gelegenen
Funkten.  Man  kann  aber  von  den  äussern  Verhältnissen  abstrahiren,
indem  man  solche  Veränderungen  mit  ihnen  vornimmt,  bei  denen  die
Länge  der  Linien  in  ihnen  ungeändert  bleibt,  d.  h.  sie  sich  beliebig
---  ohne  Dehnung  ---  gebogen  denkt,  und  alle  so  auseinander  ent-
stehenden Flächen  als  gleichartig  betrachtet.  Es  gelten  also  z.  B.  be-
liebige cylindrische  oder  conische  Flächen  einer  Ebene  gleich,  weil  sie
sich  durch  blosse  Biegung  aus  ihr  bilden  lassen,  wobei  die  innern
Massverhältnisse  bleiben,  und  sämmtliche  Sätze  über  dieselben  ---  also
die  ganze  Planimetrie  ---  ihre  Gültigkeit  behalten;  dagegen  gelten  sie
\newpage

% --- Page 297 ---
XIII.    üeber  die  Hypothesen,  welche  der  Geometrie  zu  Giundc  liegen.    281
als  wesentlich  verschieden  von  der  Kugel,  welche  sich  nicht  ohne
Dehnung  in  eine  Ebene  verwandeln  lässt.  Nach  der  vorigen  Unter-
suchung werden  in  jedem  Punkte  die  innern  Massverhältnisse  einer
zweifach  ausgedehnten  Grösse,  wenn  sich  das  Linienelement  durch  die
Quadratwurzel  aus  einem  Differentialausdruck  zweiten  Grades  ausdrücken
lässt,  wie  dies  bei  den  Flächen  der  Fall  ist,  charakterisirt  durch  das
Krümmungsmass.  Dieser  Grösse  lässt  sich  nun  bei  den  Flächen  die
anschauliche  Bedeutung  geben,  dass  sie  das  Product  aus  den  beiden
Krümmungen  der  Fläche  in  diesem  Punkte  ist,  oder  auch,  dass  das
Product  derselben  in  ein  unendlich  kleines  aus  kürzesten  Linien  ge-
bildetes Dreieck  gleich  ist  dem  halben  üeberschusse  seiner  Winkel-
summe über  zwei  Rechte  in  Theilen  des  Halbmessers.  Die  erste  De-
finition würde  den  Satz  voraussetzen,  dass  das  Product  der  beiden
Krümmungshalbmesser  bei  der  blossen  Biegung  einer  Fläche  ungeändert
bleibt,  die  zweite,  dass  an  demselben  Orte  der  Ueberschuss  der  Winkel-
summe eines  unendlich  kleinen  Dreiecks  über  zwei  Rechte  seinem
Inhalte  proportional  ist.  Um  dem  Krümmungsmass  einer  wfach  aus-
gedehnten Mannigfaltigkeit  in  einem  gegebenen  Punkte  und  einer  ge-
gebenen durch  ihn  gelegten  Flächenrichtung  eine  greifbare  Bedeutung
zu  geben,  muss  man  davon  ausgehen,  dass  eine  von  einem  Punkte
ausgehende  kürzeste  Linie  völlig  bestimmt  ist,  wenn  ihre  Anfangs-
richtung gegeben  ist.  Hienach  wird  man  eine  bestimmte  Fläche  er-
halten, wenn  man  sämmtliche  von  dem  gegebenen  Punkte  ausgehenden
und  in  dem  gegebenen  Flächenelement  liegenden  Anfangsrichtungen
zu  kürzesten  Linien  verlängert,  und  diese  Fläche  hat  in  dem  gegebenen
Punkte  ein  bestimmtes  Krümmungsmass,  welches  zugleich  das  Krüm-
mungsmass der  wfach  ausgedehnten  Mannigfaltigkeit  in  dem  gegebenen
Punkte  und  der  gegebenen  Flächenrichtung  ist.
Es  sind  nun  noch,  ehe  die  Anwendung  auf  den  Raum  gemacht
wird,  einige  Betrachtungen  über  die  ebenen  Mannigfaltigkeiten  im  All-
gemeinen nöthig,  d.  h.  über  diejenigen,  in  welchen  das  Quadrat  des
Liuienelements  durch  eine  Quadratsumme  vollständiger  Differentialien
darstellbar  ist.
In  einer  ebenen  nfach  ausgedehnten  Mannigfaltigkeit  ist  das
Krümmungsmass  in  jedem  Punkte  in  jeder  Richtung  Null;  es  reicht
aber  nach  der  frühern  Untersuchung,  um   die  Massverhältnisse  zu  be-
stimmen,  hin  zu  wissen,   dass  es  in  jedem  Punkte  in  n --- ^ ---  Flächen-
richtungen,   deren    Krümmungsmasse    von    einander    unabhängig   sind,
\newpage

% --- Page 298 ---
Null  sei.  Die  Mannigfaltigkeiten,  deren  Krümmungsmass  überall  ==  0
ist,  lassen  sich  betrachten  als  ein  besonderer  Fall  derjenigen  Mannig-
faltigkeiten, deren  Krümmungsmass  allenthalben  constant  ist.  Der
gemeinsame  Charakter  dieser  Mannigfaltigkeiten,  deren  Krümmungs-
mass constant  ist,  kann  auch  so  ausgedrückt  werden,  dass  sich  die
Figuren  in  ihnen  ohne  Dehnung  bewegen  lassen.  Denn  offenbar  wür-
den die  Figuren  in  ihnen  nicht  beliebig  verschiebbar  und  drehbar  sein
können,  wenn  nicht  in  jedem  Punkte  in  allen  Richtungen  das  Krüm-
mungsmass dasselbe  wäre.  Andererseits  aber  sind  durch  das  Krüm-
mungsmass die  Massverhältnisse  der  Mannigfaltigkeit  vollständig  be-
stimmt; es  sind  daher  um  einen  Punkt  nach  allen  Richtungen  die
Massverhältnisse  genau  dieselben,  wie  um  einen  andern,  und  also  von
ihm  aus  dieselben  Constructionen  ausführbar,  und  folglich  kann  in  den
Mannigfaltigkeiten  mit  constantem  Krümmungsmass  den  Figuren  jede
beliebige  Lage  gegeben  werden.  Die  Massverhältnisse  dieser  Mannig-
faltigkeiten hängen  nur  von  dem  Werthe  des  Krümmungsmasses  ab,
und  in  Bezug  auf  die  analytische  Darstellung  mag  bemerkt  werden,
dass,  wenn  man  diesen  Werth  durch  a  bezeichnet,  dem  Ausdruck  für
das  Linienelement  die  Form
$\blacksquare$
ysdx^
gegeben  werden  kann.
5.
Zur  geometrischen  Erläuterung  kann  die  Betrachtung  der  Flächen
mit  constantem  Krümmungsmass  dienen.  Es  ist  leicht  zu  sehen,  dass
sich  die  Flächen,  deren  Krümmungsmass  positiv  ist,  immer  auf  eine
Kugel,  deren  Radius  gleich  1  dividirt  durch  die  Wurzel  aus  dem
Krümmungsmass  ist,  wickeln  lassen  werden;  um  aber  die  ganze  Mannig-
faltigkeit dieser  Flächen  zu  übersehen,  gebe  man  einer  derselben  die
Gestalt  einer  Kugel  und  den  übrigen  die  Gestalt  von  Umdrehungs-
flächen, welche  sie  im  Aequator  berühren.  Die  Flächen  mit  grösserem
Krümmungsmass,  als  diese  Kugel,  werden  dann  die  Kugel  von  innen
berühren  und  eine  Gestalt  annehmen,  wie  der  äussere  der  Axe  ab-
gewandte Theil  der  Oberfläche  eines  Ringes;  sie  würden  sich  auf  Zonen
von  Kugeln  mit  kleinerem  Halbmesser  wickeln  lassen,  aber  mehr  als
einmal  herumreichen.  Die  Flächen  mit  kleinerem  positiven  Krümmungs-
mass wird  man  erhalten,  wenn  man  aus  Kugelflächen  mit  grösserem
Radius  ein  von  zwei  grössten  Halbkreisen  begrenztes  Stück  ausschneidet
und  die  Schnittlinien  zusammenfügt.    Die  Fläche  mit  dem  Krümmungs-
\newpage

% --- Page 299 ---
XI 11.    Ucbei-  die  Hypothesen,  welche  der  Geometrie  zu  Grunde  liegen.    28.3
mass  Null  wird  eine  auf  dem  Aequator  stehende  Cylinderfiäche  sein;
die  Flächen  mit  negativem  Krümmungsmass  aber  werden  diesen  Cylin-
der  von  aussen  berühren  und  wie  der  innere  der  Axe  zugewandte  Theil
der  Oberfläche  eines  Ringes  geformt  sein.  Denkt  man  sich  diese
Flächen  als  Ort  für  in  ihnen  bewegliche  Flächenstücke,  wie  den  Raum
als  Ort  für  Körper,  so  sind  in  allen  diesen  Flächen  die  Flächenstückc
ohne  Dehnung  beweglich.  Die  Flächen  mit  positivem  Krümmungsmass
lassen  sich  stets  so  formen,  dass  die  Flächenstücke  auch  ohne  Biegung
beliebig  bewegt  werden  können,  nämlich  zu  Kugelflächen,  die  mit  ne-
gativem aber  nicht.  Ausser  dieser  Unabhängigkeit  der  Flächenstücke
vom  Ort  findet  bei  der  Fläche  mit  dem  Krümmungsmass  Null  auch
eine  Unabhängigkeit  der  Richtung  vom  Ort  statt,  welche  bei  den
übrigen  Flächen  nicht  stattfindet.
\section*{III.    Anwendung  auf  den  Baum.}
1.
Nach  diesen  Untersuchungen  über  die  Bestimmung  der  Massver-
hältnisse einer  nfach  ausgedehnten  Grösse  lassen  sich  nun  die  Bedin-
gungen angeben,  welche  zur  Bestimmung  der  Massverhältnisse  des
Raumes  hinreichend  und  nothwendig  sind,  wenn  Unabhängigkeit  der
Linien  von  der  Lage  und  Darstellbarkeit  des  Linienelements  durch  die
Quadratwurzel  aus  einem  [Differentialausdrucke  zweiten  Grades,  also
Ebenheit  in  den  kleinsten  Theilen  vorausgesetzt  wird.
Sie  lassen  sich  erstens  so  ausdrücken,  dass  das  Krümmungsmass
in  jedem  Punkte  in  drei  Flächenrichtungen  =  0  ist,  und  es  sind  da-
her die  Massverhältnisse  des  Raumes  bestimmt,  wenn  die  Winkel-
summe im  Dreieck  allenthalben  gleich  zwei  Rechten  ist.
Setzt  man  aber  zweitens,  wie  Euklid,  nicht  bloss  eine  von  der
Lage  unabhängige  Existenz  der  Linien,  sondern  auch  der  Körper
voraus,  so  folgt,  dass  das  Krümmungsmass  allenthalben  constant  ist,
und  es  ist  dann  in  allen  Dreiecken  die  Winkelsumme  bestimmt,  wenn
sie  in  Einem  bestimmt  ist.
Endlich  könnte  man  drittens,  anstatt  die  Länge  der  Linien  als
unabhängig  von  Ort  und  Richtung  anzunehmen,  auch  eine  Unab-
hängigkeit ihrer  Länge  und  Richtung  vom  Ort  voraussetzen.  Nach
dieser  Auffassung  sind  die  Ortsänderungen  oder  Ortsverschiedenheiten
complexe  in  drei  unabhängige  Einheiten  ausdrückbare  Grössen.
2.
Im  Laufe  der  bisherigen  Betrachtungen  wurden  zunächst  die  Aus-
dehnungs-    oder  Gebietsverhältnisse  von  den  Massverhältuissen    geson-
\newpage

% --- Page 300 ---
284    XIII.    üeber  die  Hypothesen,  welche  der  Geometrie  zu  Grunde  liegen.
dert,  und  gefunden,  dass  bei  denselben  Ausdehnungsverhältnissen  ver-
schiedene Massverhältnisse  denkbar  sind;  es  vrurden  dann  die  Systeme
einfacher  Massbestimmungen  aufgesucht,  durch  welche  die  Massver-
hältnisse des  Raumes  völlig  bestimmt  sind  und  von  welchen  alle  Sätze
über  dieselben  eine  nothwendige  Folge  sind;  es  bleibt  nun  die  Frage
zu  erörtern,  wie,  in  welchem  Grade  und  in  welchem  Umfange  diese
Voraussetzungen  durch  die  Erfahrung  verbürgt  werden.  In  dieser  Be-
ziehung findet  zwischen  den  blossen  Ausdehnungsverhältnissen  und  den
Massverhältnissen  eine  wesentliche  Verschiedenheit  statt,  insofern  bei
erstem,  wo  die  möglichen  Fälle  eine  discrete  Mannigfaltigkeit  bilden,
die  Aussagen  der  Erfahrung  zwar  nie  völlig  gewiss,  aber  nicht  un-
genau sind,  während  bei  letztern,  wo  die  möglichen  Fälle  eine  stetige
Mannigfaltigkeit  bilden,  jede  Bestimmung  aus  der  Erfahrung  immer
ungenau  bleibt  ---  es  mag  die  Wahrscheinlichkeit,  dass  sie  nahe  richtig
ist,  noch  so  gross  sein.  Dieser  Umstand  wird  wichtig  bei  der  Aus-
dehnung dieser  empirischen  Bestimmungen  über  die  Grenzen  der  Beob-
achtung in's  Unmessbargrosse  und  Unmessbarkleine;  denn  die  letztern
können  offenbar  jenseits  der  Grenzen  der  Beobachtung  immer  unge-
nauer werden,  die  ersteren  aber  nicht.
Bei  der  Ausdehnung  der  Raumconstructionen  ins  Unmessbargrosse
ist  Unbegrenztheit  und  Unendlichkeit  zu  scheiden;  jene  gehört  zu  den
Ausdehnungsverhältnissen,  diese  zu  den  Massverhältnissen.  Dass  der
Raum  eine  unbegrenzte  dreifach  ausgedehnte  Mannigfaltigkeit  sei,  ist
eine  Voraussetzung,  welche  bei  jeder  Auffassung  der  Aussenwelt  an-
gewandt wird,  nach  welcher  in  jedem  Augenblicke  das  Gebiet  der
wirklichen  Wahrnehmungen  ergänzt  und  die  möglichen  Orte  eines  ge-
suchten Gegenstandes  construirt  werden  und  welche  sich  bei  diesen
Anwendungen  fortwährend  bestätigt.  Die  Unbegrenztheit  des  Raumes
besitzt  daher  eine  grössere  empirische  Gewissheit',  als  irgend  eine
äussere  Erfahrung.  Hieraus  folgt  aber  die  Unendlichkeit  keineswegs;
vielmehr  würde  der  Raum,  wenn  man  Unabhängigkeit  der  Körper  vom
Ort  voraussetzt,  ihm  also  ein  constantes  Krümmungsmass  zuschreibt,
nothwendig  endlich  sein,  so  bald  dieses  Krümmungsmass  einen  noch
so  kleinen  positiven  Werth  hätte.  Man  würde,  wenn  man  die  in  einem
Flächenelement  liegenden  Anfangsrichtungen  zu  kürzesten  Linien  ver-
längert, eine  unbegrenzte  Fläche  mit  constantem  positiven  Krümmungs-
mass, also  eine  Fläche  erhalten,  welche  in  einer  ebenen  dreifach  aus-
gedehnten Mannigfaltigkeit  die  Gestalt  einer  Kugelfläche  annehmen
würde  und  welche  folglich  endlich  ist.
\newpage

\clearpage
\chapter*{Source segment 7: riemann pages 301 350}
\addcontentsline{toc}{chapter}{Source segment 7: riemann pages 301 350}
\begin{center}
\Large\bfseries Riemann --- Gesammelte Werke (pp.~301--350)
\end{center}

\tableofcontents
\newpage

%==============================================================================
% XIII. End of "Ueber die Hypothesen, welche der Geometrie zu Grunde liegen"
%==============================================================================
\section*{XIII. Ueber die Hypothesen, welche der Geometrie zu Grunde liegen.}
\addcontentsline{toc}{section}{XIII. Ueber die Hypothesen, welche der Geometrie zu Grunde liegen.}

\subsection*{3.}
\addcontentsline{toc}{subsection}{\S~3. In wie weit im Unendlichkleinen?}

\selectlanguage{german}

Die Fragen ueber das Unmessbargrosse sind fuer die Naturerklaerung muessige Fragen. Anders verhaelt es sich aber mit den Fragen ueber das Unmessbarkleine. Auf der Genauigkeit, mit welcher wir die Erscheinungen in's Unendlichkleine verfolgen, beruht wesentlich die Erkenntniss ihres Causalzusammenhangs. Die Fortschritte der letzten Jahrhunderte in der Erkenntniss der mechanischen Natur sind fast allein bedingt durch die Genauigkeit der Construction, welche durch die Erfindung der Analysis des Unendlichen und die von Archimed, Galil\"ai und Newton aufgefundenen einfachen Grundbegriffe, deren sich die heutige Physik bedient, m\"oglich geworden ist. In den Naturwissenschaften aber, wo die einfachen Grundbegriffe zu solchen Constructionen bis jetzt fehlen, verfolgt man, um den Causalzusammenhang zu erkennen, die Erscheinungen in's r\"aumlich Kleine, so weit es das Mikroskop nur gestattet. Die Fragen ueber die Massverh\"altnisse des Raumes im Unmessbarkleinen geh\"oren also nicht zu den muessigen.

Setzt man voraus, dass die K\"orper unabh\"angig vom Ort existiren, so ist das Kr\"ummungsmass ueberall constant, und es folgt dann aus den astronomischen Messungen, dass es nicht von Null verschieden sein kann; jedenfalls m\"usste sein reciprocer Werth eine Fl\"ache sein, gegen welche das unsern Teleskopen zug\"angliche Gebiet verschwinden m\"usste. Wenn aber eine solche Unabh\"angigkeit der K\"orper vom Ort nicht stattfindet, so kann man aus den Massverh\"altnissen im Grossen nicht auf die im Unendlichkleinen schliessen; es kann dann in jedem Punkte das Kr\"ummungsmass in drei Richtungen einen beliebigen Werth haben, wenn nur die ganze Kr\"ummung jedes messbaren Raumtheils nicht merklich von Null verschieden ist; noch complicirtere Verh\"altnisse k\"onnen eintreten, wenn die vorausgesetzte Darstellbarkeit eines Linienelements durch die Quadratwurzel aus einem Differentialausdruck zweiten Grades nicht stattfindet. Nun scheinen aber die empirischen Begriffe, in welchen die r\"aumlichen Massbestimmungen gegr\"undet sind, der Begriff des festen K\"orpers und des Lichtstrahls, im Unendlichkleinen ihre G\"ultigkeit zu verlieren; es ist also sehr wohl denkbar, dass die Massverh\"altnisse des Raumes im Unendlichkleinen den Voraussetzungen der Geometrie nicht gem\"ass sind, und dies w\"urde man in der That annehmen m\"ussen, sobald sich dadurch die Erscheinungen auf einfachere Weise erkl\"aren liessen.

Die Frage ueber die G\"ultigkeit der Voraussetzungen der Geometrie im Unendlichkleinen h\"angt zusammen mit der Frage nach dem Innern Grunde der Massverh\"altnisse des Raumes. Bei

\newpage
\noindent wohl noch zur Lehre vom R\"aume gerechnet werden darf, kommt die obige Bemerkung zur Anwendung, dass bei einer discreten Mannigfaltigkeit das Princip der Massverh\"altnisse schon in dem Begriffe dieser Mannigfaltigkeit enthalten ist, bei einer stetigen aber anders woher hinzukommen muss. Es muss also entweder das dem R\"aume zu Grunde liegende Wirkliche eine discrete Mannigfaltigkeit bilden, oder der Grund der Massverh\"altnisse ausserhalb, in darauf wirkenden bindenden Kr\"aften, gesucht werden.

Die Entscheidung dieser Fragen kann nur gefunden werden, indem man von der bisherigen durch die Erfahrung bew\"ahrten Auffassung der Erscheinungen, wozu Newton den Grund gelegt, ausgeht und diese durch Thatsachen, die sich aus ihr nicht erkl\"aren lassen, getrieben allm\"ahlich umarbeitet; solche Untersuchungen, welche, wie die hier gef\"uhrte, von allgemeinen Begriffen ausgehen, k\"onnen nur dazu dienen, dass diese Arbeit nicht durch die Beschr\"anktheit der Begriffe gehindert und der Fortschritt im Erkennen des Zusammenhangs der Dinge nicht durch ueberlieferte Vorurtheile gehemmt wird.

Es f\"uhrt dies hin\"uber in das Gebiet einer andern Wissenschaft, in das Gebiet der Physik, welches wohl die Natur der heutigen Veranlassung nicht zu betreten erlaubt.

\medskip
\noindent\textbf{"Ubersicht.}

\begin{center}
\begin{tabular}{p{10cm}r}
\multicolumn{2}{r}{Seite} \\[2mm]
Plan der Untersuchung & 272 \\[1mm]
I. Begriff einer $n$fach ausgedehnten Gr\"osse\textsuperscript{*}) & 273 \\[1mm]
\quad \S.~1. Stetige und discrete Mannigfaltigkeiten. Bestimmte Theile einer Mannigfaltigkeit heissen Quanta. Eintheilung der Lehre von den stetigen Gr\"ossen in die Lehre & \\[1mm]
\quad 1) von den blossen Gebietsverh\"altnissen, bei welcher eine Unabh\"angigkeit der Gr\"ossen vom Ort nicht vorausgesetzt wird, & \\[1mm]
\quad 2) von den Massverh\"altnissen, bei welcher eine solche Unabh\"angigkeit vorausgesetzt werden muss & 273 \\[2mm]
\quad \S.~2. Erzeugung des Begriffs einer einfach, zweifach, \dots, $n$fach ausgedehnten Mannigfaltigkeit & 275 \\[1mm]
\quad \S.~3. Zur\"uckf\"uhrung der Ortsbestimmung in einer gegebenen Mannigfaltigkeit auf Quantit\"atsbestimmungen. Wesentliches Kennzeichen einer $n$fach ausgedehnten Mannigfaltigkeit & 275 \\[2mm]
II. Massverh\"altnisse, deren eine Mannigfaltigkeit von $n$ Dimensionen f\"ahig
\end{tabular}
\end{center}

\newpage
\begin{center}
\begin{tabular}{p{10cm}r}
\noindent ist\textsuperscript{*}), unter der Voraussetzung, dass die Linien unabh\"angig von der Lage eine L\"ange besitzen, also jede Linie durch jede messbar ist & 276 \\[2mm]
\quad \S.~1. Ausdruck des Linienelements. Als eben werden solche Mannigfaltigkeiten betrachtet, in denen das Linienelement durch die Wurzel aus einer Quadratsumme vollst\"andiger Differentialien ausdr\"uckbar ist & 277 \\[2mm]
\quad \S.~2. Untersuchung der $n$fach ausgedehnten Mannigfaltigkeiten, in welchen das Linienelement durch die Quadratwurzel aus einem Differentialausdruck zweiten Grades dargestellt werden kann. Mass ihrer Abweichung von der Ebenheit (Kr\"ummungsmass) in einem gegebenen Punkte und einer gegebenen Fl\"achenrichtung. Zur Bestimmung ihrer Massverh\"altnisse ist es (unter gewissen Beschr\"ankungen) zul\"assig und hinreichend, dass das Kr\"ummungsmass in jedem Punkte in $\frac{1}{2}n(n-1)$ Fl\"achenrichtungen beliebig gegeben wird & 279 \\[2mm]
\quad \S.~3. Geometrische Erl\"auterung & 280 \\[1mm]
\quad \S.~4. Die ebenen Mannigfaltigkeiten (in denen das Kr\"ummungsmass allenthalben $=0$ ist) lassen sich betrachten als einen besondern Fall der Mannigfaltigkeiten mit constantem Kr\"ummungsmass. Diese k\"onnen auch dadurch definirt werden, dass in ihnen Unabh\"angigkeit der $n$fach ausgedehnten Gr\"ossen vom Ort (Bewegbarkeit derselben ohne Dehnung) stattfindet & 281 \\[2mm]
\quad \S.~5. Fl\"achen mit constantem Kr\"ummungsmasse & 282 \\[2mm]
III. Anwendung auf den Raum & 283 \\[2mm]
\quad \S.~1. Systeme von Thatsachen, welche zur Bestimmung der Massverh\"altnisse des Raumes, wie die Geometrie sie voraussetzt, hinreichen & 283 \\[1mm]
\quad \S.~2. In wie weit ist die G\"ultigkeit dieser empirischen Bestimmungen wahrscheinlich jenseits der Grenzen der Beobachtung im Unmessbargrossen? & 283 \\[1mm]
\quad \S.~3. In wie weit im Unendlichkleinen? Zusammenhang dieser Frage mit der Naturerklaerung\textsuperscript{**}) & 286
\end{tabular}
\end{center}

\bigskip
\noindent\small\textsuperscript{*}) Art.~I. bildet zugleich die Vorarbeit f\"ur Beitr\"age zur analysis situs.

\small\textsuperscript{**}) Der \S.~3 des Art.~III bedarf noch einer Umarbeitung und weiteren Ausf\"uhrung. (Bemerkungen von Riemann.)

\newpage
%==============================================================================
% XIV. Ein Beitrag zur Elektrodynamik
%==============================================================================
\section*{XIV. Ein Beitrag zur Elektrodynamik.}
\addcontentsline{toc}{section}{XIV. Ein Beitrag zur Elektrodynamik.}

\noindent\small{(Aus Poggendorff's Annalen der Physik und Chemie, Bd.~CXXXI.)}

\medskip

\normalsize
Der K\"oniglichen Societ\"at erlaube ich mir eine Bemerkung mitzutheilen, welche die Theorie der Elektricit\"at und des Magnetismus mit der des Lichts und der strahlenden W\"arme in einen nahen Zusammenhang bringt. Ich habe gefunden, dass die elektrodynamischen Wirkungen galvanischer Str\"ome sich erkl\"aren lassen, wenn man annimmt, dass die Wirkung einer elektrischen Masse auf die uebrigen nicht momentan geschieht, sondern sich mit einer constanten (der Lichtgeschwindigkeit innerhalb der Grenzen der Beobachtungsfehler gleichen) Geschwindigkeit zu ihnen fortpflanzt. Die Differentialgleichung f\"ur die Fortpflanzung der elektrischen Kraft wird bei dieser Annahme dieselbe, wie die f\"ur die Fortpflanzung des Lichts und der strahlenden W\"arme.

Es seien $S$ und $S'$ zwei von constanten galvanischen Str\"omen durchflossene und gegen einander nicht bewegte Leiter, $s$ sei ein elektrisches Massentheilchen im Leiter $S$, welches sich zur Zeit $t$ im Punkte $(x,y,z)$ befinde, $e$ ein elektrisches Massentheilchen von $S'$ und befinde sich zur Zeit $t$ im Punkte $(x',y',z')$. Ueber die Bewegung der elektrischen Massentheilchen, welche in jedem Leitertheilchen f\"ur die positiv und negativ elektrischen entgegengesetzt ist, mache ich die Voraussetzung, dass sie in jedem Augenblicke so vertheilt sind, dass die Summen
\[\sum s f(x,y,z), \quad \sum s' f(x',y',z'),\]
\"uber s\"ammtliche Massentheilchen der Leiter ausgedehnt gegen dieselben Summen, wenn sie nur ueber die positiv elektrischen oder nur ueber die negativ elektrischen Massentheilchen ausgedehnt werden, vernachl\"assigt werden d\"urfen, sobald die Function $f$ und ihre Differentialquotienten stetig sind.

Diese Voraussetzung kann auf sehr mannigfaltige Weise erf\"ullt werden. Nimmt man z.~B. an, dass die Leiter in den kleinsten Theilen krystallinisch sind, so dass sich dieselbe relative Vertheilung der

\newpage
\noindent Elektricit\"aten in bestimmten gegen die Dimensionen der Leiter unendlich kleinen Abst\"anden periodisch wiederholt, so sind, wenn $\beta$ die L\"ange einer solchen Periode bezeichnet, jene Summen unendlich klein, wie $\beta^{n}$, wenn $f$ und ihre Derivirten bis zur $(n-1)$ten Ordnung stetig sind, und unendlich klein wie $e^{-\frac{c}{\beta}}$, wenn sie s\"ammtlich stetig sind.

\medskip
\noindent\textbf{Erfahrungsm\"assiges Gesetz der elektrodynamischen Wirkungen.}

Sind die specifischen Stromintensit\"aten nach mechanischem Mass zur Zeit $t$ im Punkte $(x,y,z)$ parallel den drei Axen $u,v,w$, und im Punkte $(x',y',z')$ $u',v',w'$, und bezeichnet $r$ die Entfernung beider Punkte, $c$ die von Kohlrausch und Weber bestimmte Constante, so ist der Erfahrung nach das Potential der von $S$ auf $S'$ ausge\"ubten Kr\"afte
\[\frac{2}{cc} \int\!\int \frac{uu' + vv' + ww'}{r}\, dS\, dS',\]
dieses Integral ueber s\"ammtliche Elemente $dS$ und $dS'$ der Leiter $S$ und $S'$ ausgedehnt. F\"uhrt man statt der specifischen Stromintensit\"aten die Producte aus den Geschwindigkeiten in die specifischen Dichtigkeiten und dann f\"ur die Producte aus diesen in die Volumelemente die in ihnen enthaltenen Massen ein, so geht dieser Ausdruck ueber in
\[\frac{1}{cc} \sum\sum \frac{ee'}{r} \frac{d d'(r^{2})}{dt\,dt'},\]
wenn die Aenderung von $r^{2}$ w\"ahrend der Zeit $dt$, welche von der Bewegung von $e$ herr\"uhrt, durch $d$, und die von der Bewegung von $e'$ herr\"uhrende durch $d'$ bezeichnet wird.

Dieser Ausdruck kann durch Hinwegnahme von
\[\sum e' \sum \frac{s}{cc} \frac{d'(r^{2})}{r} \frac{1}{dt}\, dt',\]
welches durch die Summirung nach $e$ verschwindet, in
\[\frac{1}{cc} \sum\sum \frac{e'e}{r} \frac{d'(r^{2})}{dt}\, dt\]
und dieses wieder durch Addition von
\[\sum e \sum \frac{e'}{cc} \frac{d'(r^{2})}{r} \frac{1}{dt}\, dt',\]
welches durch die Summation nach $e'$ Null wird, in
\[\frac{1}{cc} \sum\sum ee' \frac{dd'(r^{2})}{dt\, dt'}\]
verwandelt werden.

\medskip
\noindent\textbf{Ableitung dieses Gesetzes aus der neuen Theorie.}

Nach der bisherigen Annahme ueber die elektrostatische Wirkung wird die Potentialfunction $U$ beliebig vertheilter elektrischer Massen, wenn $\varrho$ ihre Dichtigkeit im Punkte $(x,y,z)$ bezeichnet, durch die Bedingung
\[\frac{\partial^{2}U}{\partial x^{2}} + \frac{\partial^{2}U}{\partial y^{2}} + \frac{\partial^{2}U}{\partial z^{2}} + 4\pi\varrho = 0\]
und durch die Bedingung, dass $U$ stetig und in unendlicher Entfernung von wirkenden Massen constant sei, bestimmt. Ein particul\"ares Integral der Gleichung
\[\frac{\partial^{2}U}{\partial x^{2}} + \frac{\partial^{2}U}{\partial y^{2}} + \frac{\partial^{2}U}{\partial z^{2}} = 0\]
welches ueberall ausser dem Punkte $(x',y',z')$ stetig bleibt, ist
\[\frac{1}{r},\]
und diese Function bildet die vom Punkte $(x',y',z')$ aus erzeugte Potentialfunction, wenn sich in demselben zur Zeit $t$ die Masse $-f(t)$ befindet.

Statt dessen nehme ich nun an, dass die Potentialfunction $U$ durch die Bedingung
\[\frac{\partial^{2}U}{\partial x^{2}} + \frac{\partial^{2}U}{\partial y^{2}} + \frac{\partial^{2}U}{\partial z^{2}} - aa \frac{\partial^{2}U}{\partial t^{2}} + 4\pi\varrho = 0\]
bestimmt wird, so dass die vom Punkte $(x',y',z')$ aus erzeugte Potentialfunction, wenn sich in demselben zur Zeit $t$ die Masse $-f(t)$ befindet,
\[\frac{f\left(t - \frac{r}{a}\right)}{r}\]
wird.

Bezeichnet man die Coordinaten der Masse $e$ zur Zeit $t$ durch $x_{t}$, $y_{t}$, $z_{t}$, und die der Masse $e'$ zur Zeit $t'$ durch $x'_{t'}$, $y'_{t'}$, $z'_{t'}$, und setzt zur Abk\"urzung
\[\frac{1}{a}\sqrt{(x_{t} - x'_{t'})^{2} + (y_{t} - y'_{t'})^{2} + (z_{t} - z'_{t'})^{2}} = \frac{r}{a} = F(t,t'),\]
so wird nach dieser Annahme das Potential von $e$ auf $e'$ zur Zeit $t$
\[= -ee' \frac{f'(t - aF(t,t'))}{r}.\]

\newpage
Das Potential der von s\"ammtlichen Massen $e$ des Leiters $S$ auf die Massen $e'$ des Leiters $S'$ von der Zeit $0$ bis zur Zeit $t$ ausge\"ubten Kr\"afte wird daher
\[P = - \int_{0}^{t} \sum\sum ee' F'(t-a, t')\, dt',\]
die Summen ueber s\"ammtliche Massen beider Leiter ausgedehnt.

Da die Bewegung f\"ur entgegengesetzt elektrische Massen in jedem Leitertheilchen entgegengesetzt ist, so erlangt die Function $F(t,t')$ durch die Derivation nach $t$ die Eigenschaft, mit $e$, und durch die Derivation nach $t'$ die Eigenschaft, mit $e'$ ihr Zeichen zu \"andern. Bei der vorausgesetzten Vertheilung der Elektricit\"aten wird daher, wenn man die Derivationen nach $t$ durch obere und nach $t'$ durch untere Accente bezeichnet, $\sum\sum ee' F'_{n}(t,t)$, ueber s\"ammtliche elektrische Massen ausgedehnt, nur dann nicht unendlich klein gegen die ueber die elektrischen Massen einer Art erstreckte Summe, wenn $n$ und $\nu$ beide ungerade sind.

Man nehme nun an, dass die elektrischen Massen w\"ahrend der Fortpflanzungszeit der Kraft von einem Leiter zum anderen nur einen sehr kleinen Weg zur\"ucklegen, und betrachte die Wirkung w\"ahrend eines Zeitraums, gegen welchen die Fortpflanzungszeit verschwindet. In dem Ausdrucke von $P$ kann man dann zun\"achst
\[F'(t-a, t') - F'(t, t')\]
durch
\[-a \frac{\partial F'(t,t')}{\partial t}\]
ersetzen, da $\sum\sum ee' F'(t,t)$ vernachl\"assigt werden darf. Man erh\"alt dadurch
\[P = \int_{0}^{t} dt' \sum\sum ee' a \frac{\partial F'(t,t')}{\partial t},\]
oder wenn man die Ordnung der Integrationen umkehrt und $t+a$ f\"ur $\tau$ setzt,
\[P = \sum\sum ee' \int_{0}^{t} d\vartheta \int_{-\vartheta}^{t} d\tau\, F'(\tau, \tau+\vartheta).\]
Verwandelt man die Grenzen des innern Integrals in $0$ und $t$, so wird dadurch an der obern Grenze der Ausdruck

\newpage
\noindent $H(t) = \sum\sum ee' \int_{0}^{t} d\vartheta \int_{-\vartheta}^{0} dt\, F'(t+\tau, t+\tau+\vartheta)$

\medskip
\noindent hinzugef\"ugt, und an der untern Grenze der Werth dieses Ausdrucks f\"ur $t=0$ hinweggenommen. Man hat also
\[P = \int_{0}^{t} dt' \sum\sum ee' \int_{0}^{t} d\vartheta\, F'(t, \tau+\vartheta) - H(t) + H(0).\]
In diesem Ausdrucke kann man $F'(t, t+\vartheta)$ durch $F'(t, t+\vartheta) - F'(\tau, \tau)$ ersetzen, da
\[\sum\sum ee' F'(\tau, \tau)\]
vernachl\"assigt werden darf. Man erh\"alt dadurch als Factor von $ee'$ einen Ausdruck, der sowohl mit $e$ als mit $e'$ sein Zeichen \"andert, so dass sich bei den Summationen die Glieder nicht gegen einander aufheben, und unendlich kleine Bruchtheile der einzelnen Glieder vernachl\"assigt werden d\"urfen. Es ergiebt sich daher, indem man
\[F'(t, \tau+\vartheta) - F'(t, \tau)\]
durch
\[\vartheta \frac{\partial F'(t,\tau)}{\partial \tau}\]
ersetzt und die Integration nach $\vartheta$ ausf\"uhrt, bis auf einen zu vernachl\"assigenden Bruchtheil
\[P = \int_{0}^{t} \frac{aa}{2} \sum\sum ee' \frac{\partial F'(t,\tau)}{\partial \tau}\, dt.\]
Es ist leicht zu sehen, dass $H(t)$ und $H(0)$ vernachl\"assigt werden d\"urfen; denn es ist
\[H(t) = \sum\sum ee' \int_{0}^{t} d\vartheta \int_{-\vartheta}^{0} d\tau\, F'(t+\tau, t+\tau+\vartheta),\]
folglich:
\[H'(t) = \sum\sum ee' \int_{0}^{t} d\vartheta \left(F'(t,t+\vartheta) - F'(t-\vartheta, t)\right).\]
Hierin aber ist nur das erste Glied des Factors von $ee'$ mit dem Factor in dem ersten Bestandtheile von $P$ von gleicher Ordnung, und dieses liefert wegen der Summation nach $e$ nur einen zu vernachl\"assigenden Bruchtheil desselben.

Der Werth von $P$, welcher sich aus unserer Theorie ergiebt, stimmt mit dem erfahrungsm\"assigen
\[\frac{1}{cc} \sum\sum ee' \frac{dd'(r^{2})}{dt\,dt'}\]
\"uberein, wenn man $aa = \frac{1}{4}cc$ annimmt.

Nach der Bestimmung von Weber und Kohlrausch ist
\[c = 439450 \cdot 10^{8} \frac{\text{mm}}{\text{Sec}},\]
woraus sich $a$ zu $41949$ geographischen Meilen in der Secunde ergiebt, w\"ahrend f\"ur die Lichtgeschwindigkeit von Busch aus Bradley's Aberrationsbeobachtungen $41994$ Meilen, und von Fizeau durch directe Messung $41882$ Meilen gefunden worden sind.

\medskip
\noindent\small Dieser Aufsatz wurde von Riemann der K\"onigl.\ Gesellschaft der Wissenschaften zu G\"ottingen am 10.\ Februar 1858 \"uberreicht, wie aus einer dem Titel des Manuscriptes hinzugef\"ugten Bemerkung des damaligen Secret\"ars der Gesellschaft hervorgeht, sp\"ater aber wieder zur\"uckgezogen. Nachdem der Aufsatz nach Riemann's Tode ver\"offentlicht worden war, wurde er durch Clausius (Poggendorff's Annalen Bd.~CXXXV p.~606) einer Kritik unterworfen, deren wesentlichster Einwand in Folgendem besteht:

Nach den Voraussetzungen hat die Summe:
\[\sum\sum \frac{ee'}{r}\]
einen verschwindend kleinen Werth. Die Operation, verm\"oge deren sp\"ater daf\"ur ein nicht verschwindend kleiner Werth gefunden wird, muss daher einen Irrthum enthalten, den Clausius in der Ausf\"uhrung einer unberechtigten Umkehrung der Integrationsfolge findet.

Der Einwand scheint mir begr\"undet und ich bin mit Clausius der Meinung, dass Riemann sich denselben selbst gemacht und deshalb die Arbeit vor der Publication zur\"uckgezogen hat.

Obwohl damit der wesentlichste Inhalt der Riemann'schen Deduction dahinfallen w\"urde, habe ich mich doch zur Aufnahme dieses Aufsatzes in die vorliegende Sammlung entschlossen, weil ich nicht zu entscheiden wagte, ob er nicht doch noch Keime zu weiteren fruchtbaren Gedanken ueber diese h\"ochst interessante Frage enth\"alt. W.

\newpage
%==============================================================================
% XV. Beweis des Satzes
%==============================================================================
\section*{XV. Beweis des Satzes, dass eine einwerthige mehr als $2n$fach periodische Function von $n$ Ver\"anderlichen unm\"oglich ist.}
\addcontentsline{toc}{section}{XV. Beweis des Satzes, dass eine einwerthige mehr als $2n$fach periodische Function unm\"oglich ist.}

\noindent\small{(Auszug aus einem Schreiben Riemann's an Herrn Weierstrass.)}

\noindent\small{(Aus Borchardt's Journal f\"ur reine und angewandte Mathematik, Bd.~71.)}

\medskip

\normalsize
\dots\ Den Beweis des Satzes, auf welchen Sie neulich die Unterhaltung lenkten, dass eine einwerthige mehr als $2n$fach periodische Function von $n$ Ver\"anderlichen unm\"oglich ist, habe ich im Gespr\"ach wohl nicht ganz klar ausgedr\"uckt, auch nur die Grundgedanken angegeben; ich theile ihn Ihnen daher hier noch einmal mit.

Es sei $f$ eine $2n$fach periodische Function von $n$ Ver\"anderlichen $x_{1}, x_{2}, \dots, x_{n}$ und --- ich darf wohl meine Ihnen bekannten Benennungen gebrauchen --- der Periodicit\"atsmodul von $x_{\nu}$ f\"ur die $\mu$te Periode $a^{\mu}_{\nu}$. Es lassen sich dann bekanntlich die Gr\"ossen $x$ in die Form
\[x_{\nu} = \sum_{\mu} a^{\mu}_{\nu} \xi_{\mu} \quad \text{f\"ur } \nu = 1, 2, \dots, n\]
setzen\textsuperscript{*}), so dass die Gr\"ossen $\xi$ reell sind. L\"asst man nun die Gr\"ossen $\xi$ die Werthe von $0$ bis $1$ mit Ausschluss eines von diesen Grenzwerthen durchlaufen, so hat das dadurch entstehende $2n$fach ausgedehnte Gr\"ossengebiet die Eigenschaft, dass jedes System von Werthen der $n$ Ver\"anderlichen einem und nur einem Werthsysteme innerhalb dieses Gr\"ossengebiets nach den $2n$ Mod\"ulsystemen congruent ist. Ich werde, um mich sp\"ater k\"urzer ausdr\"ucken zu k\"onnen, dieses Gebiet ``das bei diesen $2n$ Modulsystemen periodisch sich wiederholende Gr\"ossengebiet'' nennen.

Hat die Function nun noch ein $(2n+1)$tes Modulsystem, welches sich nicht aus den $2n$ ersten Modulsystemen zusammensetzen l\"asst, so kann man die einem Gr\"ossensysteme nach diesem Modulsysteme congruenten Gr\"ossensysteme auf innerhalb dieses Gebiets liegende nach den $2n$ ersten Modulsystemen ihnen congruente zur\"uckf\"uhren und dadurch offenbar beliebig viele innerhalb dieses Gebiets liegende und nach den $2n+1$ Modulsystemen einander congruente Gr\"ossensysteme erhalten, wenn nicht zwei von den nach dem $(2n+1)$ten Modulsysteme congruente Gr\"ossensysteme auch nach den $2n$ ersten Modulsystemen congruent sind. In diesem Falle w\"urden zwischen den $2n+1$ Modulsystemen $n$ Gleichungen von der Form
\[\xi_{\nu} = \sum_{\mu=1}^{2n+1} m^{\mu} a^{\mu}_{\nu}\]
worin die Gr\"ossen $m$ ganze Zahlen w\"aren, stattfinden, und folglich, wie ich sp\"ater zeigen werde, die $2n+1$ Modulsysteme sich aus $2n$ Modulsystemen zusammensetzen lassen.

Man theile nun f\"ur jede der Gr\"ossen $\xi$ die Strecke von $0$ bis $1$ in $q$ gleiche Theile, wodurch das bei den $2n$ ersten Modulsystemen periodisch wiederkehrende Gebiet in $q^{2n}$ Gebiete zerf\"allt, in deren jedem sich die Gr\"ossen $\xi$ nur um $\frac{1}{q}$ \"andern. Offenbar m\"ussen dann von mehr als $q^{2n}$ nach den $2n+1$ Modulsystemen einander congruenten und in jenem Gebiete liegenden Gr\"ossensystemen nothwendig zwei in dasselbe Theilgebiet fallen, so dass sich die Werthe derselben Gr\"osse $\xi$ in beiden keinesfalls um mehr als $\frac{1}{q}$ von einander unterscheiden. Die Function bleibt also dann unge\"andert, w\"ahrend keine der Gr\"ossen $\xi$ um mehr als $\frac{1}{q}$ ge\"andert wird, und ist folglich, da $q$ beliebig gross genommen werden kann, wenn sie stetig ist, eine Function von weniger als $n$ linearen Ausdr\"ucken der Gr\"ossen $x$.

Es ist nun noch zu zeigen, dass sich $2n+1$ Modulsysteme, zwischen denen die $n$ Gleichungen
\[\xi_{\nu} = \sum_{\mu=1}^{2n+1} m^{\mu} a^{\mu}_{\nu}\]
stattfinden, aus $2n$ Modulsystemen zusammensetzen lassen.

Man kann zun\"achst leicht beweisen, dass sich zu einem Modulsysteme
\[\frac{a_{\nu 1}}{\vartheta_{1}}, \frac{a_{\nu 2}}{\vartheta_{2}}, \dots, \frac{a_{\nu, 2n}}{\vartheta_{2n}},\]
worin die Gr\"ossen $m$ ganze Zahlen ohne gemeinschaftlichen Theiler sind, immer $2n-1$ andere Modulsysteme $b_{1}, b_{2}, \dots, b_{2n}$ so finden lassen,

\newpage
\noindent dass Congruenz nach den Modulsystemen $a$ mit Congruenz nach den Modulsystemen $b$ identisch ist. Es seien $\vartheta_{1}$ der gr\"osste gemeinschaftliche Theiler von $m_{1}$ und $m_{2}$ und $\alpha$, $\beta$ zwei der Gleichung
\[\beta m_{1} - \alpha m_{2} = \vartheta_{1}\]
gen\"ugende ganze Zahlen. Setzt man dann
\[\alpha a_{\nu 1} + \beta a_{\nu 2} = c_{\nu 1}\]
und
\[\alpha a'_{\nu 1} + \beta a'_{\nu 2} = b_{\nu, 2n},\]
so hat man
\[a_{\nu 1} = \beta c_{\nu 1} - \frac{m_{2}}{\vartheta_{1}} b_{\nu, 2n}, \quad a_{\nu 2} = -\alpha c_{\nu 1} + \frac{m_{1}}{\vartheta_{1}} b_{\nu, 2n}.\]
Es lassen sich also auch umgekehrt die Modulsysteme $a_{1}$ und $a_{2}$ aus den Modulsystemen $b_{\nu, 2n}$ und $c_{\nu 1}$ zusammensetzen, und folglich ist Congruenz nach jenen mit Congruenz nach diesen gleichbedeutend. Man kann daher die Modulsysteme $a_{1}$ und $a_{2}$ durch die Modulsysteme $c_{1}$ und $b_{2n}$ ersetzen. Auf dieselbe Weise kann man nun, wenn $\vartheta_{2}$ der gr\"osste gemeinschaftliche Theiler von $\vartheta_{1}$ und $m_{3}$ ist, die Modulsysteme $c_{1}$ und $a_{3}$ durch das Modulsystem
\[\frac{c_{\nu 1}}{\vartheta_{2}}\]
und durch ein Modulsystem $b_{\nu, 2n-1}$ ersetzen. Durch Fortsetzung dieses Verfahrens erh\"alt man offenbar den zu beweisenden Satz. Der Inhalt des periodisch sich wiederholenden Gebiets ist f\"ur die neuen Modulsysteme $b$ derselbe wie f\"ur die alten.

Mit H\"ulfe dieses Satzes lassen sich in den $n$ Gleichungen
\[\xi_{\nu} = \sum_{\mu=1}^{2n+1} m^{\mu} a^{\mu}_{\nu}\]
die $2n$ ersten Modulsysteme so durch $2n$ neue $b_{1}, b_{2}, \dots, b_{2n}$ ersetzen, dass diese Gleichungen die Form
\[\xi_{\nu} = \frac{p\, b^{\,2n+1}_{\nu} + q\, b^{\,1}_{\nu}}{\vartheta}\]
annehmen, worin $p$ und $q$ ganze Zahlen ohne gemeinschaftlichen Theiler sind. Sind nun $\gamma$, $\delta$ zwei der Gleichung
\[p\delta + q\gamma = 1\]
gen\"ugende ganze Zahlen, so lassen sich offenbar die beiden Modulsysteme $b^{1}$ und $a^{2n+1}$ durch das eine Modulsystem
\[\gamma\, b^{\,2n+1}_{\nu} + \delta\, b^{\,1}_{\nu}\]
ersetzen. S\"ammtliche Modulsysteme, welche sich aus den Modulsystemen $a_{1}, a_{2}, \dots, a_{2n+1}$ zusammensetzen lassen, k\"onnen also auch aus den $2n$ Modulsystemen $\frac{b}{\vartheta}$, $b_{2}, b_{3}, \dots, b_{2n}$ zusammengesetzt werden, und umgekehrt. Der Inhalt des periodisch wiederkehrenden Gebiets betr\"agt f\"ur diese $2n$ Modulsysteme nur $\frac{1}{\vartheta}$ von dem f\"ur die $2n$ ersten Modulsysteme $a$. Hat die Function nun ausser diesen Modulsystemen noch ein durch \"ahnliche ganzzahlige Gleichungen mit ihnen verbundenes, so lassen sich wieder $2n$ neue Modulsysteme finden, aus welchen sich alle diese Modulsysteme zusammensetzen lassen, und der Inhalt des periodisch sich wiederholenden Gebiets wird dabei wieder auf einen aliquoten Theil reducirt. Wenn dieses Gebiet unendlich klein wird, so wird die Function eine Function von weniger als $n$ linearen Ausdr\"ucken der Ver\"anderlichen und zwar von $n-1$ oder $n-2$ oder $n-m$, jenachdem nur eine, oder zwei oder $m$ Dimensionen dieses Gr\"ossengebiets unendlich klein werden. Soll dies aber nicht eintreten, so muss die Operation schliesslich abbrechen, und man wird also zu $2n$ Modulsystemen gelangen, aus welchen sich s\"ammtliche Modulsysteme der Function zusammensetzen lassen.

\bigskip
\noindent G\"ottingen, den 26ten October 1859.

\bigskip
\noindent\small\textsuperscript{*}) Dies ist nicht immer der Fall, sondern nur, wenn die $2n$ Gleichungen, durch welche die Gr\"ossen $\xi$ bestimmt werden, von einander unabh\"angig sind; die Ausnahmen sind aber leicht zu behandeln.

\newpage
%==============================================================================
% XVI. Extract from Italian letter to Betti
%==============================================================================
\section*{XVI. Estratto di una lettera scritta in lingua Italiana il di 21 Gennaio 1864 al Sig.\ Professore Enrico Betti.}
\addcontentsline{toc}{section}{XVI. Estratto di una lettera al Sig.\ Professore Enrico Betti.}

\noindent\small{(Annali di Matematica, Ser.~1, T.~VII.)}

\medskip

\selectlanguage{italian}

Carissimo Amico

\dots Per trovare l'attrazione di un cilindro omogeneo retto ellisoidale qualunque, io considero, introducendo coordinate rettangolari $x$, $y$, $z$, il cilindro infinito limitato della diseguaglianza:
\[\frac{x^{2}}{a^{2}} + \frac{y^{2}}{b^{2}} + \frac{z^{2}}{c^{2}} < 1,\]
ripieno di massa di densit\`a costante $+1$, se $\frac{x^{2}}{a^{2}} + \frac{y^{2}}{b^{2}} + \frac{z^{2}}{c^{2}} < 0$, e di densit\`a $-1$, se $\frac{x^{2}}{a^{2}} + \frac{y^{2}}{b^{2}} + \frac{z^{2}}{c^{2}} > 0$. Allora se poniamo, come \`e solito, il potenziale nel punto $x$, $y$, $z$ eguale a $F$ e
\[\frac{\partial F}{\partial x} = X, \quad \frac{\partial F}{\partial y} = Y, \quad \frac{\partial F}{\partial z} = Z,\]
si ha per $\frac{x^{2}}{a^{2}} + \frac{y^{2}}{b^{2}} < 0$, $F = 0$, $X = 0$, $Y = 0$.

$Z$ \`e eguale al potenziale dell' ellisse:
\[1 - \frac{x^{2}}{a^{2}} - \frac{y^{2}}{b^{2}} > 0\]
colla densit\`a $2$, e si trova col metodo di Dirichlet, se denotiamo con $\mathfrak{C}$ la radice maggiore dell' equazione:
\[\frac{x^{2}}{a^{2}+s} + \frac{y^{2}}{b^{2}+s} + \frac{z^{2}}{c^{2}+s} = 1,\]
\[Z = 2abc \int_{-\mathfrak{C}}^{\infty} \frac{\sqrt{(a^{2}+s)(b^{2}+s)}}{(c^{2}+s)^{2}} \, \frac{ds}{D},\]
con $D$:
\[D = \sqrt{(a^{2}+s)(b^{2}+s)(c^{2}+s)}.\]

$X$ ed $Y$ si possono determinare dalle equazioni:
\[\frac{\partial X}{\partial z} = \frac{\partial Z}{\partial x}, \quad \frac{\partial Y}{\partial z} = \frac{\partial Z}{\partial y}\]
e dalle condizioni:
\[X = 0, \quad Y = 0\]
per $z = 0$.

Per effettuare questa determinazione conviene di sostituire invece di $Z$ la sua espressione esteso per il contorno intero di un pezzo del Piano degli $s$, che contiene il valore $\mathfrak{C}$ senza contenere verun altro valore di diramazione o di discontinuit\`a della funzione sotto il segno integrale. Se denotiamo le radici di $F=0$ in ordine di grandezza con $\mathfrak{a}$, $\mathfrak{a}'$, $\mathfrak{a}''$ questi valori sono tutti reali e in ordine di grandezza:
\[\mathfrak{a}, \; \mathfrak{a}', \; -b^{2}, \; \mathfrak{a}'', \; -a^{2},\]
in modo che:
\[\mathfrak{a} > 0 > \mathfrak{a}' > -b^{2} > \mathfrak{a}'' > -a^{2}.\]

Posto
\[F = \mathfrak{a}' - \frac{1}{s},\]
viene
\[Z = 2abc \int_{0}^{\mathfrak{a}'} \frac{\sqrt{(a^{2}+s)(b^{2}+s)}}{(c^{2}+s)^{2}} \, \frac{ds}{D\sqrt{F}},\]
\[X = \int \frac{\partial Z}{\partial x}\, dz, \quad Y = \int \frac{\partial Z}{\partial y}\, dz.\]

Ma:
\[\frac{\partial Z}{\partial x} = \frac{\partial Z}{\partial t} \frac{\partial t}{\partial x}, \quad \frac{\partial Z}{\partial y} = \frac{\partial Z}{\partial t} \frac{\partial t}{\partial y},\]
e:
\[\frac{dt}{dz} = \sqrt{\frac{a^{2}+s}{b^{2}+s}} \, \frac{1}{D}, \quad \frac{dt}{dy} = \sqrt{\frac{b^{2}+s}{a^{2}+s}} \, \frac{1}{D}.\]

\dots Dunque si trova per integrazione parziale:
\[X = \dots\]
Se si prende la via dell' integrazione come nella espressione di $Z$ il valore dell' integrale sodisfa sempre alla condizione:
\[\frac{\partial X}{\partial z} = \frac{\partial Z}{\partial x},\]
ma pu\`o differire di funzioni di $x$ e di $y$, la funzione sotto segno integrale essendo discontinua anche per $\xi = 0$. Dunque occorre una determinazione ulteriore della via dell' integrazione.

Nella espressione di $\frac{\partial X}{\partial z}$ la funzione sotto segno integrale \`e continua per $s=0$; dunque il pezzo del piano degli $s$, per il cui contorno l'integrale \`e esteso, deve contenere $s = \mathfrak{C}$ e pu\`o contenere o no $s = 0$, ma nessuno altro dei valori sopra notati. Nella espressione di $X$ questo pezzo deve essere determinato in modo che $X$ sia $=0$ per $z=0$; e affinch\`e ci\`o avvenga, dovendo contenere $s = \mathfrak{C}$, deve anche contenere la maggiore radice di $\mathfrak{ts}=0$ (la qu\`ale \`e la maggiore radice di $\mathfrak{t}=0$, se
\[\frac{x^{2}}{a^{2}} + \frac{y^{2}}{b^{2}} - 1 > 0,\]
ed \`e $=0$, se:
\[\frac{x^{2}}{a^{2}} + \frac{y^{2}}{b^{2}} - 1 < 0)\]
ma nessun altra radice di $\mathfrak{ts}=0$. Perch\`e per $\xi=0$ le radici di $F=0$ coincidono colle radici di $\mathfrak{ts}=0$, e se la via dell' integrazione passasse tra due valori di discontinuit\`a che coincidono per $\xi=0$, doverebbe per $\xi=0$ passare per questo valore in modo che l'integrale nella espressione di $X$ diverrebbe infinito ed il valore nonostante il fattore $z$ rimarrebbe finito. ---

Vostro aff.\textsuperscript{mo} Amico Riemann.

\newpage
%==============================================================================
% XVII. Ueber die Flaeche vom kleinsten Inhalt bei gegebener Begrenzung
%==============================================================================
\selectlanguage{german}
\section*{XVII. Ueber die Fl\"ache vom kleinsten Inhalt bei gegebener Begrenzung.\textsuperscript{*})}
\addcontentsline{toc}{section}{XVII. Ueber die Fl\"ache vom kleinsten Inhalt bei gegebener Begrenzung.}

\subsection*{1.}

Eine Fl\"ache l\"asst sich im Sinne der analytischen Geometrie darstellen, indem man die rechtwinkligen Coordinaten $x$, $y$, $z$ eines in ihr beweglichen Punktes als eindeutige Functionen von zwei unabh\"angigen ver\"anderlichen Gr\"ossen $p$ und $q$ angiebt. Nehmen dann $p$ und $q$ bestimmte constante Werthe an, so entspricht dieser einen Combination immer nur ein einziger Punkt der Fl\"ache. Die unabh\"angigen Variabeln $p$ und $q$ k\"onnen in sehr mannigfacher Weise gew\"ahlt werden. F\"ur eine einfach zusammenh\"angende Fl\"ache geschieht dies zweckm\"assig wie folgt. Man l\"asst die Fl\"ache l\"angs der ganzen Begrenzung abnehmen um einen Fl\"achenstreifen, dessen Breite ueberall unendlich klein in derselben Ordnung ist. Durch Wiederholung dieses Verfahrens wird die Fl\"ache fortw\"ahrend verkleinert, bis sie in einen Punkt uebergeht. Die hierbei der Reihe nach auftretenden Begrenzungscurven sind in sich zur\"ucklaufende, von einander getrennte Linien. Man kann sie dadurch unterscheiden, dass man in jeder von ihnen der Gr\"osse $p$ einen besondern constanten Werth beilegt, der um ein Unendlichkleines zu- oder abnimmt, je nachdem man zu der benachbarten umschliessenden oder umschlossenen Curve uebergeht. Die Function $p$ hat dann einen constanten Maximalwerth in der Begrenzung der Fl\"ache und einen Minimalwerth in dem einen Punkte im Innern, in welchen die

\bigskip
\noindent\small\textsuperscript{*}) Dieser Abhandlung liegt ein Manuscript Riemann's zu Grunde, welches nach der eigenen Aeusserung des Verfassers in den Jahren 1860 und 1861 entstanden ist. Dieses Manuscript, welches in gedr\"angter K\"urze nur die Formeln und keinen Text enth\"alt, wurde mir von Riemann im April 1866 zur Bearbeitung anvertraut. Es ist daraus die Abhandlung hervorgegangen, welche ich am 6.\ Januar 1867 der K\"oniglichen Gesellschaft der Wissenschaften zu G\"ottingen eingereicht habe, und welche im 13.\ Band der Abhandlungen dieser Gesellschaft abgedruckt ist. Diese Abhandlung kommt hier in sorgf\"altiger Ueberarbeitung zum zweiten Male zum Abdruck. K.\ Hattendorff.

\newpage
\noindent allm\"ahlich abnehmende Fl\"ache zuletzt zusammenschrumpft. Den Uebergang von einer Begrenzung der abnehmenden Fl\"ache zur n\"achsten kann man dadurch hergestellt denken, dass man jeden Punkt der Curve $(p)$ in einen bestimmten unendlich nahen Punkt der Curve $(p+dp)$ uebergehen l\"asst. Die Wege der einzelnen Punkte bilden dann ein zweites System von Curven, die von dem Punkte des Minimalwerthes von $p$ strahlenf\"ormig nach der Begrenzung der Fl\"ache verlaufen. In jeder dieser Curven legt man $q$ einen besondern constanten Werth bei, der in einer beliebig gew\"ahlten Anfangscurve am kleinsten ist und von da beim Uebergange von einer Curve des zweiten Systems zur andern stetig w\"achst, wenn man zum Zweck des Ueberganges irgend eine Curve $(p)$ in bestimmter Richtung durchl\"auft. Beim Uebergange von der letzten Curve $(q)$ zur Anfangscurve \"andert sich $q$ sprungweise um eine endliche Constante.

Um eine mehrfach zusammenh\"angende Fl\"ache ebenso zu behandeln, kann man sie zuvor durch Querschnitte in eine einfach zusammenh\"angende zerlegen.

Irgend ein Punkt der Fl\"ache l\"asst sich hiernach als Durchschnitt einer bestimmten Curve des Systems $(p)$ mit einer bestimmten Curve des Systems $(q)$ auffassen. Die in dem Punkte $(p,q)$ errichtete Normale verl\"auft von der Fl\"ache aus in zwei entgegengesetzten Richtungen, der positiven und der negativen. Zu ihrer Unterscheidung hat man ueber die gegenseitige Lage der wachsenden positiven Normale, der wachsenden $p$ und der wachsenden $q$ eine Bestimmung zu treffen. Ist nichts anderes festgesetzt, so m\"oge, von der positiven $x$-Axe aus gesehen, die positive $y$-Axe auf dem k\"urzesten Wege in die positive $z$-Axe uebergef\"uhrt werden durch eine Drehung von rechts nach links. Und die Richtung der wachsenden positiven Normale liege zu den Richtungen der wachsenden $p$ und der wachsenden $q$, wie die positive $x$-Axe zur positiven $y$-Axe und zur positiven $z$-Axe. Die Seite der Fl\"ache, auf welcher die positive Normale liegt, soll die positive Seite der Fl\"ache genannt werden.

\subsection*{2.}

Ueber das Gebiet der Fl\"ache sei ein Integral zu erstrecken, dessen Element gleich ist dem Element $dp\,dq$ multiplicirt in eine Functionaldeterminante, also
\[\iint f \frac{\partial(g,h)}{\partial(p,q)}\, dp\, dq,\]
wof\"ur zur Abk\"urzung geschrieben werden soll
\[\iint f \,(dg\, dh).\]

\newpage
Denkt man sich $f$ und $g$ als unabh\"angige Variable eingef\"uhrt, so geht das Integral ueber in $\iint df\, dg$, und es l\"asst sich die Integration nach $f$ oder nach $g$ ausf\"uhren. Die wirkliche Einsetzung von $f$ und $g$ als unabh\"angigen Variabeln verursacht aber Schwierigkeiten oder wenigstens weitl\"aufige Unterscheidungen, wenn dieselbe Werthecombination von $f$ und $g$ in mehreren Punkten der Fl\"ache oder in einer Linie vorhanden ist. Sie ist ganz unm\"oglich, wenn $f$ und $g$ complex sind.

Es ist daher zweckm\"assig, zur Ausf\"uhrung der Integration nach $f$ oder $g$ das Verfahren von Jacobi (Crelle's Journal Bd.~27 p.~208) anzuwenden, bei welchem $p$ und $q$ als unabh\"angige Variable beibehalten werden. Um in Beziehung auf $f$ zu integriren, hat man die Functionaldeterminante in die Form zu bringen
\[\frac{\partial(f,h)}{\partial(p,q)} = \frac{\partial}{\partial q} \left(\frac{\partial f}{\partial p}\, h\right) - \frac{\partial}{\partial p} \left(\frac{\partial f}{\partial q}\, h\right)\]
und erh\"alt zun\"achst
\[\int\!\int \frac{\partial(f,h)}{\partial(p,q)}\, dp\, dq = 0,\]
weil die Integration durch eine in sich zur\"ucklaufende Linie erstreckt wird. Dagegen ist
\[\int\!\int \frac{\partial(g,h)}{\partial(p,q)}\, dp\, dq = \int\!\int \frac{\partial}{\partial q} \left(\frac{\partial g}{\partial p}\, h\right) dp\, dq = \int \frac{\partial g}{\partial p}\, h\, dp\]
in der Richtung der wachsenden $p$ zu nehmen, d.\ h.\ von dem Minimalpunkte im Innern durch eine Curve $(q)$ bis zur Begrenzung. Man erh\"alt $\int h\, dg$ und zwar den Werth, den dieser Ausdruck in der Begrenzung annimmt, da an der untern Grenze des Integrals $h=0$ ist. Folglich wird
\[\int\!\int f \,(dg\, dh) = \int\!\int g \,(df\, dh) = \int g\, dh\]
und das einfache Integral rechts ist in der Richtung der wachsenden $q$ durch die Begrenzung erstreckt. Andererseits hat man nach der eingef\"uhrten Bezeichnung $(df\, dg) = -(dg\, df)$, und daher
\[\int\!\int f \,(df\, dg) = -\int\!\int f \,(dg\, df) = -\int f\, dg,\]
wobei das einfache Integral rechts ebenfalls in der Richtung der wachsenden $q$ durch die Begrenzung der Fl\"ache zu nehmen ist.

\subsection*{3.}

Die Fl\"ache, deren Punkte durch die Curvensysteme $(p)$, $(q)$ festgelegt sind, soll in der folgenden Weise auf einer Kugel vom Radius $1$ abgebildet werden. Im Punkte $(p,q)$ der Fl\"ache, dessen rechtwinklige Coordinaten $x$, $y$, $z$ sind, ziehe man die positive Normale und lege zu ihr eine Parallele durch den Mittelpunkt der Kugel. Der Endpunkt dieser Parallelen auf der Kugeloberfl\"ache ist die Abbildung des Punktes $(x,y,z)$. Durchl\"auft der Punkt $(x,y,z)$ auf der stetig gekr\"ummten Fl\"ache eine zusammenh\"angende Linie, so wird auch die Abbildung derselben auf der Kugel eine zusammenh\"angende Linie sein. Auf dieselbe Weise erh\"alt man als Abbildung eines Fl\"achenst\"ucks ein Fl\"achenst\"uck, als Abbildung der ganzen Fl\"ache eine Fl\"ache, welche die Kugel oder einen Theil derselben einfach oder mehrfach bedeckt.

Der Punkt auf der Kugel, welcher die Richtung der positiven $x$-Axe angiebt, werde zum Pol gew\"ahlt und der Anfangsmeridian durch den Punkt gelegt, welcher der positiven $y$-Axe entspricht. Die Abbildung des Punktes $(x,y,z)$ wird dann auf der Kugel festgelegt durch ihre Poldistanz $r$ und den Winkel $\varphi$, welchen ihr Meridian mit dem Anfangsmeridian einschliesst. F\"ur das Vorzeichen von $\varphi$ gilt die Bestimmung, dass der der positiven $z$-Axe entsprechende Punkt die Coordinaten $r=\frac{\pi}{2}$, $\varphi = +\frac{\pi}{2}$ haben soll.

\subsection*{4.}

Hiernach erh\"alt man als Differential-Gleichung der Fl\"ache
\begin{equation}\tag{1}
\cos r\, dx + \sin r \cos\varphi\, dy + \sin r \sin\varphi\, dz = 0.
\end{equation}
Sind $y$ und $z$ die unabh\"angigen Variabeln, so ergeben sich f\"ur $r$ und $\varphi$ die Gleichungen
\[\cos r = \pm \frac{1}{\sqrt{1 + \left(\frac{\partial x}{\partial y}\right)^{2} + \left(\frac{\partial x}{\partial z}\right)^{2}}},\]
\[\sin r \cos\varphi = \frac{\frac{\partial x}{\partial y}}{\sqrt{1 + \left(\frac{\partial x}{\partial y}\right)^{2} + \left(\frac{\partial x}{\partial z}\right)^{2}}},\]
\[\sin r \sin\varphi = \frac{\frac{\partial x}{\partial z}}{\sqrt{1 + \left(\frac{\partial x}{\partial y}\right)^{2} + \left(\frac{\partial x}{\partial z}\right)^{2}}},\]
in welchen gleichzeitig entweder die oberen oder die unteren Vorzeichen gelten.

\newpage
Ein Parallelogramm auf der positiven Seite der Fl\"ache, begrenzt von den Curven $(p)$ und $(p+dp)$, $(q)$ und $(q+dq)$, projicirt sich auf der $yz$-Ebene in einem Fl\"achenelemente, dessen Inhalt gleich dem absoluten Werthe von $(dy\, dz)$ ist. Das Vorzeichen dieser Functionaldeterminante ist verschieden, je nachdem die im Punkte $(p,q)$ errichtete positive Normale mit der positiven $x$-Axe einen spitzen oder stumpfen Winkel einschliesst. In dem ersten Falle liegen n\"amlich die Projectionen von $dp$ und $dq$ in der $yz$-Ebene ebenso zu einander wie die positive $y$-Axe zur positiven $z$-Axe, im zweiten Falle umgekehrt. Daher ist die Functionaldeterminante im ersten Falle positiv, im zweiten negativ. Und der Ausdruck
\[\sqrt{(dy\, dz)^{2}}\]
ist immer positiv. Er giebt den Inhalt des unendlich kleinen Parallelogramms auf der Fl\"ache. Um also den Inhalt der Fl\"ache selbst zu erhalten, hat man das Doppelintegral
\[\iint \sqrt{(dy\, dz)^{2}}\]
\"uber die ganze Fl\"ache zu erstrecken.

Soll dieser Inhalt ein Minimum sein, so ist die erste Variation des Doppelintegrals $=0$ zu setzen. Man erh\"alt
\[\delta \iint \sqrt{(dy\, dz)^{2}} = \iint \frac{(dy\, dz)\, (d\delta y\, dz + dy\, d\delta z)}{\sqrt{(dy\, dz)^{2}}} = 0,\]
und es gilt das obere oder das untere Zeichen vor der Wurzel, je nachdem $(dy\, dz)$ positiv oder negativ ist. Die linke Seite l\"asst sich schreiben
\[\delta \iint (\sin r \sin\varphi\, dy - \sin r \cos\varphi\, dz)\, (dy\, dz)\]
\[= \iint \delta(\sin r \sin\varphi)\, (dy\, dz) - \iint \delta(\sin r \cos\varphi)\, (dy\, dz)\]
\[+ \iint \sin r \sin\varphi\, (d\delta y\, dz) - \iint \sin r \cos\varphi\, (dy\, d\delta z).\]
Die beiden ersten Integrale reduciren sich auf einfache Integrale, die in der Richtung der wachsenden $q$ durch die Begrenzung der Fl\"ache zu nehmen sind, n\"amlich
\[\int \delta x \, (-\sin r \cos\varphi\, dz + \sin r \sin\varphi\, dy).\]
Der Werth ist $=0$, da in der Begrenzung $\delta x = 0$ ist. Die Bedingung des Minimum lautet also
\[\iint \frac{\partial(\sin r \sin\varphi)}{\partial y} \delta x\, (dy\, dz) + \iint \frac{\partial(\sin r \cos\varphi)}{\partial z} \delta x\, (dy\, dz) = 0.\]
Sie wird erf\"ullt, wenn
\begin{equation}\tag{2}
-\sin r \sin\varphi\, dy + \sin r \cos\varphi\, dz = d\psi
\end{equation}
ein vollst\"andiges Differential ist.

\subsection*{5.}

Die Coordinaten $r$ und $\varphi$ auf der Kugel lassen sich ersetzen durch eine complexe Gr\"osse $\eta = \tg\frac{r}{2} e^{\varphi i}$, deren geometrische Bedeutung leicht zu erkennen ist. Legt man n\"amlich an die Kugel im Pol eine Tangentialebene, deren positive Seite von der Kugel abgekehrt ist, und zieht vom Gegenpol eine Gerade durch den Punkt $(r,\varphi)$, so trifft diese die Tangentialebene in einem Punkte, der die complexe Gr\"osse $2\eta$ repr\"asentirt. Dem Pol entspricht $\eta = 0$, dem Gegenpol $\eta = \infty$. F\"ur die Punkte, welche die Richtungen der positiven $y$- und der positiven $z$-Axe angeben, ist $i\eta = +1$ und resp.\ $=+i$.

F\"uhrt man noch die complexen Gr\"ossen
\[\eta' = \tg\frac{r}{2} e^{-\varphi i}, \quad s = y + zi, \quad s' = y - zi\]
ein, so gehen die Gleichungen (1) und (2) ueber in folgende:
\begin{equation}\tag{1$^{*}$}
(1 - \eta\eta')\, dx + \eta\, ds' + \eta'\, ds = 0,
\end{equation}
\begin{equation}\tag{2$^{*}$}
(1 + \eta\eta')\, d\psi i - \eta'\, ds + \eta\, ds' = 0.
\end{equation}
Diese lassen sich durch Addition und Subtraction verbinden. Dabei werde
\[x + \psi i = 2X, \quad x - \psi i = 2X'\]
gesetzt, so dass umgekehrt $x = X + X'$ ist. Das Problem findet dann seinen analytischen Ausdruck in den beiden Gleichungen
\begin{equation}\tag{3}
ds - \eta\, dX + \eta'\, dX' = 0,
\end{equation}
\begin{equation}\tag{4}
ds' - \eta'\, dX + \eta\, dX' = 0.
\end{equation}
Betrachtet man $X$ und $X'$ als unabh\"angige Variable und stellt die Bedingungen daf\"ur auf, das $ds$ und $ds'$ vollst\"andige Differentiale sind, so findet sich
\[\frac{\partial \eta}{\partial X'} = 0, \quad \frac{\partial \eta'}{\partial X} = 0,\]
d.\ h.\ es ist $\eta$ nur von $X$, $\eta'$ nur von $X'$ abh\"angig, und deshalb umgekehrt $X$ eine Function nur von $\eta$, $X'$ eine Function nur von $\eta'$.

\newpage
Hiernach ist die Aufgabe darauf zur\"uckgef\"uhrt, $\eta$ als Function der complexen Variabeln $X$ oder umgekehrt $X$ als Function der complexen Variabeln $\eta$ so zu bestimmen, dass zugleich den Grenzbedingungen Gen\"uge geleistet werde. Kennt man $\eta$ als Function von $X$, so ergiebt sich daraus $\eta'$, indem man in dem Ausdrucke von $\eta$ jede complexe Zahl in die conjugirte verwandelt. Alsdann hat man nur noch die Gleichungen (3) und (4) zu integriren, um die Ausdr\"ucke f\"ur $s$ und $s'$ zu erlangen. Aus diesen erh\"alt man endlich durch Elimination von $i$ eine Gleichung zwischen $x$, $y$, $z$, die Gleichung der Minimalfl\"ache.

\subsection*{6.}

Sind die Gleichungen (3) und (4) integrirt, so l\"asst sich auch der Inhalt der Minimalfl\"ache selbst leicht angeben, n\"amlich
\[S = 2 \iint \sqrt{(dy\, dz)^{2}}.\]
Die Functionaldeterminante $(dy\, dz)$ formt sich in folgender Weise um
\[(dy\, dz) = \left(\frac{\partial y}{\partial X} \frac{\partial z}{\partial X'} - \frac{\partial y}{\partial X'} \frac{\partial z}{\partial X}\right) (dX\, dX') = -(dX\, dX').\]
Danach erh\"alt man
\[S = 2 \iint (dX\, dX').\]
Zur weiteren Umformung dieses Ausdruckes kann man $y$ aus $Y$ und $Y'$, $z$ aus $Z$ und $Z'$ ebenso zusammensetzen wie $x$ aus $X$ und $X'$, so dass die Gleichungen gelten
\[x = X + X', \quad x' = X - X',\]
\[y = Y + Y', \quad y' = Y - Y',\]
\[z = Z + Z', \quad z' = Z - Z'.\quad \text{q.\ e.\ d.}\]

\newpage
Alsdann erh\"alt man schliesslich
\begin{equation}\tag{5}
\begin{aligned}
S &= -i \iint \left[(dX\, dX') + (dY\, dY') + (dZ\, dZ')\right] \\
&= \tfrac{1}{2} \iint \left[(dx\, dx') + (dy\, dy') + (dz\, dz')\right].
\end{aligned}
\end{equation}

\subsection*{7.}

Die Minimalfl\"ache und ihre Abbildungen auf der Kugel wie in den Ebenen, deren Punkte resp.\ die complexen Gr\"ossen $\eta$, $X$, $Y$, $Z$ repr\"asentiren, sind einander in den kleinsten Theilen \"ahnlich. Man erkennt dies sofort, wenn man das Quadrat des Linearelementes in diesen Fl\"achen ausdr\"uckt. Dasselbe ist

auf der Kugel $\sin r^{2}\, d\log\eta\, d\log\eta'$,

in der Ebene der $\eta$ $\frac{4}{(1+\eta\eta')^{2}}\, d\eta\, d\eta'$,

in der Ebene der $X$ $\frac{dX}{d\eta} \frac{dX'}{d\eta'}\, d\eta\, d\eta'$,

in der Ebene der $Y$ $\frac{dY}{d\eta} \frac{dY'}{d\eta'}\, d\eta\, d\eta'$,

in der Ebene der $Z$ $\frac{dZ}{d\eta} \frac{dZ'}{d\eta'}\, d\eta\, d\eta'$,

in der Minimalfl\"ache selbst
\[\begin{aligned}
dx^{2} + dy^{2} + dz^{2} &= (dX+dX')^{2} + (dY+dY')^{2} + (dZ+dZ')^{2} \\
&= 2(dX\, dX' + dY\, dY' + dZ\, dZ') \\
&= 2\left(\frac{dX}{d\eta} \frac{dX'}{d\eta'} + \frac{dY}{d\eta} \frac{dY'}{d\eta'} + \frac{dZ}{d\eta} \frac{dZ'}{d\eta'}\right) d\eta\, d\eta'.
\end{aligned}\]
Es ist n\"amlich nach den Gleichungen (3) und (4), wenn man darin $\eta$ und $\eta'$ als unabh\"angige Variable ansieht:
\[\frac{dX}{d\eta} = \frac{ds}{d\eta} \frac{1}{2\eta'}, \quad \frac{dX'}{d\eta'} = \frac{ds'}{d\eta'} \frac{1}{2\eta},\]
\[\frac{dX}{d\eta'} = \frac{ds}{d\eta'} \frac{1}{2\eta}, \quad \frac{dX'}{d\eta} = \frac{ds'}{d\eta} \frac{1}{2\eta'},\]
und deshalb
\[dX^{2} + dY^{2} + dZ^{2} = 0,\]
\[dX'^{2} + dY'^{2} + dZ'^{2} = 0.\]
Das Verh\"altniss von irgend zwei der obigen Quadrate von Linearelementen ist unabh\"angig von $d\eta$ und $d\eta'$, d.\ h.\ von der Richtung des Elementes, und darin beruht die in den kleinsten Theilen \"ahnliche Abbildung. Da die Linearvergr\"osserung bei der Abbildung in irgend einem Punkte nach allen Richtungen dieselbe ist, so erh\"alt man die Fl\"achenvergr\"osserung gleich dem Quadrat der Linearvergr\"osserung. Das Quadrat des Linearelementes in der Minimalfl\"ache ist aber gleich

\newpage
\noindent der doppelten Summe der Quadrate der entsprechenden Linearelemente in den Ebenen der $X$, der $Y$ und der $Z$. Daher ist auch das Fl\"achenelement in der Minimalfl\"ache gleich der doppelten Summe der entsprechenden Fl\"achenelemente in jenen Ebenen. Dasselbe gilt von der ganzen Fl\"ache und ihren Abbildungen in den Ebenen der $X$, $Y$, $Z$.

\subsection*{8.}

Eine wichtige Folgerung l\"asst sich noch aus dem Satze von der Aehnlichkeit in den kleinsten Theilen ziehen, wenn man eine neue complexe Variable $\eta_{1}$ dadurch einf\"uhrt, dass man auf der Kugel den Pol in einen beliebigen Punkt $(\eta = a)$ verlegt und den Anfangsmeridian beliebig w\"ahlt. Hat dann $\eta_{1}$ f\"ur das neue Coordinatensystem dieselbe Bedeutung wie $\eta$ f\"ur das alte, so kann man jetzt ein unendlich kleines Dreieck auf der Kugel sowohl in der Ebene der $\eta$ als in der der $\eta_{1}$ abbilden. Die beiden Bilder sind dann auch Abbildungen von einander und in den kleinsten Theilen \"ahnlich. F\"ur den Fall der directen Aehnlichkeit ergiebt sich ohne Weiteres, dass $\frac{d\eta_{1}}{d\eta}$ unabh\"angig ist von der Richtung der Verschiebung von $\eta$, d.\ h.\ dass $\eta_{1}$ eine Function der complexen Variabeln $\eta$ ist. Den Fall der inversen (symmetrischen) Aehnlichkeit kann man auf den vorigen zur\"uckf\"uhren, indem man statt $\eta_{1}$ die conjugirte complexe Gr\"osse nimmt. Um nun $\eta_{1}$ als Function von $\eta$ auszudr\"ucken, hat man zu beachten, dass $\eta_{1} = 0$ ist in dem einen Punkte der Kugel, f\"ur welchen $\eta = a$, und $\eta_{1} = \infty$ in dem diametral gegen\"uberliegenden Punkte, d.\ h.\ f\"ur $\eta = -\frac{1}{a'}$. Danach ergiebt sich $\eta_{1} = c\frac{\eta - a}{1 + a'\eta}$. Zur Bestimmung der Constanten $c$ dient die Bemerkung, dass, wenn $\eta_{1} = \beta$ ist f\"ur $\eta = 0$, daraus $\eta_{1} = -\frac{1}{\beta'}$ gefunden wird f\"ur $\eta = \infty$. Es ist also $\beta = -ca'$ und $-\frac{1}{\beta'} = \frac{c}{a}$, d.\ h.\ $\beta\beta' = \frac{c}{a}a'$. Hieraus ergiebt sich $cc' = 1$ und daher $c = e^{\vartheta i}$ f\"ur ein reelles $\vartheta$. Die Gr\"ossen $\alpha$ und $\vartheta$ k\"onnen beliebige Werthe erhalten: $\alpha$ h\"angt von der Lage des neuen Pols, $\vartheta$ von der Lage des neuen Anfangsmeridians ab. Diesem neuen Coordinatensystem auf der Kugel entsprechen die Richtungen der Axen eines neuen rechtwinkligen Systems. Es m\"ogen in dem neuen System $X_{1}$, $s_{1}$, $s'_{1}$ dasselbe bezeichnen wie $X$, $s$, $s'$ in dem alten. Dann erlangt man die Transformationsformeln

\newpage
\begin{equation}\tag{6}
\begin{aligned}
(1 + \alpha\alpha') X_{1} &= (1 - \alpha\alpha') X + \alpha' s + \alpha s', \\
(1 + \alpha\alpha') s'_{1} e^{\vartheta i} &= -2\alpha X - \alpha^{2} s + s.
\end{aligned}
\end{equation}

\subsection*{9.}

Aus den Transformationsformeln (6) berechnen wir
\[\left(\frac{d\eta_{1}}{d\eta}\right)^{2} \frac{dX_{1}}{d\eta_{1}} = \frac{dX}{d\eta}\]
oder
\[\left(\frac{d\eta_{1}}{d\eta}\right)^{2} dX_{1} = \left(\frac{d\log\eta_{1}}{d\log\eta}\right)^{2} dX.\]
Hiernach empfiehlt es sich, eine neue complexe Gr\"osse $u$ einzuf\"uhren, welche durch die Gleichung definirt wird
\begin{equation}\tag{7}
du = \left(\frac{d\log\eta_{1}}{d\log\eta}\right)^{2} dX,
\end{equation}
und die von der Lage des Coordinatensystems $(x,y,z)$ unabh\"angig ist. Gelingt es dann, $u$ als Function von $\eta$ zu bestimmen, so erh\"alt man
\begin{equation}\tag{8}
X = \int du \left(\frac{d\log\eta_{1}}{d\log\eta}\right)^{-2}.
\end{equation}
$X$ ist der Abstand des zu $\eta$ geh\"origen Punktes der Minimalfl\"ache von einer Ebene, die durch den Anfangspunkt der Coordinaten rechtwinklig zur Richtung $\eta = 0$ gelegt ist. Man erh\"alt den Abstand desselben Punktes der Minimalfl\"ache von einer durch den Anfangspunkt der Coordinaten gelegten Ebene, die rechtwinklig auf der Richtung $\eta = \alpha$ steht, indem man in (8) $-e^{\vartheta i}\frac{1}{\eta}$ statt $\eta$ setzt. Speciell also f\"ur $\alpha = 1$ und $\alpha = i$
\begin{equation}\tag{9}
x = \int du, \quad y = \int \frac{du}{\eta^{2}}, \quad z = \int \frac{du}{\eta^{2}}.
\end{equation}

\newpage
\subsection*{10.}

Die Gr\"osse $u$ ist als Function von $\eta$ zu bestimmen, d.\ h.\ als einwerthige Function des Ortes in derjenigen Fl\"ache, welche, ueber die $\eta$-Ebene ausgebreitet, die Minimalfl\"ache in den kleinsten Theilen \"ahnlich abbildet. Daher kommt es vor allen Dingen auf die Unstetigkeiten und Verzweigungen in dieser Abbildung an. Bei der Untersuchung derselben hat man Punkte im Innern der Fl\"ache von Begrenzungspunkten zu unterscheiden.

Handelt es sich um einen Punkt im Innern der Minimalfl\"ache, so lege man in ihn den Anfangspunkt des Coordinatensystems $(x,y,z)$, die Axe der positiven $x$ in die positive Normale, folglich die $yz$-Ebene tangential. Dann fehlen in der Entwicklung von $x$ das freie Glied und die in $y$ und $z$ multiplicirten Glieder. Durch geeignet gew\"ahlte Richtung der $y$-Axe und der $z$-Axe kann man auch das in $yz$ multiplicirte Glied verschwinden lassen. Die partielle Differentialgleichung der Minimalfl\"ache reducirt sich unter dieser Voraussetzung f\"ur unendlich kleine Werthe von $y$ und $z$ auf $\frac{\partial^{2}x}{\partial y^{2}} + \frac{\partial^{2}x}{\partial z^{2}} = 0$. Das Kr\"ummungsmass ist also negativ, die Haupt-Kr\"ummungsradien sind einander entgegengesetzt gleich. Die Tangentialebene theilt die Fl\"ache in vier Quadranten, wenn die Kr\"ummungshalbmesser nicht $\infty$ sind. Diese Quadranten liegen abwechselnd ueber und unter der Tangentialebene. Beginnt die Entwicklung von $x$ erst mit den Gliedern $n$ter Ordnung ($n>2$), so sind die Kr\"ummungsradien $\infty$, und die Tangentialebene theilt die Fl\"ache in $2n$ Sectoren, die abwechselnd ueber und unter jener Ebene liegen und von den Kr\"ummungslinien halbirt werden.\textsuperscript{(*)})

Will man nun $X$ als Function der complexen Variabeln $Y$ ansehen, so ergiebt sich in dem Falle der vier Sectoren
\[\log X = 2\log Y + \text{funct.\ cont.},\]
in dem Falle der $2n$ Sectoren
\[\log X = n\log Y + \text{f.\ c.}\]
Und da nach (8) und (9) $\frac{du}{d\log\eta} = dX$ ist, so beginnt die Entwicklung von $\eta$ im ersten Falle mit der ersten, im zweiten mit der $(n-1)$ten Potenz von $Y$. Umgekehrt wird also, wenn $Y$ als Function von $\eta$ angesehen werden soll, die Entwicklung im ersten Falle nach ganzen Potenzen von $\eta$, im zweiten nach ganzen Potenzen von $\eta^{\frac{1}{n-1}}$ fortschreiten. D.\ h.\ die Abbildung auf der $\eta$-Ebene hat an der betreffenden Stelle keinen oder einen $(n-2)$fachen Verzweigungspunkt, je nachdem der erste oder der zweite Fall eintritt.

\newpage
Was $u$ betrifft, so ergiebt sich $\frac{du}{d\log Y} = \frac{du}{d\log\eta} \frac{d\log\eta}{d\log Y}$, also mit H\"ulfe der Gleichung (9)
\[\frac{du}{d\log Y} = 2\, \frac{dX}{d\log\eta}\, \frac{d\log\eta}{d\log Y}.\]
Demnach ist in einem $(n-2)$fachen Verzweigungspunkte der Abbildung auf der $\eta$-Ebene
\[\log \frac{du}{d\log Y} = 2\log Y + \text{f.\ c.}\]
oder
\[u = \text{const.} + Y^{2}.\]

\subsection*{11.}

Die weitere Untersuchung soll zun\"achst auf den Fall beschr\"ankt werden, dass die gegebene Begrenzung aus geraden Linien besteht. Dann l\"asst sich die Abbildung der Begrenzung auf der $\eta$-Ebene wirklich herstellen. Die in irgend welchen Punkten einer geraden Begrenzungslinie errichteten Normalen liegen in parallelen Ebenen, und daher ist die Abbildung auf der Kugel ein Bogen eines gr\"ossten Kreises.

Um einen Punkt im Innern einer geraden Begrenzungslinie zu untersuchen, legt man wie vorher in ihn den Anfangspunkt der Coordinaten, die positive $x$-Axe in die positive Normale. Dann f\"allt die ganze Begrenzungslinie in die $yz$-Ebene. Der reelle Theil von $X$ ist demnach in der ganzen Begrenzungslinie $=0$. Geht man also durch das Innere der Minimalfl\"ache um den Anfangspunkt der Coordinaten herum von einem vorangehenden bis zu einem nachfolgenden Begrenzungspunkte, so muss dabei der Arcus von $X$ sich \"andern um $n\pi$, ein ganzes Vielfaches von $\pi$. Der Arcus von $Y$ \"andert sich gleichzeitig um $\pi$. Man hat also, wie vorher
\[\log X = n\log Y + \text{f.\ c.}\]
\[\log \eta = (n-1)\log Y + \text{f.\ c.}\]
Dem betrachteten Begrenzungspunkte entspricht ein $(n-2)$facher Verzweigungspunkt in der Abbildung auf der $\eta$-Ebene. In dieser Abbildung macht das auf den Punkt folgende Begrenzungsst\"uck mit dem ihm vorhergehenden den Winkel $(n-1)\pi$.

\subsection*{12.}

Bei dem Uebergange von einer Begrenzungslinie zur folgenden hat man zwei F\"alle zu unterscheiden. Entweder treffen sie zusammen in einem im Endlichen liegenden Schnittpunkte, oder sie erstrecken sich ins Unendliche.

Im ersten Falle sei $\alpha\pi$ der im Innern der Minimalfl\"ache liegende Winkel der beiden Begrenzungslinien. Legt man den Anfangspunkt der Coordinaten in den zu untersuchenden Eckpunkt, die positive $x$-Axe in die positive Normale, so ist in beiden Begrenzungslinien der reelle Theil von $X = 0$. Beim Uebergange von der ersten Begrenzungslinie zur folgenden \"andert sich also der Arcus von $X$ um $m\pi$, ein ganzes Vielfaches von $\pi$, der Arcus von $Y$ um $\alpha\pi$. Man hat daher
\[\log X = m\log Y + \text{f.\ c.}\]
\[\left(1 - \frac{\alpha}{m}\right)\log\eta = \log Y + \text{f.\ c.}\]

Erstreckt sich die Fl\"ache zwischen zwei auf einander folgenden Begrenzungsgeraden ins Unendliche, so lege man die positive $x$-Axe in ihre k\"urzeste Verbindungslinie, parallel der positiven Normalen im Unendlichen. Die L\"ange der k\"urzesten Verbindungslinie sei $A$, und $\alpha\pi$ der Winkel, welchen die Projection der Minimalfl\"ache in der $yz$-Ebene ausf\"ullt. Dann bleiben die reellen Theile von $X$ und $i\, d\log\eta$ im Unendlichen endlich und stetig und nehmen in den begrenzenden Geraden constante Werthe an. Hieraus ergiebt sich (f\"ur $y=\infty$, $z=\infty$)
\[X = \frac{A}{2\alpha\pi}\, d\log\eta + \text{f.\ c.}\]
Legt man die $x_{1}$-Axe eines Coordinatensystems in eine begrenzende Gerade, die $x_{2}$-Axe eines andern Systems in die zweite begrenzende Gerade u.\ s.\ f., so ist in der ersten Linie $\log\eta_{1}$, in der zweiten $\log\eta_{2}$ u.\ s.\ f.\ rein imagin\"ar, da die Normale zu der betreffenden Axe der $x_{1}$, der $x_{2}$ u.\ s.\ f.\ senkrecht steht. Es ist also $i\frac{du}{d\log\eta_{1}}$ in der ersten Begrenzungslinie reell, $i\frac{du}{d\log\eta_{2}}$ in der zweiten u.\ s.\ f.\ Da aber auch f\"ur ein beliebiges Coordinatensystem $(x,y,z)$ immer
\[i\, du = i\left(\frac{d\log\eta_{1}}{d\log\eta}\right)^{2} \frac{du}{d\log\eta_{1}}\]
ist, so findet sich, dass in jeder geraden Begrenzungslinie
\[i\sqrt{\frac{du}{d\log\eta}}\]
entweder reelle oder rein imagin\"are Werthe besitzt.

\newpage
\subsection*{13.}

Die Minimalfl\"ache ist bestimmt, sobald man eine der Gr\"ossen $u$, $\eta$, $X$, $Y$, $Z$ durch eine der uebrigen ausgedr\"uckt hat. Dies gelingt in vielen F\"allen. Besondere Beachtung verdienen darunter diejenigen, in welchen $\frac{du}{d\log\eta}$ eine algebraische Function von $\eta$ ist. Dazu ist n\"othig und hinreichend, dass die Abbildung auf der Kugel und ihre symmetrischen und congruenten Fortsetzungen eine geschlossene Fl\"ache bilden, welche die ganze Kugel einfach oder mehrfach bedeckt.

Im Allgemeinen aber wird es schwierig sein, direct eine der Gr\"ossen $u$, $\eta$, $X$, $Y$, $Z$ durch eine der uebrigen auszudr\"ucken. Statt dessen kann man aber auch jede von ihnen als Function einer neuen zweckm\"assig gew\"ahlten unabh\"angigen Variabeln bestimmen. Wir f\"uhren eine solche unabh\"angige Variable $t$ ein, dass die Abbildung der Fl\"ache auf der $t$-Ebene die halbe unendliche Ebene einfach bedeckt, und zwar diejenige H\"alfte, f\"ur welche der imagin\"are Theil von $t$ positiv ist. In der That ist es immer m\"oglich, $t$ als Function von $u$ (oder von irgend einer der uebrigen Gr\"ossen $\eta$, $X$, $Y$, $Z$) in der Fl\"ache so zu bestimmen, dass der imagin\"are Theil in der Begrenzung $=0$ ist, und dass sie in einem beliebigen Begrenzungspunkte ($u=b$) unendlich von der ersten Ordnung wird, d.\ h.
\[t = \frac{\text{const.}}{u-b}.\]
Der Arcus des Factors von $\frac{1}{u-b}$ ist durch die Bedingung bestimmt, dass der imagin\"are Theil von $t$ in der Begrenzung $=0$, im Innern der Fl\"ache positiv sein soll. Es bleibt also in dem Ausdrucke von $t$ nur der Modul dieses Factors und eine additive Constante willk\"urlich.

Es sei $t = a_{1}, a_{2}, \dots$ f\"ur die Verzweigungspunkte im Innern der Abbildung auf der $\eta$-Ebene, $t = b_{1}, b_{2}, \dots$ f\"ur die Verzweigungspunkte in der Begrenzung, die nicht Eckpunkte sind, $t = c_{1}, c_{2}, \dots$ f\"ur die Eckpunkte, $t = e_{1}, e_{2}, \dots$ f\"ur die ins Unendliche sich erstreckenden Sectoren. Wir wollen der Einfachheit wegen voraussetzen, dass die s\"ammtlichen Gr\"ossen $a$, $b$, $c$, $e$ im endlichen Gebiete der $t$-Ebene liegen.

Dann hat man

f\"ur $t = a$: $\log\frac{du}{d\log\eta} = (\gamma - 1)\log(t-a) + \text{f.\ c.}$,

f\"ur $t = b$: $\log\frac{du}{d\log\eta} = \left(\frac{\gamma}{2} - 1\right)\log(t-b) + \text{f.\ c.}$,

f\"ur $t = c$: $\log\frac{du}{d\log\eta} = \left(\frac{\gamma}{2} - 1\right)\log(t-c) + \text{f.\ c.}$,

f\"ur $t = e$: $u = \gamma\sqrt{t-e}\, \log(t-e) + \text{f.\ c.}$

\newpage
Man kann die Untersuchung auf den Fall $n=3$, $m=1$ beschr\"anken, d.\ h.\ auf einfache Verzweigungspunkte, und den allgemeinen Fall aus diesem dadurch ableiten, dass man mehrere einfache Verzweigungspunkte zusammenfallen l\"asst.

Um den Ausdruck f\"ur $\frac{du}{dt}$ zu bilden, hat man zu beachten, dass l\"angs der Begrenzung $dt$ reell, $du$ entweder reell oder rein imagin\"ar ist. Demnach ist $i\frac{du}{dt}$ reell, wenn $t$ reell ist. Diese Function kann man ueber die Linie der reellen Werthe von $t$ hin\"uber stetig fortsetzen, indem man die Bestimmung trifft, dass f\"ur conjugirte Werthe $t$ und $t'$ der Variabeln auch die Function conjugirte Werthe haben soll. Alsdann ist $i\frac{du}{dt}$ f\"ur die ganze $t$-Ebene bestimmt und zeigt sich einwerthig.

Es seien $a'_{1}, a'_{2}, \dots$ die conjugirten Werthe zu $a_{1}, a_{2}, \dots$; und das Product $(t-a_{1})(t-a_{2})\cdots$ werde mit $\prod(t-a)$ bezeichnet. Alsdann ist
\begin{equation}\tag{11}
du = \text{const.} \cdot i \sqrt{ \frac{ \prod(t-a)\prod(t-a')\prod(t-b) }{ \prod(t-c)^{2}\prod(t-e)^{2}} } \, dt.
\end{equation}
Die Constanten $a$, $b$, $c$ etc.\ m\"ussen so bestimmt werden, dass f\"ur $t=e$: $u = \sqrt{t-e}\, \log(t-e) + \text{f.\ c.}$ wird. Damit $u$ f\"ur alle Werthe von $t$ ausser $a$, $b$, $c$, $e$ endlich und stetig bleibe, muss f\"ur die Anzahl dieser letztgenannten Werthe eine Relation bestehen. Es muss die Differenz der Anzahl der Eckpunkte und der in der Begrenzung liegenden Verzweigungspunkte um $4$ gr\"osser sein als die doppelte Differenz der Anzahl der innern Verzweigungspunkte und der ins Unendliche verlaufenden Sectoren. Setzt man zur Abk\"urzung
\[\prod(t-a)\prod(t-a')\prod(t-b) = \varphi(t),\]
\[\prod(t-c)\prod(t-e) = \chi(t),\]
\newpage
\noindent d.\ h.
\[\frac{du}{dv} = \text{const.} \cdot i \sqrt{\frac{\varphi(t)}{\chi(t)^{2}}},\]
so ist die ganze Function $\varphi(t)$ vom Grade $\nu - 4$, wenn $\chi(t)$ vom Grade $\nu$ ist. Hier bedeutet $\nu$ die Anzahl der Eckpunkte vermehrt um die doppelte Anzahl der ins Unendliche verlaufenden Sectoren.

\subsection*{14.}

Es ist noch $\eta$ als Function von $t$ auszudr\"ucken. Direct gelangt man dazu nur in den einfachsten F\"allen. Im Allgemeinen ist der folgende Weg einzuschlagen. Es sei $v$ eine noch n\"aher zu bestimmende Function von $t$, die als bekannt vorausgesetzt wird. In den Gleichungen (8), (9), (10) kommt es wesentlich an auf $\frac{dX}{d\log\eta}$, wof\"ur man auch schreiben kann $\frac{dX}{dv} \cdot \frac{dv}{d\log\eta}$.

Der letzte Factor l\"asst sich ansehen als Product der beiden Factoren
\begin{equation}\tag{12}
k_{1} = \sqrt{\frac{dv}{d\eta}}, \quad k_{2} = \eta\sqrt{\frac{dv}{d\eta}},
\end{equation}
die der Differentialgleichung erster Ordnung gen\"ugen
\begin{equation}\tag{13}
k_{2} \frac{dk_{1}}{dv} - k_{1} \frac{dk_{2}}{dv} = 1,
\end{equation}
sowie der Differentialgleichung zweiter Ordnung
\begin{equation}\tag{14}
\frac{1}{k} \frac{d^{2}k}{dv^{2}} = \frac{1}{k_{1}} \frac{d^{2}k_{1}}{dv^{2}} = \frac{1}{k_{2}} \frac{d^{2}k_{2}}{dv^{2}}.
\end{equation}
Gelingt es also, die eine oder die andere Seite dieser letzten Gleichung als Function von $t$ auszudr\"ucken, so l\"asst sich eine homogene lineare Differentialgleichung zweiter Ordnung herstellen, von welcher $k_{1}$ und $k_{2}$ particul\"are Integrale sind. Es sei $k$ das vollst\"andige Integral. Wir ersetzen $\frac{d^{2}k}{dv^{2}}$ durch das ihm gleichbedeutende
\[\frac{dv}{dt} \frac{d^{2}k}{dt^{2}} - \frac{dk}{dt} \frac{d^{2}v}{dt^{2}} \bigg/ \left(\frac{dv}{dt}\right)^{2}\]
und erhalten f\"ur $k$ die Differentialgleichung
\begin{equation}\tag{15}
\frac{dv}{dt} \frac{d^{2}k}{dt^{2}} - \frac{d^{2}v}{dt^{2}} \frac{dk}{dt} = \left(\frac{dv}{dt}\right)^{3} \left(\frac{1}{k_{1}} \frac{d^{2}k_{1}}{dv^{2}}\right) k.
\end{equation}
Von der Gleichung (15) seien zwei von einander unabh\"angige particul\"are Integrale $K_{1}$ und $K_{2}$ gefunden, deren Quotient $K_{2}:K_{1} = H$

\newpage
\noindent ein von Bogen gr\"osster Kreise begrenztes Abbild der positiven $t$-Halbebene auf der Kugelfl\"ache liefert. Dasselbe leistet dann jeder Ausdruck von der Form
\begin{equation}\tag{16}
\eta = e^{\vartheta i} \frac{K_{2} - \alpha K_{1}}{K_{2} - \alpha' K_{1}},
\end{equation}
worin $\vartheta$ reell und $\alpha$, $\alpha'$ conjugirte complexe Gr\"ossen sind.

Die Function $v$ ist so zu w\"ahlen, dass f\"ur endliche Werthe von $t$ die Unstetigkeiten von $\frac{1}{k} \frac{d^{2}k}{dv^{2}}$ nicht ausserhalb der Punkte $a$, $a'$, $b$, $c$, $e$ liegen.

Setzt man
\begin{equation}\tag{17}
\frac{dv}{dt} = \frac{1}{\sqrt{\psi(t)\chi(t)}} = \frac{1}{\Pi(t)},
\end{equation}
so wird die Function $\frac{1}{k} \frac{d^{2}k}{dv^{2}}$ im Endlichen unstetig nur f\"ur die Punkte $a$, $a'$, $b$, $c$, und zwar f\"ur jeden unendlich in erster Ordnung. Man erh\"alt n\"amlich f\"ur $t = c$
\[\frac{1}{k} \frac{d^{2}k}{dv^{2}} = \frac{\text{const.}}{t-c}.\]
Folglich:
\[k = \text{const.} \cdot (t-c)^{\frac{1}{2}}, \quad \eta = \text{const.} \cdot (t-c),\]
und hieraus:
\[\eta - \eta_{c} = \text{const.} \cdot (t-c)^{\frac{\gamma}{2}}.\]
F\"ur $t = c$:
\[k \frac{dv^{2}}{dt^{2}} : \frac{1}{t-c}.\]
Entsprechende Ausdr\"ucke erh\"alt man f\"ur $t = a$, $a'$, $b$, in denen $c$ resp.\ durch $a$, $a'$, $b$, und $\gamma$ durch $2$ zu ersetzen ist.

Eine \"ahnliche Betrachtung lehrt, dass f\"ur $t = e$ die Function $\frac{1}{k} \frac{d^{2}k}{dv^{2}}$ stetig bleibt.

F\"ur $t = \infty$ ergiebt sich
\[\frac{1}{k} \frac{d^{2}k}{dv^{2}} = \frac{\text{const.}}{t^{2}}.\]
\[\frac{1}{k} \frac{d^{2}k}{dv^{2}} = \sum \frac{\frac{1}{2}(\frac{\gamma}{2}-1)}{t-g} + \frac{F(t)}{\psi(t)\chi(t)}.\]
Die Summe bezieht sich auf alle Punkte $g = a$, $a'$, $b$, $c$, und bei $a$, $a'$, $b$ ist $2$ statt $\gamma$ zu setzen. $F(t)$ ist eine ganze Function vom Grade $(2\nu - 6)$, in der die ersten beiden Coefficienten sich folgendermassen bestimmen. Man bringe $dv$ in die Form

\newpage
\noindent $dv = \sqrt{a}\, dv'$

oder k\"urzer $dv^{2} = a\, dv'^{2}$.

Dann ergiebt sich durch Differentiation
\[2\, dv \frac{d^{2}v}{dt^{2}} = \frac{da}{dt}\, dv'^{2} + 2a\, dv' \frac{d^{2}v'}{dt^{2}},\]
folglich
\[\frac{1}{dv} \frac{d^{2}v}{dt^{2}} = \frac{1}{2} \frac{da}{dt}\frac{1}{a} + \frac{dv'}{dv} \frac{d^{2}v'}{dt^{2}}.\]
\[\left(\frac{d^{2}v}{dv^{2}}\right)^{2} \frac{d}{dv'} \left\{\left(\frac{dv}{dv'}\right)^{2} \frac{d^{2}v}{dt^{2}}\right\} = a^{2} \frac{d^{2}v}{dv^{2}}\]
\[= \sum \frac{\frac{1}{2}(\frac{\gamma}{2}-1)}{t-g} + \frac{F(t)}{\psi(t)\chi(t)}.\]
Die Function auf der linken Seite ist endlich f\"ur $t = \infty$. Folglich hat man rechts in der Entwicklung von $\frac{1}{\psi(t)\chi(t)} F(t)$ und von $\frac{dv}{dt}$ die Coefficienten von $t^{2}$ und resp.\ von $t$ einander gleich zu setzen. Die Entwickelung von $\frac{dv}{dt}$ giebt nach einfacher Rechnung

Hiernach bleiben in $F(t)$ noch $2\nu - 7$ unbestimmte Coefficienten. Es ist aber wichtig zu bemerken, dass dieselben reell sein m\"ussen. Denn wir haben in \S.~12 gefunden, dass $du$ reell oder rein imagin\"ar ist in allen geraden Begrenzungslinien der Minimalfl\"ache und folglich auch an jeder Stelle in der Begrenzung der Abbildungen. Verm\"oge der Gleichung (17) gilt dasselbe von $dv$. Daraus l\"asst sich beweisen, dass f\"ur reelle Werthe von $t$ die Function $\frac{1}{k} \frac{d^{2}k}{dv^{2}}$ nothwendigerweise reelle Werthe besitzt.

Um diesen Beweis zu f\"uhren, betrachten wir die Abbildung auf der Kugel vom Radius $1$ und nehmen irgend einen Theil der Begrenzung, also den Bogen eines gewissen gr\"ossten Kreises. Im Pole dieses gr\"ossten Kreises legen wir die Tangential-Ebene an und bezeichnen

\newpage
\noindent sie als die Ebene der $\eta_{1}$. Dann lassen sich die constanten Gr\"ossen $\alpha_{1}$, $\alpha'_{1}$, $\vartheta_{1}$ so bestimmen, dass
\[\eta_{1} = e^{\vartheta_{1}i} \frac{\eta - \alpha_{1}}{1 + \alpha'_{1}\eta}\]
ist, und wir erhalten zwei Functionen $k'_{1} = \sqrt{\frac{dv}{d\eta_{1}}}$ und $k'_{2} = \eta_{1}\sqrt{\frac{dv}{d\eta_{1}}}$, die particul\"are Integrale der Differentialgleichung (15) sind. Folglich haben wir
\[\frac{1}{k_{1}} \frac{d^{2}k_{1}}{dv^{2}} = \frac{1}{k'_{1}} \frac{d^{2}k'_{1}}{dv^{2}}.\]
Der eben betrachtete Theil der Begrenzung bildet sich in der $\eta_{1}$-Ebene ab durch die Gleichung
\[\eta_{1} + \eta'_{1} = 0\]
und wenn man dies in $k'_{1}$ einf\"uhrt, so erkennt man leicht, dass in dem fraglichen Begrenzungstheile $\frac{1}{k'_{1}} \frac{d^{2}k'_{1}}{dv^{2}}$ reell ausf\"allt. Folglich gilt dasselbe von $\frac{1}{k_{1}} \frac{d^{2}k_{1}}{dv^{2}}$, und da diese Betrachtung f\"ur jedes einzelne Begrenzungsst\"uck angestellt werden kann, so ist $\frac{1}{k} \frac{d^{2}k}{dv^{2}}$ reell in der ganzen Begrenzung.

Nun f\"allt aber bei einem reellen oder rein imagin\"aren $dv$ die Function $\frac{1}{k_{1}} \frac{d^{2}k_{1}}{dv^{2}}$ auch dann reell aus, wenn man allgemeiner
\[k'_{1} = \varrho_{1} e^{\psi_{1}i}\]
setzt und den Modul $\varrho_{1}$ constant nimmt. Damit also die Axe der reellen $t$ sich auf der Kugel vom Radius $1$ wirklich in B\"ogen gr\"osster Kreise abbilde, muss f\"ur jeden Begrenzungstheil $\varrho_{1} = 1$ sein. Dies liefert ebenso viele Bedingungsgleichungen, als einzelne Begrenzungslinien gegeben sind.

Bei dieser Untersuchung ist, wie schon im vorigen Paragraphen, vorausgesetzt, dass die Werthe $a$, $b$, $c$, $e$ s\"ammtlich endlich seien. Trifft dies nicht zu, so bedarf die Betrachtung einer geringen Modification.

\medskip
\noindent\textbf{Anmerkung.} Die Aufgabe ist hiermit vollst\"andig formulirt. Im einzelnen Falle kommt es nur darauf an, die Differentialgleichung (15) wirklich aufzustellen und zu integriren. Uebrigens ist es nicht unwichtig, zu bemerken, dass die Anzahl der in der L\"osung auftretenden willk\"urlichen reellen Constanten ebenso gross ist wie die Anzahl der Bedingungsgleichungen, welche verm\"oge der Natur der Aufgabe und verm\"oge der Daten des Problems erf\"ullt sein m\"ussen. Wir bezeichnen die Anzahl der Punkte $a$, $b$, $c$, $e$ resp.\ mit $A$, $B$, $C$, $E$ und beachten, dass $2A + B + 4 = C + 2E = \nu$ ist. In der Differentialgleichung (15) treten $2A + B + 4C + 5E - 10$ willk\"urliche reelle Constanten auf, n\"amlich: die

\newpage
\noindent Winkel $\gamma$, deren Anzahl $C$ ist; die $2\nu - 7$ Constanten der Function $F(t)$; die reellen Gr\"ossen $b$, $c$, $e$, von denen man dreien beliebige Werthe geben kann, indem man f\"ur $t$ eine lineare Substitution mit reellen Coefficienten macht; die reellen und imagin\"aren Theile der Gr\"ossen $a$. Zu diesen willk\"urlichen Constanten kommen bei der Integration noch $10$ hinzu, n\"amlich, wenn $\eta = \frac{K_{2} - \alpha K_{1}}{K_{2} - \alpha' K_{1}}$ ist, die drei complexen Verh\"altnisse $\alpha:\beta:\gamma:\delta$, die f\"ur $6$ reelle Constanten zu z\"ahlen sind, ein (reeller oder rein imagin\"arer) Factor von $du$, und je eine additive reelle Constante in den Ausdr\"ucken f\"ur $x$, $y$, $z$. Diese Constanten m\"ussen aber noch Bedingungsgleichungen unterworfen werden, die erf\"ullt sein m\"ussen, wenn unsere Formeln wirklich eine Minimalfl\"ache darstellen sollen. Von diesen Bedingungsgleichungen sind $2A + B$ den Punkten $a$, $a'$, $b$ entsprechend, die aussagen, dass in den in der Umgebung dieser Punkte g\"ultigen Entwicklungen der L\"osungen der Differentialgleichung (15) keine Logarithmen auftreten (vergl.\ Anmerkung (4)), und $C + E$, die besagen, dass die zwischen den einzelnen Punkten $c$, $e$ gelegenen St\"ucke der Axe der reellen $t$ sich auf der Kugel mit dem Radius $1$ in $C + E$ B\"ogen gr\"osster Kreise abbilden. Sonach ist die Anzahl der in der L\"osung uebrig bleibenden unbestimmten Constanten $3C + 4E$.

Die Daten des Problems bestehen in den Coordinaten der Eckpunkte, und in den Winkeln, die die Richtungen der ins Unendliche verlaufenden Begrenzungslinien festlegen. Diese Daten sprechen sich in $3C + 4E$ Gleichungen aus, zu deren Erf\"ullung man eine ebenso grosse Zahl verf\"ugbarer Constanten hat.\textsuperscript{(*)})

\bigskip
\noindent\textbf{Beispiele.}

\subsection*{15.}

Die Begrenzung bestehe aus zwei unendlichen geraden Linien, die nicht in einer Ebene liegen. Ihre k\"urzeste Verbindungslinie habe die L\"ange $A$, und es sei $\alpha\pi$ der Winkel, welchen die Projection der Fl\"ache auf der rechtwinklig gegen jene Verbindungslinie gelegten Ebene ausf\"ullt.

Nimmt man die k\"urzeste Verbindungslinie zur $x$-Axe, so hat in jeder der beiden Begrenzungsgeraden $x$ einen constanten Werth. Ebenso ist $\varphi$ in jeder der beiden Begrenzungsgeraden constant. In unendlicher Entfernung ist die positive Normale f\"ur den einen Sector parallel der positiven, f\"ur den andern Sector parallel der negativen $x$-Axe. Die Begrenzung bildet sich auf der Kugel in zwei gr\"ossten Kreisen ab, die durch die Pole $\eta = 0$ und $\eta = \infty$ gehen und den Winkel $\alpha\pi$ einschliessen.

Hiernach hat man
\[x = \frac{iA}{2\alpha\pi} \log\eta,\]
\[s = \frac{iA}{2\alpha\pi\sqrt{\eta}} \left(\frac{1}{\eta} - \eta'\right),\]
\[s' = \frac{iA}{2\alpha\pi\sqrt{\eta'}} \left(\frac{1}{\eta'} - \eta\right),\]
folglich
\[x = -\frac{iA}{2\alpha\pi} \log(\eta\eta'),\]
\[y = \frac{A}{2\alpha\pi} \cdot \frac{1}{2i} \log\frac{s'}{s},\]
worin man die Gleichung der Schraubenfl\"ache erkennt.

\subsection*{16.}

Die Begrenzung bestehe aus drei geraden Linien, von denen zwei sich schneiden und die dritte zur Ebene der beiden ersten parallel l\"auft. Legt man den Anfangspunkt der Coordinaten in den Schnittpunkt der beiden ersten Geraden, die positive $x$-Axe in die negative Normale, so bildet jener Schnittpunkt auf der Kugel sich ab im Punkte $\eta = \infty$. Die Abbildung der beiden ersten Geraden sind gr\"osste Halbkreise, die von $\eta = \infty$ bis $\eta = 0$ laufen. Ihr Winkel sei $\alpha\pi$. Die Abbildung der dritten Linie ist der Bogen eines gr\"ossten Kreises, der von $\eta = 0$ ausgeht, an einer gewissen Stelle umkehrt und in sich selbst bis zum Punkte $\eta = 0$ zur\"uckl\"auft. Dieser Bogen bilde mit den beiden ersten gr\"ossten Halbkreisen die Winkel $-\beta\pi$ und $\gamma\pi$, so dass $\beta$ und $\gamma$ absolute Zahlen sind und $\beta + \gamma = \alpha$ sich ergiebt. Um die Abbildung auf der halben $t$-Ebene zu erhalten, setzen wir fest, dass $t = \infty$ sein soll f\"ur $\eta = \infty$, dass dem unendlichen Sector zwischen der ersten und dritten Linie $t = b$, dem unendlichen Sector zwischen der zweiten und dritten Linie $t = c$, dem Umkehrpunkte der Normalen auf der dritten Linie $t = a$ entsprechen soll. Dabei sind $a$, $b$, $c$ reell und $c > a > b$.

Diesen Bestimmungen entspricht $\eta = (t-b)^{\beta}(t-c)^{\gamma}$. Der Werth $a$ h\"angt von $b$ und $c$ ab. Man hat n\"amlich
\[\frac{d\log\eta}{dt} = \frac{\beta(t-c) + \gamma(t-b)}{(t-b)(t-c)}\]
und dieses muss f\"ur den Umkehrpunkt $=0$ sein, also $a = \frac{\beta c + \gamma b}{\beta + \gamma}$.

Man hat weiter nach Art.~12 und 13
\[\frac{du}{dt} = i\sqrt{(c-b)(\beta+\gamma)} \frac{(t-a)\, dt}{(t-b)(t-c)},\]
oder wenn man $c - b = -r$ annimmt
\[\frac{du}{dt} = i\sqrt{(c-b)(\beta+\gamma)} \frac{(t-a)\, dt}{(t-b)(t-c)}.\]

\newpage
Folglich
\[\text{(b)} \quad X = i\sqrt{\beta+\gamma} \int \frac{(t-a)\, dt}{\sqrt{(t-b)(t-c)}},\]
\[\frac{ds}{dt} = i\sqrt{\beta+\gamma} \cdot \frac{1}{\eta^{2}} \frac{(t-a)\, dt}{(t-b)(t-c)},\]
\[\frac{ds'}{dt} = i\sqrt{\beta+\gamma} \cdot \eta^{2} \frac{(t-a)\, dt}{(t-b)(t-c)},\]
\[s = \frac{i}{2} \int \frac{(t-b)^{\beta}(t-c)^{\gamma} - (t-b)^{-\beta}(t-c)^{-\gamma}}{(t-b)(t-c)}\, (t-a)\, dt,\]
\[\quad + \frac{i}{2} \int \frac{(t'-b)^{\beta}(t'-c)^{\gamma} - (t'-b)^{-\beta}(t'-c)^{-\gamma}}{(t'-b)(t'-c)}\, (t'-a)\, dt',\]
\[\frac{(t-b)^{\beta}(t-c)^{\gamma} + (t-b)^{-\beta}(t-c)^{-\gamma}}{(t-b)(t-c)}\, (t-a)\, dt.\]

\subsection*{17.}

Die Begrenzung bestehe aus drei einander kreuzenden geraden Linien, deren k\"urzeste Abst\"ande $A$, $B$, $C$ sein m\"ogen. Zwischen je zwei begrenzenden Linien erstreckt sich die Fl\"ache ins Unendliche. Es seien $\alpha\pi$, $\beta\pi$, $\gamma\pi$ die Winkel der Richtungen, in welchen die Grenzlinien des ersten, des zweiten, des dritten Sectors ins Unendliche verlaufen. Setzt man fest, dass f\"ur die drei Sectoren der Minimalfl\"ache im Unendlichen die Gr\"osse $t$ resp.\ $= 0$, $\infty$, $1$ sein soll, so erh\"alt man
\[\frac{du}{dt} = \frac{\sqrt{\varphi(t)}}{t(1-t)},\]
$\varphi(t)$ ist eine ganze Function zweiten Grades. Ihre Coefficienten bestimmen sich daraus, dass
\[\text{f\"ur } t = 0: \quad \frac{du}{d\log t} = \frac{A}{2\pi},\]
\[\text{f\"ur } t = \infty: \quad \frac{du}{d\log t} = \frac{B}{2\pi},\]
\[\text{f\"ur } t = 1: \quad \frac{du}{d\log(1-t)} = \frac{C}{2\pi}\]
sein muss.

Danach ergiebt sich
\[\varphi(t) = \frac{1}{4\pi^{2}} \left\{A^{2} + (B^{2} - A^{2} - C^{2})t + C^{2}t^{2}\right\}.\]
Je nachdem die Wurzeln der Gleichung $\varphi(t) = 0$ imagin\"ar oder reell sind, hat die Abbildung auf der Kugel einen Verzweigungspunkt im Innern oder zwei Umkehrpunkte der Normalen auf der Begrenzung.

\newpage
Die Functionen $k_{1} = \sqrt{\frac{dv}{d\eta}}$ und $k_{2} = \eta\sqrt{\frac{dv}{d\eta}}$ werden nur f\"ur die drei Sectoren unstetig, wenn man $\frac{dv}{dt} = \varphi(t)$ nimmt, und zwar ist die Unstetigkeit von $k_{1}$ der Art, dass
\[\text{f\"ur } t = 0: \quad k_{1} = t^{-\frac{\alpha}{2}},\]
\[\text{f\"ur } t = \infty: \quad k_{1} = t^{\frac{\beta}{2}},\]
\[\text{f\"ur } t = 1: \quad k_{1} = (1-t)^{\frac{\gamma}{2}}\]
ein\"andrig und verschieden von $0$ und $\infty$ wird. $k_{1}$ und $k_{2}$ sind particul\"are Integrale einer homogenen linearen Differentialgleichung zweiter Ordnung, die sich ergiebt, wenn man $\frac{1}{k} \frac{d^{2}k}{dv^{2}}$ aus seinen Unstetigkeiten als Function von $t$ darstellt und $t$ statt $v$ als unabh\"angige Variable in $\frac{d^{2}k}{dv^{2}}$ einf\"uhrt. Hat man das particul\"are Integral $k_{1}$ gefunden, so ergiebt sich $k_{2}$ aus der Differentialgleichung erster Ordnung
\[k_{2} \frac{dk_{1}}{dv} - k_{1} \frac{dk_{2}}{dv} = 1.\]
Das vollst\"andige Integral der homogenen linearen Differentialgleichung zweiter Ordnung werde mit
\[\text{(d)} \quad k = P \begin{pmatrix} 0 & \infty & 1 \\ -\frac{\alpha}{2} & +\frac{\beta}{2} & \frac{\gamma}{2} & t \\ \ast & \ast & \ast \end{pmatrix}\]
bezeichnet. Diese Function gen\"ugt wesentlich denselben Bedingungen, die in der Abhandlung ueber die Gauss'sche Reihe $F(\alpha, \beta, \gamma, x)$ als Definition der $P$-Function ausgesprochen sind\textsuperscript{*}). Sie weicht von der $P$-Function darin ab, dass die Summe der Exponenten $-1$ ist, nicht $+1$ wie bei $P$.

Man kann die Function $Q$ mit H\"ulfe einer Function $P$ und ihrer ersten Derivirten ausdr\"ucken. Zun\"achst ist n\"amlich
\[k_{1} = P \begin{pmatrix} 0 & \infty & 1 & \\ -\frac{\alpha}{2} & +\frac{\beta}{2} & \frac{\gamma}{2} & t \\ \frac{\alpha}{2} & -\frac{\beta}{2}-\gamma-1 & -\frac{\gamma}{2} \end{pmatrix},\]
\[Q = P \begin{pmatrix} 0 & \infty & 1 \\ 0 & -\alpha-\beta-\gamma+1 & 0 & t \\ \alpha & -\alpha+\beta-\gamma+1 & \gamma \end{pmatrix}.\]

\newpage
Setzt man nun
\[k_{1} = t^{-\frac{\alpha}{2}} (1-t)^{\frac{\gamma}{2}} \vartheta,\]
\[Q = P \begin{pmatrix} 0 & \infty & 1 \\ 0 & -\alpha-\beta-\gamma+1 & 0 & t \\ \alpha & -\alpha+\beta-\gamma+1 & \gamma \end{pmatrix},\]
so lassen sich die Constanten $a$, $b$, $c$ so bestimmen, dass
\[\text{(e)} \quad k_{1} = t^{-\frac{\alpha}{2}} (1-t)^{\frac{\gamma}{2}} \left((a+bt)\vartheta + ct(1-t) \frac{d\vartheta}{dt}\right)\]
wird. In der That hat man nur diesen Ausdruck in die Differentialgleichung (c) einzusetzen und die Differentialgleichung zweiter Ordnung f\"ur $\vartheta$ zu beachten, um zu der Gleichung zu gelangen
\[\varphi(t) = a^{2} t^{-\alpha} (1-t)^{-\gamma} (\vartheta, \vartheta)\, dt,\]
\[F(t) = a(a+\alpha)(1-t) + (a+b)(a+b-\gamma c)t\]
\[= (1-t)\left(1 - \frac{a+b+c}{a}t\right)\left(1 - \frac{a+b-c}{a}t\right).\]
Verm\"oge der Eigenschaften der Function $\vartheta$ kann man setzen
\[t^{-\alpha}(1-t)^{-\gamma}(\vartheta, \vartheta) = 1.\]
und folglich muss $F(t) = \varphi(t)$ sein. Hieraus ergeben sich drei Bedingungsgleichungen f\"ur $a$, $b$, $c$, die eine sehr einfache Form annehmen, wenn man
\[a+b+c = p, \quad a+b-c = q, \quad \frac{a+b+c}{a} = r\]
setzt. Die Bedingungsgleichungen lauten dann
\[pp - \alpha\alpha(p+q+r)^{2} = \frac{A^{2}}{2\pi^{2}},\]
\[qq - \beta\beta(p+q+r)^{2} = \frac{B^{2}}{2\pi^{2}},\]
\[rr - \gamma\gamma(p+q+r)^{2} = \frac{C^{2}}{2\pi^{2}}.\]
Mit H\"ulfe der Function
\[\varDelta = P \begin{pmatrix} 0 & \infty & 1 \\ \frac{\alpha}{2} & \frac{\beta}{2} & \frac{\gamma}{2} & t \\ -\frac{\alpha}{2} & -\frac{\beta}{2}+2 & -\frac{\gamma}{2}+2 \end{pmatrix},\]
deren Zweige $\varDelta_{1}$ und $\varDelta_{2}$ der Differentialgleichung gen\"ugen

\newpage
\noindent den Coordinatenaxen parallel laufen, ist $\alpha = \beta = \gamma = \frac{1}{2}$. Dann erh\"alt man
\[\text{(f)} \quad k = t^{\frac{1}{4}} \left((p+qt)^{2} + ct(1-t)^{2}\right).\]
Es w\"urde nicht schwer sein, die einzelnen Zweige der Function $k$ in der Form von bestimmten Integralen herzustellen. Der Weg dazu ist in Art.~7 der Abhandlung ueber die Function $P$ vorgezeichnet.

In dem besondern Falle, dass die drei begrenzenden geraden Linien den Coordinatenaxen parallel laufen, ist $\alpha = \beta = \gamma = \frac{1}{2}$. Dann erh\"alt man

Der Zweig $K_{1}$ dieser Function ist
\[= \left(\frac{c}{p}\right)^{\frac{1}{2}} \sqrt{t^{2} + (q-1)^{2}} \cdot \text{const.},\]
und daraus ergiebt sich
\[\begin{aligned}
k_{1} &= \sqrt{2} t^{\frac{1}{4}} (t-1)^{\frac{1}{4}} \sqrt{t^{2} + (q-1)^{2}} \left[p + qt + i\sqrt{c}\, t\sqrt{t^{2}+(q-1)^{2}}\right], \\
k_{2} &= \sqrt{2} t^{\frac{1}{4}} (t-1)^{\frac{1}{4}} \sqrt{t^{2} + (q-1)^{2}} \left[p + qt - i\sqrt{c}\, t\sqrt{t^{2}+(q-1)^{2}}\right].
\end{aligned}\]

Mit H\"ulfe dieser beiden Functionen lassen sich $dX$, $dY$, $dZ$ folgendermassen ausdr\"ucken
\[dX = -i\, k_{1}k_{2}\, \frac{du}{d\log\eta}\, dt,\]
\[dZ = -\frac{i}{2} (k_{1}^{2} + k_{2}^{2}) \frac{du}{d\log\eta}\, dt,\]
\[\text{(g)} \quad iY = -(p-q+r)\sqrt{t} + (p+q-r)\sqrt{t^{2}}\]
\[\quad - i(p+3q+3r)(p+q-r)\log\sqrt{t+1},\]
\[iZ = (p-q+r)^{2}(1-t)^{\frac{1}{2}} + (p+q-r)^{2}(1-t)^{-\frac{1}{2}}\]
\[\quad + i(3p+q+r)(-p+q+r)\log\sqrt{1+t^{2}}.\]

\newpage
Wenn $p$, $q$, $r$ reell sind, so geben die doppelten Coefficienten von $i$ in den drei Gr\"ossen rechts die rechtwinkligen Coordinaten eines Punktes der Fl\"ache.

\subsection*{18.}

Die Begrenzung bestehe aus vier sich schneidenden geraden Linien, die man erh\"alt\textsuperscript{**}) wenn von den Kanten eines beliebigen Tetraeders zwei nicht zusammenstossende weggelassen werden. Die Abbildung auf der Kugeloberfl\"ache ist ein sph\"arisches Viereck, dessen Winkel $\alpha\pi$, $\beta\pi$, $\gamma\pi$, $\delta\pi$ sein m\"ogen. Es ergiebt sich
\[\frac{du}{dt} = \frac{G\, dt}{\sqrt{(t-a)(t-b)(t-c)(t-d)}} = \frac{G\, dt}{\sqrt{\varDelta(t)}},\]
wenn die reellen Werthe $t = a$, $b$, $c$, $d$ die Punkte der $t$-Ebene bezeichnen, in welchen sich die Eckpunkte des Vierecks abbilden.

Soll die in \S.~14 entwickelte Methode zur Bestimmung von $\eta$ angewandt werden, so hat man hier speciell $\varphi(t) = 1$, $\chi(t) = \varDelta(t)$, folglich $\nu = 4$ und
\[\frac{dv}{dt} = \frac{1}{\sqrt{\varDelta(t)}}.\]
Die Functionen $k_{1}$ und $k_{2}$ gen\"ugen der Differentialgleichung
\[k_{2} \frac{dk_{1}}{dv} - k_{1} \frac{dk_{2}}{dv} = 1\]
und sind particul\"are Integrale der Differentialgleichung zweiter Ordnung
\[\frac{1}{k} \frac{d^{2}k}{dv^{2}} = \frac{(\alpha\alpha-1)\varDelta'(a)}{4(t-a)} + \frac{(\beta\beta-1)\varDelta'(b)}{4(t-b)}\]
\[+ \frac{(\gamma\gamma-1)\varDelta'(c)}{4(t-c)} + \frac{(\delta\delta-1)\varDelta'(d)}{4(t-d)}.\]
Die Function $F(t)$ des \S.~14 ist hier vom zweiten Grade, aber die Coefficienten von $t^{2}$ und von $t$ sind gleich Null, also $h$ eine Constante.

In der letzten Gleichung hat man auf der linken Seite $t$ als unabh\"angige Variable einzuf\"uhren und erh\"alt
\[\frac{1}{k} \frac{d^{2}k}{dt^{2}} = \frac{(\alpha\alpha-1)\varDelta'(a)}{t-a} + \frac{(\beta\beta-1)\varDelta'(b)}{t-b}\]
\[+ \frac{(\gamma\gamma-1)\varDelta'(c)}{t-c} + \frac{(\delta\delta-1)\varDelta'(d)}{t-d}\]
als die Differentialgleichung zweiter Ordnung, welcher $k$ Gen\"uge leisten muss.

Sind $X$, $Y$, $Z$ als Functionen von $t$ wirklich ausgedr\"uckt, so treten in der L\"osung noch $16$ unbestimmte reelle Constanten auf, n\"amlich die vier Gr\"ossen $a$, $b$, $c$, $d$, von denen wie oben drei beliebig angenommen werden k\"onnen, die vier Gr\"ossen $\alpha$, $\beta$, $\gamma$, $\delta$, die Gr\"osse $h$,

\newpage
\noindent ferner $6$ reelle Constanten in dem Ausdrucke f\"ur $\eta$, ein constanter Factor in $du$ und je eine additive Constante in $x$, $y$, $z$. Zur Bestimmung dieser $16$ Gr\"ossen sind $16$ Bedingungsgleichungen vorhanden, n\"amlich $4$ Gleichungen, welche ausdr\"ucken, dass die vier Begrenzungslinien in der Ebene der $\eta$ sich auf der Kugel in gr\"ossten Kreisen abbilden, und $12$ Gleichungen, welche aussagen, dass $x$, $y$, $z$ in den $4$ Eckpunkten gegebene Werthe haben.

In dem speciellen Falle eines regul\"aren Tetraeders ist die Abbildung auf der Kugel ein regelm\"assiges Viereck, in welchem jeder Winkel $= \frac{2}{3}\pi$. Die Diagonalen halbiren sich und stehen rechtwinklig aufeinander. Die den Eckpunkten diametral gegen\"uberliegenden Punkte der Kugeloberfl\"ache sind die Ecken eines congruenten Vierecks. Zwischen beiden liegen vier dem urspr\"unglichen ebenfalls congruente Vierecke, die je zwei Eckpunkte mit dem urspr\"unglichen, zwei mit dem gegen\"uberliegenden gemein haben. Diese sechs Vierecke f\"ullen die Kugeloberfl\"ache einfach aus. Es wird also $\frac{du}{d\log\eta}$ eine algebraische Function von $\eta$ sein.

Man kann die gesuchte Minimalfl\"ache ueber ihre urspr\"ungliche Begrenzung dadurch stetig fortsetzen, dass man sie um jede ihrer Grenzlinien als Drehungsaxe um $180^{\circ}$ dreht. L\"angs einer solchen Grenzlinie haben dann die urspr\"ungliche Fl\"ache und die Fortsetzung gemeinschaftliche Normalen. Wiederholt man die Construction an den neuen Fl\"achentheilen, so l\"asst sich die urspr\"ungliche Fl\"ache beliebig weit fortsetzen. Welche Fortsetzung man aber auch betrachte, immer bildet sie sich auf der Kugel in einem der sechs congruenten Vierecke ab. Und zwar haben die Abbildungen von zwei Fl\"achentheilen eine Seite gemein oder sie liegen einander gegen\"uber, je nachdem die Fl\"achentheile selbst in einer Grenzlinie an einander stossen oder an gegen\"uberliegenden Grenzlinien eines mittleren Fl\"achentheils gelegen sind. In dem letzteren Falle k\"onnen die betreffenden Fl\"achentheile durch parallele Verschiebung zur Deckung gebracht werden. Daher muss $\left(\frac{du}{d\log\eta}\right)^{2}$ unver\"andert bleiben, wenn $\eta$ mit $\frac{1}{\eta'}$ vertauscht wird.

Legt man den Pol ($\eta = 0$) in den Mittelpunkt eines Vierecks, den Anfangsmeridian durch die Mitte einer Seite, so ist f\"ur die Eckpunkte dieses Vierecks
\[\eta = \tg\frac{\pi}{8} \cdot e^{\frac{\pi i}{4}}, \quad \tg\frac{\pi}{8} \cdot e^{\frac{3\pi i}{4}}, \quad \tg\frac{\pi}{8} \cdot e^{-\frac{3\pi i}{4}}, \quad \tg\frac{\pi}{8} \cdot e^{-\frac{\pi i}{4}},\]
und f\"ur die Eckpunkte des gegen\"uberliegenden Vierecks die entgegengesetzten Werthe.

\newpage
Punkte, denen entgegengesetzte Werthe von $\eta$ angeh\"oren, haben dieselbe $X$-Coordinate. Es muss also $\left(\frac{du}{d\log\eta}\right)^{2}$ bei der Vertauschung von $\eta$ mit $-\eta$ unver\"andert bleiben. Hiernach erh\"alt man
\[\left(\frac{du}{d\log\eta}\right)^{2} = C_{1}^{2} \frac{(\eta^{4}+\eta'^{4}+1)}{(\eta^{4}\eta'^{4}+\eta^{4}+\eta'^{4})}.\]
Die Constante $C_{1}^{2}$ muss reell sein, damit $du^{2}$ in der Begrenzung reelle Werthe besitze.

Zu demselben Resultate gelangt man auf dem folgenden Wege. Die Substitution
\[\xi^{2} = \frac{\eta^{4}+\eta'^{4}+1}{\eta^{4}\eta'^{4}+\eta^{4}+\eta'^{4}}\]
liefert auf der $\xi$-Ebene eine Abbildung, die von einer geschlossenen ueberall stetig gekr\"ummten Linie begrenzt wird. Die Rechnung zeigt, dass $d\log\xi$ in der Begrenzung rein imagin\"ar ist. Folglich ist die Abbildung der Begrenzung in der $\xi$-Ebene ein Kreis um den Mittelpunkt $\xi = 0$. Der Radius dieses Kreises ist $=1$. Den Eckpunkten
\[\eta = \tg\frac{\pi}{8} \cdot e^{\pm\frac{\pi i}{4}}\]
entspricht $\xi = +1$, den Eckpunkten
\[\eta = \tg\frac{\pi}{8} \cdot e^{\pm\frac{3\pi i}{4}}\]
entspricht $\xi = i$. Geht man an irgend einer dieser vier Stellen durch das Innere der Minimalfl\"ache von einer Grenzlinie zur folgenden, so \"andert sich dabei der Arcus von $d\xi$ um $\pi$. Daher kann man, wie in \S.~13, auch hier setzen
\[\frac{du}{dt} = \frac{C_{2}}{\sqrt{(t-a)(t-b)(t-c)}},\]
und es muss $C_{2}$ rein imagin\"ar sein, damit $du^{2}$ in der Begrenzung reell ausfalle. Es findet sich $C_{2} = 3\sqrt[4]{C_{1}^{2}} \cdot i$.

Dieser Ausdruck stimmt mit dem vorher aufgestellten f\"ur $\left(\frac{du}{d\log\eta}\right)^{2}$.

Zur weitern Vereinfachung nehme man
\[\eta^{2} = \frac{\sqrt{t^{2}+1} + \sqrt{2}\, t}{\sqrt{t^{2}+1} - \sqrt{2}\, t}\]
und beachte, dass
\[\left(\frac{du}{d\log\eta}\right)^{2} = 2C_{1}^{2} \frac{dt}{d\log\eta}.\]
Dann ergiebt eine sehr einfache Rechnung

\newpage
\begin{equation}\tag{h}
\frac{1}{2} \left(\frac{du}{d\log\eta}\right)^{2} = C \int \frac{dt}{\sqrt{1-t^{4}}},
\end{equation}
\[\frac{1}{2} \left(\frac{du}{d\log\eta}\right)^{2} = C \int \frac{dt}{\sqrt{(1-\varrho t^{2})(1-\varrho^{2}t^{2})}},\]
wenn $\varrho = -\frac{1}{2}(1-i\sqrt{3})$ eine dritte Wurzel der Einheit bezeichnet.

Die reelle Constante $C = \sqrt[3]{6} C_{1}^{2}$ bestimmt sich aus der gegebenen L\"ange der Tetraederkanten.

\subsection*{19.}

Endlich soll noch die Aufgabe der Minimalfl\"ache f\"ur den Fall behandelt werden, dass die Begrenzung aus zwei beliebigen Kreisen besteht, die in parallelen Ebenen liegen. Dann kennt man die Richtung der Normalen in der Begrenzung nicht. Daher l\"asst sich diese auch nicht auf der Kugel abbilden. Man gelangt aber zur L\"osung der Aufgabe durch die Annahme, dass alle zu den Ebenen der Grenzkreise parallel gelegten ebenen Schnitte Kreise seien, und es wird sich zeigen, dass unter dieser Annahme der Minimalbedingung Gen\"uge geleistet werden kann.

Legt man die $x$-Axe rechtwinklig gegen die Ebenen der Grenzkreise, so ist die Gleichung der Schnittcurve in einer parallelen Ebene
\[y^{2} + z^{2} + 2\alpha y + 2\beta z + \gamma = 0,\]
und $\alpha$, $\beta$, $\gamma$ sind als Functionen von $x$ zu bestimmen. Zur Abk\"urzung werde
\[F = y^{2} + z^{2} + 2\alpha y + 2\beta z + \gamma\]
gesetzt, so dass
\[\cos r = w \frac{\partial F}{\partial x}, \quad \sin r \cos\varphi = w \frac{\partial F}{\partial y}, \quad \sin r \sin\varphi = w \frac{\partial F}{\partial z}\]
ist. Dann l\"asst sich die Bedingung des Minimum in die Form bringen
\[\frac{\partial}{\partial x} \left(\frac{1}{w}\right) + \frac{\partial}{\partial y} \left(\frac{\sin r \cos\varphi}{w}\right) + \frac{\partial}{\partial z} \left(\frac{\sin r \sin\varphi}{w}\right) = 0\]
oder nach Ausf\"uhrung der Differentiation
\[\frac{\partial^{2}F}{\partial y^{2}} + \frac{\partial^{2}F}{\partial z^{2}} + \frac{w'}{w} \frac{\partial F}{\partial x} + 4w^{2}(F + \alpha^{2} + \beta^{2} - \gamma) = 0.\]
Schreibt man $\alpha^{2} + \beta^{2} - \gamma = -q$ und beachtet, dass $F = 0$ ist, so geht die letzte Gleichung ueber in

\newpage
\noindent $\frac{w'}{w} + 2q + \text{const.} = 0$

und giebt nach einmaliger Integration
\[\left|\frac{dx}{w}\right| + 2\int q\, dx + \text{const.} = 0.\]
Die Integrationsconstante ist von $x$ unabh\"angig. Nimmt man andererseits $\frac{dx}{w}$ unabh\"angig von $y$ und $z$, so muss die Integrationsconstante eine lineare Function von $y$ und $z$ sein, weil $2\int q\, dx$ eine solche ist. Man hat also
\[\frac{1}{w} + 2\int q\, dx + 2ay + 2bz + \text{const.} = 0.\]
Vergleicht man damit das Resultat der directen Differentiation von $F$, n\"amlich
\[\frac{\partial F}{\partial x} = 2y \frac{d\alpha}{dx} + 2z \frac{d\beta}{dx} + \frac{d\gamma}{dx},\]
so ergiebt sich
\[\frac{d\alpha}{dx} = -aw, \quad \frac{d\beta}{dx} = -bw,\]
und wenn man $\int q\, dx = m$ setzt:
\[\alpha = -am + d, \quad \beta = -bm + e.\]
Hiernach hat man
\[q = -2a\varrho y - 2b\varrho z + f(m),\]
\[\frac{d^{2}q}{dx^{2}} = -2a\varrho \frac{dy}{dx} - 2b\varrho \frac{dz}{dx} + f'(m)\frac{dm}{dx} + f''(m)\left(\frac{dm}{dx}\right)^{2}\]
und diese Ausdr\"ucke sind in die Gleichung (1) einzuf\"uhren. Nach geh\"origer Hebung erh\"alt man
\[\frac{d^{2}q}{dx^{2}} + 2q \frac{dm}{dx} + 2q = 0\]
eine Gleichung, die sich weiter vereinfacht, wenn man beachtet, dass
\[f(m) = (a^{2} + b^{2})m^{2} - 2(ad + be)m + d^{2} + e^{2}.\]
Nimmt man hieraus $\frac{dq}{dx}$ und $\frac{d^{2}q}{dx^{2}}$, so geht die Differentialgleichung, welche die Bedingung des Minimum ausdr\"uckt, ueber in folgende:
\[\frac{d^{2}m}{dx^{2}} + 2m + 2(a^{2} + b^{2})m^{2} = 0.\]
Zur Ausf\"uhrung der Integration setze man $\frac{dm}{dx} = p$ und betrachte

\newpage
\noindent $q$ als unabh\"angige Variable. Dadurch erh\"alt man f\"ur $p^{2}$ als Function von $q$ eine lineare Differentialgleichung erster Ordnung, n\"amlich
\[\frac{d(p^{2})}{dq} - p^{2} + 2q + 2(a^{2} + b^{2})q^{2} = 0\]
oder
\[\frac{d(p^{2})}{dq} - p^{2} = -2q - 2(a^{2} + b^{2})q^{2}.\]
Das Integral lautet
\[p^{2} = e^{q} \left\{ \text{const.} + \int e^{-q} \left[2q + 2(a^{2} + b^{2})q^{2}\right] dq \right\}.\]
Darin ist f\"ur $p$ wieder $\frac{dm}{dx}$ zu setzen, wodurch man erh\"alt
\[\frac{dx}{dm} = \frac{\sqrt{2q + 2cq^{2} - (a^{2} + b^{2})2q^{3}}}{1},\]
also ergiebt sich
\[\text{(o)} \quad m = \int \frac{dq}{\sqrt{2q + 2cq^{2} - (a^{2} + b^{2})2q^{3}}},\]
\[2y = am - d + \sqrt{-q}\, \cos\psi,\]
\[2z = bm - e + \sqrt{-q}\, \sin\psi.\]
Man hat demnach $x$, $y$, $z$ als Functionen von zwei reellen Variabeln $q$ und $\psi$ ausgedr\"uckt. Die Ausdr\"ucke sind, abgesehen von algebraischen Gliedern, elliptische Integrale mit der obern Grenze $q$. Nach der oben entwickelten allgemeinen Methode h\"atte man $x$, $y$, $z$ erhalten als Summen von zwei conjugirten Functionen zweier conjugirter complexer Variabeln. Danach liegt die Vermuthung nahe, dass diese complexen Ausdr\"ucke mit H\"ulfe der Additionstheoreme der elliptischen Functionen sich je in einen einzigen Integralausdruck mit der Variabeln $q$ zusammenziehen lassen.

Und dies ist leicht zu best\"atigen. Man hat n\"amlich aus den Formeln f\"ur die Richtungscoordinaten $r$ und $\varphi$ der Normalen
\[\sqrt{\frac{\partial F}{\partial y} + \frac{\partial F}{\partial z} i} = \sqrt{y + zi + \alpha + \beta i} = \sqrt{\eta}\, e^{\frac{r}{2}i}.\]
Verbindet man damit die Definitionsgleichung von $q$, n\"amlich:
\[\sqrt{(y + zi + \alpha + \beta i)(y - zi + \alpha - \beta i)} = -q,\]
so ergiebt sich

\newpage
\noindent $(y + zi) + (\alpha + \beta i) = (-q)^{\frac{1}{2}} \eta^{\frac{1}{2}} e^{\frac{r}{2}i},$
$(y - zi) + (\alpha - \beta i) = (-q)^{\frac{1}{2}} \eta^{-\frac{1}{2}} e^{-\frac{r}{2}i}.$

Ferner hat man
\[\cot\frac{r}{2} \frac{dx}{\sqrt{-(dy^{2} + dz^{2})}} = \frac{1}{\sqrt{-q}} \frac{dx}{d\log\eta}\]
oder
\[\sqrt{nn'} = \frac{1}{\sqrt{-q}} \frac{dx}{d\log\eta},\]
\[2\sqrt{-2}\left\{p - 2aq(y + \alpha) - 2bq(z + \beta)\right\}\]
\[= \sqrt{\eta}\left(\cos\frac{r}{2} - \sin\frac{r}{2}\right) - \frac{1}{\sqrt{\eta}}\left(\sin\frac{r}{2} + \cos\frac{r}{2}\right)\]
\[\cdot \sqrt{-2}\left\{p - 2a^{2}(y + \alpha) - 2b^{2}(z + \beta)\right\}.\]
Auf der rechten Seite sind f\"ur $y + \alpha$ und $z + \beta$ die eben gefundenen Ausdr\"ucke in $\eta$ und $\eta'$ einzuf\"uhren. Dadurch geht die Gleichung ueber in folgende:
\[\sqrt{-q}\, \frac{d\log\eta}{dx} = \sqrt{2c + (a+bi)\eta^{-1} - (a-bi)\eta}.\]
Quadrirt man beide Seiten dieser Gleichung und setzt f\"ur $\frac{dx}{d\log\eta}$ seinen Werth aus (n), so ergiebt sich nach geh\"origer Reduction
\[\text{(p)} \quad (1 + \eta\eta')^{2} \left\{2c - (a+bi)\frac{\eta'}{\eta} - (a-bi)\frac{\eta}{\eta'}\right\}\]
\[= 8c - 2(a+bi)(\eta' - \eta^{-1}) - 2(a-bi)(\eta - \eta'^{-1}).\]
Die so gefundene Gleichung, welche den Zusammenhang von $q$, $\eta$, $\eta'$ angiebt, kann man als Integral einer Differentialgleichung f\"ur $\eta$ und $\eta'$ ansehen und $q$ als Integrationsconstante auffassen. Die Differentialgleichung ergiebt sich durch unmittelbare Differentiation in folgender Form
\[0 = \frac{d\eta}{\eta} \sqrt{2c + (a+bi)\eta^{-1} - (a-bi)\eta}\]
\[- \frac{d\eta'}{\eta'} \sqrt{2c + (a-bi)\eta'^{-1} - (a+bi)\eta'}.\]
Mit H\"ulfe der primitiven Gleichung (p) lassen sich aber die Factoren von $\frac{d\eta}{\eta}$ und $\frac{d\eta'}{\eta'}$ anders ausdr\"ucken. Man braucht nur die linke Seite

\newpage
\noindent von (p) in zweifacher Weise zu einem vollst\"andigen Quadrat zu erg\"anzen, indem man das fehlende doppelte Product das eine mal positiv, das andere mal negativ hinzuf\"ugt. Dadurch erh\"alt man
\[= +2\sqrt{[2c + (a+bi)\eta^{-1}][2c - (a-bi)\eta]}\]
\[-2\sqrt{[2c - (a+bi)\eta^{-1}][2c + (a-bi)\eta]}.\]
Nimmt man die Quadratwurzeln mit gleichen Vorzeichen, so geht die Differentialgleichung ueber in
\[0 = \frac{d\eta}{2\eta\sqrt{2c + (a+bi)\eta^{-1} - (a-bi)\eta}}\]
\[\text{(q)} \quad - \frac{d\eta'}{2\eta'\sqrt{2c + (a-bi)\eta'^{-1} - (a+bi)\eta'}}.\]
Ihr Integral in algebraischer Form ist in der Gleichung (p) ausgesprochen oder, was auf dasselbe hinauskommt, in den beiden Gleichungen
\[\sqrt{-q}\, (1 + \eta\eta') = \sqrt{[2c + (a+bi)\eta + (a-bi)\eta']}\]
\[\text{(r)} \quad + \sqrt{[2c - (a+bi)\eta - (a-bi)\eta']},\]
\[\sqrt{-q}\, i[(a+bi)\eta - (a-bi)\eta'] = \sqrt{[2c(a+bi) + 2c\eta' - (a-bi)\eta'^{2}]}\]
\[- \sqrt{[2c(a-bi) + 2c\eta - (a+bi)\eta^{2}]}.\]
In transcendenter Form lautet das Integral
\[\text{const.} = \int \frac{d\eta}{2\eta\sqrt{(a+bi) + 2c\eta - (a-bi)\eta^{2}}}\]
\[+ \int \frac{d\eta'}{2\eta'\sqrt{(a-bi) + 2c\eta' - (a+bi)\eta'^{2}}},\]
und die Integrationsconstante l\"asst sich ausdr\"ucken
\[\text{const.} = \int \frac{dq}{2\sqrt{q[1 + 2cq - (a^{2} + b^{2})q^{2}]}},\]
was aus der Gleichung (r) leicht hervorgeht, wenn man $\eta$ oder $\eta'$ constant und zwar $= 0$ nimmt. Man erkennt darin das Additionstheorem der elliptischen Integrale erster Gattung.

\newpage
%==============================================================================
% ANMERKUNGEN (Notes)
%==============================================================================
\section*{Anmerkungen.}
\addcontentsline{toc}{section}{Anmerkungen.}

Die erste Ausgabe der Abhandlung ueber Fl\"achen kleinsten Inhalts durch Hattendorff in den Abhandlungen der Gesellschaft der Wissenschaften zu G\"ottingen vom Jahre 1867 war eingeleitet durch eine historische Betrachtung, die bereits in der ersten Auflage dieser Gesammtausgabe als nicht von Riemann herr\"uhrend in Uebereinstimmung mit dem leider fr\"uh verstorbenen Hattendorff unterdr\"uckt wurde. Sie ist auch in der zweiten Auflage nicht mit aufgenommen. Zu erw\"ahnen ist aber, dass ungef\"ahr gleichzeitig mit Riemann Weierstrass Untersuchungen ueber die Fl\"achen vom kleinsten Inhalt angestellt hat, deren Resultate in den Monatsberichten der Berliner Akademie vom October und December 1866 zuerst ver\"offentlicht sind. Die Weierstrass'schen Arbeiten haben den Anstoss gegeben zu den weitergehenden Untersuchungen von H.~A.~Schwarz, von denen die erste Mittheilung im Monatsberichte der Berliner Akademie vom April 1865 gemacht ist. Die ausf\"uhrliche Ver\"offentlichung der gekr\"onten Preisschrift ``Bestimmung einer speciellen Minimalfl\"ache'' erschien im Jahre 1871. Darin ist neben Anderem das in Art.~18 der Riemann'schen Abhandlung behandelte Problem der durch ein regelm\"assiges r\"aumliches Viereck begrenzten Minimalfl\"ache bis zur wirklichen Aufstellung einer Gleichung zwischen den Coordinaten $x$, $y$, $z$ der Fl\"ache durchgef\"uhrt. Zu erw\"ahnen ist noch die an Riemann direct ankn\"upfende, im Jahre 1867 von der philosophischen Facult\"at zu G\"ottingen gekr\"onte Preisschrift von Arthur Schondorff ``Ueber die Minimalfl\"ache, deren Begrenzung von einem doppeltgleichschenkligen r\"aumlichen Viereck gebildet wird''.

Die Riemann'sche Abhandlung ist hier, mit wenigen nothwendigen Aenderungen, in der letzten Hattendorff'schen Bearbeitung abgedruckt. Einige Erl\"auterungen und Zus\"atze gebe ich in den nachstehenden Anmerkungen, wobei ich Mittheilungen und Winke von H.~A.~Schwarz mit Dank benutzt habe.

\medskip
\noindent\textbf{(1)} (Zu Seite 307.) Es ist hierbei zu beachten, dass nach (3) und (4)
\[\frac{dX}{d\eta} = \frac{1}{\eta'} \frac{ds}{d\eta}, \quad \frac{dX}{d\eta'} = \frac{1}{\eta} \frac{ds}{d\eta'}\]
Functionen von $\eta$ allein sind, und dass also auch $Y$, $Z$ als Functionen von $\eta$ allein betrachtet werden k\"onnen. $Y'$, $Z'$ sind dann die conjugirten Functionen, die nur von $\eta'$ abh\"angen.

\medskip
\noindent\textbf{(2)} (Zu Seite 311.) F\"ur unendlich kleine Werthe von $\eta$ (also auch von $r$) ergiebt sich aus (1) und (2)
\[\frac{dx}{dy} = -\sin r \cos\varphi, \quad \frac{dx}{dz} = -\sin r \sin\varphi,\]
\[\frac{dx}{dy} = -\sin r \sin\varphi, \quad \frac{dx}{dz} = \sin r \cos\varphi.\]

\clearpage
\chapter*{Source segment 8: riemann pages 351 400}
\addcontentsline{toc}{chapter}{Source segment 8: riemann pages 351 400}
\begin{center}
{\LARGE\bfseries Riemann --- Gesammelte Werke (pp.~351--400)}\\[1em]
\end{center}

\tableofcontents
\newpage

\section*{Anmerkungen}
\addcontentsline{toc}{section}{Anmerkungen}
\selectlanguage{ngerman}

\subsection*{(3) (Zu Seite 313.)}

Wenn der Winkel $\alpha = 0$ und also zwei Begrenzungsgerade einander parallel werden, so treten an Stelle dieser Entwickelungen die folgenden:
\[
y = \frac{\eta}{\sqrt{d\eta}} \log \xi + \text{f.c.}, \qquad
x = -\frac{\xi}{\sqrt{d\xi}} \eta + \text{f.c.}
\]

\subsection*{(4) (Zu Seite 320.)}

Der Schl\"ussel zu Art.~14 ist in den Entwicklungen des Fragments XXV zu suchen. Es handelt sich bei der Bestimmung von $\eta$ als Function von $t$ um die conforme Abbildung einer von Kreisbogen begrenzten Figur in der $\eta$-Ebene auf die positive $t$-Halbebene.

Setzt man
\[
\eta = \frac{y_1}{y_2},
\]
so ist nach dem genannten Fragment
\[
\frac{d^2y}{dt^2} + \Phi(t) \cdot y = 0
\]
eine rationale Function von $t$, die f\"ur reelle Werthe von $t$ reelle Werthe hat, und $y_1$, $y_2$ sind zwei particulare L\"osungen der linearen Differentialgleichung
\begin{equation}\label{eq:diffeq2}
\frac{d^2y}{dt^2} + \Phi(t) \cdot y = 0
\end{equation}
$\eta$ ist der Quotient zweier Particularl\"osungen dieser Gleichung. Man kann diese Gleichung dadurch transformiren, dass man $y_1$, $y_2$ mit einem und demselben Factor multiplicirt, ohne ihm dadurch die Eigenschaft zu nehmen, dass der Quotient zweier Particularl\"osungen die Function $\eta$ giebt. Setzt man, indem man unter $f$ zun\"achst eine willk\"urliche Function von $t$ versteht,
\[
y_1 = k\sqrt{f}, \qquad y_2 = \frac{1}{k}\sqrt{f},
\]
so erh\"alt man f\"ur $k$ die Differentialgleichung
\begin{equation}\label{eq:diffeq4}
\frac{d^2k}{dt^2} + 2\frac{df}{dt}\frac{1}{f}\frac{dk}{dt} = \Phi(t) \cdot k.
\end{equation}

worin
\begin{equation}\label{eq:phi5}
\Phi(t) = \frac{1}{4f^2}\left[4f'f'''(t) - 3f''(t)^2\right],
\end{equation}
und wenn also $f$ eine rationale Function von $t$ ist, so gilt das gleiche von $\Phi$.

Die Formeln des Art.~14 erh\"alt man, wenn man unter $f(t)$ die ganze rationale Function $(iv - \xi)^{2\nu}$ Grades
\[
f(t) = \Pi(t-a)\,\Pi(t-d)\,\Pi(t-b)\,\Pi(t-c)\,\Pi(t-e)
\]
versteht, und die Betrachtung der Unstetigkeiten ergiebt
\begin{equation}\label{eq:phi6}
\Phi(t) - \frac{1}{2}\frac{(\eta\eta'+\xi\xi')f'(t)}{(iv-\xi)^2f(t)} = F(t),
\end{equation}
wenn $F(t)$ eine ganze rationale Function vom Grade $2\nu-6$ mit reellen Coefficienten ist. Um auf andere Weise als im Text die Coefficienten der beiden h\"ochsten Potenzen in der Function $F(t)$ zu finden, beachte man, dass, wenn $t=\infty$ ein gew\"ohnlicher Punkt der Fl\"ache ist, die Entwicklung von $\frac{d\eta}{dt}$ nach fallenden Potenzen von $t$ mit $t^{\nu-1}$ anf\"angt, und dass folglich die Entwicklung $\frac{1}{y_2^2}$ mit $t$ beginnt.

Es muss sich also die Differentialgleichung~(\ref{eq:diffeq4}) durch eine Reihe von folgender Form integriren lassen
\[
c_0 t^\alpha + c_1 t^{\alpha-1} + c_2 t^{\alpha-2} + \cdots
\]
und dies hat man in~(\ref{eq:diffeq4}) einzusetzen, um Recursionsformeln zur successiven Berechnung der Coefficienten $c_\nu$ zu erhalten. Setzt man
\[
\alpha = \frac{2\nu-6}{2}, \qquad
\alpha(\alpha-1) = 2\cdot\frac{(2\nu-6)(2\nu-7)}{4} = (\nu-3)(2\nu-7),
\]
\[
f = t^{2\nu-4}, \qquad \frac{f'}{f} = \frac{2\nu-4}{t},
\]
so erh\"alt man aus~(\ref{eq:diffeq4})
\[
c_0(\alpha(\alpha-1) - (2\nu-4)\alpha + \Phi) = 0.
\]
Die Gleichsetzung der von $t$ unabh\"angigen Glieder ergiebt
\[
K = \alpha(\alpha-1) - (2\nu-4)\alpha = (\nu-3)(\nu-4),
\]
und die Vergleichung der Glieder mit $t^{-1}$
\[
K_1 = \alpha_1(1-\alpha)(\nu-3).
\]
$c_0$ und $c_1$ k\"onnen nicht bestimmt werden und die Vergleichung der h\"oheren Glieder ergiebt die Coefficienten $c_2$, $c_3$, \ldots

Auf ganz \"ahnliche Weise kann man auch die in der Anmerkung zu Art.~14 erw\"ahnten, den Punkten $a$, $a$, $b$ entsprechenden Bedingungsgleichungen erhalten.

Setzt man n\"amlich $z = t-a$, $t-a$, $t-b$, so muss sich die Differentialgleichung~(\ref{eq:diffeq4}) durch eine Reihe der Form
\[
z^{\lambda_0}(c_0 + c_1 z + c_2 z^2 + \cdots)
\]
integriren lassen. Ist nun
\[
\lambda_0 = \frac{1}{2}, \qquad l_0 = \lambda_0(\lambda_0-1),
\]
so ergiebt sich aus der Differentialgleichung~(\ref{eq:diffeq4})
\[
\frac{1}{z^2}\left(l_0 c_0 + (l_0+1)c_1 z + (l_0+2)c_2 z^2 + \cdots\right) - \frac{2}{z}\left(\frac{\lambda_0 c_0}{z} + \frac{(\lambda_0+1)c_1}{1} + \cdots\right) + \Phi(c_0 + c_1 z + c_2 z^2 + \cdots) = 0.
\]
Die Vergleichung der von $z$ unabh\"angigen Glieder giebt in Uebereinstimmung mit der Formel~(\ref{eq:phi6})
\[
l_0 - 2\lambda_0 = -\frac{1}{4},
\]
und die beiden n\"achsten Potenzen von $z$ ergeben
\[
c_1(l_0 + 2\lambda_0 + 1) + c_0(l_1 - 2(\lambda_0+1)\alpha_1) = 0
\]
woraus durch Elimination von $c_1$, $c_2$ eine Relation zwischen $l_0$, $l_1$, $l_2$ folgt.

$c_1$ kann aus~(\ref{eq:diffeq4}) nicht bestimmt werden. Die h\"oheren Glieder geben die Bestimmung der Coefficienten $c_2$, $c_3$, \ldots

\newpage

\section{Mechanik des Ohres}
\addcontentsline{toc}{section}{XVIII. Mechanik des Ohres}

\begin{center}
\small
(Aus Henle und Pfeuffer's Zeitschrift f\"ur rationelle Medicin, dritte Reihe, Bd.~29.)
\end{center}

\subsection*{1. Ueber die in der Physiologie der feineren Sinnesorgane anzuwendende Methode.}

F\"ur die Physiologie eines Sinnesorganes sind ausser den allgemeinen Naturgesetzen zwei besondere Grundlagen n\"othig, eine psychophysische, die erfahrungsm\"assige Feststellung der Leistungen des Organes, und eine anatomische, die Erforschung seines Baues.

Es sind demnach zwei Wege m\"oglich, um zur Kenntniss seiner Functionen zu gelangen. Man kann entweder vom Baue des Organes ausgehen und hieraus die Gesetze der Wechselwirkung seiner Theile und den Erfolg \"ausserer Einwirkungen zu bestimmen suchen, oder man kann von den Leistungen des Organes ausgehen und diese zu erkl\"aren versuchen.

Bei dem ersten Wege schliesst man von gegebenen Ursachen auf die Wirkungen, bei dem zweiten sucht man zu gegebenen Wirkungen die Ursachen.

Man kann mit Newton und Herbart den ersten Weg den \emph{synthetischen}, den zweiten den \emph{analytischen} nennen.

\textbf{Synthetischer Weg.}

Der erste Weg liegt dem Anatomen am n\"achsten. Mit der Untersuchung der einzelnen Bestandtheile des Organs besch\"aftigt, f\"uhlt er sich veranlasst, bei jedem einzelnen Theile zu fragen, welchen Einfluss er auf die Th\"atigkeit des Organs haben m\"oge. Dieser Weg w\"urde auch in der Physiologie der Sinnesorgane mit demselben Erfolg eingeschlagen werden k\"onnen, wie in der Physiologie der Bewegungsorgane, wenn die physikalischen Eigenschaften der einzelnen Theile sich bestimmen liessen. Die Bestimmung dieser Eigenschaften aus den Beobachtungen bleibt aber bei mikroskopischen Objecten immer mehr oder weniger ungewiss und jedenfalls im h\"ochsten Grade ungenau.

Man ist daher zu einer Erg\"anzung nach Gr\"unden der Analogie oder Teleologie gen\"othigt, wobei die gr\"osste Willk\"ur unvermeidlich ist, und aus diesem Grunde f\"uhrt das synthetische Verfahren in der Physiologie der Sinnesorgane selten zu richtigen und jedenfalls nicht zu sichern Ergebnissen.

\textbf{Analytischer Weg.}

Bei dem zweiten Wege sucht man zu den Leistungen des Organes die Erkl\"arung.

Das Gesch\"aft zerf\"allt in drei Theile.

\begin{enumerate}
\item Das Aufsuchen einer Hypothese, welche zur Erkl\"arung der Leistungen gen\"ugt.
\item Die Untersuchung, in wie weit sie zur Erkl\"arung nothwendig ist.
\item Die Vergleichung mit der Erfahrung, um sie zu best\"atigen oder zu berichtigen.
\end{enumerate}

I. Man muss das Instrument gleichsam nacherfinden und in so fern die Leistungen des Organs als Zweck, seine Sch\"opfung als Mittel zu diesem Zweck betrachten. Aber der Zweck ist kein vermutheter, sondern ein durch die Erfahrung gegebener, und wenn man von der Herstellung des Organs absieht, kann der Begriff der Endursachen ganz ausser dem Spiele bleiben.

Zu den thats\"achlichen Leistungen des Organs sucht man in dem Baue des Organs die Erkl\"arung. Bei dem Aufsuchen dieser Erkl\"arung hat man zuv\"orderst die Aufgabe des Organs zu analysiren; hieraus werden sich eine Reihe von secund\"aren Aufgaben ergeben, und erst nachdem man sich \"uberzeugt hat, dass sie gel\"ost sein m\"ussen, sucht man die Art und Weise, wie sie gel\"ost sind, aus dem Baue des Organs zu schliessen.

II. Nachdem aber eine Vorstellung gewonnen worden ist, welche zur Erkl\"arung des Organs ausreicht, darf man nicht unterlassen zu untersuchen, in wie weit sie zur Erkl\"arung nothwendig ist. Man muss sorgf\"altig unterscheiden, welche Voraussetzungen unbedingt oder vielmehr in Folge unbezweifelter Naturgesetze nothwendig sind, und welche Vorstellungsarten vielleicht durch andere ersetzt werden k\"onnen, das ganz willk\"urlich Hinzugedachte aber ausscheiden. Nur auf diese Weise k\"onnen die nachtheiligen Folgen der Benutzung von Analogien bei dem Aufsuchen der Erkl\"arung beseitigt werden, und auf diese Weise wird auch die Pr\"ufung der Erkl\"arung an der Erfahrung (durch Aufstellung von zu beantwortenden Fragen) wesentlich erleichtert.

III. Zur Pr\"ufung der Erkl\"arung an der Erfahrung k\"onnen theils die Folgerungen dienen, die sich aus ihr f\"ur die Leistungen des Organs ergeben, theils die bei dieser Erkl\"arung vorauszusetzenden physikalischen Eigenschaften der Bestandtheile des Organs. Was die Leistungen des Organs betrifft, so ist eine genaue Vergleichung mit der Erfahrung \"ausserst schwierig, und man muss die Pr\"ufung der Theorie meist auf die Frage beschr\"anken, ob kein Ergebniss eines Versuchs oder einer Beobachtung ihr widerspricht. Was dagegen die Folgerungen \"uber die physikalischen Eigenschaften der Bestandtheile betrifft, so k\"onnen diese von allgemeiner Tragweite sein und zu Fortschritten in der Erkenntniss der Naturgesetze Anlass geben, wie dies z.~B. bei dem Aufsuchen der Erkl\"arung der Achromasie des Auges durch Euler der Fall war.

F\"ur die beiden eben einander gegen\"ubergestellten Forschungsweisen gelten \"ubrigens die Bezeichnungen synthetisch und analytisch nur \emph{a potiori}. Genau genommen ist weder eine rein synthetische, noch eine rein analytische Forschung m\"oglich. Denn jede Synthese st\"utzt sich auf das Ergebniss einer vorausgehenden Analyse und jede Analyse bedarf zu ihrer Best\"atigung oder Berichtigung durch die Erfahrung der nachfolgenden Synthese. Bei dem ersten Verfahren bilden die allgemeinen Bewegungsgesetze das vorausgesetzte Ergebniss einer fr\"heren Analyse.

Das erste vorzugsweise synthetische Verfahren ist f\"ur die Theorie der feineren Sinnesorgane deshalb zu verwerfen, weil die Voraussetzungen f\"ur die Anwendbarkeit des Verfahrens zu unvollst\"andig erf\"ullt sind, die Erg\"anzung der Voraussetzungen durch Analogie und Teleologie hier aber v\"ollig willk\"urlich bleibt.

Bei dem zweiten vorzugsweise analytischen Verfahren kann die H\"ulfe der Teleologie und Analogie zwar auch nicht ganz entbehrt, wohl aber bei ihrer Benutzung die Willk\"urlichkeit vermieden werden, indem man

\begin{enumerate}
\item die Anwendung der Teleologie auf die Frage beschr\"ankt, durch welche Mittel die thats\"achlichen Leistungen des Organs ausgef\"uhrt werden, nicht aber bei den einzelnen Bestandtheilen des Organs die Frage nach dem Nutzen aufwirft;
\item die Anwendung von Analogien (das ``Dichten von Hypothesen'') sich zwar nicht, wie Newton will, g\"anzlich versagt, aber hinterher die Bedingungen, die zur Erkl\"arung der Leistungen des Organs erf\"ullt sein m\"ussen, heraushebt, und die zur Erkl\"arung nicht n\"othigen Vorstellungen, welche durch Benutzung der Analogie herbeigef\"uhrt worden sind, davon absondert.
\end{enumerate}

Nach diesen Principien m\"ussen nun f\"ur unsern Zweck zuv\"orderst die Leistungen des Geh\"ororgans festgestellt werden. Mit welcher Sch\"arfe, Feinheit und Treue das Ohr die Wahrnehmung des Schalles, seines Klanges und Tones, seiner St\"arke und Richtung vermittelt, dieses muss durch Beobachtung und Versuch so genau, wie irgend m\"oglich, bestimmt werden.

Ich setze diese Thatsachen als bekannt voraus. In dem Buche ``die Lehre von den Tonempfindungen als physiologische Grundlage f\"ur die Theorie der Musik'' von Helmholtz, findet man die Fortschritte zusammengestellt, welche in der so \"ausserst schwierigen Ermittelung der Thatsachen, die die Wahrnehmung der T\"one betreffen, in neuester Zeit gemacht worden sind und zwar vorz\"uglich von Helmholtz selbst.

Da ich den Folgerungen, welche Helmholtz aus den Versuchen und Beobachtungen zieht, entgegen zu treten vielfach gen\"othigt bin, so glaube ich um so mehr gleich hier aussprechen zu m\"ussen, wie sehr ich die grossen Verdienste seiner Arbeiten \"uber unsern Gegenstand anerkenne. Sie sind aber meiner Ansicht nach nicht in seinen Theorien von den Bewegungen des Ohres zu suchen, sondern in der Verbesserung der erfahrungsm\"assigen Grundlage f\"ur die Theorie dieser Bewegungen.

Ebenso muss ich auch den Bau des Ohres hier als bekannt voraussetzen, und bitte den geneigten Leser, n\"othigenfalls ein mit Abbildungen versehenes Handbuch der Anatomie zur H\"ulfe zu nehmen. Die Ergebnisse der neuesten Forschungen \"uber den Bau der Schnecke und des Ohres \"uberhaupt findet man dargestellt in der vor Kurzem erschienenen dritten Lieferung des zweiten Bandes von Henle's Handbuch der Anatomie des Menschen.

Ich betrachte es hier allein als meine Aufgabe, jene psychophysischen Thatsachen aus diesen anatomischen Thatsachen zu erkl\"aren. Die Theile des Ohres, die f\"ur unsern Zweck in Betracht kommen, sind die Paukenh\"ohle und das Labyrinth, welches aus dem Vorhofe, den Bogeng\"angen und der Schnecke besteht. Wir verfahren nun so, dass wir zun\"achst aus dem Baue dieser Theile zu schliessen suchen, was jeder derselben zu den Leistungen des Ohres beitragen m\"oge, dann aber bei jedem einzelnen Theile wieder von der durch ihn zu l\"osenden Aufgabe ausgehen und zun\"achst die Bedingungen aufsuchen, deren Erf\"ullung zu einer gen\"ugenden L\"osung der Aufgabe erforderlich ist.

\subsection*{2. Paukenh\"ohle.}

Man hat l\"angst erkannt, dass der Apparat in der Paukenh\"ohle die Wirkung hat, den Druck der Luft auf das Labyrinthwasser verst\"arkt zu \"ubertragen.

Nach den oben entwickelten Principien m\"ussen wir nun aus den in der Erfahrung gegebenen Leistungen des Organs die Bedingungen ableiten, welche bei dieser Uebertragung erf\"ullt werden m\"ussen. Es ergeben sich diese vorz\"uglich aus der Feinheit des Ohres in der Wahrnehmung des Klanges und aus der grossen Sch\"arfe, welche das Ohr, zumal das unverk\"ummerte Ohr des Wilden und des W\"ustenbewohners, besitzt. Versteht man unter Klang die Beschaffenheit des Schalles, welche von St\"arke und Richtung desselben unabh\"angig ist, so wird diese offenbar durch den Apparat vollkommen treu mitgetheilt, wenn er die Druck\"anderung der Luft in jedem Augenblick in constantem Verh\"altniss vergr\"ossert auf das Labyrinthwasser \"ubertr\"gt.

Es ist unverf\"anglich, dies als Zweck des Apparats anzusehen, wenn man nur dabei nicht unterl\"asst, zugleich aus den Leistungen des Ohres zu bestimmen, wie weit man durch die Erfahrung berechtigt d.\ h. gen\"othigt ist, die wirkliche Erf\"ullung dieses Zwecks vorauszusetzen.

Wir wollen dies sogleich thun, vorher jedoch f\"ur die Beschaffenheit der Druck\"anderung, von welcher der Klang abh\"angt, einen mathematischen Ausdruck suchen. Die Curve, welche die Geschwindigkeit der Druck\"anderung als Function der Zeit darstellt, bestimmt die Schallwelle vollst\"andig bis auf ihre Richtung, also auch St\"arke und Klang des Schalles. Nimmt man nun statt dieser Geschwindigkeit den Logarithmus von dieser Geschwindigkeit, oder wenn man lieber will, von deren Quadrat, so erh\"alt man eine Curve, deren Form von Richtung und St\"arke des Schalles unabh\"angig ist, die aber den Klang vollst\"andig bestimmt und daher ``Klangcurve'' heissen m\"oge.

L\"oste der Apparat seine Aufgabe vollkommen, so w\"urden die Klangcurven des Labyrinthwassers mit den Klangcurven der Luft v\"ollig \"ubereinstimmen. Durch die Feinheit des Ohres in der Wahrnehmung des Klanges halten wir uns nun zu der Annahme berechtigt, dass die Klangcurve durch die Uebertragung nur sehr wenig ge\"andert werde und also das Verh\"altniss zwischen den gleichzeitigen Druck\"anderungen der Luft und des Labyrinthwassers w\"ahrend eines Schalles sehr nahe constant bleibe.

Eine langsame Ver\"anderlichkeit dieses Verh\"altnisses ist damit sehr wohl vereinbar und wahrscheinlich. Sie w\"urde nur eine Ver\"anderlichkeit des Ohres in der Sch\"atzung der Schallst\"arke zur Folge haben, deren Annahme die Erfahrung durchaus nicht verbietet. W\"urde die Klangcurve merklich ge\"andert, so scheint eine solche Feinheit des Geh\"ors, wie sie sich z.~B. in der Wahrnehmung geringer Verschiedenheiten der Aussprache zeigt, mir kaum denkbar. Die unmittelbare Beurtheilung der Feinheit der Klangwahrnehmungen und besonders die Sch\"atzung der den Klangverschiedenheiten entsprechenden Verschiedenheiten der Klangcurve bleibt freilich immer sehr subjectiv.

Die Verschiedenheit des Klanges dient uns aber auch, die Entfernung der Schallquelle zu sch\"atzen. Von dieser Klangverschiedenheit k\"onnen wir die mechanische Ursache, die Ver\"anderung der Klangcurve bei der Fortpflanzung des Schalles in der Luft durch Rechnung bestimmen.

Wir k\"onnen indess dies hier nicht weiter verfolgen und wollen von dem Uebertragungsapparat nur fordern, dass er keine groben Entstellungen des Klanges bewirke, obgleich wir glauben, dass seine Treue viel gr\"osser ist, als man gew\"ohnlich annimmt.

I. Der Apparat in der Paukenh\"ohle (im unverk\"ummerten Zustande) ist ein mechanischer Apparat von einer Empfindlichkeit, die Alles, was wir von Empfindlichkeit mechanischer Apparate kennen, himmelweit hinter sich l\"asst.

In der That ist es durchaus nicht unwahrscheinlich, dass durch denselben Schallbewegungen treu mitgetheilt werden, die so klein sind, dass sie mit dem Mikroskop nicht wahrgenommen werden k\"onnten.

Die mechanische Kraft der schw\"achsten Sch\"alle, welche das Ohr noch wahrnimmt, l\"asst sich freilich kaum direct sch\"atzen; aber man kann mit H\"ulfe des Gesetzes, nach welchem die St\"arke des Schalles bei seiner Verbreitung in der Luft abnimmt, zeigen, dass das Ohr Sch\"alle wahrnimmt, deren mechanische Kraft Millionen Mal kleiner ist, als die der Sch\"alle von gew\"ohnlicher St\"arke.

In Ermangelung anderer von Fehlerquellen freier Beobachtungen berufe ich mich auf die Angabe von Nicholson, nach welcher das Rufen der Schildwachen von Portsmouth 4 bis 5 englische Meilen weit zu Ryde auf der Insel Wight bei Nacht deutlich geh\"ort wird. Wenn man erw\"agt, welche Vorrichtungen Colladon n\"othig hatte, um die Verbreitung des Schalles im Wasser wahrzunehmen, so wird man zugeben, dass von einer erheblichen Verst\"arkung des Schalles durch Fortpflanzung im Wasser nicht die Rede sein kann und dass hier in der That die mechanische Kraft des Schalles umgekehrt proportional dem Quadrat der Entfernung und wahrscheinlich noch schneller abnimmt.

Da die Entfernung von 4 bis 5 Meilen etwa 2000 Mal so gross ist als die Entfernung von 8 bis 10 Fuss, so ist die mechanische Kraft der das Trommelfell treffenden Schallwellen hier vier Millionen Mal kleiner, als in der Entfernung von 8 bis 10 Fuss von der Schildwache und die Bewegungen sind 2000 Mal kleiner. Man muss zugeben, dass bei den Schall-Empfindungen durchaus nichts von Verh\"altnissen, wie 1 zu 1000 Millionen oder 1 zu Tausend bemerkt wird. Nach den neueren Untersuchungen \"uber das Verh\"altniss der psychischen Sch\"atzung der Schallst\"arken zum physischen oder mechanischen Mass der Schallst\"arke bildet dies jedoch durchaus keinen Einwand gegen die eben erhaltenen Resultate. Wahrscheinlich ist dies Abh\"angigkeitsverh\"altniss gerade so, wie das unserer Sch\"atzung der Lichtst\"arke oder Gr\"osse der Fixsterne zu der mechanischen Kraft des uns von ihnen zugesandten Lichtes.

Hier hat man bekanntlich aus den Stern-Eichungen geschlossen, dass die mechanische Kraft des Lichtes im geometrischen Verh\"altnisse abnimmt, wenn die Gr\"osse des Fixsternes in arithmetischer Reihe steigt. Theilte man dem analog die Sch\"alle, von denen von gew\"ohnlicher St\"arke bis zu den eben noch wahrnehmbaren, in Sch\"alle von der ersten bis zur achten Gr\"osse, so w\"urde die mechanische Kraft f\"ur die Sch\"alle zweiter Gr\"osse etwa $\frac{1}{\sqrt{10}}$, f\"ur die dritter $\frac{1}{\sqrt{100}}$, \ldots, f\"ur die achter $\frac{1}{\sqrt{10000000}}$, d.\ h. $2^{\text{te}}$ Millionten Theil so gross sein, als f\"ur die Sch\"alle erster Gr\"osse, und die Weite der Bewegungen w\"urde f\"ur die Sch\"alle erster, dritter, f\"unfter, siebenter Gr\"osse sich wie $1 : \frac{1}{\sqrt{10}} : \frac{1}{\sqrt{100}} : \frac{1}{\sqrt{1000}}$ verhalten.

Ich habe oben bei der Betrachtung der das Ohr treffenden Schallwellen vor dem Trommelfell Halt gemacht, weil Einige eine D\"ampfung der st\"arkeren Sch\"alle (durch Spannung des Trommelfells?) annehmen. Ich muss jedoch gestehen, dass mir diese Meinung als eine v\"ollig willk\"urliche Vermuthung erscheint. Es m\"ogen allerdings, wenn ein starker Knall die Membranen des Labyrinths zu verletzen droht, Schutzvorrichtungen wirksam werden; aber ich finde in der Beschaffenheit der Geh\"orseindr\"ucke durchaus nichts Analoges mit dem Beleuchtungsgrad des Gesichtsfeldes beim Auge, und w\"usste durchaus nicht, was eine fortw\"ahrend ver\"anderliche Reflexth\"atigkeit des M.~tensor tympani f\"ur das genaue Auffassen eines Musikst\"ucks n\"utzen sollte. Meiner Ansicht nach hat man durchaus keinen Grund, bei dem Schalle in 10 Fuss Entfernung von der Schildwache ein anderes Verh\"altniss zwischen den Bewegungen der Luft vor dem Trommelfell und den Bewegungen der Steigb\"ugelplatte anzunehmen, als in der Entfernung von 20{,}000 Fuss; aber selbst wenn man eine ziemlich starke Ver\"anderlichkeit der Spannung des Trommelfells annimmt, werden unsere Schl\"usse dadurch nicht beeintr\"achtigt.

Wenn nun die Bewegungen der Steigb\"ugelplatte in der Entfernung von 10 Fuss von der Schildwache wahrscheinlich zu den eben mit blossen Augen noch wahrnehmbaren geh\"oren, so werden die Bewegungen in der Entfernung von 20{,}000 Fuss bei einer 2000 fachen Vergr\"osserung eben wahrnehmbar sein.

II. Soll der Paukenapparat so kleine Bewegungen treu mittheilen, wie er es der Erfahrung nach thut, so m\"ussen die festen K\"orper, aus denen er besteht, an den Stellen, wo sie auf einander wirken sollen, v\"ollig genau auf einander schliessen; denn offenbar kann ein K\"orper einem anderen eine Bewegung nicht mittheilen, sobald er um mehr als die Weite der Bewegung von ihm absteht.

Es wird ferner nur ein kleiner Theil der mechanischen Kraft der Schallbewegung durch anderweitige Arbeit, wie Spannung von Gelenkkapseln und Membranen, f\"ur das Labyrinth verloren gehen d\"urfen.

Ein solcher Verlust wird vermieden durch die \"ausserst geringe Breite des freien Randes der Membran des Vorhofsfensters. W\"are dieser Rand breiter, so w\"urden die Schwingungen der Steigb\"ugelplatte beinahe ganz durch Schwingungen dieses Randes ausgeglichen werden, und auf die Membranen der Schnecke und des Schneckenfensters nur eine geringe Wirkung stattfinden.

Die Wirkung dieses Membranenrandes auf die Steigb\"ugelplatte wird wegen der geringen Breite des Randes f\"ur die verschiedenen Lagen der Steigb\"ugelplatte w\"ahrend der Schallbewegung sehr verschieden sein. Man muss daher, wenn sie den Klang nicht entstellen soll, annehmen, dass die Elasticit\"at der Membran sehr gering ist, und die Steigb\"ugelplatte nicht durch sie, sondern durch andere Kr\"afte in die richtige Gleichgewichtslage gebracht wird.

III. Da die Theile des Paukenapparates, um die erfahrungsm\"asse Sch\"arfe des Ohres m\"oglich zu machen, fortw\"ahrend mit mehr als mikroskopischer Genauigkeit in einander greifen m\"ussen, so scheinen Correctionsvorrichtungen wegen der Ausdehnung und Zusammenziehung der K\"orper durch die W\"arme durchaus unentbehrlich. Die Temperatur\"anderungen m\"ogen innerhalb der Paukenh\"ohle nur sehr klein sein; dass sie aber stattfinden, ist nicht zu bezweifeln. F\"ur die Temperaturvertheilung im menschlichen K\"orper gilt, wenn die \"aussere Temperatur hinreichend lange constant gewesen ist, nahe das Gesetz, dass der Abstand der Temperatur an einer beliebigen Stelle des K\"orpers von der Hirntemperatur proportional ist dem Abst\"ande der \"ausseren Temperatur von der Hirntemperatur. Dieses Gesetz ergiebt sich aus dem Newton'schen und der Voraussetzung, dass der W\"armeleitungscoefficient und die specifische W\"arme innerhalb der in Betracht kommenden Temperaturen constant sei, eine Voraussetzung, die wahrscheinlich nahe erf\"ullt ist.

Man kann durch dieses Gesetz aus dem Abst\"ande der Temperatur der Paukenh\"ohle von der Hirntemperatur auf die Temperatur\"anderungen schliessen. Wenn sich nun auch der Temperaturunterschied zwischen Paukenh\"ohle und Hirn nicht bestimmen l\"asst, so kann man doch aus mehreren Gr\"unden, aus den Communicationen mit der \"ausseren Luft durch den \"ausseren Geh\"organg und die Tuba, auch wohl aus der Art und Weise der Blutversorgung der Paukenh\"ohle, mit grosser Wahrscheinlichkeit schliessen, dass ein merklicher Temperaturunterschied stattfindet.

Dagegen hat der Pyramidenknochen, weil er den Can.~caroticus enth\"alt, wahrscheinlich sehr nahe die Temperatur des Hirns, und wir m\"ussen daher annehmen, dass die innere Auskleidung der Paukenh\"ohle ein sehr schlechter W\"armeleiter und Strahler ist.

Von den \"ubrigen die Paukenh\"ohle umgebenden Knochen l\"asst sich freilich wohl nicht behaupten, dass sie eine so hohe Temperatur besitzen, wie das Hirn oder die Pyramide. Doch enthalten sie bedeutende W\"armequellen in Blutleitern, grossen Arterien und Venen und sind, wie die Pyramide, durch Schleimhaut und Periost gegen die Ausstrahlung in die Paukenh\"ohle gesch\"utzt. Wir d\"urfen daher annehmen, dass ihre Temperatur merklich h\"oher ist als die der Paukenh\"ohle.

Wenn nun die \"aussere Temperatur sinkt, so wird nach dem oben angef\"uhrten Gesetze der Abstand von der Hirntemperatur allenthalben im K\"orper in demselben Verh\"altniss (auf das Doppelte) steigen, die Paukenh\"ohle wird sich in Folge dessen merklich, die umgebenden Knochen nur sehr wenig abk\"uhlen, und die Geh\"orkn\"ochelchen werden sich merklich zusammenziehen, w\"ahrend die W\"ande der Paukenh\"ohle fast unge\"andert bleiben.

Viel mehr als dieses, dass die Geh\"orkn\"ochelchen sich beim Sinken der \"ausseren Temperatur viel st\"arker abk\"uhlen und zusammenziehen, als die W\"ande der Paukenh\"ohle, d\"urfte sich \"uber den Einfluss der Temperatur auf den Paukenapparat bei unserer g\"anzlichen Unbekanntschaft mit den thermischen Eigenschaften seiner Bestandtheile nicht feststellen lassen.

IV. Ich werde nun zun\"achst die Ver\"anderungen zu bestimmen suchen, welche bei einem Sinken der \"ausseren Temperatur in der Lage der Geh\"orkn\"ochelchen eintreten, damit alle zur Ber\"uhrung bestimmten Theile des Apparates fortfahren, genau auf einander zu schliessen. Der Theil des Geh\"orknochelsystems, der am unver\"anderlichsten mit der Wand der Paukenh\"ohle verbunden ist, ist das Ambos-Paukengelenk. Durch Abk\"uhlung werden alle Entfernungen in festen K\"orpern kleiner, also auch die Entfernung des Ambos-Steigb\"ugelgelenks von dieser Gelenkfl\"ache. Vom Hammer ist wahrscheinlich das obere Griffende derjenige Theil, welcher, wenigstens parallel dem Paukenfellring, die geringsten Verschiebungen zul\"asst.

Da nun bei der Abk\"uhlung die Entfernung des Ambos-Paukengelenks von dem am unver\"anderlichsten befestigten Punkt des oberen Hammergriffs im Paukenfell nahe unge\"andert bleibt, die Entfernungen dieser Punkte vom Ambos-Hammergelenk aber beide abnehmen, so muss sich am Ambos-Hammergelenk der Winkel zwischen den nach diesen Punkten gehenden Linien etwas weiter \"offnen.

Bei diesen beiden Aenderungen in der Lage der Geh\"orkn\"ochelchen wird der Hammer ein wenig in der Richtung vorn-median-hinten und gleichzeitig (um das Gelenkn\"opfchen des Amboses in seiner H\"ohe zu erhalten) sehr wenig in der Richtung vorn-oben-hinten gedreht. Der lange Fortsatz des Hammers w\"urde dabei in der Fissur nach oben und medianw\"arts bewegt werden, wenn er gegen Griff und Kopf des Hammers eine und dieselbe Lage behielte. Durch die Wirkung der Abk\"uhlung wird er aber st\"arker gekr\"ummt und dem Hammergriff gen\"ahert, so dass er sich w\"ahrend der Temperatur\"anderung wahrscheinlich nur allm\"ahlich ein wenig aus der Fissur herausbewegt.

V. Wir haben eben die Bedingungen aufgestellt, denen die Lage der Geh\"orkn\"ochelchen wahrscheinlich gen\"ugt, damit sie fortw\"ahrend genau auf einander schliessen und dabei weder im Rande der Vorhofsmembran, noch im Paukenfell eine merklich ungleichm\"assige Spannung erzeugen. Wir fragen nun nach den Mitteln, durch welche den Geh\"orkn\"ochelchen jederzeit die richtige Lage gegeben und gesichert wird. (Es wird dies meist durch einander entgegengesetzte Kr\"afte geschehen, welche bei der richtigen Lage des Kn\"ochelchens sich das Gleichgewicht halten und es, wenn es aus ihr entfernt w\"urde, in sie zur\"ucktreiben w\"urden.)

Es ist klar, dass diese in den beiden die Lage der Geh\"orkn\"ochelchen regulirenden Muskeln, in den Gelenkkapseln, Ligamenten, den Schleimhautfalten und den beiden Membranen, mit denen die Geh\"orkn\"ochelchen verwachsen sind, gesucht werden m\"ussten. Bei diesem Aufsuchen der Ursachen einer bestimmten Wirkung auf die Geh\"orkn\"ochelchen ergeben sich jedoch, namentlich wenn man die Schleimhautfalten mit in Betracht zieht, oft mehrere Wege zur Erzielung der Wirkung als m\"oglich. Um aus diesen verschiedenen M\"oglichkeiten die wahrscheinlichste herauszufinden, ist es vor allen Dingen n\"othig, sich durch anatomische Untersuchungen an frischen Pr\"aparaten ein ungef\"ahres Urtheil \"uber die Elasticit\"at und Spannung der B\"ander, H\"aute etc. zu verschaffen, was mir unm\"oglich ist. Man darf jedoch auch hoffen, durch sorgf\"altige Entwicklung der Consequenzen der verschiedenen Hypothesen bei den falschen auf Unwahrscheinlichkeiten zu stossen und diese so zu excludiren.

\newpage
Es ist f\"ur unsere jetzige Untersuchung zweckm\"assig, zu unterscheiden zwischen dem lauschenden, zum genauen H\"oren adaptirten Ohr und dem nicht lauschenden Ohr, und f\"ur bestimmte Fragen zwischen dem Ohr des Neugeborenen und des Erwachsenen. Wir machen die Unterscheidung zwischen dem lauschenden und nicht lauschenden Ohr, wenn die Steigb\"ugelplatte durch den Zug des M.~tensor tympani ein wenig gegen das Labyrinthwasser gedr\"uckt wird, so dass der Druck im Labyrinthwasser ein wenig st\"arker ist als in der Luft der Paukenh\"ohle; es werden dabei die Theile der festen K\"orper, deren Ber\"uhrung gesichert werden soll, ein wenig gegen einander gedr\"uckt. Diejenigen nun, welche eine solche fortw\"ahrende Spannung des Apparates (das Paukenfell etwa ausgenommen) f\"ur unwahrscheinlich halten, m\"ogen annehmen, dass bei den Temperatur\"anderungen die Geh\"orkn\"ochelchen durch die Wirkung der Haft- und Gelenkb\"ander und die allm\"ahliche Aenderung der Contraction der Muskeln ihre Lage \"andern, ohne gegen einander gedr\"uckt zu werden, weil wir gefunden haben, dass nur dann das genaue Ineinandergreifen aller Theile des Apparats gesichert ist.

Es bleibt dann unsere Untersuchung f\"ur das lauschende, zum genauen H\"oren absichtlich vorgerichtete Ohr g\"ultig, w\"ahrend daneben doch immer die M\"oglichkeit bestehen bleibt, dass das Ohr (des Wachenden?) fortw\"ahrend, wenn auch vielleicht in geringerem Grade, adaptirt ist.

Der Geh\"orkn\"ochelapparat besteht aus einem aus zwei Theilen (Hammer und Ambos) zusammengesetzten, um eine Axe drehbaren K\"orper und aus einem mit diesem K\"orper articulirenden, auf das Wasser des Vorhofsfensters dr\"uckenden Stempel (dem Steigb\"ugel). Das eine Ende der Umdrehungsaxe, der kurze Fortsatz des Amboses, ist mittelst des Ambos-Paukengelenks an der hintern Wand der Paukenh\"ohle befestigt, das andere Ende, der lange Fortsatz des Hammers, ragt, nur von Weichtheilen umgeben, in eine Spalte zwischen dem vordem obern Ende des Paukenfellrings und dem Felsenbein und legt sich in eine Furche dieses Ringes. (Wenigstens ist es so beim Ohr des Neugeborenen.)

Die Bestimmung der relativen Lage der Geh\"orkn\"ochelchen gegen die Paukenh\"ohle wird sehr erleichtert durch das Verfahren von Henle, die Paukenh\"ohle sich so gedreht zu denken, dass die Umdrehungsaxe horizontal von hinten nach vorn geht und das Vorhofsfenster vertical steht.

Wird der Stiel des Hammers durch Steigerung des Druckes der Luft auf das mit ihm verwachsene Trommelfell nach innen getrieben, so wird die Basis des Steigb\"ugels gegen die Membran des (ovalen) Vorhofsfensters gedr\"uckt, der Druck im Labyrinthwasser gesteigert und dadurch die Membran des (runden) Schneckenfensters nach aussen getrieben.

Damit der Apparat die kleinsten Druck\"anderungen der Luft, in stets gleichem Verh\"altniss vergr\"ossert, dem Labyrinthwasser mittheilen k\"onne, ist es vor allen Dingen n\"othig, dass der Druck des Steigb\"ugels stets in v\"ollig gleicher Weise auf das Labyrinthwasser wirke. Zu diesem Ende muss

\begin{enumerate}
\item der Druck der Basis stets Eine und dieselbe Fl\"ache treffen und die Richtung der Bewegung unver\"anderlich sein;
\item es d\"urfen keine Anheftungen des Steigb\"ugels an die Wand des Vorhofsfensters stattfinden, wenigstens keine solchen, die irgend einen merklichen Einfluss auf seine Lage und Bewegung aus\"uben k\"onnten;
\item der Steigb\"ugel darf nie aufh\"oren, gegen die Membran des Vorhofsfensters zu dr\"ucken.
\end{enumerate}

Wie man bei einiger Ueberlegung leicht finden wird, w\"urden die Druck\"anderungen der Luft entweder gar nicht oder nach v\"ollig ver\"anderten Gesetzen auf das Labyrinthwasser wirken, sobald Eine dieser Bedingungen verletzt w\"urde.

Um die Erf\"ullung der 3.\ Bedingung zu sichern, muss durch den M.~tensor tympani, welcher den Hammerstiel nach innen zieht, der Druck gegen die Membran des Vorhofsfensters stets auf einer solchen H\"ohe erhalten werden, dass er die gr\"ossten, beim H\"oren zu erwartenden Druck\"anderungen betr\"achtlich \"ubertrifft. Wahrscheinlich wird am Schnecken- oder Vorhofsfenster eine Wirkung dieses Druckes, sei es die Spannung oder Kr\"ummung (Ausdehnung, Form\"anderung) der Membran empfunden und durch den M.~tensor tympani der f\"ur das genaue H\"oren g\"unstigste Druck hergestellt.

Der Druck h\"angt nur von der Lage des Hammerstiels ab, und um die erforderliche Einstellung dieses Stiels zu bewirken, muss der Zug des Muskels gerade so stark sein, dass er der Wirkung der Spannung des Paukenfells bei dieser Einstellung das Gleichgewicht h\"alt. Ob die Spannung des Paukenfells dabei gr\"osser oder kleiner ist, darauf kommt gar nichts an; nur muss sie, wie wir jetzt zeigen wollen, so gross bleiben, dass nur ein sehr kleiner Theil der mechanischen Kraft der das Ohr treffenden Wellen an die Luft im Innern der Paukenh\"ohle verloren geht.

Wenn eine in freier Luft ausgespannte Membran von einer Schallwelle getroffen wird, so entstehen eine Schwingung der Membran, eine zur\"uckgeworfene Luftwelle und eine weitergehende (gebrochene) Luftwelle. Wie sich die mechanische Kraft der Schallwelle auf diese drei Wirkungen vertheilt, h\"angt von der Spannung der Membran ab. Ist diese Spannung sehr gering, so sind die beiden ersten Wirkungen sehr schwach, und es geht die Schallwelle fast unver\"andert weiter. Ist dagegen die Membran so stark gespannt, dass ihre Bewegungen nur sehr klein sind gegen die Schwingungen der Lufttheilchen in der auf sie treffenden Schallwelle, so kann sie der Luft auf der hintern Seite nur sehr kleine Bewegungen mittheilen und folglich auch ihren Druck nur wenig ver\"andern, und es wird fast die ganze Druck\"anderung auf der vordern Seite zur Spannung der Membran verwandt. Ausserdem aber entsteht, wenn die Membran in freier Luft ausgespannt ist, eine zur\"uckgeworfene Welle.

\newpage

Die Lage des Linsenbeines gegen das Vorhofsfenster kann also nicht unver\"andert bleiben; aber es kann durch Drehung des Amboses um seinen Befestigungspunkt (das Paukengelenk) bewirkt werden, dass das Linsenbein sich nur parallel der L\"angsaxe des Vorhofsfensters verschiebt, und also nur in dieser Richtung eine Drehung des Steigb\"ugels um das Centrum der Ambosgelenkfl\"ache n\"othig ist, um die Steigb\"ugelplatte an ihrem Platze zu erhalten. Da nun nur f\"ur diese Richtung eine Vorrichtung (der M.~stapedius) vorhanden ist, den Steigb\"ugel um das Ambosgelenkn\"opfchen willk\"urlich zu drehen, f\"ur die darauf senkrechte aber nicht, so darf man wohl vermuthen, dass die letztere Vorrichtung eben dadurch \"uberfl\"ussig gemacht worden ist, dass das Ambosgelenkn\"opfchen fortw\"ahrend in derselben H\"ohe erhalten wird.

VI. Dem Zuge der Sehne des M.~tensor tympani wird zum Theil das Gleichgewicht gehalten durch die Befestigung des Hammergriffs im Paukenfell und des Paukenfells im Sulcus tympanicus. Die Anheftung des Paukenfells an dem Hammergriff reicht aber (nach v.~Tr\"oltsch und Gerlach) nur wenig h\"oher, als der Insertionspunkt der Sehne, und ihr Endpunkt liegt selbst schon h\"oher als die Endigungen des Sulcus tympanicus. Offenbar kann also die Befestigung des Paukenfells im S.~t.\ dem M.~tensor tympani allein nicht das Gleichgewicht halten. Zum Gleichgewicht des Hammers ist vielmehr erforderlich, dass auf den oberhalb des Insertionspunkts gelegenen Theil ein gleich grosses entgegengesetzt gerichtetes Drehungsmoment wirke, wie auf den unterhalb gelegenen Griff.

Man kann diese zur Herstellung des Gleichgewichts n\"othige Kraft suchen

\begin{enumerate}
\item entweder in der Verbindung des Paukenfells mit den oberfl\"achlichen Schichten der Haut des \"ausseren Geh\"organgs,
\item oder in der Wirkung der hinteren Paukenfelltasche,
\item oder vielleicht in dem Zusammenwirken der Anheftungen des Hammerkopfes an die Paukenh\"ohlenwand durch den Ambos einerseits und anderseits durch das Lig.\ superius Arnoldi. Diese Anheftungen bilden einen etwa gegen die Spitze des kurzen Fortsatzes gerichteten Winkel und dr\"ucken, wenn sie gespannt sind, diese Spitze gegen das Paukenfell.
\end{enumerate}

\newpage

\begin{center}
\textbf{Dritte Abtheilung.}
\end{center}

\newpage

\section{Versuch einer allgemeinen Auffassung der Integration und Differentiation}
\addcontentsline{toc}{section}{XIX. Versuch einer allgemeinen Auffassung der Integration und Differentiation}

\begin{footnotesize}
\begin{flushright}
(20.\ Febr.\ 1847.)
\end{flushright}
\end{footnotesize}

In dem folgenden Aufsatze ist der Versuch gemacht, ein Verfahren aufzustellen, mittelst dessen man aus einer gegebenen Function einer Ver\"anderlichen eine andere Function derselben Ver\"anderlichen ableiten k\"onne, deren Abh\"angigkeit von jener urspr\"unglichen sich durch eine Zahl ausdr\"ucken l\"asst und die f\"ur den Fall, dass diese Zahl eine ganze positive, negative oder null ist, beziehungsweise mit den Differentialquotienten, Integralen und der urspr\"unglichen Function \"ubereinstimmt.

Die Resultate der Differential- und Integral-Rechnung werden zwar als Grundlage hier vorausgesetzt, aber nicht in der Weise, dass diejenigen derselben, die f\"ur alle Differentiale und Integrale, deren Ordnung durch eine ganze Zahl ausgedr\"uckt wird, gelten, auch auf die gebrochenen Ordnungen ausgedehnt w\"urden; sondern sie sollen nur einerseits zur Begr\"undung des oben angedeuteten Verfahrens benutzt werden und andrerseits als Wegweiser dienen dasselbe zu finden.

Zu diesem letzteren Zwecke wollen wir einmal die Reihe der Differentialquotienten etwas n\"aher betrachten. Es ist klar, dass man hierbei nicht von der gew\"ohnlichen Definition derselben ausgehen kann, die sich auf ihr recurrentes Bildungsgesetz gr\"undet, da man ja durch dasselbe unm\"oglich auf andere Glieder der Reihe, als auf solche, die ganzen Indices entsprechen, gelangen kann; man muss sich also nach einer independenten Bestimmung derselben umsehen. Ein Mittel dazu bietet uns die Entwicklung der Function, welche aus der urspr\"unglichen durch Vermehrung der Ver\"anderlichen um einen beliebigen Zuwachs entsteht, nach ganzen positiven Potenzen dieses Zuwachses dar.

Denn da die bekannte Entwicklung
\begin{equation}\label{eq:entwicklung1}
z(x+h) = \sum_{\mu=0}^{\infty} \frac{d^\mu z}{dx^\mu} \frac{h^\mu}{1 \cdot 2 \cdots \mu}
\end{equation}
(wo $z^{(\mu)}(x+h)$ das bedeutet, was aus $z(x)$ wird, wenn man darin statt $x$ setzt $x+h$) f\"ur jeden beliebigen Werth von $h$ g\"ultig ist, so m\"ussen die Coefficienten in derselben einen ganz bestimmten Werth haben; man kann dieselben also zur Definition der Differentialquotienten verwenden. Demgem\"ass stellen wir folgende Definition auf: der $n$te Differentialquotient der Function $z$ ist gleich dem Coefficienten von $h^n$ in der Entwicklung von $z(x+h)$ nach ganzen positiven Potenzen von $h$, multiplicirt in einen nach $x$ constanten, nur von $n$ abh\"angigen Factor, n\"amlich in $1 \cdot 2 \cdots n$. Diese Betrachtungsweise der Differentialquotienten f\"uhrt sehr leicht zur Feststellung einer allgemeinen Operation, in welcher die Differentiation und Integration enthalten ist und welche wir (da die Bezeichnung und Benennung derselben als die Grenze des Quotienten verschwindender Gr\"ossen bei dieser Betrachtungsweise keinen Sinn hat) durch $d^\nu$ bezeichnen und nach dem Vorgange von Lagrange in der Benennung ``fonctions deriv\'{e}es'' \emph{Ableitung} benennen wollen.

Wir verstehen n\"amlich unter $d^\nu z$ oder unter dem Ausdruck ``$\nu$te Ableitung von $z$ nach $x$'' den Coefficienten von $h^\nu$ in einer nach Potenzen von $h$, deren Exponenten um eine ganze Zahl von einander abstehen, r\"uckw\"arts und vorw\"arts ins unendliche fortlaufenden Entwicklung von $z(x+h)$, multiplicirt in einen nach $x$ constanten, nur von $\nu$ abh\"angigen Factor, d.\ h.\ wir definiren $d^\nu z$ durch die Gleichung
\begin{equation}\label{eq:abl2}
z(x+h) = \sum_{\nu=-\infty}^{\nu=\infty} \frac{h^\nu}{\Pi(\nu)} d^\nu z.
\end{equation}
In dieser Definition muss nun nat\"urlich der von $\nu$ allein abh\"angige Factor $k_\nu$ so bestimmt werden, dass f\"ur den Fall, dass die Exponenten von $h$ ganze Zahlen sind, die Reihe~(\ref{eq:abl2}) in die~(\ref{eq:entwicklung1}) \"ubergeht, weil nur dann die Differentialquotienten wirklich als besondere F\"alle in den Ableitungen enthalten sind; sollte dies nicht m\"oglich sein, so w\"are diese Definition unserm Zwecke, eine Operation, welche die Differentiation als besondern Fall in sich schliesst, festzustellen, nicht entsprechend, und wir m\"ussten uns also nach einem anderen Wege, ihn zu erreichen, umsehen.

\newpage

Bevor wir aber diesen Factor zu bestimmen suchen, wollen wir erst Einiges \"uber die Reihen von der angegebenen Form vorausschicken, da sie, wie man sieht, die Grundlagen dieses ganzen Versuchs einer Theorie der Ableitungen bilden.

Man hat wohl die Behauptung aufgestellt, man k\"onne auf die Reihen im Allgemeinen gar keine sicheren Schl\"usse gr\"unden, sondern nur unter der Bedingung, dass man den darin vorkommenden Gr\"ossen solche Zahlenwerthe beilege, dass die Reihe convergire, d.\ h.\ dass sich ihr (wenigstens gen\"ahrter) Werth durch eine wirkliche Ziffernaddition finden lasse. Nun k\"onnen wir aber, wenn, wie hier immer vorausgesetzt wird, die Coefficienten einem bestimmten Gesetze gehorchen, jeden einzelnen Theil derselben genau angeben; sie ist folglich eine in allen ihren Theilen genau begrenzte, also bestimmte Gr\"osse; und ich sehe darin, dass der Mechanismus der Ziffernaddition nicht ausreicht, diesen ihren bestimmten Werth zu finden, keinen Grund, warum wir nicht die Gesetze, die f\"ur die Zahlengr\"ossen als solche erwiesen sind, auf sie anwenden und die Resultate, die wir dadurch erhalten, als richtig ansehen sollten.

Um an einem Beispiele zu zeigen, dass man f\"ur eine Reihe von der Form~(\ref{eq:abl2}) wirklich einen Werth finden kann, wollen wir durch ein Verfahren, das in vielen F\"allen f\"ur diesen Zweck anwendbar ist, die Function $x^\mu$ in eine nach gebrochenen Potenzen von $(x-b)$ fortlaufende Reihe entwickeln, eine Entwicklung, deren wir ohnehin im Lauf der Untersuchung bed\"urfen.

Die Reihe, die $x^\mu$ gleich sein soll und die wir der K\"urze wegen durch $z$ bezeichnen, sei
\[
z = \sum_{\alpha} c_\alpha (x-b)^\alpha.
\]
Wenn $z = x^\mu$, so ist
\[
\frac{dz}{dx} = \mu x^{\mu-1}, \qquad \frac{d^2z}{dx^2} = \mu(\mu-1)x^{\mu-2},
\]
folglich
\[
(x-b)\frac{d^2z}{dx^2} - (\mu-1)\frac{dz}{dx} = 0.
\]
es muss also, auch
\[
\sum \left[ (\alpha - \mu)c_\alpha - b(\alpha+1)c_{\alpha+1} \right](x-b)^\alpha = 0
\]
sein. Dieser Bedingung ist offenbar Gen\"uge geleistet, sobald
\[
(\mu - \alpha)c_\alpha - b(\alpha+1)c_{\alpha+1} = 0.
\]
Nun sind aber alle Ausdr\"ucke, welche dieser Differentialgleichung gen\"ugen, in den verschiedenen Werthen von $k \cdot x^\mu$ enthalten, es muss also die Reihe, in der das Gesetz
\[
(\mu - \alpha)c_\alpha - b(\alpha+1)c_{\alpha+1} = 0
\]
stattfindet, nothwendig einem derselben gleich sein; um diesen zu finden, machen wir
\[
c_{\alpha-1}(x-b)^{\alpha-1} + c_\alpha(x-b)^\alpha = p,
\]
\[
p' = c_{\alpha+1}(x-b)^{\alpha+1} + c_{\alpha+2}(x-b)^{\alpha+2} + \cdots; \quad \text{also}
\]
\[
p + p' = z = k \cdot x^\mu; \quad \text{folglich}
\]
\[
x \frac{dp}{dx} - (\mu - \alpha)p = (\mu - \alpha)c_{\alpha-1}(x-b)^{\alpha-1} = \alpha c_\alpha(x-b)^{\alpha-1},
\]
\[
x \frac{dp'}{dx} - (\mu - \alpha)p' = -(\mu - \alpha)c_{\alpha+1}(x-b)^{\alpha+1} - \cdots = -b \cdot c_\alpha (x-b)^\alpha.
\]
Diese Differentialgleichungen haben zum allgemeinen Integral
\[
- \int x^{\mu-\alpha-1}dx + k_1 = p \cdot x^{-\mu} = c_\alpha(x-b)^\alpha x^{-\mu} + c_{\alpha-1}(x-b)^{\alpha-1}x^{-\mu},
\]
\[
\int x^{\mu-\alpha-1}dx + k_2 = p' \cdot x^{-\mu} = c_{\alpha+1}(x-b)^{\alpha+1}x^{-\mu} + c_{\alpha+2}(x-b)^{\alpha+2}x^{-\mu} + \cdots
\]
Substituirt man hierin f\"ur $x$ seinen Werth und $\frac{x-b}{x}$ f\"ur $y$, so erh\"alt man
\[
p \cdot x^{-\mu} = c_\alpha b^{\alpha-\mu}(1-y)^\alpha y^{\mu-\alpha} + c_{\alpha-1}b^{\alpha-1-\mu}(1-y)^{\alpha-1}y^{\mu-\alpha+1} + \cdots,
\]
\[
\frac{1}{x}(\mu-\alpha)b^{\alpha-\mu}\int y^{\mu-\alpha-1}(1-y)^\alpha \, dy + k_1 = p \cdot x^{-\mu},
\]
\[
\frac{1}{x}(\mu-\alpha)b^{\alpha-\mu}\int y^{\mu-\alpha-1}(1-y)^\alpha \, dy + k_2 = c_{\alpha+1}b^{\alpha+1-\mu}(1-y)^{\alpha+1}y^{\mu-\alpha-1} + \cdots
\]

In dem Falle, dass $\mu > \alpha > -1$, verschwinden nun offenbar die Ausdr\"ucke rechts beziehungsweise f\"ur $y=0$ und $y=1$, und die beiden Integrale werden ihnen also, das erste von $0$ bis $y$, das zweite von $1$ bis $y$ genommen, genau gleich sein, wenn dieselben zwischen diesen Grenzen continuirlich sind. Es k\"onnte scheinen, als ob diese Bedingung verletzt w\"are, so bald einige oder alle Glieder einer Reihe ins Positive oder Negative \"uber alle Grenzen hinaus wachsen; daraus w\"urde aber, da sich dieselben gegenseitig aufheben k\"onnen, nur folgen, dass sich durch eine wirkliche Addition ein bestimmter Werth f\"ur die Reihe nicht finden l\"asst. Da wir nun den Schluss, als ob die Reihe in einem solchen Falle \"uberhaupt keinen bestimmten Werth habe, nach dem Obigen nicht zugeben, so k\"onnen wir die Continuit\"at oder Discontinuit\"at der Reihen $p \cdot x^{-\mu}$ und $p' \cdot x^{-\mu}$ nur durch die Betrachtung der ihnen gleichen Integrale erfahren\footnote{Behandelt man die Integrale vor der Substitution von $\frac{x-b}{x}$ statt $x$, so werden sie f\"ur $x=0$ discontinuirlich. Man erkennt aber auch unter dieser Form leicht, dass die ihnen zugeh\"origen Constanten f\"ur positive und negative Werthe von $x$ dieselben Werthe haben m\"ussen, da der Werth der Integrale bei dem Uebergange des $x$ von $+\infty$ zu $-\infty$ sich stetig \"andert.}. Bekanntlich kann nun aber ein Ausdruck nur discontinuirlich werden, wenn sein Differential unendlich wird; der Ausdruck $(1-y)^{\mu-\alpha-1}y^{\alpha}$ hat aber f\"ur alle endlichen Werthe von $y$ einen endlichen Werth, wenn die Exponenten $\mu-\alpha-1$ und $\alpha$ positiv sind; die Integrale \"andern sich also dann stetig, und aus der Betrachtung der singul\"aren Integrale f\"ur $y=1$ und $y=0$ ersieht man, dass dies auch noch stattfindet, so lange beide Exponenten gr\"osser als $-1$ bleiben. Es ist demnach f\"ur den Fall, dass $\mu > \alpha > -1$ und $y$ endlich ist\footnote{F\"ur den Fall, dass $y = +\infty$, also $x=0$, ist der Werth beider Integrale $\infty$; folglich $k = \infty - \infty$, d.\ h.\ beliebig, was offenbar aus der blossen Betrachtung dieses Falles hervorgeht.},
\begin{equation}\label{eq:intresult}
k = z \cdot x^{-\mu} = p \cdot x^{-\mu} + p' \cdot x^{-\mu} = (\mu - \alpha) c_\alpha b^{\alpha-\mu} \int_0^1 (1-y)^{\mu-\alpha-1} y^\alpha \, dy
\end{equation}
(wo $\Pi$ das bekannte bestimmte Integral bezeichnet). Dies Resultat gilt, wie bemerkt, nur, wenn $\mu > \alpha > -1$; es l\"asst sich aber auf alle Werthe von $\mu$ und $\alpha$ ausdehnen, wenn man das $\Pi$ einer negativen Zahl (wie im Lauf dieser Untersuchung immer angenommen werden soll) als durch das Gesetz $\Pi(n) = \frac{1}{n+1}\Pi(n+1)$ aus den positiven abgeleitet definirt. Denn erstens muss es nach dem Gesetz, welches angenommener Massen zwischen den Coefficienten der Reihe stattfindet, f\"ur jeden Werth von $\alpha$ gelten, wenn nur einer derselben $>-1$ ist; es ist also, wenn $\mu$ positiv ist,
\[
k = \frac{\Pi(\alpha)\Pi(\mu-\alpha-1)}{\Pi(\mu)}
\]
oder
\[
k \cdot x^\mu = \frac{\Pi(\alpha)\Pi(\mu-\alpha-1)}{\Pi(\mu)} b^{\alpha-\mu} x^\mu,
\]
daraus aber erh\"alt man durch $n$ malige Differentiation nach $x$
\[
\frac{d^n(k \cdot x^\mu)}{dx^n} = \frac{\Pi(\alpha)\Pi(\mu-\alpha-1)}{\Pi(\mu-n)} b^{\alpha-\mu} x^{\mu-n},
\]
wodurch das Gesetz auch f\"ur negative Werthe von $\mu$ erwiesen ist.

Es ist also ganz allgemein
\begin{equation}\label{eq:allgemein3}
x^\mu = \sum \frac{\Pi(\alpha)}{\Pi(\mu)} \frac{(x-b)^\alpha}{b^{\alpha-\mu}}.
\end{equation}
Bemerkenswerth ist es, dass man durch diese Formel eine Reihe f\"ur $x^\mu$ nicht erh\"alt, wenn $\mu$ eine negative ganze Zahl ist, da der Ausdruck links dann $0$ wird, worauf wir sp\"ater zur\"uckkommen werden. Man sieht auch, dass es Reihen von dieser Form giebt, die der Null oder einer Constanten, f\"ur jeden Werth von $x$, gleich sind.

Nach dieser Protestation gegen das Verdammungsurtheil, welches man den divergirenden Reihen gesprochen hat, wollen wir jetzt den eingeschlagenen Weg zur Feststellung des Begriffs der Ableitungen weiter verfolgen. Man sieht, dass der Zweck, den wir uns gesetzt haben, dass n\"amlich die Differentiation als besonderer Fall in der Ableitung enthalten sein soll, erf\"ullt ist, so bald nur die Function $k_\nu$ f\"ur alle ganzen positiven Werthe von $\nu = \frac{1}{\Pi(\nu)}$ und f\"ur alle ganzen negativen Werthe $=0$ ist; denn dann geht die Reihe~(\ref{eq:abl2}) in die~(\ref{eq:entwicklung1}) \"uber; dieser Bedingung kann aber offenbar durch unendlich viele verschiedene Functionen von $\nu$ gen\"ugt werden; man kann ferner durchaus nicht annehmen, dass es nur Eine Entwicklung derselben Function nach denselben Potenzen von $h$ gebe, d.\ h.\ dass nur Ein System von Coefficienten einer Reihe von einer bestimmten Form einen bestimmten Werth gebe; man muss vielmehr unendlich viele verschiedene Systeme als m\"oglich voraussetzen; wir haben also, unbeschadet unseres Zweckes, sowohl unter den verschiedenen m\"oglichen Functionen von $\nu$ f\"ur $k_\nu$ als unter verschiedenen m\"oglichen Systemen von Coefficienten die Wahl, und es ist offenbar am zweckm\"assigsten, diese Wahl wom\"oglich so zu treffen, dass die Ableitungen noch mehreren Gesetzen gehorchen, die bei einer andern Wahl nur f\"ur Ableitungen mit ganzen Indices g\"ultig sein w\"urden.

Hierzu dienen folgende Betrachtungen.

Da der Ausdruck $\Sigma k_\nu \, dl^\nu z \, h^\nu$ alle in dieser Form m\"oglichen Entwicklungen $z(x+h)$ umfassen soll, so muss
\[
\frac{\partial}{\partial x}\Sigma k_\nu \, dl^\nu z \, h^\nu = \Sigma k_\nu \, dl^{\nu+1} z \, h^\nu \cdot \frac{1}{h}
\]
alle in dieser Form m\"oglichen Entwicklungen von $\frac{\partial z(x+h)}{\partial x}$ umfassen, und ebenso
\[
\frac{\partial}{\partial h}\Sigma k_\nu \, dl^\nu z \, h^\nu = \Sigma k_\nu \, \nu \, dl^\nu z \, h^{\nu-1}
\]
alle Entwicklungen dieser Form von $\frac{\partial z(x+h)}{\partial h}$. Bekanntlich sind nun $\frac{\partial z(x+h)}{\partial x}$ und $\frac{\partial z(x+h)}{\partial h}$ identisch; beide Ausdr\"ucke umfassen also genau dieselben Reihen; es m\"ussen also auch $k_{\nu+1}\,dl^{\nu+1}z$ und $k_\nu \cdot \nu \, dl^{\nu+1}z$ genau dieselben Werthe haben, d.\ h.\ sie sind einander gleich; setzt man nun $k_{\nu+1} = k_\nu \cdot \nu \cdot l(\nu)$, was der obigen Hauptbedingung offenbar nicht widerspricht, da f\"ur ganze Werthe von $\nu$ verm\"oge derselben dies Gesetz stattfinden muss, so erreicht man dadurch, dass auch f\"ur die Ableitungen mit gebrochenen Indices
\begin{equation}
dl^{\nu+1}z = \frac{d\,dl^\nu z}{dx}
\end{equation}
ist und folglich allgemein, wenn $n$ eine ganze Zahl ist,
\begin{equation}
dl^{\nu+n}z = \frac{d^n\,dl^\nu z}{dx^n}.
\end{equation}
Aus dem angenommenen Gesetze f\"ur $k_\nu$ folgt, dass
\[
\Pi(\nu) \cdot k_\nu = \Pi(\nu+1) \cdot k_{\nu+1}
\]
ist, es hat also die Function $\Pi(\nu) \cdot k_\nu$, die wir durch $l(\nu)$ bezeichnen wollen, f\"ur alle Werthe von $\nu$, die um ganze Zahlen von einander abstehen, stets denselben Werth. Wir k\"onnen daher f\"ur die zweckm\"assigste Wahl der Function $k_\nu$ nicht mehr aus der Betrachtung einer einzelnen Entwicklungsform, sondern nur aus der Combination verschiedener Schl\"usse ziehen; demgem\"ass wollen wir versuchen, ob wir sie so w\"ahlen k\"onnen, dass $dl^\mu dl^\nu z = dl^{\mu+\nu}z$ ist.

L\"asst man zu diesem Zwecke $x$ in der Formel~(\ref{eq:abl2}) noch einmal wachsen, und bezeichnet man diesen Zuwachs durch $k$, so ist
\begin{equation}\tag{$\alpha$}
z(x+h+k) = \sum_{\mu+\nu=-\infty}^{\infty} \frac{h^\mu}{\Pi(\mu)} \frac{k^\nu}{\Pi(\nu)} dl^\mu dl^\nu z,
\end{equation}
und dieser Ausdruck bezeichnet alle nach denselben Potenzen von $h$ und $k$ m\"oglichen Entwicklungen von $z(x+h+k)$. Es ist aber auch
\begin{multline}\tag{$\beta$}
z(x+h+k) = z((x+h)+k) = \sum_{\nu=-\infty}^{\infty} \frac{k^\nu}{\Pi(\nu)} dl^\nu z(x+h) \\
= \sum_{\nu=-\infty}^{\infty} \sum_{\mu=-\infty}^{\infty} \frac{k^\nu}{\Pi(\nu)} \frac{h^\mu}{\Pi(\mu)} dl^\mu dl^\nu z \cdot \frac{l(\mu+\nu)}{l(\mu)l(\nu)}.
\end{multline}
Nun bezeichnet der letzte Ausdruck ($\beta$) zwar nicht alle m\"oglichen Entwicklungen dieser Form von $z(x+h+k)$, da die Gleichung~(\ref{eq:allgemein3}) nur Eine Entwicklung von $x^\mu$ giebt, ohne dass dies die einzig m\"ogliche zu sein brauchte; es m\"ussen aber alle in ihm enthaltenen Entwicklungen auch in ($\alpha$) enthalten sein; stellt man also f\"ur die Function $l$ das Gesetz $l(\mu+\nu) = l(\mu) \cdot l(\nu)$ auf, so werden alle Werthe von $dl^{\mu+\nu}z$ auch Werthe von $dl^\mu dl^\nu z$ sein, obgleich der letzte Ausdruck auch noch andere Werthe haben kann.

Es ist also
\begin{equation}
dl^\mu dl^\nu z = dl^{\mu+\nu} z
\end{equation}
unter der ausgesprochenen Beschr\"ankung.

Aus $l(\mu+\nu) = l(\mu) \cdot l(\nu)$ folgt aber
\[
l(\mu+\nu+\pi) = l(\mu+\nu) \cdot l(\pi) = l(\mu) \cdot l(\nu) \cdot l(\pi)
\]
und allgemein, dass das Product der $l$ verschiedener Zahlen gleich ist dem $l$ ihrer Summe, oder wenn man die einzelnen Factoren einander gleich setzt $l(m\nu) = l(\nu)^m$, so oft $m$ eine ganze Zahl ist; bezeichnet man nun $\frac{1}{m}$ durch $\mu$, so ist
\[
l(m\nu)^{\frac{1}{m}} = l(\nu) \quad \text{und} \quad l(\mu\nu)^{\frac{1}{\mu}} = l(\nu), \quad \text{oder} \quad l(\mu\nu) = l(\nu)^\mu.
\]
Das Gesetz $l(\mu\nu) = l(\nu)$ ist also f\"ur alle rationalen Werthe von $\mu$, und folglich (nach dem bekannten Gesetz der Interpolation) allgemein g\"ultig.

Da nun f\"ur ganze Werthe von $\nu$ $l(\nu) = 1$ sein muss, so ist $l(\nu) = 1$.

Sollen demnach die Gesetze~(4) und~(5) f\"ur die Ableitungen im Allgemeinen gelten, und die Differentiation in der Ableitung als besonderer Fall enthalten sein, so m\"ussen wir die Ableitungen unter denjenigen Functionen von $x$ w\"ahlen, die der Gleichung
\[
\Pi(\nu) \cdot k_\nu = \text{const.}
\]
gen\"ugen. Diese Wahl wird am zweckm\"assigsten auf diejenigen unter ihnen fallen, welche am geschmeidigsten f\"ur die Rechnung sind; versucht man aber die Entwicklung einiger Functionen von $x+h$ in Reihen, die nach gebrochenen Potenzen von $h$ fortlaufen, so wird man sehen, dass am leichtesten und einfachsten Entwicklungen in solche Reihen sind, in denen der Coefficient von $\frac{h^\nu}{\Pi(\nu)}$ das Differential des Coefficienten von $\frac{h^{\nu+1}}{\Pi(\nu+1)}$ ist: wir wollen also obige Begrenzung der Ableitungen dahin beschr\"anken, dass das Zeichen $dl^\nu$ den Coefficienten von $\frac{h^\nu}{\Pi(\nu)}$ nicht in allen m\"oglichen Entwicklungen von $z(x+h)$ bezeichnen soll, sondern nur in solchen, in denen der Coefficient von $\frac{h^{\nu+1}}{\Pi(\nu+1)}$ das Differential des Coefficienten von $\frac{h^\nu}{\Pi(\nu)}$ ist\footnote{Aus~(4) folgt zwar, dass wenn $\sum \frac{h^\nu}{\Pi(\nu)} dl^\nu z$ eine Entwicklung von $z(x+h)$ ist, auch $\frac{d}{dx} \sum \frac{h^{\nu+1}}{\Pi(\nu+1)} dl^\nu z = \sum \frac{h^\nu}{\Pi(\nu)} \frac{d\,dl^\nu z}{dx}$ eine Entwicklung von $\frac{dz}{dx}$ ist, aber nicht dass diese beiden Entwicklungen identisch sind. Durch die gemachte Annahme erreicht man auch, dass die Ableitungen mit ganzen negativen Indices, die nach dem Bisherigen noch gar keinen Sinn hatten, mit den Integralen zusammenfallen, wie weiter unten bewiesen werden wird.}.

\newpage

Hieraus folgt zun\"achst, dass Ein Werth von $dl^\nu z$ nur einer Entwicklung angeh\"oren kann; denn gesetzt, ein Werth von $dl^\nu z$, $p_\nu$, geh\"orte zwei Entwicklungen, $a$ und $b$, an, so m\"ussten diese beiden Entwicklungen in allen folgenden Gliedern \"ubereinstimmen, da diese durch Differentiation aus $p_\nu$ entstehen. Bezeichnen wir nun die vorhergehenden Glieder in $a$ durch $p_{\nu-1}$, $p_{\nu-2}$, \ldots, in $b$ durch $q_{\nu-1}$, $q_{\nu-2}$, \ldots, so m\"ussen $p_{\nu-1}$ und $q_{\nu-1}$ beide zum Differential $p_\nu$ haben; sie k\"onnen also nur um eine Constante verschieden sein, d.\ h.
\[
q_{\nu-1} = p_{\nu-1} + K_1,
\]
ebenso muss
\[
q_{\nu-2} = p_{\nu-2} + K_1 x + K_2, \quad q_{\nu-3} = p_{\nu-3} + K_1 \frac{x^2}{1 \cdot 2} + K_2 x + K_3
\]
sein. Die Entwicklung $b$ ist also
\[
\sum_{m=1}^{n=\infty} \frac{K_m}{\Pi(n)} \frac{h^{\nu-n}}{\Pi(\nu-n)} = \sum_{m=1}^{\infty} K_m \frac{h^{\nu-m}}{\Pi(\nu-m)}.
\]
nun soll aber f\"ur alle Werthe von $(x+h)$ $a=b$ sein, was bekanntlich nur stattfinden kann, wenn alle Constanten null sind; dann aber sind beide Entwicklungen identisch.

Ist $p_\nu$ ein Werth von $dl^\nu z$, so ist $p_\nu + K \frac{x^{\nu-n}}{\Pi(\nu-n)}$ (wo $n$ positiv und ganz und $K$ eine endliche Constante ist) ebenfalls ein Werth desselben; denn die Reihe
\[
\sum \frac{h^{\nu+n}}{\Pi(\nu+n)} dl^{\nu+n}z + \frac{h^\nu}{\Pi(\nu)}(p_\nu + K) + \sum \frac{h^{\nu-n}}{\Pi(\nu-n)}\left(p_{\nu-n} + K\frac{x^n}{\Pi(n)}\right)
\]
und es findet in ihr das Gesetz statt
\[
\frac{d}{dx}\left(p_{\nu-n} + K\frac{x^n}{\Pi(n)}\right) = p_{\nu-n+1} + K\frac{x^{n-1}}{\Pi(n-1)}.
\]

Den Inbegriff aller Werthe von $dl^\nu z$, die sich durch Addition von Ausdr\"ucken von der Form $K \frac{x^n}{\Pi(n)}$ aus einander ableiten lassen, wollen wir ein \emph{System} von Werthen nennen; es sind also alle Werthe von $dl^\nu z$, die demselben Systeme angeh\"oren, in dem Ausdruck
\begin{equation}
p_\nu + \sum_{n=0}^{\infty} K_n \frac{x^n}{\Pi(n)}
\end{equation}
enthalten (wo $K_n$ endliche Constanten bedeuten).

Wir wollen nun einen Werth von $dl^\nu z$ zu bestimmen suchen. Bekanntlich ist
\[
z(x) = z(k) + \left(\frac{dz}{dx}\right)_{\!k}(x-k) + \frac{1}{1 \cdot 2}\left(\frac{d^2z}{dx^2}\right)_{\!k}(x-k)^2 + \cdots,
\]
sobald $z$ zwischen den Grenzen $x$ und $k$ continuirlich ist; setzt man hierin $x+h$ f\"ur $x$ und entwickelt die Glieder der Reihe mittels~(\ref{eq:allgemein3}) nach Potenzen von $h$, so erh\"alt man
\[
z(x+h) = \sum_{\mu=-\infty}^{\infty} \frac{h^\mu}{\Pi(\mu)} \left[ \frac{(x-k)^{-\mu}}{\Pi(-\mu)}z(k) + \frac{(x-k)^{-\mu+1}}{\Pi(-\mu+1)}\left(\frac{dz}{dx}\right)_{\!k} + \frac{(x-k)^{-\mu+2}}{\Pi(-\mu+2)}\left(\frac{d^2z}{dx^2}\right)_{\!k} + \cdots \right]
\]
und in dieser Reihe ist der Coefficient von $\frac{h^\mu}{\Pi(\mu)}$ das Differential des Coefficienten von $\frac{h^{\mu-1}}{\Pi(\mu-1)}$; er ist folglich ein Werth von $dl^\nu z$, den wir durch $p_\nu$ bezeichnen wollen. Differentiirt man nach $k$, so erh\"alt man
\[
\frac{dp_\nu}{dk} = \frac{(x-k)^{-\mu-1}}{\Pi(-\mu-1)}z(k) + \frac{(x-k)^{-\mu}}{\Pi(-\mu)}\left(\frac{dz}{dx}\right)_{\!k} + \cdots = p_{\nu+1}.
\]
Nun verschwinden alle Glieder der obigen Reihe f\"ur $k=x$; das Integral wird also von $h$ bis $x$ genommen $= p_{\nu+1}$ sein, wenn es zwischen den Grenzen continuirlich ist; dies ist aber, da $z$ zwischen den Grenzen $x$ und $k$ continuirlich sein soll und $-\mu-1 > -1$, offenbar der Fall und es ist also
\begin{equation}
\int_x^k (x-t)^{-\mu-1}z(t)\,dt = \frac{p_{\nu+1}}{\Pi(-\mu-1)}
\end{equation}
ein Werth von $dl^\nu z$, sobald $z$ zwischen den Grenzen $x$ und $k$ continuirlich und $\mu$ negativ ist. Der derselben Entwicklung angeh\"orige Werth von $dl^{\nu-n}z$ ist gleich
\[
\Pi(-\mu+n-1) \int_x^k (x-t)^{-\mu+n-1}dl^n z\,dt.
\]

Man sieht leicht, dass, je nachdem man dem $k$ verschiedene Werthe giebt, verschiedene Entwicklungen von $z(x+h)$ daraus hervorgehen, aber alle diese Entwicklungen geh\"oren demselben Systeme an. Denn aus dem Werth
\[
\int_x^k (x-t)^{-\mu-1}z(t)\,dt
\]
geht offenbar
\[
\int_{k_1}^k (x-t)^{-\mu-1}z(t)\,dt
\]
hervor durch Addition von
\[
\int_x^{k_1} (x-t)^{-\mu-1}z(t)\,dt = \Pi(-\mu-1)\Pi(\mu)\,dl^\nu z;
\]
da nun $z$ zwischen $x$ und $k$, und also auch zwischen $k$ und $k_1$ continuirlich ist, so sind alle jene Integrale endliche und zwar nach $x$ constante Gr\"ossen. Man wird demnach durch das angewandte Verfahren stets auf dasselbe System von Werthen gelangen: beschr\"anken wir also den Begriff der Ableitungen auf dies System von Werthen, so haben wir die Bestimmung derselben auf bekannte Werthe zur\"uckgef\"uhrt und werden mittels dieser Definition die Eigenschaften derselben und ihre Werthe f\"ur bestimmte Functionen ableiten k\"onnen.

Es ist demnach
\begin{equation}
dl^\nu z = \Pi(-\nu-1) \int_x^k (x-t)^{-\nu-1}z(t)\,dt + \sum_{n=0}^{\infty} K_n \frac{x^n}{\Pi(n)}
\end{equation}
wenn $K_n$ endliche willk\"urliche Constanten sind\footnote{Alle diese willk\"urlichen Functionen wollen wir durch $\varphi_\nu$ bezeichnen; wir machen zugleich darauf aufmerksam, dass (wenn $n$ positiv und ganz) jede Function $\varphi_\nu$ auch eine Function $\varphi_{\nu-n}$ ist.}, $\nu$ negativ, und $z$ zwischen den Grenzen $x$ und $k$ continuirlich ist; f\"ur einen Werth von $\nu$ aber, der $>0$ ist, bezeichnet $dl^\nu z$ dasjenige, was aus $dl^{\nu-m}z$ (wo $m>\nu$) durch $m$ malige Differentiation nach $x$ hervorgeht\footnote{Die Definition
\[
dl^\nu z = \sum_{n=-\infty}^{\infty} \frac{h^n}{\Pi(n)}dl^{\nu+n}z,
\]
welche mit der gegebenen identisch ist, w\"urde zwar f\"ur alle Werthe von $\nu$ gelten; wir haben ihr aber die gew\"ahlte ihrer gr\"osseren Geschmeidigkeit wegen vorgezogen.}, ein Werth, welcher stets auch der Gleichung
\[
dl^{\nu+n}z = \frac{d^n\,dl^\nu z}{dx^n}
\]
gen\"ugen muss\footnote{Ob die obige Formel~1.\ alle Werthe enth\"alt, die dieser Gleichung gen\"ugen, h\"angt offenbar davon ab, ob die Functionen $\varphi_\nu$ die einzigen sind, welche, statt $dl^\nu z$ substituirt, die Reihe~2.\ zu Null machen. Nun l\"asst sich zwar ohne Schwierigkeit zeigen, dass keine algebraische Function von $x$, die nicht in $\varphi_\nu$ enthalten ist, dies leistet; ob aber \"uberhaupt keine Function dieser Bedingung gen\"ugt, dar\"uber konnte ich bis jetzt zu keinem Resultat gelangen.}.

Hieraus folgt
\begin{enumerate}
\item[$1$.] $dl^\nu z = \Pi(-\nu-1) \int_x^k (x-t)^{-\nu-1}z(t)\,dt + \varphi_\nu$,
\item[$2$.] $dl^\nu z = \sum_{n=-\infty}^{\infty} \frac{h^n}{\Pi(n)} dl^{\nu+n} z$,
\item[$3$.] $dl^\nu(dl^\mu z) = dl^{\nu+\mu}z$,
\item[$4$.] $dl^0 z = z$,
\item[$5$.] $dl^n z = \frac{d^n z}{dx^n}$, wenn $n$ positiv und ganz,
\item[$6$.] $dl^{-n} z = \Pi(n-1) \int_x^k (x-t)^{n-1}z(t)\,dt + \varphi_{-n}$,
\item[$7$.] $dl^\nu(e^{cx}) = c^\nu e^{cx}$,
\item[$8$.] $dl^\nu(cz) = c \cdot dl^\nu z$,
\item[$9$.] $dl^\nu(z_1 + z_2) = dl^\nu z_1 + dl^\nu z_2$,
\item[$10$.] $dl^\nu z \cdot dl^\mu z \neq dl^{\nu+\mu}z$ im Allgemeinen.
\end{enumerate}
Zwei Werthe von $dl^\nu z$ und $dl^\nu z$, in denen die Constanten $K$, $K_1$, etc. s\"ammtlich einander gleich sind, sollen \emph{correspondirende} Werthe heissen. Alle derselben Entwicklung von $z(x+h)$ angeh\"origen Werthe sind correspondirende.

Wir wollen nun zu der Bestimmung der Ableitungen bestimmter Functionen von $x$ \"ubergehen. Dabei kann es nat\"urlich nur darauf ankommen, einen Werth Einer Ableitung zu finden, da sich aus diesem ihr allgemeiner Werth durch Addition der Function $\varphi_\nu$ sofort ergiebt, und zwar wird dieser Werth, wenn die Umformung des Ausdrucks 1.\ \"uberhaupt etwas n\"utzen soll, ein einfacherer, als dieser Ausdruck, also eine explicite Function von $x$ in endlicher Form sein m\"ussen. Diese Umformung wird also im Allgemeinen darin bestehen, dass man das $x$ aus dem Integralzeichen herauszuschaffen sucht.

Betrachten wir nun zuerst die Function $x^\mu$.

Ist $\mu$ positiv, so ist $x^\mu$ f\"ur alle Werthe von $x$ continuirlich; es wird also
\[
\frac{dl^\nu(x^\mu)}{\Pi(-\nu-1)} = \int_x^k (x-t)^{-\nu-1}t^\mu\,dt
\]
immer ein Werth von $dl^\nu(x^\mu)$ sein; dies Integral ist aber
\[
\frac{x^{\mu-\nu}}{\Pi(-\nu-1)} \int_0^{\frac{k}{x}}(1-y)^{-\nu-1}y^\mu\,dy = \frac{\Pi(\mu)}{\Pi(\mu-\nu)}x^{\mu-\nu}
\]
(nach~(\ref{eq:intresult})). Da das $n$te Differential hiervon $\frac{\Pi(\mu)}{\Pi(\mu-\nu-n)}x^{\mu-\nu-n} = dl^{\nu+n}(x^\mu)$ ist (nach~4), so ist f\"ur jeden Werth von $\nu$
\begin{equation}
dl^\nu(x^\mu) = \frac{\Pi(\mu)}{\Pi(\mu-\nu)}x^{\mu-\nu} + \varphi_\nu.
\end{equation}

Ist $\mu$ negativ, so ist $x^\mu$ f\"ur $x=0$ discontinuirlich, f\"ur alle andern Werthe aber continuirlich; in dem Ausdrucke~1.\ m\"ussen also $x$ und $k$ stets gleiches Zeichen haben. Nun erh\"alt man aber durch $m$ malige partielle Integration
\[
\int_x^k (x-t)^{-\nu-1}t^{\mu+m}\,dt = \left[\frac{(x-t)^{-\nu}}{-\nu}t^{\mu+m}\right]_x^k + \frac{\mu+m}{-\nu}\int_x^k(x-t)^{-\nu}t^{\mu+m-1}\,dt + \cdots
\]
so lange $-\nu-m > 0$ ist, wodurch sich also, wenn $-\nu > -\mu$ ist, diejenigen Integrale, worin $\mu < -1$ ist, auf solche zur\"uckf\"uhren lassen, in denen der Exponent von $t > -1$ ist; ist er $> -1$, so geh\"ort $\int_x^k(x-t)^{-\nu-1}t^{\mu+m}\,dt$ zu den Functionen $\varphi_\nu$, und es ist also
\[
\frac{\Pi(\mu+m)}{\Pi(-\nu-1-m)\Pi(\mu+m)}\int_x^k(x-t)^{-\nu-1-m}t^{\mu+m}\,dt = \frac{\Pi(\mu)}{\Pi(\mu-\nu)}x^{\mu-\nu}
\]
ein Werth von $dl^\nu(x^\mu)$, wenn $-\nu > -\mu$, welches Resultat nach dem Gesetze $dl^{\nu+n}z = \frac{d^n\,dl^\nu z}{dx^n}$ f\"ur jedes $\nu$ gelten muss.

Ist aber $\mu+m = -1$, so ist
\[
\int_x^k(x-t)^{-\nu-1}t^{-1}\,dt = \frac{1}{x}\int_0^{\frac{k}{x}}(1-y)^{-\nu-1}y^{-1}\,dy
\]
\[
= \frac{1}{x}\left[\log\frac{x}{k} + \int_0^1\frac{(1-y)^{-\nu-1}-1}{y}\,dy + \int_1^{\frac{k}{x}}\frac{(1-y)^{-\nu-1}}{y}\,dy\right].
\]
\[
= \frac{1}{x}\log\frac{x}{k} \cdot k_1 + \frac{1}{x}\int_0^1\frac{(1-y)^{-\nu-1}-1}{y}\,dy.
\]
Verallgemeinert man auch das hieraus erhaltene Resultat durch Differentiation, so hat man folgende Werthe f\"ur $dl^\nu(x^\mu)$:

\begin{enumerate}
\item[$11$.] $\displaystyle dl^\nu(x^\mu) = \frac{\Pi(\mu)}{\Pi(\mu-\nu)}x^{\mu-\nu}$, wenn $\mu$ nicht eine negative ganze Zahl ist;
\item[$12$.] $\displaystyle dl^\nu(x^\mu) = \frac{\Pi(\mu)}{\Pi(\mu-\nu)}\left[\log x \cdot x^{\mu-\nu} - (\Psi(\mu-\nu) - \Psi(0))x^{\mu-\nu}\right]$, wenn $\mu$ eine ganze negative Zahl ist.
\end{enumerate}
Es ist zu bemerken, dass aus der Formel~12.\ die Formel~11.\ hervorgeht, sobald man nur die Constanten, die f\"ur diesen Fall $\frac{0}{0}$ werden, einer geeigneten Behandlung unterwirft, was auch in dem Falle geschehen muss, wo $(\mu-\nu)$ und $\mu$ beide ganze negative Zahlen sind.

Man \"ubersieht leicht, dass die aus diesen Formeln f\"ur verschiedene Werthe von $\nu$ hervorgehenden Werthe correspondirende sind; dies ist auch der Grund, warum wir in~12.\ nicht, wie wir es f\"ur den Fall $\mu = $ einer negativen ganzen Zahl konnten, den blos $x^{\mu-\nu}$ enthaltenden Theil in die Function $\varphi_\nu$ einschlossen.

Wendet man ein \"ahnliches Verfahren auf $e^x$ an, so erh\"alt man
\begin{enumerate}
\item[$13$.] $\displaystyle dl^\nu(e^x) = e^x$.
\end{enumerate}
Die Ableitungen von $\log x$ ergeben sich durch dieselbe Methode, noch leichter aber und zwar sogleich f\"ur alle Werthe von $\nu$ aus~6.\ und~12.:
\begin{enumerate}
\item[$14$.] $\displaystyle dl^\nu(\log x) = dl^{\nu+1}dl^{-1}\log x = \frac{\Pi(0)}{\Pi(-\nu-1)}\left[\log x - (\Psi(-\nu) - \Psi(0))\right]x^{-\nu-1}$.
\end{enumerate}
Durch Anwendung der Regeln~7 bis~10 findet man aus~13.\ und~14.\ mit der gr\"ossten Leichtigkeit auch die Ableitungen von $\sin x$, $\cos x$, $\tg x$ und $\arctg x$.

Schliesslich bemerken wir noch, dass sich die aufgestellte Theorie mit derselben Sicherheit auch auf den Fall ausdehnen l\"asst, wo man den in Rede stehenden Gr\"ossen imagin\"are Werthe beilegt.

\newpage

\section{Neue Theorie des R\"uckstandes in electrischen Bindungsapparaten}
\addcontentsline{toc}{section}{XX. Neue Theorie des R\"uckstandes in electrischen Bindungsapparaten}

\begin{footnotesize}
\begin{flushright}
(1854.)
\end{flushright}
\end{footnotesize}

\subsection*{Vorbemerkung.}

Herrn Professor Kohlrausch ist es gelungen, die Bildung des R\"uckstandes in electrischen Bindungsapparaten scharfen Messungen zu unterwerfen und darauf eine den Beobachtungen gen\"ugende Theorie dieser Erscheinung zu gr\"unden, welche in Poggendorff's Annalen\footnote{Bd.~91.\ pag.~56.} ver\"offentlicht worden ist. Die Genauigkeit dieser Messungen reizte mich, ein aus andern Gr\"unden wahrscheinliches Gesetz f\"ur die Bewegungen der Electricit\"at an denselben zu pr\"ufen; in der Form, welche ihm f\"ur diesen Zweck gegeben wurde, ist es auf die Bewegungen der Electricit\"at in allen ponderabeln K\"orpern anwendbar, jedoch nur unter der Voraussetzung, dass die in Betracht kommenden ponderabeln K\"orper gegen einander ruhen und keine merklichen thermischen und magnetischen (oder voltinductorischen) Wirkungen und Einfl\"usse stattfinden. Behufs unbeschr\"ankter Anwendbarkeit bedarf es noch einer Umarbeit und Erg\"anzung, mit welcher ich mich an einem andern Orte besch\"aftigen werde.

Im folgenden Aufsatze, welcher einem Schreiben an Herrn Professor Kohlrausch entnommen ist, ist diese neue Theorie des electrischen R\"uckstandes indess nicht selbstst\"andig, sondern im Anschlusse an seine Theorie entwickelt worden; ich war bestrebt, jene Theorie, nicht geradeswegs die Erscheinungen auf sie zur\"uckzuf\"uhren. Ich habe daher die von Herrn Professor Kohlrausch in seiner Abhandlung gebrauchten Begriffe: electrisches Moment der isolirenden Wand, Spannung, Gesammtladung, disponible Ladung, R\"uckstand, \"uberall durch die hier zu Grunde gelegten Begriffe ausgedr\"uckt und auch sonst in mancher Hinsicht die dortige Betrachtungsweise ber\"ucksichtigt.

\subsection*{2. Das der Rechnung zu Grunde gelegte Gesetz.}

Es bezeichne $t$ die Zeit, $x$, $y$, $z$ rechtwinklige Coordinaten, $q$ die Dichtigkeit der Spannungselectricit\"at zur Zeit $t$ im Punkte $(x,y,z)$, $u$ den $4\pi$ten Theil des (Gauss'schen) Potentials aller wirkenden electrischen Massen im Punkte $(x,y,z)$ zur Zeit $t$, also die Gr\"osse
\[
\frac{1}{4\pi} \int \frac{q'\,dx'\,dy'\,dz'}{\sqrt{(x-x')^2 + (y-y')^2 + (z-z')^2}},
\]
wenn $q\,dx'\,dy'\,dz'$ die Spannungselectricit\"at des Elements $dx'\,dy'\,dz'$ zur Zeit $t$ bedeutet. Man hat dann
\[
\frac{d^2u}{dx^2} + \frac{d^2u}{dy^2} + \frac{d^2u}{dz^2} = -q.
\]

Die hier anzuwendenden Gesetze f\"ur die Bewegungen der Electricit\"at im Innern eines homogenen ponderabeln K\"orpers unter den erw\"ahnten Umst\"anden sind nun folgende:

I.\ Die electromotorische Kraft im Punkte $(x,y,z)$ zur Zeit $t$ setzt sich zusammen aus zwei Bestandtheilen, aus einem dem Coulomb'schen Gesetz gem\"assen, dessen Componenten proportional
\[
\frac{\partial u}{\partial x}, \qquad \frac{\partial u}{\partial y}, \qquad \frac{\partial u}{\partial z}
\]
sind, und einem andern, dessen Componenten proportional sind
\[
-\frac{\partial q}{\partial x}, \qquad -\frac{\partial q}{\partial y}, \qquad -\frac{\partial q}{\partial z},
\]
so dass ihre Componenten gleichgesetzt werden k\"onnen
\[
-\frac{\partial u}{\partial x} - \beta\beta\frac{\partial q}{\partial x}, \qquad
-\frac{\partial u}{\partial y} - \beta\beta\frac{\partial q}{\partial y}, \qquad
-\frac{\partial u}{\partial z} - \beta\beta\frac{\partial q}{\partial z},
\]
wo $\beta\beta$ nur von der Natur des ponderabeln K\"orpers abh\"angt.

II.\ Die Stromintensit\"at ist der electromotorischen Kraft proportional, also
\[
\xi = -\alpha\alpha\left(\frac{\partial u}{\partial x} + \beta\beta\frac{\partial q}{\partial x}\right), \quad
\eta = -\alpha\alpha\left(\frac{\partial u}{\partial y} + \beta\beta\frac{\partial q}{\partial y}\right), \quad
\zeta = -\alpha\alpha\left(\frac{\partial u}{\partial z} + \beta\beta\frac{\partial q}{\partial z}\right),
\]
wenn $\alpha$ eine von der Natur des ponderabeln K\"orpers abh\"angige Constante und $\xi$, $\eta$, $\zeta$ die Componenten der Stromintensit\"at sind.

Mit Zuziehung der phoronomischen Gleichung
\[
\frac{\partial q}{\partial t} + \frac{\partial\xi}{\partial x} + \frac{\partial\eta}{\partial y} + \frac{\partial\zeta}{\partial z} = 0
\]
erh\"alt man daher f\"ur $u$ die Gleichungen
\[
\frac{\partial^2 u}{\partial x^2} + \frac{\partial^2 u}{\partial y^2} + \frac{\partial^2 u}{\partial z^2} = -q
\]
und
\[
\frac{\partial q}{\partial t} = \alpha\alpha\left(\frac{\partial^2 u}{\partial x^2} + \frac{\partial^2 u}{\partial y^2} + \frac{\partial^2 u}{\partial z^2}\right) + \alpha\alpha\beta\beta\left(\frac{\partial^2 q}{\partial x^2} + \frac{\partial^2 q}{\partial y^2} + \frac{\partial^2 q}{\partial z^2}\right),
\]
oder, wenn man die L\"ange $\beta$ und die Zeit $\alpha$ zur Einheit nimmt,
\[
q + \frac{\partial q}{\partial t} + \Delta\Delta u = 0, \qquad \text{und} \qquad q + \frac{\partial u}{\partial t} + \Delta\Delta u = 0.
\]
Dies giebt f\"ur $u$ eine partielle Differentialgleichung, welche in Bezug auf $t$ vom ersten, in Bezug auf die Raumcoordinaten vom vierten Grade ist, und um von einem bestimmten Zeitpunkte an $u$ allenthalben im Innern des ponderabeln K\"orpers zu bestimmen, werden ausser dieser Gleichung noch eine Bedingung in jedem Punkte desselben f\"ur die Anfangszeit und f\"ur die Folge in jedem Oberfl\"achenpunkte zwei Bedingungen erforderlich sein.

\subsection*{3. Plausible Auffassung dieses Gesetzes.}

Das Bewegungsgesetz der Electricit\"at ist unter voriger Nummer durch Begriffe, welche jetzt in der Lehre von der Electricit\"at gebr\"auchlich sind, ausgedr\"uckt worden. Diese Auffassung desselben ist jedoch einer Umarbeitung f\"ahig, durch welche, wie es scheint, ein etwas treueres und vollst\"andigeres Bild des wirklichen Zusammenhangs gewonnen wird.

Statt eine Ursache anzunehmen, welche im Punkte $(x,y,z)$ die positive Electricit\"at in den Richtungen der drei Axen mit den Kr\"aften $-\frac{\partial u}{\partial x}$, $-\frac{\partial u}{\partial y}$, $-\frac{\partial u}{\partial z}$ und die negative mit den entgegengesetzten treibt, kann man auch eine Ursache annehmen, welche im Punkte $(x,y,z)$ die positive Electricit\"at mit der Intensit\"at $\beta\beta q$ zu vermindern und die negative zu vermehren strebt, und diese Ursache kann man in einem Widerstreben des Ponderabile gegen das Enthalten von Spannungselectricit\"at oder den electrischen Zustand suchen.

Ebenso kann man auch die electromotorische Kraft, deren Componenten $\frac{\partial u}{\partial x}$, $\frac{\partial u}{\partial y}$, $\frac{\partial u}{\partial z}$ sind, durch eine Ursache von der Intensit\"at $u$ im Punkte $(x,y,z)$ ersetzen, welche die Dichtigkeit der Electricit\"at gleichen Zeichens zu vermindern und die der entgegengesetzten zu vermehren strebt.

Es ist aber dann, um der Gr\"osse $q$ eine reelle Bedeutung zu geben, nicht n\"othig zweierlei Electricit\"aten anzunehmen und $q\,dx\,dy\,dz$ als den Ueberschuss der positiven Electricit\"at des Elements $dx\,dy\,dz$ \"uber die negative zu betrachten, sondern man kann im Wesentlichen zu der Franklin'schen Auffassung der electrischen Erscheinungen zur\"uckkehren, am einfachsten wohl durch folgende Annahme:

Das Ponderabile, welches Sitz der Electricit\"at ist, erf\"ullt den Raum stetig und mit gleichm\"assiger electrischer Capacit\"at, welche seinem Leitungswiderstande umgekehrt proportional ist, und von welcher die Dichtigkeit der wirklich in ihm enthaltenen Electricit\"at immer nur um einen unmerklich kleinen Bruchtheil abweicht. Bei \"ubersch\"ussiger oder fehlender Electricit\"at (positiver oder negativer Spannungselectricit\"at) ger\"ath das Ponderabile in einen positiv oder negativ electrischen Zustand, verm\"oge dessen es die Dichtigkeit der in ihm enthaltenen Electricit\"at zu vermindern oder zu vermehren strebt und zwar mit einem Drucke, welcher gleich ist der Dichtigkeit seiner Spannungselectricit\"at, $q$, multiplicirt in einen von der Natur des Ponderabile abh\"angigen Factor (seine antelectrische Kraft). Ihrerseits ger\"ath bei auftretender Spannungselectricit\"at die Electricit\"at in einen Zustand, Spannung, verm\"oge dessen sie ihre Dichtigkeit zu vermindern (oder bei negativer Spannung zu vermehren) strebt und dessen Gr\"osse $u$ in jedem Augenblicke abh\"angt von s\"ammtlichen Massen Spannungselectricit\"at nach der Formel
\[
\frac{1}{4\pi}\int \frac{q'\,dx'\,dy'\,dz'}{\sqrt{(x-x')^2 + (y-y')^2 + (z-z')^2}}
\]
oder auch vermittelst des Gesetzes
\[
\frac{\partial^2 u}{\partial x^2} + \frac{\partial^2 u}{\partial y^2} + \frac{\partial^2 u}{\partial z^2} = -q
\]
und der Bedingung, dass $u$ in unendlicher Entfernung von Spannungselectricit\"at unendlich klein bleibt. Die Electricit\"at bewegt sich gegen die ponderabeln K\"orper mit einer Geschwindigkeit, welche in jedem Augenblicke der aus diesen Ursachen hervorgehenden electromotorischen Kraft gleich ist.

Uebrigens m\"ussen diese Bewegungsgesetze der Electricit\"at, wenn deren Verh\"altniss zu W\"arme und Magnetismus in Rechnung gezogen werden soll, vorbemerktermassen selbst noch abge\"andert und umgeformt werden, und dann wird eine ver\"anderte Auffassung dieser Erscheinungen n\"othig.

\newpage

\subsection*{4. Behandlung des Problems der R\"uckstandsbildung. Ausdruck der zu bestimmenden Gr\"ossen durch das Potential.}

Indem ich mich nun zur Untersuchung der R\"uckstandsbildung wende, besch\"aftige ich mich zun\"achst damit, die zu bestimmenden Gr\"ossen durch das Potential, oder vielmehr, was die Rechnung vereinfacht, durch die ihm proportionale Function $u$ auszudr\"ucken. Zu gr\"osserer Bequemlichkeit f\"ur die an abstracte Gr\"ossenbetrachtung minder gew\"ohnten Physiker habe ich das Potential als das Maass einer Ursache, Spannung, betrachtet, welche die Dichtigkeit der Electricit\"at im Punkte $(x,y,z)$ zu vermindern strebt, und diese im Punkte $(x,y,z) = u$, also die Componenten der durch sie bewirkten electromotorischen Kraft
\[
\frac{\partial u}{\partial x}, \qquad \frac{\partial u}{\partial y}, \qquad \frac{\partial u}{\partial z}
\]
gesetzt. Man muss dann als Spannungseinheit die im Innern einer Kugel vom Radius~1 durch auf der Oberfl\"ache vertheilte Electricit\"at von der Dichtigkeit~1 entstehende Spannung annehmen oder als Einheit der electromotorischen Kr\"afte die von der Masse $4\pi$ in der Entfernungseinheit erzeugte.

Zur Vereinfachung der Rechnung ist ferner als Zeiteinheit $\alpha$, als L\"angeneinheit $\beta$ eingef\"uhrt worden; macht man die Einheit der electromotorischen Kr\"afte auf die hier angenommene Weise von der electrischen Masseneinheit abh\"angig, so sind $\alpha$ und $\beta\beta$ die Maasse f\"ur den Leitungswiderstand $\left(\frac{\text{electromotorische Kraft}}{\text{Stromintensit\"at}}\right)$ und die antelectrische Kraft $\left(\frac{\text{Druck des Ponderabile}}{\text{Dichtigkeit der Spannungselectricit\"at}}\right)$ des ponderabeln Sitzes.

Zur Discussion der vorliegenden Beobachtungen gen\"ugt die L\"osung der Aufgabe: die Aenderungen der Spannungselectricit\"at im Innern einer \"uberall gleich dicken homogenen Wand zu bestimmen, wenn die Oberfl\"achen mit vollkommenen Leitern belegt sind, gleiche Mengen entgegengesetzter Electricit\"at empfangen und keine electromotorische Kraft besitzen (keine Contactwirkung in ihnen stattfindet), und ihre Dimensionen gegen die Dicke der Wand als unendlich gross betrachtet werden d\"urfen (d.\ h.\ der Einfluss des Randes und der Kr\"ummung vernachl\"assigt werden darf).

Legt man den Anfangspunkt der Coordinaten in die Mitte der Wand, die $x$-Axe auf ihre Oberfl\"achen senkrecht und bezeichnet ihre halbe Dicke durch $a$, so wird der Ausdruck f\"ur die Wand $a > x > -a$, $u$ eine blosse Function von $x$ und folglich
\[
q = -\frac{d^2u}{dx^2}.
\]
Die zwischen zwei Werthen von $x$ \"uber der Fl\"acheneinheit enthaltene Electricit\"atsmenge ist also, geometrisch ausgedr\"uckt, gleich der Differenz zwischen den Tangenten der Neigungen der Spannungscurve, d.\ h.\ der Curve, deren Ordinate f\"ur die Abscisse $x$ gleich $u$ ist; diese Curve ist gerade, wo keine Spannungselectricit\"at vorhanden ist, nach oben (oder f\"ur Orte mit gr\"osseren Ordinaten) convex, wo positive, nach unten, wo negative stetig vertheilt ist, und gebrochen f\"ur einen Werth von $x$, bei welchem eine endliche Menge angeh\"auft ist.

Die durch eine Ladung erzeugte oder durch eine Entladung vernichtete Spannung wird daher stets dargestellt durch eine Curve von der Form~A, d.\ h.\ ist sie in den Belegungen $u_a$, $u_{-a}$ und folglich in der Mitte $=u_0$, so ist sie im Innern
\[
u = u_0 + \frac{x}{a}(u_a - u_0).
\]

Durch das Eindringen der Electricit\"at ins Innere erh\"alt die Spannungscurve die Form~B. F\"ur die Fl\"acheneinheit ist die Gesammtmenge der geschiedenen Electricit\"aten gleich der Tangente ihrer Neigung in der Mitte
\[
\left(\frac{du}{dx}\right)_{\!0};
\]
das electrische Moment
\[
\int_{-a}^{+a} qx\,dx = u_a - u_{-a} - a\left(\left(\frac{du}{dx}\right)_{\!a} + \left(\frac{du}{dx}\right)_{\!-a}\right) = u_a - u_{-a},
\]
also gleich der Spannungsdifferenz der Oberfl\"achen.

Durch eine Entladung wird die Spannung in den Belegungen aufgehoben. Die vernichtete Spannung ist daher in den Belegungen $=u_a$, $u_{-a}$, im Innern
\[
= u_0 + \frac{x}{a}(u_a - u_0),
\]
die disponible Ladung f\"ur die Fl\"acheneinheit
\[
= -(u_a - u_0),
\]
die bleibende Spannung im Innern
\[
= u - u_0 - \frac{x}{a}(u_a - u_0),
\]
und f\"ur die Fl\"acheneinheit der verborgene R\"uckstand
\[
A = -\frac{du}{dx} + \frac{u_a - u_0}{a}.
\]

\begin{center}
\begin{tabular}{ccc}
B & C & D \\
\multicolumn{3}{c}{1) Spannungscurve der Gesammtladung} \\
\multicolumn{3}{c}{2) \quad"\quad\quad der disponiblen Ladung} \\
\multicolumn{3}{c}{3) \quad"\quad\quad des R\"uckstandes.} \\
\end{tabular}
\end{center}
Gesammtladung: $=ac$; disponible Ladung: $=ab$; R\"uckstand: $=bc$.

\newpage

\subsection*{5. L\"osung der Aufgabe im einfachsten Falle, wo kein Ab- und Zufluss durch die Oberfl\"achen stattfindet.}

Nach dieser Uebersicht und geometrischen Darstellung der gesuchten Gr\"ossen gehe ich zu ihrer Bestimmung durch Rechnung nach dem angegebenen Gesetze \"uber. Ich behandle zun\"achst den Fall, wo anfangs im Innern keine freie Electricit\"at vorhanden ist, und den Oberfl\"achen auf der Fl\"acheneinheit die Masseneinheit mitgetheilt wird, sp\"ater aber kein Ab- und Zufluss durch die Oberfl\"achen stattfindet.

Die Bedingungen zur Bestimmung von $u$ sind:

f\"ur $t > 0$, $a > x > -a$:
\[
\frac{d^4u}{dx^4} + \frac{d^2u}{dt^2} + \frac{dq}{dt} = 0, \qquad q = -\frac{d^2u}{dx^2};
\]

$f = 0$, $a > x > -a$:
\[
\frac{du}{dx} = 0;
\]

$t > 0$, $x = \pm a$:
\[
\frac{du}{dx} = 0, \qquad \frac{d^3u}{dx^3} + \frac{dq}{dx} = 0,
\]
welche letzteren ausdr\"ucken, dass in den Oberfl\"achen sowohl die Electricit\"atsmengen, als der Durchfluss, und folglich die electromotorische Kraft $=0$ sein soll.

Diesen Bedingungen gen\"ugen zwei Ausdr\"ucke, der eine f\"ur kleine, der andere f\"ur grosse Werthe von $t$ brauchbar.

Setzt man zur Abk\"urzung
\[
\int_0^\infty e^{-\lambda^2}\,d\lambda = \varphi(\lambda)
\]
und
\[
\chi(\lambda) = \frac{2}{\sqrt{\pi}}\varphi(\lambda),
\]
so gen\"ugt erstens
\[
u - u_a = -\left[1 + \frac{x}{a}\chi\left(\frac{a+x}{2\sqrt{t}}\right) + \frac{a-x}{2a}\chi\left(\frac{a-x}{2\sqrt{t}}\right)\right],
\]
zweitens
\[
u - u_a = e^{-t}\sqrt{\frac{2}{\pi}}\int_0^\infty \frac{\cos(x\lambda)}{\cos(a\lambda)}e^{-\lambda^2t}\,d\lambda - \cdots
\]
Die hieraus sich ergebenden Bestimmungen sind:

f\"ur die Vertheilung der Electricit\"at\footnote{Vergl.\ Jacobi.\ Fundamenta nova theoriae functionum ellipticarum. \S\S.~61,~63.}
\[
q = \frac{2}{a}\sum_{n=1}^{\infty}(-1)^{n-1}e^{-\frac{(2n-1)^2\pi^2t}{4a^2}}\cos\left(\frac{(2n-1)\pi x}{2a}\right);
\]
f\"ur die Gesammtladung
\[
L_t = \frac{4}{\pi}\sum_{n=1}^{\infty}\frac{(-1)^{n-1}}{2n-1}e^{-\frac{(2n-1)^2\pi^2t}{4a^2}};
\]
f\"ur die disponible Ladung
\[
D_t = \frac{2}{a}\sum_{n=1}^{\infty}\frac{(-1)^{n-1}}{\lambda_n}e^{-\lambda_n^2t}\cos(\lambda_n a);
\]
f\"ur den R\"uckstand
\[
R_t = \frac{1}{a}\int_0^a q\,dx \cdot L_t = a^{-1}\sqrt{\frac{2}{\pi}}\sum_{n=1}^{\infty}\left[(-1)^{n-1} - \frac{2}{(2n-1)\pi}\right]e^{-\frac{(2n-1)^2\pi^2t}{4a^2}}\cdot\sin\frac{(2n-1)\pi}{2}.
\]

\subsection*{6. Zur\"uckf\"uhrung der allgemeinen Aufgabe auf diesen einfachsten Fall.}

Um auf diesen einfachsten Fall den Fall zur\"uckzuf\"uhren, wo Ab- und Zufluss durch die Oberfl\"achen stattfindet, bezeichne $\chi(t)$ den Ausdruck f\"ur die Spannungsdifferenz $u - u_0$ zur Zeit $t$ in diesem einfachsten Falle; f\"ur negative Werthe von $t$ sei $\chi(t) = 0$.

Soll nun die Spannung bestimmt werden, welche entsteht, wenn den Oberfl\"achen ($x = \pm a$) zur Zeit $0$ die Mengen $\pm f$, darauf zur Zeit $t'$ die Mengen $\pm f'$, zur Zeit $t''$ die Mengen $\pm f''$, \ldots\ mitgetheilt werden, so hat man
\[
u - u_0 = f\chi(t) + f'\chi(t-t') + f''\chi(t-t'') + \cdots;
\]
denn dieser Werth gen\"ugt s\"ammtlichen zu seiner Bestimmung gegebenen Bedingungen.

Findet ein stetiger Ab- und Zufluss von Electricit\"at statt, so wird
\[
u - u_0 = \int_0^t \chi(t-\tau)\,df,
\]
wenn $\pm df$ die im Zeitelement $d\tau$ durch die Oberfl\"ache ($x = \pm a$) nach Innen str\"omende Electricit\"atsmenge bezeichnet.

Beide Ausdr\"ucke kann man zusammenfassen in dem Ausdruck
\[
u - u_0 = \int_0^t \chi(t-\tau)\,df,
\]
wenn man durch $\pm df$ die im Zeitelement $d\tau$ auf der Oberfl\"ache ($x = \pm a$) hinzukommende Electricit\"atsmenge bezeichnet, wo diese dann einen endlichen Werth hat oder $d\tau$ proportional ist, je nachdem eine pl\"otzliche Ladung oder Entladung, oder ein stetiger Ab- oder Zufluss stattfindet.

Aus diesem Ausdrucke f\"ur die Spannung folgt
\[
q_t = \int_0^t q_{t-\tau}\,df, \qquad L_t = \int_0^t L_{t-\tau}\,df, \qquad R_t = \int_0^t R_{t-\tau}\,df.
\]
In diesen Formeln sind die Zeiten in Theilen von $\alpha$, die L\"angen in Theilen von $\beta$ ausgedr\"uckt; um bekannte Maasse einzuf\"uhren, hat man nur $a$ und $x$ durch $\frac{a}{\beta}$, $\frac{x}{\beta}$, $t$ und $\tau$ durch $\frac{t}{\alpha}$, $\frac{\tau}{\alpha}$ zu ersetzen.

\subsection*{7. Vergleichung der Rechnung mit den Beobachtungen.}

Um nun die erhaltenen Formeln mit dem wirklichen Verlaufe der R\"uckstandsbildung zu vergleichen, wie er durch die in Poggendorff's Annalen ver\"offentlichten Messungen des Herrn Professor Kohlrausch mit so grosser Genauigkeit festgestellt worden ist, geht man wohl am zweckm\"assigsten von der Thatsache aus, dass die Ladungscurve einer Parabel nahe kommt mit allm\"ahlich abnehmendem Parameter, d.\ h.\ dass die Gr\"osse $\frac{L_0 - L_t}{\sqrt{t}}$ langsam abnimmt.

Zufolge der f\"ur $L_t$ abgeleiteten Formel ist $L_0 - L_t$ f\"ur sehr kleine Werthe von $t$ proportional $\sqrt{t}$ und zwar
\[
\frac{L_0 - L_t}{\sqrt{t}} = \frac{2}{a\sqrt{\pi}}.
\]
Zufolge der Messungen muss man annehmen, dass diese Proportionalit\"at n\"aherungsweise noch w\"ahrend der Beobachtungen stattfindet.

Man wird daher die Zeit $\frac{a^2}{\beta\beta}$ in roher Ann\"aherung aus den Beobachtungen bestimmen k\"onnen, und dann ist in der That
\[
\frac{L_t}{L_0 - L_t}\sqrt{t}
\]
eine Function, welche mit wachsendem $t$ langsam abnimmt. Nichtsdestoweniger w\"urde $\frac{L_0 - L_t}{\sqrt{t}}$ mit wachsendem $t$ zunehmen, wenn man $\frac{a^2}{\beta\beta}$ einen merklichen Werth beilegte. Dasselbe scheint sich auch zu ergeben, wenn man einen betr\"achtlichen Verlust durch die Luft annimmt, wenigstens wenn man daf\"ur das Coulomb'sche Gesetz\footnote{Pogg.\ Ann.\ Bd.~71.\ pag.~359; Bd.~72.\ pag.~374.} zu Grunde legt.

Man wird daher f\"ur die erste Bearbeitung der Beobachtungen die Zeit $\alpha$ (d.\ h.\ den Leitungswiderstand des Glases f\"ur die dem Coulomb'schen Gesetz gem\"assen electromotorischen Kr\"afte) unendlich gross annehmen, den Verlust durch die Luft vernachl\"assigen und sich zun\"achst darauf beschr\"anken m\"ussen, zu untersuchen, in wie weit sich durch geh\"orige Bestimmung von $\frac{a^2}{\beta\beta}$ den Beobachtungen gen\"ugen l\"asst.

Sobald man sich \"uberzeugt hat, dass die Voraussetzungen der Rechnung n\"aherungsweise richtig sind, ist eine sch\"arfere Vergleichung der Rechnung mit den Beobachtungen verlorene Arbeit, wenn man nicht die Gelegenheit hat, die Quellen der Differenzen zwischen Rechnung und Beobachtung an der Hand der Erfahrung aufzusuchen, um die wegen der Abweichungen von den Voraussetzungen der Rechnung n\"othigen Correctionen anzubringen. Da mir nun zu einem experimentellen Studium des Gegenstandes die Mittel fehlen, so musste ich von einer weiteren Verfolgung desselben vorl\"aufig abstehen.

\subsection*{Verh\"altniss dieses Problems zur Electrometrie und zur Theorie verwandter Erscheinungen.}

Die Gr\"osse $\frac{\beta\beta}{a}$, bei der Flasche etwa $\frac{1}{2000}$, giebt den Quotienten $\frac{\text{antelectrische Kraft}}{\text{Leitungswiderstand}}$ des Glases der Flasche in absolutem Maass, wenn als L\"angeneinheit die Flaschendicke, als Zeiteinheit die Secunde angenommen wird. F\"ur diese Bestimmung ist es gleichg\"ultig, wie man die Einheit der electromotorischen Kr\"afte von der Einheit der electrischen Massen abh\"angig macht; die Constanten $\alpha$ und $\beta\beta$ w\"urden aber den Leitungswiderstand und die antelectrische Kraft in einem andern Maasse als dem Weber'schen geben, wo die Einheit der electromotorischen Kr\"afte durch die dem Ampere'schen Gesetz gem\"assen Wirkungen der Masseneinheit festgesetzt wird.

Zur Vergleichung des hier untersuchten Falles mit den Erscheinungen an guten Leitern kann die Betrachtung des Beharrungszustandes bei constant erhaltener Spannungsdifferenz der Oberfl\"achen (oder constantem Zufluss) dienen. F\"ur diesen ist

die Dichtigkeit im Innern: $q = -\dfrac{d^2u}{dx^2} = e^x - e^{-x}$,

die Spannung: $u = u_0 - e^x + e^{-x} + x(e^x + e^{-x})$,

die Spannungsdifferenz der Oberfl\"achen:
\[
u_a - u_{-a} = 2\left(a(e^a + e^{-a}) - (e^a - e^{-a})\right),
\]
die Gesammtladung: $\left(\dfrac{du}{dx}\right)_{\!0} = e^a + e^{-a} - 2$,

der R\"uckstand: $\left(\dfrac{du}{dx}\right)_{\!0} - \dfrac{u_a - u_{-a}}{2a} = 2 - \dfrac{e^a - e^{-a}}{a}$,

die in der Zeiteinheit durchfliessende Menge:
\[
-\frac{du}{dx}\bigg|_{a} = \text{const.},
\]
oder gleich proportionalen Gr\"ossen, wobei zur Vereinfachung, wie oben, als Zeiteinheit $\alpha$, als L\"angeneinheit $\beta$, als Spannungseinheit die Spannung im Innern einer Kugel vom Radius~1 bei auf der Oberfl\"ache vertheilter Electricit\"at von der Dichtigkeit~1 angenommen ist.

Besonders wichtig scheint mir die Pr\"ufung des vermutheten Gesetzes und eventuell die Bestimmung der Constanten $\alpha$ und $\beta$ bei den Gasen zu sein. Die Beobachtungen von Riess\footnote{Pogg.\ Ann.\ Bd.~71.\ pag.~359.} und Kohlrausch\footnote{Pogg.\ Ann.\ Bd.~72.\ pag.~374.}, nach welchen f\"ur den Electricit\"atsverlust an die Luft in einem geschlossenen Raume das Gesetz Coulomb's nicht gilt, k\"onnen vielleicht als Ausgangspunkt f\"ur diese Untersuchung dienen und es w\"are f\"ur dieselben wohl zun\"achst ein System von Messungen \"uber den Electricit\"atsverlust im Innern eines einigermassen regelm\"assigen geschlossenen Raumes zu w\"unschen.

\newpage

\section{Zwei allgemeine S\"atze \"uber lineare Differentialgleichungen mit algebraischen Coefficienten}
\addcontentsline{toc}{section}{XXI. Zwei allgemeine S\"atze \"uber lineare Differentialgleichungen mit algebraischen Coefficienten}

\begin{footnotesize}
\begin{flushright}
(20.\ Febr.\ 1857.)
\end{flushright}
\end{footnotesize}

Bekanntlich l\"asst sich jede L\"osung einer linearen homogenen Differentialgleichung $n$ter Ordnung in $n$ von einander unabh\"angige particulare L\"osungen linear mit constanten Coefficienten ausdr\"ucken. Sind die Coefficienten der Differentialgleichung rationale Functionen der unabh\"angigen Ver\"anderlichen $x$, so wird jeder Zweig der, allgemein zu reden, vielwerthigen Functionen, welche ihr gen\"ugen, sich linear mit constanten Coefficienten in $n$ f\"ur jeden Werth von $x$ eindeutig bestimmte Functionen ausdr\"ucken lassen, welche freilich dann l\"angs eines gewissen Liniensystems unstetig sein m\"ussen. Sind die Coefficienten aber algebraische Functionen von $x$, welche sich rational in $X$ und eine $\mu$-werthige algebraische Function von $x$ ausdr\"ucken lassen, so geh\"ort zu jedem Zweige dieser $\mu$-werthigen Function eine Gruppe von $n$ von einander unabh\"angigen particularen L\"osungen, so dass in diesem Falle jeder Zweig einer L\"osung der Differentialgleichung als ein linearer Ausdruck von h\"ochstens $\mu n$ eindeutigen Functionen sich darstellen l\"asst, welcher aber von ihnen immer nur $n$ einer Gruppe angeh\"orige enthalten wird. Aus diesen Vorbemerkungen wird man, da sich jede nicht homogene lineare Differentialgleichung leicht in eine homogene von der n\"achst h\"ohern Ordnung verwandeln l\"asst, ersehen, dass die folgenden S\"atze alle linearen Differentialgleichungen mit algebraischen Coefficienten umfassen.

Es seien $y_1$, $y_2$, \ldots, $y_n$ Functionen von $x$, welche f\"ur alle complexen Werthe dieser Gr\"osse ein\"andrig und endlich sind, ausser f\"ur $a$, $b$, $c$, \ldots, $g$, und welche durch einen Umlauf des $x$ um einen dieser Verzweigungswerthe in lineare Functionen mit constanten Coefficienten von ihren fr\"heren Werthen \"ubergehen.

Zu ihrer n\"aheren Bestimmung scheide man die Gesammtheit der complexen Werthe in zwei Gebiete durch eine in sich zur\"ucklaufende Linie, die der Reihe nach durch s\"ammtliche Verzweigungswerthe $(g, \ldots, c, b, a)$ geht, so dass in jedem dieser Gebiete die Functionen v\"ollig gesondert und stetig verlaufen, und betrachte die Werthe der Functionen in dem auf der positiven Seite dieser Linie liegenden Gebiete als gegeben. Durch einen positiven Umlauf des $x$ um $a$ gehe nun $y_1$ in $\sum A_i^{(1)}y_i$, $y_2$ in $\sum A_i^{(2)}y_i$, \ldots, $y_n$ in $\sum A_i^{(n)}y_i$ \"uber und \"ahnlich durch einen positiven Umlauf um $b$: $y_1$ in $\sum B_i^{(1)}y_i$, etc., durch einen positiven Umlauf um $g$: $y_1$ in $\sum G_i^{(1)}y_i$.

Bezeichnet man nun zur Abk\"urzung das System der $n$ Werthe $(y_1, y_2, \ldots, y_n)$ durch $(y)$, das System der $nn$ Coefficienten
\[
\begin{array}{cccc}
A_1^{(1)}, & A_2^{(1)}, & \ldots, & A_n^{(1)} \\
\vdots & \vdots & & \vdots \\
A_1^{(n)}, & A_2^{(n)}, & \ldots, & A_n^{(n)}
\end{array}
\]
durch $(A)$, das System der $B$ durch $(B)$, \ldots, der $G$ durch $(G)$, und die aus $(y)$ mittelst des Coefficientensystems $(A)$ gebildeten Werthe $\sum A_i^{(1)}y_i$, $\sum A_i^{(2)}y_i$, \ldots, $\sum A_i^{(n)}y_i$ durch $(A)(y_1, y_2, \ldots, y_n) = (A)(y)$, so findet zwischen diesen Coefficientensystemen die Gleichung
\begin{equation}
(G)(F)\cdots(B)(A) = (0)
\end{equation}
statt, wenn man durch $(0)$ ein Coefficientensystem bezeichnet, das nichts \"andert, oder in welchem die Coefficienten der abw\"arts nach rechts gehenden Diagonale $=1$ und alle \"ubrigen $=0$ sind. In der That, durchl\"auft $x$ die ganze Grenzlinie so, dass es sich von einem Verzweigungswerth zum folgenden auf der positiven Seite bewegt, dann aber jedesmal um diesen Verzweigungswerth positiv herum, so gehen die Functionen $(y)$ nach und nach in $(G)(y)$, $(G)(F)(y)$, schliesslich in $(G)(F)\cdots(B)(A)(y)$ \"uber. Es hat aber denselben Erfolg, wenn $x$ die negative Seite der Grenzlinie oder die ganze Begrenzung des negativ erseits liegenden Gebiets durchl\"auft, wobei $(y_1, y_2, \ldots, y_n)$ ihre fr\"heren Werthe wieder annehmen m\"ussen, da sie in diesem Gebiet allenthalben ein\"andrig sind.

Ein System von $n$ Functionen, welches die eben angegebenen Eigenschaften hat, werde durch
\[
\left(\begin{array}{ccccc}
a & b & c & \ldots & g \\
(A) & (B) & (C) & \ldots & (G)\\
\multicolumn{5}{c}{y}
\end{array}\right)
\]
bezeichnet.

Man betrachte nun als zu einer \emph{Klasse} geh\"orig s\"ammtliche Systeme, f\"ur welche die Verzweigungswerthe und die um sie stattfindenden Substitutionen gegebene der Gleichung~(1) gen\"ugende Werthe haben, was, wie sich bald ergeben wird, f\"ur unendlich viele Systeme der Fall ist. Nach einem leicht zu beweisenden, von Jacobi vielfach angewandten Satze l\"asst sich jede Substitution, allgemein zu reden, in drei Substitutionen zerlegen, von denen die letzte die inverse der ersten ist, und in der mittleren die Coefficienten ausser der Diagonale s\"ammtlich $=0$ sind, so dass durch sie jede von den Gr\"ossen, auf welche sie angewandt wird, nur einen Factor erh\"alt. Es l\"asst sich also z.~B.
\[
(A) = (a)
\begin{pmatrix}
\lambda_1 & 0 & \ldots & 0 \\
0 & \lambda_2 & \ldots & 0 \\
\vdots & \vdots & \ddots & \vdots \\
0 & 0 & \ldots & \lambda_n
\end{pmatrix}
(a)^{-1}
\]
setzen, wenn $(a)^{-1}$ die inverse Substitution von $(a)$ bezeichnet. Die Gr\"ossen $\lambda$ werden dabei die $n$ Wurzeln einer durch $(A)$ v\"ollig bestimmten Gleichung $n$ten Grades. F\"ur den Fall, dass diese Gleichung gleiche Wurzeln h\"atte, m\"usste man der mittleren Substitution eine etwas abge\"anderte Form geben; wir wollen aber zur Vereinfachung diesen Fall vorl\"aufig ausschliessen und annehmen, dass er bei der Zerlegung der Substitutionen $(A)$, $(B)$, \ldots, $(G)$ nicht eintritt. Die Substitution $(a)$ kann in
\[
(a) =
\begin{pmatrix}
1 & 0 & \ldots & 0 \\
0 & 1 & \ldots & 0 \\
\vdots & \vdots & \ddots & \vdots \\
0 & 0 & \ldots & 1
\end{pmatrix}
\]
durch Hinzuf\"ugung einer nur multiplicirenden Substitution verwandelt werden; in dieser Form aber sind, wie die Gleichungen, durch welche sie bestimmt wird, zeigen, alle m\"oglichen Werthe derselben enthalten.

Durch einen positiven Umlauf des $x$ um $a$ gehen die Werthe der Functionen $y$ aus $(y_1, y_2, \ldots, y_n)$ in $(A)(y)$ \"uber. Die Werthe der durch die Substitution $(a)^{-1}$ aus $(y)$ gebildeten Functionen
\[
(z_1, z_2, \ldots, z_n) = (a)^{-1}(y)
\]
gehen daher aus $(a)^{-1}(y)$ in
\[
(a)^{-1}(A)(y) =
\begin{pmatrix}
\lambda_1 & 0 & \ldots & 0 \\
0 & \lambda_2 & \ldots & 0 \\
\vdots & \vdots & \ddots & \vdots \\
0 & 0 & \ldots & \lambda_n
\end{pmatrix}
(a)^{-1}(y)
\]
\"uber, oder $(z_1, z_2, \ldots, z_n)$ in $(\lambda_1 z_1, \lambda_2 z_2, \ldots, \lambda_n z_n)$.

Wenn eine Function $z$ durch einen positiven Umlauf des $x$ um $a$ den Constanten Factor $\lambda$ erh\"alt, so kann sie durch Multiplication mit einer Potenz von $(x-a)$ in eine Function verwandelt werden, die in der Umgebung von $a$ ein\"andrig ist. In der That erh\"alt $(x-a)^\mu$ durch einen positiven Umlauf des $x$ um $a$ den Factor $e^{-2\pi i\mu}$; bestimmt man also $\mu$ so, dass $e^{-2\pi i\mu} = \lambda$, oder setzt man $\mu = \frac{\log\lambda}{-2\pi i}$, so wird $z(x-a)^{-\mu}$ eine f\"ur $x=a$ ein\"andrig e Function. Diese Function l\"asst sich also nach ganzen Potenzen von $(x-a)$ entwickeln, und $z$ selbst nach Potenzen, die sich von $\mu$ um ganze Zahlen unterscheiden.

Demnach sind $z_1$, $z_2$, \ldots, $z_n$ nach Potenzen von $x-a$ entwickelbar, deren Exponenten in der Form
\[
\mu_1 + \frac{m}{2\pi i}, \quad \mu_2 + \frac{m}{2\pi i}, \quad \ldots, \quad \mu_n + \frac{m}{2\pi i}
\]
enthalten sind, wenn $m$ eine ganze Zahl bedeutet. Wir wollen nun annehmen, dass die Functionen $y$ nirgends unendlich von unendlich grosser Ordnung werden, so dass diese Reihen auf der Seite der fallenden Potenzen abbrechen m\"ussen, und bezeichnen durch $\mu_1$, $\mu_2$, \ldots, $\mu_n$ die niedrigsten Potenzen in diesen Reihen, so dass
\[
z_1(x-a)^{-\mu_1}, \quad z_2(x-a)^{-\mu_2}, \quad \ldots, \quad z_n(x-a)^{-\mu_n}
\]
endliche von $0$ verschiedene Werthe haben. Offenbar kann die Differenz zweier von den Gr\"ossen $\mu_1$, $\mu_2$, \ldots, $\mu_n$ nie eine ganze Zahl sein, da die Werthe der Gr\"ossen $\lambda_1$, $\lambda_2$, \ldots, $\lambda_n$ s\"ammtlich von einander verschieden sind; dagegen werden die Werthe der entsprechenden Exponenten bei zwei zu derselben Klasse geh\"origen Systemen sich nur um ganze Zahlen unterscheiden k\"onnen, da die Gr\"ossen $\lambda_1$, $\lambda_2$, \ldots, $\lambda_n$ durch $(A)$ v\"ollig bestimmt sind. Diese Exponenten k\"onnen dazu dienen, die verschiedenen Functionensysteme derselben Klasse von einander zu unterscheiden, oder doch sie zu gruppiren, und es gen\"ugt, wenn sie bekannt sind, statt $(A)$ die Substitution $(a)$ anzugeben, da die Gr\"ossen $\lambda_1$, $\lambda_2$, \ldots, $\lambda_n$ schon durch sie bestimmt sind: wir werden uns daher zur genaueren Charakteristik des Systems $(y_1, y_2, \ldots, y_n)$ des Ausdrucks
\[
\left(\begin{array}{cccc|cccc}
(a) & (b) & \ldots & (g) & \mu_1 & \nu_1 & \ldots & \varrho_1 \\
&&&& \mu_2 & \nu_2 & \ldots & \varrho_2 \\
&&&& \vdots & \vdots & & \vdots \\
&&&& \mu_n & \nu_n & \ldots & \varrho_n
\end{array}\right)x
\]
bedienen, in welchem die Gr\"ossen der \"ubrigen Verticalreihen f\"ur die Verzweigungswerthe $b$, \ldots, $g$ die analoge Bedeutung haben sollen, wie die der ersten f\"ur $a$. Es liegt dabei auf der Hand, dass jedes System als ein specieller Fall eines andern betrachtet werden kann, in welchem die entsprechenden Exponenten zum Theil oder s\"ammtlich niedriger sind.

Es ist nun nicht schwer zu beweisen, dass zwischen je $n+1$ Systemen, die derselben Klasse angeh\"oren, eine lineare homogene Gleichung mit ganzen Functionen von $x$ als Coefficienten stattfindet. Wir unterscheiden die entsprechenden Gr\"ossen in diesen $n+1$ Systemen durch obere Indices. Nehmen wir an, dass zwischen ihnen die $n$ Gleichungen stattfinden:
\[
a_0 y_n + a_1 y_n^{(1)} + \cdots + a_n y_n^{(n)} = 0,
\]
so m\"ussen die Gr\"ossen $a_0$, $a_1$, \ldots, $a_n$ proportional sein den Determinanten der Systeme, welche man erh\"alt, wenn man in dem Systeme der $n(n+1)$ Gr\"ossen $y$ der Reihe nach die $1$te, $2$te, \ldots, $(n+1)$te Verticalreihe wegl\"asst. Eine solche Determinante $\sum \pm y_1^{(i_1)} y_2^{(i_2)} \cdots y_n^{(i_n)}$ erh\"alt durch einen positiven Umlauf des $x$ um $a$ den Factor $\det(A)$ und kann f\"ur $x=a$ nicht unendlich von unendlich grosser Ordnung werden; sie l\"asst sich also nach um $1$ steigenden Potenzen von $x-a$ entwickeln. Um den niedrigsten Exponenten in dieser Entwicklung zu bestimmen, kann diese Determinante in die Form gesetzt werden $\det(a) \cdot \sum \pm z_1^{(i_1)} z_2^{(i_2)} \cdots z_n^{(i_n)}$. In letzterer Determinante ist das erste Glied multiplicirt in eine Function, die f\"ur $x=a$ einen endlichen und von $0$ verschiedenen Werth hat. Der niedrigste Exponent in der Entwicklung dieses Gliedes nach Potenzen von $(x-a)$ ist daher
\[
\mu_1^{(i_1)} + \mu_2^{(i_2)} + \cdots + \mu_n^{(i_n)},
\]
und hieraus erh\"alt man durch Permutation der oberen Indices die niedrigsten Exponenten in den Entwicklungen der \"ubrigen Glieder. Offenbar ist der gesuchte Exponent, allgemein zu reden, gleich dem kleinsten von diesen Werthen und jedenfalls nicht kleiner. Bezeichnen wir den kleinsten dieser Werthe durch $\mu$, den \"ahnlichen Werth f\"ur den zweiten Verzweigungswerth durch $\nu$, \ldots, f\"ur den letzten durch $\varrho$, so ist
\[
(x-a)^\mu(x-b)^\nu\cdots(x-g)^\varrho \cdot \det(y_1^{(i_1)} \cdots y_n^{(i_n)})
\]
eine Function von $x$, welche f\"ur alle endlichen complexen Werthe ein\"andrig und endlich bleibt und f\"ur $x=\infty$ unendlich gross h\"ochstens von der Ordnung $-(\mu+\nu+\cdots+\varrho)$ wird, folglich eine ganze Function h\"ochstens vom Grade $-(\mu+\nu+\cdots+\varrho)$. Diese Gr\"osse muss daher, wenn die Function nicht identisch verschwindet, eine ganze, nicht negative Zahl sein.

Die partiellen Determinanten, welchen die Gr\"ossen $a_0$, $a_1$, \ldots, $a_n$ proportional sind, verhalten sich demnach wie ganze Functionen, multiplicirt mit Potenzen von $x-a$, $x-b$, \ldots, $x-g$, deren Exponenten in den verschiedenen Determinanten sich um ganze Zahlen unterscheiden. Die Gr\"ossen $a_0$, $a_1$, \ldots, $a_n$ verhalten sich daher selbst wie ganze Functionen und k\"onnen in den Gleichungen~(2) durch diese ersetzt werden, wodurch man den zu beweisenden Satz erh\"alt.

Die Derivirten der Functionen $y_1$, $y_2$, \ldots, $y_n$ nach $x$ bilden offenbar ein derselben Klasse angeh\"origes System, denn die Differentialquotienten der Functionen $(A)(y_1, y_2, \ldots, y_n)$, in welche $(y_1, y_2, \ldots, y_n)$ durch einen positiven Umlauf des $x$ um $a$ \"ubergehen, sind
\[
\frac{d}{dx}\sum A_i^{(1)}y_i = \sum A_i^{(1)}\frac{dy_i}{dx}, \quad \ldots,
\]
da die Coefficienten in $(A)$ constant sind. Durch diese Bemerkung erh\"alt man aus dem eben bewiesenen Satz die beiden \emph{Corollare}:

\begin{quote}
``Die Functionen $y$ eines Systems gen\"ugen einer Differentialgleichung $n$ter Ordnung, deren Coefficienten ganze Functionen von $x$ sind.''
\end{quote}
und:
\begin{quote}
``Jedes derselben Klasse angeh\"orige System l\"asst sich in diese Functionen und ihre $n-1$ ersten Differentialquotienten linear mit rationalen Coefficienten ausdr\"ucken.''
\end{quote}

Mit H\"ulfe des letzteren l\"asst sich ein allgemeiner Ausdruck f\"ur s\"ammtliche Systeme einer Klasse bilden, aus welchem man sofort sehen w\"urde, dass die Anzahl s\"ammtlicher Systeme, wie oben behauptet, unendlich ist; es soll indess hier nur angewandt werden zur Aufsuchung aller Systeme, in welchen nicht bloss die Substitutionen, sondern auch die Exponenten dieselben sind. F\"ur ein beliebiges System $Y_1$, $Y_2$, \ldots, $Y_n$ derselben Klasse hat man nach diesem Corollar
\[
Y_\nu = \sum_{\mu=0}^{n-1} f_{\nu,\mu}(x) \frac{d^\mu y_1}{dx^\mu},
\]
und wenn man diese Werthe in die Ausdr\"ucke der Substitutionen einf\"uhrt, so erh\"alt man f\"ur die Functionen $f$ Bedingungsgleichungen, welche zeigen, dass sie rational von $x$ abh\"angen und sich durch eine endliche Anzahl von willk\"urlichen Constanten ausdr\"ucken lassen.

\newpage

\clearpage
\chapter*{Source segment 9: riemann pages 401 450}
\addcontentsline{toc}{chapter}{Source segment 9: riemann pages 401 450}
\documentclass[12pt,a4paper]{article}
\usepackage[utf8]{inputenc}
\usepackage[T1]{fontenc}
\usepackage[german,latin,italian]{babel}
\usepackage{amsmath,amssymb,amsthm}
\usepackage{geometry}
\usepackage{fancyhdr}
\usepackage{titlesec}
\usepackage{enumitem}

\geometry{margin=2.5cm}

% Macro definitions
\newcommand{\cO}{\mathcal{O}}
\newcommand{\cF}{\mathcal{F}}
\newcommand{\cE}{\mathcal{E}}
\newcommand{\cL}{\mathcal{L}}
\newcommand{\cH}{\mathcal{H}}
\newcommand{\cA}{\mathcal{A}}
\newcommand{\cB}{\mathcal{B}}
\newcommand{\cC}{\mathcal{C}}
\newcommand{\cD}{\mathcal{D}}
\newcommand{\cG}{\mathcal{G}}
\newcommand{\cI}{\mathcal{I}}
\newcommand{\cJ}{\mathcal{J}}
\newcommand{\cK}{\mathcal{K}}
\newcommand{\cM}{\mathcal{M}}
\newcommand{\cN}{\mathcal{N}}
\newcommand{\cP}{\mathcal{P}}
\newcommand{\cS}{\mathcal{S}}
\newcommand{\cT}{\mathcal{T}}
\newcommand{\cU}{\mathcal{U}}
\newcommand{\cV}{\mathcal{V}}
\newcommand{\cW}{\mathcal{W}}
\newcommand{\bA}{\mathbb{A}}
\newcommand{\bC}{\mathbb{C}}
\newcommand{\bF}{\mathbb{F}}
\newcommand{\bG}{\mathbb{G}}
\newcommand{\bN}{\mathbb{N}}
\newcommand{\bP}{\mathbb{P}}
\newcommand{\bQ}{\mathbb{Q}}
\newcommand{\bR}{\mathbb{R}}
\newcommand{\bZ}{\mathbb{Z}}
\newcommand{\fm}{\mathfrak{m}}
\newcommand{\fp}{\mathfrak{p}}
\newcommand{\fq}{\mathfrak{q}}
\newcommand{\fP}{\mathfrak{P}}
\newcommand{\fa}{\mathfrak{a}}
\DeclareMathOperator{\Gal}{Gal}
\DeclareMathOperator{\Aut}{Aut}
\DeclareMathOperator{\End}{End}
\DeclareMathOperator{\Hom}{Hom}
\DeclareMathOperator{\Ext}{Ext}
\DeclareMathOperator{\Tor}{Tor}
\DeclareMathOperator{\Pic}{Pic}
\DeclareMathOperator{\Ker}{Ker}
\let\Im\relax
\DeclareMathOperator{\Im}{Im}
\DeclareMathOperator{\codim}{codim}
\DeclareMathOperator{\supp}{supp}
\DeclareMathOperator{\Spec}{Spec}
\DeclareMathOperator{\Spf}{Spf}
\DeclareMathOperator{\Proj}{Proj}
\DeclareMathOperator{\Div}{Div}
\DeclareMathOperator{\ord}{ord}
\DeclareMathOperator{\card}{card}
\DeclareMathOperator{\rg}{rg}
\DeclareMathOperator{\tr}{tr}
\DeclareMathOperator{\Tr}{Tr}
\DeclareMathOperator{\Res}{Res}
\DeclareMathOperator{\Ind}{Ind}
\DeclareMathOperator{\Cor}{Cor}
\DeclareMathOperator{\Br}{Br}
\DeclareMathOperator{\Alb}{Alb}
\DeclareMathOperator{\Jac}{Jac}
\DeclareMathOperator{\Prym}{Prym}
\DeclareMathOperator{\Sym}{Sym}
\DeclareMathOperator{\Alt}{Alt}
\DeclareMathOperator{\Num}{Num}
\DeclareMathOperator{\GL}{GL}

\pagestyle{fancy}
\fancyhf{}
\fancyhead[C]{\small Riemann --- Gesammelte Werke (pp.~401--450)}
\fancyfoot[C]{\thepage}
\renewcommand{\headrulewidth}{0.4pt}

\titleformat{\section}{\Large\bfseries}{\thesection.}{0.5em}{}
\titleformat{\subsection}{\large\bfseries}{\thesubsection.}{0.5em}{}

\begin{document}

\begin{center}
{\LARGE\bfseries Riemann --- Gesammelte Mathematische Werke}\\[0.5em]
{\Large (Seiten 401--450)}
\end{center}
\vspace{1em}

%====================================================================
\section*{XXII. Commentatio mathematica, qua respondere tentatur
quaestioni ab Ill\textsuperscript{ma} Academia Parisiensi propositae}
%====================================================================
\selectlanguage{latin}

\ldots manifesto obtinebuntur casus omnes, in quibus $u$ functio temporis et
 duarum tantum variabilium fieri possit. Unde quaestio prior soluta erit.

Superest ut quaeramus, quando expressio $\Sigma \beta_{\iota,\iota'}\, dS_{\iota}\, dS_{\iota'}$
in formam datam $\Sigma \alpha_{\iota,\iota'}\, dx_{\iota}\, dx_{\iota'}$ transformari possit.

\subsection*{Pars secunda}

De transformatione expressionis $\Sigma b_{\iota,\iota'}\, dS_{\iota} dS_{\iota'}$ in formam
 datam $\Sigma a_{\iota,\iota'}\, dX_{\iota}\, dX_{\iota'}$.

Quum quaestio ab Ill\textsuperscript{ma} Academia ad corpora homogenea restricta
 sit, in quibus conductibilitates resultantes sint constantes, evolvamus
primum conditiones, ut expressio $\Sigma b_{\iota,\iota'}\, ds_{\iota}\, ds_{\iota'}$, aequando
 quantitates $s$ functionibus ipsarum $x$, in formam
 $\Sigma a_{\iota,\iota'}\, dx_{\iota}\, dx_{\iota'}$, constantibus coefficientibus
 $a_{\iota,\iota'}$ affectam transformari possit. Deinde de transformatione
in formam quamlibet datam pauca adjiciemus.

Expressionem $\Sigma a_{\iota,\iota'}\, dx_{\iota}\, dx_{\iota'}$, si est, id quod supponimus, forma
positiva ipsarum $dx$, semper in formam $\Sigma dx_{\iota}^2$ redigi posse constat.
Unde si $\Sigma b_{\iota,\iota'}\, dS_{\iota}\, dS_{\iota'}$ in formam
$\Sigma a_{\iota,\iota'}\, dx_{\iota}\, dx_{\iota'}$ transformari potest, redigi etiam
 potest in formam $\Sigma dx_{\iota}^2$ et vice versa. Quaeramus igitur,
quando in formam $\Sigma dx_{\iota}^2$ transformari possit.

Sit determinans $\Sigma \pm b_{\iota_1,\iota_1'}\, b_{\iota_2,\iota_2'} \cdots b_{\iota_n,\iota_n'} = B$
et determinantes partiales $= \beta_{\iota,\iota'}$, quo pacto erit
$\Sigma \beta_{\iota,\iota'}\, b_{\iota,\iota'} = B$ et
$\Sigma \beta_{\iota,\iota''}\, b_{\iota,\iota'} = 0$, si $\iota' \neq \iota''$.

Si $\Sigma b_{\iota,\iota'}\, dS_{\iota}\, dS_{\iota'} = \Sigma dX_{\iota}^2$ pro valoribus
 quibuslibet ipsarum $dx$, substituendo $d + \delta$ pro $d$ invenitur etiam
$\Sigma b_{\iota,\iota'}\, dS_{\iota}\, \delta S_{\iota'} = \Sigma dx_{\iota}\, \delta x_{\iota}$
 pro valoribus quibuslibet ipsarum $d$ et $\delta$.

Hinc si quantitates $dS_{\iota}$ per $dX_{\iota}$ et quantitates $\delta x_{\iota}$ per
 quantitates $dS_{\iota}$ exprimuntur, sequitur
\[
\sum_{\iota,\iota'} b_{\iota,\iota'}\, dS_{\iota}\, dS_{\iota'} =
\sum_{\iota} \left( \sum_{\iota'} \frac{\partial x_{\iota}}{\partial S_{\iota'}}\, dS_{\iota'} \right)^2,
\]
et proinde
\[
\sum_{\iota''} b_{\iota'',\iota'}\, dS_{\iota''} =
\sum_{\iota} \frac{\partial x_{\iota}}{\partial S_{\iota'}}\, dx_{\iota},
\qquad
\sum_{\iota} \frac{\partial x_{\iota}}{\partial S_{\iota'}}\,
\frac{\partial x_{\iota}}{\partial S_{\iota''}} = b_{\iota',\iota''}.
\]

Unde porro deducitur, quoniam sit
\[
\sum_{\iota} \frac{\partial x_{\iota}}{\partial S_{\iota'}}\,
 dx_{\iota} = \sum_{\iota,\iota''} \frac{\partial x_{\iota}}{\partial S_{\iota'}}\,
\frac{\partial x_{\iota}}{\partial S_{\iota''}}\, dS_{\iota''}
= \sum_{\iota''} b_{\iota'',\iota'}\, dS_{\iota''},
\]
et differentiando formulam (3)
\[
\sum_{\iota} dx_{\iota}\, d\frac{\partial x_{\iota}}{\partial S_{\iota'}} =
\sum_{\iota''} db_{\iota'',\i'}\, dS_{\iota''}.
\]

Iam ex his ipsarum $b_{\iota',\iota''}$ expressionibus eruitur
\begin{equation}\label{eq:5}
2\sum_{\iota} \frac{d^2 x_{\iota}}{dS_{\iota'}\, dS_{\iota''}}\, dx_{\iota} =
 db_{\iota',\iota''} + \frac{\partial b_{\iota'',\iota'}}{\partial S_{\iota'''}}\, dS_{\iota'''}
 + \frac{\partial b_{\iota',\iota'''}}{\partial S_{\iota''}}\, dS_{\iota'''}
\end{equation}
et si haec quantitas per $P_{\iota',\iota'',\iota'''}$ designatur
\begin{equation}\label{eq:6}
2P_{\iota',\iota'',\iota'''} =
 \frac{\partial b_{\iota',\iota''}}{\partial S_{\iota'''}}
 + \frac{\partial b_{\iota'',\iota'''}}{\partial S_{\iota'}}
 + \frac{\partial b_{\iota''',\iota'}}{\partial S_{\iota''}}.
\end{equation}
Quantitatibus $b_{\iota,\iota'}$, $\iota''$ iterum differentiatis obtinetur
\[
2\sum_{\iota} \frac{d^2 x_{\iota}}{dS_{\iota'}\, dS_{\iota''}}\,
 d\frac{\partial x_{\iota}}{\partial S_{\iota'''}} = dP_{\iota',\iota'',\iota'''}
\]
umde tandem prodit, substitutis valoribus modo inventis (6) et (4)
\begin{equation}\label{eq:I}
\tag{I}
\begin{aligned}
&\frac{\partial^2 b_{\iota',\iota''}}{\partial S_{\iota'''}\,\partial S_{\iota^{IV}}}
+ \frac{\partial^2 b_{\iota''',\iota^{IV}}}{\partial S_{\iota'}\,\partial S_{\iota''}}
- \frac{\partial^2 b_{\iota',\iota'''}}{\partial S_{\iota''}\,\partial S_{\iota^{IV}}}
- \frac{\partial^2 b_{\iota'',\iota^{IV}}}{\partial S_{\iota'}\,\partial S_{\iota'''}} \\
&\quad = \sum_{\kappa,\kappa'} \beta_{\kappa,\kappa'}\,
 P_{\kappa,\iota',\iota'''}\, P_{\kappa',\iota'',\iota^{IV}}.
\end{aligned}
\end{equation}

Hujus modi igitur aequationibus functiones $b$ satisfaciant necesse
est, quando $\Sigma b_{\iota,\iota'}\, dS_{\iota}\, dS_{\iota'}$ in formam
$\Sigma dX_{\iota}^2$ transformari potest: partes laevas harum aequationum
 designabimus per
\[
\bigl(\iota'\ \iota''\ \iota'''\ \iota^{IV}\bigr).
\]

Ut indoles harum aequationum melius perspiciatur, formetur expressio
\[
\delta\delta \sum b_{\iota,\iota'}\, ds_{\iota}\, ds_{\iota'}
- 2\, d\delta \sum b_{\iota,\iota'}\, ds_{\iota}\, \delta s_{\iota'}
+ dd \sum b_{\iota,\iota'}\, \delta s_{\iota}\, \delta s_{\iota'},
\]
determinatis variationibus secundi ordinis $d^2$, $d\delta$, $\delta^2$ ita, ut sit
\[
\begin{aligned}
d^2 \sum b_{\iota,\iota'}\, ds_{\iota}\, ds_{\iota'}
 - \delta\delta \sum b_{\iota,\iota'}\, ds_{\iota}\, d's_{\iota'}
 - dd \sum b_{\iota,\iota'}\, \delta s_{\iota}\, \delta s_{\iota'} &= 0,\\
d^2 \sum b_{\iota,\iota'}\, ds_{\iota}\, ds_{\iota'}
 - 2\, d\delta \sum b_{\iota,\iota'}\, ds_{\iota}\, \delta's_{\iota'} &= 0,\\
\delta^2 \sum b_{\iota,\iota'}\, \delta s_{\iota}\, \delta s_{\iota'}
 - 2\, \delta\delta \sum b_{\iota,\iota'}\, \delta s_{\iota}\, \delta's_{\iota'} &= 0,
\end{aligned}
\]
(denotante $\delta'$ variationem quamcunque). Quo pacto haec expressio invenitur
\begin{equation}\label{eq:II}
\tag{II}
= \sum_{(\iota'\, \iota''\, \iota''')}
 (\iota'\, \iota''\, \iota''')\,
 (ds_{\iota}\, ds_{\iota'''} - ds_{\iota'''}\, \delta s_{\iota})
 (ds_{\iota''}\, \delta s_{\iota^{IV}} - ds_{\iota^{IV}}\, \delta s_{\iota''}).
\end{equation}

Iam ex hac formatione hujus expressionis sponte patet, mutatis
variabilibus independentibus transmutari eam in expressionem a nova
forma ipsius $\Sigma b_{\iota,\iota'}\, ds_{\iota}\, ds_{\iota'}$ eadem lege dependentem.
At si quantitates $b$ sunt constantes, omnes coefficientes expressionis
(II) cifrae aequales evadunt. Unde si
$\Sigma b_{\iota,\iota'}\, dS_{\iota}\, dS_{\iota'}$ in expressionem similem
 constantibus coefficientibus affectam transformari potest, expressio (II)
 identice evanescat necesse est.

Perinde patet, si expressio (II) non evanescat, expressionem
\begin{equation}\label{eq:III}
\tag{III}
\frac{\sum_{(\iota'\, \iota''\, \iunit''')}
 (\iota'\, \iota''\, \iunit''')\,
 (ds_{\iota}\, ds_{\iunit'''} - ds_{\iunit'''}\, \delta s_{\iota})
 (ds_{\iunit''}\, \delta s_{\iunit^{IV}} - ds_{\iunit^{IV}}\, \delta s_{\iunit''})}
 {\sum b_{\iota,\iota'}\, ds_{\iota}\, ds_{\iunit'} \cdot
  \sum b_{\iota,\iota'}\, \delta s_{\iota}\, \delta s_{\iunit'}
  - \bigl(\sum b_{\iota,\iota'}\, ds_{\iota}\, \delta s_{\iunit'}\bigr)^2}
\end{equation}
mutatis variabilibus independentibus non mutari, insuperque immutatam
 manere, si in locos variationum $ds_{\iota}$, $\delta s_{\iota}$ expressiones
ipsarum lineares quaelibet $\alpha\, dS_{\iota} + \beta\, \delta S_{\iota}$,
$\gamma\, dS_{\iota} + \delta\, \delta S_{\iota}$ substituantur. Valores
 autem maximi et minimi hujus functionis (III) ipsarum $dS_{\iota}$,
$\delta S_{\iota}$ neque a forma expressionis
$\Sigma b_{\iota,\iota'}\, dS_{\iota}\, dS_{\iota'}$ neque a valoribus variationum
$ds_{\iota}$, $\delta s_{\iota}$ pendebunt, unde ex his valoribus dignosci
 poterit, an duae hujusmodi expressiones in se transformari possint.

Disquisitiones haece interpretatione quadam geometrica illustrari
 possunt, quae quamquam conceptibus inusitatis nitatur, tamen obiter
 eam addigavisse juvabit.

Expressio $\sqrt{\sum b_{\iota,\iota'}\, dS_{\iota}\, dS_{\iota'}}$ spectari potest tanquam
elementum lineare in spatio generatiore $n$ dimensionum nostrum
 intuitum transcendente. Quodsi in hoc spatio a puncto
 $(s_1, s_2, \ldots, s_n)$ ducantur omnes lineae brevissimae, in quarum
elementis initialibus variationes ipsarum $s$ sunt ut
 $\alpha\, ds_1 + \beta\, \delta s_1 :
 \alpha\, ds_2 + \beta\, \delta s_2 : \cdots :
 \alpha\, ds_n + \beta\, \delta s_n$, denotantibus
$\alpha$ et $\beta$ quantitates quaslibet, hae lineae superficiem constituent,
quam in spatium vulgare nostro intuitui subjectum exolvere licet.
Quo pacto expressio (III) erit mensura curvaturae hujus superficiei
in puncto $(s_1, s_2, \ldots, s_n)$\footnote{Riemann'sche \selectlanguage{german}Werke,\selectlanguage{latin} S.~272.}.

Si jam ad casum $n=3$ redimus, expressio (II) est forma secundi
 gradus ipsarum
\[ds_2\, \delta s_3 - ds_3\, \delta s_2,\quad
  ds_3\, \delta s_1 - ds_1\, \delta s_3,\quad
  ds_1\, \delta s_2 - ds_2\, \delta s_1,\]
umde in hoc casu sex obtinemus aequationes, quibus functiones $b$
satisfacere debent, ut
$\Sigma b_{\iota,\iota'}\, dS_{\iota}\, dS_{\iota'}$ in formam constantibus coefficientibus
gaudentem transformari possit. Nec difficile, ope notionum modo
traditarum, est demonstratu, has sex conditiones, ut hoc fieri possit,
sufficere. Observandum tamen est ternas tantum esse a se indepen-
dentes.

Iam ut quaestionem ab Ill\textsuperscript{ma} Academia propositam persolvamus,
in his sex aequationibus formae functionum $b$, methodo supra exposita
inventae, sunt substituendae, quo pacto omnes casus invenientur, in
quibus temperatura $u$ in corporibus homogeneis functio temporis et
duarum tantum variabilium fieri possit.

Sed angustia temporis non permisit hos calculos perscribere.
Contenti igitur esse debemus, postquam methodos quibus usi sumus
exposuimus, solutiones singulas quaestionis propositae enumerasse.

Si brevitatis causa casum simplicissimum, quando temperatura $u$
secundum legem
\[a^2\frac{d^2 u}{dx^2} + b^2\frac{d^2 u}{dy^2} + c^2\frac{d^2 u}{dz^2} =
 k\, \frac{du}{dt}\]
variatur, solum respicimus, ad quem casus reliquos facile reduci posse
constat: casus $m=1$ tum tantum evenire potest, quando $u$ est
 constans aut in lineis rectis parallelis, aut in circulis helicibusve,
ita ut coordinatis rectangularibus $\varrho$, $\varrho\cos\varphi$, $\varrho\sin\varphi$
rite electis, poni possit
$\alpha = r$, $\beta = \theta + \varphi \cdot \text{const.}$

Casus $m=2$ locum inveniet si $u = f(\alpha) + g(\beta)$, casus $m=3$
si $u = \alpha e^{\lambda t} + f(\beta)$, denotante $\lambda$ constantem realem,
casus denique $m=4$, ut jam supra invenimus, si $u$ est aut
$= \alpha e^{\lambda t} + \beta e^{-\lambda t} + \text{const.}$,
aut $= (\alpha + \beta t)e^{\lambda t} + \text{const.}$,
aut $= f(\alpha)$.

Iam ut formae functionis $u$ penitus innotescant, annotari tantum
opus est, temperaturam $u$, nisi sit formae $\alpha e^{\lambda t}$, tum tantum
 functionem temporis et unius variabilis esse posse, quando sit constans
aut in planis parallelis, aut in cylindris eadem axi gaudentibus, aut
in sphaeris concentricis. Si $u$ est formae $\alpha e^{\lambda t}$, ex aequatione
differentiali (I) sequitur
\[
\frac{d^2 \alpha}{da^2} = \lambda\alpha, \qquad
\frac{d^2 \beta}{db^2} = -\lambda\beta,
\]
et perinde in casu quarto substituendo valores ipsius $u$ in aequatione
differentiali (I), functiones $\alpha$ et $\beta$ facile determinantur, dummodo
animadvertas, in hoc casu $\alpha e^{\lambda t}$ et $\beta e^{-\lambda t}$ esse posse
 quantitates complexas conjugatas.\footnote{Cf.\ Seite~370.}

\newpage
%====================================================================
\selectlanguage{german}
\section*{Anmerkungen}
%====================================================================

\textbf{(1) (Zu Seite 403.)} Diese Untersuchungen enthalten die analytischen Ausf\

\end{document}
